% Ensamblaje causal de la Parte III del sucesor de 102 capítulos.
% Los propietarios rectores se consumen completos.  Los títulos de capítulo
% se sustituyen únicamente por la jerarquía aprobada; las ampliaciones se
% añaden dentro de la residencia científica correspondiente.

% Punto de montaje redefinible por el compilador único.
% La ruta por defecto es válida cuando el módulo se consume desde
% ``manuscrito/``; el editor único puede redefinirla antes del \input.
\providecommand{\HMTParteIIIRoot}{../colaboracion/parte_iii}
\providecommand{\HMTParteIIISource}{\HMTParteIIIRoot/source/public_final}
\providecommand{\HMTParteIIICapitulos}{\HMTParteIIIRoot/tex/capitulos_finales}
\providecommand{\APP}{\mathrm{APP}}
\providecommand{\TRIT}{\mathrm{TRIT}}
\providecommand{\TPK}{\mathrm{TPK}}
\providecommand{\etaRet}{\eta_{\mathrm{ret}}^{(\mathrm{deg})}}
\providecommand{\etaRetrad}{\eta_{\mathrm{ret}}^{(\mathrm{rad})}}
\providecommand{\alphaAngle}{\alpha_{\angle}}


L'ordre des chapitres suivants est expositif, et non causal. La connexion nonadique et le prolongement complet

\[
w_6\longrightarrow w_{12}\longrightarrow w_{18}\longrightarrow w_{24}
\longrightarrow w_{30}\longrightarrow R_{36}\longrightarrow G_9
\]

stabilisent un état coinductif unique d'APP--TRIT--TPK. À partir de cet état, \(\pi\), \(e\), \(\varphi\) et \(\alpha\) sont publiés de manière corrélée : les chapitres successifs séparent leurs lecteurs afin de les étudier avec précision, sans transformer cette séparation typographique en génération sérielle. La clôture de \(\pi\), la propagation de \(e\), l'auto-échelle de \(\varphi\) et la publication signée et dodécaphasique de \(\alpha\) conservent leurs statuts propres au sein d'une même génération coinductive corrélée.

% Literal expansion of the unchanged chapter-input wrapper.

\HMTFlushCurrentChapter
\chapter{Algèbre et classification généalogique des constantes fondamentales}%
\begingroup
\renewcommand{\chapter}[2][]{}%
\chapter{Algèbre et classification généalogique des constantes fondamentales}
\label{chap:p3-final:27:algebra_clasificacion}
\label{chap:p3-83-23-algebra_clasificacion}


\label{int:chap:producto-nativo}
\index{produit natif}

L'extension euclidienne résidu--quotient d'APP contient deux applications
locales : l'application additive normalise les collisions et la multiplicative
produit des produits partiels. Dans cette section, on construit la multiplication
globale des mots en utilisant exclusivement ces deux applications. La correspondance avec la multiplication entière sera une
conséquence de l'entrelaceur, et non la définition de l'opération.

\section{Coordonnées euclidiennes de la cellule locale}

\begin{definition}[Applications additive et multiplicative sur un même support]
Soit $\Dnine=\{1,\ldots,9\}$. Pour $a,b\in\Dnine$, on définit par division
euclidienne avec reste dans $\Dnine$ :
\begin{align}
 a+b&=9q^+(a,b)+R^+(a,b),\label{int:eq:local-plus}\\
 ab&=9q^\times(a,b)+R^\times(a,b).\label{int:eq:local-times}
\end{align}
En particulier, le reste nul est représenté par $9$ et le quotient est réduit
d'une unité. Le quadruplet
\[
 \CellRQ(a,b)=
 (R^+(a,b),q^+(a,b);R^\times(a,b),q^\times(a,b))
\]
est l'extension arithmétique locale par résidu et quotient.
\end{definition}

L'agrégation selon les trois classes modulo trois du tableau visible
\(R^\times\) définit une matrice distincte des applications locales :
\[
(M_\Pi)_{rs}
=\frac19
\sum_{\substack{a\equiv r\;(3)\\b\equiv s\;(3)}}R^\times(a,b),
\qquad r,s\in\{1,2,0\},
\]
et, dans cet ordre des classes,
\[
M_\Pi=
\begin{pmatrix}4&5&6\\5&4&6\\6&6&9\end{pmatrix}.
\]
Ainsi, \(q^+\) et \(q^\times\) sont des coordonnées de quotient d'une cellule,
tandis que \(M_\Pi\) est un opérateur agrégé sur les classes résiduelles.

Les $81$ cellules satisfont exactement
\[
\begin{array}{c|rrrr}
 &\sum R&\sum q&\sum(9q+R)\\\hline
 +&405&45&810\\
 \times&459&174&2025
\end{array}
\]
et donc
\[
 2025-810=(459-405)+9(174-45)=54+9\cdot129=1215.
\]
Le quotient $q^\times$ est nécessaire à la reconstruction. Par exemple,
\[
 2\cdot5=9\cdot1+1.
\]
Si l'on ne conserve que $R^\times$, le produit $10$ se confond avec la valeur
visible $1$.

\subsection{Cocycles d'associativité et de distributivité}

La cohérence du transducteur peut déjà être décidée dans la cellule. Pour tout
$a,b,c\in\Dnine$, les trois identités sont vérifiées
\begin{align}
q^+(a,b)+q^+(R^+(a,b),c)
&=q^+(b,c)+q^+(a,R^+(b,c)),\label{int:eq:cocycle-plus}\\
q^\times(R^\times(a,b),c)+cq^\times(a,b)
&=q^\times(a,R^\times(b,c))+aq^\times(b,c),
\label{int:eq:cocycle-times}\\
q^\times(a,R^+(b,c))+aq^+(b,c)
&=q^+(R^\times(a,b),R^\times(a,c))\notag\\
&\quad+q^\times(a,b)+q^\times(a,c).
\label{int:eq:cocycle-distributive}
\end{align}

\begin{theorem}[Cohérence locale de l'extension euclidienne]
Les identités \eqref{int:eq:cocycle-plus}--\eqref{int:eq:cocycle-distributive},
avec les égalités correspondantes de résidus, expriment
respectivement l'associativité additive, l'associativité multiplicative et la
distributivité dans l'extension résidu--quotient. Les trois familles présentent
zéro écart sur les $9^3=729$ triplets.
\end{theorem}

\begin{proof}
En développant $(a+b)+c$ au moyen de \eqref{int:eq:local-plus}, le coefficient de
neuf est le membre de gauche de \eqref{int:eq:cocycle-plus} ; en développant
$a+(b+c)$, on obtient celui de droite. L'unicité de la division avec reste dans
$\Dnine$ égalise aussi bien les résidus que les quotients. Le même argument appliqué à
$(ab)c=a(bc)$ donne \eqref{int:eq:cocycle-times}. Enfin, les deux développements de
$a(b+c)=ab+ac$ produisent \eqref{int:eq:cocycle-distributive}. L'énumération
directe des \(729\) cas de chaque identité confirme le calcul
algébrique.
\end{proof}

\section{Bijectivité de la numération}

Le symbole $9$ n'est pas identifié à zéro. Dans $\CellRQ$, $9$ représente un
franchissement complet de la frontière nonadique. Le remplacement par les chiffres
$0,\ldots,8$ éliminerait le quotient euclidien qui intervient dans
l'entrelacement.

\begin{definition}[Mots canoniques]
Soit
\[
 \Wnine^+=\bigsqcup_{\ell\geq1}\Dnine^\ell,
 \qquad \Wnine=\{\varnothing\}\sqcup\Wnine^+.
\]
Les mots s'écrivent à partir du chiffre le moins significatif. Nous définissons
\[
 \operatorname{ev}_9(u_0,\ldots,u_{\ell-1})
 =\sum_{j=0}^{\ell-1}u_j9^j,
 \qquad \operatorname{ev}_9(\varnothing)=0.
\]
\end{definition}

\begin{theorem}[Unicité de la forme normale]
L'évaluation $\operatorname{ev}_9$ est une bijection de $\Wnine^+$ sur
$\N_{\geq1}$, et de $\Wnine$ sur $\N_0$.
\end{theorem}

\begin{proof}
Les mots de longueur $\ell$ couvrent l'intervalle
\[
 \frac{9^\ell-1}{8}
 \leq \operatorname{ev}_9(u)
 \leq 9\frac{9^\ell-1}{8}.
\]
Le minimum de l'intervalle suivant est supérieur d'une unité au maximum du
précédent. À l'intérieur de chaque intervalle, la division de $n-1$ par neuf détermine
de manière unique le chiffre
\[
 d=(n-1)\bmod9+1
\]
et le quotient qui est codé récursivement. Le processus se termine parce que le
quotient décroît strictement.
\end{proof}

Exemples :
\[
 9\longleftrightarrow(9),\qquad
 10\longleftrightarrow(1,1),\qquad
 81\longleftrightarrow(9,8),\qquad
 1000\longleftrightarrow(1,3,3,1).
\]

\section{Fibres des projections résiduelles dans les relèvements entiers}

Pour ce bloc, nous distinguons deux alphabets. La forme normale
\(\Wnine\) utilise les chiffres \(\{1,\ldots,9\}\), écrits à partir du moins
significatif. En revanche, \(\mathscr B_9\) désigne les écritures positionnelles
conventionnelles en base neuf, avec les chiffres \(\{0,\ldots,8\}\) et le plus
significatif à gauche. Les deux types ne sont pas identifiés.

Soit
\[
 \mathscr W_{10}^{(2)}=\{0,\ldots,9\}^2,\qquad
 \rho_{10}(a,b)=(b,a),\qquad
 \operatorname{ev}_{10}(a,b)=10a+b,
\]
et soit \(q_9:\mathbb N\to\mathbb Z/9\mathbb Z\) la réduction résiduelle.

\begin{lemma}[Descente résiduelle du renversement décimal]
\label{int:pf84:SYN30-RHO-DESC}
Le renversement sur les mots décimaux de longueur deux est involutif et son
ombre modulo neuf est l'identité :
\[
 \rho_{10}^2=I,\qquad
 q_9\!\left(\operatorname{ev}_{10}(\rho_{10}(a,b))\right)
 =
 q_9\!\left(\operatorname{ev}_{10}(a,b)\right).
\]
La longueur fixe conserve les zéros initiaux.
\end{lemma}

\begin{proof}
La double permutation retrouve \((a,b)\). Comme \(10\equiv1\pmod9\),
\[
 10b+a\equiv a+b\equiv10a+b\pmod9.
\]
Un mot dont le double renversement ne serait pas lui-même ou dont la somme des chiffres
changerait modulo neuf réfuterait le lemme.
\end{proof}

\begin{lemma}[Descente résiduelle du carré]
\label{int:pf84:SYN30-SQ-DESC}
L'application \(s(n)=n^2\) induit l'endoapplication résiduelle
\[
 [n]_9\longmapsto[n]_9^2,
 \qquad q_9(s(n))=q_9(n)^2.
\]
\end{lemma}

\begin{proof}
La compatibilité d'une congruence avec le produit donne
\(n\equiv m\pmod9\Rightarrow n^2\equiv m^2\pmod9\). Un entier pour lequel
\(q_9(n^2)\ne q_9(n)^2\) réfuterait l'énoncé.
\end{proof}

\begin{proposition}[Fibres résiduelles du renversement entier]
\label{int:pf84:SYN30-RHO-NONFAITH}
L'ombre résiduelle ne reconstruit pas le relèvement entier
\(\widehat\rho_{10}=\operatorname{ev}_{10}\circ\rho_{10}\). En effet,
\[
 18\equiv27\pmod9,\qquad
 \widehat\rho_{10}(1,8)=81\ne72=\widehat\rho_{10}(2,7).
\]
\end{proposition}

\begin{proof}
Les deux entiers appartiennent à la classe résiduelle zéro, mais leurs renversements sont
distincts. Il n'existe donc pas de
\(\overline\rho:\mathbb Z/9\mathbb Z\to\mathbb N\) qui factorise
simultanément ces deux valeurs.
\end{proof}

\begin{proposition}[Fibres résiduelles du carré entier]
\label{int:pf84:SYN30-SQ-NONFAITH}
La même ombre ne reconstruit pas non plus le carré entier :
\[
 18\equiv27\pmod9,\qquad
 18^2=324\ne729=27^2.
\]
\end{proposition}

\begin{proof}
Une fonction \(\overline s:\mathbb Z/9\mathbb Z\to\mathbb N\) qui
factoriserait \(s\) devrait attribuer à la même classe zéro les deux valeurs
\(324\) et \(729\), ce qui est impossible.
\end{proof}

\begin{corollary}[Témoin croisé \(27\)--\(72\)--\(729\)]
\label{int:pf84:SYN30-CROSS-WITNESS}
\[
 \widehat\rho_{10}(2,7)=72,\qquad
 27^2=729=9^3=3^6,
\]
et, dans l'alphabet conventionnel \(\mathscr B_9\),
\[
 27=30_9,\qquad72=80_9,\qquad729=1000_9.
\]
Les trois entiers ont une ombre résiduelle zéro et une racine numérique positive neuf.
\end{corollary}

\begin{proof}
Les deux opérations se calculent directement ; les divisions successives par
neuf produisent les trois écritures. Le résultat échouerait si l'on confondait
\(\mathscr B_9\) avec \(\Wnine\), si l'une des écritures était incorrecte
ou si l'on affirmait que la descente résiduelle fidèle reconstruit l'opération et sa
généalogie complète.
\end{proof}

\section{Additionneur natif et propagation du quotient}

\begin{algorithm}[Somme de mots]
On parcourt deux mots par positions, avec une retenue initiale \(c_0=0\). Si, à
une position, il y a deux chiffres, on applique \eqref{int:eq:local-plus} ; s'il n'y en a qu'un,
ce chiffre est pris comme résidu provisoire et le quotient provisoire est zéro.
Lorsque \(c_j\ne0\), on applique une seconde cellule additive au résidu provisoire
et à \(c_j\). On publie le nouveau résidu et l'on transporte la somme des
quotients à la position suivante. S'il reste à la fin une retenue non nulle, elle est
publiée comme dernier chiffre. De cette manière, les cellules locales ne reçoivent jamais le
symbole absent \(0\), qui n'appartient pas à \(\Dnine\).
\end{algorithm}

\begin{lemma}[Invariant de l'additionneur]
Si $u\boxplus_\#v$ est le mot produit par l'algorithme précédent,
\[
 \operatorname{ev}_9(u\boxplus_\#v)
 =\operatorname{ev}_9(u)+\operatorname{ev}_9(v).
\]
La retenue intermédiaire est dans $\{0,1,2\}$.
\end{lemma}

\begin{proof}
Après traitement des positions $0,\ldots,k-1$, soit $c_k$ le quotient en attente.
Les identités locales impliquent
\[
 \sum_{j<k}(u_j+v_j)9^j
 =\sum_{j<k}w_j9^j+c_k9^k,
\]
en incluant comme zéro les chiffres absents. Chaque nouvelle cellule conserve cette
égalité. Le maximum local est $9+9+2=20$, de sorte que $c_{k+1}\leq2$. À la
fin, le dernier quotient est publié, et l'identité globale demeure.
\end{proof}

\section{Multiplication par un chiffre}

\begin{algorithm}[Transducteur scalaire]
$d\in\Dnine$ étant fixé, chaque chiffre $u_j$ produit
\[
 u_jd=9q_j+r_j,
 \qquad
 (r_j,q_j)=\bigl(R^\times(u_j,d),q^\times(u_j,d)\bigr).
\]
Soit $c_j\in\{0,\ldots,9\}$ le quotient entrant. Comme $0\notin\Dnine$,
la transition est définie par branches :
\[
 (w_j,c_{j+1})=
 \begin{cases}
 (r_j,q_j),&c_j=0,\\[1mm]
 \bigl(R^+(r_j,c_j),q_j+q^+(r_j,c_j)\bigr),&c_j\in\Dnine.
 \end{cases}
\]
Ainsi, aucune application locale n'est évaluée hors de son domaine. La première
branche exprime l'absence de quotient en attente ; la seconde normalise la
collision au moyen de l'application additive.
\end{algorithm}

\begin{lemma}[Invariant scalaire]
Le mot $d\odot_\#u$ produit par le transducteur satisfait
\[
 \operatorname{ev}_9(d\odot_\#u)
 =d\operatorname{ev}_9(u).
\]
Le quotient multiplicatif appartient à $\{0,\ldots,9\}$.
\end{lemma}

\begin{proof}
L'argument est celui de l'additionneur, en remplaçant chaque terme local par
$u_jd$. Si $c_j=0$, l'identité
$u_jd=9q_j+r_j$ conserve directement l'égalité partielle. Si
$c_j\in\Dnine$, la seconde division euclidienne donne
\[
 r_j+c_j=9q^+(r_j,c_j)+R^+(r_j,c_j),
\]
et conserve la même égalité après transport de
$q_j+q^+(r_j,c_j)$. Le maximum $9\cdot9+9=90$ et la forme du reste dans
$\Dnine$ maintiennent le quotient dans la borne déclarée. Le dernier quotient est
publié comme chiffre.
\end{proof}

\section{Produit global par Horner}

\begin{definition}[Produit natif]
Soit $v=(v_0,\ldots,v_{m-1})$. En partant de $a_m=\varnothing$, on définit pour
$j=m-1,\ldots,0$
\[
 a_j=(9\odot_\#a_{j+1})\boxplus_\#(v_j\odot_\#u).
\]
Le résultat est
\[
 u\boxtimes_\#v:=a_0.
\]
\end{definition}

Cette définition ne contient pas
\[
 \operatorname{dec}_9(\operatorname{ev}_9(u)\operatorname{ev}_9(v)).
\]
L'évaluation n'apparaîtra que dans la démonstration de l'entrelacement.

\begin{theorem}[Multiplicativité native]
Pour tous les mots $u,v\in\Wnine$,
\begin{equation}
 \label{int:eq:ev-multiplicative}
 \boxed{
 \operatorname{ev}_9(u\boxtimes_\#v)
 =\operatorname{ev}_9(u)\operatorname{ev}_9(v).}
\end{equation}
Le mot vide est absorbant et $(1)$ est l'unité.
\end{theorem}

\begin{proof}
Par les deux lemmes précédents,
\[
 \operatorname{ev}_9(a_j)
 =9\operatorname{ev}_9(a_{j+1})
  +v_j\operatorname{ev}_9(u).
\]
L'itération de Horner donne
\[
 \operatorname{ev}_9(a_0)
 =\operatorname{ev}_9(u)
  \sum_{j=0}^{m-1}v_j9^j.
\]
Le dernier facteur est $\operatorname{ev}_9(v)$. L'unicité de la forme normale
identifie le mot résultant.
\end{proof}

\begin{corollary}[Monoïde multiplicatif de mots]
La structure
\[
\WordMonoid=(\Wnine,\boxtimes_\#,(1))
\]
est un monoïde commutatif ; le mot vide est son élément absorbant. Sa
construction n'utilise pas la multiplication globale des entiers.
\end{corollary}

L'exemple de frontière le plus court est
\[
 (9)\boxtimes_\#(9)=(9,8),
 \qquad 9+8\cdot9=81.
\]
Les \(17\) cellules locales de la multiplication produisent
\[
 (9,8)\boxtimes_\#(9,8)=(9,8,8,8),
 \qquad \operatorname{ev}_9(9,8,8,8)=6561.
\]

\section{Linéarisation matricielle de l'état généalogique et construction de l'entrelaceur complètement positif}

\begin{definition}[Secteurs de Hilbert]
Soient
\[
 \Krad=\ell^2(\Wnine^+),
 \qquad
 \Hplus=\ell^2(\N_{\geq1}).
\]
Nous définissons sur les bases canoniques
\[
 W_\times e_u=e_{\operatorname{ev}_9(u)}.
\]
Dans \(\Hplus\), l'opérateur nombre est noté
\[
 \mathsf N e_n=n e_n,
 \qquad \Spec(\mathsf N)=\N_{\geq1}.
\]
Le vide peut être ajouté séparément comme vecteur de vide $e_0$ ; il ne se
confond pas avec $\Hplus$, dont le spectre commence à un.
\end{definition}

\begin{theorem}[Unitarité et carré multiplicatif]
$W_\times$ est unitaire. Dans le cœur algébrique
$\mathbb C[\Wnine^+\times\Wnine^+]$, soit
\[
 \mu_\#(e_u\otimes e_v)=e_{u\boxtimes_\#v},
 \qquad
 \mu_{\N}(e_a\otimes e_b)=e_{ab}.
\]
Alors on a le carré
\[
\begin{CD}
\Krad\otimes_{\mathrm{alg}}\Krad @>{W_\times\otimes W_\times}>>
\Hplus\otimes_{\mathrm{alg}}\Hplus\\
@V{\mu_\#}VV @VV{\mu_{\N}}V\\
\Krad @>>{W_\times}> \Hplus
\end{CD}
\]
et, par conséquent, le carré commute.
\end{theorem}

\begin{proof}
La bijection des formes normales envoie une base orthonormale sur une autre et prouve
l'unitarité. Sur un tenseur de base, les deux chemins produisent
$e_{\operatorname{ev}_9(u)\operatorname{ev}_9(v)}$ par
\eqref{int:eq:ev-multiplicative}. La linéarité donne l'égalité dans le cœur.
\end{proof}

Les applications $\mu_\#$ et $\mu_{\N}$, définies initialement sur les cœurs
algébriques respectifs, sont densément définies et non bornées. Pour
\(c=(c_{u,v})\in\ell^2(\Wnine^+\times\Wnine^+)\), chaque somme sur une fibre de produit
est finie. L'adjoint de \(\mu_\#\) contient tous les vecteurs à support
fini dans son domaine, de sorte que \(\mu_\#\) est fermable. Le domaine de sa
fermeture est
\[
 \Dom(\overline{\mu_\#})=
 \left\{c\in\ell^2(\Wnine^+\times\Wnine^+):
 \sum_{w\in\Wnine^+}
 \left|\sum_{u\boxtimes_\#v=w}c_{u,v}\right|^2<\infty
 \right\}.
\]
Chaque fibre est finie parce que $\operatorname{ev}_9$ l'identifie aux couples
ordonnés de diviseurs d'un entier. La même description vaut pour
\(\overline{\mu_{\N}}\) après application de \(W_\times\otimes W_\times\), et
$W_\times$ entrelace les graphes des deux fermetures.

\section{Coproduit et sélecteur natifs}

L'ouverture multiplicative est maintenant définie avant de regarder le codomaine :
\[
 \Delta_\#e_w
 =\sum_{u\boxtimes_\#v=w}e_u\otimes e_v.
\]
La somme s'obtient en énumérant les produits du transducteur. On ne demande pas si
un entier ordinaire en divise un autre.
Dans le secteur entier, nous définissons, sur le cœur algébrique correspondant,
\[
 \Delta_\times e_n
 =\sum_{\substack{a,b\in\N_{\geq1}\\ab=n}}e_a\otimes e_b.
\]
Les deux coproduits énumèrent des couples ordonnés.

\begin{theorem}[Coassociativité native]
Dans le cœur algébrique,
\[
 (\Delta_\#\otimes I)\Delta_\#
 =(I\otimes\Delta_\#)\Delta_\#.
\]
De plus,
\[
 (W_\times\otimes W_\times)\Delta_\#
 =\Delta_\times W_\times.
\]
\end{theorem}

\begin{proof}
Les deux membres de la première égalité énumèrent une fois chaque triplet natif
$a\boxtimes_\#b\boxtimes_\#c=w$. La seconde égalité découle de la bijection
entre fibres induite par \eqref{int:eq:ev-multiplicative}.
\end{proof}

\begin{theorem}[Coproduit arithmétique complet]
\label{int:pf84:C06}
Dans le cœur algébrique
\(c_{00}(\mathbb N_{\geq1})\subset\ell^2(\mathbb N_{\geq1})\), soit
\[
 \Delta_\times e_n=\sum_{ab=n}e_a\otimes e_b.
\]
Sa fermeture a pour domaine
\[
 \operatorname{Dom}(\overline{\Delta_\times})
 =
 \left\{x=(x_n):
 \sum_{n\geq1}d(n)|x_n|^2<\infty\right\},
\]
est coassociative dans le cœur et, pour \(Qe_n=ne_n\), satisfait
\[
 (\log Q\otimes I+I\otimes\log Q)\Delta_\times
 =\Delta_\times\log Q.
\]
\end{theorem}

\begin{proof}
Les supports de \(\Delta_\times e_n\) sont orthogonaux pour des entiers distincts
et leur norme au carré est le nombre \(d(n)\) de diviseurs, ce qui donne le domaine
pondéré de la fermeture. Les deux coassociations réindexent la même somme
sur \(abc=n\). Enfin, \(\log a+\log b=\log n\) dans chaque terme
prouve l'entrelacement. Une triple factorisation comptée de manière différente
ou l'échec de cette identité réfuterait le théorème.
\end{proof}

\begin{theorem}[Coproduit réduit et sélecteur d'irréductibles]
\label{int:pf84:C07}
Dans le secteur entier, le coproduit réduit
\[
 \Delta_\times^{\rm red}e_n
 =\sum_{\substack{ab=n\\a,b>1}}e_a\otimes e_b
\]
est fermable. L'opérateur
\[
 N_\times^{\rm red}
 =\bigl(\overline{\Delta_\times^{\rm red}}\bigr)^*
  \overline{\Delta_\times^{\rm red}}
\]
est positif autoadjoint et diagonal :
\[
 N_\times^{\rm red}e_n=d_{\rm red}(n)e_n.
\]
Son noyau est engendré par \(e_1\) et par les \(e_p\) avec \(p\) irréductible ;
par conséquent,
\[
 P_{\Irr}
 =\mathbf1_{\{0\}}(N_\times^{\rm red})
  -\lvert e_1\rangle\langle e_1\rvert
\]
sélectionne les irréductibles et exclut l'unité.

Dans le secteur natif, le coproduit correspondant est
\[
 \widetilde\Delta_\#e_w
 =\sum_{\substack{u\boxtimes_\#v=w\\u,v\ne(1)}}e_u\otimes e_v.
\]
L'entrelaceur \(W_\times\) transporte
\(\Delta_\times^{\rm red}\), \(N_\times^{\rm red}\) et \(P_{\Irr}\) vers
\(\widetilde\Delta_\#\),
\[
 A_\#:=
 \bigl(\overline{\widetilde\Delta_\#}\bigr)^*
 \overline{\widetilde\Delta_\#},
 \qquad
 P_{\Irr,\#}
 =(I-P_{(1)})\mathbf1_{\{0\}}(A_\#),
 \qquad
 P_{(1)}=\lvert e_{(1)}\rangle\langle e_{(1)}\rvert.
\]
\end{theorem}

\begin{proof}
Les supports des images de deux vecteurs de base distincts sont orthogonaux.
La norme au carré de
\(\Delta_\times^{\rm red}e_n\) est \(d_{\rm red}(n)\), qui s'annule
exactement pour \(n=1\) ou \(n\) irréductible. Cela prouve la diagonalisation
et le noyau déclaré ; soustraire le projecteur de \(e_1\) élimine l'unité.
La multiplicativité de \(W_\times\) identifie les fibres entières aux
fibres natives et transporte les trois opérateurs. Un composé dans le noyau,
un irréductible en dehors de celui-ci ou l'inclusion de l'unité dans \(P_{\Irr}\)
réfuteraient le théorème.
\end{proof}

\begin{proposition}[Coïncidence nonadique des nombres premiers et composés]
\label{int:pf84:B05}
Soit \(\gcd(n,90)=1\),
\(\tau_n=\operatorname{ord}_n(10)\) et
\[
 M_n=\frac{10^{\tau_n}-1}{n}.
\]
Alors \(M_n\equiv0\pmod9\), aussi bien pour les nombres premiers que pour les composés. En
particulier, la clôture nonadique et la période décimale prises isolément ne sont pas un
sélecteur d'irréductibles.
\end{proposition}

\begin{proof}
Comme \(9\mid10^{\tau_n}-1\) et \(n\) est inversible modulo neuf, la
divisibilité par neuf survit au quotient. Les nombres premiers \(7,13\) et les
composés \(77,91,143\) partagent la période six ; les nombres de Carmichael
apportent des contrôles supplémentaires. L'énoncé ne serait réfuté que si la
congruence ou la période séparaient tous les nombres premiers de tous les composés
dans le domaine déclaré.
\end{proof}

Cette affirmation est strictement arithmétique : elle identifie des fibres de
factorisation et les opérateurs qu'elles induisent sur le secteur des formes
normales.

\ifdefined\HMTInternalTraceability
\section{Extension monoïdale du produit natif aux histoires compatibles}
\label{int:pf84:SYN-05D-HIST}

Le résultat distingue l'espace \(\Krad\) des formes normales canoniques et
l'espace \(\mathcal K_{\#,\mathrm{hist}}\) des histoires complètes de
cellules et de quotients propagés.
De nombreuses histoires peuvent produire la même forme normale. Par conséquent,
$W_\times$ est unitaire sur $\Krad$, mais la projection d'une histoire sur
sa valeur empêche qu'il soit injectif sur toutes les classes de transport.

Soit \(\mathscr X_n^{\rm hist}\) l'ensemble des histoires enrichies de
profondeur \(n\) et soit
\(\mathcal H_n=\ell^2(\mathscr X_n^{\rm hist})\). Notons \(P_n\) la
projection orthogonale sur le sous-espace engendré par les histoires
survivantes et supposons définie une multiplication \(\mu_n\) sur un
domaine algébrique de
\(\mathcal H_n\otimes_{\mathrm{alg}}\mathcal H_n\). Une extension monoïdale
du produit de mots au système inverse consiste en une famille
d'applications \(\iota_n:\Krad\to\mathcal H_n\) qui commute avec les
restrictions de profondeur et satisfait, pour chaque \(n\),
\[
 P_n\,\iota_n\mu_\#
 =P_n\,\mu_n
 (\iota_n\otimes\iota_n)
\]
dans le cœur algébrique. Cette égalité est la condition de compatibilité qui
définit l'extension. Les théorèmes précédents n'établissent ni l'existence ni
l'unicité d'une famille \((\iota_n)_n\) ; ils démontrent le produit dans le
monoïde \(\WordMonoid\) et déterminent l'équation qu'une extension à l'espace
des histoires doit satisfaire. Le transfert conserve donc le
statut de \emph{programme ouvert} : il n'hérite pas de l'exactitude du produit dans
\(\WordMonoid\). L'inexistence d'une famille compatible, l'échec de la
commutation avec les troncatures, la perte de survie, de quotient,
de retenue, d'orientation ou de mémoire, ou le non-respect de l'identité précédente
sont ses falsificateurs.
\fi

\section{Énumération exhaustive et nécessité des coordonnées structurelles}

L'énumération finie comprend :
\[
\begin{array}{@{}lr@{}}
\toprule
\text{Test}&\text{Résultat}\\\midrule
\text{Cellules par application locale}&81\\
\text{Parcours réversibles de numérotation}&59050\\
\text{Paires de l’entrelaceur}&532900\\
\text{Erreurs de l’entrelaceur}&0\\
\text{Triplets associatifs}&531441\\
\text{Erreurs associatives}&0\\
\text{Fenêtre de coproduit}&1\text{ à }500\\
\text{Irréductibles obtenus}&95\\
\bottomrule
\end{array}
\]

\begin{definition}[Quotient multiplicatif]
Conserver seulement $R^\times$ transforme $2\cdot5=10$ en la valeur visible $1$.
Perturber $q^\times(9,9)$ de $8$ à $7$ transforme $81$ en $72$. Le quotient
local est nécessaire.
\end{definition}

\begin{definition}[Application additive]
L'application multiplicative produit des produits partiels, mais ne résout pas leurs collisions.
Éliminer $q^+$ rompt la normalisation globale ; par exemple, elle cesse de reproduire
les quotients intermédiaires de $10\cdot18=180$.
\end{definition}

\begin{definition}[Coordonnées et symbole de bord]
Le zéro ordinaire n'est pas un chiffre de la construction. Le mot vide est
$0$ et le symbole $9$ conserve sa valeur et son quotient. Mélanger les coordonnées
positionnelles indexées à partir de zéro avec les étiquettes $1,\ldots,9$ modifie les
cellules.
\end{definition}

\begin{proposition}[Portée exacte des ablations et du recensement fini]
\label{int:pf84:D10}
Les ablations précédentes prouvent, dans les conventions déclarées :
\begin{enumerate}
  \item sans \(q^\times\), \(2\cdot5=10\) se réduit au résidu visible \(1\) ;
  \item remplacer \(q^\times(9,9)=8\) par \(7\) transforme \(81\) en \(72\) ;
  \item sans \(q^+\), la normalisation des collisions additives échoue.
\end{enumerate}
L'énumération exhaustive comprend les \(2\cdot81=162\) cellules des deux
applications, les \(729\) triplets locaux, \(532\,900\) paires de mots et
les \(95\) irréductibles inférieurs ou égaux à \(500\). Dans ces domaines finis,
les identités indiquées sont démontrées, mais la survie
coinductive de toute histoire TPK n'est pas prouvée.
\end{proposition}

\begin{proof}
Les trois contre-exemples se calculent directement avec
\eqref{int:eq:local-plus} et \eqref{int:eq:local-times}. L'énumération des quatre
domaines finis produit les cardinaux écrits dans l'énoncé et n'ajoute
aucune identité coinductive à celles qui ont déjà été démontrées dans le monoïde
de mots.
\end{proof}

\section{Structure algébrique du monoïde de mots}

La construction démontre, dans le monoïde \(\WordMonoid\) :
\begin{enumerate}
 \item une forme normale unique pour chaque entier non négatif ;
 \item une somme et un produit globaux construits uniquement avec les deux applications locales ;
 \item la conservation symbolique de la valeur sous la propagation du quotient ;
 \item un unitaire $W_\times$ et le carré multiplicatif exact ;
 \item un coproduit diviseur et un sélecteur irréductible définis nativement ;
 \item la nécessité de conserver la classe de transport lorsque l'on quitte le
 secteur des formes normales.
\end{enumerate}

Le passage à l'espace enrichi se réduit maintenant à un problème d'existence :
construire une famille \((\iota_n)_n\) compatible avec les restrictions de
profondeur et démontrer qu'elle satisfait l'équation de la section précédente.
Cette question est distincte de l'entrelaceur exact sur
\(\WordMonoid\), qui est déjà déterminé par la forme normale et les
applications locales.



\label{p3-83-c23-s1:ch:foundation}

\section{Problème de détermination des constantes à partir de l’état généalogique}

Une constante physique est représentée comme une section d’une droite de
grandeur munie d’une règle d’évaluation, d’une orientation et d’une normalisation.
Sa coordonnée numérique dépend de la carte choisie, tandis que sa
généalogie reste invariante. Par conséquent, le problème de
détermination comprend :

\begin{enumerate}
\item l’identification de l’état HMT antécédent ;
\item la spécification de l’opérateur ou du caractère de sélection ;
\item la détermination de l’invariant conservé ;
\item la droite de grandeur correspondante ;
\item la carte dans laquelle s’exprime la valeur terminale.
\end{enumerate}

La chaîne causale du travail est

\begin{equation}
\begin{aligned}
\APP&\longrightarrow\TRIT\longrightarrow\TPK
\longrightarrow\text{réalisation interne}
\longrightarrow\text{invariant sans dimension}\\
&\longrightarrow\text{section de grandeur}
\longrightarrow\R.
\end{aligned}
\label{p3-83-c23-s1:eq:master-chain}
\end{equation}

Le terme \(\R\) de la chaîne représente l’évaluation terminale de la
section ; il ne constitue pas le domaine originaire à partir duquel la route est
reconstruite rétrospectivement.

\section{Chaîne antécédente APP--TRIT--TPK et domaine des opérateurs d’extraction}

APP fournit la géométrie des positions, des restes, des quotients et des
frontières. TRIT incorpore l’orientation \((-1,0,+1)\) et distingue
l’ouverture, le seuil et le retour. TPK complète l’état par la retenue, la route,
le feuillet, la mémoire et la survie. En particulier, la transition \(9\to1\)
conserve la phase et modifie l’état de mémoire ; elle ne constitue donc pas
une réinitialisation dépourvue de généalogie.

Les réalisations développées ci-dessous conservent cet antécédent sans le remplacer. Le vide s’obtient par une fonction spectrale de l’opérateur bicouche ; la gravité sélectionne une section au moyen de la structure chronologique \(108\) ; la criticité utilise une famille dodécaphasique ; les constantes de Gauss proviennent de l’évaluation de cylindres ; et le Hamiltonien modulaire est défini sur une frontière qui conserve l’incidence \(6\to8\). La production de valeurs scalaires n’établit pas une identité fonctionnelle entre ces réalisations.

\begin{proposition}[Insuffisance de la coïncidence scalaire]
Soient \(F_1:X_1\to L_1\) et \(F_2:X_2\to L_2\) deux fonctionnelles typées.
L’égalité des coordonnées \(c_1(F_1(x_1))=c_2(F_2(x_2))\) n’induit à elle
seule ni une identification \(X_1\simeq X_2\) ni un entrelaceur
\(F_2I=JF_1\).
\end{proposition}

En particulier, une proximité numérique telle que
\(I_{\rm or}\approx\gammaH a_0\) ne peut être promue en identité sans
construire l’entrelaceur correspondant.

\section{Droites de grandeur, cartes orientées et systèmes de coordonnées dimensionnelles}

Soit \(L_Q\) une droite réelle orientée représentant la grandeur \(Q\).
Une unité \(u_Q\in L_Q\setminus\{0\}\) détermine une coordonnée
provisoire \(L_Q\simeq\R\), sans remplacer la structure intrinsèque de
la droite. La relation constitutive

\[
L_\mu=L_S L_C^{-1}L_Q^{-2},\qquad
L_\varepsilon=L_S^{-1}L_C^{-1}L_Q^2
\]

est antérieure au choix SI. De même,

\[
L_Z=L_SL_Q^{-2},\qquad L_G=L_C^5L_T^2L_S^{-1}
\]

typent l’impédance et la constante gravitationnelle avant d’exprimer leurs grandeurs en
ohms ou en unités SI. Dans cet ouvrage, nous écrirons

\[
Q_{\rm int}=\widehat Q\,u_Q
\]

pour distinguer le caractère adimensionnel \(\widehat Q\) de la section
de référence \(u_Q\).

\section{Torseur de repères et covariance des évaluations dimensionnelles}

Les cartes admissibles forment un torseur et, par conséquent, aucune ne possède un caractère
canonique par rapport aux autres. Le changement de repère
\(u_Q\mapsto\lambda_Q u_Q\) induit la transformation contragrédiente
de la coordonnée. Les identités physiques doivent être homogènes par rapport
à cette action. Par exemple,

\[
\mu\varepsilon=c^{-2},\qquad \frac\mu\varepsilon=Z^2
\]

sont valides dans tout repère compatible. L’unification mathématique et
dimensionnelle n’élimine pas les unités, mais les reconstruit par la
covariance des droites et permet de placer l’invariant avant son expression
en coordonnées.

\section{Évaluation métrologique terminale et interdiction de sélection rétrospective}

Chaque entrée de l’atlas final prend la forme

\begin{equation}
\boxed{\text{antécédent}\to\text{opérateur}\to\text{invariant}
\to\text{signification}\to\text{valeur terminale}.}
\label{p3-83-c23-s1:eq:atlas-rule}
\end{equation}

Cet ordre exprime une dépendance causale : la valeur numérique n’est
déterminée qu’après fixation de l’antécédent, de l’opérateur, de l’invariant et de
l’interprétation. L’omission de l’un de ces éléments réduit le
résultat à une reconstruction numérique et non à une dérivation
généalogique.

\section{Domaine exact de la génération généalogique et de sa carte métrologique terminale}

On adopte comme résultats de départ la génération de
\(\pi,e,\varphi\), la clôture dodécaphasique de \(\alpha_{\HMT}\),
la section d’action et la sortie \(\gammaH\) de Barbero--Immirzi. Le
présent développement ne recalibre pas ces constructions ; il détermine les constantes
et les invariants qui résultent de leur application aux réalisations
constitutive, thermique, gravitationnelle, modulaire et spectrale.

La masse de l’électron n’est pas non plus incluse comme problème directeur, puisqu’elle
relève de la Loi générale des masses. Toute apparition de \(m_e\) dans une
identité atomique est traitée comme une donnée provenant de ce domaine et non comme
une nouvelle sélection du présent développement.

\subsection*{Portée logique et domaine de validité}
L’architecture APP--TRIT--TPK et les droites dimensionnelles constituent la
structure de départ. La règle uniforme à cinq champs de l’atlas explicite
sa discipline causale dans chaque entrée.



\section{Taxonomie fonctionnelle de \(114\) expressions numériques et de \(106\) lectures opératorielles}
\label{p3-83-c23-s2:sec:atlas-taxonomico-114-constantes-rev11}
\index{constantes!atlas taxonomique de 114 entrées}

Soit l’ensemble fini indexé
\[
 E_{114}=\{\varepsilon_1,\ldots,\varepsilon_{114}\}
\]
d’expressions numériques et soit \(\mathcal O\) l’espace des lectures du
lecteur triadique. L’application
\[
 T:E_{114}\longrightarrow\mathcal O,
 \qquad \varepsilon_i\longmapsto T_{\rm tri}(\varepsilon_i),
 \qquad |\operatorname{im}T|=106,
\]
classe les expressions selon leur rôle opératoire. Chaque
\(\varepsilon_i\), y compris les cibles comparées, est fixé comme
entrée de \(T\) ; cette application ne génère donc pas les valeurs qu’elle
évalue. La partition fonctionnelle est
\[
\begin{array}{l|r}
\text{famille fonctionnelle}&\text{entrées}\\ \hline
\text{noyau structurel}&5\\
\text{couplage et clôture}&1\\
\text{ponts entre modes}&7\\
\text{angulaire et clôture}&22\\
\text{exponentielle--logarithmique}&12\\
\text{spectral et régularisation}&14\\
\text{hyperbolique et modulaire}&2\\
\text{algébrique et métrique}&50\\
\text{divers typés}&1
\end{array}
\qquad
\sum n_i=114.
\]
La centralité de réseau, les chemins minimaux et les communautés calculées sur
le graphe induit par \(T\) décrivent la connectivité au sein de la classification ;
ils n’établissent ni causalité mathématique ni généalogie HMT.

Sept fibres de \(T\) réunissent quinze expressions ayant une lecture commune. Une
fibre partagée n’identifie pas non plus leurs valeurs. En particulier,
\(\varphi^{-1}\), \(\varphi\) et \(\varphi^2\) sont des réels distincts que
l’application \(T\) regroupe selon une même orbite d’auto-échelle ; la projection
commune conserve la fonction dynamique et oublie la coordonnée archimédienne qui les
sépare.

L’énumération exhaustive a parcouru les \(114\) entrées de \(E_{114}\),
quatre cent soixante-huit candidats directs \(U_{12}\) et cent quatre
mille neuf cent soixante-seize germes pondérés sans trouver, dans
ce domaine fini, de rival reproduisant la correction dodécaphasique
orientée de \(\alpha_{\rm HMT}\). Le domaine de la recherche est ainsi
borné par ces trois ensembles. Le recensement démontre l’absence d’un
rival dans la classe énumérée ; il ne démontre pas l’unicité parmi toutes les
fonctionnelles possibles et n’intervient pas comme prémisse du théorème de
\(\alpha_{\rm HMT}\).

Pour chaque \(\varepsilon_i\), la taxonomie attribue le quadruplet
\((\text{objet},\text{lecteur},\text{domaine},\text{fonction structurelle})\)
utilisé dans l’énumération.

\section{Confluence spectrale, informationnelle et thermodynamique de \(\log3\)}

L’identité

\[
\log3=C_1+C_2-2C_0
=S_{\TRIT}/k_B
=E_P/(k_B\Theta_{\rm clk}),
\]

relie trois réalisations distinctes. Le premier membre est un défaut
résiduel, le deuxième une entropie de décision et le troisième un rapport
entre énergie et température. Leur égalité démontre la conservation du
scalaire sous le changement de réalisation, sans identifier les espaces
d'états.

La première égalité se déduit des trois canaux
résiduels
\[
Z_r(s)=\sum_{\substack{n\ge1\\n\equiv r\ ({\rm mod}\ 3)}}n^{-s},
\qquad r=0,1,2.
\]
Alors
\[
Z_0(s)=3^{-s}\zeta(s),\quad
Z_1(s)+Z_2(s)=(1-3^{-s})\zeta(s),\quad
Z_1(s)-Z_2(s)=L(s,\chi_{-3}).
\]
Si \(C_r\) désigne la partie finie de \(Z_r\) en \(s=1\), le développement
de \(3^{-s}\) donne
\begin{equation}
C_0+C_1+C_2=\gamma,qquad
-2C_0+C_1+C_2=\log3,qquad
C_1-C_2=\frac{\pi}{3\sqrt3}.
\label{eq:residual-three-covectors}
\end{equation}
Les trois covecteurs sont linéairement indépendants. Ils retrouvent, à partir d'une seule
décomposition résiduelle, le terme fini total, la dilatation de la
classe zéro et l'orientation entre les deux classes non nulles. \(\log3\) est
donc un invariant de changement d'échelle avant d'être interprété comme
entropie ou température.

\section{Classification épistémologique des constantes dérivées}

Les contributions du chapitre sont classées selon trois modalités :

\begin{enumerate}
\item attribution d'un opérateur interne à une constante connue, comme dans
le cas de Catalan ;
\item évaluation composée exacte de deux caractères internes, comme
\(L(1,\chi_{12})\) ;
\item obtention d'une constante à partir d'un corps ou d'une famille
construits en interne, comme \(\gamma_{\mathbb Q(\sqrt3)}\) ou \(A_k\).
\end{enumerate}

La provenance mathématique détermine le statut de récupération ou de
nouveauté ; la proximité décimale n'intervient pas dans cette classification.

\begin{remark}[Contrôle négatif]
Le caractère \(\chi_{12}\) doit avoir exactement le motif
\((+,-,-,+)\) dans \((1,5,7,11)\) ; tout autre motif invalide
\cref{eq:L1chi12,eq:L2chi12}. L'identité
\cref{eq:odd-zeta-bendersky} et la factorisation
\cref{eq:dedekind-sqrt3} sont des contrôles indépendants.
\end{remark}

\subsection*{Portée logique et domaine de validité}
\(J^2=-I\), Catalan, \(\chi_{12}\), ses valeurs \(L\), Euler--Kronecker et la tour de Bendersky--Glaisher : \emph{Théorème interne exact}. La réunion en un réseau \(3/4/12\) est une \emph{Formalisation nouvelle} qui conserve les provenances mathématiques.


Tous ces résultats sont unis parce qu’ils transforment la valeur terminale en l’expression quantitative d’une structure préalablement construite. L’exactitude numérique accompagne une conception différente de ce qu’est une constante fondamentale. La connexion nonadique n’engendre pas d’abord trois constantes pour en ajouter ensuite une quatrième : le prolongement \(w_6\to w_{12}\to w_{18}\to w_{24}\to w_{30}\to R_{36}\to G_9\) stabilise un état coinductif commun à partir duquel \(\pi\), \(e\), \(\varphi\) et \(\alpha\) sont publiés de manière corrélée. Leurs lecteurs distinguent clôture, propagation, auto-échelle et couplage ; la coordonnée \(\alpha\) conserve, au sein de cette génération conjointe, sa réalisation signée et dodécaphasique et ses deux voies internes de dérivation.\par

Dans la description habituelle, une constante est souvent présentée comme un nombre placé devant une équation : une quantité dont la théorie a besoin, mais dont la valeur doit être reçue de l’extérieur. Dans l’holographie modulaire triadique, c’est l’inverse qui se produit. La constante apparaît à la fin, lorsqu’un même état a été parcouru par plusieurs lecteurs et que l’on vérifie quelle relation leur a tous survécu.\par

C’est pourquoi les constantes du document sont des grandeurs dotées d’une généalogie : des restes invariants d’une transformation. Chacune dit ce qui a été conservé lorsque quelque chose a changé :\par

\begin{itemize}[leftmargin=*,itemsep=0.35em]
\item un feuillet a été inversé ;
\item un cycle est revenu à sa phase initiale ;
\item une mémoire a traversé la couture ;
\item une description intérieure a été publiée au bord ;
\item une grandeur adimensionnelle a reçu une réalisation physique ;
\item un système a changé d’échelle en conservant intégralement sa généalogie.
\end{itemize}

La mécanique dimensionnelle ajoute ensuite les droites de grandeur et les unités. Mais la structure qui sélectionne la constante est antérieure. La valeur réelle apparaît à la fin parce que \(\mathbb R\) fonctionne ici comme évaluation terminale : elle est l’ombre quantitative d’une construction qui était déjà complète avant de devenir un nombre décimal.\par

Dans cette perspective, les neuf développements considérés forment une seule narration sur l’orientation, la mémoire, la géométrie, l’information et l’échelle.\par

%
\endgroup
% Rectificación quirúrgica: generación conjunta, profundidad arbitraria y
% estatuto tipado de los transportes ternarios.  Este bloque pertenece al
% capítulo 27 y precede a la separación expositiva de los cuatro lectores.

\section{Théorème conjoint de génération coinductive et de profondeur arbitraire}
\label{sec:c27-generacion-coinductiva-profundidad-arbitraria}

La structure de départ de ce théorème n'est pas l'ensemble nu de neuf
symboles. C'est le système typé
\[
 \mathfrak H_9=
 \bigl(\mathrm{APP},\mathrm{TRIT},\mathrm{TPK};
       \widehat{\mathcal X}^{\mathrm{enr}},\mathcal U,
       \Gamma_9,\operatorname{Ext},\operatorname{tr}\bigr),
\]
construit sur la double réalisation positionnelle et arithmétique de
\(\mathcal D_9=\{1,\ldots,9\}\). APP accumule et replie au moyen de ses feuillets
\(\Sigma\) et \(\Pi\), en conservant résidu et quotient ; TRIT oriente chaque
transition dans l'un de ses trois régimes ; et TPK sélectionne, transporte,
mémorise et revient sur la fibre enrichie. En particulier,
\[
 M_{\rm ph}:\mathcal D_9\longrightarrow\{+1,-1,0\},
 \qquad
 m_{\rm ph}:=\bigl(M_{\rm ph}(1),\ldots,M_{\rm ph}(9)\bigr)^{\mathsf T},
\]
\[
 \operatorname{Tra}_t(y,d):=\operatorname{Lift}(T_d)(y),
 \qquad
 \operatorname{Upd}_t:=\operatorname{Rec}_t\circ\operatorname{Mem}_t,
\]
et, pour tout \(x\in\widehat{\mathcal X}^{\mathrm{enr}}\),
\[
 \mathcal U_t(x)
 =\operatorname{Upd}_t\!\left(
   \operatorname{Tra}_t\!\left(
    \operatorname{Sel}_t\!\left(x,M_{\rm ph}(g(t))\right),d(t)
   \right)\right).
\]
Cette composition est la donnée dynamique effective : la carte à neuf symboles
est son support fini, non un substitut de la dynamique. En particulier,
\(M_{\rm ph}(g(t))\) est la valeur scalaire que reçoit le sélecteur, tandis
que \(m_{\rm ph}\) est le vecteur des neuf valeurs de la fonction ; aucun
ne se compose directement avec un opérateur d'état.

Pour \(r\in6\mathbb N_{>0}\), soient
\(\widehat{\mathcal X}^{\mathrm{enr}}_r\) les états tronqués avec feuillet,
orientation, résidu, quotient, retenue, voie, bord et mémoire, et soient
\[
 \operatorname{Ext}_{r+6,r}\subseteq
 \widehat{\mathcal X}^{\mathrm{enr}}_r\times
 \widehat{\mathcal X}^{\mathrm{enr}}_{r+6},
 \qquad
 \operatorname{tr}_{r+6,r}:
 \widehat{\mathcal X}^{\mathrm{enr}}_{r+6}
 \longrightarrow
 \widehat{\mathcal X}^{\mathrm{enr}}_r
\]
la relation de prolongement admissible et la troncature. Pour
\(x_r\in\widehat{\mathcal X}^{\mathrm{enr}}_r\), la fibre locale
\[
 \operatorname{Ext}_{r+6,r}(x_r)
 :=\{x_{r+6}:(x_r,x_{r+6})\in\operatorname{Ext}_{r+6,r}\}
\]
peut contenir plusieurs enfants. La compatibilité correcte est
\[
 (x_r,x_{r+6})\in\operatorname{Ext}_{r+6,r}
 \Longrightarrow
 \operatorname{tr}_{r+6,r}(x_{r+6})=x_r.
\]
Dans \(R_{36}\), les trois lignes retrouvées
\(222220,021222,102011\), transportées au calibre actif par la permutation
\(p_0=(0,1,2,4,5,3)\) en indices commençant à zéro ---de manière équivalente,
\(p_1=(1,2,3,5,6,4)\) en indices commençant à un---, portent avant toute
réalisation réelle les caractères
internes de clôture, de propagation et d'auto-échelle. Leur triplet, avec le
registre signé \(U_{12}\) et le sceau \(K\), porte la coordonnée interne de
\(\alpha\). Pour chaque
\(\chi\in\{\pi,\varphi,e,\alpha\}\), nous écrivons l'état ainsi enrichi
comme
\[
 x_r^\chi=(\widetilde x_r,\chi_{\rm HMT},
            \varepsilon_r^\chi,N_r^\chi).
\]
Ici \(\varepsilon_r^\chi\) est le résidu interne produit par le quotient
APP dans l'état, non \(\operatorname{frac}(P_{\rm ar}(\chi))\). La division
euclidienne interne fixe
\[
 d_r^\chi=\operatorname{Dig}_{729}(\varepsilon_r^\chi),
 \quad 0\le d_r^\chi<729,
 \quad
 \varepsilon_{r+6}^\chi=729\varepsilon_r^\chi-d_r^\chi,
 \quad
 N_{r+6}^\chi=729N_r^\chi+d_r^\chi.
\]
La sélection, le transport et la mise à jour sur la section
\(\chi\) sont typés par
\[
 \operatorname{Sel}_r^\chi(x)
 :=\operatorname{Sel}\!\left(
 x,M_{\rm ph}(\rho_9(r/6+1))\right),
 \qquad
 \operatorname{Tra}_r^\chi(y,d):=\operatorname{Lift}(T_d)(y),
 \qquad
 \operatorname{Upd}_r^\chi:=\operatorname{Rec}_r^\chi\circ
 \operatorname{Mem}_r^\chi.
\]
La transition complète est
\[
 \mathcal U_r^\chi(x)
 :=\operatorname{Upd}_r^\chi\!\left(
  \operatorname{Tra}_r^\chi\!\left(
   \operatorname{Sel}_r^\chi(x),d_r^\chi
  \right)\right),
\]
et ne lit que l'état enrichi précédent : phase, mémoire, retenue, feuillet,
orientation, voie, bord, signature, registre, caractère et résidu internes.
Elle ne lit pas un chiffre futur, \(P_{\rm ar}\), ni un encadrement de Machin,
factoriel ou de Perron. Nous définissons la sous-relation en avant
\[
 \operatorname{Ext}^{\chi,\rightarrow}_{r+6,r}
 :=\{(x_r^\chi,\mathcal U_r^\chi(x_r^\chi))\}
 \subseteq\operatorname{Ext}_{r+6,r}.
\]
Ce n'est pas la sélection rétrospective
\(\operatorname{Ext}\cap\{P_{\rm ar}=\chi\}\) : le caractère est déjà une
coordonnée causale de \(x_{36}^\chi\). Les équations précédentes donnent
\(\lfloor N_{r+6}^\chi/729\rfloor=N_r^\chi\) et conservent feuillet, orientation,
retenue, bord, voie, signature, registre et généalogie.

La relation complète \(\operatorname{Ext}\) demeure multivoque.
L'unicité requise est celle de chaque section en avant sur son germe R36
produit,
\[
 \#\Bigl\{(x_r^\chi)_{r\ge36}:x_{36}^\chi\text{ fixé},\ 
 (x_r^\chi,x_{r+6}^\chi)
 \in\operatorname{Ext}^{\chi,\rightarrow}_{r+6,r}\Bigr\}=1,
\]
non la fausse égalité \(\#\operatorname{Ext}_{r+6,r}(x_r)=1\) à chaque porte
locale. Cette compatibilité se prolonge à tout niveau. La notation
\[
 w_6\longrightarrow w_{12}\longrightarrow w_{18}\longrightarrow w_{24}
 \longrightarrow w_{30}\longrightarrow R_{36}\longrightarrow G_9
\]
désigne ces applications et leurs changements de type, non une liste d'étiquettes.
L'holonomie \(\Gamma_9\) satisfait
\[
 g(k+9)=g(k),
 \qquad q^{(9)}(k+9)=q^{(9)}(k)+1,
\]
de sorte que la phase revient tandis qu'avancent le quotient, la voie et la
mémoire de l'état.

Une histoire canonique générique basée sur \(a_{(1,1)}\) et
\(t_j=M_{\rm ph}(\rho_9(j+1))\) atteste la non-vacuité de
\(\operatorname{Ext}\) ; elle n'individualise pas à elle seule les quatre sections
précédentes. Celles-ci sont raccordées par leurs coordonnées internes R36 et par la
règle en avant qui vient d'être définie.

\begin{theorem}[Construction inaugurale à partir des germes]
\label{thm:c27-construccion-inaugural-semillas}
Soit
\[
 I_9=\{0,\ldots,8\},\qquad D=\{N,E,S,O\},\qquad
 \mathcal S=I_9^2\times D,
\]
et soit \(\Omega_{\mathrm{seed}}=\mathcal S_+\times\mathcal S_\times\) le
domaine des deux curseurs qui réalisent les feuillets additif et multiplicatif
d'APP. Alors
\[
 |\Omega_{\mathrm{seed}}|=(9^2\cdot4)^2=104\,976,
\]
et la récurrence de \(54\) pas et la réduction ternaire définissent
\[
 \Omega_{\mathrm{seed}}
 \xrightarrow{\ T_{54}^{\rm seed}\ }
 \mathcal U_6^{\mathrm{cat}}
 \xrightarrow{\ \rho_3\ }
 \mathcal W_6
 \hookrightarrow\mathbb F_3^6,
\]
avec
\[
 |\mathcal U_6^{\mathrm{cat}}|=468,
 \qquad |\mathcal W_6|=243,
 \qquad |\mathbb F_3^6|=729.
\]
Les fibres satisfont exactement
\[
 432\cdot144+18\cdot432+18\cdot1944=104\,976.
\]
La relation d'extension du TPK prolonge les mots survivants à la
famille des sections compatibles de profondeur arbitraire. Les sous-relations
en avant ancrées dans R36 individualisent la génération coinductive
corrélée de \(\pi,\varphi,e\) et de \(\alpha\) ; la relation complète peut
continuer à posséder plusieurs prolongements locaux. Sur cette famille sont
construites conjointement l'ombre archimédienne et la réalisation
exceptionnelle du \cref{thm:c27-tesis-rectora-doble-proyeccion}.
\end{theorem}

\begin{proof}
Le produit cartésien précédent énumère exhaustivement les positions et
orientations des deux curseurs. L'application \(T_{54}^{\rm seed}\) est obtenue en
composant à chaque pas la sélection de feuillet, le transport cardinal, la
mise à jour du quotient et la mémoire de retenue. L'énumération de son
image contient \(468\) sextuplets distincts ; la somme des cardinalités de
ses fibres est l'identité exposée. L'application composante par composante
\(\rho_3\) contient \(243\) mots distincts. Enfin, la non-vacuité à
tout horizon et la compatibilité de \(\operatorname{Ext}\) avec les
troncatures produisent la limite inverse. La fonctionnalité de chaque
\(\operatorname{Ext}^{\chi,\rightarrow}\), à partir de l'état R36 qui
porte déjà \(\chi_{\rm HMT}\), démontre par récurrence l'unicité de cette section
globale sans exiger que chaque fibre locale de \(\operatorname{Ext}\) soit
un singleton. Chaque affirmation est une identité interne ou une conséquence de la
compatibilité projective ; aucune n'utilise en entrée un nombre réel, une
constante reconnue ou une table décimale.
\end{proof}

\begin{theorem}[Thèse directrice : un continu et deux projections]
\label{thm:c27-tesis-rectora-doble-proyeccion}
Soit \(\Omega_\Gamma^{\mathrm{cal}}\subset
X_\infty^{\mathrm{cal}}\) l'espace des sections survivantes des
tours enrichis du TPK. La retenue des trois tours
\(\gamma_-,\gamma_0,\gamma_+\) définit, sans évaluation archimédienne préalable,
l'homéomorphisme
\[
 \beta:\Omega_\Gamma^{\mathrm{cal}}\xrightarrow{\ \sim\ }
 (\mathbb F_3^6)^{\mathbb N}=: \mathcal W_\infty.
\]
Les formes équilibrées normalisées
\[
 a_k+c_k=\rho_k+3c_{k+1},\qquad
 \rho_k\in\{-1,0,+1\},
\]
construisent l'anneau ordonné \(\mathbb Z_{\rm HMT}\), son corps des
fractions \(\mathbb Q_{\rm HMT}\) et la complétion
\[
 \mathbb R_{\rm HMT}:=
 \operatorname{Cau}(\mathbb Q_{\rm HMT})/{\asymp}.
\]
La structure discrète conjointe est le produit fibré
\[
\mathcal C_{\mathrm{cont}}^{\mathrm{disc}}
=\bigl(\mathbb Z_{\rm HMT}\times
\Omega_\Gamma^{\mathrm{cal}}\bigr)
\times_{\mathcal W_\infty}X_\infty^{\mathrm{cal}}.
\]
Ses deux applications naturelles sont
\[
 P_{\mathrm{ar}}:\mathbb Z_{\rm HMT}\times
 \Omega_\Gamma^{\mathrm{cal}}\longrightarrow\mathbb R_{\rm HMT},
 \qquad
 Q_\infty:X_\infty^{\mathrm{cal}}\longrightarrow
 (\mathscr C_W^{\mathrm{cal}})^{\mathbb N}.
\]
Si \(\beta(\xi)=(b_q)_{q\geq0}\),
\(\nu(b)=\sum_{i=1}^{6}b_i3^{6-i}\) et
\[
 q_n(m,\xi)=m+\sum_{q=0}^{n-1}
 \frac{\nu(b_q)}{729^{q+1}},
\]
alors \((q_n)\) est de Cauchy et \(P_{\mathrm{ar}}(m,\xi)\) est sa classe
dans \(\mathbb R_{\rm HMT}\). Les cylindres rationnels
\[
 I_n(m;b_0,\ldots,b_{n-1})=
 \left[q_n,q_n+729^{-n}\right]
\]
sont compatibles, ont un diamètre tendant vers zéro et la partition récurrente
en \(729\) sous-cylindres construit, pour toute classe de Cauchy, un ruban
\((b_q)\) qui la représente. En conséquence,
\[
 P_{\mathrm{ar}}\ \text{est surjective},\qquad
 (\mathbb Z_{\rm HMT}\times\Omega_\Gamma^{\mathrm{cal}})
 /{\sim_{\rm ar}}\ \cong\ \mathbb R_{\rm HMT}.
\]
Ce n'est qu'après cette construction interne que l'unicité du corps ordonné
complet archimédien reconnaît \(\mathbb R_{\rm HMT}\cong\mathbb R\).

La seconde projection conserve l'incidence calibrée de chaque niveau ; dans les
coordonnées \(x=((w_j,f_j))_{j\geq1}\), elle agit par
\[
 Q_\infty(x)
 =\bigl(f_j,\Gamma_{f_jA_Wf_j^{-1}}(w_j)\bigr)_{j\geq1}.
\]
Ainsi, \(\mathbb R_{\rm HMT}\), reconnu ultérieurement comme \(\mathbb R\),
est l'ombre archimédienne obtenue en contractant la généalogie complète des
cylindres. La projection exceptionnelle conserve cette généalogie comme incidence
de bord et admet la réalisation successive
Paley--Witt--Golay--Leech--\(V^\natural\)--Monstre. Les deux branches proviennent
du même état et aucune n'est reconstruite rétrospectivement à partir de l'autre.
\end{theorem}

\begin{proof}
Les trois tours enrichis existent après chaque préfixe et leur retenue
les distingue par \(-1,0,+1\). La concaténation et la troncature sont
inverses, ce qui prouve la bijectivité et la continuité de \(\beta\). Pour une
classe de Cauchy de \(\mathbb Q_{\rm HMT}\), on choisit d'abord son cylindre
unitaire ordonné, puis, récursivement, l'unique sous-cylindre semi-ouvert
de la partition interne en \(729\) morceaux qui contient la classe ; sur un
bord commun, les deux écritures sont identifiées par \(\sim_{\rm ar}\).
La suite des extrémités gauches est précisément \((q_n)\), converge vers la
classe donnée et prouve la surjectivité sans présupposer un développement de
\([0,1]\) ni un point de \(\mathbb R\). L'identification au corps réel
est le théorème d'unicité appliqué ensuite à \(\mathbb R_{\rm HMT}\).

Les sous-espaces qui satisfont la fermeture d'un caractère précis sont des filtres
de survie au sein de la capacité globale. Leur fonction est d'individualiser
les sections de \(\pi,e,\varphi,\alpha\) ; la surjectivité du continu
provient de la complétion de \(\mathbb Q_{\rm HMT}\). Cette séparation des
domaines empêche
de transformer une condition de sélection de constantes en prémisse de la
couverture de \(\mathbb R\).

Dans la seconde branche, la compatibilité
\(r_{n,m}Q_n=Q_m p_{n,m}\) permet de passer à la limite et définit \(Q_\infty\).
La conjugaison par \(f_j\) transporte le calibre sans effacer l'incidence.
Les foncteurs ultérieurs de code, de réseau et d'algèbre de vertex réalisent
la chaîne exceptionnelle. L'égalité des rubans dans le produit fibré
construit la source conjointe et conserve des informations complémentaires : valeur
archimédienne dans la première ; incidence, bord et généalogie dans la seconde.
\end{proof}

\begin{corollary}[Développement infini comme objet mathématique]
\label{cor:c27-expansion-infinita-objeto}
Pour chaque caractère numérique \(\chi\) publié par l'état coinductif, la
lecture décimale est une séquence complète
\[
 \operatorname{Dec}_\chi(x_\infty)
 =(d_{\chi,1},d_{\chi,2},\ldots)
 \in\{0,\ldots,9\}^{\mathbb N}.
\]
Son préfixe de longueur \(N\) est la projection
\(\rho_N\operatorname{Dec}_\chi(x_\infty)\). Ainsi, mille, dix mille ou un
million de chiffres sont des coupes du même objet infini ; la génération ne
consiste pas à relever successivement une borne expérimentale.
\end{corollary}

\begin{theorem}[Génération coinductive corrélée de \(\pi\), \(\varphi\), \(e\) et \(\alpha\)]
\label{thm:c27-cuadrirrelacion-generada}
Le profil complet du système \(\mathfrak H_9\) détermine la famille
corrélée de sections en avant compatibles
\[
 \boldsymbol x_\infty=
 (x_\infty^\pi,x_\infty^\varphi,x_\infty^e,x_\infty^\alpha)
 \in
 \widehat\Omega_\pi\times_{R_{36}}
 \widehat\Omega_\varphi\times_{R_{36}}
 \widehat\Omega_e\times_{R_{36}}\widehat\Omega_\alpha,
 \qquad
 \widehat\Omega_\chi:=
 \varprojlim_{r\ge36}
 (\widehat{\mathcal X}^{\mathrm{enr},\chi}_r,
  \operatorname{Ext}^{\chi,\rightarrow}_{r+6,r}),
\]
et quatre caractères corrélés
\[
 \bigl(\Pi_{\mathrm{cl}},\Pi_{\mathrm{aut}},
       \Pi_{\mathrm{prop}},\Pi_{\mathrm{dod}}\bigr)
       (\boldsymbol x_\infty)
 =\bigl(\pi_{\mathrm{HMT}},\varphi_{\mathrm{HMT}},
        e_{\mathrm{HMT}},\alpha_{\mathrm{HMT}}\bigr).
\]
Les trois premiers caractères proviennent, respectivement, de l'unique
orbite de clôture, de l'unique orbite d'auto-échelle et de l'unique
orbite de propagation du catalogue enrichi. Le quatrième provient de la
coordonnée primitive signée de la famille terminale corrélée et de sa
fermeture dodécaphasique. Les quatre coordonnées sont donc des sorties de la
génération coinductive corrélée de profondeur arbitraire gouvernée par
la même dynamique APP--TRIT--TPK, avec retour de phase, avancée de mémoire et
conservation généalogique, bien que leurs applications de publication soient
distinctes.
\end{theorem}

\begin{proof}
Le recensement exhaustif
\(104,976\to468\to243\subset\mathbb F_3^6\) décompose l'image réalisée
en \(43\) orbites diédriques. Les invariants de fibre laissent une seule orbite
dans chacune des classes de clôture, de propagation et d'auto-échelle ; le calibre
interne y oriente, sans consultation décimale, les représentants
\[
 w_6^{\mathrm{cl}}=010211,
 \qquad w_6^{\mathrm{prop}}=201101,
 \qquad w_6^{\mathrm{aut}}=121200.
\]
Les sous-relations en avant conservent le caractère interne et les prédicats qui
définissent les trois classes ; Machin, la récurrence factorielle et Perron ne choisissent
aucun enfant de \(\operatorname{Ext}\). Dans le
canal de clôture, la composition à cinq régions avec retour unitaire produit
les rapports \(1/5\), \(1/239\) et la relation \(120/119\) ; l'identité de
Machin reconnaît exactement le demi-tour. Dans le canal de propagation, la
loi de semi-groupe normalisée impose
\((n+1)a_{n+1}=a_n\), donc \(a_n=1/n!\) et \(P_1=e\). Dans le canal
d'auto-échelle, l'incidence
\[
 F_{\mathrm{av}}=\begin{pmatrix}0&1\\1&1\end{pmatrix}
\]
possède un unique rayon positif et son équation propre impose
\(x^2-x-1=0\), \(x>0\), c'est-à-dire \(x=\varphi\).

Dans le profil TPK complet, la lecture terminale
\(x_{\mathrm{term}}\mapsto(B_{\mathrm{term}},U_{12}^{\mathrm{sgn}})\)
publie simultanément la carte des trois canaux et la coordonnée
primitive signée. La
seconde est identifiée réversiblement au sceau \(K\) et à ses différences
dodécaphasiques. La récurrence euclidienne avec retenue
\[
 s_m=P_m+E_m-\Phi_m-K_m+c_{m+1},\qquad
 A_m=s_m\bmod1000,\qquad
 c_m=(s_m-A_m)/1000
\]
produit \(\alpha_{\mathrm{HMT}}\) sans la recevoir en entrée. La
caractérisation analytique indépendante commence par le noyau fini
\(E_{21/22}^{(9)}\), qui détermine une racine simple \(\xi_9\) dans
\((0,10^{-2})\) sans consulter la sortie dodécaphasique. L'itération du même
lecteur gradué produit le système inverse
\(E_{21/22}^{(6)}\leftarrow E_{21/22}^{(9)}\leftarrow
E_{21/22}^{(12)}\leftarrow\cdots\). Ses troncatures commutent avec celles de la
voie dodécaphasique et les deux familles déterminent à chaque profondeur le même
cylindre survivant. Comme leurs diamètres tendent vers zéro, l'intersection est
un singleton et l'on obtient exactement
\(\alpha_A=\alpha_B=\alpha_{\mathrm{HMT}}\). L'égalité
\(\operatorname{Tr}_9(\xi_9^{-1})=
\operatorname{Tr}_9(\alpha_{12}^{-1})\) n'est que le certificat fini
d'ordre neuf de cette identité de sections complètes. Ainsi, la séparation
des lecteurs conserve l'origine commune des quatre coordonnées sans introduire
un prolongement en suspens ni une valeur cible.
\end{proof}

\begin{proposition}[Registre dodécaphasique de l'état complet antérieur à \(\alpha\)]
\label{prop:c27-registro-u-k-producido}
L'état terminal enrichi possède la double publication
\[
x_{\mathrm{term}}
\xrightarrow{(\pi_{\mathrm{term}},R_{\mathrm{sgn}})}
(B_{\mathrm{term}},U),
\]
où
\[
U=(2378,1406,2479,-452,998,-551,
-668,-204,-371,-322,-28,-997).
\]
En ordonnant les coordonnées selon les trois orbites de pas trois, la transformée
\[
\mathcal H_{12}
=P_{\mathrm{orb}}^{-1}(H_4\oplus H_4\oplus H_4)P_{\mathrm{orb}},
\qquad H_4^2=4I_4,
\]
possède l'inverse \(\mathcal H_{12}/4\) sur son image entière et détermine
de manière univoque
\[
K=(234,543,140,729,659,824,621,058,914,794,146,601).
\]
Par conséquent, au sein du profil TPK complet,
\[
x_{\mathrm{term}}\xrightarrow{R_{\mathrm{sgn}}}U
\xrightarrow{\mathcal H_{12}^{-1}}K
\]
est une composition exacte antérieure à la retenue qui publie \(\alpha\). La matrice
visible \(B_{\mathrm{term}}\) et le registre signé \(U\) sont deux
coordonnées du même état ; la seconde conserve les champs d'orientation,
d'opposition, de quotient et de registre que la première contracte. La composition
\[
 \Omega_{\mathrm{seed}}\longrightarrow
 \mathcal U_6^{\mathrm{cat}}\longrightarrow
 \widehat{\mathcal X}^{\mathrm{enr}}_{\mathrm{term}}
 \xrightarrow{R_{\mathrm{sgn}}}U
 \xrightarrow{\mathcal H_{12}^{-1}}K
\]
construit chaque coordonnée par APP--TRIT--TPK et conserve leurs registres. Ses
entrées sont les germes, les feuillets et les règles opératorielles ; \(U\), \(K\),
leurs différences, \(\alpha\) et les développements publiés appartiennent au
codomaine d'étapes successives.
\end{proposition}

\begin{theorem}[Clôture bidirectionnelle de la seconde voie de \(\alpha\)]
\label{thm:c27-alpha-segunda-via-bidireccional}
Le prolongement du noyau des lacunes jusqu'aux degrés \(7{:}9\) détermine
\[
\begin{aligned}
B_9(\alpha):={}&\log_{10}\!\left(\frac{10}{9}\right)
+\frac74\alpha^2-\frac{\alpha^4}{20}
+\frac{2\alpha^5}{21}+\frac{\alpha^6}{46}\\
&+\frac{\alpha^7}{120}-\frac{\alpha^8}{45}
-\frac{2\alpha^9}{495}.
\end{aligned}
\]
En relevant la rotation à l'inconnue positive \(T\), la résolvante équivaut à
\[
\boxed{\alpha T^2-B_9(\alpha)T+\alpha^3=0.}
\]
Avec
\(\xi=\operatorname{arcosh}(B_9(\alpha)/(2\alpha^2))\), les solutions sont
\[
T_\pm=\alpha e^{\pm\xi},
\qquad T_+T_-=\alpha^2,
\qquad
J_\alpha(T)=\frac{\alpha^2}{T},
\qquad J_\alpha(T_\pm)=T_\mp.
\]
La chambre \(6<T<7\) sélectionne \(T_+\) ; la clôture finie
\(\pi_9=T_+/2\) se prolonge au moyen de \(\Gamma_9\) jusqu'à
\(\pi_{\mathrm{HMT}}\). L'identité
\(\alpha=\sqrt{T_+T_-}\) fixe la moyenne géométrique des branches conjuguées.
Le jet \(7{:}9\), la racine simple, l'involution et le retour constituent la
seconde voie interne complète.
\end{theorem}

\begin{theorem}[Couplage exact du secteur \(\alpha\) avec le Monstre]
\label{thm:c27-alpha-rama-excepcional}
Dans la carte d'incidence du même registre dodécaphasique, les racines radiales
de \(A_2^{12}\) déterminent
\[
v_\alpha=4\rho_1+\rho_2+\cdots+\rho_{12},
\qquad v_\alpha^2=54,
\]
et le voisin d'indice trois
\[
N_{\alpha,0}
=\{x\in N(C_W):\langle x,v_\alpha\rangle\equiv0\pmod3\},
\qquad
N_\alpha=N_{\alpha,0}+\mathbb Z\frac{v_\alpha}{3}.
\]
Le réseau \(N_\alpha\) est pair, unimodulaire, de rang vingt-quatre et sans
racines, de sorte que \(N_\alpha\cong\Lambda_{24}\). La chaîne
\[
C_W\longrightarrow N(A_2^{12})
\xrightarrow{v_\alpha}\Lambda_{24}
\longrightarrow V^\natural
\xrightarrow{\operatorname{Aut}}\mathbb M
\]
expose la charnière exceptionnelle. Le scalaire \(\alpha_{\mathrm{HMT}}\) et le
vecteur \(v_\alpha\) appartiennent à des codomaines distincts et partagent l'état
dodécaphasique d'origine ; cette corrélation typée est le raccordement aux
symétries du Monstre.
\end{theorem}

\begin{proposition}[Statut des transports \(L_0,L_1\)]
\label{prop:c27-estatuto-l0-l1}
Les endomorphismes \(L_0,L_1\in\operatorname{End}(\mathbb F_3^6)\) et le
calendrier \((L_0,L_0,L_1,L_1)\) appartiennent à la loi de transport de l'état
TPK enrichi. Les six transitions publiées de rang plein fixent
chaque matrice de manière unique par \(L=X^{-1}Y\) dans \(\mathbb F_3\), et le
calendrier est le seul des \(16\) calendriers binaires qui conserve
simultanément les trois biographies structurelles. Sa fonction est de prolonger
les représentants déjà sélectionnés ; l'évaluation archimédienne et la
nomenclature conventionnelle sont appliquées ensuite.
\end{proposition}

Le domaine de cette proposition est \(\mathfrak H_9\), le système complet
APP--TRIT--TPK avec état enrichi. Le recensement APP constitue son étape
initiale ; TRIT apporte l'orientation et TPK construit le prolongement, la
mémoire et le retour. Cette séquence fixe exactement le domaine de chaque
flèche et l'emplacement de chaque sortie.

\begin{theorem}[Profondeur arbitraire et détermination des développements complets]
\label{thm:c27-profundidad-arbitraria}
Pour chaque
\(\chi\in\{\pi_{\mathrm{HMT}},e_{\mathrm{HMT}},
\varphi_{\mathrm{HMT}}\}\) et chaque \(N\ge1\), il existe un unique préfixe
décimal \(p_{\chi,N}\) de longueur \(N\), produit par une troncature
finie de l'état enrichi, tel que
\[
 \rho_{N+1,N}(p_{\chi,N+1})=p_{\chi,N}.
\]
La famille compatible est la famille des projections d'un unique élément
\[
 p_{\chi,\infty}
 \in\varprojlim_N\{0,\ldots,9\}^{N}
 \cong\{0,\ldots,9\}^{\mathbb N},
\]
et l'intersection de ses cylindres archimédiens est constituée d'un unique réel.
Par conséquent, la connexion engendre le développement complet ; chaque précision
finie est une projection de cette sortie infinie.
\end{theorem}

\begin{proof}
La preuve comporte deux étapes qui ne peuvent être inversées. Dans la première, entièrement
antérieure à la carte décimale, la section survivante de
\(\mathfrak H_9\) publie trois caractères exacts. Le caractère de clôture
est fixé par la pentafibre orientée, sa pente primitive \(1/5\), la
composition de ses quatre compagnes \(120/119\) et le compensateur entier
\(239\) ; le produit gaussien correspondant se termine sur la demi-droite
réelle négative et fixe de manière univoque le demi-tour orienté. Le caractère de
propagation est l'unique série formelle normalisée par
\(a_0=1\) et \((n+1)a_{n+1}=a_n\). Le caractère d'auto-échelle est l'unique rayon
expansif positif de
\[
 F_{\rm av}=\begin{pmatrix}0&1\\1&1\end{pmatrix},
 \qquad X^2-X-1=0.
\]
Ces trois conditions sont des sorties de l'état enrichi ; elles ne contiennent pas de
préfixes décimaux et ne consultent pas de valeurs conventionnelles.

Ce n'est qu'à la seconde étape qu'agit la publication archimédienne. Pour chaque caractère
déjà généré et chaque \(N\), soit \(\operatorname{Dec}_N\) l'unique cylindre
décimal canonique de diamètre \(10^{-N}\) qui contient sa réalisation exacte.
Ce cylindre est une coordonnée publiée du caractère : ce n'est pas une fibre locale
de \(\operatorname{Ext}\) et il n'est pas utilisé pour la résoudre.
L'unicité de la réalisation et la convention non terminale impliquent
\[
 \rho_{N+1,N}\operatorname{Dec}_{N+1}(\chi)
 =\operatorname{Dec}_N(\chi).
\]
Les identités de Machin, les sommes partielles de la série exponentielle et les
quotients consécutifs de Fibonacci--Perron fournissent des intervalles rationnels
certifiés pour calculer \(\operatorname{Dec}_N\) ; ils n'interviennent pas dans
\(\operatorname{Ext}\), ne choisissent ni orbite ni caractère et ne sélectionnent pas
d'enfants de ses fibres locales. Ils calculent bien le bloc de coordonnées de publication,
qui demeure typé comme sortie ultérieure. Les
cylindres publiés sont emboîtés et leurs diamètres tendent vers zéro, de sorte
que leur limite inverse détermine un unique point. Comme \(N\) est arbitraire et
que la relation structurelle d'extension est dépourvue de profondeur terminale, la
conclusion vaut simultanément pour toute longueur finie et pour le développement
infini.
\end{proof}

Le porteur exécutable impose matériellement cette séparation au moyen de trois
fonctions disjointes : il engendre d'abord les caractères exacts, puis les évalue
et ne publie qu'à la fin les blocs archimédiens. Son propre audit syntaxique
échoue si un évaluateur entre dans le générateur. L'exécution de \(396\) niveaux,
\(2376\) trits, \(43\) tours, \(387\) couples \((K,K+9)\) par canal et
\(1000\) chiffres vérifie une vaste instance matérielle du théorème ; celle de
\(10,000\) chiffres en vérifie une autre. Aucune de ces deux quantités ne limite la
profondeur. Le quantificateur \(\forall N\) et la compatibilité
\(\rho_{N+1,N}p_{N+1}=p_N\), avec l'existence de l'élément
\(p_\infty\) de la limite inverse, démontrent la sortie décimale infinie ; les
calculs finis sont des certificats reproductibles de ses projections.

La coordonnée \(\alpha_{\mathrm{HMT}}\) appartient à la même famille limite
corrélée. Sa voie entière se prolonge au moyen de
\(\mathsf T_{\alpha,\infty}\). Le lecteur
analytique interne construit exactement \(E_{21/22}^{(9)}\) et sa racine simple
\(\xi_9\) dans le quotient des jets d'ordre neuf. La stabilité vérifiée à
\(10,000\) chiffres certifie la propagation distante de la retenue de la voie
entière, mais n'est qu'une projection finie. Les carrés de troncature de
la famille dodécaphasique et de la famille de jets commutent exactement ;
l'unicité de leurs sections compatibles identifie
\(\alpha_A=\alpha_B=\alpha_{\mathrm{HMT}}\) dans
\(\mathbb R_{\rm HMT}\) avant toute comparaison métrologique. Les exécutions à
neuf, mille ou dix mille chiffres sont des témoins finis de cette identité et non sa
portée.

\begin{proposition}[Prolongement complet de la voie entière de \(\alpha\)]
\label{prop:c27-alpha-prolongacion-completa}
Soit \(B=1000\), soient \((P_k),(E_k),(\Phi_k)\) les trois suites de blocs
publiées par l'état coinductif et soit \((K_k)\) la suite
dodécaphasique prolongée. Les quatre séries
\[
 p=\sum_{k\geq1}P_kB^{-k},\qquad
 e_0=\sum_{k\geq1}E_kB^{-k},\qquad
 \phi_0=\sum_{k\geq1}\Phi_kB^{-k},\qquad
 \kappa=\sum_{k\geq1}K_kB^{-k}
\]
convergent absolument et déterminent
\[
 \alpha_A=p+e_0-\phi_0-\kappa.
\]
La récurrence de retenue est l'écriture en base \(B\) de cette égalité.
Pour chaque \(M\), une troncature suffisamment profonde détermine les
premiers \(M\) chiffres du développement canonique de \(\alpha_A\) ; l'erreur de
queue est bornée par une constante multipliée par \(B^{-N}\). Par
conséquent, \(\mathsf T_{\alpha,\infty}\) publie une suite décimale
complète, et ses exécutions à \(3000\) ou \(10\,000\) chiffres sont des projections
finies de cette sortie.
\end{proposition}

\begin{proof}
Chaque bloc appartient à un ensemble entier borné, de sorte que les quatre
séries sont dominées par une série géométrique. La somme partielle jusqu'à \(N\)
diffère de la somme complète de \(O(B^{-N})\). La division euclidienne de
droite à gauche transporte exactement cette queue au moyen de la retenue ;
par unicité du développement canonique, tout préfixe qui reste séparé d'un
bord de cylindre se stabilise lorsque \(N\) augmente, et sur un bord
on applique la convention canonique non terminale. Cela définit simultanément
tous les chiffres, non une suite d'objectifs numériques indépendants.
\end{proof}

\begin{corollary}[Séparation entre génération et reconnaissance]
\label{cor:c27-generacion-reconocimiento}
Les identités de Machin, la série exponentielle, le rayon de Perron et la carte
décimale sont des réalisations exactes de caractères préalablement sélectionnés
par \(\mathfrak H_9\). Les noms conventionnels \(\pi,e,\varphi,\alpha\)
et les tables métrologiques reconnaissent les sorties ; ils ne déterminent ni orbites,
ni caractères, ni transports, ni retenues. Les préfixes qu'ils calculent sont uniquement
des projections coordonnées ultérieures des caractères déjà générés.
\end{corollary}

\paragraph{Évaluation des caractères construits.}
À l'étape d'évaluation archimédienne, le caractère de propagation
détermine la condition initiale \(a_0=1\) et la récurrence
\((n+1)a_{n+1}=a_n\). Le calcul des termes applique cette récurrence.
Pour l'auto-échelle, le vecteur ligne évolue au moyen de l'incidence
préalablement construite,
\[
 (x_{n+1},y_{n+1})=(x_n,y_n)F_{\rm av},
 \qquad
 F_{\rm av}=\begin{pmatrix}0&1\\1&1\end{pmatrix},
 \qquad
 \det(tI-F_{\rm av})=t^2-t-1.
\]
Les quotients de composantes fournissent les bornes rationnelles du
lecteur d'auto-échelle. L'implémentation vérifie la concordance entre
incidence, polynôme caractéristique, orientation et normalisation avant
d'effectuer cette évaluation. Une altération ou une absence de ces données
interrompt le calcul au lieu de conserver tacitement la valeur
précédente. Ce contrôle relie les caractères déjà produits à leurs
évaluateurs ; sa place dans la construction est postérieure au générateur.

% El inventario completo, el registro integral y la carta firmada son los
% propietarios upstream ya demostrados. Se imprimen aquí, sin reescritura,
% antes de cualquier publicación decimal o reconocimiento convencional.
\begin{HMTScientificOwnerSubordinated}
% Unidad U006F. Inventario operatorio completo del Topological Phase Kernel.

\chapter{Inventaire opératoriel complet et chaîne causale du Topological Phase Kernel}
\label{chap:inventario-operatorio-completo-tpk}

Le Topological Phase Kernel n’est ni une réduction modulo trois ni une horloge à
cent huit positions. C’est une catégorie d’états enrichis équipée
d’opérateurs de sélection, de transport, de lecture, de mémoire, de périodicité,
de prolongement et de publication. Pour empêcher une abréviation de remplacer
l’objet, chaque opérateur est enregistré selon son type et l’information qu’il
conserve.

\section{Fiche typée d'un opérateur}

\begin{definition}[Fiche opératoire]
\label{def:u006f-ficha-operatoria}
Une réalisation canonique du TPK est un sextuplet
\begin{equation}
 \mathcal O=
 (D_{\mathcal O},C_{\mathcal O},f_{\mathcal O},
 I_{\mathcal O},L_{\mathcal O},M_{\mathcal O}),
 \label{eq:u006f-ficha}
\end{equation}
où :
\begin{enumerate}
 \item \(f_{\mathcal O}:D_{\mathcal O}\to C_{\mathcal O}\) est sa règle ;
 \item \(I_{\mathcal O}\) est la famille des relations conservées ;
 \item \(L_{\mathcal O}\) déclare l’information publiée ou perdue ;
 \item \(M_{\mathcal O}\) énumère la mémoire qui demeure dans l’état
 enrichi.
\end{enumerate}
\end{definition}

Deux réalisations ne représentent le même opérateur que s’il existe une
équivalence typée qui entrelace les règles et les invariants. Partager une période,
une cardinalité ou une sortie visible ne suffit pas.

\section{Huit classes structurelles}

L’inventaire est partitionné en :
\begin{enumerate}
 \item[\textbf{C1.}] supports et états ;
 \item[\textbf{C2.}] sélection interne ;
 \item[\textbf{C3.}] transporte;
 \item[\textbf{C4.}] lecteurs et quotients ;
 \item[\textbf{C5.}] mémoire et frontière ;
 \item[\textbf{C6.}] périodicité et retour ;
 \item[\textbf{C7.}] prolongement ;
 \item[\textbf{C8.}] salidas.
\end{enumerate}
La classe est déterminée par le domaine, le codomaine et la fonction, non par le nom
d’un chiffre associé.

\section{Quatorze regroupements fonctionnels}

\begingroup
\small
\renewcommand{\arraystretch}{1.12}
\begin{longtable}{@{}p{.20\textwidth}p{.29\textwidth}p{.43\textwidth}@{}}
\toprule
Regroupement&Opérateurs ou espaces&Fonction et invariants\\
\midrule
\endfirsthead
\toprule
Regroupement&Opérateurs ou espaces&Fonction et invariants\\
\midrule
\endhead
État et mémoire de route&
\(\mathcal X_{\rm TPK}^{\rm enr}\), \(\mathcal F_q\)&
Ils conservent la position, le feuillet, l’orientation, le résidu, le quotient, la retenue, la signature,
la frontière, le cylindre, le préfixe, la provenance et la mémoire.\\

Sélection interne&
\(M_{\rm ph}:\{1,\ldots,9\}\to\{-1,0,+1\}\)&
Active l’évaluation additive, multiplicative ou le pas de seuil ; ne déplace pas
l’état.\\

Transporte espacial&
\(T_N,T_E,T_S,T_O,F_{\rm loc},F_{\rm obs}\)&
Réalise des translations cardinales, une rétroaction de feuillet et une chronologie à deux
curseurs.\\

Lectores modulares&
\(\rho_9,q_9,\rho_3^{\rm or},c_3,q_{1000},r_{1000}\)&
Publient la phase, l’orientation et le bloc décimal en conservant les quotients de
relèvement.\\

Bloc et mémoire&
\(\mathcal K_{27}=F_{\rm obs}^{27}\), \(\mathcal N_{27}\)&
Le premier compose vingt-sept transports ; le second filtre des événements de
mémoire. Un indice identique n’implique pas un opérateur identique.\\

État dodécaphasique&
\(\mathcal R_{12},C_{12},\Theta_{108},\Gamma_{108}\)&
Distingue le registre enrichi, le secteur, la loi cyclique et le support ordonné.\\

Périodicité et profondeurs&
\(C_{12}\hookrightarrow C_{108}\twoheadrightarrow C_9\),
\(w_{108},w_{1080}\)&
Conserve l’extension non scindée, les mots de longueurs différentes et
la mémoire verticale.\\

Canaux \(90/120\)&
\(\mathcal F_6,\mathcal F_8,\mathcal I_{6\to8},
S_{90},S_{120}\)&
Sépare les incidences finies, les queues analytiques et leurs normalisations.\\

Bidirectionnalité&
\(P_\pm,\operatorname{rev},\operatorname{Ext}_\pm,
N_\pm,\beta,C_M\)&
Réalise l’échange, l’inversion de route, le prolongement futur et passé,
la valence, le bilan et la conjugaison.\\

Émission et coinduction&
\(L_0,L_1,R_{36},G_9,\varprojlim\operatorname{Cyl}_m\)&
Relève le germe, ouvre la frontière et organise des prolongements finis et
compatibles.\\

Lecteurs numériques&
\(\mathscr C_\pi,\mathscr P_e,F_{\rm av}\)&
Construisent des encadrements rationnels, une propagation factorielle et une auto-échelle de
Fibonacci.\\

Projection exceptionnelle&
\(\Gamma_{A_W}\), poids, support projectif&
Transporte la même émission vers une incidence ternaire sans sélectionner de chiffres
archimédiens.\\

Atlas et réalisation&
signature de route, atlas \(81/56/13/324\), \(C_M\)&
Classe les histoires et les caractères avant d’appliquer une carte dimensionnelle.\\

Validation finie&
recensements, ablations, empreintes de hachage et preuves d'indépendance&
Vérifie la couverture à la profondeur exécutée et distingue générateur, vérificateur et données de comparaison.\\
\bottomrule
\end{longtable}
\endgroup

Les huit classes décrivent la nature mathématique ; les quatorze regroupements,
la fonction d’exécution. Ce sont deux gradations du même inventaire.

\section{Registre nominal de trente-quatre entrées}

\begin{equation}
 \mathfrak O_{\rm TPK}:=(\mathfrak o_1,\ldots,\mathfrak o_{34}).
 \label{eq:u006f-registro-operadores}
\end{equation}

\begingroup
\scriptsize
\renewcommand{\arraystretch}{1.12}
\begin{longtable}{@{}r p{.19\textwidth}p{.28\textwidth}p{.42\textwidth}@{}}
\toprule
\#&Symbole&Domaine et codomaine&Règle ou fonction\\
\midrule
\endfirsthead
\toprule
\#&Symbole&Domaine et codomaine&Règle ou fonction\\
\midrule
\endhead
1&\(\mathcal A_9\)&\((\mathbb Z/9)^2\), matrice et loi interne&
Cellule arithmétique additive--multiplicative.\\
2&\((\rho_9,q_9)\)&\(\mathbb Z_{>0}\to
\{1,\ldots,9\}\times\mathbb Z_{\geq0}\)&
Division nonadique avec représentant positif et quotient.\\
3&\(\rho_3^{\rm or},c_3\)&\(\mathbb Z\to
\{-1,0,+1\}\times\mathbb Z\)&
Relèvement ternaire orienté \(n=\rho_3^{\rm or}(n)+3c_3(n)\).\\
4&\(\iota_{1000}:=(q_{1000},r_{1000})\)&\(\mathbb Z\to
\mathbb Z\times\{0,\ldots,999\}\)&
Division euclidienne \(N=1000q_{1000}+r_{1000}\).\\
5&\(T_d\)&\(X_9\to X_9\), \(d\in\{N,E,S,O\}\)&
Translation cardinale orientée.\\
6&\(F_{\rm loc}\)&\(\mathcal Q_{\rm loc}\times
\mathscr D_{\rm card}\to\mathcal Q_{\rm loc}\)&
Rétroaction locale de position et de feuillet.\\
7&\(T_{\min}\)&\(\Omega_{\rm seed}\to\mathcal U_6\)&
Émetteur minimal des deux curseurs, y compris la signature coordonnée.\\
8&\(G\)&\(\operatorname{im}s_\Sigma\times
\operatorname{im}s_\Pi\to\{0,\ldots,999\}^6\)&
Couplage décimal coordonnée par coordonnée.\\
9&\(\Psi_6\)&\(\mathcal U_6\to\mathbb F_3^6\)&
Réduction ternaire du catalogue.\\
10&\(\mathcal R_3\)&\(\mathcal Q_{\rm loc}\times
\mathscr D_{\rm card}^3\to\{-1,0,+1\}\)&
Facteur observable local de trois pas.\\
11&\(\mathcal K_{27}\)&\(\mathcal Q_{\rm obs}\to\mathcal Q_{\rm obs}\)&
Composition \(F_{\rm obs}^{27}\).\\
12&\(\mathcal X_{\rm TPK}^{\rm enr}\)&catégorie des états enrichis&
Conserve toutes les coordonnées typées du noyau.\\
13&\(\mathcal H\)&relation sur les états enrichis&
Transition de Hensel qui retient le résidu et le quotient.\\
14&\(\mathcal C_{3,1000}\)&états croisés
\((N,L,D,K,A,B)\)&
Actualisation exacte des cylindres ternaires--décimaux.\\
15&\(\Sigma_{\rm B}\)&mémoires de frontière \(\to\) signature conjointe&
Exprime l'orientation du bloc en conservant son registre.\\
16&\({\rm EF}\text{--}G_9\)&relation d’extension&
Filtre les extensions compatibles et organise la généalogie nonadique.\\
17&\(\Gamma_9\)&histoires enrichies \(\longrightarrow\) histoires&
Holonomie d’un tour de phase avec mémoire de frontière.\\
18&\(X_\chi\)&\(\varprojlim_m\operatorname{Surv}^{(\chi)}_m\)&
Section globale lorsque la non-vacuité, la compatibilité et l’unicité sont
démontrées à toute profondeur.\\
19&\(\mathcal P_{\rm APP}^{(2)}\)&catégorie bivariée de chemins&
Transporte simultanément l’évaluation multiplicative et l’accumulation
additive.\\
20&\(F_{\rm obs}\)&\(\mathcal Q_{\rm obs}\to\mathcal Q_{\rm obs}\)&
Chronologie observable déterministe à deux curseurs.\\
21&\(\mathcal A_{42}\)&alphabet de quarante-deux triades&
Support décimal élargi de l’émission finie.\\
22&\(\operatorname{Ext}_r\)&
\(\operatorname{Ext}_r\subseteq\mathcal F_q\times\mathcal F_{q'}\)&
Relation typée de prolongement des fibres.\\
23&\(\mathcal K_{\rm ph}\)&\(\mathcal U_6\times C_9\) sur lui-même&
Transfert de continuations compatibles avec phase explicite.\\
24&\((P,H,Q)\)&\(\mathcal X_\times\to\mathcal X_\times\)&
Avancée, demi-tour et quart de tour des germes.\\
25&\((c_{\rm mode},c_{\rm or})\)&caminos \(\to\mathbb Z^2\)&
Cocycles de changement de mode et d’orientation.\\
26&\(E_{20}\)&états finis \(\to\) vingt blocs&
Prolongement triadique fini avec livre de provenance.\\
27&\(\mathcal R_{12}\)&histoire enrichie \(\to\) registre de douze
secteurs&
Système temporel orienté avec parcours, signature, retenue et registre.\\
28&\(R_{\rm sgn}\)&état terminal \(\to U_{12}^{\rm sgn}
\subseteq\mathbb Z^{12}\)&
Lecteur de l’observable harmonique signée.\\
29&\(\mathcal H_{12}\)&\(\mathbb Z^{12}\to\mathbb Z^{12}\)&
Transformation par trois orbites de pas trois ; inverse
\(\mathcal H_{12}/4\) sur l’image pertinente.\\
30&\(\mathcal E_{\rm dod}=(D_3,D_4,Q)\)&
\(\mathbb Z^{12}\to\mathbb Z^{25}\)&
Extracteur transversal conjoint d’arités trois et quatre.\\
31&\(\mathcal C_\alpha\)&blocs et retenues \(\to\) blocs&
Récurrence dodécaphasique de clôture en base mille.\\
32&\((L_{\rm tick},\chi_{\rm mass})\)&histoire enrichie \(\to\)
caractère temporel&
Interface dimensionnelle ultérieure ; elle ne participe pas au constructeur numérique de
cet ouvrage.\\
33&\(\mathcal W_\pm\)&réseau bidirectionnel d’indice douze&
Lecture dans les deux sens en conservant l’orientation et la provenance.\\
34&\(\mathcal A_{\rm species}\)&familles d’extension \(\to\) atlas&
Interface de réalisation ultérieure, non un quatrième primitif.\\
\bottomrule
\end{longtable}
\endgroup

\section{Chaîne causale d'actualisation}

\begin{definition}[État canonique complet]
\label{def:u006f-esquema-estado}
Soit
\[
 \mathcal E:=
 \mathbb Z_{\rm quo}\times\mathbb Z_{\rm car}\times
 \mathsf{Mem}\times\mathsf{Term}\times\mathsf{Cyl}\times
 \mathsf{Fr}\times\mathsf{Pref}\times\Lambda .
\]
L’état enrichi canonique à l’instant \(t\) est
\begin{equation}
 x_t=(p_t,m_t,\sigma_t,b_t,r_t;
 q_t,c_t,\eta_t,s_t,C_t,\partial_t,N_t,\lambda_t)
 \in
 X_9\times\{\Sigma,\Pi\}\times\mathsf{TRIT}\times
 C_{12}\times C_9\times\mathcal E.
 \label{eq:u006f-esquema-estado}
\end{equation}
Les abréviations d’état utilisées dans d’autres unités sont des projections de
ce tuple. En particulier, ni le quotient \(q_t\), ni la retenue
\(c_t\), ni le terminal \(s_t\), ni le cylindre \(C_t\), ni la frontière
\(\partial_t\), ni le préfixe \(N_t\), ni le livre \(\lambda_t\) ne sont omis.
\end{definition}

Soit \(\rho_t\) la flèche observable déterminée par \(x_t\). Les lois de transport de U005 et la transition \(\mathcal C_{3,1000}\) déterminent un endomorphisme de mémoire. Pour \(z=(q,c,\eta,s,C,\partial,N,\lambda)\in\mathcal E\), on définit
\begin{align}
 \mathbf E_t:\mathcal E&\longrightarrow\mathcal E,\notag\\
 \mathbf E_t(z)&=\bigl(
 \mathsf H(\rho_t)q,
 \mathsf C(\rho_t)c+\kappa(\rho_t),
 \mathsf M(\rho_t)\eta,
 \mathsf T(\rho_t)s,
 \mathsf{Cyl}(\rho_t)C,
 \mathsf{Fr}(\rho_t)\partial,
 \mathsf{Pref}(\rho_t)N,
 \Lambda(\rho_t)\lambda
 \bigr).
 \label{eq:u006f-actualizacion-memoria}
\end{align}
Chaque composante publiée dans \eqref{eq:u006f-actualizacion-memoria} possède le
domaine indiqué par sa fibre ; \(\kappa(\rho_t)\) est la cochaîne de retenue.

Nous définissons trois endomorphismes du produit canonique. Si
\(\sigma'=M_{\rm ph}(g(t))\) et \(m'=\mu(\sigma',m)\), alors
\begin{align}
 \mathsf{Sel}_t(p,m,\sigma,b,r;e)
  &:=(p,m',\sigma',b,r;e),\label{eq:u006f-selector-tipado}\\
 \mathsf{Tra}_t(p,m,\sigma,b,r;e)
  &:=(T_{d(t)}p,m,\sigma,b,r;e),\label{eq:u006f-transporte-tipado}\\
 \mathsf{Upd}_t(p,m,\sigma,b,r;e)
  &:=(p,m,\sigma,b',r';\mathbf E_t(e)),
 \label{eq:u006f-memoria-tipada}
\end{align}
où \((b',r')\) est la cible de \(\rho_t:(b,r)\to(b',r')\). La chaîne causale est une composition d'applications, non la composition d'une application avec une valeur tritique :
\begin{equation}
 \boxed{\mathcal U_t:=
 \mathsf{Upd}_t\circ\mathsf{Tra}_t\circ\mathsf{Sel}_t.}
 \label{eq:u006f-tuberia}
\end{equation}

\begin{theorem}[Clôture de la chaîne]
\label{thm:u006f-cierre-tuberia}
L’application \(\mathcal U_t\) est un endomorphisme de l’état canonique
complet. Les lecteurs, prolongements et projections sont appliqués après
\(\mathcal U_t\) uniquement lorsque leur domaine a été produit.
\end{theorem}

\begin{proof}
\(\mathsf{Sel}_t\) ne change que le mode et le TRIT ;
\(\mathsf{Tra}_t\) ne change que la position APP ; et
\(\mathsf{Upd}_t\) applique le produit typé
\eqref{eq:u006f-actualizacion-memoria} et la flèche de phase. Les trois
codomaines sont le même produit de la
\cref{def:u006f-esquema-estado} ; leur composition est donc définie et
conserve toutes les coordonnées de l’état.
\end{proof}

\section{Règles de non-identification}

L’inventaire impose
\begin{equation}
 \mathcal S_4\neq D_4\neq\operatorname{sd}_4,
 \qquad
 \mathcal R_{12}\neq C_{12}\neq\Theta_{108}\neq\Gamma_{108},
 \qquad
 F_{\rm obs}\neq\mathcal K_{\rm ph},
 \label{eq:u006f-no-identificacion-uno}
\end{equation}
et
\begin{equation}
 U_{12}^{\rm cat}\neq U_{12}^{\rm sgn},
 \qquad
 \mathcal E_{\rm dod}\neq D_{\rm raw}^{90,120}
 \neq\mathcal K_{6\to8}.
 \label{eq:u006f-no-identificacion-dos}
\end{equation}
Ce sont des incompatibilités de type, non de simples inégalités de valeur.

\begin{criterion}[Audit de complétude opératorielle]
\label{crit:u006f-completitud}
Une exécution n’est opératoriellement complète que si :
\begin{enumerate}
 \item elle déclare les entrées de \(\mathfrak O_{\rm TPK}\) employées ;
 \item elle enregistre le domaine et le codomaine de chaque composition ;
 \item elle conserve l’ordre causal de \eqref{eq:u006f-tuberia} ;
 \item elle distingue dynamique, transfert, lecture, prolongement et sortie ;
 \item elle permet de reconstruire à partir de \(\lambda_t\) chaque changement de bloc et
 chaque retenue ;
 \item elle ne promeut pas une vérification finie en une section infinie sans
 prouver la non-vacuité, la compatibilité et l’unicité à toute profondeur.
\end{enumerate}
Tout symbole utilisé sans type ou toute fusion interdite par
\eqref{eq:u006f-no-identificacion-uno} et
\eqref{eq:u006f-no-identificacion-dos} constitue un échec vérifiable.
\end{criterion}

% Unidad U006D. Registro dodecafásico integral y descomposición 1+5+6.

\chapter{Registre dodécaphasique intégral, extension cyclique et décomposition \(1+5+6\)}
\label{chap:registro-dodecafasico-integral}

La période observable de cent huit positions contient douze secteurs de
neuf positions. Cette égalité de cardinalités ne suffit pas à décrire
l’état temporel : il faut distinguer le support ordonné, la loi cyclique,
l’indice sectoriel et le registre enrichi des transports, orientations,
retenues et incidences. Cette unité construit les quatre objets et démontre
la reconstruction entière des douze coordonnées de mémoire.

\section{Support ordonné et coordonnée cyclique}

Pour \(m\in\{1,\ldots,12\}\), nous définissons
\begin{equation}
 E_m:=\{9(m-1)+1,\ldots,9m\}.
 \label{eq:u006d-bloques-nonadicos}
\end{equation}
Le support ordonné et le groupe cyclique sont
\begin{equation}
 \Gamma_{108}:=
 E_1\sqcup\cdots\sqcup E_{12},
 \qquad
 \Theta_{108}:=\mathbb Z/108\mathbb Z.
 \label{eq:u006d-gamma-theta}
\end{equation}
\(\Gamma_{108}\) conserve la position, l’ordre et l’appartenance à un bloc ;
\(\Theta_{108}\) conserve la loi de translation. Ce ne sont pas le même objet.

Pour \(\theta\in\Theta_{108}\), soit
\(\widetilde\theta\in\{0,\ldots,107\}\) son représentant. Nous définissons
\begin{equation}
 \beta(\theta):=
 \left[
 \left\lfloor\frac{\widetilde\theta}{9}\right\rfloor
 \right]_{12},
 \qquad
 r(\theta):=1+(\widetilde\theta\bmod9).
 \label{eq:u006d-coordenadas-bloque-residuo}
\end{equation}
La première coordonnée localise le secteur dodécaphasique ; la seconde publie
la position positive \(1,\ldots,9\) à l’intérieur du secteur.

\section{Extension cyclique non scindée}

Soient \(C_n=\mathbb Z/n\mathbb Z\). Les applications
\begin{equation}
 \iota:C_{12}\longrightarrow C_{108},
 \quad \iota([b]_{12})=[9b]_{108},
 \qquad
 p:C_{108}\longrightarrow C_9,
 \quad p([\theta]_{108})=[\theta]_9
 \label{eq:u006d-mapas-extension}
\end{equation}
forment la suite
\begin{equation}
 0\longrightarrow C_{12}
 \xrightarrow{\;\iota\;}C_{108}
 \xrightarrow{\;p\;}C_9
 \longrightarrow0.
 \label{eq:u006d-extension}
\end{equation}
Pour la section ensembliste
\[
 s([r]_9):=[\widetilde r]_{108},
 \qquad 0\leq\widetilde r<9,
\]
le défaut d’additivité est
\begin{equation}
 c(r,r'):=
 \left[
 \left\lfloor
 \frac{\widetilde r+\widetilde r'}9
 \right\rfloor
 \right]_{12},
 \qquad
 s(r)+s(r')
 =s(r+r')+\iota(c(r,r')).
 \label{eq:u006d-cociclo-extension}
\end{equation}

\begin{theorem}[Exactitude, non-scindement et classe cohomologique]
\label{thm:u006d-extension-no-escindida}
La suite \eqref{eq:u006d-extension} est exacte et ne se scinde pas. Pour
l’action triviale,
\begin{equation}
 H^2(C_9,C_{12})
 \simeq C_{12}/9C_{12}
 \simeq C_3,
 \label{eq:u006d-cohomologia-extension}
\end{equation}
et le cocycle \eqref{eq:u006d-cociclo-extension} représente sa classe
génératrice.
\end{theorem}

\begin{proof}
L’image de \(\iota\) est le sous-groupe des multiples de neuf, qui coïncide
avec \(\ker p\) ; \(p\) est surjective. Cela prouve l’exactitude. Si
l’extension se scindait, \(C_{108}\) serait isomorphe à
\(C_{12}\times C_9\), mais
\[
 \exp(C_{108})=108,
 \qquad
 \exp(C_{12}\times C_9)=\operatorname{mcm}(12,9)=36.
\]
L’identité de cocycle est la division euclidienne de
\(\widetilde r+\widetilde r'\) par neuf. Pour un groupe cyclique à action
triviale,
\(H^2(C_9,C_{12})\simeq C_{12}/9C_{12}\) ; la retenue unitaire se projette sur
\([1]\), qui engendre le quotient d’ordre trois.
\end{proof}

Le non-scindement type le retour :
\[
 g_0(\theta+9)=g_0(\theta),
 \qquad
 \beta(\theta+9)=\beta(\theta)+1\pmod{12}.
\]
Neuf itérations ferment la phase visible et translatent d’un secteur ; cent
huit ferment les deux coordonnées de l’horloge. L’état enrichi conserve
en outre la mémoire verticale, le cylindre, la frontière et la généalogie.

\section{Registre temporel enrichi}

\begin{definition}[Registre dodécaphasique orienté]
\label{def:u006d-registro-dodecafasico}
Pour une histoire enrichie, soit \(\tau_m\) la sous-route ordonnée sur
\(E_m\), \(\sigma_m\) son orientation, \(\kappa_m\) la retenue qui franchit la
frontière du bloc et \(\Lambda_m\) son livre local d’incidences. Le
registre complet est
\begin{equation}
 \mathcal R_{12}:=
 \bigl(
 (E_m,\tau_m,\sigma_m,\kappa_m,\Lambda_m)
 \bigr)_{m=1}^{12}.
 \label{eq:u006d-registro-r12}
\end{equation}
\end{definition}

On distingue ainsi quatre niveaux :
\begin{equation}
 \boxed{
 \mathcal R_{12}
 \neq C_{12}
 \neq\Theta_{108}
 \neq\Gamma_{108}.}
 \label{eq:u006d-cuatro-objetos}
\end{equation}
Le premier est un registre enrichi ; le deuxième, un indice de secteur ; le
troisième, un groupe de translations ; le quatrième, un support ordonné. Il existe
des projections entre eux, mais pas d’identifications.

Soit \(q_{\rm mem}\in\mathbb Z_{\geq0}\) le nombre non borné de tours
nonadiques. La coordonnée dodécaphasique n’est que son résidu
\[
 \beta=[q_{\rm mem}]_{12}.
\]
Après cent huit pas,
\[
 \beta\longmapsto\beta,
 \qquad
 q_{\rm mem}\longmapsto q_{\rm mem}+12.
\]
Le calendrier revient ; la mémoire verticale n’est pas réinitialisée.

\section{Extracteur signé d’événements}

Soit \((d_t)_{t=1}^{108}\) le registre de chiffres sur
\(\Gamma_{108}\). Nous définissons
\begin{equation}
 \mathcal D_{\rm evt}:=\{2,4,6,8\},
 \qquad
 \Sigma_{12}:=(+,+,+,-,-,-,+,+,+,-,-,-).
 \label{eq:u006d-datos-eventos}
\end{equation}
Pour \(m=0,\ldots,11\), soit
\begin{equation}
 e_m:=
 \Sigma_{12}(m+1)\,
 \#\{t\in E_{m+1}:d_t\in\mathcal D_{\rm evt}\}
 \in\mathbb Z.
 \label{eq:u006d-eventos-firmados}
\end{equation}
L’application conserve le recensement local et l’orientation sectorielle comme
facteurs séparés.

Nous apparions les deux demi-couronnes par
\begin{equation}
 p_i:=e_i+e_{i+6},
 \qquad
 a_i:=e_i-e_{i+6},
 \qquad 0\leq i\leq5.
 \label{eq:u006d-pares-antipodales}
\end{equation}
Nous définissons le recensement total, les cinq différences par rapport à la sixième paire et
les six antisymétries :
\begin{equation}
 s_C:=\sum_{i=0}^{5}p_i,
 \qquad
 M_i:=p_i-p_5\quad(0\leq i\leq4),
 \qquad
 A_6:=(a_0,\ldots,a_5).
 \label{eq:u006d-coordenadas-uno-cinco-seis}
\end{equation}

\begin{definition}[Lecteur entier \(1+5+6\)]
\label{def:u006d-lector-integral}
Le lecteur linéaire est
\begin{equation}
 \mathcal L_{12}:\mathbb Z^{12}\longrightarrow\mathbb Z^{12},
 \qquad
 \mathcal L_{12}(e)
 :=
 (s_C,M_0,\ldots,M_4,a_0,\ldots,a_5).
 \label{eq:u006d-lector-l12}
\end{equation}
La distribution des coordonnées est
\begin{equation}
 \underbrace{1}_{s_C}
 +\underbrace{5}_{M_0,\ldots,M_4}
 +\underbrace{6}_{a_0,\ldots,a_5}
 =12.
 \label{eq:u006d-uno-cinco-seis}
\end{equation}
\end{definition}

\begin{theorem}[Reconstruction exacte et image entière]
\label{thm:u006d-inversa-l12}
Le lecteur \(\mathcal L_{12}\) est injectif. Sur son image,
\begin{align}
 p_5&=
 \frac{s_C-\sum_{i=0}^{4}M_i}{6},
 &
 p_i&=p_5+M_i\quad(0\leq i\leq4),
 \label{eq:u006d-inversa-p}\\
 e_i&=\frac{p_i+a_i}{2},
 &
 e_{i+6}&=\frac{p_i-a_i}{2}
 \quad(0\leq i\leq5).
 \label{eq:u006d-inversa-e}
\end{align}
Un tuple de sortie appartient à \(\operatorname{im}\mathcal L_{12}\) si et
seulement si
\begin{equation}
 s_C-\sum_{i=0}^{4}M_i\equiv0\pmod6
 \label{eq:u006d-condicion-divisibilidad}
\end{equation}
et, pour les \(p_i\) reconstruits,
\begin{equation}
 p_i\equiv a_i\pmod2
 \qquad(0\leq i\leq5).
 \label{eq:u006d-condicion-paridad}
\end{equation}
\end{theorem}

\begin{proof}
De \(M_i=p_i-p_5\), on obtient
\[
 s_C=\sum_{i=0}^{4}(p_5+M_i)+p_5
 =6p_5+\sum_{i=0}^{4}M_i,
\]
ce qui donne la première équation de \eqref{eq:u006d-inversa-p} ; les autres sont
immédiates. Additionner et soustraire
\(p_i=e_i+e_{i+6}\) et \(a_i=e_i-e_{i+6}\) produit
\eqref{eq:u006d-inversa-e}. La divisibilité par six et les six
congruences de parité sont précisément les conditions nécessaires et
suffisantes pour que toutes les coordonnées reconstruites soient entières.
\end{proof}

La division par six apparaît après le recensement orienté, non pendant celui-ci. Les
conditions d’intégralité de cette unité sont exactement
\eqref{eq:u006d-condicion-divisibilidad} et
\eqref{eq:u006d-condicion-paridad} ; on n’attribue pas ici de forme normale de
Smith à une matrice qui n’aurait pas été préalablement définie.

\section{Deux mots dodécaphasiques de provenances distinctes}

La distinction doit être conservée
\begin{equation}
 U_{12}^{\rm cat}:=U_6\Vert U_6,
 \qquad
 U_{12}^{\rm sgn}:=\mathcal H_{12}(K).
 \label{eq:u006d-dos-palabras-doce}
\end{equation}
La première concatène deux émissions du catalogue de longueur six. La seconde
est une coordonnée signée reliée réversiblement à \(K\) par l’opérateur
\(\mathcal H_{12}\). Leur longueur commune n’implique pas une même provenance,
sémantique ou loi de transport.

\begin{criterion}[Critères de réfutation]
\label{crit:u006d-refutacion}
La construction est réfutée si la suite
\eqref{eq:u006d-extension} n’est pas exacte, si sa classe cohomologique est
triviale, si un tuple de l’image viole
\eqref{eq:u006d-condicion-divisibilidad} ou
\eqref{eq:u006d-condicion-paridad}, si les formules inverses ne retrouvent pas
les douze entiers initiaux, ou si le retour de \(\beta\) est interprété comme
une réinitialisation de \(q_{\rm mem}\).
\end{criterion}

% Unidad U016. Registro firmado, lectura terminal y reconstruccion de K.
% Redaccion autonoma desde la autoridad material 1063 y su sucesor V2.

\chapter{Registre dodécaphasique signé et reconstruction intégrale du sceau}
\label{chap:registro-dodecafasico-hadamard-k}

La construction du sceau entier ne commence ni par \(\alpha\), ni par son inverse
ni par un tableau métrologique. Son domaine est l’état terminal enrichi
du Topological Phase Kernel obtenu après la chaîne déjà construite
\begin{equation}
 \boxed{
 \mathrm{APP}\longrightarrow\mathrm{TRIT}\longrightarrow\mathrm{TPK}
 \longrightarrow\Omega_{\rm seed}
 \xrightarrow{T_{\min}}\mathcal U_6^{\rm cat}
 \xrightarrow{\rho_3}W_{243}
 \longrightarrow w_{30}\longrightarrow R_{36}
 \longrightarrow G_9\longrightarrow x_{\rm term}.}
 \label{eq:cadena-upstream-registro-dodecafasico}
\end{equation}
Les flèches précédentes ont été définies avec leurs domaines respectifs :
paires de germes, signatures, émissions, mots ternaires, frontières,
histoires compatibles et holonomie nonadique. Cette unité ne choisit à nouveau
aucune de ces données. Elle construit les cartes entières du même état
terminal et démontre leur équivalence.

\section{Fenêtres de la trace orientée de longueur cent huit}

Soit
\[
 \Theta_{108}=\mathbb Z/108\mathbb Z
\]
le calendrier observable. Ses douze fenêtres nonadiques sont
\[
 E_m=\{9m+1,\ldots,9m+9\},
 \qquad 0\leq m<12.
\]
La partition est disjointe et exhaustive :
\[
 \{1,\ldots,108\}=\bigsqcup_{m=0}^{11}E_m.
\]
L’ensemble des codes de direction qui publient un événement est fixé comme
\[
 D_{\rm evt}=\{2,4,6,8\}.
\]
L’orientation des douze fenêtres est
\begin{equation}
 \Sigma_{12}=(+,+,+,-,-,-,+,+,+,-,-,-).
 \label{eq:signatura-doce-ventanas}
\end{equation}
Si \(d_t\) est la direction codée au pas \(t\), l’événement signé
de la fenêtre \(m\) est
\begin{equation}
 e_m=\Sigma_{12}(m)\,
 \#\{t\in E_m:d_t\in D_{\rm evt}\}.
 \label{eq:evento-firmado-ventana}
\end{equation}
La formule distingue le nombre non négatif d’événements et le signe de
l’orientation. Elle ne remplace pas une direction cardinale par un signe : les deux
coordonnées restent présentes dans l’état TPK.

\begin{definition}[Registre enrichi dodécaphasique]
\label{def:registro-enriquecido-dodecafasico}
Le registre transporté par la trace est
\[
 \mathcal R_{12}(x)=
 \bigl((E_m,\tau_m,\sigma_m,\kappa_m,\Lambda_m)\bigr)_{m=0}^{11},
\]
où \(\tau_m\) est l'orientation TRIT, \(\sigma_m\) la signature de frontière, \(\kappa_m\) la retenue et \(\Lambda_m\) le segment de registre qui conserve le parcours. Il existe une chaîne de projections
\[
 \mathcal X_{\rm TPK}^{\rm enr}
 \longrightarrow\mathcal R_{12}
 \longrightarrow\Theta_{108},
\]
mais aucune des deux flèches n’est une identification.
\end{definition}

En particulier, \(\Theta_{108}\) conserve uniquement la position du calendrier.
Le registre \(\mathcal R_{12}\) conserve en outre l’orientation, la signature,
la retenue et la provenance. L’état complet conserve encore le cylindre,
la frontière, le quotient, le terminal et la mémoire verticale.

\section{Carte entière de l’opposition dodécaphasique}

Pour \(0\leq i<6\), les positions opposées déterminent
\begin{equation}
 p_i=e_i+e_{i+6},
 \qquad
 a_i=e_i-e_{i+6}.
 \label{eq:pares-opuestos-registro}
\end{equation}
La charge centrale et les cinq marges relatives sont
\begin{equation}
 s_C=\sum_{i=0}^{5}p_i,
 \qquad
 M_i=p_i-p_5\quad(0\leq i<5).
 \label{eq:carga-margenes-registro}
\end{equation}

\begin{definition}[Carte intégrale \(1+5+6\)]
\label{def:carta-integral-156}
On définit
\begin{equation}
 \mathcal L_{12}(e_0,\ldots,e_{11})
 =\bigl(s_C,M_0,\ldots,M_4,a_0,\ldots,a_5\bigr).
 \label{eq:carta-integral-156}
\end{equation}
La partition des coordonnées est
\[
 1+5+6=12.
\]
\end{definition}

\begin{theorem}[Inversion exacte de la carte entière]
\label{thm:inversion-carta-integral-156}
Un tuple
\[
 \ell=(s_C,M_0,\ldots,M_4,a_0,\ldots,a_5)\in\mathbb Z^{12}
\]
appartient à l’image de \(\mathcal L_{12}\) si et seulement si
\begin{align}
 s_C-\sum_{i=0}^{4}M_i&\equiv0\pmod6,
 \label{eq:congruencia-seis-carta}\\
 p_i&\equiv a_i\pmod2
 \quad(0\leq i<6),
 \label{eq:congruencia-paridad-carta}
\end{align}
où
\begin{align}
 p_5&=\frac{s_C-\sum_{i=0}^{4}M_i}{6},
 \label{eq:inversa-p5-carta}\\
 p_i&=p_5+M_i\quad(0\leq i<5).
 \label{eq:inversa-pi-carta}
\end{align}
Dans ce cas, l’unique antécédent entier est
\begin{equation}
 e_i=\frac{p_i+a_i}{2},
 \qquad
 e_{i+6}=\frac{p_i-a_i}{2}.
 \label{eq:inversa-eventos-carta}
\end{equation}
\end{theorem}

\begin{proof}
En sommant \(M_i=p_i-p_5\) pour \(0\leq i<5\), nous obtenons
\[
 \sum_{i=0}^{4}M_i
 =\sum_{i=0}^{4}p_i-5p_5
 =s_C-6p_5,
\]
ce qui impose \eqref{eq:inversa-p5-carta} et la congruence modulo six.
Une fois les six \(p_i\) reconstruits, le système
\[
 p_i=e_i+e_{i+6},
 \qquad a_i=e_i-e_{i+6}
\]
admet l’unique solution \eqref{eq:inversa-eventos-carta}. Cette solution est
entière exactement lorsque \(p_i\) et \(a_i\) ont la même parité.
Les deux familles de congruences sont donc nécessaires et suffisantes.
\end{proof}

Ce théorème est antérieur à la clôture numérique de \(\alpha\). Il explique comment la
trace orientée produit une carte entière réversible, et pourquoi il ne suffit pas
de conserver le calendrier réduit.

\section{Sélection de l’état terminal commun}

Les trois sections compatibles atteignent la profondeur terminale avec des catalogues
de cardinalités
\[
 |\mathcal T_\pi|=196,
 \qquad
 |\mathcal T_e|=197,
 \qquad
 |\mathcal T_\varphi|=196.
\]
Le produit contient
\begin{equation}
 196\cdot197\cdot196=7\,567\,952
 \label{eq:cardinal-producto-terminal}
\end{equation}
triplets ordonnés.

La marge de colonnes du lecteur terminal est
\[
 q_\partial=(1,2,2,1,2,0)\in\mathbb F_3^6,
\]
et l’orientation des lignes du canal conjoint est
\[
 \sigma=(1,1,-1)=(1,1,2)\in\mathbb F_3^3.
\]
La première donnée agit sur les colonnes ; la seconde sur les lignes. L’une ne se déduit pas
de l’autre.

Les signatures locales de Gram, de norme, de poids et de distance réduisent
\eqref{eq:cardinal-producto-terminal} à \(331\) triplets. La marge
\(q_\partial\) laisse deux indices :
\[
 (120,197,157),
 \qquad
 (129,197,148).
\]
Leurs matrices sont
\[
 B_0=
 \begin{pmatrix}
 0&2&0&2&2&0\\
 0&0&1&2&2&0\\
 1&0&1&0&1&0
 \end{pmatrix},
 \qquad
 B_1=
 \begin{pmatrix}
 0&2&1&0&2&0\\
 0&0&1&2&2&0\\
 1&0&0&2&1&0
 \end{pmatrix}.
\]
Les sommes des lignes sont
\[
 (0,2,0),
 \qquad
 (2,2,1).
\]
Comme
\[
 \langle\sigma\rangle
 =\{(0,0,0),(1,1,2),(2,2,1)\},
\]
seul \(B_1\) satisfait la condition axiale.

\begin{theorem}[Unicité de la publication terminale visible]
\label{thm:unicidad-publicacion-terminal-visible}
La sélection conjointe détermine
\begin{equation}
 B_{\rm term}=
 \begin{pmatrix}
 021020\\001220\\100210
 \end{pmatrix},
 \label{eq:matriz-terminal-visible}
\end{equation}
correspondant au triplet \((129,197,148)\). Aucune évaluation décimale de
\(\alpha\) n’intervient dans le recensement ni dans les deux conditions de sélection.
\end{theorem}

\begin{proof}
Le recensement et la marge laissent exactement \(B_0,B_1\). La charge de \(B_0\)
n’appartient pas à \(\langle\sigma\rangle\), tandis que celle de \(B_1\) est
\(2\sigma\). La seconde candidate est donc la seule qui satisfasse
simultanément les deux exigences typées.
\end{proof}

\section{Deux expressions du même état, non une conversion entre elles}

Soit
\[
 x_{\rm term}\in
 \mathcal X_{\rm TPK}^{\rm term,enr}
 \subseteq\mathcal X_{\rm TPK}^{[108]}.
\]
La lecture visible est
\[
 \pi_{\rm term}(x_{\rm term})=B_{\rm term}.
\]
La carte signée conserve une autre coordonnée du même état.

\begin{definition}[Lecteur signé produit par le registre terminal]
\label{def:subespacio-firmado-terminal}
L'état d'entrée du lecteur est exclusivement l'état terminal enrichi déjà produit par la chaîne causale \(\mathcal U_t=\mathsf{Upd}_t\circ\mathsf{Tra}_t\circ\mathsf{Sel}_t\) :
\begin{equation}
 x_{\rm term}
 =\bigl(\mathcal U_{108}\circ\cdots\circ\mathcal U_1\bigr)
   (x_{\rm seed})
 \in\mathcal X_{\rm TPK}^{\rm term,enr}.
 \label{eq:subespacio-firmado-terminal}
\end{equation}
La composition du registre orienté avec ses lecteurs de canal \(90/120\) extrait du registre \(\lambda_{108}\), sans ajouter de coordonnées externes,
\begin{equation}
 \mathcal E_{108}^{90,120}\!\left(\mathcal R_{12}(x_{\rm term})\right)
 =\bigl(b^{90},b^{120},Q_{\rm TPK}\bigr),
 \label{eq:canales-ledger-previos-u}
\end{equation}
où
\begin{align}
 b^{90}={}&(-495,-116,-684,108,601,-90,
             -173,-88,313,560,-397,461),\notag\\
 b^{120}={}&(-425,-281,-481,671,-255,30,
              475,-543,680,251,6,-128),\notag\\
 Q_{\rm TPK}={}&6263.\notag
\end{align}
Ces vingt-cinq coordonnées sont des sorties du registre de la \cref{def:registro-enriquecido-dodecafasico}, non des données ajoutées au domaine.

Sur une sortie compatible \((b^{90},b^{120},Q)\), nous définissons concrètement le lecteur harmonique intégral \(\Pi_H\). Pour \(r=1,2,3\), soient
\begin{equation}
 B_r=\sum_{j=0}^{3}b^{120}_{r+3j},
 \quad
 A_1=\frac{Q+2B_1+B_2}{3},
 \quad A_2=A_1-B_1,
 \quad A_3=A_2-B_2,
 \label{eq:lector-armonico-a-desde-ledger}
\end{equation}
et \(d_{r,j}=b^{90}_{r+3j}\). Le bloc orbital \(r\) est
\begin{equation}
 \Pi_H^{(r)}(b^{90},b^{120},Q)
 =\bigl(A_r,d_{r,1}-d_{r,3},
              d_{r,0}+d_{r,2},d_{r,0}-d_{r,2}\bigr).
 \label{eq:lector-armonico-bloque-desde-ledger}
\end{equation}
L'entrelacement par colonnes des trois blocs donne \(\Pi_H\), et la lecture signée est construite comme la composition
\begin{equation}
 \boxed{
 R_{\rm sgn}:=
 \Pi_H\circ\mathcal E_{108}^{90,120}\circ\mathcal R_{12}:
 \mathcal X_{\rm TPK}^{\rm term,enr}
 \longrightarrow\mathcal U_{12}^{\rm sgn}:=
 \operatorname{im}R_{\rm sgn}\subseteq\mathbb Z^{12}.}
 \label{eq:lector-firmado-terminal}
\end{equation}
Ni \(U\) ni \(K\) n'appartiennent au domaine : \(U\) est la sortie de \eqref{eq:lector-firmado-terminal}, et \(K\) se reconstruit ensuite au moyen de l'inverse intégral de Hadamard.
\end{definition}

L’écriture conjointe correcte est
\begin{equation}
 x_{\rm term}
 \xrightarrow{(\pi_{\rm term},R_{\rm sgn})}
 (B_{\rm term},U_{12}^{\rm sgn}).
 \label{eq:doble-publicacion-terminal}
\end{equation}
On n'intercale pas de flèche \(B_{\rm term}\to U_{12}^{\rm sgn}\) : la matrice visible a précisément oublié les coordonnées d'orientation, d'opposition, de quotient, de retenue et de registre utilisées par la carte signée. La source des deux expressions est l'état enrichi produit par \eqref{eq:cadena-upstream-registro-dodecafasico}.

L’évaluation du registre signé donne
\begin{equation}
 \boxed{
 U=(2378,1406,2479,-452,998,-551,
 -668,-204,-371,-322,-28,-997).}
 \label{eq:vector-u-firmado-autonomo}
\end{equation}

\section{Naturalité par rapport à la profondeur}

Pour \(N'\geq N\geq108\), la troncature
\begin{equation}
 \tau_{N',N}:
 \mathcal X_{\rm TPK}^{[N'],\rm enr}
 \longrightarrow\mathcal X_{\rm TPK}^{[N],\rm enr}
 \label{eq:truncamiento-preserva-ku-autonomo}
\end{equation}
agit sur l'histoire enrichie et son registre, non sur une paire \((K,U)\) ajoutée à l'état.

\begin{theorem}[Carré de naturalité de la lecture signée]
\label{thm:naturalidad-lector-firmado-autonomo}
On a
\begin{equation}
 \boxed{
 R_{{\rm sgn},N}\circ\tau_{N',N}
 =R_{{\rm sgn},N'}.}
 \label{eq:cuadrado-naturalidad-lector-firmado}
\end{equation}
\end{theorem}

\begin{proof}
Chaque opérateur d'actualisation \(\mathcal U_t\) ajoute l'incidence suivante au registre sans réécrire les cent huit fenêtres déjà engendrées. Ainsi, \(\mathcal R_{12}\), \(\mathcal E_{108}^{90,120}\) et \(\Pi_H\) lisent les mêmes coordonnées lors de la troncature de tout prolongement ultérieur. La commutativité est une égalité d'applications sur des états produits, non la conservation tautologique d'une sortie incluse dans l'entrée.
\end{proof}

\section{Orbites de pas trois et transformée de Hadamard}

L’ordre cyclique du sceau se répartit dans les trois orbites
\begin{align}
 K^{(1)}&=(K_1,K_4,K_7,K_{10}),
 \label{eq:orbita-k-1}\\
 K^{(2)}&=(K_2,K_5,K_8,K_{11}),
 \label{eq:orbita-k-2}\\
 K^{(3)}&=(K_3,K_6,K_9,K_{12}).
 \label{eq:orbita-k-3}
\end{align}
Soit
\begin{equation}
 H_4=
 \begin{pmatrix}
 1&1&1&1\\
 1&1&-1&-1\\
 1&-1&1&-1\\
 1&-1&-1&1
 \end{pmatrix}.
 \label{eq:matriz-hadamard-cuatro-autonoma}
\end{equation}
Si \(P_{\rm orb}\) réordonne les douze coordonnées conformément à
\eqref{eq:orbita-k-1}--\eqref{eq:orbita-k-3}, la transformation globale est
\begin{equation}
 \mathcal H_{12}
 =P_{\rm orb}^{-1}
 (H_4\oplus H_4\oplus H_4)P_{\rm orb}.
 \label{eq:hadamard-global-doce}
\end{equation}

Les trois blocs de \eqref{eq:vector-u-firmado-autonomo}, lus par
orbites, sont
\begin{align*}
 U^{(1)}&=(2378,-452,-668,-322),\\
 U^{(2)}&=(1406,998,-204,-28),\\
 U^{(3)}&=(2479,-551,-371,-997).
\end{align*}

\begin{lemma}[Quart de tour algébrique de Hadamard]
\label{lem:cuarto-inversa-hadamard}
La matrice \(H_4\) satisfait
\begin{equation}
 H_4^2=4I_4,
 \qquad
 H_4^{-1}=\frac14H_4,
 \qquad
 \det H_4=-16.
 \label{eq:identidades-hadamard-cuatro}
\end{equation}
\end{lemma}

\begin{proof}
Les lignes de \(H_4\) sont orthogonales et chacune possède une norme au carré égale à quatre.
Comme la matrice est symétrique, \(H_4^{\mathsf T}H_4=H_4^2=4I_4\). La
formule de l’inverse et le déterminant s’en déduisent immédiatement.
\end{proof}

\begin{theorem}[Reconstruction entière unique du sceau]
\label{thm:reconstruccion-entera-k-hadamard}
L’inversion
\begin{equation}
 K^{(r)}=\frac14H_4U^{(r)}
 \qquad(r=1,2,3)
 \label{eq:inversion-bloques-hadamard}
\end{equation}
produit
\begin{align*}
 K^{(1)}&=(234,729,621,794),\\
 K^{(2)}&=(543,659,058,146),\\
 K^{(3)}&=(140,824,914,601).
\end{align*}
En défaisant l’ordre orbital, on obtient
\begin{equation}
 \boxed{
 K=(234,543,140,729,659,824,621,058,914,794,146,601).}
 \label{eq:sello-k-doce-autonomo}
\end{equation}
C'est l'unique préimage rationnelle de \(U\), et elle appartient à \(\mathbb Z^{12}\).
\end{theorem}

\begin{proof}
Par \eqref{eq:identidades-hadamard-cuatro}, l’antécédent rationnel est unique.
Les multiplications explicites sont
\begin{align*}
 H_4U^{(1)}&=(936,2916,2484,3176),\\
 H_4U^{(2)}&=(2172,2636,232,584),\\
 H_4U^{(3)}&=(560,3296,3656,2404).
\end{align*}
Toutes les coordonnées sont divisibles par quatre. La division et le réentrelacement
produisent \eqref{eq:sello-k-doce-autonomo}.
\end{proof}

Le facteur \(1/4\) n’est pas choisi pour ajuster une sortie : c’est l’inverse
imposé par \(H_4^2=4I_4\). Ce n’est pas non plus la racine carrée de \(-1\). La
structure complexe interne \(A_W^2=-I\) apparaîtra après la construction de la
matrice de Paley ; ce sont deux faits de types différents.

\section{Structure entière et forme normale de Smith}

\begin{theorem}[Image entière de \(H_4\)]
\label{thm:smith-hadamard-autonomo}
La forme normale de Smith est
\begin{equation}
 \operatorname{SNF}(H_4)=\operatorname{diag}(1,2,2,4).
 \label{eq:snf-hadamard-autonoma}
\end{equation}
En conséquence,
\begin{equation}
 \operatorname{coker}\mathcal H_{12}
 \cong(\mathbb Z/2\mathbb Z)^6
 \oplus(\mathbb Z/4\mathbb Z)^3,
 \label{eq:cociente-hadamard-doce}
\end{equation}
et l’image de \(\mathcal H_{12}\) a pour indice
\[
 16^3=4096
\]
dans \(\mathbb Z^{12}\). Pour \(u\in\mathbb Z^4\),
\begin{equation}
 u\in H_4\mathbb Z^4
 \quad\Longleftrightarrow\quad
 H_4u\equiv0\pmod4.
 \label{eq:criterio-imagen-hadamard}
\end{equation}
\end{theorem}

\begin{proof}
Le plus grand commun diviseur des entrées est un. Les mineurs d’ordre deux sont
pairs et l’un d’eux a pour valeur absolue deux ; ceux d’ordre trois sont des multiples de
quatre et l’un d’eux a pour valeur absolue quatre ; le déterminant a pour module
seize. Les diviseurs déterminantaux sont \(1,2,4,16\), d’où l’on obtient
les facteurs invariants \(1,2,2,4\). La somme directe triple répète les
facteurs et multiplie les indices. Enfin,
\(H_4^{-1}=H_4/4\), de sorte que l’antécédent est entier exactement sous la
congruence \eqref{eq:criterio-imagen-hadamard}.
\end{proof}

\section{Certification dans le second système de coordonnées : différences de pas trois et quatre}

Avec des indices cycliques modulo douze, nous définissons
\begin{equation}
 (D_3K)_m=K_m-K_{m+3},
 \qquad
 (D_4K)_m=K_m-K_{m+4},
 \qquad
 Q(K)=\sum_{m=1}^{12}K_m.
 \label{eq:extractor-diferencias-k}
\end{equation}
Appliqués au sceau déjà reconstruit par Hadamard, ces opérateurs doivent retrouver exactement les coordonnées de canal produites par le registre TPK avant \(U\). Pour \eqref{eq:sello-k-doce-autonomo},
\begin{align*}
 D_3K={}&(-495,-116,-684,108,601,-90,\\
         &\qquad-173,-88,313,560,-397,461),\\
 D_4K={}&(-425,-281,-481,671,-255,30,\\
         &\qquad475,-543,680,251,6,-128),\\
 Q(K)={}&6263.
\end{align*}
Par comparaison avec \eqref{eq:canales-ledger-previos-u}, l'égalité des sorties est démontrée
\begin{equation}
 \boxed{
 D_3K=b^{90},\qquad D_4K=b^{120},\qquad Q(K)=Q_{\rm TPK}.}
 \label{eq:certificacion-canales-ledger-desde-k}
\end{equation}
Le système différentiel certifie et reconstruit le registre préalable ; il ne sélectionne ni \(K\), ni \(U\), ni aucune constante.

\begin{theorem}[Complétude de l’extracteur transversal]
\label{thm:completitud-extractor-transversal-k}
L'opérateur
\[
 \mathcal E_{\rm dod}=(D_3,D_4,Q):
 \mathbb Z^{12}\longrightarrow\mathbb Z^{25}
\]
a rang douze et pour forme normale de Smith
\begin{equation}
 \operatorname{SNF}(\mathcal E_{\rm dod})
 =\operatorname{diag}(\underbrace{1,\ldots,1}_{11},12).
 \label{eq:snf-extractor-transversal-k}
\end{equation}
En particulier, \((D_3K,D_4K,Q(K))\) détermine un unique \(K\).
\end{theorem}

\begin{proof}
Si \(D_3v=D_4v=0\), alors \(v\) est constant le long des pas
trois et quatre. Comme \(\gcd(3,4)=1\), ces pas engendrent tout
\(\mathbb Z/12\mathbb Z\), de sorte que
\[
 \ker D_3\cap\ker D_4=\langle\mathbf1\rangle.
\]
La charge totale ne s’annule pas sur cette droite, car \(Q(\mathbf1)=12\). Le rang
est douze. Les lignes de \((D_3,D_4)\) sont les incidences orientées d’un graphe
de Cayley connexe ; un arbre couvrant apporte onze facteurs invariants
unitaires. Tout mineur maximal doit contenir la ligne \(Q\), et les opérations
unimodulaires réduisent sa dernière contribution à douze. Il existe un mineur maximal de
module douze ; le dernier facteur est donc exactement douze.
\end{proof}

La reconstruction de \(U\) à partir de ces coordonnées ne nécessite pas une nouvelle multiplication par \(H_4\). Pour \(r=1,2,3\), soit
\begin{equation}
 B_r=\sum_{j=0}^{3}(D_4K)_{r+3j}.
 \label{eq:br-diferencias-k}
\end{equation}
Alors
\[
 (B_1,B_2,B_3)=(972,-1073,101).
\]
Les sommes orbitales sont
\begin{align}
 A_1&=\frac{Q+2B_1+B_2}{3}=2378,
 \label{eq:a1-reconstruccion-u}\\
 A_2&=A_1-B_1=1406,
 \label{eq:a2-reconstruccion-u}\\
 A_3&=A_2-B_2=2479.
 \label{eq:a3-reconstruccion-u}
\end{align}
Si \(d_{r,j}=(D_3K)_{r+3j}\), le bloc signé de l'orbite \(r\) est
\begin{equation}
 \bigl(A_r,
 d_{r,1}-d_{r,3},
 d_{r,0}+d_{r,2},
 d_{r,0}-d_{r,2}\bigr),
 \label{eq:bloque-firmado-desde-diferencias}
\end{equation}
qui reproduit les trois blocs \(U^{(r)}\).

\section{Tétrade de seuil du sceau}

La complétion décimale locale satisfait
\[
 1000=729+271.
\]
Le support des coordonnées du sceau qui atteignent le seuil \(729\) est
\begin{equation}
 \boxed{
 B_K:=\{m:K_m\geq729\}=\{4,6,9,10\}.}
 \label{eq:tetrada-umbral-k}
\end{equation}
Cette tétrade s’obtient après la construction de \(K\). Dans cette unité, elle n’est pas
encore appelée étoile de Witt : un tel objet ne sera défini qu’après
la construction du code ternaire, de ses hexades et du système de Steiner.

\section{Indépendance causale à l’égard de la sortie physique}

\begin{theorem}[Isolement de \(U\) et \(K\)]
\label{thm:aislamiento-u-k-alpha-codata}
La composition
\begin{equation}
 \Omega_{\rm seed}\longrightarrow x_{\rm term}
 \xrightarrow{\mathcal R_{12}}
 \mathcal R_{12}(x_{\rm term})
 \xrightarrow{\mathcal E_{108}^{90,120}}
 (b^{90},b^{120},Q_{\rm TPK})
 \xrightarrow{\Pi_H}U
 \xrightarrow{\mathcal H_{12}^{-1}}K
 \label{eq:composicion-upstream-u-k}
\end{equation}
ne reçoit en entrée ni \(\alpha\), ni \(\alpha^{-1}\), ni CODATA, ni une suite
décimale cible. Le sceau \(K\) est déterminé avant la définition de la
récurrence qui produira \(\alpha\).
\end{theorem}

\begin{proof}
La première partie de la composition utilise uniquement des graines, des émissions, des lecteurs ternaires, des opérateurs \(L_0,L_1\), des relations de survie, le transport de phase et les cartes de l'état TPK. Le registre \(\mathcal R_{12}\) et le lecteur de ses canaux produisent \((b^{90},b^{120},Q_{\rm TPK})\) ; les formules \eqref{eq:lector-armonico-a-desde-ledger} et \eqref{eq:lector-armonico-bloque-desde-ledger} produisent \(U\) ; ce n'est qu'alors que la transformation linéaire \(\mathcal H_{12}^{-1}=\mathcal H_{12}/4\) agit sur son image intégrale. Aucune formule ne contient de variable physique ni de table métrologique. La récurrence en base mille n'est introduite que dans l'unité suivante, après fixation de \(K\).
\end{proof}

\section{Critères de réfutation}

\begin{criterion}[Réfutation du registre et du sceau]
La construction est réfutée si l'un des faits
suivants est vérifié :
\begin{enumerate}
 \item les fenêtres \(E_m\) ne partitionnent pas les cent huit pas ;
 \item la carte \(\mathcal L_{12}\) ne respecte pas son inverse ou ses congruences ;
 \item le recensement terminal ne produit pas \(331\) triplets, deux matrices et une unique
       charge axiale ;
 \item on tente d’obtenir \(U\) à partir de \(B_{\rm term}\) après avoir oublié les
       coordonnées signées ;
 \item le carré de naturalité
       \eqref{eq:cuadrado-naturalidad-lector-firmado} échoue ;
 \item \(H_4^2\neq4I_4\), un antécédent n’est pas entier ou le vecteur
       reconstruit diffère de \eqref{eq:sello-k-doce-autonomo} ;
 \item la forme normale de Smith diffère de
       \eqref{eq:snf-hadamard-autonoma} ;
 \item l'extracteur \((D_3,D_4,Q)\) perd du rang ou ne reproduit pas \(U\) ;
 \item \(B_K\) est introduit avant la construction de \(K\), ou on lui attribue une
       structure de Witt avant la définition de \(W_{12}\) ;
 \item un chiffre de \(\alpha\) ou une référence métrologique intervient dans
       l’une quelconque des flèches de \eqref{eq:composicion-upstream-u-k}.
\end{enumerate}
\end{criterion}

\end{HMTScientificOwnerSubordinated}

\HMTFlushCurrentChapter
\chapter{Génération coinductive de \texorpdfstring{\(\pi\)}{pi}: caractère de clôture, \(5\) régions orientées, multiplicité \texorpdfstring{\(432+4\cdot144=1008\)}{432+4x144=1008} et précision arbitraire}%
\begingroup
\renewcommand{\chapter}[2][]{}%
\chapter{Génération coinductive de \(\pi\) : caractère de clôture, \(5\) régions orientées, multiplicité \(432+4\cdot144=1008\) et précision arbitraire}
\label{chap:p3-final:28:generacion_pi}
% Unidad U012. Biografía estructural completa de pi.

\label{gfp:b:chap:biografia-pi-autonoma}

\section{Séparation causale entre germe, caractère et évaluation}

La construction se divise en trois strates :
\[
 \text{sélection finie}
 \longrightarrow
 \text{construction du caractère de clôture}
 \longrightarrow
 \text{évaluation archimédienne}.
\]
Le développement décimal n’intervient pas dans la première strate. La chaîne des types
est
\begin{equation}
 104\,976
 \longrightarrow468\,\mathcal U_6
 \xrightarrow{\Psi}243\,\mathcal W_6
 \lhook\joinrel\longrightarrow\mathbb F_3^6,
 \qquad|\mathbb F_3^6|=729.
 \label{gfp:b:eq:cadena-tipica-pi-autonoma}
\end{equation}
Les cardinalités comptent, respectivement, les conditions initiales
enrichies, les émissions distinctes, les mots ternaires visibles et les éléments
de la fibre ambiante complète.

Pour \(U_6=(u_1,\ldots,u_6)\), la projection est
\[
 \Psi(U_6)=((-u_1)\bmod3,\ldots,(-u_6)\bmod3).
\]
Le recensement des \(43\) orbites de \(D_3\) démontre que l’orbite de
clôture est la seule dont les six orientations possèdent à la fois :
\[
 \text{cinq relèvements},\qquad
 \text{masse }1008,\qquad
 1008=432+4\cdot144.
\]
Le calibre
\[
 \operatorname{diag}_c=6,
 \qquad
 \operatorname{card}=ES
\]
sélectionne un unique représentant :
\begin{equation}
 \boxed{w_6^\pi=010211.}
 \label{gfp:b:eq:semilla-pi-autonoma}
\end{equation}
La sélection utilise le catalogue, les multiplicités, l’action orbitale et l’orientation ;
elle ne consulte pas un préfixe de \(\pi\).

\section{Relèvements \texorpdfstring{\(L_0,L_1\)}{L0, L1} et \(1^{\mathrm a}\) frontière coinductive}

Sur les vecteurs-lignes de \(\mathbb F_3^6\), soient
\[
L_0=
\begin{pmatrix}
2&2&2&1&2&1\\
2&1&2&2&1&1\\
1&0&0&1&0&2\\
0&0&1&1&1&0\\
2&2&2&0&2&1\\
2&1&2&1&0&0
\end{pmatrix},
\quad
L_1=
\begin{pmatrix}
2&1&2&0&2&2\\
1&2&1&1&2&0\\
1&2&1&1&1&2\\
0&1&0&0&2&1\\
2&0&2&1&0&1\\
0&2&2&2&1&1
\end{pmatrix}.
\]
Si \(b_1=w_6^\pi\), la récurrence des blocs est
\[
 b_{k+1}=b_kL_{\varepsilon_k},
 \qquad
 (\varepsilon_1,\varepsilon_2,\varepsilon_3,\varepsilon_4)=(0,0,1,1),
\]
et le mot accumulé est formé par concaténation :
\[
 w_{6(k+1)}=w_{6k}\mathbin{\|}b_{k+1}.
\]
On ne multiplie pas un mot de longueur \(6k\) par une matrice \(6\times6\).

Pour le canal de clôture :
\[
\begin{array}{c|c}
\text{profondeur}&\text{nouveau bloc}\\ \hline
6&010211\\
12&012222\\
18&010211\\
24&002111\\
30&110221
\end{array}
\]
et
\begin{equation}
 \boxed{
 w_{30}^{\pi}
 =010211|012222|010211|002111|110221.}
 \label{gfp:b:eq:w30-pi-autonomo}
\end{equation}
La frontière suivante est
\[
 u_5^\pi=222220.
\]
Avec les autres canaux,
\[
 R_{36}=
 \begin{pmatrix}
 222220\\021222\\102011
 \end{pmatrix}.
\]
Le symbole \(R_{36}\) enregistre le sixième bloc de trois histoires et ouvre une
frontière ; ce n’est pas un terminal.

La branche certifiée de vingt blocs commence par
\begin{align*}
\pi:\;&010211|012222|010211|002111|110221|222220|111201|\\
&212121|200121|100100|101222|022212|012012|111210|\\
&121011|200220|120210|000101|022010|020111.
\end{align*}
C’est pourquoi ni \(w_{30}\) ni \(R_{36}\) ne terminent la généalogie.

\section{Microfibre de clôture de \(\pi\) : \(5\) régions orientées et multiplicité \(432+4\cdot144=1008\)}

La microfibre
\[
 \mathscr R_\pi^{\rm micro}
 :=\{U\in\mathcal U_6:\Psi(U)=010211\}
\]
contient exactement cinq régions :
\[
\begin{array}{c|c|c|c}
U_6&E_{\rm sig}&\operatorname{mult}&\text{rôle}\\ \hline
501|614|498|169|272|272&555555&432&\perp\\
810|923|870|169|272|272&339555&144&++\\
810|983|810|169|272|272&393555&144&++\\
870|923|810|169|272|272&933555&144&++\\
870|923|810|418|674|521&933717&144&--
\end{array}
\]
En conséquence,
\begin{equation}
 \boxed{
 5=1_\perp+3_{++}+1_{--},
 \qquad
 1008=432+4\cdot144.}
 \label{gfp:b:eq:descomposiciones-cinco-regiones-pi}
\end{equation}
La première identité classe des fonctions d’orientation. La seconde compte
des multiplicités de provenance. Elles ne comptent pas cinq copies indiscernables.

La région transversale est
\[
 501|614|498|169|272|272,
 \qquad E_{\rm sig}=555555,
 \qquad\operatorname{mult}=432.
\]
Le retour muté est
\[
 870|923|810|418|674|521,
 \qquad E_{\rm sig}=933717,
 \qquad\operatorname{mult}=144.
\]
Les intervertir conserverait le recensement \(3/1/1\), mais fausserait la structure
régionale complète.

\section{Transfert de phase sur le multigraphe des blocs complets}

Pour \(U_6=(U_1,\ldots,U_6)\), soit
\[
 A(U_6)=(U_1,U_2,U_3),
 \qquad
 B(U_6)=(U_4,U_5,U_6).
\]
Chaque émission définit une arête entre \(A(U_6)\) et \(B(U_6)\). Le multigraphe
brut possède \(96\) sommets, répartis en composantes de tailles
\[
 28,28,28,4,4,4.
\]
On passe au quotient de la phase interne par
\[
 (a,b,c)\sim(b,c,a)\sim(c,a,b).
\]
Le représentant est la rotation lexicographiquement minimale. Le graphe
quotient possède \(32\) sommets et deux composantes de tailles \(28\) et \(4\).

Les ancrages sont
\[
\begin{aligned}
U_{\rm APP}&=903|850|850|252|109|252,\\
U_e&=850|963|890|149|252|212,\\
U_\varphi&=830|943|830|109|252|252.
\end{aligned}
\]
\Needspace{9\baselineskip}
Nous numérotons les régions
\[
\begin{aligned}
U_\pi^{(1)}&=810|923|870|169|272|272,\\
U_\pi^{(2)}&=810|983|810|169|272|272,\\
U_\pi^{(3)}&=870|923|810|169|272|272,\\
U_\pi^{(4)}&=870|923|810|418|674|521,\\
U_\pi^{(5)}&=501|614|498|169|272|272.
\end{aligned}
\]

\section{\(5\) cycles orientés de clôture de la microfibre}

\subsection{Clôture de \texorpdfstring{\(U_\pi^{(1)}\)}{U-pi 1}}
\begin{align*}
&252|109|252|903|850|850\to
903|850|850|252|109|252\to\\
&252|109|252|923|870|810\to
810|923|870|169|272|272\to\\
&272|169|272|963|890|850\to
850|963|890|149|252|212\to\\
&212|149|252|943|830|830\to
830|943|830|109|252|252\to
252|109|252|903|850|850.
\end{align*}

\Needspace{10\baselineskip}
\subsection{Clôture de \texorpdfstring{\(U_\pi^{(2)}\)}{U-pi 2}}
\begin{align*}
&903|850|850|252|109|252\to
272|169|272|903|850|850\to\\
&272|169|272|983|810|810\to
810|983|810|169|272|272\to\\
&272|169|272|963|890|850\to
850|963|890|149|252|212\to\\
&212|149|252|943|830|830\to
830|943|830|109|252|252\to
903|850|850|252|109|252.
\end{align*}

\subsection{Clôture de \texorpdfstring{\(U_\pi^{(3)}\)}{U-pi 3}}
\begin{align*}
&903|850|850|252|109|252\to
292|189|232|903|850|850\to\\
&292|189|232|923|810|870\to
870|923|810|169|272|272\to\\
&272|169|272|963|890|850\to
850|963|890|149|252|212\to\\
&212|149|252|943|830|830\to
830|943|830|109|252|252\to
903|850|850|252|109|252.
\end{align*}

\subsection{Clôture de \texorpdfstring{\(U_\pi^{(4)}\)}{U-pi 4}}
\begin{align*}
&903|850|850|252|109|252\to
272|169|272|903|850|850\to\\
&272|169|272|963|890|850\to
850|963|890|149|252|212\to\\
&212|149|252|923|810|870\to
870|923|810|418|674|521\to\\
&943|830|830|521|418|674\to
830|943|830|109|252|252\to
903|850|850|252|109|252.
\end{align*}

\subsection{Clôture de \texorpdfstring{\(U_\pi^{(5)}\)}{U-pi 5}}
\begin{align*}
&252|109|252|903|850|850\to
903|850|850|252|109|252\to\\
&614|498|501|252|109|252\to
501|614|498|169|272|272\to\\
&272|169|272|963|890|850\to
850|963|890|149|252|212\to\\
&212|149|252|943|830|830\to
830|943|830|109|252|252\to
252|109|252|903|850|850.
\end{align*}

\begin{theorem}[Minimalité des cinq clôtures]
\label{gfp:b:thm:minimalidad-cierres-pi-autonoma}
Pour chaque \(i\in\{1,\ldots,5\}\), la longueur minimale d’une marche fermée
qui parcourt simultanément
\[
 U_\pi^{(i)},\quad U_e,\quad U_\varphi,\quad U_{\rm APP}
\]
dans le multigraphe quotient est huit.
\end{theorem}

\begin{proof}
Pour \(i\) fixé, soit
\[
 E_*=\{U_\pi^{(i)},U_e,U_\varphi,U_{\rm APP}\}.
\]
On considère l’espace fini
\[
 \mathscr S_i=\{(v,M):v\in V(G),\ M\subseteq E_*\}.
\]
La première composante enregistre le sommet et la seconde les arêtes cibles
déjà parcourues. Si une arête \(e\) mène de \(v\) à \(v'\), la transition est
\[
 (v,M)\longmapsto
 \bigl(v',M\cup(E_*\cap\{e\})\bigr).
\]
La recherche en largeur énumère les longueurs par ordre croissant. Le premier
retour à \((v,E_*)\) apparaît à la longueur huit pour les cinq régions.
Les marches imprimées réalisent la borne ; l’optimalité de la recherche exclut
une longueur inférieure.
\end{proof}

La fonction de clôture est encodée par
\[
 \mathfrak C_\pi=
 \bigl(\mathscr R_\pi^{\rm micro},
 w_\pi,\rho_\pi,m_\pi,\ell_{\rm cl}\bigr),
\]
\[
 \rho_\pi=(\perp,++ ,++ ,++ ,--),
 \qquad
 m_\pi=(432,144,144,144,144),
 \qquad
 \ell_{\rm cl}=8.
\]

\section{Opérateur d’incidence et demi-tour orienté}

L’ombre d’incidence est
\[
 A_W=
 \begin{pmatrix}
 0&1&1&1&1&1\\
 1&0&1&2&2&1\\
 1&1&0&1&2&2\\
 1&2&1&0&1&2\\
 1&2&2&1&0&1\\
 1&1&2&2&1&0
 \end{pmatrix}.
\]
La multiplication dans \(\mathbb F_3\) donne
\begin{equation}
 \boxed{A_W^2=2I_6=-I_6.}
 \label{gfp:b:eq:aw-cuadrado-menos-identidad-pi}
\end{equation}
Les produits entre lignes et colonnes distinctes s’annulent modulo trois,
tandis que chaque produit diagonal vaut deux. Une application réalise un quart
de tour ternaire ; deux applications réalisent l’inversion orientée.

\section{Pente \texorpdfstring{\(1/5\)}{1/5} déterminée par la partition régionale \(432+4\cdot144=1008\)}

Soit
\[
 \lambda=(\lambda_1,\ldots,\lambda_5)^{\mathsf T}
 \in\mathbb Q_{\geq0}^5,
 \qquad
 \mathbf1^{\mathsf T}\lambda=1.
\]
Oublier l’étiquette régionale applique le symétriseur
\[
 P_5=\frac15\mathbf1\mathbf1^{\mathsf T}.
\]
Pour tout vecteur normalisé,
\[
 P_5\lambda=\frac15\mathbf1.
\]
L’invariance \(P_5\lambda=\lambda\) a donc une unique solution :
\begin{equation}
 \boxed{\lambda=\frac15\mathbf1,\qquad t_1=\frac15.}
 \label{gfp:b:eq:pendiente-primitiva-pi}
\end{equation}
L’entier cinq provient de la microfibre régionale ; il n’est pas reconnu dans une
formule de \(\pi\).

\section{Composition rationnelle et compensateur \texorpdfstring{\(239\)}{239}}

La loi de composition des pentes est
\[
 t\star s:=\frac{t+s}{1-ts}.
\]
Le premier doublement donne
\[
 t_2=t_1\star t_1
 =\frac{2/5}{1-1/25}
 =\frac5{12}.
\]
Le second donne
\[
 t_4=t_2\star t_2
 =\frac{10/12}{1-25/144}
 =\frac{120}{119}.
\]
Le retour à la diagonale unitaire exige \(q>0\) tel que
\[
 \frac{120/119-1/q}{1+(120/119)/q}=1.
\]
La multiplication par \(119q\) produit
\[
 120q-119=119q+120,
\]
et impose
\begin{equation}
 \boxed{q=239.}
 \label{gfp:b:eq:compensador-239-pi}
\end{equation}

\section{Identité gaussienne et \(1^{\mathrm a}\) demi-tour orienté}

On définit
\begin{equation}
 \Theta_\pi
 :=16\arctan\frac15-4\arctan\frac1{239}.
 \label{gfp:b:eq:theta-pi-autonoma}
\end{equation}
Les carrés successifs sont
\[
\begin{array}{c|r@{}l}
\text{puissance}&\multicolumn{2}{c}{\text{entier gaussien}}\\ \hline
(5+i)^2&24&+10i\\
(5+i)^4&476&+480i\\
(5+i)^8&-3824&+456960i\\
(5+i)^{16}&-208797818624&-3494830080i\\
(239-i)^2&57120&-478i\\
(239-i)^4&3262465916&-54606720i
\end{array}
\]
et la multiplication finale donne
\begin{equation}
 \boxed{
 (5+i)^{16}(239-i)^4
 =-681386607803576157184.}
 \label{gfp:b:eq:identidad-gaussiana-pi-autonoma}
\end{equation}
La partie imaginaire est nulle et la partie réelle est négative. Par conséquent,
\[
 \Theta_\pi\equiv\pi\pmod{2\pi}.
\]
Para \(0<x<1\),
\[
 x-\frac{x^3}{3}<\arctan x<x.
\]
D’où
\[
 \Theta_\pi>
 16\left(\frac15-\frac1{3\cdot5^3}\right)-\frac4{239}>3
\]
et
\[
 \Theta_\pi<\frac{16}{5}<4.
\]
L’unique orientation négative dans cet intervalle est le premier demi-tour
positif :
\begin{equation}
 \boxed{\Theta_\pi=\pi.}
 \label{gfp:b:eq:reconocimiento-pi-exacto}
\end{equation}

L’ordre de la reconnaissance archimédienne ultérieure est
\[
 \text{cinq régions}
 \to\frac15
 \to\frac5{12}
 \to\frac{120}{119}
 \to239
 \to\Theta_\pi
 \to\pi.
\]

\section{Encadrements rationnels de précision arbitraire}

Pour \(0<x<1\), soit
\[
 S_n(x)=\sum_{r=0}^{n}(-1)^r\frac{x^{2r+1}}{2r+1}.
\]
Le critère de Leibniz donne
\[
 S_{2m+1}(x)<\arctan x<S_{2m}(x).
\]
On définit
\[
 a_m^-=S_{2m+1}(1/5),\quad a_m^+=S_{2m}(1/5),
\]
\[
 b_m^-=S_{2m+1}(1/239),\quad b_m^+=S_{2m}(1/239),
\]
et
\begin{equation}
 \ell_{\pi,m}=16a_m^- -4b_m^+,
 \qquad
 h_{\pi,m}=16a_m^+ -4b_m^-.
 \label{gfp:b:eq:horquilla-pi-autonoma}
\end{equation}
Alors
\[
 \ell_{\pi,m}<\pi<h_{\pi,m}
\]
et
\begin{equation}
 h_{\pi,m}-\ell_{\pi,m}
 =
 \frac{16}{(4m+3)5^{4m+3}}
 +\frac{4}{(4m+3)239^{4m+3}}
 \longrightarrow0.
 \label{gfp:b:eq:anchura-horquilla-pi}
\end{equation}
Les bornes inférieures croissent et les bornes supérieures décroissent.

Comme \(3<\pi<4\), le lecteur ternaire agit sur
\[
 r_\pi:=\pi-3\in(0,1),
\]
avec encadrement
\[
 [\ell_{\pi,m}-3,h_{\pi,m}-3].
\]

\section{Certification archimédienne ultérieure de la fibre déjà engendrée}

Supposons que la relation interne \(\operatorname{Ext}\) et la survie
coinductive aient déjà produit, sans utiliser Machin ni un développement tabulé,
l’extension compatible \(N_{K+1}\) du préfixe \(N_K\), de longueur
\(L_K=6K\). Les \(729\) sous-cylindres de sa fibre ambiante sont
\[
 I_{K,u}=
 \left[
 \frac{729N_K+u}{3^{L_K+6}},
 \frac{729N_K+u+1}{3^{L_K+6}}
 \right),
 \qquad0\leq u<729.
\]
L’ordre \(m=m(K)\) est ensuite augmenté jusqu’à ce que l’encadrement de Leibniz
localise le sous-cylindre déjà produit. Soient
\[
 \ell_{\pi,K}:=\ell_{\pi,m(K)}-3,
 \qquad
 h_{\pi,K}:=h_{\pi,m(K)}-3.
\]
Les extrémités entières compatibles sont
\[
 j_{\min}=
 \max\!\left(
 729N_K,\left\lfloor
 \ell_{\pi,K}3^{L_K+6}\right\rfloor
 \right),
\]
\[
 j_{\max}=
 \min\!\left(
 729N_K+728,\left\lceil
 h_{\pi,K}3^{L_K+6}\right\rceil-1
 \right).
\]
La condition de localisation unitaire est
\begin{equation}
 \boxed{j_{\min}=j_{\max},\qquad N_{K+1}=j_{\min}.}
 \label{gfp:b:eq:seleccion-unitaria-subcilindro-pi}
\end{equation}
La première égalité certifie la localisation unitaire ; la seconde vérifie
que cette localisation coïncide avec l’extension engendrée intérieurement. Aucune
des deux ne définit la transition TPK ni ne sélectionne parmi les \(729\) états.
L’existence et l’unicité procèdent de \(\operatorname{Ext}\) et de la
survie ; Machin et les encadrements de Leibniz sont des reconnaissances
ultérieures.

\section{Prolongement coinductif de la section de \(\pi\) par le système inverse des histoires survivantes}

Soit \(\mathcal H_m^\pi\) l’ensemble des histoires enrichies à la profondeur
\(m\), et soit
\[
 p_{n,m}:\mathcal H_n^\pi\longrightarrow\mathcal H_m^\pi
\]
la troncature. Nous définissons
\[
 \operatorname{Surv}_m^{(N),\pi}
 :=p_{N,m}(\mathcal H_N^\pi),
\]
\[
 \operatorname{Surv}_m^\pi
 :=\bigcap_{N\geq m}\operatorname{Surv}_m^{(N),\pi}.
\]
Lorsque la relation \(\operatorname{Ext}\) est non vide, univoque et
compatible à tous les niveaux, les ensembles sont des singletons et satisfont
\[
 p_{m+1,m}(\operatorname{Surv}_{m+1}^{\pi})
 =\operatorname{Surv}_{m}^{\pi}.
\]
Dans ce domaine de prolongeabilité, il existe donc une unique section
\[
 x_\pi\in
 X_\pi:=\varprojlim_m\operatorname{Surv}_m^\pi.
\]
Les cylindres sont emboîtés et leurs diamètres sont \(3^{-6K}\), qui tend vers
zéro. Leur intersection contient un unique point, identifié par le caractère
angulaire à \(r_\pi=\pi-3\).

Pour toute profondeur décimale \(M\), on peut choisir un niveau \(K(M)\) dont le
cylindre est contenu dans une seule cellule décimale de diamètre \(10^{-M}\).
Les préfixes profonds se tronquent exactement aux précédents.

\section{Caractérisation structurelle, conditions nécessaires de validité et obstructions}

\begin{theorem}[Caractérisation du canal de clôture]
\label{gfp:b:thm:caracterizacion-final-pi-autonoma}
À partir du catalogue, de l’action \(D_3\), de la jauge orientée, des
opérateurs \(L_0,L_1,A_W\), de la loi des pentes et de la filtration rationnelle :
\begin{enumerate}
 \item on sélectionne \(w_6^\pi=010211\) ;
 \item on construit cinq régions avec les identités de
       \eqref{gfp:b:eq:descomposiciones-cinco-regiones-pi} ;
 \item chaque région appartient à une clôture minimale de longueur huit ;
 \item \(A_W^2=-I_6\) fixe l’inversion orientée ;
 \item la symétrisation régionale force \(1/5\) ;
 \item la composition force \(120/119\) ;
 \item le retour unitaire force \(239\) ;
 \item l’identité gaussienne reconnaît le premier demi-tour ;
 \item les encadrements et la fibre de \(729\) éléments produisent une section
       compatible à toute profondeur finie.
\end{enumerate}
Son évaluation est
\[
 \boxed{\pi=16\arctan\frac15-4\arctan\frac1{239}.}
\]
\end{theorem}

La chaîne complète est
\[
 \mathrm{APP}\to\mathrm{TRIT}\to\mathrm{TPK}
 \to104\,976\to468\to243\to w_6^\pi
 \to w_{30}^\pi\to R_{36}\to\Gamma_9
 \to X_\pi\to\pi.
\]

\begin{remark}[Certificat et contrôle de validité : réfutation du canal de clôture]
La dérivation est réfutée si :
\begin{enumerate}
 \item l’orbite de clôture cesse d’être unique ;
 \item l’une des cinq régions disparaît ou son identité complète change ;
 \item l’une des cinq clôtures a une longueur inférieure à huit ;
 \item \(A_W^2\ne2I_6\);
 \item le symétriseur admet une autre distribution normalisée ;
 \item les compositions ne produisent pas \(5/12\), \(120/119\) et \(239\) ;
 \item l’identité gaussienne échoue ;
 \item un encadrement cesse de contenir le caractère angulaire ;
 \item un niveau ne parcourt pas les \(729\) éléments de la fibre ambiante ou ne
       localise pas exactement une extension ;
 \item un développement cible intervient avant la reconnaissance.
\end{enumerate}
\end{remark}

% Unidad U015. Acoplamiento de horquillas racionales, refinamiento TPK y
% certificación a profundidad arbitraria.

\label{gfp:b:chap:precision-arbitraria-constantes}

Cette unité a le statut de certification ultérieure. Les préfixes et leurs
extensions compatibles procèdent de la relation interne
\(\operatorname{Ext}\), de la survie coinductive et des lois de
transition APP--TRIT--TPK. Les encadrements rationnels qui suivent ne choisissent pas
parmi les \(729\) éléments d’une fibre : ils localisent dans la carte archimédienne
l’extension déjà engendrée et certifient sa précision à une profondeur arbitraire.

\section{Séparation typée des \(4\) indices du prolongement nonadique}

Pour chaque caractère \(\chi\in\{\pi,e,\varphi\}\) interviennent quatre
paramètres de natures différentes :
\begin{enumerate}
 \item \(n\), ordre de l’encadrement analytique rationnel ;
 \item \(K\), niveau de raffinement du Topological Phase Kernel ;
 \item \(L_K=6K\), longueur ternaire du préfixe publié ;
 \item \(M\), nombre de positions décimales demandé.
\end{enumerate}
On ne postule ni \(n=K\), ni \(n=L_K\), ni \(K=M\), ni une proportionnalité fixe
entre eux. Le couplage est construit par des inclusions d’intervalles
décidées par arithmétique entière.

Les parties entières sont
\[
 m_\pi=3,
 \qquad
 m_e=2,
 \qquad
 m_\varphi=1,
\]
et les caractères fractionnaires sont
\[
 \rho_\pi=\pi-3,
 \qquad
 \rho_e=e-2,
 \qquad
 \rho_\varphi=\varphi-1.
\]
Les trois familles d’encadrements semi-ouverts s’écrivent
\[
 J_{\chi,n}=[\ell_{\chi,n},h_{\chi,n}).
\]
Leurs extrémités sont rationnelles, contiennent \(\rho_\chi\) et satisfont
\[
 \ell_{\chi,n}\nearrow\rho_\chi,
 \qquad
 h_{\chi,n}\searrow\rho_\chi,
 \qquad
 h_{\chi,n}-\ell_{\chi,n}\longrightarrow0.
\]
Pour \(\pi\), les extrémités proviennent des sommes alternées de l’identité
de Machin démontrée dans le \cref{gfp:b:chap:biografia-pi-autonoma} ; pour \(e\),
des sommes factorielles et de leur reste rationnel du
\cref{gfp:b:chap:biografia-completa-e} ; pour \(\varphi\), des quotients de
Fibonacci et de Cassini du \cref{gfp:b:chap:biografia-completa-phi}.

\section{Cylindres ternaires et fibre ambiante immédiate}

Un préfixe ternaire de longueur \(L_K=6K\) est représenté par un entier
\(N_K\) avec
\[
 0\leq N_K<3^{L_K}.
\]
Son cylindre semi-ouvert est
\begin{equation}
 C(N_K,L_K)=
 \left[
 \frac{N_K}{3^{L_K}},
 \frac{N_K+1}{3^{L_K}}
 \right).
 \label{gfp:b:eq:cilindro-nivel-K}
\end{equation}
L’émission suivante ajoute six trits. La fibre ambiante immédiate contient
\(3^6=729\) éléments, indexés par \(u\in\{0,\ldots,728\}\) :
\begin{equation}
 C_{K,u}=
 \left[
 \frac{729N_K+u}{3^{L_K+6}},
 \frac{729N_K+u+1}{3^{L_K+6}}
 \right).
 \label{gfp:b:eq:729-subcilindros}
\end{equation}
Les \(C_{K,u}\) sont disjoints et
\[
 C(N_K,L_K)=\bigsqcup_{u=0}^{728}C_{K,u}.
\]
Le nombre \(729\) compte les coordonnées de cette fibre ambiante ; il ne compte
ni les germes, ni les émissions distinctes, ni les orbites, ni les états survivants.

\subsection{Comparaison croisée des numérateurs et des dénominateurs dans l’encadrement rationnel}

Soit \(J=[\ell,h)\subset C(N_K,L_K)\), avec
\(\ell,h\in\mathbb Q\). Nous écrivons \(S_K=3^{L_K+6}\) et définissons
\begin{align}
 j_{\min}(J,K)
 &=\max\!\left(729N_K,\left\lfloor S_K\ell\right\rfloor\right),
 \label{gfp:b:eq:jmin-diagonal}\\
 j_{\max}(J,K)
 &=\min\!\left(729N_K+728,\left\lceil S_Kh\right\rceil-1\right).
 \label{gfp:b:eq:jmax-diagonal}
\end{align}

\begin{lemma}[Critère de localisation unitaire]
\label{gfp:b:lem:criterio-localizacion-unitaria}
L’encadrement \(J\) est contenu dans un unique élément de la fibre
\eqref{gfp:b:eq:729-subcilindros} si et seulement si
\begin{equation}
 \boxed{j_{\min}(J,K)=j_{\max}(J,K).}
 \label{gfp:b:eq:jmin-igual-jmax}
\end{equation}
Dans ce cas, si \(j\) est la valeur commune,
\[
 u=j-729N_K,
 \qquad
 N_{K+1}=j.
\]
\end{lemma}

\begin{proof}
Les entiers \(j\) dont les intervalles élémentaires de longueur \(S_K^{-1}\)
rencontrent \(J\) forment l’intervalle entier
\[
 \left[
 \lfloor S_K\ell\rfloor,
 \lceil S_Kh\rceil-1
 \right]\cap
 [729N_K,729N_K+728].
\]
Celui-ci contient exactement un élément si et seulement si ses extrémités coïncident.
L’égalité identifie le numérateur du sous-cylindre et, en tronquant ses six
derniers trits, on retrouve \(N_K\).
\end{proof}

\section{Certification diagonale minimale}

La génération ne fixe pas à l’avance un même ordre analytique pour tous les
niveaux. Une fois l’extension compatible produite, l’ordre de
l’encadrement n’est augmenté que jusqu’à sa certification unitaire.

\begin{definition}[Récurrence diagonale de certification]
\label{gfp:b:def:recurrencia-diagonal-constantes}
Une fois fixée par \(\operatorname{Ext}\) une famille de cylindres compatibles
\(C(N_K,L_K)\) dont la publication contient \(\rho_\chi\), on choisit un ordre
initial \(n_\chi(K_0-1)\geq1\) et nous définissons récursivement
\begin{equation}
 n_\chi(K)=
 \min\left\{
 n\geq n_\chi(K-1):
 j_{\min}(J_{\chi,n},K)=j_{\max}(J_{\chi,n},K)
 \right\},
 \label{gfp:b:eq:nchi-K-minimo}
\end{equation}
et nous vérifions
\begin{equation}
 N_{\chi,K+1}
 =j_{\min}(J_{\chi,n_\chi(K)},K).
 \label{gfp:b:eq:Nchi-K-mas-uno}
\end{equation}
L’ensemble
\[
 \Delta_\chi=
 \{(n_\chi(K),K):K\geq K_0\}
\]
est la diagonale analytique--nonadique du caractère \(\chi\).
\end{definition}

\begin{theorem}[Existence et unicité de la diagonale de certification]
\label{gfp:b:thm:existencia-unicidad-diagonal}
Pour \(\chi\in\{\pi,e,\varphi\}\), la récurrence
\eqref{gfp:b:eq:nchi-K-minimo} est définie à tout niveau fini. L’indice
\(n_\chi(K)\) et la coordonnée \(N_{\chi,K+1}\) sont uniques.
\end{theorem}

\begin{proof}
Supposons engendrée intérieurement l’extension cylindrique de niveau \(K\) qui contient
\(\rho_\chi\). Comme \(\rho_\chi\) est irrationnel, il n’appartient à aucune
frontière ternaire \(3^{-L}\mathbb Z\). Il est donc à l’intérieur d’un
unique sous-cylindre \(C_{K,u_*}\), à une distance positive de ses deux extrémités.
Les encadrements \(J_{\chi,n}\) contiennent \(\rho_\chi\) et leur diamètre
tend vers zéro ; il existe donc un \(n\) pour lequel l’encadrement est
contenu dans \(C_{K,u_*}\). L’ensemble sur lequel est pris le minimum dans
\eqref{gfp:b:eq:nchi-K-minimo} est non vide et, par le bon ordre de
\(\mathbb N\), possède un plus petit élément. Le
\cref{gfp:b:lem:criterio-localizacion-unitaria} prouve que l’encadrement la
localise de manière unitaire et que le numérateur localisé coïncide avec celui déjà
engendré. La récurrence sur \(K\) complète l’argument.
\end{proof}

\begin{corollary}[Absence de paramètres décimaux]
La diagonale est calculée avec des sommes, des produits, des puissances, des divisions euclidiennes,
des parties entières inférieures et supérieures d’entiers. Elle ne reçoit aucun chiffre décimal de \(\pi\), de \(e\) ou de
\(\varphi\) comme donnée de certification.
\end{corollary}

\section{Compatibilité avec la connexion nonadique conjointe}

Soit \(\widehat x_{\chi,K}\) l’état enrichi du canal \(\chi\) au
niveau \(K\). Ses coordonnées incluent le préfixe, le cylindre, la phase,
le quotient, la retenue, la mémoire verticale non bornée, le terme terminal
et le registre de provenance. La phase résiduelle est
\[
 g_0(K)=[K]_9\in\mathbb Z/9\mathbb Z,
\]
tandis que l’étiquette publique est
\[
 g(K)=1+((K-1)\bmod9)\in\{1,\ldots,9\}.
\]
La mémoire verticale \(q_{\rm mem}(K)\in\mathbb Z_{\geq0}\) n’est pas réduite
modulo \(12\) ; seule l’est sa projection dodécaphasique
\[
 \beta(K)=q_{\rm mem}(K)\bmod12
\]
le fait.

Pour chaque \(K\), soit
\[
 \Gamma_{9,K}:\widehat X_K\longrightarrow\widehat X_{K+9}
\]
l’holonomie d’un tour de phase. Alors
\[
 g_0(K+9)=g_0(K),
 \qquad
 q_{\rm mem}(K+9)=q_{\rm mem}(K)+1,
 \qquad
 \beta(K+9)=\beta(K)+1\pmod{12}.
\]

\begin{theorem}[Retour de phase avec conservation généalogique]
\label{gfp:b:thm:retorno-fase-conservacion-genealogica}
L’égalité de phase ne réinitialise pas l’état :
\[
 g_0(K+9)=g_0(K),
 \qquad
 \widehat x_{\chi,K+9}\ne\widehat x_{\chi,K}.
\]
Les profondeurs \(27\), \(54\) et \(108\) sont les compositions typées
\begin{align*}
 \Gamma_{27,K}
 &=\Gamma_{9,K+18}\circ\Gamma_{9,K+9}\circ\Gamma_{9,K},\\
 \Gamma_{54,K}
 &=\Gamma_{9,K+45}\circ\cdots\circ\Gamma_{9,K},\\
 \Gamma_{108,K}
 &=\Gamma_{9,K+99}\circ\cdots\circ\Gamma_{9,K}.
\end{align*}
En aucun cas il ne s’agit de réinitialisations du préfixe, de la retenue, du terminal ou de la
mémoire.
\end{theorem}

\begin{proof}
La première coordonnée est calculée dans \(\mathbb Z/9\mathbb Z\), tandis que
\(q_{\rm mem}\) appartient à un monoïde non borné. Chaque application de
\(\Gamma_{9,K}\) ajoute une unité de mémoire, six blocs par niveau
intermédiaire et la composition correspondante des cocycles. C’est pourquoi la phase
peut coïncider sans que l’état enrichi coïncide. Les formules de
profondeur montrent les indices de chaque application et évitent de les composer comme si toutes
étaient des endomorphismes du même espace.
\end{proof}

\subsection{Compression nonadique, évaluation archimédienne et terme terminal du cylindre}

La réalisation opératorielle d’un tour satisfait
\begin{equation}
 U_\Gamma J=JT+\eta,
 \qquad
 J^*\eta=0,
 \qquad
 I-T^*T=\eta^*\eta.
 \label{gfp:b:eq:compresion-expresion-precision}
\end{equation}
La contraction visible stricte est réservée à
\[
 M_9=\frac59V_9;
\]
n’est pas identifié à l’opérateur de compression \(T\). Le télescopage de
profondeur \(L\) conserve le terme terminal :
\begin{equation}
 S_0=
 \sum_{r=0}^{L-1}M_9^{*r}D_9^*D_9M_9^r
 +M_9^{*L}S_0M_9^L.
 \label{gfp:b:eq:telescopia-terminal-precision}
\end{equation}
Omettre le dernier terme avant de prouver sa disparition effacerait
de l’information de frontière et transformerait le retour de phase en une fausse
réinitialisation.

\section{Système inverse des histoires compatibles}

Soit \(\mathcal H_K^\chi\) l’ensemble des histoires enrichies admissibles
du canal \(\chi\) à la profondeur \(K\), et soit
\[
 p_{N,K}:\mathcal H_N^\chi\longrightarrow\mathcal H_K^\chi
 \qquad(N\geq K)
\]
la troncature. Nous définissons
\[
 \operatorname{Surv}_K^{(N)}(\chi)
 =p_{N,K}(\mathcal H_N^\chi)
\]
et
\begin{equation}
 \operatorname{Surv}_K(\chi)
 =\bigcap_{N\geq K}\operatorname{Surv}_K^{(N)}(\chi).
 \label{gfp:b:eq:supervivencia-todas-profundidades}
\end{equation}
L’objet coinductif est
\begin{equation}
 X_\chi=
 \varprojlim_K
 \bigl(\operatorname{Surv}_K(\chi),p_{K+1,K}\bigr).
 \label{gfp:b:eq:limite-inverso-caracteres}
\end{equation}

\begin{theorem}[Section compatible sous prolongeabilité enrichie]
\label{gfp:b:thm:seccion-compatible-caracteres}
Supposons en outre que, sur chaque sous-cylindre numérique engendré, la
relation \(\operatorname{Ext}\) soit non vide, univoque et compatible avec les
troncatures. Alors \(\operatorname{Ext}\) et la survie produisent une famille
\[
 x_\chi=(\widehat x_{\chi,K})_{K\geq K_0}
\]
qui satisfait
\[
 p_{K+1,K}(\widehat x_{\chi,K+1})
 =\widehat x_{\chi,K}.
\]
Par conséquent, \(x_\chi\in X_\chi\).
\end{theorem}

\begin{proof}
Le sous-cylindre de niveau \(K+1\) a pour numérateur
\(N_{\chi,K+1}=729N_{\chi,K}+u\). La troncature élimine ses six derniers
trits et retrouve \(N_{\chi,K}\). Les autres coordonnées sont obtenues par les
lois fonctorielles de retenue, de mémoire, de terminal et de phase. L’hypothèse
d’univocité de \(\operatorname{Ext}\) fixe le sous-cylindre et l’extension
enrichie ; l’unicité locale du
\cref{gfp:b:lem:criterio-localizacion-unitaria} certifie ensuite sa
localisation archimédienne, et la compatibilité de \(\operatorname{Ext}\)
fournit l’égalité de troncature.
\end{proof}

\section{Cellules décimales et précision prescrite}

Pour \(M\geq1\) et \(q\in\{0,\ldots,10^M-1\}\), soit
\[
 D_{M,q}=
 \left[
 \frac{q}{10^M},
 \frac{q+1}{10^M}
 \right).
\]
Un encadrement \(J=[\ell,h)\) est contenu dans une unique cellule décimale si
et seulement si
\begin{equation}
 \boxed{
 \left\lfloor10^M\ell\right\rfloor
 =\left\lceil10^Mh\right\rceil-1=q.}
 \label{gfp:b:eq:criterio-celda-decimal}
\end{equation}

\begin{definition}[Niveau minimal de certification décimale]
Pour \(\chi\in\{\pi,e,\varphi\}\), nous définissons \(K_\chi(M)\) comme le plus petit
niveau pour lequel il existe un ordre analytique \(n\) et un entier \(q\) tels
que
\begin{enumerate}
 \item \(J_{\chi,n}\) satisfait
       \eqref{gfp:b:eq:criterio-celda-decimal} ;
 \item \(J_{\chi,n}\) est contenu dans le cylindre engendré par
       l’état \(\widehat x_{\chi,K}\) ;
 \item la transition vers le niveau suivant satisfait
       \eqref{gfp:b:eq:jmin-igual-jmax}.
\end{enumerate}
La paire \((K,n)\) est minimisée lexicographiquement.
\end{definition}

\begin{theorem}[Certification à précision décimale arbitraire des sorties coinductives]
\label{gfp:b:thm:precision-decimal-arbitraria}
Pour chaque \(\chi\in\{\pi,e,\varphi\}\) et chaque \(M\geq1\), il existe
\(K_\chi(M)<\infty\). La section engendrée détermine un unique préfixe décimal
de \(M\) positions, et l’encadrement reconnaît et certifie ce préfixe. Si
\(M'<M\), le préfixe d’ordre \(M\) se tronque à celui
d’ordre \(M'\), et les cellules correspondantes sont emboîtées.
\end{theorem}

\begin{proof}
Les trois caractères sont irrationnels : \(\varphi\) l’est par son polynôme
quadratique à discriminant non carré ; l’irrationalité de \(e\) et de
\(\pi\) est tirée de leurs théorèmes classiques de caractérisation, postérieurs à
la construction finie. Aucun n’appartient donc à une frontière décimale ou
ternaire rationnelle. Pour un \(M\) fixé, la distance du caractère à l’ensemble
fini des frontières décimales pertinentes est positive. Les encadrements
rationnels contractent leur diamètre jusqu’à être contenus dans une unique cellule décimale.
Le \cref{gfp:b:thm:existencia-unicidad-diagonal} couple cet encadrement à un
cylindre TPK et à une extension unique. Le minimum lexicographique existe par
bon ordre. La compatibilité de la section prouve l’emboîtement lorsque
\(M\) augmente.
\end{proof}

\begin{corollary}[Unicité réelle]
Les cellules emboîtées ont des diamètres \(10^{-M}\to0\). Leur intersection est
un singleton, identifié par les caractérisations précédentes à
\(\pi\), \(e\) ou \(\varphi\), respectivement.
\end{corollary}

\section{Équivalence de \(2\) réalisations finies de la loi uniforme}

Les témoins de \(1,000\) et de \(10,000\) positions ne définissent pas les
caractères et ne limitent pas le prolongement. Ce sont des exécutions de la même récurrence
uniforme à deux profondeurs différentes.

\begin{table}[htbp]
\centering
\caption{Dimensions exactes des exécutions certifiées, par canal.}
\label{gfp:b:tab:dimensiones-ejecuciones-1000-10000}
\begin{tabular}{lrr}
\toprule
Objet compté & \(1,000\) positions & \(10,000\) positions\\
\midrule
Partitions exhaustives de la fibre ambiante & 396 & 3501\\
Trits transportados & 2376 & 21006\\
Paires de même phase \((K,K+9)\) & 387 & 3492\\
Tours nonadiques complets & 43 & 388\\
Éléments examinés par partition & 729 & 729\\
Extensions compatibles conservées par partition & 1 & 1\\
\bottomrule
\end{tabular}
\end{table}

Para \(10,000\) posiciones,
\[
 B=3501=389\cdot9,
 \qquad
 6B=21006,
 \qquad
 3^{6B}>10^{10000}.
\]
Les \(396\) premiers blocs de l’exécution longue coïncident avec
l’exécution de \(1,000\) positions. Dans chacune des \(3492\) paires de
même phase, la coordonnée de phase est préservée et le préfixe, la
profondeur, le cylindre et la mémoire sont prolongés. Le certificat enregistre aussi les
répétitions éventuelles de projections locales sans les confondre avec un retour
de l’état complet.

\section{Séparation entre génération et vérification}

\begin{definition}[Générateur structurel]
Le générateur reçoit le catalogue fini, \(L_0,L_1\), l’état initial, les
lois de connexion et la relation interne \(\operatorname{Ext}\). À chaque niveau :
\begin{enumerate}
 \item il énumère les \(729\) éléments de la fibre ambiante ;
 \item il applique \(\operatorname{Ext}\) et conserve l’unique extension compatible ;
 \item il actualise la phase, le quotient, la retenue, la mémoire et le terminal ;
 \item il écrit le bloc ternaire, le reçu de sélection et le registre généalogique.
\end{enumerate}
\end{definition}

\begin{definition}[Vérificateur ultérieur]
Le vérificateur reçoit les sorties closes du générateur et les trois familles
d’encadrements rationnels ; il calcule \(j_{\min}\) et \(j_{\max}\), certifie leur
égalité avec l’extension déjà produite, reconstruit les
partitions, vérifie les carrés de troncature, les paires
\((K,K+9)\), les sommes de partition
\[
 \operatorname{killed}_{\mathrm{izq}}+1+\operatorname{killed}_{\mathrm{der}}=729,
\]
les identités de contenu et, à la fin, la coïncidence avec des développements de référence.
\end{definition}

\begin{theorem}[Isolement des données cibles]
\label{gfp:b:thm:aislamiento-datos-objetivo-precision}
Les développements décimaux de référence et les données métrologiques
n’interviennent pas dans la sélection d’une extension compatible. La comparaison est
effectuée après l’écriture des trois préfixes.
\end{theorem}

\begin{proof}
L’entrée mathématique de chaque décision générative consiste dans le catalogue
fini, \(L_0,L_1\), l’état hérité et la relation interne
\(\operatorname{Ext}\). Les extrémités rationnelles, les puissances de \(3\), les parties entières inférieures et
supérieures appartiennent au vérificateur ultérieur. Le générateur termine avant que
celui-ci ouvre les encadrements ou les contrôles décimaux. Le sens du flux de
données est
\[
 \text{structure}\longrightarrow
 \text{préfixes et registres}\longrightarrow
 \text{vérification ultérieure};
\]
il n’existe pas de flèche en sens contraire.
\end{proof}

\section{Réalisation finie du certificat de précision arbitraire}

Une réalisation matérielle de la procédure produit, pour chacun des trois caractères, un préfixe de \(10^4\) positions. Cette coupe finie est un contrôle ultérieur de la construction et n’intervient pas dans le choix des cylindres, des retenues ni des extensions compatibles. La proposition démontrée conserve une portée arbitraire en \(M\) ; la précision \(10^4\) certifie une exécution particulière, mais ne limite pas le générateur coinductif.

\section{Obstructions différenciées de la génération et de sa certification à précision arbitraire}

\begin{remark}[Certificat et contrôle de validité : réfutation de la génération à précision arbitraire]
Les échecs d’orbite, d’extension, de troncature ou de retour réfutent la génération ;
les échecs d’encadrement, d’ordre ou de contrôle décimal réfutent le certificat
analytique correspondant ; et une divergence entre les deux réfute
l’identification conjointe de la sortie publiée. Avec cette séparation des types,
l’affirmation complète est réfutée si :
\begin{enumerate}
 \item on identifie l’un quelconque des indices \(n,K,L_K,M\) sans un
       théorème de liaison ;
 \item un encadrement cesse de contenir son caractère ou ne contracte pas son diamètre ;
 \item une étape satisfait \(j_{\min}\ne j_{\max}\) après l’ordre
       déclaré minimal ;
 \item une extension ne se tronque pas au cylindre du niveau précédent ;
 \item le retour de phase réinitialise le préfixe, la retenue, le terminal ou la mémoire ;
 \item le terme terminal est omis dans
       \eqref{gfp:b:eq:telescopia-terminal-precision} sans que son annulation soit démontrée ;
 \item les \(396\) premiers blocs de l’exécution longue diffèrent de
       l’exécution courte ;
 \item un contrôle décimal ou métrologique est lu avant la clôture de la sortie ;
 \item un chiffre certifié diffère de la sortie matérielle du générateur ;
 \item la preuve dépend du fait que la précision demandée soit \(1,000\) ou
       \(10,000\), au lieu d’un \(M\) arbitraire.
\end{enumerate}
\end{remark}

%
\endgroup

\HMTFlushCurrentChapter
\chapter{Génération coinductive de \texorpdfstring{\(e\)}{e}: caractère exponentiel, semi-groupe factoriel et retour de phase à mémoire d'état}%
\begingroup
\renewcommand{\chapter}[2][]{}%
\chapter{Génération coinductive de \(e\) : caractère exponentiel, semi-groupe factoriel et retour de phase avec mémoire d’état}
\label{chap:p3-final:29:generacion_e}
% Unidad U013. Biografía estructural completa del carácter exponencial.

\label{gfp:b:chap:biografia-completa-e}

\section{Sélecteur orbital fini : existence et unicité de l’émission compatible}

Soit \(\mathcal W=\mathbb F_3^6\). Pour une émission
\(U=(u_1,\ldots,u_6)\), le mot visible est
\[
 \Psi(U)=((-u_1)\bmod3,\ldots,(-u_6)\bmod3).
\]
Sur \(\mathcal W\), l’action diédrale est donnée par
\[
\begin{aligned}
 r(u_1,u_2,u_3,u_4,u_5,u_6)
 &=(u_3,u_1,u_2,u_5,u_6,u_4),\\
 s(u_1,u_2,u_3,u_4,u_5,u_6)
 &=(u_4,u_5,u_6,u_1,u_2,u_3).
\end{aligned}
\]
Pour \(w\in\mathcal W\), nous définissons
\[
 n(w)=\#\Psi^{-1}(w),
 \qquad
 M(w)=\sum_{U\in\Psi^{-1}(w)}\operatorname{mult}(U),
\]
et
\[
 \nu(w)=
 \bigl(
 \#\{j:w_j=0\},
 \#\{j:w_j=1\},
 \#\{j:w_j=2\}
 \bigr).
\]

\begin{theorem}[Sélection structurelle du canal exponentiel]
\label{gfp:b:thm:seleccion-estructural-e}
Parmi les \(43\) orbites du catalogue visible, il existe une unique orbite dont les
fibres satisfont simultanément :
\[
 n(w)=1,\qquad M(w)=144,\qquad
 \mathfrak D(w)=\{0,3,6\},\qquad
 \nu(w)=(2,3,1).
\]
Le calibre
\[
 \operatorname{diag}_c=6,
 \qquad \operatorname{card}=ES
\]
y sélectionne
\[
 \boxed{w_6^e=201101.}
\]
\end{theorem}

\begin{proof}
Les quatre conditions sont évaluées sur les \(243\) mots du catalogue,
non sur un échantillon. Une unique orbite satisfait le prédicat conjoint. Dans
cette orbite, une seule émission possède à la fois la diagonale \(6\) et l’orientation
\(ES\). L’énumération utilise uniquement des champs du catalogue APP--TRIT et ne
consulte pas de développement de \(e\).
\end{proof}

\section{Relèvement par blocs et compatibilité des frontières successives}

La récurrence \(L_0,L_0,L_1,L_1\) produit
\[
 w_{30}^{e}
 =201101|121221|102011|012222|102011,
\]
et la frontière conjointe apporte
\[
 u_5^e=021222.
\]
Le registre de trente-six trits devient
\[
 w_{36}^{e}
 =201101|121221|102011|012222|102011|021222.
\]

\section{Signature positionnelle décimale et carré de compatibilité}

La lecture initiale de douze chiffres est
\[
 \mathbf d_e^{(12)}
 =718281828459
 =7\,\underbrace{1828\,1828}_{\text{carré local}}\,459.
\]
L’égalité centrale est la concaténation
\[
 1828\mathbin{\|}1828,
\]
non un produit arithmétique ni le début supposé d’une période.

\begin{definition}[Carré positionnel]
Un mot \(d_1\cdots d_{12}\) contient un carré de longueur \(\ell\)
à la position \(p\) lorsque
\[
 d_p\cdots d_{p+\ell-1}
 =d_{p+\ell}\cdots d_{p+2\ell-1}.
\]
La position et les contextes latéraux appartiennent à la signature.
\end{definition}

\begin{table}[htbp]
\centering
\caption{Recensement exhaustif des carrés dans les lectures de douze chiffres.}
\label{gfp:b:tab:censo-cuadrados-e}
\begin{tabular}{c*{6}{r}}
\toprule
Longueur \(\ell\)&1&2&3&4&5&6\\
\midrule
Occurrences&240&21&3&2&0&0\\
Mots qui les contiennent&153&19&3&2&0&0\\
\bottomrule
\end{tabular}
\end{table}

L’autre mot contenant un carré de longueur quatre est
\[
 575157518472
 =\underbrace{5751\,5751}_{\text{position initiale }1}\,8472.
\]

\begin{proposition}[Unicité positionnelle]
\label{gfp:b:prop:cuadrado-posicional-e}
La lecture du canal \(e\) est la seule qui possède un carré de longueur
quatre après un préfixe de longueur un et avec les contextes \(7\) et \(459\).
\end{proposition}

\begin{proof}
Le recensement ne laisse que deux mots contenant des carrés de longueur quatre. L’un
commence son carré à la première position ; l’autre après le préfixe \(7\). La
position et les contextes séparent les deux possibilités.
\end{proof}

\section{Retour de phase avec changement d’état et conservation de la mémoire parcourue}

Dans le relèvement
\[
 201101|121221|102011|012222|102011
\]
les troisième et cinquième blocs coïncident :
\[
 B_3=B_5=102011.
\]
Leurs préétats modulo \(27\) sont cependant
\[
 p_3=(25,11,7,17,8,16),
 \qquad
 p_5=(7,8,4,23,26,7),
\]
et les quotients ultérieurs :
\[
 q_3=(6,13,9,2,13,1),
 \qquad
 q_5=(15,0,3,19,23,25).
\]
On vérifie \(p_3\ne p_5\) et \(q_3\ne q_5\).

\begin{figure}[htbp]
\centering
\begin{tikzpicture}[
 node distance=11mm and 7mm,
 every node/.style={font=\small},
 block/.style={draw,rounded corners,minimum width=17mm,minimum height=8mm},
 state/.style={draw,dashed,rounded corners,align=center,inner sep=3pt},
 >=stealth
]
\node[block] (b1) {\(201101\)};
\node[block,right=of b1] (b2) {\(121221\)};
\node[block,right=of b2] (b3) {\(102011\)};
\node[block,right=of b3] (b4) {\(012222\)};
\node[block,right=of b4] (b5) {\(102011\)};
\draw[->] (b1)--(b2);
\draw[->] (b2)--(b3);
\draw[->] (b3)--(b4);
\draw[->] (b4)--(b5);
\node[state,below=of b3] (p3)
 {\(p_3=(25,11,7,17,8,16)\)\\
  \(q_3=(6,13,9,2,13,1)\)};
\node[state,below=of b5] (p5)
 {\(p_5=(7,8,4,23,26,7)\)\\
  \(q_5=(15,0,3,19,23,25)\)};
\draw[->] (p3)--(b3);
\draw[->] (p5)--(b5);
\draw[<->,thick] (b3) to[bend left=28]
 node[above]{même mot visible} (b5);
\end{tikzpicture}
\caption{Retour visible de \(102011\) avec renouvellement du préétat et du
quotient.}
\label{gfp:b:fig:ruptura-memoria-e}
\end{figure}

\begin{theorem}[Retour visible sans périodicité de l’état]
\label{gfp:b:thm:retorno-visible-no-periodico-e}
\(B_3=B_5\) est un retour de la projection visible, non une orbite périodique
de l’état enrichi.
\end{theorem}

\begin{proof}
Le retour de l’état complet exigerait l’égalité de toutes ses coordonnées.
Les inégalités du préétat et du quotient l’excluent. De même,
\(1828|1828\) ne commence pas une période décimale : après le deuxième bloc apparaît
\(4\), tandis qu’une continuation périodique exigerait \(1\).
\end{proof}

\section{Certification factorielle ultérieure de la section exponentielle coinductive}

Soit \(\mathcal A\) une algèbre commutative unitaire sur \(\mathbb Q\), et
soit \(P_t\in\mathcal A[[t]]\) le propagateur formel. La concaténation
temporelle impose
\[
 P_{s+t}=P_sP_t,
 \qquad
 P_0=1.
\]
Écrivons
\[
 P_t=\sum_{n=0}^{\infty}a_nt^n.
\]
L’unité fixe \(a_0=1\), et le transport élémentaire normalisé fixe
\(a_1=1\).

\begin{theorem}[Récurrence factorielle]
\label{gfp:b:thm:recurrencia-factorial-e}
\[
 (n+1)a_{n+1}=a_n
 \qquad(n\geq0),
\]
et, par conséquent,
\[
 a_n=\frac1{n!},
 \qquad
 P_t=\sum_{n=0}^{\infty}\frac{t^n}{n!}.
\]
En \(t=1\),
\[
 \boxed{P_1=e.}
\]
\end{theorem}

\begin{proof}
Le coefficient de \(s^nt\) dans \(P_{s+t}\) est
\[
 \binom{n+1}{n}a_{n+1}=(n+1)a_{n+1}.
\]
Dans \(P_sP_t\), il est \(a_na_1=a_n\). L’égalité des séries produit la
récurrence. Une récurrence à partir de \(a_0=1\) donne \(a_n=1/n!\).
\end{proof}

\begin{corollary}[Équation différentielle]
\[
 \frac{d}{dt}P_t=P_t,
 \qquad P_0=1.
\]
\end{corollary}

\begin{proof}
\[
 P_t'=\sum_{n\geq0}(n+1)a_{n+1}t^n
 =\sum_{n\geq0}a_nt^n=P_t.
\]
\end{proof}

\section{Emboîtement cofinal des encadrements rationnels et convergence vers la valeur publiée}

Soit
\[
 S_N=\sum_{n=0}^{N}\frac1{n!}.
\]

\begin{theorem}[Borne factorielle explicite]
\label{gfp:b:thm:cota-factorial-e}
Pour \(N\geq1\),
\begin{equation}
 \boxed{
 S_N<e<
 S_N+\frac{N+2}{(N+1)(N+1)!}.}
 \label{gfp:b:eq:horquilla-factorial-e}
\end{equation}
\end{theorem}

\begin{proof}
Le reste est
\[
 R_N=\sum_{k=0}^{\infty}\frac1{(N+1+k)!}>0.
\]
Para \(k\geq0\),
\[
 (N+1+k)!
 \geq(N+1)!(N+2)^k.
\]
Par conséquent,
\[
 R_N
 \leq\frac1{(N+1)!}
 \sum_{k=0}^{\infty}\frac1{(N+2)^k}
 =\frac{N+2}{(N+1)(N+1)!}.
\]
L’inégalité est stricte à partir de \(k=2\).
\end{proof}

On définit
\[
 J_N^{(e)}=
 \left[
 S_N,\,
 S_N+\frac{N+2}{(N+1)(N+1)!}
 \right]
\]
et, pour la partie fractionnaire,
\[
 \widetilde J_N^{(e)}=J_N^{(e)}-2.
\]
La largeur converge vers zéro. Pour chaque \(M\), il existe \(N\) tel que
\[
 \operatorname{diam}J_N^{(e)}<10^{-M}.
\]
Si les extrémités appartiennent à une même cellule décimale, les \(M\) premiers
chiffres sont certifiés par arithmétique rationnelle.

\Needspace{30\baselineskip}
\section{Prolongement coinductif de la section exponentielle par les sous-cylindres ternaires du Topological Phase Kernel}

Au niveau \(K\), TPK et la relation \(\operatorname{Ext}\) ont déjà engendré
un unique sous-cylindre compatible. On prend le plus petit \(N=N_e(K)\) pour lequel
l’encadrement factoriel est contenu dans ce sous-cylindre. \(N_e(K)\) est un
ordre de certification analytique, non un indice générateur ni une règle de
sélection de la fibre.

Les cylindres TPK sont compatibles par troncature indépendamment du
semi-groupe. Le théorème du semi-groupe identifie ensuite leur évaluation
archimédienne à \(e-2\) ; la partie entière retrouve \(e\). Si l’orbite,
\(\operatorname{Ext}\) ou la troncature échouent, la génération échoue ; si la
récurrence ou la borne factorielle échouent, c’est ce certificat analytique qui échoue, non la
généalogie APP--TRIT--TPK.

\section{Certificat interne d’unicité orbitale, de récurrence factorielle, de convergence et de prolongement cylindrique}

\begin{remark}[Certificat et contrôle de validité : réfutation du canal exponentiel]
La construction est réfutée si :
\begin{enumerate}
 \item une autre orbite satisfait le prédicat fini complet ;
 \item la signature \(7[1828\,1828]459\) n’est pas positionnellement unique ;
 \item les préétats ou les quotients des deux blocs \(102011\) coïncident ;
 \item le semi-groupe ne force pas \((n+1)a_{n+1}=a_n\) ;
 \item la borne \eqref{gfp:b:eq:horquilla-factorial-e} échoue ;
 \item un niveau ne localise pas de manière unitaire une extension compatible ;
 \item un développement tabulé intervient avant la reconnaissance.
\end{enumerate}
\end{remark}

%
\endgroup

\HMTFlushCurrentChapter
\chapter{Génération coinductive de \texorpdfstring{\(\varphi\)}{phi}: incidence modale, rayon positif de Perron et unicité de l'auto-échelle}%
\begingroup
\renewcommand{\chapter}[2][]{}%
\chapter{Génération coinductive de \(\varphi\) : incidence modale, demi-droite positive de Perron et unicité de l’auto-échelle}
\label{chap:p3-final:30:generacion_phi}
% Unidad U014. Biografía estructural completa del carácter de autoescala.

\label{gfp:b:chap:biografia-completa-phi}

\section{Sélection finie du germe d’auto-échelle}

Soit \(\mathcal U_6\subset\{0,\ldots,999\}^6\) le catalogue des
\(468\) émissions distinctes construit dans le
\cref{gfp:b:chap:semillas-emisor-54}, et soit
\[
 \Psi:\mathcal U_6\longrightarrow\mathbb F_3^6,
 \qquad
 \Psi(U)=((-U_j)\bmod3)_{j=1}^{6}.
\]
Pour \(w\in\Psi(\mathcal U_6)\), rappelons les invariants
\[
 n(w)=\#\Psi^{-1}(w),
 \qquad
 M(w)=\sum_{U\in\Psi^{-1}(w)}\operatorname{mult}(U),
\]
et le vecteur d’occupation
\[
 \nu(w)=
 \bigl(\#\{j:w_j=0\},\#\{j:w_j=1\},\#\{j:w_j=2\}\bigr).
\]

\begin{theorem}[Sélection orbitale de l’auto-échelle]
\label{gfp:b:thm:seleccion-orbital-phi}
L’action diédrale sur \(\Psi(\mathcal U_6)\) contient une unique orbite
qui satisfait le prédicat d’auto-échelle
\[
 n(w)=1,
 \qquad
 M(w)=144,
 \qquad
 \nu(w)=(2,2,2).
\]
La jauge interne
\[
 \operatorname{diag}_c=6,
 \qquad
 \operatorname{card}=ES
\]
sélectionne dans cette orbite le mot
\[
 \boxed{w_6^{\varphi}=121200.}
\]
Son unique émission est
\[
 \boxed{U_{\varphi}=(830,943,830,109,252,252),}
\]
de multiplicité \(144\), de signature multiplicative \(555933\), d’orientation
additive \(ES\), de support multiplicatif de cardinal six et de type résiduel
\(A\).
\end{theorem}

\begin{proof}
La projection directe donne
\[
 U_{\varphi}\bmod3=(2,1,2,1,0,0),
 \qquad
 -U_{\varphi}\bmod3=(1,2,1,2,0,0).
\]
L’énumération est effectuée sur les \(243\) mots visibles et leurs
\(43\) orbites, non sur une sélection d’exemples. Le prédicat conjoint
\(n=1\), \(M=144\), \(\nu=(2,2,2)\) laisse une orbite. En son sein, la
diagonale et l’orientation déclarées laissent une émission. À cette étape,
ni l’équation \(x^2-x-1=0\), ni une table décimale, ni le nom
conventionnel de la constante n’ont été introduits.
\end{proof}

\section{Relèvement par blocs et \(1^{\mathrm a}\) frontière coinductive}

Soit \(b_1=w_6^{\varphi}\), et soit
\(\varepsilon=(0,0,1,1)\). Avec les matrices \(L_0,L_1\) construites dans le
\cref{gfp:b:chap:orbitas-elevacion-r36}, nous définissons
\[
 b_{k+1}=b_kL_{\varepsilon_k}\pmod3
 \qquad(1\leq k\leq4),
\]
et concaténons
\[
 w_{6(k+1)}^{\varphi}=w_{6k}^{\varphi}\mathbin{\|}b_{k+1}.
\]
Le calcul matriciel produit, bloc par bloc,
\[
\begin{aligned}
 b_1&=121200,\\
 b_2&=112202,\\
 b_3&=121020,\\
 b_4&=010210,\\
 b_5&=010200,
\end{aligned}
\]
et, par conséquent,
\begin{equation}
 \boxed{
 w_{30}^{\varphi}
 =121200|112202|121020|010210|010200.}
 \label{gfp:b:eq:w30-phi}
\end{equation}
La neutralité duale de la frontière ajoute
\[
 u_5^{\varphi}=102011,
\]
de sorte que
\begin{equation}
 \boxed{
 w_{36}^{\varphi}
 =121200|112202|121020|010210|010200|102011.}
 \label{gfp:b:eq:w36-phi}
\end{equation}

\begin{remark}[Deux usages distincts d’un mot visible]
Le mot \(102011\) est également visible dans la biographie de \(e\). L’égalité
d’une projection n’identifie pas les états enrichis : le canal,
le préfixe, le quotient, la retenue, la phase et la mémoire appartiennent au type
de l’objet. Cette séparation est obligatoire avant de construire tout
caractère réel.
\end{remark}

\section{Du calendrier tritique au multigraphe d’incidence}

\subsection{Règle modale primitive}

Le calendrier tritique primitif est
\[
 (+1,-1,0).
\]
Nous fixons deux modes visibles, \(\Sigma\) et \(\Pi\), et la règle
\begin{equation}
 +1:m\longmapsto\Sigma,
 \qquad
 -1:m\longmapsto\Pi,
 \qquad
 0:m\longmapsto m.
 \label{gfp:b:eq:regla-modal-primitiva-phi}
\end{equation}
Avec la condition de clôture cyclique, l’orbite modale produite par
\eqref{gfp:b:eq:regla-modal-primitiva-phi} est
\[
 \boxed{(\Sigma,\Pi,\Pi).}
\]
Les trois paires consécutives sont
\[
 \Sigma\longrightarrow\Pi,
 \qquad
 \Pi\longrightarrow\Pi,
 \qquad
 \Pi\longrightarrow\Sigma.
\]

\begin{definition}[Matrice d’incidence chronologique]
Pour \(i,j\in\{\Sigma,\Pi\}\), soit
\[
 N_3(i,j)=
 \#\{r\in\mathbb Z/3\mathbb Z:(m_r,m_{r+1})=(i,j)\}.
\]
L’indice \(3\) compte les instants du calendrier primitif ; il ne désigne pas une
puissance matricielle.
\end{definition}

\begin{theorem}[Descente nécessaire de l’incidence modale]
\label{gfp:b:thm:descenso-incidencia-modal-phi}
Dans l’ordre \((\Sigma,\Pi)\),
\begin{equation}
 \boxed{
 N_3=
 \begin{pmatrix}
 0&1\\
 1&1
 \end{pmatrix}.}
 \label{gfp:b:eq:N3-phi}
\end{equation}
Par répétition exacte du calendrier,
\begin{equation}
 \boxed{
 N_9=3N_3,
 \qquad
 N_{27}=9N_3,
 \qquad
 N_{108}=36N_3.}
 \label{gfp:b:eq:incidencias-9-27-108-phi}
\end{equation}
Ces trois égalités sont des égalités de dénombrements d’arêtes ; elles n’affirment
ni \(N_9=N_3^3\), ni \(N_{27}=N_3^9\), ni \(N_{108}=N_3^{36}\).
\end{theorem}

\begin{proof}
\(\Sigma\to\Sigma\) n’apparaît pas ; chacune des trois autres paires apparaît
une fois. Cela détermine les quatre entrées de \(N_3\). Les calendriers de
longueurs \(9\), \(27\) et \(108\) contiennent respectivement \(3\), \(9\)
et \(36\) copies du bloc primitif. Chaque répétition multiplie tous les
dénombrements par le même entier, ce qui prouve
\eqref{gfp:b:eq:incidencias-9-27-108-phi}.
\end{proof}

\subsection{Publication aréale--volumétrique}

L’étiquette modale et le feuillet géométrique sont des objets distincts. La carte
de publication est
\[
 \mathfrak g_{\rm av}(\Sigma)=\mathscr H_{\rm ar},
 \qquad
 \mathfrak g_{\rm av}(\Pi)=\mathscr H_{\rm vol}.
\]
En transportant \eqref{gfp:b:eq:N3-phi} par \(\mathfrak g_{\rm av}\), nous obtenons
l’opérateur d’incidence
\begin{equation}
 \boxed{
 F_{\rm av}=
 \begin{pmatrix}
 0&1\\
 1&1
 \end{pmatrix}.}
 \label{gfp:b:eq:Fav-derivada-phi}
\end{equation}
Ainsi, les deux transitions croisées, l’unique autorétention volumétrique et
l’absence d’autorétention aréale découlent du calendrier ; ce ne sont pas quatre
entrées indépendantes.

\begin{remark}[Portée de la descente]
L’oubli de la phase élémentaire ne produit pas, à lui seul, un quotient de Markov
fort du Topological Phase Kernel, parce que le mode visible ne détermine pas
lequel des trois opérateurs internes agira au pas suivant. Ce qui
descend en revanche sans hypothèse supplémentaire est le multigraphe des arêtes d’un
bloc complet. Son enveloppe minimale de mémoire un a pour matrice
d’adjacence \(F_{\rm av}\).
\end{remark}

\section{Demi-droite positive invariante et unicité de l’auto-échelle}

\begin{theorem}[Demi-droite de Perron dérivée]
\label{gfp:b:thm:rayo-perron-phi}
Le cône positif de \(F_{\rm av}\) contient une unique demi-droite propre. Si nous la
normalisons comme \(v=(1,x)^{\mathsf T}\), alors
\[
 x^2-x-1=0,
 \qquad
 x>0,
\]
et, par conséquent,
\begin{equation}
 \boxed{x=\frac{1+\sqrt5}{2}=\varphi.}
 \label{gfp:b:eq:phi-perron}
\end{equation}
\end{theorem}

\begin{proof}
L’équation \(F_{\rm av}v=\lambda v\) donne, coordonnée par coordonnée,
\[
 x=\lambda,
 \qquad
 1+x=\lambda x.
\]
En éliminant \(\lambda\), on obtient \(x^2-x-1=0\). Ses racines sont
\((1\pm\sqrt5)/2\), et seule la positive appartient à l’intérieur du cône. La
matrice \(F_{\rm av}\) est primitive, puisque \(F_{\rm av}^2\) a toutes ses
entrées positives ; le théorème de Perron--Frobenius garantit que cette demi-droite
positive est simple et unique. Le nom \(\varphi\) est attribué après cette
construction.
\end{proof}

\begin{corollary}[Équation scalaire d’auto-échelle]
\[
 \varphi=1+\frac1\varphi,
 \qquad
 \varphi^2=\varphi+1,
 \qquad
 \varphi^{-1}=\varphi-1.
\]
Ces identités sont des conséquences de \eqref{gfp:b:eq:phi-perron} ; elles ne sélectionnent
ni le germe ni la matrice.
\end{corollary}

\section{Publication de l’auto-échelle \(\varphi\) par les puissances de Fibonacci et la convergence diophantienne}

Soit \(F_0=0\), \(F_1=1\) et
\[
 F_{n+1}=F_n+F_{n-1}\qquad(n\geq1).
\]

\begin{theorem}[Puissances exactes de l’incidence]
\label{gfp:b:thm:potencias-fibonacci-phi}
Pour tout \(n\geq1\),
\begin{equation}
 \boxed{
 F_{\rm av}^{\,n}=
 \begin{pmatrix}
 F_{n-1}&F_n\\
 F_n&F_{n+1}
 \end{pmatrix}.}
 \label{gfp:b:eq:potencias-Fav-fibonacci}
\end{equation}
\end{theorem}

\begin{proof}
Pour \(n=1\), l’identité est la définition de \(F_{\rm av}\). Si elle vaut
pour \(n\), la multiplication à droite par \(F_{\rm av}\) donne
\[
 \begin{pmatrix}
 F_n&F_{n-1}+F_n\\
 F_{n+1}&F_n+F_{n+1}
 \end{pmatrix}
 =
 \begin{pmatrix}
 F_n&F_{n+1}\\
 F_{n+1}&F_{n+2}
 \end{pmatrix},
\]
ce qui est le cas \(n+1\).
\end{proof}

\begin{theorem}[Identité de Cassini]
\label{gfp:b:thm:cassini-phi}
Pour tout \(n\geq1\),
\begin{equation}
 \boxed{F_{n+1}F_{n-1}-F_n^2=(-1)^n.}
 \label{gfp:b:eq:cassini-phi}
\end{equation}
\end{theorem}

\begin{proof}
Le déterminant de \(F_{\rm av}\) est \(-1\). En prenant les déterminants dans
\eqref{gfp:b:eq:potencias-Fav-fibonacci},
\[
 F_{n-1}F_{n+1}-F_n^2
 =\det(F_{\rm av}^{\,n})=(-1)^n.
\]
\end{proof}

Nous définissons \(r_n=F_{n+1}/F_n\) pour \(n\geq1\). De Cassini découle
\[
 r_n-r_{n-1}
 =\frac{F_{n+1}F_{n-1}-F_n^2}{F_nF_{n-1}}
 =\frac{(-1)^n}{F_nF_{n-1}}.
\]
Par conséquent, les approximations alternent autour de la demi-droite de Perron.

\begin{proposition}[Encadrements rationnels de Perron]
\label{gfp:b:prop:horquillas-perron-phi}
Pour \(m\geq1\),
\begin{equation}
 \frac{F_{2m+2}}{F_{2m+1}}
 <\varphi<
 \frac{F_{2m+1}}{F_{2m}},
 \label{gfp:b:eq:horquilla-fibonacci-phi}
\end{equation}
et la largeur est
\begin{equation}
 \boxed{
 \frac{F_{2m+1}}{F_{2m}}
 -\frac{F_{2m+2}}{F_{2m+1}}
 =\frac{1}{F_{2m}F_{2m+1}}.}
 \label{gfp:b:eq:anchura-fibonacci-phi}
\end{equation}
\end{proposition}

\begin{proof}
L’alternance précédente place les quotients pairs en dessous et les impairs
au-dessus de \(\varphi\). La formule de la largeur est
\eqref{gfp:b:eq:cassini-phi} avec \(n=2m+1\), après réduction des deux fractions au
même dénominateur. Comme \(F_n\to\infty\), la largeur tend vers zéro.
\end{proof}

\section{Incidence de mémoire d’ordre \(1\) et mesure de Parry de l’enveloppe symbolique}

Soit \(X_{\rm av}\) le plus petit sous-décalage de type fini sur
\(\{\mathscr H_{\rm ar},\mathscr H_{\rm vol}\}\) qui admet exactement
les trois paires observées dans le calendrier modal. Sa matrice d’adjacence est
\(F_{\rm av}\).

\begin{theorem}[Transition et mesure de Parry]
\label{gfp:b:thm:parry-phi}
Soit \(r=(1,\varphi)^{\mathsf T}\). La matrice stochastique de Parry est
\begin{equation}
 \boxed{
 P_{ij}=
 \frac{(F_{\rm av})_{ij}r_j}{\varphi r_i}
 =
 \begin{pmatrix}
 0&1\\
 \varphi^{-2}&\varphi^{-1}
 \end{pmatrix}.}
 \label{gfp:b:eq:parry-matriz-phi}
\end{equation}
La mesure stationnaire des sommets est proportionnelle à
\((1,\varphi^2)\), donc
\begin{equation}
 \boxed{
 \frac{\mu_{\rm ar}}{\mu_{\rm vol}}=\varphi^{-2}.}
 \label{gfp:b:eq:parry-razon-phi}
\end{equation}
\end{theorem}

\begin{proof}
La somme de la ligne \(i\) de \(P\) est
\[
 \frac{(F_{\rm av}r)_i}{\varphi r_i}=1.
\]
Comme \(F_{\rm av}\) est symétrique, les vecteurs propres gauche et droit
coïncident. Les poids stationnaires sont proportionnels à leurs produits
coordonnée par coordonnée, c’est-à-dire à \((1,\varphi^2)\). Le rapport entre les
deux poids est \(\varphi^{-2}\).
\end{proof}

\begin{remark}[Trois mesures à ne pas confondre]
La fréquence chronologique de l’orbite primitive est
\[
 \nu_{\rm cyc}(\mathscr H_{\rm ar})=\frac13,
 \qquad
 \nu_{\rm cyc}(\mathscr H_{\rm vol})=\frac23.
\]
La mesure uniforme \(\tau_{468}\) vit sur le catalogue des émissions. La
mesure de Parry vit sur les histoires admissibles de longueur arbitraire. La
première compte les instants, la deuxième compte les émissions et la troisième pondère
les cylindres dynamiques. Aucune n’est une approximation des autres.
\end{remark}

\section{Caractère multiplicatif et conservation cylindrique}

Pour une route admissible \(\gamma:i\to j\) de longueur \(n\), nous définissons
\begin{equation}
 \chi_P(\gamma)=\varphi^n\frac{r_i}{r_j}.
 \label{gfp:b:eq:caracter-parry-phi}
\end{equation}

\begin{proposition}[Multiplicativité par concaténation]
\label{gfp:b:prop:multiplicatividad-parry-phi}
Si \(\gamma_1:i\to j\) et \(\gamma_2:j\to k\), alors
\[
 \chi_P(\gamma_2\circ\gamma_1)
 =\chi_P(\gamma_2)\chi_P(\gamma_1).
\]
Sur une route fermée de longueur \(n\),
\[
 \chi_P(\gamma)=\varphi^n.
\]
\end{proposition}

\begin{proof}
La longueur s’additionne et la coordonnée intermédiaire se simplifie par télescopage :
\[
 \varphi^{n_1+n_2}\frac{r_i}{r_k}
 =\left(\varphi^{n_1}\frac{r_i}{r_j}\right)
  \left(\varphi^{n_2}\frac{r_j}{r_k}\right).
\]
Si \(i=k\), le quotient terminal est égal à un.
\end{proof}

Nous normalisons le vecteur gauche comme
\[
 \ell=\frac{r}{1+\varphi^2},
 \qquad
 \ell^{\mathsf T}r=1.
\]
La mesure d’un cylindre est
\[
 \mu[x_0\cdots x_n]
 =\ell_{x_0}r_{x_n}\varphi^{-n}.
\]
Par conséquent,
\begin{equation}
 V_n(x)
 :=\frac{\mu[x_0]}{\mu[x_0\cdots x_n]}
 =\varphi^n\frac{r_{x_0}}{r_{x_n}}
 =\chi_P(x_0\cdots x_n).
 \label{gfp:b:eq:volumen-cilindrico-phi}
\end{equation}
Pour une dimension typée \(D>0\), l’échelle
\[
 a_n^{(D)}(x)=V_n(x)^{1/D}
\]
satisfait la loi exacte
\begin{equation}
 \boxed{
 \bigl(a_n^{(D)}(x)\bigr)^D
 \mu[x_0\cdots x_n]=\mu[x_0].}
 \label{gfp:b:eq:conservacion-cilindrica-phi}
\end{equation}

En dimension trois et sur des retours fermés,
\[
 \frac{a_9}{a_0}=\varphi^3,
 \qquad
 \frac{a_{27}}{a_0}=\varphi^9,
 \qquad
 \frac{a_{108}}{a_0}=\varphi^{36}.
\]
Si \(A_k=a_{9k}\), alors
\begin{equation}
 A_{k+1}-2A_k+A_{k-1}
 =A_{k-1}(\varphi^3-1)^2>0.
 \label{gfp:b:eq:convexidad-retornos-phi}
\end{equation}
C’est une convexité en profondeur de retour ; elle n’est pas interprétée comme une
grandeur cinématique sans une carte supplémentaire de durée.

\section{Mémoire quadratique et branche contractée}

Les valeurs propres de \(F_{\rm av}\) sont
\[
 \varphi,
 \qquad
 -\varphi^{-1}.
\]
La branche contractée change d’orientation à chaque itération. Pour conserver
ce signe avant d’élever au carré, nous passons à
\(\operatorname{Sym}^2(\mathbb R^2)\). Dans la base
\((x^2,2xy,y^2)\),
\begin{equation}
 \operatorname{Sym}^2(F_{\rm av})=
 \begin{pmatrix}
 0&0&1\\
 0&1&2\\
 1&1&1
 \end{pmatrix}.
 \label{gfp:b:eq:sym2-Fav-phi}
\end{equation}

\begin{theorem}[Spectre de la mémoire de second ordre]
\label{gfp:b:thm:espectro-sym2-phi}
\[
 \det\bigl(tI-\operatorname{Sym}^2(F_{\rm av})\bigr)
 =(t+1)(t^2-3t+1),
\]
et
\begin{equation}
 \boxed{
 \operatorname{Spec}\bigl(\operatorname{Sym}^2(F_{\rm av})\bigr)
 =\{\varphi^2,-1,\varphi^{-2}\}.}
 \label{gfp:b:eq:espectro-sym2-phi}
\end{equation}
L’unique mode stable positif est \(\varphi^{-2}\).
\end{theorem}

\begin{proof}
Les valeurs propres d’une puissance symétrique de degré deux sont les produits
\(\lambda_1^2\), \(\lambda_1\lambda_2\) et \(\lambda_2^2\). Avec
\(\lambda_1=\varphi\) et \(\lambda_2=-\varphi^{-1}\), on obtient
\(\varphi^2\), \(-1\) et \(\varphi^{-2}\). La première dépasse un, la
deuxième a pour module un et la troisième appartient à \((0,1)\).
\end{proof}

Sur une route fermée avec \(V_n=\varphi^n\), le mode stable satisfait
\begin{equation}
 \boxed{M_n=\varphi^{-2n}=V_n^{-2}.}
 \label{gfp:b:eq:modo-estable-phi}
\end{equation}

\section{Réalisation hyperbolique normalisée de l’auto-échelle : \(\exp(2\log\varphi\,J_h)=F_{\mathrm{av}}^2\)}

Soit
\[
 A=F_{\rm av}^{\,2}=
 \begin{pmatrix}1&1\\1&2\end{pmatrix},
 \qquad
 J_h=\frac{2A-3I}{\sqrt5}.
\]
Alors
\[
 J_h^2=I,
 \qquad
 A=\frac32I+\frac{\sqrt5}{2}J_h.
\]
Comme
\[
 \cosh(2\log\varphi)=\frac32,
 \qquad
 \sinh(2\log\varphi)=\frac{\sqrt5}{2},
\]
on obtient
\begin{equation}
 \boxed{
 \exp(2\log\varphi\,J_h)=F_{\rm av}^{\,2}.}
 \label{gfp:b:eq:realizacion-hiperbolica-phi}
\end{equation}
La constante d’auto-échelle est donc simultanément la valeur propre de
Perron de l’incidence et le paramètre exponentiel de sa réalisation
hyperbolique normalisée.

\section{Caractérisation, encadrements et prolongement}

La construction finie détermine la demi-droite positive de \(F_{\rm av}\). La
caractérisation analytique ultérieure est
\[
 x^2-x-1=0,
 \qquad x>0.
\]
Les encadrements \eqref{gfp:b:eq:horquilla-fibonacci-phi} ont des extrémités rationnelles
et une largeur explicite \eqref{gfp:b:eq:anchura-fibonacci-phi}. Pour chaque niveau
nonadique \(K\), nous choisissons le plus petit \(m=m_{\varphi}(K)\) dont l’encadrement est
contenu dans un unique sous-cylindre parmi les \(729\) éléments de la fibre
ambiante. La minimalité fait de la localisation une fonction, non un
choix tacite.

Les cylindres sélectionnés sont compatibles avec les troncatures et leurs
diamètres tendent vers zéro. Leur intersection contient un unique réel ; le
\cref{gfp:b:thm:rayo-perron-phi} l’identifie à \(\varphi\). Les quotients de
Fibonacci sont donc des certificats rationnels de la localisation, non des données
décimales employées pour décider du germe.

\Needspace{22\baselineskip}
\section{Certificat interne d’unicité, de compatibilité et de convergence de l’auto-échelle de Perron}

\begin{remark}[Certificat et contrôle de validité : réfutation de la biographie d’auto-échelle]
La construction est réfutée si l’un quelconque des faits
suivants se produit :
\begin{enumerate}
 \setlength{\itemsep}{0.15\baselineskip}
 \item le prédicat fini sélectionne plus d’une orbite ou plus d’une émission ;
 \item la récurrence \(L_0,L_0,L_1,L_1\) ne produit pas
       \eqref{gfp:b:eq:w30-phi} ;
 \item le calendrier ne produit pas exactement les trois paires utilisées dans
       \eqref{gfp:b:eq:N3-phi} ;
 \item un multiple de dénombrement \(N_{9},N_{27},N_{108}\) est confondu avec une
       puissance de \(N_3\) ;
 \item \(F_{\rm av}\) possède une autre demi-droite propre positive ;
 \item \eqref{gfp:b:eq:potencias-Fav-fibonacci} ou l’identité de Cassini échoue ;
 \item les encadrements rationnels ne sont pas emboîtés ou ne contractent pas leur largeur ;
 \item la fréquence chronologique, la mesure uniforme du
       catalogue ou la mesure de Parry sont identifiées ;
 \item une table décimale intervient avant la construction de la demi-droite de Perron ;
 \item un cylindre sélectionné ne se tronque pas au cylindre précédent.
\end{enumerate}
\end{remark}

%
\endgroup

% Propietario dodecafásico exacto y reunión transversal exacta. El primero
% deriva la coordenada firmada y el sello; el segundo hace visible, antes del
% capítulo de publicación, la bifurcación numérica/excepcional del mismo
% estado. Ninguno recibe \alpha ni valores convencionales como entrada.
\begin{HMTScientificOwnerSubordinated}
\chapter{La fermeture dodécaphasique de \texorpdfstring{\(\alpha\)}{alpha}}
\label{chap:alpha}
\index{alpha@\(\alpha\)}
\index{cycle dodécaphasique}

Cette section établit la fermeture dodécaphasique au moyen de l'arithmétique entière en
base mille. La construction explicite les douze récurrences, les quotients
euclidiens propagés et les conditions de frontière qui déterminent la sortie.

\section{Le sélecteur terminal des trois canaux}
\label{sec:selector-terminal-alpha-rev6}

La fermeture réunit trois générations déjà construites, mais ne les identifie pas par
leur nom décimal. À profondeur terminale, les catalogues de queues compatibles
possèdent des cardinaux
\[
 |\mathcal T_\pi|=196,\qquad
 |\mathcal T_e|=197,\qquad
 |\mathcal T_\varphi|=196.
\]
L'objet de sélection est donc
\(\mathcal T_\pi\times\mathcal T_e\times\mathcal T_\varphi\), avec
\[
 196\cdot197\cdot196=7\,567\,952
\]
triplets possibles.

La frontière de coordonnées fixée par le lecteur terminal est
\[
 q=(1,2,2,1,2,0)\in\mathbb F_3^6.
\]
La loi de canal \(\pi+e-\varphi\) possède le vecteur de signes
\[
 \sigma=(1,1,-1)=(1,1,2)\in\mathbb F_3^3
\]
et la droite duale non nulle
\[
 \langle2\sigma\rangle
 =\langle(2,2,1)\rangle.
\]
La première donnée agit sur les colonnes de la frontière ; la seconde, sur
ses lignes. Ce ne sont pas deux traits interchangeables.

\begin{theorem}[Sélecteur matriciel terminal de \(\alpha\)]
\label{thm:selector-terminal-alpha-rev6}
Dans le catalogue terminal précédent, les signatures locales de Gram, de norme, de poids et de
distance de Hamming réduisent les \(7\,567\,952\) triplets à \(331\). En exigeant
la marge de colonnes \(q\), il reste deux candidates :
\[
 (120,197,157),\qquad(129,197,148).
\]
La condition que la charge des lignes appartienne à la droite du canal
\(\langle\sigma\rangle\) sélectionne exactement la seconde. Ses blocs sont
\[
 \boxed{
 b_\pi=021020,\qquad
 b_e=001220,\qquad
 b_\varphi=100210.}
\]
\end{theorem}

\begin{proof}
Les deux matrices terminales qui survivent à la marge de colonnes sont
\[
B_0=
\begin{pmatrix}
0&2&0&2&2&0\\
0&0&1&2&2&0\\
1&0&1&0&1&0
\end{pmatrix},
\qquad
B_1=
\begin{pmatrix}
0&2&1&0&2&0\\
0&0&1&2&2&0\\
1&0&0&2&1&0
\end{pmatrix}.
\]
Leurs sommes de lignes dans \(\mathbb F_3\) sont, respectivement,
\[
 (0,2,0),\qquad(2,2,1).
\]
La première n'appartient pas à
\(\langle(1,1,2)\rangle=\{0,(1,1,2),(2,2,1)\}\) ; la seconde, si, car
\((2,2,1)=2\sigma\). Ainsi, seule \(B_1\) satisfait simultanément la
frontière de coordonnées et la charge du canal. Ses trois lignes sont les blocs
présentés et correspondent au triplet \((129,197,148)\).
\end{proof}

Le théorème n'obtient pas \(q\) à partir de la charge des lignes : \(q\) est la donnée
de coordonnées du lecteur terminal et \(\sigma\) est l'orientation du canal.
Ce qui est démontré, c'est qu'une fois déclarées les deux frontières du même objet,
la branche réelle est unique. Cette séparation empêche de transformer une signature trouvée
a posteriori en une loi prétendument dérivée.

\begin{definition}[Deux lectures de l'état terminal enrichi]
Soit
\[
 x_{\rm term}\in\mathcal X_{\rm TPK}^{\rm term,enr}
 \subseteq\mathcal X_{\rm TPK}^{[108]}.
\]
Cet état est la sortie
\[
 x_{\rm term}
 =(\mathcal U_{108}\circ\cdots\circ\mathcal U_1)(x_{\rm seed})
\]
de la chaîne complète de sélection, de transport et de mise à jour du TPK.
Le sélecteur précédent est la lecture
\[
 \pi_{\rm term}(x_{\rm term})
 =B_{\rm term}
 :=\begin{pmatrix}021020\\001220\\100210\end{pmatrix}.
\]
Le registre orienté et le lecteur des canaux de \(90^\circ\) et
\(120^\circ\) produisent, à partir du livre du même état,
\begin{equation}\label{eq:alpha-subespacio-firmado}
 \mathcal E_{108}^{90,120}\!
 \left(\mathcal R_{12}(x_{\rm term})\right)
 =(b^{90},b^{120},Q_{\rm TPK}),
\end{equation}
avec
\[
\begin{aligned}
 b^{90}={}&(-495,-116,-684,108,601,-90,
             -173,-88,313,560,-397,461),\\
 b^{120}={}&(-425,-281,-481,671,-255,30,
              475,-543,680,251,6,-128),\\
 Q_{\rm TPK}={}&6263.
\end{aligned}
\]
Pour \(r=1,2,3\), soient
\[
 B_r=\sum_{j=0}^{3}b^{120}_{r+3j},\qquad
 A_1=\frac{Q_{\rm TPK}+2B_1+B_2}{3},\qquad
 A_2=A_1-B_1,\qquad A_3=A_2-B_2,
\]
et \(d_{r,j}=b^{90}_{r+3j}\). Le lecteur harmonique \(\Pi_H\) forme les
blocs
\begin{equation}\label{eq:alpha-modulo-imagen-H12}
 \Pi_H^{(r)}
 =\bigl(A_r,d_{r,1}-d_{r,3},
              d_{r,0}+d_{r,2},d_{r,0}-d_{r,2}\bigr)
\end{equation}
et les entrelace par colonnes. La lecture signée est donc la
composition produite
\begin{equation}\label{eq:alpha-Rsgn-definido}
 \boxed{
 R_{\rm sgn}:=
 \Pi_H\circ\mathcal E_{108}^{90,120}\circ\mathcal R_{12}:
 \mathcal X_{\rm TPK}^{\rm term,enr}
 \longrightarrow
 \mathcal U_{12}^{\rm sgn}:=\operatorname{im}R_{\rm sgn}.}
\end{equation}
Ni \(U\) ni \(K\) ne sont des coordonnées ajoutées au domaine : \(U\) est la sortie
de \eqref{eq:alpha-Rsgn-definido}, et \(K\) sera son unique antécédent entier
par Hadamard.
On n'affirme pas une factorisation
\(B_{\rm term}\to U_{12}^{\rm sgn}\) : les deux sorties conservent des composantes
distinctes du même état terminal.

Pour \(N'\geq N\geq108\), les troncatures agissent uniquement sur les histoires
enrichies :
\begin{equation}\label{eq:alpha-truncamiento-firmado-preserva-KU}
 \tau_{N',N}:
 \mathcal X_{\rm TPK}^{[N'],\rm enr}
 \longrightarrow\mathcal X_{\rm TPK}^{[N],\rm enr}.
\end{equation}
Si \(R_{\rm sgn,N}\) désigne la même composition à la profondeur \(N\), la
naturalité est le carré concret
\begin{equation}\label{eq:alpha-cuadrado-naturalidad-Rsgn}
 \boxed{
 R_{\rm sgn,N}\circ\tau_{N',N}
 =R_{\rm sgn,N'}.}
\end{equation}
\end{definition}

\begin{theorem}[Construction réversible de la coordonnée signée]
\label{thm:alpha-Rsgn-construido}
La composée \eqref{eq:alpha-Rsgn-definido}, appliquée à l'état terminal
produit, donne exactement
\begin{equation}\label{eq:alpha-Rsgn-evaluacion}
 U=R_{\rm sgn}(x_{\rm term})
 =(2378,1406,2479,-452,998,-551,-668,-204,-371,-322,-28,-997).
\end{equation}
Son inversion par trois blocs de Hadamard d'ordre quatre produit ensuite
\[
 \boxed{
 K=(234,543,140,729,659,824,621,058,914,794,146,601).}
\]
La transformation est inversible ; elle satisfait le carré
\eqref{eq:alpha-cuadrado-naturalidad-Rsgn} sous les troncatures
\eqref{eq:alpha-truncamiento-firmado-preserva-KU} ; et elle est indépendante de
\(\alpha\), de CODATA et de toute valeur métrologique.
\end{theorem}

\begin{proof}
Les formules \eqref{eq:alpha-subespacio-firmado} et
\eqref{eq:alpha-modulo-imagen-H12} donnent
\[
 (B_1,B_2,B_3)=(972,-1073,101),\qquad
 (A_1,A_2,A_3)=(2378,1406,2479),
\]
et produisent les trois blocs signés
\[
\begin{aligned}
 U^{(1)}&=(2378,-452,-668,-322),\\
 U^{(2)}&=(1406,998,-204,-28),\\
 U^{(3)}&=(2479,-551,-371,-997).
\end{aligned}
\]
Leur entrelacement est \eqref{eq:alpha-Rsgn-evaluacion}. Comme
\(H_4^2=4I_4\) et
\(\det(H_4\oplus H_4\oplus H_4)=(-16)^3=-4096\), on a
\begin{equation}\label{eq:alpha-Rsgn-inversa}
 K=\mathcal H_{12}^{-1}U
 =\frac14\mathcal H_{12}U
\end{equation}
après avoir défait la permutation des orbites. En effet,
\[
 H_4U^{(1)}=(936,2916,2484,3176),\quad
 H_4U^{(2)}=(2172,2636,232,584),
\]
et
\[
 H_4U^{(3)}=(560,3296,3656,2404).
\]
La division forcée par quatre donne les orbites
\((234,729,621,794)\), \((543,659,058,146)\) et
\((140,824,914,601)\), dont l'entrelacement est le sceau déclaré.

À titre de contrôle indépendant, les opérateurs cycliques
\((D_3K)_m=K_m-K_{m+3}\) et
\((D_4K)_m=K_m-K_{m+4}\), avec \(Q(K)=\sum_mK_m\), satisfont
\begin{equation}\label{eq:alpha-extractor-90-120}
 \operatorname{rank}(D_3,D_4)=11,
 \qquad
 \operatorname{rank}(D_3,D_4,Q)=12.
\end{equation}
Ainsi \((D_3K,D_4K,Q(K))\) reconstruit un sceau unique ; de plus, le
calcul retrouve exactement
\[
 D_3K=b^{90},\qquad D_4K=b^{120},\qquad Q(K)=Q_{\rm TPK},
\]
ce qui certifie le lecteur préalable du registre, sans transformer ses sorties
en entrées. La troncature
\eqref{eq:alpha-truncamiento-firmado-preserva-KU} ne réécrit pas les cent huit premières
fenêtres du livre ; les deux membres de
\eqref{eq:alpha-cuadrado-naturalidad-Rsgn} lisent donc le même registre et
produisent \(U\). Aucune des matrices utilisées ne contient \(\alpha\) ni
une donnée expérimentale.
\end{proof}

Le calendrier observable de 108 pas est une projection strictement plus
pauvre : réinitialiser ses curseurs tous les 9 pas efface précisément les
coordonnées bidirectionnelles que le registre enrichi conserve. C'est
pourquoi le registre brut ne remplace pas l'état enrichi et n'est pas utilisé pour
reconstruire \(R_{\rm sgn}\). Le théorème de sélection précédent construit
\(B_{\rm term}\) ; le théorème \ref{thm:alpha-Rsgn-construido} construit, de
manière séparée, la seconde publication du même état.

\section{État dodécaphasique et applications réversibles}

Le théorème de cette section reçoit du
théorème~\ref{thm:alpha-Rsgn-construido} le vecteur signé
\[
\Usgn=(2378,1406,2479,-452,998,-551,-668,-204,-371,-322,-28,-997).
\]
La voie typée démontrée à partir de cette entrée est
\[
\Usgn\longleftrightarrow K,
\]
\[
K\longleftrightarrow(D_3K,D_4K,Q),
\]
\[
(P,E,\Phi,K,c_{13}=0)
\longleftrightarrow A=\alpha_{12}.
\]
L'évaluation directe ne reçoit ni \(\alpha\) ni une valeur métrologique en
entrée. Elle reçoit \(\Usgn\), les trois représentations archimédiennes et la
condition aux limites. Le
\cref{thm:selector-terminal-alpha-rev6} construit auparavant la sélection des
trois blocs terminaux ; la présente couche commence par le vecteur signé qui
les transporte à la fermeture dodécaphasique. La trace réduite d'une seule
coordonnée positionnelle n'intervient pas dans cette détermination.

\begin{remark}[Deux représentations dodécaphasiques de types distincts]
\[
U_{12}^{\rm cat}=U_6\Vert U_6
\]
est un mot périodique du catalogue, tandis que
\[
\Usgn=(2378,1406,2479,-452,998,-551,-668,-204,-371,-322,-28,-997)
\]
est un vecteur entier signé de l'espace complet des états de TPK\@. Le premier appartient à une
représentation par blocs et a une période de six. Le second conserve l'orientation, n'est pas
un mot en base mille et n'a pas cette période. Ils ne sont pas interchangeables.
\end{remark}

\section{Le vecteur canonique et les trois orbites}

Le vecteur dodécaphasique est
\[
\boxed{
K=(234,543,140,729,659,824,621,058,914,794,146,601).}
\]
Sa somme est
\[
Q(K)=6263.
\]
Le pas trois divise \(\Cdoce\) en trois orbites de quatre secteurs :
\[
\begin{aligned}
K^{(1)}&=(K_1,K_4,K_7,K_{10})=(234,729,621,794),\\
K^{(2)}&=(K_2,K_5,K_8,K_{11})=(543,659,058,146),\\
K^{(3)}&=(K_3,K_6,K_9,K_{12})=(140,824,914,601).
\end{aligned}
\]
Chaque orbite parcourt quatre secteurs séparés par \(90^\circ\). Les étiquettes
cardinales nord--est--sud--ouest apparaissent après fixation de l'origine et de l'orientation ; l'ordre
primaire est \((m,m+3,m+6,m+9)\).

\section{Première représentation : transformée de Hadamard}

Soit
\[
H_4=
\begin{pmatrix}
1&1&1&1\\
1&1&-1&-1\\
1&-1&1&-1\\
1&-1&-1&1
\end{pmatrix}.
\]
L'action sur les trois orbites donne
\[
\begin{aligned}
H_4K^{(1)}&=(2378,-452,-668,-322),\\
H_4K^{(2)}&=(1406,998,-204,-28),\\
H_4K^{(3)}&=(2479,-551,-371,-997).
\end{aligned}
\]
En entrelaçant par colonnes, nous retrouvons
\[
\mathcal H_{12}(K)=\Usgn.
\]
Comme
\[
H_4^2=4I_4,
\qquad
\det H_4=-16,
\]
la transformation globale est inversible sur \(\mathbb Q^{12}\) :
\[
K^{(r)}=\frac14H_4U^{(r)},
\qquad
\det\mathcal H_{12}=(-16)^3=-4096\ne0.
\]
Dans la représentation canonique, l'inverse produit un unique vecteur entier de
\([0,999]^{12}\).

\begin{theorem}[Structure entière de la transformation de Hadamard]
\label{thm:smith-hadamard}
La forme normale de Smith du bloc \(H_4\), notée
\(\operatorname{SNF}(H_4)\), est
\[
 \operatorname{SNF}(H_4)=\operatorname{diag}(1,2,2,4).
\]
Par conséquent, pour la somme directe des trois blocs,
\[
 \operatorname{coker}\mathcal H_{12}
 \cong(\mathbb Z/2\mathbb Z)^6\oplus(\mathbb Z/4\mathbb Z)^3,
\]
et l'image de \(\mathbb Z^{12}\) a pour indice \(4096\). Un vecteur
\(u\in\mathbb Z^4\) appartient à \(H_4\mathbb Z^4\) exactement lorsque
\[
 H_4u\equiv0\pmod 4.
\]
\end{theorem}

\begin{proof}
Le plus grand commun diviseur des entrées de \(H_4\) est un. Les mineurs
d'ordre deux sont pairs et l'un a pour valeur absolue deux ; les mineurs d'ordre
trois sont des multiples de quatre et l'un a pour valeur absolue quatre ;
enfin, \(\lvert\det H_4\rvert=16\). Les diviseurs déterminantaux sont
 donc \(1,2,4,16\), d'où l'on obtient les facteurs invariants
\(1,2,2,4\). La somme directe de trois blocs répète ces facteurs et a pour
indice \(16^3=4096\). Puisque \(H_4^{-1}=H_4/4\), la dernière
congruence équivaut à l'intégralité de l'antécédent.
\end{proof}

\section{Transporteur cyclotomique relatif entre les secteurs trois et quatre}
\label{sec:transportador-ciclotomico-a34-rev8}

Soit \(\mathsf R\) la rotation régulière du cycle \(C_{12}\) sur
\(\mathbb Q^{12}\), et soient
\[
 W_d^{\rm cyc}:=\ker\Phi_d(\mathsf R)
\]
ses secteurs cyclotomiques. Les projecteurs rationnels sur les secteurs
d'ordres trois et quatre peuvent s'écrire
\[
 \Pi_3=
 \frac14(I+\mathsf R^3+\mathsf R^6+\mathsf R^9)
 -\frac1{12}\sum_{k=0}^{11}\mathsf R^k,
\]
\[
 \Pi_4=
 \frac16(I-\mathsf R^2+\mathsf R^4-\mathsf R^6
              +\mathsf R^8-\mathsf R^{10}).
\]
Pour fixer l'ordre des coordonnées, soit
\[
 \mathsf P_{\rm orb}(z_1,\ldots,z_{12})
 =\bigl((z_1,z_4,z_7,z_{10}),
        (z_2,z_5,z_8,z_{11}),
        (z_3,z_6,z_9,z_{12})\bigr).
\]
La transformée globale utilisée ci-dessus est, dans l'ordre naturel de \(C_{12}\),
\[
 \mathcal H_{12}
 :=\mathsf P_{\rm orb}^{-1}
   (H_4\oplus H_4\oplus H_4)\mathsf P_{\rm orb}.
\]
Autrement dit, \(H_4\oplus H_4\oplus H_4\) agit dans l'ordre des trois
orbites de pas trois, et non sur trois quadruplets contigus de l'ordre naturel.
Nous définissons alors le transporteur relatif
\[
 A_{34}:=\left.\Pi_4\mathcal H_{12}\Pi_3
          \right|_{W_3^{\rm cyc}}:
 W_3^{\rm cyc}\longrightarrow W_4^{\rm cyc}.
\]

\begin{theorem}[Transporteur cyclotomique dodécaphasique]
\label{thm:transportador-ciclotomico-a34-rev8}
Dans les bases
\[
 u_1=(1,-1,0)^{\times4},\qquad
 u_2=(0,1,-1)^{\times4}
\]
de \(W_3^{\rm cyc}\), et
\[
 v_1=(1,0,-1,0)^{\times3},\qquad
 v_2=(0,1,0,-1)^{\times3}
\]
de \(W_4^{\rm cyc}\), on a
\[
 A_{34}u_1=\frac23(v_1-v_2),\qquad
 A_{34}u_2=\frac23(v_1+v_2),
\]
et, par conséquent,
\[
 \boxed{\det A_{34}=\frac89}.
\]
En particulier, \(A_{34}\) est un isomorphisme rationnel entre les deux secteurs.
\end{theorem}

\begin{proof}
On substitue les quatre suites périodiques précédentes dans les expressions
de \(\Pi_3,\Pi_4\) et dans les trois blocs de Hadamard. La matrice de \(A_{34}\)
dans les bases déclarées est
\[
 \begin{pmatrix}
  2/3&2/3\\
 -2/3&2/3
 \end{pmatrix},
\]
dont le déterminant est \(8/9\).
\end{proof}

L'involution rationnelle qui échange \(W_3^{\rm cyc}\) et
\(W_4^{\rm cyc}\) au moyen de \(A_{34}\) et de \(A_{34}^{-1}\), et fixe
\(W_{12}^{\rm cyc}\), transporte la somme
\(W_4^{\rm cyc}\oplus W_{12}^{\rm cyc}\) vers
\(W_3^{\rm cyc}\oplus W_{12}^{\rm cyc}\). Dans cette dernière somme, le polynôme
cyclotomique de la rotation est
\[
 \Phi_3(x)\Phi_{12}(x),
\]
qui coïncide avec le polynôme du Coxeter de type \(E_6\). C'est une signature
spectrale exacte qui prépare le corridor géométrique ultérieur.

\begin{statusbox}[title={Transport rationnel et continuation entière}]
L'opérateur \(A_{34}\) réalise le transport cyclotomique rationnel entre les
secteurs d'ordres trois et quatre. En fixant \(W_{12}^{\rm cyc}\), il situe la
dodécaphase dans le secteur spectral exact
\(\Phi_3\Phi_{12}\) de type \(E_6\). La continuation entière est construite
au moyen de deux sources ultérieures et compatibles : le
\cref{thm:entrelazador-dodecafase-witt} fournit l'isométrie
\(M_{12}\)-équivariante vers le module de Witt, et le
\cref{thm:identificacion-isometrica-tres-proyectores} enchaîne les trois
projecteurs d'incidence. Sur cette réalisation entière, le
\cref{p12-83-c20-s6:thm:grafo-27-rectas-nuevo} construit le graphe des
vingt-sept droites et le
\cref{p12-83-c20-s6:chap:holonomia-schlafli-e6-nuevo} identifie son action de
Weyl. Ainsi, \(A_{34}\) fixe l'étape rationnelle de la généalogie et les
opérateurs d'entrelacement de Witt--Schläfli réalisent son prolongement réticulaire et
incidenciel.
\end{statusbox}

La séparation des signes est
\[
U_{12}^{+}
=(2378,1406,2479,0,998,0,0,0,0,0,0,0),
\]
\[
U_{12}^{-}
=(0,0,0,452,0,551,668,204,371,322,28,997),
\]
et
\[
\Usgn=U_{12}^{+}-U_{12}^{-}.
\]

\section{Seconde représentation : différences transversales}

Avec des indices dans \(\{1,\ldots,12\}\) et la convention cyclique
\[
 K_{m+r}:=K_{1+((m+r-1)\bmod12)},
\]
nous définissons
\[
(D_3K)_m=K_m-K_{m+3},
\qquad
(D_4K)_m=K_m-K_{m+4}.
\]
L'extracteur transversal dodécaphasique est l'opérateur
\[
 \mathcal E_{\mathrm{dod}}
 :=(D_3,D_4,Q).
\]
L'indice ``dod'' déclare son domaine. Les décalages trois et quatre
correspondent à des séparations angulaires de \(90^\circ\) et \(120^\circ\) dans le
cycle orienté \(C_{12}\), mais \(\mathcal E_{\mathrm{dod}}\) n'est ni l'opérateur
d'incidence hexade--octade ni le noyau angulaire utilisé ultérieurement dans
les tours.
Les données sont
\[
\begin{aligned}
D_3K={}&(-495,-116,-684,108,601,-90,\\
       &\qquad-173,-88,313,560,-397,461),\\[1mm]
D_4K={}&(-425,-281,-481,671,-255,30,\\
       &\qquad475,-543,680,251,6,-128),\\[1mm]
Q={}&6263.
\end{aligned}
\]
Si \(D_3v=D_4v=0\), le vecteur est constant aux pas trois et quatre. Comme
\(\gcd(3,4)=1\), ces pas relient tout le cycle et
\[
\ker D_3\cap\ker D_4=\langle\mathbf1\rangle.
\]
Par conséquent
\[
\operatorname{rank}(D_3,D_4)=11.
\]
La charge totale élimine la translation constante, car
\(Q(\mathbf1)=12\ne0\). Donc
\[
\boxed{\operatorname{rank}(D_3,D_4,Q)=12,}
\]
et
\[
K\longleftrightarrow(D_3K,D_4K,Q)
\]
est un système de coordonnées exact.

\begin{theorem}[Forme normale entière du système transversal]
\label{thm:smith-dodecafase}
Soit
\[
 \mathcal D=(D_3,D_4,Q):
 \mathbb Z^{12}\longrightarrow\mathbb Z^{25}.
\]
Ses facteurs invariants non nuls sont
\[
 \operatorname{SNF}(\mathcal D)
 =\operatorname{diag}(\underbrace{1,\ldots,1}_{11},12),
\]
et, par conséquent,
\[
 \operatorname{coker}\mathcal D
 \cong\mathbb Z^{13}\oplus\mathbb Z/12\mathbb Z.
\]
\end{theorem}

\begin{proof}
Les lignes de \((D_3,D_4)\) sont des incidences orientées du graphe
\[
 \operatorname{Cay}
 \bigl(\mathbb Z/12\mathbb Z,\{\pm3,\pm4\}\bigr).
\]
Comme les pas trois et quatre engendrent \(\mathbb Z/12\mathbb Z\), le graphe
est connexe.
Les incidences d'un arbre couvrant fournissent onze mineurs réduits
unimodulaires ; les onze premiers facteurs invariants sont donc égaux à
un.

Tout mineur d'ordre douze contient la ligne \(Q\), puisque le bloc
d'incidences est de rang onze. En ajoutant les onze premières colonnes à la
dernière, opération unimodulaire, la dernière colonne s'annule dans les lignes
d'incidence et vaut douze dans la ligne \(Q\). Tous les mineurs maximaux sont
divisibles par douze. Le choix des onze incidences d'un arbre donne un
mineur restant de valeur absolue un et, par conséquent, un déterminant maximal
de valeur absolue douze. Le dernier facteur invariant est exactement douze.
\end{proof}

La composante \(\mathbb Z/12\mathbb Z\) enregistre la condition de
compatibilité entière de la charge totale avec le potentiel reconstruit.
En particulier, l'unicité rationnelle prouvée par le rang est affinée ici en
un énoncé sur le réseau entier et son unique obstruction torsionnelle.

L'observable signé peut aussi être reconstruit à partir de ces coordonnées.
Pour \(r=1,2,3\), soit
\[
B_r=\sum_{j=0}^{3}(D_4K)_{r+3j}.
\]
On obtient
\[
(B_1,B_2,B_3)=(972,-1073,101).
\]
Les sommes orbitales sont
\[
A_1=\frac{Q+2B_1+B_2}{3}=2378,
\quad
A_2=A_1-B_1=1406,
\quad
A_3=A_2-B_2=2479.
\]
Si
\(d_{r,j}=(D_3K)_{r+3j}\) pour \(j=0,1,2,3\), avec la même convention
cyclique, le bloc signé de l'orbite \(r\) est
\[
(A_r,d_{r,1}-d_{r,3},d_{r,0}+d_{r,2},d_{r,0}-d_{r,2}),
\]
exactement le bloc de Hadamard correspondant.

\begin{remark}[Codomaines de \(\mathcal H_{12}\) et
\(\mathcal E_{\mathrm{dod}}\)]
L'écriture correcte est
\[
\mathcal H_{12}(K)=\Usgn,
\qquad
\mathcal E_{\mathrm{dod}}(K)=(D_3K,D_4K,Q).
\]
Les deux représentations sont compatibles car elles reconstruisent le même vecteur,
mais aucun diagramme entre leurs codomaines n'a été défini ici. Leurs formules et
codomaines sont distincts. Partager les pas \(3/4\) ou les indices \(90/120\) n'autorise pas
à les identifier.
\end{remark}

\section{Les trois coordonnées archimédiennes}

Les vecteurs de douze triades sont
\[
\begin{aligned}
P={}&(141,592,653,589,793,238,462,643,383,279,502,884),\\
E={}&(718,281,828,459,045,235,360,287,471,352,662,497),\\
\Phi={}&(618,033,988,749,894,848,204,586,834,365,638,117).
\end{aligned}
\]
Sous forme concaténée :
\[
\begin{aligned}
\iota(P)&=141592653589793238462643383279502884,\\
\iota(E)&=718281828459045235360287471352662497,\\
\iota(\Phi)&=618033988749894848204586834365638117.
\end{aligned}
\]
\(P\) est la coordonnée décimale partagée par les cinq régions de
\(\pi\) ; le calcul n'affirme pas que ces régions sont identiques.

\section{Récurrence de retenue en base mille}

Nous fixons
\[
c_{13}=0
\]
et calculons de droite à gauche, \(m=12,11,\ldots,1\) :
\[
s_m=P_m+E_m-\Phi_m-K_m+c_{m+1},
\]
\[
A_m=s_m\bmod1000,
\qquad
c_m=\frac{s_m-A_m}{1000},
\qquad0\le A_m<1000.
\]
Le tableau est imprimé dans l'ordre croissant des indices, mais chaque retenue provient de la ligne
inférieure.

\begingroup
\scriptsize
\setlength{\tabcolsep}{3.3pt}
\begin{longtable}{@{}r|rrrr|rr|rr@{}}
\caption{Fermeture dodécaphasique complète en base mille.}\label{tab:acarreo-alpha}\\
\toprule
\(m\)&\(P_m\)&\(E_m\)&\(\Phi_m\)&\(K_m\)&\(c_{m+1}\)&\(s_m\)&\(A_m\)&\(c_m\)\\
\midrule
\endfirsthead
\toprule
\(m\)&\(P_m\)&\(E_m\)&\(\Phi_m\)&\(K_m\)&\(c_{m+1}\)&\(s_m\)&\(A_m\)&\(c_m\)\\
\midrule
\endhead
1&141&718&618&234&0&7&007&0\\
2&592&281&033&543&0&297&297&0\\
3&653&828&988&140&\(-1\)&352&352&0\\
4&589&459&749&729&\(-1\)&\(-431\)&569&\(-1\)\\
5&793&045&894&659&\(-2\)&\(-717\)&283&\(-1\)\\
6&238&235&848&824&\(-1\)&\(-1200\)&800&\(-2\)\\
7&462&360&204&621&0&\(-3\)&997&\(-1\)\\
8&643&287&586&058&\(-1\)&285&285&0\\
9&383&471&834&914&\(-1\)&\(-895\)&105&\(-1\)\\
10&279&352&365&794&0&\(-528\)&472&\(-1\)\\
11&502&662&638&146&0&380&380&0\\
12&884&497&117&601&0&663&663&0\\
\bottomrule
\end{longtable}
\endgroup

Les retenues sont
\[
(c_1,\ldots,c_{13})
=(0,0,0,-1,-1,-2,-1,0,-1,-1,0,0,0).
\]
Les valeurs négatives de \(s_m\) sont des emprunts ordinaires en base mille. La
division euclidienne fixe de manière unique le reste et le quotient.

La sortie est
\[
\boxed{
A=(007,297,352,569,283,800,997,285,105,472,380,663),}
\]
et donc
\[
\boxed{
\begin{aligned}
\alpha_{12}
&:=\frac{7297352569283800997285105472380663}{10^{36}},\\
&=0.007297352569283800997285105472380663.
\end{aligned}}
\]
Ce rationnel est la fenêtre dodécaphasique initiale de la section compatible
engendrée par le transport TPK. Si
\(\rho_{N',N}:C_{N'}^\alpha\to C_N^\alpha\) sont ses troncatures, la
sortie complète est
\[
 \boxed{
 \alpha_{\mathrm{HMT}}
 :=\varprojlim_N\alpha_N,\qquad
 \rho_{N',N}(\alpha_{N'})=\alpha_N,\qquad
 \alpha_{36}=\alpha_{12}.}
\]
Le retour de phase prolonge la mémoire et resserre les cylindres
survivants à une profondeur arbitraire ; il ne réinitialise pas l'état et n'achève pas la
construction à trente-six chiffres.

\section{L'identité affine complète}

Coordonnée par coordonnée,
\[
\boxed{
P_m+E_m-\Phi_m-K_m+c_{m+1}
=A_m+1000c_m.}
\]
Si
\[
C=(c_1,\ldots,c_{12}),
\qquad
C^+=(c_2,\ldots,c_{13}),
\]
la forme vectorielle correcte est
\[
\boxed{P+E-\Phi-K+C^+-1000C=A.}
\]
Pour éviter d'identifier les queues finies aux constantes complètes,
nous définissons
\[
 \tau_{36}(x)=10^{-36}\left\lfloor10^{36}
 \bigl(x-\lfloor x\rfloor\bigr)\right\rfloor,
 \qquad
 \kappa(K)=10^{-36}\iota(K).
\]
La notation décimale exacte de l'identité concaténée est alors
\[
 \tau_{36}(\pi)+\tau_{36}(e)-\tau_{36}(\varphi)-\kappa(K)
 =\alpha_{12}.
\]
Sans la troncature explicite, le plongement \(\kappa\) et le cocycle
\(C^+-1000C\), l'égalité en coordonnées serait fausse.

On définit
\[
\iota(v_1,\ldots,v_{12})
=\sum_{m=1}^{12}v_m1000^{12-m}.
\]
En multipliant chaque récurrence par \(1000^{12-m}\) et en sommant, les retenues
se télescopent :
\[
\iota(P)+\iota(E)-\iota(\Phi)-\iota(K)-\iota(A)
=1000^{12}c_1-c_{13}.
\]
Comme \(c_1=c_{13}=0\) :
\[
\boxed{
\iota(P)+\iota(E)-\iota(\Phi)-\iota(K)=\iota(A).}
\]
Les deux entiers restants sont
\[
\iota(K)=234543140729659824621058914794146601,
\]
\[
\iota(A)=007297352569283800997285105472380663.
\]

\section{Inversion et équivalence des trois représentations}

La récurrence s'inverse par
\[
K_m=P_m+E_m-\Phi_m+c_{m+1}-A_m-1000c_m,
\]
et la représentation entière concaténée par
\[
\iota(K)=\iota(P)+\iota(E)-\iota(\Phi)-\iota(A).
\]
Le diagramme complet est
\[
\boxed{
K\longleftrightarrow\Usgn,
\qquad
K\longleftrightarrow(D_3K,D_4K,Q),
\qquad
K\longleftrightarrow A\mid(P,E,\Phi,c_{13}=0).}
\]
Les deux premières sont des représentations linéaires ; la troisième, une application
affine réversible dans la fibre
conditionnée par \(P,E,\Phi\) et la frontière.

\begin{remark}[Commutativité de la fermeture]
L'application directe obtient \(K\) à partir de la représentation signée et produit
\(A\). La formule inverse démontre que la sortie et la frontière retrouvent le
même vecteur sans liberté résiduelle. L'inversibilité ne constitue pas une
preuve statistique indépendante ; la commutativité des trois
représentations est une propriété algébrique de la fermeture.
\end{remark}

L'unicité s'exprime à trois niveaux de la même construction :
\begin{enumerate}
  \item une fois fixés le vecteur signé, \(P,E,\Phi\) et la frontière, le vecteur et la
  sortie sont uniques ;
  \item les représentations sont réversibles et commutatives ;
  \item la généalogie complète
  \(\mathrm{APP}\to\mathrm{TRIT}\to\mathrm{TPK}\) produit ces primitives
  avant la fermeture locale, et ses troncatures compatibles prolongent
  l'identité à la section coinductive \(\alpha_{\mathrm{HMT}}\).
\end{enumerate}

\section{Condition aux limites dodécaphasique}

La première fenêtre dodécaphasique utilise \(c_{13}=0\) et se termine en
\[
\ldots|380|663.
\]
Cette frontière fait partie du domaine de l'application affine réversible. Le
prolongement n'introduit pas d'autre système : il itère le même transport gradué
du TPK sur
\[
 w_6\to w_{12}\to w_{18}\to w_{24}\to w_{30}\to R_{36}\to G_9
\]
et construit les fenêtres successives avec retour de phase, avancée de la mémoire et
retenue terminale compatible. Le
\cref{sec:l9s102:alpha-dos-vias} démontre que les carrés de troncature
de cette voie et du lecteur analytique commutent et que leurs cylindres
survivants possèdent une intersection réduite à un singleton.

\subsection{Classe entière de la retenue}

La condition terminale peut être comparée au problème périodique. Soit
\(S\) le décalage cyclique sur \(\mathbb Z^{12}\) et, pour une base
entière \(b\geq2\), définissons
\[
 \partial_b=S-bI.
\]
L'identité affine périodique prend la forme
\[
 A=P+E-\Phi-K+\partial_bC.
\]
Pour tout \(h\in\mathbb Z^{12}\), la transformation
\[
 K\longmapsto K+\partial_bh,
 \qquad
 C\longmapsto C+h
\]
conserve \(A\). Ainsi, le choix des retenues est un choix de
représentant au sein d'une classe entière.

\begin{proposition}[Quotient périodique de la retenue]
\label{prop:cociente-periodico-acarreo}
Sur un cycle de longueur douze,
\[
 \operatorname{coker}(S-bI)\cong
 \mathbb Z/(b^{12}-1)\mathbb Z.
\]
La concaténation en base \(b\), prise modulo \(b^{12}-1\), détermine la
classe périodique.
\end{proposition}

\begin{proof}
Le module cyclique s'identifie à
\(\mathbb Z[t]/(t^{12}-1)\). Imposer la relation \(t=b\) produit
\(\mathbb Z/(b^{12}-1)\mathbb Z\). L'évaluation du polynôme des
coefficients est la concaténation en base \(b\), avec l'orientation fixée
par l'ordre des douze secteurs.
\end{proof}

La fermeture de \(\alpha\) n'utilise pas le quotient périodique : elle utilise une chaîne
orientée avec une extrémité droite et exige \(c_{13}=0\). La division euclidienne
sélectionne alors un représentant unique. Cette différence explique
pourquoi la frontière ouverte fait partie de la définition de l'application et ne
peut pas être remplacée par une identification cyclique des retenues.

\section{Comparaison avec la référence métrologique}

Après la génération, l'arithmétique entière de la première fenêtre
vérifie la reconnaissance
\[
\boxed{
\alpha_{12}
=10^{-36}
\left\lfloor\frac{10^{36}}{137.035999084}\right\rfloor.}
\]
Son inverse est
\[
\alpha_{12}^{-1}
=137.03599908400000000000000000000001066588\ldots,
\]
et reproduit exactement la chaîne centrale publiée dans l'ajustement de 2018 du
Committee on Data for Science and Technology (CODATA) \cite{Tiesinga2021CODATA} :
\[
\boxed{137.035999084.}
\]
Il ne s'agit pas d'une coïncidence en virgule flottante. Les trente-six chiffres de
\(\alpha_{12}\) constituent la troncature exacte de l'inverse de la
chaîne centrale imprimée. Cette vérification finie ne définit pas
\(\alpha_{\mathrm{HMT}}\) et ne limite pas sa profondeur : la section coinductive
déjà produite détermine toutes ses fenêtres compatibles, tandis que CODATA ne
reconnaît qu'ensuite le préfixe que sa résolution expérimentale permet
de confronter.

La reconnaissance métrologique ne fait pas partie de la récurrence. Le résultat
mathématique est l'identité exacte des deux voies
\[
 \alpha_A=\alpha_B=\alpha_{\mathrm{HMT}},
\]
dont les fenêtres dodécaphasiques sont déterminées par le vecteur signé, les
trois mots engendrés et les frontières compatibles. Le tableau de CODATA
ne compare qu'ensuite une sortie déjà fixée.

La comparaison avec des ajustements métrologiques ultérieurs est un contrôle externe
séparé. La sortie HMT reste figée et peut donc être confrontée sans
modifier le générateur.

\section{Indépendance des coordonnées et rigidité}

Les tests de la fermeture distinguent les modifications profondes des perturbations
locales :
\begin{center}
\small
\begin{tabularx}{\textwidth}{@{}lYY@{}}
\toprule
Variante & Résultat & Ce qu'il contrôle \\
\midrule
protocole canonique droite--gauche & igualdad exacta & référence \\
utiliser \(U_6\Vert U_6\) comme \(K\) & échoue macroscopiquement & distinction des types \\
rotar \(K\) & rompt la fermeture & origine dodécaphasique \\
reflejar \(K\) & inverse le bord & orientation \\
retenue gauche--droite & change le mot & sens du transport \\
perturber \(K_i\) d'une unité & le mot exact échoue & rigidez profunda \\
\bottomrule
\end{tabularx}
\end{center}
Une perturbation locale peut ne modifier qu'un chiffre profond ; elle reste une
réfutation de l'égalité exacte. Les variantes structurelles produisent des échecs
plus importants et montrent que la fermeture n'est pas invariante sous des ordres arbitraires.

\begin{remark}[Structure de la fermeture dodécaphasique]
La fermeture de \(\alpha\) comprend un cycle orienté, deux représentations
exactes du vecteur \(K\), douze entrées typées, une récurrence de droite à
gauche, treize retenues, une condition aux limites et une identité entière
télescopique. Ces données déterminent la sortie de manière reproductible.
\end{remark}

\end{HMTScientificOwnerSubordinated}
\begingroup
\let\HMTRecoveredOriginalPart\part
\let\part\section
\begin{HMTScientificOwnerSubordinated}
% Adaptador local de aristas para el propietario exacto
% propietarios_exactos/alpha/03a_alpha_y_dependencias_transversales.tex.
% El propietario conserva sus once llamadas \HMTInputLive. En este ensamblado
% sus cuerpos tienen residencias exactas o finales propias; la llamada verifica
% que la arista esté registrada y la deja en el log, pero no vuelve a incluir el
% cuerpo. La existencia y el orden de ambas puntas se auditan contra el grafo de
% inputs. Así no se anteponen lectores ni se publica la copia histórica débil de
% 08a.

\makeatletter
\newcommand{\HMTRegisterLiveInputEdge}[3]{%
  \expandafter\def\csname hmt@liveinput@residence@\detokenize{#1}\endcsname{#2}%
  \expandafter\def\csname hmt@liveinput@order@\detokenize{#1}\endcsname{#3}%
}
\newcommand{\HMTInputLive}[1]{%
  \ifcsname hmt@liveinput@residence@\detokenize{#1}\endcsname
    \edef\HMTLiveInputResidence{%
      \csname hmt@liveinput@residence@\detokenize{#1}\endcsname}%
    \edef\HMTLiveInputOrder{%
      \csname hmt@liveinput@order@\detokenize{#1}\endcsname}%
    \typeout{HMT-INPUT-LIVE-EDGE: #1.tex =>
      \HMTLiveInputResidence\space [\HMTLiveInputOrder]; corps non rechargé}%
  \else
    \PackageError{hmt-alpha-live-adapter}
      {Arête sans résidence effective : #1.tex}
      {Enregistrez un propriétaire exact ou final ; une copie réductrice n'est pas admise.}%
  \fi
}

% Todas las aristas del propietario 03a, en orden de aparición.
\HMTRegisterLiveInputEdge
  {sections/hmt/08_dodecafase_alpha}
  {sections/hmt/08_dodecafase_alpha.tex}
  {UPSTREAM-CH31-PREOWNER}
\HMTRegisterLiveInputEdge
  {sections/hmt/08a_alpha_dos_vias_t1_rev8}
  {../colaboracion/parte_iii/source/public_final/alpha/08a_alpha_dos_vias_t1_rev8_body.tex}
  {DOWNSTREAM-CH31-FINAL}
\HMTRegisterLiveInputEdge
  {sections/hmt/08b_alpha_t1_10000_rev9}
  {../colaboracion/parte_iii/source/public_final/alpha/08b_alpha_t1_10000_rev9.tex}
  {DOWNSTREAM-CH31-FINAL}
\HMTRegisterLiveInputEdge
  {sections/ampliacion_20260818/03_extension_holografica_euler}
  {sucesor_102/deltas_ley9/c57_sistema_operatorio_euler_hmt.tex}
  {DOWNSTREAM-CH57-SUCCESSOR}
\HMTRegisterLiveInputEdge
  {sections/hmt/07b_realizacion_analitica_contraangulo}
  {sucesor_102/wrappers/capitulo_38_jerarquia_unica.tex}
  {DOWNSTREAM-CH38-FINAL}
\HMTRegisterLiveInputEdge
  {sections/hmt/16c_transportador_positivo_accion}
  {../colaboracion/parte_iii/source/public_final/ascensos_20260826/16d_medida_accion_positiva.tex + sucesor_102/wrappers/capitulo_38_jerarquia_unica.tex}
  {DOWNSTREAM-CH35-AND-CH38-FINAL}
\HMTRegisterLiveInputEdge
  {sections/hmt/delta_297/16c_entrelazador_dodecafase_witt}
  {sections/hmt/delta_297/16c_entrelazador_dodecafase_witt.tex}
  {UPSTREAM-PART-I-CH20}
\HMTRegisterLiveInputEdge
  {sections/hmt/delta_297/16d_medida_accion_positiva}
  {../colaboracion/parte_iii/source/public_final/ascensos_20260826/16d_medida_accion_positiva.tex}
  {DOWNSTREAM-CH35-FINAL}
\HMTRegisterLiveInputEdge
  {sections/md/04b_banda_particula_virtual}
  {integracion_83/parte_iv/chapters/75_definicion_particula.tex}
  {DOWNSTREAM-PART-IV-CH94}
\HMTRegisterLiveInputEdge
  {sections/md/06f_biografia_electron_rev6}
  {integracion_83/parte_iv/body/B001_ch48_electron_primera_realizacion_a3f11d97b332_06f_biografia_electron_rev6.tex}
  {DOWNSTREAM-PART-IV-CH58}
\HMTRegisterLiveInputEdge
  {sections/ampliacion_20260818/07_transduccion_metrologica_estados_genealogicos}
  {integracion_83/continuo_constantes/tex/transiciones/umbral_09_metrologia_y_programa_cuerpo.tex}
  {DOWNSTREAM-TRANSITION-III-IV}

% Al no reincluir cuerpos, las antiguas operaciones de descenso/restauración
% deben ser neutras. El grupo exterior impide que estas definiciones se filtren.
\renewcommand{\HMTDemoteHeadings}{}
\renewcommand{\HMTRestoreHeadings}{}

% El propietario conserva su cabecera primaria. Sus otras dos \chapter* eran
% rótulos de cuerpos que ahora residen después; se suprimen junto con las dos
% entradas de índice asociadas, evitando capítulos vacíos sin tocar el dueño.
\let\HMTAlphaOwnerChapter\chapter
\let\HMTAlphaOwnerAddContentsLine\addcontentsline
\newif\ifHMTAlphaOwnerPrimaryChapterSeen
\newif\ifHMTAlphaOwnerSuppressNextContentsLine
\HMTAlphaOwnerPrimaryChapterSeenfalse
\HMTAlphaOwnerSuppressNextContentsLinefalse
\long\def\chapter{%
  \@ifstar{\HMTAlphaOwnerStarChapter}{%
    \@ifnextchar[{\HMTAlphaOwnerOptionalChapter}{\HMTAlphaOwnerPlainChapter}}%
}
\long\def\HMTAlphaOwnerOptionalChapter[#1]#2{%
  \HMTAlphaOwnerPrimaryChapterSeentrue
  \HMTAlphaOwnerChapter[#1]{#2}%
}
\long\def\HMTAlphaOwnerPlainChapter#1{%
  \HMTAlphaOwnerPrimaryChapterSeentrue
  \HMTAlphaOwnerChapter{#1}%
}
\long\def\HMTAlphaOwnerStarChapter#1{%
  \ifHMTAlphaOwnerPrimaryChapterSeen
    \HMTAlphaOwnerSuppressNextContentsLinetrue
    \typeout{HMT-ALPHA-OWNER-EMPTY-HEADING-SKIP: \detokenize{#1}}%
  \else
    \HMTAlphaOwnerPrimaryChapterSeentrue
    \HMTAlphaOwnerChapter*{#1}%
  \fi
}
\renewcommand{\addcontentsline}[3]{%
  \ifHMTAlphaOwnerSuppressNextContentsLine
    \HMTAlphaOwnerSuppressNextContentsLinefalse
    \typeout{HMT-ALPHA-OWNER-EMPTY-TOC-SKIP: #1/#2}%
  \else
    \HMTAlphaOwnerAddContentsLine{#1}{#2}{#3}%
  \fi
}
\makeatother

\part{Clôture dodécaphasique et dépendances transversales de l'état}

\chapter[Dérivation interne de la constante de structure fine]{Dérivation interne de \texorpdfstring{\(\alpha_{\rm HMT}\)}{alpha HMT} : généalogie de survie, clôture dodécaphasique, deux constructions compatibles et bifurcation de Witt}
\label{chap:alpha-rev7}
\HMTClaim{HMT-R-ALPHA-DODECAPHASE}
\HMTClaim{HMT-R-DODEC-EXTRACTOR}

\paragraph{Quatre strates de \(\alpha_{\rm HMT}\).}
\begin{enumerate}
 \item Le constructeur généalogique est l'état terminal enrichi muni de ses deux expressions parallèles \(B_{\rm term}\) et \(U_{12}^{\rm sgn}\).
 \item La caractérisation analytique est la chaîne réversible \(U_{12}^{\rm sgn}\leftrightarrow K\to A\to\alpha_{12}\), suivie de son prolongement coinductif.
 \item Étant adimensionnelle, sa transduction commune n'introduit pas d'unité : elle fixe la carte angulaire \(A_{\rm deg}=1000\alpha_{\rm HMT}\) utilisée par les canaux ultérieurs.
 \item La comparaison avec CODATA ne s'ouvre qu'après l'obtention de \(\alpha_{\rm HMT}\) ; elle ne sélectionne ni \(K\), ni la retenue, ni le prolongement.
\end{enumerate}
La généalogie complète conserve explicitement la chaîne d'entrée et atteint un état terminal enrichi. Ses deux expressions sont parallèles :
\[
\begin{aligned}
 w_6\longrightarrow w_{30}\longrightarrow R_{36}\longrightarrow G_9
 \longrightarrow x_{\rm term}
 &\xrightarrow{\ \pi_{\rm term}\ }B_{\rm term},\\
 x_{\rm term}
 &\xrightarrow{\ R_{\rm sgn}\ }U_{12}^{\rm sgn}
 \xleftrightarrow{\ \mathcal H_{12}\ }K
 \longrightarrow A\longrightarrow\alpha_{12}.
\end{aligned}
\]
\(B_{\rm term}\) n'engendre pas \(U_{12}^{\rm sgn}\) : tous deux conservent des coordonnées différentes du même état. La transformation \(\mathcal H_{12}\) est inversible, et l'extracteur \((D_3K,D_4K,Q(K))\) est de rang \(12\). Dans ce domaine de primitives terminales explicites, la chaîne ne reçoit en entrée ni \(\alpha\), ni CODATA, ni une masse, ni aucune autre coordonnée métrologique. L'évaluation archimédienne forme \((P,E,\Phi,K,c_{13}=0)\), applique le cocycle de retenue et dérive \(\alpha_{\rm HMT}\). L'évaluation d'incidence transporte les mêmes mots et le même sceau vers \(B_K\subset H_\alpha^-\), qui constituera l'entrée de la branche Witt--Golay. La construction E21/22 et le prolongement distant \(\mathsf T_{\alpha,\infty}\) ---appelé \(T1\) dans la notation originale--- conservent leurs propres domaines et applications comme développements de la généalogie dodécaphasique principale.

\HMTDemoteHeadings
\HMTInputLive{sections/hmt/08_dodecafase_alpha}
\HMTInputLive{sections/hmt/08a_alpha_dos_vias_t1_rev8}
\HMTInputLive{sections/hmt/08b_alpha_t1_10000_rev9}
\HMTRestoreHeadings

\HMTInputLive{sections/ampliacion_20260818/03_extension_holografica_euler}

\HMTInputLive{sections/hmt/07b_realizacion_analitica_contraangulo}
\HMTInputLive{sections/hmt/16c_transportador_positivo_accion}
\HMTInputLive{sections/hmt/delta_297/16c_entrelazador_dodecafase_witt}
\begingroup
\def\HMTInternalTraceability{1}
\HMTInputLive{sections/hmt/delta_297/16d_medida_accion_positiva}
\endgroup

\HMTInputLive{sections/md/04b_banda_particula_virtual}

\chapter*{Construction holographique complète de l'électron comme dépendance du lecteur}
\addcontentsline{toc}{chapter}{Construction holographique complète de l'électron comme dépendance du lecteur}
\HMTInputLive{sections/md/06f_biografia_electron_rev6}

\chapter*{Transduction métrologique des états généalogiques}
\addcontentsline{toc}{chapter}{Transduction métrologique des états généalogiques}
\HMTInputLive{sections/ampliacion_20260818/07_transduccion_metrologica_estados_genealogicos}

\end{HMTScientificOwnerSubordinated}
\let\part\HMTRecoveredOriginalPart
\endgroup

\HMTFlushCurrentChapter
\chapter{Clôture dodécaphasique de \texorpdfstring{\(\alpha_{\mathrm{HMT}}\)}{alpha HMT} et lecteur analytique \texorpdfstring{\(E_{21/22}^{(9)}\)}{E21/22 d'ordre 9}: génération réversible, racine simple et comparateur typé}%
\begingroup
\renewcommand{\chapter}[2][]{}%
\chapter{Dérivations internes de la constante de structure fine \(\alpha\) : clôture entière dodécaphasique réversible et racine analytique unique du noyau \(E_{21/22}\)}
\label{chap:p3-final:31:dos_derivaciones_alpha}
\label{chap:p3-83-25-alpha}

\section[Dérivation interne de \texorpdfstring{\(\alpha_{\rm HMT}\)}{alpha HMT}]
{Dérivation interne de \texorpdfstring{\(\alpha_{\rm HMT}\)}{alpha HMT} :
état terminal commun, sceau dodécaphasique et cocycle de retenue}
\label{p3-83-c25-s1:chap:parte-iii-alpha-hmt}

Le générateur précédent atteint un état terminal enrichi avec deux
publications coordonnées :
\[
\begin{aligned}
w_6\longrightarrow w_{30}\longrightarrow R_{36}\longrightarrow G_9
\longrightarrow x_{\rm term}
&\xrightarrow{\ \pi_{\rm term}\ }B_{\rm term},\\
x_{\rm term}
&\xrightarrow{\ R_{\rm sgn}\ }U_{12}^{\rm sgn}
\xleftrightarrow{\ \mathcal H_{12}\ }K
\longrightarrow A\longrightarrow\alpha_{12}.
\end{aligned}
\]
Les deux flèches terminales publient des coordonnées distinctes du même antécédent.
La transformée \(\mathcal H_{12}\) est réversible ; le système transversal
\((D_3K,D_4K,Q(K))\) atteint le rang douze ; la division euclidienne orientée conserve
le cocycle de retenue et les conditions aux frontières. La chaîne utilise
l’état signé, le sceau et les primitives déclarées, et obtient
\(\alpha_{\rm HMT}\) avant d’ouvrir la reconnaissance métrologique.

\subsection{Clôture dodécaphasique et constructions compatibles}
La clôture dodécaphasique fournit l’énoncé et la preuve directrice. Les
constructions E21/22 et \(\mathsf T_{\alpha,\infty}\) conservent leurs domaines et leurs
fonctions propres dans la généalogie principale.


\section{Registre signé, reconstruction du sceau et retenue entière}
Le registre signé et le cocycle de retenue déploient l’état terminal
commun, les cartes de Hadamard et de différences transversales, la forme normale de
Smith, l’identité affine, le prolongement coinductif et les frontières
\(662/663\).


\section{Réversibilité, indépendance et non-circularité de la clôture dodécaphasique}
L’état \(108\), le constructeur prospectif de \(K\), la transformée de
Hadamard bidirectionnelle, la retenue entière et le prolongement infini
s’articulent au moyen de contrôles explicites de réversibilité, d’indépendance et de
non-circularité. Chaque formulation conserve son domaine et son critère de
réfutation.


% Unidad U016. Registro firmado, lectura terminal y reconstruccion de K.
% Redaccion autonoma desde la autoridad material 1063 y su sucesor V2.

\section{Factorisation du registre dodécaphasique signé et clôture entière du sceau \(K\)}
\label{p3-83-c25-s4:chap:registro-dodecafasico-hadamard-k}

La construction du sceau entier ne commence ni par \(\alpha\), ni par son inverse
ni par un tableau métrologique. Son domaine est l’état terminal enrichi
du Topological Phase Kernel obtenu après la chaîne déjà construite
\begin{equation}
 \boxed{
 \mathrm{APP}\longrightarrow\mathrm{TRIT}\longrightarrow\mathrm{TPK}
 \longrightarrow\Omega_{\rm seed}
 \xrightarrow{T_{\min}}\mathcal U_6^{\rm cat}
 \xrightarrow{\rho_3}W_{243}
 \longrightarrow w_{30}\longrightarrow R_{36}
 \longrightarrow G_9\longrightarrow x_{\rm term}.}
 \label{p3-83-c25-s4:eq:cadena-upstream-registro-dodecafasico}
\end{equation}
Les flèches précédentes ont été définies avec leurs domaines respectifs :
paires de germes, signatures, émissions, mots ternaires, frontières,
histoires compatibles et holonomie nonadique. Cette unité ne choisit à nouveau
aucune de ces données. Elle construit les cartes entières du même état
terminal et démontre leur équivalence.

\subsection{Fenêtres de la trace orientée de longueur \(108\)}

Soit
\[
 \Theta_{108}=\mathbb Z/108\mathbb Z
\]
le calendrier observable. Ses douze fenêtres nonadiques sont
\[
 E_m=\{9m+1,\ldots,9m+9\},
 \qquad 0\leq m<12.
\]
La partition est disjointe et exhaustive :
\[
 \{1,\ldots,108\}=\bigsqcup_{m=0}^{11}E_m.
\]
L’ensemble des codes de direction qui publient un événement est fixé comme
\[
 D_{\rm evt}=\{2,4,6,8\}.
\]
L’orientation des douze fenêtres est
\begin{equation}
 \Sigma_{12}=(+,+,+,-,-,-,+,+,+,-,-,-).
 \label{p3-83-c25-s4:eq:signatura-doce-ventanas}
\end{equation}
Si \(d_t\) est la direction codée au pas \(t\), l’événement signé
de la fenêtre \(m\) est
\begin{equation}
 e_m=\Sigma_{12}(m)\,
 \#\{t\in E_m:d_t\in D_{\rm evt}\}.
 \label{p3-83-c25-s4:eq:evento-firmado-ventana}
\end{equation}
La formule distingue le nombre non négatif d’événements et le signe de
l’orientation. Elle ne remplace pas une direction cardinale par un signe : les deux
coordonnées restent présentes dans l’état TPK.

\begin{definition}[Registre enrichi dodécaphasique]
\label{p3-83-c25-s4:def:registro-enriquecido-dodecafasico}
Le registre transporté par la trace est
\[
 \mathcal R_{12}(x)=
 \bigl((E_m,\tau_m,\sigma_m,\kappa_m,\Lambda_m)\bigr)_{m=0}^{11},
\]
où \(\tau_m\) est l’orientation TRIT, \(\sigma_m\) la signature de
frontière, \(\kappa_m\) la retenue et \(\Lambda_m\) le segment de registre généalogique
qui conserve la route. Il existe une chaîne de projections
\[
 \mathcal X_{\rm TPK}^{\rm enr}
 \longrightarrow\mathcal R_{12}
 \longrightarrow\Theta_{108},
\]
mais aucune des deux flèches n’est une identification.
\end{definition}

En particulier, \(\Theta_{108}\) conserve uniquement la position du calendrier.
Le registre \(\mathcal R_{12}\) conserve en outre l’orientation, la signature,
la retenue et la provenance. L’état complet conserve encore le cylindre,
la frontière, le quotient, le terminal et la mémoire verticale.

\subsection{Carte entière de l’opposition dodécaphasique}

Pour \(0\leq i<6\), les positions opposées déterminent
\begin{equation}
 p_i=e_i+e_{i+6},
 \qquad
 a_i=e_i-e_{i+6}.
 \label{p3-83-c25-s4:eq:pares-opuestos-registro}
\end{equation}
La charge centrale et les cinq marges relatives sont
\begin{equation}
 s_C=\sum_{i=0}^{5}p_i,
 \qquad
 M_i=p_i-p_5\quad(0\leq i<5).
 \label{p3-83-c25-s4:eq:carga-margenes-registro}
\end{equation}

\begin{definition}[Carte intégrale \(1+5+6\)]
\label{p3-83-c25-s4:def:carta-integral-156}
On définit
\begin{equation}
 \mathcal L_{12}(e_0,\ldots,e_{11})
 =\bigl(s_C,M_0,\ldots,M_4,a_0,\ldots,a_5\bigr).
 \label{p3-83-c25-s4:eq:carta-integral-156}
\end{equation}
La partition des coordonnées est
\[
 1+5+6=12.
\]
\end{definition}

\begin{theorem}[Inversion exacte de la carte entière]
\label{p3-83-c25-s4:thm:inversion-carta-integral-156}
Un tuple
\[
 \ell=(s_C,M_0,\ldots,M_4,a_0,\ldots,a_5)\in\mathbb Z^{12}
\]
appartient à l’image de \(\mathcal L_{12}\) si et seulement si
\begin{align}
 s_C-\sum_{i=0}^{4}M_i&\equiv0\pmod6,
 \label{p3-83-c25-s4:eq:congruencia-seis-carta}\\
 p_i&\equiv a_i\pmod2
 \quad(0\leq i<6),
 \label{p3-83-c25-s4:eq:congruencia-paridad-carta}
\end{align}
où
\begin{align}
 p_5&=\frac{s_C-\sum_{i=0}^{4}M_i}{6},
 \label{p3-83-c25-s4:eq:inversa-p5-carta}\\
 p_i&=p_5+M_i\quad(0\leq i<5).
 \label{p3-83-c25-s4:eq:inversa-pi-carta}
\end{align}
Dans ce cas, l’unique antécédent entier est
\begin{equation}
 e_i=\frac{p_i+a_i}{2},
 \qquad
 e_{i+6}=\frac{p_i-a_i}{2}.
 \label{p3-83-c25-s4:eq:inversa-eventos-carta}
\end{equation}
\end{theorem}

\begin{proof}
En sommant \(M_i=p_i-p_5\) pour \(0\leq i<5\), nous obtenons
\[
 \sum_{i=0}^{4}M_i
 =\sum_{i=0}^{4}p_i-5p_5
 =s_C-6p_5,
\]
ce qui impose \eqref{p3-83-c25-s4:eq:inversa-p5-carta} et la congruence modulo six.
Une fois les six \(p_i\) reconstruits, le système
\[
 p_i=e_i+e_{i+6},
 \qquad a_i=e_i-e_{i+6}
\]
admet l’unique solution \eqref{p3-83-c25-s4:eq:inversa-eventos-carta}. Cette solution est
entière exactement lorsque \(p_i\) et \(a_i\) ont la même parité.
Les deux familles de congruences sont donc nécessaires et suffisantes.
\end{proof}

Ce théorème est antérieur à la clôture numérique de \(\alpha\). Il explique comment la
trace orientée produit une carte entière réversible, et pourquoi il ne suffit pas
de conserver le calendrier réduit.

\subsection{Sélection de l’état terminal commun}

Les trois sections compatibles atteignent la profondeur terminale avec des catalogues
de cardinalités
\[
 |\mathcal T_\pi|=196,
 \qquad
 |\mathcal T_e|=197,
 \qquad
 |\mathcal T_\varphi|=196.
\]
Le produit contient
\begin{equation}
 196\cdot197\cdot196=7\,567\,952
 \label{p3-83-c25-s4:eq:cardinal-producto-terminal}
\end{equation}
triplets ordonnés.

La marge de colonnes du lecteur terminal est
\[
 q_\partial=(1,2,2,1,2,0)\in\mathbb F_3^6,
\]
et l’orientation des lignes du canal conjoint est
\[
 \sigma=(1,1,-1)=(1,1,2)\in\mathbb F_3^3.
\]
La première donnée agit sur les colonnes ; la seconde sur les lignes. L’une ne se déduit pas
de l’autre.

Les signatures locales de Gram, de norme, de poids et de distance réduisent
\eqref{p3-83-c25-s4:eq:cardinal-producto-terminal} à \(331\) triplets. La marge
\(q_\partial\) laisse deux indices :
\[
 (120,197,157),
 \qquad
 (129,197,148).
\]
Leurs matrices sont
\[
 B_0=
 \begin{pmatrix}
 0&2&0&2&2&0\\
 0&0&1&2&2&0\\
 1&0&1&0&1&0
 \end{pmatrix},
 \qquad
 B_1=
 \begin{pmatrix}
 0&2&1&0&2&0\\
 0&0&1&2&2&0\\
 1&0&0&2&1&0
 \end{pmatrix}.
\]
Les sommes des lignes sont
\[
 (0,2,0),
 \qquad
 (2,2,1).
\]
Comme
\[
 \langle\sigma\rangle
 =\{(0,0,0),(1,1,2),(2,2,1)\},
\]
seul \(B_1\) satisfait la condition axiale.

\begin{theorem}[Unicité de la publication terminale visible]
\label{p3-83-c25-s4:thm:unicidad-publicacion-terminal-visible}
La sélection conjointe détermine
\begin{equation}
 B_{\rm term}=
 \begin{pmatrix}
 021020\\001220\\100210
 \end{pmatrix},
 \label{p3-83-c25-s4:eq:matriz-terminal-visible}
\end{equation}
correspondant au triplet \((129,197,148)\). Aucune évaluation décimale de
\(\alpha\) n’intervient dans le recensement ni dans les deux conditions de sélection.
\end{theorem}

\begin{proof}
Le recensement et la marge laissent exactement \(B_0,B_1\). La charge de \(B_0\)
n’appartient pas à \(\langle\sigma\rangle\), tandis que celle de \(B_1\) est
\(2\sigma\). La seconde candidate est donc la seule qui satisfasse
simultanément les deux exigences typées.
\end{proof}

\subsection{2 publications du même état et absence de conversion rétrospective}

Soit
\[
 x_{\rm term}\in
 \mathcal X_{\rm TPK}^{\rm term,enr}
 \subseteq\mathcal X_{\rm TPK}^{[108]}.
\]
La lecture visible est
\[
 \pi_{\rm term}(x_{\rm term})=B_{\rm term}.
\]
La carte signée conserve une autre coordonnée du même état.

\begin{definition}[Lecture causale signée terminale]
\label{p3-83-c25-s4:def:subespacio-firmado-terminal}
Le domaine n’incorpore comme coordonnées d’entrée ni \(K\) ni \(U\).
L’état \(x_{\rm term}\), produit par la composition complète
\(\mathrm{APP}\to\mathrm{TRIT}\to\mathrm{TPK}\), contient le registre
terminal enrichi, et l’extracteur de ses douze fenêtres orientées est
\begin{equation}
 R_{12}:\mathcal X_{\rm TPK}^{\rm term,enr}
 \longrightarrow\mathcal R_{12}^{\rm or},
 \qquad
 R_{12}(x_{\rm term})
 =\bigl((E_m,\tau_m,\sigma_m,\kappa_m,\Lambda_m)\bigr)_{m=0}^{11}.
 \label{p3-83-c25-s4:eq:subespacio-firmado-terminal}
\end{equation}
Le lecteur de canaux \(90/120\) du registre émet, sans recevoir le résultat
comme donnée,
\[
 E_{108}^{90,120}:\mathcal R_{12}^{\rm or}
 \longrightarrow\mathbb Z^{12}\times\mathbb Z^{12}\times\mathbb Z,
 \qquad
 E_{108}^{90,120}\bigl(R_{12}(x_{\rm term})\bigr)
 =(b_{90},b_{120},Q_{\rm TPK}),
\]
avec
\begin{align*}
 b_{90}={}&(-495,-116,-684,108,601,-90,
             -173,-88,313,560,-397,461),\\
 b_{120}={}&(-425,-281,-481,671,-255,30,
              475,-543,680,251,6,-128),\\
 Q_{\rm TPK}={}&6263.
\end{align*}
Pour \(r=1,2,3\) et \(j=0,1,2,3\), on définit
\[
\begin{aligned}
 B_r&=\sum_{j=0}^{3}b_{120,r+3j},
 &d_{r,j}&=b_{90,r+3j},\\
 A_1&=\frac{Q_{\rm TPK}+2B_1+B_2}{3},
 &A_2&=A_1-B_1,
 &A_3&=A_2-B_2.
\end{aligned}
\]
La publication \(\Pi_H\) forme les trois blocs
\[
 U^{(r)}=
 \bigl(A_r,d_{r,1}-d_{r,3},d_{r,0}+d_{r,2},
 d_{r,0}-d_{r,2}\bigr),
\]
et les réentrelace dans l’ordre dodécaphasique, c’est-à-dire
\[
 \Pi_H(b_{90},b_{120},Q_{\rm TPK})
 :=\operatorname{reint}_3\bigl(U^{(1)},U^{(2)},U^{(3)}\bigr)=U.
\]
Soit \(\mathcal H_{12}\) la
transformation entière qui sera construite dans la section suivante. Le module
de publication est
\begin{equation}
 \mathcal U_{12}^{\rm sgn}
 :=\operatorname{im}
 \bigl(\mathcal H_{12}:\mathbb Z^{12}\to\mathbb Z^{12}\bigr).
 \label{p3-83-c25-s4:eq:modulo-publicacion-firmada}
\end{equation}
La lecture signée est, causalement, la composition
\begin{equation}
 R_{\rm sgn}:=
 \Pi_H\circ E_{108}^{90,120}\circ R_{12}:
 \mathcal X_{\rm TPK}^{\rm term,enr}
 \longrightarrow\mathcal U_{12}^{\rm sgn},
 \qquad
 R_{\rm sgn}(x_{\rm term})=U.
 \label{p3-83-c25-s4:eq:lector-firmado-terminal}
\end{equation}
Ce n’est qu’après cette sortie que l’on reconstruit
\(K=\mathcal H_{12}^{-1}U\) ; l’identité
\(U=\mathcal H_{12}K\) certifie cette reconstruction et ne définit pas le domaine
du lecteur.
\end{definition}

L’écriture conjointe correcte est
\begin{equation}
 x_{\rm term}
 \xrightarrow{(\pi_{\rm term},R_{\rm sgn})}
 (B_{\rm term},U_{12}^{\rm sgn}).
 \label{p3-83-c25-s4:eq:doble-publicacion-terminal}
\end{equation}
On n’interpose pas une flèche
\(B_{\rm term}\to U_{12}^{\rm sgn}\) : la matrice visible a précisément oublié
les coordonnées d’orientation, d’opposition, de quotient, de retenue et de
registre généalogique utilisées par la carte signée. La source des deux publications est
l’état enrichi produit par \eqref{p3-83-c25-s4:eq:cadena-upstream-registro-dodecafasico}.

L’évaluation du registre signé donne
\begin{equation}
 \boxed{
 U=(2378,1406,2479,-452,998,-551,
 -668,-204,-371,-322,-28,-997).}
 \label{p3-83-c25-s4:eq:vector-u-firmado-autonomo}
\end{equation}

\subsection{Naturalité par rapport à la profondeur}

Pour \(N'\geq N\geq108\), les troncatures du système signé sont
\begin{equation}
 \tau_{N',N}^{\rm sgn}
 x_{\rm term}^{[N']}
 =\tau_{N',N}x_{\rm term}^{[N']}.
 \label{p3-83-c25-s4:eq:truncamiento-preserva-ku-autonomo}
\end{equation}
Par conséquent, la troncature agit uniquement sur l’histoire et l’état
enrichi ; elle ne transporte comme entrées indépendantes ni \(K\) ni \(U\).

\begin{theorem}[Carré de naturalité de la lecture signée]
\label{p3-83-c25-s4:thm:naturalidad-lector-firmado-autonomo}
On a
\begin{equation}
 \boxed{
 R_{{\rm sgn},N}\circ\tau_{N',N}^{\rm sgn}
 =\operatorname{id}_{\mathcal U_{12}^{\rm sgn}}
  \circ R_{{\rm sgn},N'}.}
 \label{p3-83-c25-s4:eq:cuadrado-naturalidad-lector-firmado}
\end{equation}
\end{theorem}

\begin{proof}
La naturalité des états conserve littéralement les cent huit premières
fenêtres du registre :
\[
 R_{12,N}\circ\tau_{N',N}^{\rm sgn}=R_{12,N'}.
\]
La troncature peut modifier uniquement l’histoire ultérieure. Comme
\(E_{108}^{90,120}\) et \(\Pi_H\) dépendent exclusivement de ces fenêtres,
les deux membres renvoient le même \(U\). La commutativité est une égalité
d’applications, et non une analogie entre diagrammes.
\end{proof}

\subsection{Orbites de pas \(3\) et transformée de Hadamard}

L’ordre cyclique du sceau se répartit dans les trois orbites
\begin{align}
 K^{(1)}&=(K_1,K_4,K_7,K_{10}),
 \label{p3-83-c25-s4:eq:orbita-k-1}\\
 K^{(2)}&=(K_2,K_5,K_8,K_{11}),
 \label{p3-83-c25-s4:eq:orbita-k-2}\\
 K^{(3)}&=(K_3,K_6,K_9,K_{12}).
 \label{p3-83-c25-s4:eq:orbita-k-3}
\end{align}
Soit
\begin{equation}
 H_4=
 \begin{pmatrix}
 1&1&1&1\\
 1&1&-1&-1\\
 1&-1&1&-1\\
 1&-1&-1&1
 \end{pmatrix}.
 \label{p3-83-c25-s4:eq:matriz-hadamard-cuatro-autonoma}
\end{equation}
Si \(P_{\rm orb}\) réordonne les douze coordonnées conformément à
\eqref{p3-83-c25-s4:eq:orbita-k-1}--\eqref{p3-83-c25-s4:eq:orbita-k-3}, la transformation globale est
\begin{equation}
 \mathcal H_{12}
 =P_{\rm orb}^{-1}
 (H_4\oplus H_4\oplus H_4)P_{\rm orb}.
 \label{p3-83-c25-s4:eq:hadamard-global-doce}
\end{equation}

Les trois blocs de \eqref{p3-83-c25-s4:eq:vector-u-firmado-autonomo}, lus par
orbites, sont
\begin{align*}
 U^{(1)}&=(2378,-452,-668,-322),\\
 U^{(2)}&=(1406,998,-204,-28),\\
 U^{(3)}&=(2479,-551,-371,-997).
\end{align*}

\begin{lemma}[Quart de tour algébrique de Hadamard]
\label{p3-83-c25-s4:lem:cuarto-inversa-hadamard}
La matrice \(H_4\) satisfait
\begin{equation}
 H_4^2=4I_4,
 \qquad
 H_4^{-1}=\frac14H_4,
 \qquad
 \det H_4=-16.
 \label{p3-83-c25-s4:eq:identidades-hadamard-cuatro}
\end{equation}
\end{lemma}

\begin{proof}
Les lignes de \(H_4\) sont orthogonales et chacune possède une norme au carré égale à quatre.
Comme la matrice est symétrique, \(H_4^{\mathsf T}H_4=H_4^2=4I_4\). La
formule de l’inverse et le déterminant s’en déduisent immédiatement.
\end{proof}

\begin{theorem}[Reconstruction entière unique du sceau]
\label{p3-83-c25-s4:thm:reconstruccion-entera-k-hadamard}
L’inversion
\begin{equation}
 K^{(r)}=\frac14H_4U^{(r)}
 \qquad(r=1,2,3)
 \label{p3-83-c25-s4:eq:inversion-bloques-hadamard}
\end{equation}
produit
\begin{align*}
 K^{(1)}&=(234,729,621,794),\\
 K^{(2)}&=(543,659,058,146),\\
 K^{(3)}&=(140,824,914,601).
\end{align*}
En défaisant l’ordre orbital, on obtient
\begin{equation}
 \boxed{
 K=(234,543,140,729,659,824,621,058,914,794,146,601).}
 \label{p3-83-c25-s4:eq:sello-k-doce-autonomo}
\end{equation}
C'est l'unique préimage rationnelle de \(U\), et elle appartient à \(\mathbb Z^{12}\).
\end{theorem}

\begin{proof}
Par \eqref{p3-83-c25-s4:eq:identidades-hadamard-cuatro}, l’antécédent rationnel est unique.
Les multiplications explicites sont
\begin{align*}
 H_4U^{(1)}&=(936,2916,2484,3176),\\
 H_4U^{(2)}&=(2172,2636,232,584),\\
 H_4U^{(3)}&=(560,3296,3656,2404).
\end{align*}
Toutes les coordonnées sont divisibles par quatre. La division et le réentrelacement
produisent \eqref{p3-83-c25-s4:eq:sello-k-doce-autonomo}.
\end{proof}

Le facteur \(1/4\) n’est pas choisi pour ajuster une sortie : c’est l’inverse
imposé par \(H_4^2=4I_4\). Ce n’est pas non plus la racine carrée de \(-1\). La
structure complexe interne \(A_W^2=-I\) apparaîtra après la construction de la
matrice de Paley ; ce sont deux faits de types différents.

\subsection{Structure entière et forme normale de Smith}

\begin{theorem}[Image entière de \(H_4\)]
\label{p3-83-c25-s4:thm:smith-hadamard-autonomo}
La forme normale de Smith est
\begin{equation}
 \operatorname{SNF}(H_4)=\operatorname{diag}(1,2,2,4).
 \label{p3-83-c25-s4:eq:snf-hadamard-autonoma}
\end{equation}
En conséquence,
\begin{equation}
 \operatorname{coker}\mathcal H_{12}
 \cong(\mathbb Z/2\mathbb Z)^6
 \oplus(\mathbb Z/4\mathbb Z)^3,
 \label{p3-83-c25-s4:eq:cociente-hadamard-doce}
\end{equation}
et l’image de \(\mathcal H_{12}\) a pour indice
\[
 16^3=4096
\]
dans \(\mathbb Z^{12}\). Pour \(u\in\mathbb Z^4\),
\begin{equation}
 u\in H_4\mathbb Z^4
 \quad\Longleftrightarrow\quad
 H_4u\equiv0\pmod4.
 \label{p3-83-c25-s4:eq:criterio-imagen-hadamard}
\end{equation}
\end{theorem}

\begin{proof}
Le plus grand commun diviseur des entrées est un. Les mineurs d’ordre deux sont
pairs et l’un d’eux a pour valeur absolue deux ; ceux d’ordre trois sont des multiples de
quatre et l’un d’eux a pour valeur absolue quatre ; le déterminant a pour module
seize. Les diviseurs déterminantaux sont \(1,2,4,16\), d’où l’on obtient
les facteurs invariants \(1,2,2,4\). La somme directe triple répète les
facteurs et multiplie les indices. Enfin,
\(H_4^{-1}=H_4/4\), de sorte que l’antécédent est entier exactement sous la
congruence \eqref{p3-83-c25-s4:eq:criterio-imagen-hadamard}.
\end{proof}

\subsection{Certification et reconstruction par différences de pas \(3\) et \(4\)}

Une fois que le registre a émis \((b_{90},b_{120},Q_{\rm TPK})\), que
\(\Pi_H\) a produit \(U\) et que Hadamard a reconstruit \(K\), les
coordonnées suivantes agissent en aval comme certificat indépendant de
ces sorties préalables. Avec des indices cycliques modulo douze, nous définissons
\begin{equation}
 (D_3K)_m=K_m-K_{m+3},
 \qquad
 (D_4K)_m=K_m-K_{m+4},
 \qquad
 Q(K)=\sum_{m=1}^{12}K_m.
 \label{p3-83-c25-s4:eq:extractor-diferencias-k}
\end{equation}
Pour \eqref{p3-83-c25-s4:eq:sello-k-doce-autonomo}, on vérifie
\begin{align*}
 D_3K=b_{90}={}&(-495,-116,-684,108,601,-90,\\
         &\qquad-173,-88,313,560,-397,461),\\
 D_4K=b_{120}={}&(-425,-281,-481,671,-255,30,\\
         &\qquad475,-543,680,251,6,-128),\\
 Q(K)=Q_{\rm TPK}={}&6263.
\end{align*}

\begin{theorem}[Complétude de l’extracteur transversal]
\label{p3-83-c25-s4:thm:completitud-extractor-transversal-k}
L'opérateur
\[
 \mathcal E_{\rm dod}=(D_3,D_4,Q):
 \mathbb Z^{12}\longrightarrow\mathbb Z^{25}
\]
a rang douze et pour forme normale de Smith
\begin{equation}
 \operatorname{SNF}(\mathcal E_{\rm dod})
 =\operatorname{diag}(\underbrace{1,\ldots,1}_{11},12).
 \label{p3-83-c25-s4:eq:snf-extractor-transversal-k}
\end{equation}
En particulier, \((D_3K,D_4K,Q(K))\) détermine un unique \(K\). Cette
unicité certifie et reconstruit les données déjà émises par le registre ; elle ne les
engendre ni ne les sélectionne.
\end{theorem}

\begin{proof}
Si \(D_3v=D_4v=0\), alors \(v\) est constant le long des pas
trois et quatre. Comme \(\gcd(3,4)=1\), ces pas engendrent tout
\(\mathbb Z/12\mathbb Z\), de sorte que
\[
 \ker D_3\cap\ker D_4=\langle\mathbf1\rangle.
\]
La charge totale ne s’annule pas sur cette droite, car \(Q(\mathbf1)=12\). Le rang
est douze. Les lignes de \((D_3,D_4)\) sont les incidences orientées d’un graphe
de Cayley connexe ; un arbre couvrant apporte onze facteurs invariants
unitaires. Tout mineur maximal doit contenir la ligne \(Q\), et les opérations
unimodulaires réduisent sa dernière contribution à douze. Il existe un mineur maximal de
module douze ; le dernier facteur est donc exactement douze.
\end{proof}

La reconstruction certificatrice de \(U\) à partir de ces coordonnées reproduit la
publication préalable de \(\Pi_H\) et ne nécessite pas de multiplier à nouveau par
\(H_4\). Pour \(r=1,2,3\), soit
\begin{equation}
 B_r=\sum_{j=0}^{3}(D_4K)_{r+3j}
 =\sum_{j=0}^{3}b_{120,r+3j}.
 \label{p3-83-c25-s4:eq:br-diferencias-k}
\end{equation}
Alors
\[
 (B_1,B_2,B_3)=(972,-1073,101).
\]
Les sommes orbitales sont
\begin{align}
 A_1&=\frac{Q+2B_1+B_2}{3}=2378,
 \label{p3-83-c25-s4:eq:a1-reconstruccion-u}\\
 A_2&=A_1-B_1=1406,
 \label{p3-83-c25-s4:eq:a2-reconstruccion-u}\\
 A_3&=A_2-B_2=2479.
 \label{p3-83-c25-s4:eq:a3-reconstruccion-u}
\end{align}
Si \(d_{r,j}=(D_3K)_{r+3j}=b_{90,r+3j}\), le bloc signé de l’orbite
\(r\) est
\begin{equation}
 \bigl(A_r,
 d_{r,1}-d_{r,3},
 d_{r,0}+d_{r,2},
 d_{r,0}-d_{r,2}\bigr),
 \label{p3-83-c25-s4:eq:bloque-firmado-desde-diferencias}
\end{equation}
ce qui reproduit les trois blocs \(U^{(r)}\) déjà publiés par \(\Pi_H\).

\subsection{Tétrade de seuil du sceau}

La complétion décimale locale satisfait
\[
 1000=729+271.
\]
Le support des coordonnées du sceau qui atteignent le seuil \(729\) est
\begin{equation}
 \boxed{
 B_K:=\{m:K_m\geq729\}=\{4,6,9,10\}.}
 \label{p3-83-c25-s4:eq:tetrada-umbral-k}
\end{equation}
Cette tétrade s’obtient après la construction de \(K\). Dans cette unité, elle n’est pas
encore appelée étoile de Witt : un tel objet ne sera défini qu’après
la construction du code ternaire, de ses hexades et du système de Steiner.

\subsection{Indépendance causale à l’égard de la sortie physique}

\begin{theorem}[Isolement de \(U\) et \(K\)]
\label{p3-83-c25-s4:thm:aislamiento-u-k-alpha-codata}
La composition
\begin{equation}
 \Omega_{\rm seed}\longrightarrow x_{\rm term}
 \xrightarrow{R_{\rm sgn}}U
 \xrightarrow{\mathcal H_{12}^{-1}}K
 \label{p3-83-c25-s4:eq:composicion-upstream-u-k}
\end{equation}
ne reçoit en entrée ni \(\alpha\), ni \(\alpha^{-1}\), ni CODATA, ni une suite
décimale cible. Le sceau \(K\) est déterminé avant la définition de la
récurrence qui produira \(\alpha\).
\end{theorem}

\begin{proof}
La première partie de la composition utilise uniquement des germes, des émissions,
des lecteurs ternaires, des opérateurs \(L_0,L_1\), des relations de survie,
le transport de phase et les cartes de l’état TPK. La lecture
\(R_{\rm sgn}=\Pi_H\circ E_{108}^{90,120}\circ R_{12}\) est la composition
causale du registre terminal, et la dernière flèche est la transformation linéaire
\(\mathcal H_{12}^{-1}=\mathcal H_{12}/4\)
sur son image entière. Aucune des formules ne contient une variable
physique ou un tableau métrologique. La récurrence en base mille n’est introduite que
dans l’unité suivante, après avoir fixé \(K\).
\end{proof}

\subsection{Obstructions de la clôture dodécaphasique : sceau, retenue, réversibilité et non-circularité}

\begin{remark}[Certificat et contrôle de validité : réfutation du registre et du sceau]
La construction est réfutée si l’un quelconque des faits
suivants est vérifié :
\begin{enumerate}
 \item les fenêtres \(E_m\) ne partitionnent pas les cent huit pas ;
 \item la carte \(\mathcal L_{12}\) ne respecte pas son inverse ou ses congruences ;
 \item le recensement terminal ne produit pas \(331\) triplets, deux matrices et une unique
       charge axiale ;
 \item on tente d’obtenir \(U\) à partir de \(B_{\rm term}\) après avoir oublié les
       coordonnées signées ;
 \item le carré de naturalité
       \eqref{p3-83-c25-s4:eq:cuadrado-naturalidad-lector-firmado} échoue ;
 \item \(H_4^2\neq4I_4\), un antécédent n’est pas entier ou le vecteur
       reconstruit diffère de \eqref{p3-83-c25-s4:eq:sello-k-doce-autonomo} ;
 \item la forme normale de Smith diffère de
       \eqref{p3-83-c25-s4:eq:snf-hadamard-autonoma} ;
 \item le certificat \((D_3,D_4,Q)\) perd du rang, ne coïncide pas avec les
       émissions préalables du registre ou ne reproduit pas \(U\) ;
 \item \(B_K\) est introduit avant la construction de \(K\), ou on lui attribue une
       structure de Witt avant la définition de \(W_{12}\) ;
 \item un chiffre de \(\alpha\) ou une référence métrologique intervient dans
       l’une quelconque des flèches de \eqref{p3-83-c25-s4:eq:composicion-upstream-u-k}.
\end{enumerate}
\end{remark}


% Unidad U017. Cierre afín, prolongación y semántica de frontera de alpha.
% Redacción autónoma desde la autoridad material 1063 y su sucesor V2.

\section{Prolongement coinductif du cocycle dodécaphasique et compatibilité de la frontière \texorpdfstring{\(662/663\)}{662/663}}
\label{p3-83-c25-s5:chap:acarreo-alpha-frontera-662-663}

\subsection{Domaine arithmétique de la clôture}

La construction reçoit quatre coordonnées dodécaphasiques obtenues dans les
unités précédentes :
\[
\begin{aligned}
 P={}&(141,592,653,589,793,238,462,643,383,279,502,884),\\
 E={}&(718,281,828,459,045,235,360,287,471,352,662,497),\\
 \Phi={}&(618,033,988,749,894,848,204,586,834,365,638,117),\\
 K={}&(234,543,140,729,659,824,621,058,914,794,146,601).
\end{aligned}
\]
Les trois premiers vecteurs sont les douze premières triades des lectures
archimédiennes fractionnaires de \(\pi\), \(e\) et \(\varphi\). Le quatrième est
la coordonnée dodécaphasique de l’état terminal enrichi du Topological
Phase Kernel, construite avant l’ouverture de l’opération de retenue.

Nous fixons la base
\[
 B:=10^3=1000
\]
et la concaténation
\begin{equation}
 \iota(v_1,\ldots,v_{12})
 :=\sum_{m=1}^{12}v_mB^{12-m}.
 \label{p3-83-c25-s5:eq:concatenacion-base-mil-alpha}
\end{equation}
\Needspace{8\baselineskip}
En particulier,
\[
\begin{aligned}
 \iota(P)&=141592653589793238462643383279502884,\\
 \iota(E)&=718281828459045235360287471352662497,\\
 \iota(\Phi)&=618033988749894848204586834365638117,\\
 \iota(K)&=234543140729659824621058914794146601.
\end{aligned}
\]

L’opération étudiée dans cette unité est l’application
\begin{equation}
 \mathfrak C_{12}:
 (P,E,\Phi,K,c_{13})
 \longmapsto(A,C),
 \label{p3-83-c25-s5:eq:mapa-cierre-dodecafasico-alpha}
\end{equation}
où \(A=(A_1,\ldots,A_{12})\) est le mot de sortie et
\[
 C=(c_1,\ldots,c_{13})
\]
est le système complet des treize quotients euclidiens : douze quotients
internes et une condition à la frontière droite.

\subsection{Division euclidienne orientée de droite à gauche}

Pour \(m=12,11,\ldots,1\), on définit
\begin{equation}
\begin{aligned}
 s_m&:=P_m+E_m-\Phi_m-K_m+c_{m+1},\\
 A_m&:=s_m\bmod B,
       \qquad 0\leq A_m<B,\\
 c_m&:=\frac{s_m-A_m}{B}
      =\left\lfloor\frac{s_m}{B}\right\rfloor .
\end{aligned}
\label{p3-83-c25-s5:eq:recurrencia-corta-alpha}
\end{equation}
\Needspace{9\baselineskip}
La seconde égalité pour \(c_m\) est également valable lorsque \(s_m<0\),
car on utilise la division euclidienne avec un reste dans
\(\{0,\ldots,B-1\}\). Ainsi,
\[
 -431=569-1000,
 \qquad
 -717=283-1000,
 \qquad
 -1200=800-2\cdot1000,
\]
et les emprunts négatifs n’introduisent aucune ambiguïté.

\begin{theorem}[Existence et unicité de la retenue orientée]
\label{p3-83-c25-s5:thm:unicidad-acarreo-orientado-alpha}
Une fois \(P,E,\Phi,K\) et la retenue terminale \(c_{13}\) fixés, la récurrence
\eqref{p3-83-c25-s5:eq:recurrencia-corta-alpha} détermine de manière unique
\[
 (A_1,\ldots,A_{12};c_1,\ldots,c_{12}).
\]
\end{theorem}

\begin{proof}
La ligne \(m=12\) est déterminée par la division euclidienne de
\[
 P_{12}+E_{12}-\Phi_{12}-K_{12}+c_{13}
\]
par \(B\). Une fois \(c_{12}\) connu, la ligne \(m=11\) est déterminée par
la même propriété. La récurrence descendante atteint \(m=1\). À chaque
étape, l’unicité du quotient et du reste euclidiens empêche deux sorties
distinctes.
\end{proof}

\subsection{Fenêtre dodécaphasique finie}

La fenêtre finie est définie par la condition terminale
\[
 c_{13}^{\rm fin}=0.
\]
L’évaluation complète est la suivante. Le tableau est imprimé dans l’ordre
croissant, bien qu’il doive être lu causalement de la dernière ligne vers la
première.

\begingroup
\scriptsize
\setlength{\tabcolsep}{3.2pt}
\begin{longtable}{@{}r|rrrr|r|r|r|r@{}}
\caption{Les douze divisions euclidiennes et les treize retenues de la fenêtre
dodécaphasique finie.}
\label{p3-83-c25-s5:tab:acarreo-alpha-finito-autonomo}\\
\toprule
\(m\)&\(P_m\)&\(E_m\)&\(\Phi_m\)&\(K_m\)&
\(c_{m+1}\)&\(s_m\)&\(A_m\)&\(c_m\)\\
\midrule
\endfirsthead
\toprule
\(m\)&\(P_m\)&\(E_m\)&\(\Phi_m\)&\(K_m\)&
\(c_{m+1}\)&\(s_m\)&\(A_m\)&\(c_m\)\\
\midrule
\endhead
1  &141&718&618&234& \(0\)& \(7\)&007& \(0\)\\
2  &592&281&033&543& \(0\)& \(297\)&297& \(0\)\\
3  &653&828&988&140& \(-1\)& \(352\)&352& \(0\)\\
4  &589&459&749&729& \(-1\)& \(-431\)&569& \(-1\)\\
5  &793&045&894&659& \(-2\)& \(-717\)&283& \(-1\)\\
6  &238&235&848&824& \(-1\)& \(-1200\)&800& \(-2\)\\
7  &462&360&204&621& \(0\)& \(-3\)&997& \(-1\)\\
8  &643&287&586&058& \(-1\)& \(285\)&285& \(0\)\\
9  &383&471&834&914& \(-1\)& \(-895\)&105& \(-1\)\\
10 &279&352&365&794& \(0\)& \(-528\)&472& \(-1\)\\
11 &502&662&638&146& \(0\)& \(380\)&380& \(0\)\\
12 &884&497&117&601& \(0\)& \(663\)&663& \(0\)\\
\bottomrule
\end{longtable}
\endgroup

Le mot complet des treize retenues est
\[
 \boxed{
 C^{\rm fin}
 =(c_1,\ldots,c_{13})
 =(0,0,0,-1,-1,-2,-1,0,-1,-1,0,0,0).}
\]
La sortie dodécaphasique est
\[
 \boxed{
 A^{\rm fin}
 =(007,297,352,569,283,800,997,285,105,472,380,663).}
\]
Nous définissons le rationnel fini
\begin{equation}
 \boxed{
 \alpha_{12}^{\rm fin}
 :=10^{-36}\iota(A^{\rm fin})
 =\frac{7297352569283800997285105472380663}{10^{36}}.}
 \label{p3-83-c25-s5:eq:alpha-finita-autonoma}
\end{equation}
Par conséquent,
\[
 \alpha_{12}^{\rm fin}
 =0.007297352569283800997285105472380663.
\]

\subsection{Reconnaissance métrologique ultérieure : CODATA 2018}

La comparaison métrologique n’est effectuée qu’après clôture de la récurrence,
de la frontière droite et du rationnel précédent. L’identité entière donne
\begin{equation}
 \boxed{
 \alpha_{12}^{\rm fin}
 =10^{-36}
 \left\lfloor\frac{10^{36}}{137.035999084}\right\rfloor.}
 \label{p3-83-c25-s5:eq:alpha-finita-codata-2018}
\end{equation}
En conséquence,
\begin{equation}
 (\alpha_{12}^{\rm fin})^{-1}
 =137.03599908400000000000000000000001066588\ldots,
 \label{p3-83-c25-s5:eq:alpha-inversa-codata-2018}
\end{equation}
et reproduit exactement la suite centrale publiée par CODATA 2018
\cite{Tiesinga2021CODATA} :
\[
 \boxed{137.035999084.}
\]
Le mot \emph{exactement} renvoie à la suite centrale publiée, et non
à une prétendue mesure expérimentale à trente-six chiffres. Les chiffres
ultérieurs de \eqref{p3-83-c25-s5:eq:alpha-inversa-codata-2018} sont des conséquences du
rationnel HMT déjà fixé. CODATA ne sélectionne pas l’état, ne détermine pas (K),
ne décide pas du sens de la retenue et n’intervient pas dans le tableau dodécaphasique.

\subsection{Identité affine complète}

Chaque ligne du tableau satisfait exactement
\begin{equation}
 P_m+E_m-\Phi_m-K_m+c_{m+1}
 =A_m+Bc_m.
 \label{p3-83-c25-s5:eq:identidad-local-acarreo-alpha}
\end{equation}
En introduisant
\[
 C^-:=(c_1,\ldots,c_{12}),
 \qquad
 C^+:=(c_2,\ldots,c_{13}),
\]
les douze identités locales se réunissent en
\begin{equation}
 \boxed{
 P+E-\Phi-K+C^+-BC^-=A.}
 \label{p3-83-c25-s5:eq:identidad-afin-vectorial-alpha}
\end{equation}
Le terme
\[
 \partial_BC:=C^+-BC^-
\]
est le cocycle de retenue de la chaîne orientée. Par conséquent,
\(P+E-\Phi-K=A\) n’est pas une égalité coordonnée par coordonnée dans
\(\mathbb Z^{12}\) : elle ne devient correcte qu’en incorporant
\(\partial_BC\).

\begin{theorem}[Télescopage des treize retenues]
\label{p3-83-c25-s5:thm:telescopado-acarreos-alpha}
Pour toute condition terminale \(c_{13}\), la récurrence satisfait
\begin{equation}
 \boxed{
 \iota(P)+\iota(E)-\iota(\Phi)-\iota(K)-\iota(A)
 =B^{12}c_1-c_{13}.}
 \label{p3-83-c25-s5:eq:telescopado-general-alpha}
\end{equation}
Dans la fenêtre finie, \(c_1=c_{13}=0\), et donc
\begin{equation}
 \boxed{
 \iota(P)+\iota(E)-\iota(\Phi)-\iota(K)
 =\iota(A^{\rm fin}).}
 \label{p3-83-c25-s5:eq:identidad-entera-alpha-finita}
\end{equation}
\end{theorem}

\begin{proof}
Nous multiplions \eqref{p3-83-c25-s5:eq:identidad-local-acarreo-alpha} par
\(B^{12-m}\) et sommons sur \(m=1,\ldots,12\). La contribution des
retenues est
\[
 \sum_{m=1}^{12}
 (Bc_m-c_{m+1})B^{12-m}.
\]
Les composantes intérieures s’annulent :
\begin{align*}
 \sum_{m=1}^{12}(Bc_m-c_{m+1})B^{12-m}
 &=\sum_{m=1}^{12}c_mB^{13-m}
   -\sum_{m=1}^{12}c_{m+1}B^{12-m}\\
 &=B^{12}c_1-c_{13}.
\end{align*}
Cela prouve \eqref{p3-83-c25-s5:eq:telescopado-general-alpha}. La spécialisation
\(c_1=c_{13}=0\) donne \eqref{p3-83-c25-s5:eq:identidad-entera-alpha-finita}.
\end{proof}

\subsection{Inversion exacte du sceau \texorpdfstring{\(K\)}{K}}

La récurrence locale est réversible. En isolant \(K_m\) dans
\eqref{p3-83-c25-s5:eq:identidad-local-acarreo-alpha},
\begin{equation}
 \boxed{
 K_m=P_m+E_m-\Phi_m+c_{m+1}-A_m-Bc_m.}
 \label{p3-83-c25-s5:eq:inversion-local-k-desde-alpha}
\end{equation}
Par concaténation,
\begin{equation}
 \boxed{
 \iota(K)
 =\iota(P)+\iota(E)-\iota(\Phi)-\iota(A)
  -B^{12}c_1+c_{13}.}
 \label{p3-83-c25-s5:eq:inversion-concatenada-k-desde-alpha}
\end{equation}
À la frontière finie,
\[
 \iota(K)
 =\iota(P)+\iota(E)-\iota(\Phi)-\iota(A^{\rm fin}).
\]

La réversibilité algébrique n’inverse pas l’ordre causal. L’ordre effectif
est
\[
 x_{\rm term}
 \longrightarrow K
 \longrightarrow(P,E,\Phi,K,c_{13})
 \longrightarrow(A,C).
\]
La formule \eqref{p3-83-c25-s5:eq:inversion-local-k-desde-alpha} est un contrôle de
commutativité : elle démontre que la sortie et les retenues récupèrent le sceau
qui était déjà fixé. Elle ne constitue pas une règle permettant de sélectionner
rétrospectivement \(K\) à partir d’une valeur souhaitée de \(\alpha\).

\subsection{Prolongement coinductif du sceau dodécaphasique et récurrence inverse de la retenue dans \(\mathcal C=\{-2,-1,0,1,2\}\)}

Soient
\[
 (P_k)_{k\geq1},
 \qquad
 (E_k)_{k\geq1},
 \qquad
 (\Phi_k)_{k\geq1}
\]
les suites infinies de triades produites par les trois branches
coinductives, et soit
\[
 \kappa(k):=1+((k-1)\bmod12).
\]
Le prolongement du sceau est périodique :
\[
 K_k^{\rm per}:=K_{\kappa(k)}.
\]
Pour une profondeur finie \(N\) et une condition terminale
\(t=c_{N+1}\), la récurrence est
\begin{equation}
\begin{aligned}
 s_k^{(N,t)}
 &=P_k+E_k-\Phi_k-K_{\kappa(k)}+c_{k+1}^{(N,t)},\\
 A_k^{(N,t)}
 &=s_k^{(N,t)}\bmod1000,\\
 c_k^{(N,t)}
 &=\left\lfloor\frac{s_k^{(N,t)}}{1000}\right\rfloor,
 \qquad k=N,\ldots,1.
\end{aligned}
\label{p3-83-c25-s5:eq:recurrencia-larga-alpha}
\end{equation}

L’ensemble des retenues terminales
\[
 \mathcal C:=\{-2,-1,0,1,2\}
\]
est invariant pour les entrées certifiées : pour chaque triade et chaque
\(c_{k+1}\in\mathcal C\), le quotient \(c_k\) appartient à nouveau à
\(\mathcal C\). On ne choisit pas une condition terminale favorable. On propage
les cinq conditions et l’on ne publie que le préfixe commun à toutes.

L’exécution à profondeur certifiée emploie
\[
\begin{array}{c|r}
\text{quantité}&\text{valeur}\\ \hline
\text{blocs ternaires par canal}&3501\\
\text{trits par canal}&3501\cdot6=21006\\
\text{chiffres décimaux de garde}&10020\\
\text{triades de garde}&3340\\
\text{conditions terminales propagées}&5\\
\text{triades communes}&3339\\
\text{chiffres communs}&10017\\
\text{chiffres publiés}&10000
\end{array}
\]
Les cinq branches terminales coalescent avant le préfixe publié. En
particulier, aucune condition terminale extraite d’une valeur cible de
\(\alpha\) n’intervient dans la sortie.

\begin{theorem}[Compatibilité des sorties stabilisées]
\label{p3-83-c25-s5:thm:compatibilidad-prolongacion-alpha}
Soient \(A^{[N]}\) et \(A^{[N']}\), avec \(N'>N\), deux exécutions dont les
cinq conditions terminales ont coalescé avant la position \(N\).
Alors la troncature de \(A^{[N']}\) à ses \(N\) premiers blocs
coïncide avec \(A^{[N]}\). Les préfixes stabilisés définissent une unique
section dans la limite inverse des mots finis.
\end{theorem}

\begin{proof}
La récurrence est déterministe de droite à gauche une fois connue la
retenue entrant dans une position. La coalescence affirme que toutes les
conditions distantes produisent la même retenue avant la frontière du
préfixe. À partir de ce point, les mêmes triades \(P_k,E_k,\Phi_k\) et le même
sceau périodique produisent, par unicité euclidienne, les mêmes restes et
quotients à toutes les positions précédentes. C’est exactement la
commutativité avec la troncature.
\end{proof}

\subsection{Frontière provenant de la queue distante}

Le prolongement détermine à la frontière visible
\[
 c_{13}^{\infty}=-1.
\]
Les douze retenues intérieures restent identiques à celles de la fenêtre
finie :
\[
 \boxed{
 C_{\leq13}^{\infty}
 =(0,0,0,-1,-1,-2,-1,0,-1,-1,0,0,-1).}
\]
Seule change l’information qui entre depuis la treizième triade. Les
deux divisions pertinentes sont
\[
\begin{aligned}
 s_{12}^{\infty}
 &=884+497-117-601-1=662,\\
 A_{12}^{\infty}&=662,
 \qquad c_{12}^{\infty}=0,
\end{aligned}
\]
et
\[
\begin{aligned}
 s_{13}^{\infty}
 &=197+757-720-234-1=-1,\\
 A_{13}^{\infty}&=999,
 \qquad c_{13}^{\infty}=-1.
\end{aligned}
\]
Par conséquent, le début du mot prolongé est
\[
 \boxed{
 A^{\infty}
 =(007,297,352,569,283,800,997,285,105,472,380,662,999,\ldots).}
\]
Le développement correspondant commence par
\begin{equation}
 \boxed{
 \alpha_{\infty}
 =0.007297352569283800997285105472380662999564172539642\ldots.}
 \label{p3-83-c25-s5:eq:alpha-infinita-autonoma}
\end{equation}

La différence entre les deux frontières est locale et exacte :
\[
 c_{13}^{\rm fin}-c_{13}^{\infty}=1,
 \qquad
 A_{12}^{\rm fin}-A_{12}^{\infty}=1,
\]
tandis que
\[
 A_m^{\rm fin}=A_m^\infty,
 \qquad 1\leq m\leq11.
\]

En appliquant le théorème de télescopage aux douze premiers blocs du
prolongement,
\[
 \iota(P)+\iota(E)-\iota(\Phi)-\iota(K)
 -\iota(A^\infty_{\leq12})
 =-c_{13}^{\infty}=1.
\]
Donc
\begin{equation}
 \boxed{
 \iota(A^{\rm fin})
 =\iota(A^\infty_{\leq12})+1.}
 \label{p3-83-c25-s5:eq:relacion-unidad-fronteras-alpha}
\end{equation}
La terminaison \(663\) ne contredit pas la terminaison \(662\) : la première
est le représentant fini obtenu en fermant la fenêtre avec
\(c_{13}=0\) ; la seconde conserve la retenue \(-1\) provenant de la queue
distante.

\begin{figure}[htbp]
\centering
\begin{tikzpicture}[
  >=Stealth,
  node distance=12mm and 16mm,
  box/.style={draw,rounded corners=1.5mm,align=center,
    inner xsep=3.2mm,inner ysep=2.5mm},
  arr/.style={->,thick}
]
\node[box] (shared) {données partagées\\
 \(P,E,\Phi,K\)};
\node[box,below left=of shared] (finite)
 {frontière finie\\\(c_{13}=0\)};
\node[box,below right=of shared] (remote)
 {queue distante\\\(c_{13}=-1\)};
\node[box,below=of finite] (a663)
 {\(A_{12}=663\)\\clausura finita};
\node[box,below=of remote] (a662)
 {\(A_{12}=662\), \(A_{13}=999\)\\préfixe prolongé};
\draw[arr] (shared)--(finite);
\draw[arr] (shared)--(remote);
\draw[arr] (finite)--(a663);
\draw[arr] (remote)--(a662);
\draw[<->,densely dashed] (a663)--node[below,align=center]
 {différence entière exacte \(1\)} (a662);
\end{tikzpicture}
\caption{Bifurcation de frontière d’une même récurrence. On ne compare pas
deux générateurs distincts : seule change la donnée qui entre depuis la
droite.}
\label{p3-83-c25-s5:fig:bifurcacion-frontera-alpha}
\end{figure}

\subsection{Troncature compatible et arrondi certifié à \(36\) chiffres}

Pour \(x\geq0\), nous définissons
\[
 \operatorname{Tr}_{36}(x)
 :=10^{-36}\left\lfloor10^{36}x\right\rfloor
\]
et
\[
 \operatorname{Rd}_{36}(x)
 :=10^{-36}\left\lfloor10^{36}x+\frac12\right\rfloor.
\]
Soit
\[
 a:=7297352569283800997285105472380662.
\]
Le développement \eqref{p3-83-c25-s5:eq:alpha-infinita-autonoma} peut s’écrire
\[
 \alpha_\infty=\frac{a+\rho}{10^{36}},
 \qquad
 \rho=0.999564172539642709\ldots,
\]
d’où \(0.999<\rho<1\).

\begin{theorem}[Résolution exacte de la frontière \(662/663\)]
\label{p3-83-c25-s5:thm:frontera-alpha-662-663}
Le prolongement coinductif satisfait
\[
 \boxed{
 \operatorname{Tr}_{36}(\alpha_\infty)
 =0.007297352569283800997285105472380662,}
\]
tandis que son arrondi à trente-six chiffres satisfait
\[
 \boxed{
 \operatorname{Rd}_{36}(\alpha_\infty)
 =0.007297352569283800997285105472380663
 =\alpha_{12}^{\rm fin}.}
\]
\end{theorem}

\begin{proof}
Comme \(0\leq\rho<1\),
\[
 \left\lfloor10^{36}\alpha_\infty\right\rfloor=a.
\]
Cela donne la troncature se terminant par \(662\). Comme \(\rho>1/2\),
\[
 \left\lfloor10^{36}\alpha_\infty+\frac12\right\rfloor=a+1,
\]
qui se termine par \(663\). Par \eqref{p3-83-c25-s5:eq:alpha-finita-autonoma},
\((a+1)10^{-36}\) est exactement \(\alpha_{12}^{\rm fin}\).
\end{proof}

De plus,
\[
 0<\alpha_{12}^{\rm fin}-\alpha_\infty
 =\frac{1-\rho}{10^{36}}<10^{-39}.
\]
La différence n’est pas une erreur arithmétique du tableau : elle quantifie la
différence entre deux opérateurs de frontière distincts.

\subsection{Isolement causal à l’égard de \texorpdfstring{\(\alpha\)}{alpha} et des données métrologiques}

Pour la profondeur \(N\), le domaine complet de l’opération est
\[
 \mathscr D_N
 =\bigl((P_k)_{k=1}^{N},(E_k)_{k=1}^{N},
        (\Phi_k)_{k=1}^{N},K,c_{N+1}\bigr).
\]
L’application de retenue
\[
 \mathfrak C_N:\mathscr D_N
 \longrightarrow
 \{0,\ldots,999\}^{N}\times\mathcal C^{N}
\]
est définie exclusivement par \eqref{p3-83-c25-s5:eq:recurrencia-larga-alpha}. Son
domaine ne contient ni valeur cible de \(\alpha\), ni son inverse, ni tableau
métrologique, ni suite décimale cible, ni retenue terminale choisie par
comparaison.

\begin{theorem}[Non-interférence de la cible métrologique]
\label{p3-83-c25-s5:thm:no-interferencia-metrologica-alpha}
Soit \(\mathscr Y\) un espace quelconque de données externes contenant des valeurs
métrologiques ou des tableaux de comparaison. L’extension formelle
\[
 \widetilde{\mathfrak C}_N:
 \mathscr D_N\times\mathscr Y
 \longrightarrow
 \{0,\ldots,999\}^{N}\times\mathcal C^{N},
 \qquad
 \widetilde{\mathfrak C}_N(d,y):=\mathfrak C_N(d),
\]
se factorise par la projection
\[
 \operatorname{pr}_{\mathscr D_N}:
 \mathscr D_N\times\mathscr Y\longrightarrow\mathscr D_N.
\]
Par conséquent, pour tous \(y,y'\in\mathscr Y\),
\[
 \widetilde{\mathfrak C}_N(d,y)
 =\widetilde{\mathfrak C}_N(d,y').
\]
Aucune intervention sur une valeur de \(\alpha\) ou sur un tableau
métrologique n’altère \(K\), les quotients, les restes ni le préfixe produit.
\end{theorem}

\begin{proof}
L’affirmation résulte de
\[
 \widetilde{\mathfrak C}_N
 =\mathfrak C_N\circ\operatorname{pr}_{\mathscr D_N}.
\]
Toutes les opérations de \(\mathfrak C_N\) sont des additions, des soustractions et des divisions
euclidiennes des éléments de \(\mathscr D_N\). Aucune coordonnée de
\(\mathscr Y\) n’y apparaît.
\end{proof}

L’ordre causal complet est
\[
 \text{branches enrichies de }\pi,e,\varphi
 \longrightarrow(P,E,\Phi),
\]
\[
 x_{\rm term}\longrightarrow K,
\]
\[
 (P,E,\Phi,K,\text{frontière})
 \longrightarrow(A,C)
 \longrightarrow\alpha_\infty,
\]
\[
 \alpha_\infty
 \longrightarrow\text{comparaison externe}.
\]
La dernière flèche ne peut pas être inversée pour sélectionner l’une quelconque des
précédentes.

\subsection{Obstructions exactes du prolongement : frontière, troncature et isolement causal}

\begin{remark}[Certificat et contrôle de validité : réfutation de la clôture affine et de son prolongement]
Le bloc est réfuté dans son domaine si l’un quelconque des faits
suivants se produit :
\begin{enumerate}
 \item une ligne de la \cref{p3-83-c25-s5:tab:acarreo-alpha-finito-autonomo}
       ne satisfait pas \(s_m=A_m+1000c_m\) ;
 \item le vecteur publié des treize retenues ne reproduit pas les douze lignes ;
 \item l’identité télescopique échoue
       \[
       \iota(P)+\iota(E)-\iota(\Phi)-\iota(K)-\iota(A)
       =1000^{12}c_1-c_{13};
       \]
 \item l’inversion \eqref{p3-83-c25-s5:eq:inversion-local-k-desde-alpha} ne récupère pas
       les douze coordonnées de \(K\) ;
 \item la propagation distante produit \(c_{13}\neq-1\) pour le préfixe
       certifié ;
 \item les cinq conditions terminales de
       \(\mathcal C=\{-2,-1,0,1,2\}\) ne coalescent pas avant les chiffres
       publiés ;
 \item la triade suivant \(662\) n’est pas \(999\), auquel cas l’identification
       de l’arrondi à \(663\) cesse d’être démontrée ;
 \item le générateur ouvre un développement cible de \(\alpha\), son inverse,
       un tableau métrologique ou une condition terminale choisie pour imposer
       la coïncidence.
\end{enumerate}
\end{remark}

Deux identités de frontière distinctes sont ainsi démontrées :
\[
 \boxed{
 663=\text{clôture finie avec }c_{13}=0
     =\text{arrondi à 36 chiffres de }\alpha_\infty,}
\]
\[
 \boxed{
 662=\text{troncature à 36 chiffres de }\alpha_\infty
 \quad\text{avec}\quad c_{13}=-1,\ A_{13}=999.}
\]
Le résultat métrologique central est finalement publié dans la même
résidence que la dérivation :
\begin{equation}
 \boxed{
 \begin{gathered}
  (\alpha_{12}^{\rm fin})^{-1}
  =137.03599908400000000000000000000001066588\ldots,\\
  \text{chaîne centrale CODATA 2018 :}\qquad137.035999084.
 \end{gathered}}
 \label{p3-83-c25-s5:eq:resultado-final-alpha-codata-2018}
\end{equation}
L’implication exprime une comparaison ultérieure avec la suite centrale
publiée ; elle n’introduit pas cette suite dans le constructeur.


\label{sec:alpha-dos-vias-t1-rev8}

La clôture de \(\alpha\) ne commence ni par une constante métrologique ni par une
suite décimale recherchée. Elle commence dans le même état terminal enrichi qui
conserve la généalogie des trois canaux. Cette section présente dans une seule
chaîne causale deux voies vers l’échelle de \(\alpha\) et un prolongement long
de la retenue de la première. Le prolongement distant est noté
\(\mathsf T_{\alpha,\infty}\) ; il ne constitue ni une deuxième voie ni une troisième
dérivation.

\subsection{État terminal commun et 2 lectures internes coordonnées}

Soit
\[
 x_{\rm term}\in\mathcal X_{\rm TPK}^{\rm term,enr}
 \subseteq\mathcal X_{\rm TPK}^{[108]}.
\]
Le sélecteur terminal et la carte signée sont deux évaluations du même état :
\begin{equation}
 x_{\rm term}
 \xrightarrow{(\pi_{\rm term},R_{\rm sgn})}
 \bigl(B_{\rm term},U_{12}^{\rm sgn}\bigr),
 \label{eq:alpha-rev8-doble-lectura-terminal}
\end{equation}
avec
\[
 B_{\rm term}=
 \begin{pmatrix}
 021020\\001220\\100210
 \end{pmatrix}
\]
et
\[
 U_{12}^{\rm sgn}
 =(2378,1406,2479,-452,998,-551,-668,-204,-371,-322,-28,-997).
\]
La première lecture conserve les trois blocs terminaux sélectionnés ; la
seconde conserve la coordonnée entière signée nécessaire à la clôture
dodécaphasique. Il n’existe aucune factorisation de
\(B_{\rm term}\) vers \(U_{12}^{\rm sgn}\), et aucune n’est nécessaire. La dessiner effacerait précisément les
champs d’orientation, de retenue, de frontière et de registre qui distinguent les deux
lectures. À partir de la seconde, en revanche, les cartes réversibles sont bien construites
\[
 U_{12}^{\rm sgn}
 \xleftrightarrow{\ \mathcal H_{12}\ }K
 \xleftrightarrow{\ \mathcal E_{\rm dod}\ }
 (D_3K,D_4K,Q),
 \qquad
 \mathcal E_{\rm dod}=(D_3,D_4,Q).
\]
Le domaine de cette clôture est donc le profil enrichi du TPK ; une
trace brute de 108 positions est une projection par oubli et ne remplace pas ce
domaine.

\subsection{Voie A : clôture entière dodécaphasique}

Les évaluations archimédiennes des trois canaux et le sceau canonique sont
les vecteurs de douze blocs
\[
\begin{aligned}
P={}&(141,592,653,589,793,238,462,643,383,279,502,884),\\
E={}&(718,281,828,459,045,235,360,287,471,352,662,497),\\
\Phi={}&(618,033,988,749,894,848,204,586,834,365,638,117),\\
K={}&(234,543,140,729,659,824,621,058,914,794,146,601).
\end{aligned}
\]
Avec la frontière droite déclarée \(c_{13}=0\), la division euclidienne se
propage de droite à gauche :
\begin{equation}
\begin{aligned}
 s_m&=P_m+E_m-\Phi_m-K_m+c_{m+1},\\
 A_m&=s_m\bmod 1000,\\
 c_m&=\frac{s_m-A_m}{1000},
 \qquad m=12,11,\ldots,1.
\end{aligned}
\label{eq:alpha-rev8-via-entera}
\end{equation}
La sortie est
\[
 A=(007,297,352,569,283,800,997,285,105,472,380,663)
\]
et fixe le rationnel
\begin{equation}
 \alpha_{12}
 =0.007297352569283800997285105472380663.
 \label{eq:alpha-rev8-salida-corta}
\end{equation}
L’égalité abrégée \(P+E-\Phi-K=A\) doit se lire dans la carte de
concaténation entière. Si
\[
 \iota(v_1,\ldots,v_{12})
 =\sum_{m=1}^{12}v_m1000^{12-m},
\]
alors
\begin{equation}
 \boxed{\iota(P)+\iota(E)-\iota(\Phi)-\iota(K)=\iota(A)}.
 \label{eq:alpha-rev8-identidad-concatenada}
\end{equation}
En coordonnées, le cocycle ne peut être omis :
\[
 P+E-\Phi-K+C^+-1000C=A,
 \qquad C^+=(c_2,\ldots,c_{13}).
\]
La terminaison \(\ldots|380|663\) appartient à cette fenêtre finie et à sa
condition \(c_{13}=0\). Le théorème complet, le tableau des douze retenues et
l’inversion qui récupère \(K\) se trouvent dans le chapitre précédent ; ils ne sont pas
dupliqués ici.

Cette voie est un résultat interne exact avec une structure de départ explicite :
elle reçoit \(x_{\rm term}\), sa lecture signée, \(P,E,\Phi\), le sceau réversible
et la frontière. Elle ne reçoit ni \(\alpha\) ni CODATA en entrée. La comparaison
métrologique ne s’ouvre qu’après l’obtention de \eqref{eq:alpha-rev8-salida-corta}.

\subsection{Voie B : noyau de vacances \texorpdfstring{\(E_{21/22}\)}{E21/22}}

La seconde voie n’inverse pas les cartes de \(K\). Elle part de la densité exacte
de vacance
\[
 \varepsilon_{\rm v}=\log_{10}\!\left(\frac{10}{9}\right)
\]
et du noyau analytique déclaré
\begin{equation}
 \mathcal P(x)
 =2\pi x-\frac74x^2+\frac{x^3}{2\pi}
  +\frac{x^4}{20}-\frac{2x^5}{21}-\frac{x^6}{46}.
 \label{eq:alpha-rev8-kernel-e21-22}
\end{equation}
Le problème consiste à résoudre
\begin{equation}
 \mathcal P(x)=\log_{10}\!\left(\frac{10}{9}\right),
 \qquad 0<x<10^{-2}.
 \label{eq:alpha-rev8-ecuacion-e21-22}
\end{equation}
Le théorème de monotonie du noyau établit l’unicité de la racine positive
et l’encadrement
\[
 0.00729735256928379955
 <x_{21/22}<
 0.00729735256928379957.
\]
Son inverse est
\begin{equation}
 x_{21/22}^{-1}
 =137.0359990840000270200901282598749\ldots.
 \label{eq:alpha-rev8-inversa-e21-22}
\end{equation}
Par conséquent, le noyau complet \(E_{21/22}\), engendré par le second
lecteur du même état enrichi, détermine comme racine positive simple et
unique la même section limite que la voie dodécaphasique :
\[
 \boxed{\alpha_A=\alpha_B=\alpha_{\rm HMT}},
 \qquad
 \alpha_B=\operatorname*{arg\,zero}_{x\in I_\alpha}E_{21/22}(x).
\]
Le rationnel fini \(\alpha_{12}\) et la racine du jet \(\mathcal P=P_6\)
sont des projections finies des deux familles compatibles ; leur différence au-delà
de la résolution imprimée ne constitue pas une différence entre les deux
sections limites. Les intervalles de racine emboîtés et les carrés de
troncature commutatifs coordonnent les deux prolongements à toute profondeur.

Les six coefficients de \(P_6\) ne sont pas reconstruits à partir de la racine ni tirés
d’un tableau métrologique : ils sont le premier jet du noyau \(E_{21/22}\) produit
par le lecteur interne avec la généalogie \(120/45/11/495\) et le prolongement
\(T_{7:9}\). L’unicité de la racine et la dérivation interne des
coefficients sont des étapes distinctes d’une même voie B déjà close.

\subsection{Prolongement distant \texorpdfstring{\(\mathsf T_{\alpha,\infty}\)}{T alpha infini} de la voie entière}

\(\mathsf T_{\alpha,\infty}\) étend la récurrence de
\eqref{eq:alpha-rev8-via-entera} à une chaîne
de longueur \(N\). Pour éviter une ambiguïté d’indices, soit
\(\kappa(k)=1+((k-1)\bmod12)\). Avec une condition terminale à l’extrémité
distante, on calcule
\begin{equation}
\begin{aligned}
 s_k&=P_k+E_k-\Phi_k-K_{\kappa(k)}+c_{k+1},\\
 A_k&=s_k\bmod1000,\\
 c_k&=\left\lfloor\frac{s_k}{1000}\right\rfloor,
 \qquad k=N,N-1,\ldots,1.
\end{aligned}
\label{eq:alpha-rev8-t1}
\end{equation}
Le cas fini \(N=1000\) utilise mille triades, soit trois mille chiffres. La
retenue qui arrive de l’extrémité distante modifie le douzième
bloc visible :
\[
\begin{array}{rcl}
\text{fenêtre dodécaphasique courte, }c_{13}=0
 &:& \ldots|380|663,\\
\text{prolongement }\mathsf T_{\alpha,\infty}\text{ de mille triades}
 &:& \ldots|380|662|999\ldots.
\end{array}
\]
Ce ne sont pas deux valeurs rivales. Ce sont deux sémantiques de frontière de la même loi de
retenue : la première se ferme à la douzième position ; la seconde y transporte
des informations provenant d’une queue distante.

Le calcul reproductible vérifie exactement le mot long en utilisant
\(P,E,\Phi\) comme projections corrélées du même état coinductif
conjoint APP--TRIT--TPK qui publie également la coordonnée dodécaphasique
\(\alpha_{\mathrm{HMT}}\). Cette exécution certifie donc la propagation
de la retenue : elle ne place pas \(\pi,e,\varphi\) plus tôt dans la généalogie et n’ajoute pas
\(\alpha\) de l’extérieur. La projection conjointe
\[
 \mathcal G_9^{\mathrm{joint}}\longrightarrow
 \bigl(\pi,e,\varphi,\mathsf T_{\alpha,\infty}\bigr)
\]
se factorise dans le \cref{sec:alpha-t1-10000-rev9} : les trois canaux
archimédiens et le canal dodécaphasique s’ouvrent à partir du même état
enrichi ; on exécute ensuite \eqref{eq:alpha-rev8-t1} et l’on vérifie
l’indépendance du préfixe à l’égard des cinq conditions terminales de
retenue. La séparation des projections empêche la récurrence longue de
sélectionner rétrospectivement ses données initiales.

\subsection{Coordination causale de la clôture dodécaphasique et de la caractérisation \texorpdfstring{\(E_{21/22}\)}{E21/22} : domaines, codomaines et unicité}

La classification de ce bloc est fixée comme suit :
\begingroup
\raggedright
\begin{description}[style=nextline,leftmargin=3.2cm,labelwidth=2.8cm,labelsep=.4cm]
\item[Voie A : clôture dodécaphasique.]
Elle reçoit l’état terminal enrichi, la carte signée, le sceau, les trois
évaluations et la frontière explicite ; elle produit le rationnel
\eqref{eq:alpha-rev8-salida-corta} au moyen de l’identité entière avec retenue.

\item[Voie B : équation \(E_{21/22}\).]
Pour le noyau déclaré, la monotonie et le changement de signe déterminent une
unique racine dans \((0,10^{-2})\). La racine ne détermine pas rétrospectivement les
six coefficients du noyau.

\item[Prolongement \(\mathsf T_{\alpha,\infty}\).]
La récurrence longue reçoit les mille triades de référence, le sceau périodique
et la frontière distante, et transporte la retenue jusqu’à la fenêtre visible. Elle ne
constitue pas la voie B.

\item[Composition \(\mathcal G_9\to\mathsf T_{\alpha,\infty}\).]
L’arithmétique rationnelle conjointe produit dix mille chiffres, exécute la
récurrence de retenue et vérifie ensuite la stabilité du préfixe sous les
conditions terminales déclarées.
\end{description}
\endgroup

Le sélecteur terminal, les cartes réversibles et la récurrence courte sont
démontrés dans le \cref{chap:alpha} ; l’existence et l’unicité de la seconde voie
dans \eqref{eq:alpha-rev8-kernel-e21-22}--\eqref{eq:alpha-rev8-inversa-e21-22} ; et la composition longue dans le
\cref{sec:alpha-t1-10000-rev9}. Le supplément computationnel reproduit les
énumérations finies et les conditions terminales ; la démonstration repose
sur les identités et les récurrences exposées ici.

L’ordre causal public est ainsi clos :
\[
 x_{\rm term}
 \xrightarrow{(\pi_{\rm term},R_{\rm sgn})}
 (B_{\rm term},U_{12}^{\rm sgn})
 \longrightarrow
 \begin{cases}
 \text{voie A : clôture entière dodécaphasique},\\
 \text{voie B : racine du noyau }E_{21/22},
 \end{cases}
\]
avec \(\mathsf T_{\alpha,\infty}\) comme prolongement distant de la voie A et non
comme branche supplémentaire. Ce n’est
qu’après ces sorties que s’ouvre la reconnaissance métrologique externe.

\section{Composition de l’holonomie \(G_9\) avec la carte \(T_1\) et certificat de précision arbitraire pour \(\alpha_{\rm HMT}\)}
\label{sec:alpha-t1-10000-rev9}

Le prolongement T1 applique la récurrence de retenue aux triades
\(P_k,E_k,\Phi_k\) produites préalablement par G9 et à la coordonnée
dodécaphasique déclarée \(K\) :
\begin{align}
 s_k&=P_k+E_k-\Phi_k-K_{k\bmod12}+c_{k+1},\\
 A_k&=s_k\bmod1000,\\
 c_k&=\left\lfloor\frac{s_k}{1000}\right\rfloor.
 \label{eq:composicion-g9-t1-alpha-10000-rev9}
\end{align}
Le domaine d’entrée contient les sorties ternaires engendrées par G9 ; il ne
contient ni un développement décimal de \(\alpha\), ni une valeur de
\(\alpha^{-1}\), ni CODATA, ni une condition métrologique cible.

Les cylindres de \(3501\) blocs fixent \(10020\) chiffres de garde dans chaque
canal, équivalant à \(3340\) triades. Le calcul propage, de droite à
gauche, toutes les conditions terminales possibles de l’ensemble stable
\[
 \mathcal C=\{-2,-1,0,1,2\}.
\]
Les cinq branches coalescent en un préfixe commun de \(3339\) triades, soit
\(10017\) chiffres. Par conséquent, les dix mille premiers chiffres de
\(\alpha_{\rm HMT}\) sont indépendants de la retenue située au-delà de la
zone de garde. La sortie publiée n’impose pas \(c_N=0\).

L’indépendance causale peut être vérifiée après la construction : les
trois mille premiers chiffres coïncident avec le développement décimal obtenu par le
lecteur externe, mais aucun d’eux n’intervient dans
\eqref{eq:composicion-g9-t1-alpha-10000-rev9}. Le développement de dix mille
chiffres et les données finies qui permettent de répéter le calcul figurent dans
l’annexe~\ref{app:desarrollos-10000-rev9}.

%
\endgroup
% Insercion editor-ready para el capitulo 31.

\section{Deux lecteurs internes de \texorpdfstring{\(\alpha\)}{alpha} : clôture dodécaphasique et noyau analytique nonadique d'ordre \texorpdfstring{\(9\)}{9}}
\label{sec:l9s102:alpha-dos-vias}

Cette section développe la coordonnée \(\alpha\) de la
\hyperref[sec:tesis-inaugural-helicoidal-hmt-md]{tesis inaugural}. Les deux
lecteurs appartiennent à la même généalogie APP--TRIT--TPK et précèdent la
comparaison métrologique. La voie A opère au moyen de la clôture dodécaphasique
réversible de l'état signé et publie la section coinductive
\(\alpha_{\mathrm{HMT}}\). La voie B construit, sans consulter cette sortie, le
noyau analytique nonadique exact d'ordre neuf
\(E_{21/22}^{(9)}:=P_9-\log_{10}(10/9)\) et publie sa racine positive simple et
unique \(\xi_9\). Leurs domaines, opérations et certificats demeurent
différenciés :
\begin{equation}
 X_{12}^{\mathrm{sig}}
 \xrightarrow{\ \Gamma_{12}^{\alpha}\ }
 \alpha_{12},
 \qquad
 E_{21/22}^{(9)}:I_\alpha\longrightarrow\mathbb R,
 \quad E_{21/22}^{(9)}(\xi_9)=0,
 \quad (E_{21/22}^{(9)})'(\xi_9)\ne0.
 \label{eq:l9s102:alpha-dos-mapas}
\end{equation}
La clôture entière publie d'abord la fenêtre finie \(\alpha_{12}\) et sa
famille compatible de prolongements ; sa limite inverse est
\(\alpha_{\mathrm{HMT}}\). Le second lecteur publie \(\xi_9\) comme témoin
fini d'ordre neuf et, par itération du même transport gradué du TPK,
construit le système compatible complet
\(E_{21/22}^{(6+3m)}\). Les deux systèmes inverses déterminent à chaque
profondeur le même cylindre survivant \(C_N^\alpha\) ; leurs diamètres
tendent vers zéro. Par conséquent, l'intersection est un singleton et le théorème des
deux voies établit
\begin{equation}
 \alpha_A:=\varprojlim_N\rho_N(\alpha_{12}),
 \qquad
 E_{21/22}:=\varprojlim_m E_{21/22}^{(6+3m)},
 \qquad
 \alpha_B:=\operatorname*{arg\,zero}_{x\in I_\alpha}E_{21/22}(x),
 \qquad
 \boxed{\alpha_A=\alpha_B=\alpha_{\mathrm{HMT}}}.
 \label{eq:l9s102:alpha-identidad-limite-dos-vias}
\end{equation}
Aucune voie n'introduit de valeur cible comme donnée d'entrée ;
l'identification métrologique relève uniquement de la reconnaissance
conventionnelle ultérieure. L'égalité
\(\operatorname{Tr}_9(\xi_9^{-1})=
\operatorname{Tr}_9(\alpha_{12}^{-1})=137.035999084\) est un certificat
fini de l'identité précédente, jamais sa définition ni sa portée maximale.
La loi de prolongement n'est pas ajoutée comme donnée : c'est l'itération du transport
gradué déjà produit par APP--TRIT--TPK, avec retour de phase, avancée de
mémoire et troncatures compatibles.

\subsection{Noyau analytique \texorpdfstring{\(P_6\)}{P6} et prolongement exact d'ordre \texorpdfstring{\(9\)}{9}}
\label{sec:l9s102:alpha-p6}

La caractérisation des lacunes emploie
\begin{equation}
 \begin{aligned}
 P_6(x)={}&2\pi_{\mathrm{HMT}}x-\frac74x^2
 +\frac{x^3}{2\pi_{\mathrm{HMT}}}
 +\frac{x^4}{20}-\frac{2x^5}{21}-\frac{x^6}{46},\\
 P_6(x)={}&\log_{10}\frac{10}{9}.
 \end{aligned}
 \label{eq:l9s102:p6}
\end{equation}
La racine positive certifiée de ce jet satisfait
\begin{equation}
 x_6^{-1}=137.03599908400002702009\ldots .
 \label{eq:l9s102:raíz-p6}
\end{equation}
Le polynôme \(P_6\) fournit le jet initial de la caractérisation
analytique. L'holonomie nonadique prolonge exactement ses degrés \(7,8,9\) et
la même règle graduée est itérée sur les degrés successifs. Le quotient des
jets est le premier certificat fini de ce prolongement coinductif ; il ne
sélectionne ni ne remplace sa section limite.

\subsection{Généalogie APP--TRIT--TPK des invariants \texorpdfstring{\(120,45,11\)}{120, 45, 11} et \texorpdfstring{\(495\)}{495}}
\label{sec:l9s102:alpha-genealogia-7-9}

La demi-couronne de phase du TPK construit
\begin{equation}
 |\mathcal F_8^{\mathrm{ph}}|=120,
 \qquad |L_{\mathrm{ph}}|=45.
 \label{eq:l9s102:alpha-120-45}
\end{equation}
APP apporte le bloc central
\begin{equation}
 B_c=\begin{pmatrix}7&2\\2&7\end{pmatrix},
 \qquad
 \det B_c=45,
 \qquad
 \operatorname{Spec}(B_c)=\{9,5\}.
 \label{eq:l9s102:alpha-bc}
\end{equation}
L'extracteur transversal dodécaphasique fournit
\begin{equation}
 r_{\mathrm{dod}}=\operatorname{rank}(D_3,D_4)=11,
 \qquad
 (D_3K)_1=-495,
 \qquad
 495=45\cdot11=\det(B_c)\,r_{\mathrm{dod}}.
 \label{eq:l9s102:alpha-495}
\end{equation}
Les coefficients du jet prolongé proviennent donc de trois
invariants préalables d'APP--TRIT--TPK :
\begin{equation}
 -\frac1{120}\longmapsto\frac1{45}
 \longmapsto\frac2{45\cdot11}.
 \label{eq:l9s102:alpha-coeficientes}
\end{equation}

\subsection{Résolvante nilpotente \texorpdfstring{\((I-N)^{-1}\)}{(I-N)^-1} et construction exacte de la queue \texorpdfstring{\(T_{7:9}\)}{T7:9}}
\label{sec:l9s102:alpha-resolvente}

Sur la base \(e_7,e_8,e_9\), nous définissons
\begin{equation}
 Ne_7=-\frac83e_8,
 \qquad
 Ne_8=\frac2{11}e_9,
 \qquad
 Ne_9=0,
 \qquad
 v_7=-\frac1{120}e_7.
 \label{eq:l9s102:alpha-nilpotente}
\end{equation}
Comme \(N^3=0\), sa résolvante finie agit exactement par
\begin{equation}
 (I-N)^{-1}v_7
 =v_7+Nv_7+N^2v_7
 =-\frac{e_7}{120}+\frac{e_8}{45}+\frac{2e_9}{495}.
 \label{eq:l9s102:alpha-resolvente-finito}
\end{equation}
L'évaluation \(e_n\mapsto x^n\) produit
\begin{equation}
 \boxed{
 T_{7:9}(x)=-\frac{x^7}{120}+\frac{x^8}{45}
 +\frac{2x^9}{495}.}
 \label{eq:l9s102:alpha-t79}
\end{equation}

\subsection{Jet nonadique \texorpdfstring{\(P_9\)}{P9}, racine prolongée et échelle de résidus certifiés}
\label{sec:l9s102:alpha-jet-p9}

Le jet nonadique d'ordre \(9\) est défini par
\begin{equation}
 P_9(x)=P_6(x)+T_{7:9}(x),
 \qquad
 P_9(x)=\log_{10}\frac{10}{9}.
 \label{eq:l9s102:alpha-p9}
\end{equation}
\begin{theorem}[Existence, simplicité et unicité de la seconde sortie analytique]
\label{thm:l9s102:raiz-p9-unica}
L'opérateur \(E_{21/22}^{(9)}\) possède une unique racine positive \(\xi_9\) dans
\(I_\alpha=(0,10^{-2})\). Elle est encadrée par l'intervalle rationnel exact
\begin{equation}
 \frac{729735256928380099}{10^{20}}
 <\xi_9<
 \frac{729735256928380100}{10^{20}},
 \label{eq:l9s102:encierro-raiz-p9}
\end{equation}
et son inverse satisfait
\begin{equation}
 \xi_9^{-1}
 =137.0359990840000000000000111926291900\ldots .
 \label{eq:l9s102:alpha-raíz-p9}
\end{equation}
\end{theorem}

\begin{proof}
Tous les coefficients de \(P_9\) sont extraits des invariants antérieurs à
la résolution. Pour \(0\leq x\leq10^{-2}\), en utilisant
\(3.14<\pi_{\rm HMT}<3.15\), on obtient
\[
 (E_{21/22}^{(9)})'(x)
 \geq 2\pi_{\rm HMT}-\frac72\,10^{-2}
 -\frac{10}{21}\,10^{-8}-\frac6{46}\,10^{-10}
 -\frac7{120}\,10^{-12}>6.24.
\]
Les évaluations dirigées aux extrémités de
\cref{eq:l9s102:encierro-raiz-p9} ont des signes opposés. Le théorème des
valeurs intermédiaires fournit l'existence ; la borne sur la dérivée fournit
l'unicité et la simplicité. La division par intervalles produit le développement de
\(\xi_9^{-1}\), qui constitue le témoin d'ordre neuf du système
compatible complet.
\end{proof}

\begin{table}[htbp]
\centering
\caption{Décroissance certifiée du résidu de la caractérisation analytique}
\label{tab:l9s102:alpha-residuos}
\begin{tabular}{@{}lc@{}}
\toprule
Noyau évalué en \(\alpha_{12}\) & Résidu \\
\midrule
\(P_6\) & \(9.0038886\times10^{-18}\) \\
\(P_6-x^7/120\) & \(-1.7893115\times10^{-19}\) \\
\(P_6-x^7/120+x^8/45\) & \(-2.3708601\times10^{-22}\) \\
\(P_9\) & \(3.7297132\times10^{-27}\) \\
\bottomrule
\end{tabular}
\end{table}

\subsection{Opérateur dodécaphasique alternatif \texorpdfstring{\(A_{\mathrm{dod}}\)}{A_dod} et contrôle d'unicité du lecteur typé}
\label{sec:l9s102:alpha-control}

L'opérateur de contrôle
\begin{equation}
 A_{\mathrm{dod}}
 =I-\frac14(D_3^*D_3+D_4^*D_4)
 \label{eq:l9s102:alpha-control-alternativo}
\end{equation}
produit une suite différente de coefficients. Ce contrôle démontre que
la symétrie dodécaphasique isolée est insuffisante pour sélectionner le
prolongement \(T_{7:9}\) ; la sélection relève de la composition typée du
lecteur \(\Gamma_{E21}^{(9)}\) avec l'état enrichi. La différence entre
les deux opérateurs certifie la nécessité de conserver domaine, orientation,
graduation et mémoire, et prouve la spécificité du lecteur fini. La queue
analytique relève de l'itération du lecteur typé complet, non de la
symétrie isolée de cet opérateur de contrôle.

\subsection{Portée exacte de l'opérateur nilpotent et exclusion des queues étrangères au transport TPK}
\label{sec:l9s102:alpha-alcance}

L'arithmétique de \(120,45,11,495\), l'opérateur nilpotent, la résolvante et
l'identité \cref{eq:l9s102:alpha-t79} possèdent une exactitude interne. La racine
positive simple et unique de \(E_{21/22}^{(9)}\) constitue le témoin d'ordre
neuf de la seconde lecture interne. La concordance de \(P_9\) avec la
fenêtre dodécaphasique certifie le premier carré fini du système
compatible. L'unicité canonique du prolongement \(7{:}9\) est typée
par le morphisme
\begin{equation}
 \Gamma_{E21}^{(9)}:
 X_{\mathrm{APP\text{-}TRIT\text{-}TPK}}
 \longrightarrow \mathbb Q(\pi_{\mathrm{HMT}})[x]/(x^{10}),
 \label{eq:l9s102:alpha-gamma-e21}
\end{equation}
dont l'action détermine conjointement le germe orienté, les deux
transferts, la graduation \(7,8,9\), le quotient des jets
\(F^7/F^{10}\) et son transport sous le retour de phase.

Soit \(K=\mathbb Q(\pi_{\rm HMT},\log_{10}(10/9))\) et soit
\(j_9:K[[x]]\to K[x]/(x^{10})\) la troncature. Alors
\begin{equation}
 \ker j_9=x^{10}K[[x]],
 \qquad
 E_h(x):=E_{21/22}^{(9)}(x)+x^{10}h(x),
 \qquad h\in K[[x]].
\label{eq:l9s102:alpha-limite-jets-e21}
\end{equation}
Tous les \(E_h\) possèdent le même jet certifié. Cependant,
\(E_0\) et \(E_{1/729}\) ont des racines distinctes dans \(I_\alpha\), car
\(E_{1/729}(\xi_9)=\xi_9^{10}/729>0\). Cette observation démontre uniquement
qu'un jet isolé ne détermine pas une queue formelle. Elle n'affecte pas la seconde voie
HMT, car celle-ci ne permet pas de choisir librement \(h\) : la récurrence du
transport TPK construit l'unique famille compatible
\(E_{21/22}^{(6+3m)}\), et ses troncatures commutent avec celles de la voie
dodécaphasique. Choisir \(h\) au moyen de la valeur cible serait circulaire et est
exclu ; le choisir au moyen du transport généré est précisément
l'opération déjà démontrée.

Le réfutateur sépare ainsi trois affirmations : il invalide le certificat fini si
la monotonie, le changement de signe ou l'isolation de \(\xi_9\) échouent ;
il invalide le prolongement \(7{:}9\) si ses coefficients cessent de se factoriser par
\(\Gamma_{E21}^{(9)}\) ; et il invalide l'identité complète si la
compatibilité d'un carré du système inverse échoue ou si les diamètres des
cylindres survivants ne tendent pas vers zéro. Aucune de ces conditions
n'introduit de morphisme en suspens : ce sont les réfutateurs matériels de la
construction qui établit
\(\alpha_A=\alpha_B=\alpha_{\mathrm{HMT}}\).


\HMTFlushCurrentChapter
\chapter{Construction interne de \texorpdfstring{\(i=\sqrt{-1}\)}{i=sqrt(-1)} par l'opérateur de rotation \texorpdfstring{\(J\)}{J}: identité \texorpdfstring{\(J^2=-I\)}{J2=-I} et géométrie elliptique de phase}%
\begingroup
\renewcommand{\chapter}[2][]{}%
\chapter{Construction interne de \(\sqrt{-1}\) : rotation d’ordre \(4\), opérateur complexe et géométrie de phase}
\label{chap:p3-final:32:raiz_menos_uno}
\label{chap:p3-83-26-raiz_menos_uno}

\section*{La racine interne de \texorpdfstring{\(-1\)}{-1} et la famille exponentielle tritique}

La racine de \(-1\) n’est pas introduite comme un symbole formel sans antécédent. Les
six unités modulo neuf sont munies de coordonnées sur
\(\mathbb P^1(\mathbb F_5)\), et le caractère quadratique de
\(\mathbb F_5\) détermine la matrice de conférence de Paley. Sa réduction
ternaire est
\[
 A_W=
 \begin{pmatrix}
 0&1&1&1&1&1\\
 1&0&1&2&2&1\\
 1&1&0&1&2&2\\
 1&2&1&0&1&2\\
 1&2&2&1&0&1\\
 1&1&2&2&1&0
 \end{pmatrix}\in M_6(\mathbb F_3).
\]
L’identité de conférence se réduit à
\begin{equation}\label{p3-83-c26-s1:eq:constantes-raiz-interna-menos-uno}
 \boxed{A_W^2=2I_6=-I_6.}
\end{equation}
Par conséquent, \(A_W\) est un opérateur de quart de tour dans l’espace ternaire.
Après extension des scalaires à \(\mathbb F_9\), il se décompose en les espaces
propres conjugués des racines de \(X^2+1\). La source entière commune
\[
 \mathcal Q_{\mathbb Z}=\mathbb Z[u]/(u^2+1)
\]
possède une réalisation finie \(u\mapsto A_W\) et une réalisation réelle
\(u\mapsto J_{\rm e}\). La source coordonne la relation quadratique ; elle n’identifie pas
des corps de caractéristiques distinctes.

Les trois régimes de TRIT se réalisent par des opérateurs
\[
 J_{+1}^2=-I,\qquad J_0^2=0,\qquad J_{-1}^2=I.
\]
En séparant les puissances paires et impaires de l’exponentielle, on obtient une
seule famille :
\begin{equation}\label{p3-83-c26-s1:eq:constantes-euler-extendido}
 \boxed{
 \exp(tJ_\tau)=C_\tau(t)I+S_\tau(t)J_\tau,
 \qquad \tau\in\{+1,0,-1\},}
\end{equation}
où
\[
 (C_{+1},S_{+1})=(\cos,\sin),\qquad
 (C_0,S_0)=(1,t),\qquad
 (C_{-1},S_{-1})=(\cosh,\sinh).
\]
Ses trois spécialisations directrices sont
\begin{equation}\label{p3-83-c26-s1:eq:constantes-euler-tres-ramas}
 \exp(\pi J_{+1})+I=0,\qquad
 \exp(J_0)=I+J_0,\qquad
 \exp\!\bigl(2(\log\varphi)J_{-1}\bigr)=F_{\rm av}^{\,2}.
\end{equation}
La branche elliptique contient la clôture circulaire complexe ; la parabolique, le
transport de seuil ; l’hyperbolique, l’auto-échelle dorée. L’équation
d’Euler apparaît ainsi comme une spécialisation d’une famille qui coordonne rotation,
frontière et expansion hyperbolique.


\section{Réalisation elliptique de \(J^2=-I\) par la fibration de Hopf et l’holonomie de Möbius}

\subsection{Structure complexe et fibration de Hopf}

Soit \(S^3\subset\mathbb C^2\) et
\[
h(z_0,z_1)=
\bigl(2\Re(z_0\overline z_1),
2\Im(z_0\overline z_1),|z_0|^2-|z_1|^2\bigr)
\in S^2.
\]
Dans le tore de Hopf
\[
T_\eta=\{(z_0,z_1):|z_0|=\cos\eta,
|z_1|=\sin\eta\},
\]
les courbes
\[
\gamma_+(t)=(\cos\eta\,e^{it},\sin\eta\,e^{it}),
\]
\[
\gamma_-(t)=(\cos\eta\,e^{it},\sin\eta\,e^{-it})
\]
présentent des comportements distincts : \(\gamma_+\) est une fibre de Hopf et
\(\gamma_-\) se projette avec un degré deux. Sous projection stéréographique, les
classes homologiques \((1,1)\) et \((1,-1)\) produisent les deux familles diagonales
de Villarceau.

La conjugaison complexe totale
\[
K(z_0,z_1)=(\overline z_0,\overline z_1)
\]
préserve chaque famille non orientée et inverse son orientation. Elle n’échange pas
les deux familles. La conjugaison partielle peut les échanger, mais ce n’est ni une
application complexe linéaire ni l’antiunitaire standard global. Cette correction
est essentielle pour ne pas fabriquer une symétrie particule--antiparticule à partir du
dessin toroïdal.

\subsection{Revêtement double et retour spinoriel}

L’application \(SU(2)\to SO(3)\) est un revêtement double. Un tour de
\(2\pi\) dans \(SO(3)\) peut se relever en \(-I\) dans \(SU(2)\) ; le tour de
\(4\pi\) récupère \(I\). La structure de signe ne transforme pas un cercle en
ruban de Möbius. Pour obtenir un ruban, il faut un fibré réel en droites
dont le caractère d’holonomie soit \(-1\).

Soit \(L_\chi\to S^1\) le fibré associé à
\[
\chi:\pi_1(S^1)\longrightarrow\{\pm1\},
\qquad \chi(1)=-1.
\]
Son espace total est un ruban de Möbius. La route \(C_M=-\operatorname{rev}\)
apporte le candidat de signe ; le cercle de Villarceau apporte le lacet ;
l’entrelaceur physique doit prouver que le premier est l’holonomie du second.

\subsection{Réalisation toroïdale de l’opérateur \texorpdfstring{\(B_{\mathrm{mem}}^{\varepsilon}\)}{Bmem} et détermination de l’angle de Villarceau}

Avec \(B_c=7I+2S\), définissons l’opérateur orienté de mémoire
\[
B_{\rm mem}^{\varepsilon}=5I+2\varepsilon S,
\qquad\varepsilon=\pm1.
\]
Pour \(\varepsilon=+1\),
\[
\Spec(B_{\rm mem}^{+})=\{7,3\},
\qquad\det B_{\rm mem}^{+}=21.
\]
Interpréter les valeurs \((7,3)\) comme les bords équatoriaux extérieur et intérieur
dans une carte positive commune produit
\[
R+r=7\ell,\qquad R-r=3\ell,
\]
et donc
\[
R=5\ell,\qquad r=2\ell,\qquad
\sin\beta=\frac rR=\frac25.
\]
L’angle de Villarceau est
\[
\beta=\arcsin\frac25
=23.5781784782\ldots^\circ.
\]

Le théorème algébrique \(\{7,3\}\leftrightarrow(5,2)\) est exact. La prémisse
géométrique qui interprète le spectre comme les deux bords d’un tore est une
interface déclarée. Elle ne doit pas être dissimulée et l’angle ne doit pas être employé comme cible.

La réalisation du bloc complet \(B_c\) conserve une autre lecture :
\(R:r=7:2\) et \(\beta=\arcsin(2/7)\). Ce ne sont pas des valeurs rivales ; elles relèvent de
deux foncteurs : bloc total et mode de mémoire. La priorité physique du second
doit être décidée par l’entrelaceur avec la dynamique, et non par une proximité décimale avec
une autre constante.

\subsection{Action de l’opérateur \texorpdfstring{\(J^2=-I\)}{J2=-I} sur les coordonnées angulaires conjuguées}

L’angle natif de HMT est une phase \(u\in\mathbb R/\mathbb Z\). Les lecteurs
publient
\[
\theta_{\rm rad}=2\pi u,
\qquad
\theta_{\rm deg}=360u.
\]
Radians et degrés sont des cartes distinctes du même élément circulaire. Le radian
n’est pas plus « ontologique » parce qu’il est adimensionnel ; sa naturalité provient de la
longueur d’arc. Les degrés fournissent une autre coordonnée absolue du
tour. Toute relation avec \(\alpha^{-1}\), \(\varphi\) ou CKM doit
d’abord être formulée comme un quotient ou une holonomie invariante, puis seulement
être évaluée dans une carte angulaire.

L’identité exacte déjà disponible
\[
\tanh^2(2\log\varphi)=\frac59
\]
relie le canal Perron--Fibonacci au quotient spectral de \(B_c\). Son
complément est \(4/9\), le saut normalisé. Cette relation est structurelle.
Elle n’implique pas que tout angle construit avec \(\varphi\) soit un angle physique.

\subsection{Brouwer, Möbius et degré}

Le théorème du point fixe de Brouwer, les transformations fractionnaires de
Möbius et le ruban de Möbius sont trois objets distincts. Ils ne sont reliés
qu’après fixation du domaine, de l’action et du fibré. Les orbites elliptiques d’avance et de
retard peuvent être modélisées comme deux sections d’un fibré à holonomie de
signe ; la transformation de Möbius décrit l’action conforme dans la carte ;
le théorème de Brouwer garantit un point fixe pour les applications continues du disque.
Aucune de ces phrases ne remplace le diagramme de réalisation.

\begin{remark}[Résultat et frontière]
On démontre le degré double de la courbe de Hopf, l’orientation de la
conjugaison totale et la construction algébrique \(7,3\mapsto5,2\). La lecture
toroïdale (5{:}2), l’holonomie de Möbius et l’interprétation
particule--antiparticule forment une chaîne d’interfaces qui doit partager un
unique entrelaceur. Le volume n’impose pas cet entrelaceur au moyen d’une
coïncidence angulaire.
\end{remark}

%
\endgroup

\HMTFlushCurrentChapter
\chapter{Dérivation Hensel--Witt de \texorpdfstring{\(-1/12\)}{-1/12}: prolongement compatible, régularisation et mémoire modulaire}%
\begingroup
\renewcommand{\chapter}[2][]{}%
\chapter{Dérivation Hensel--Witt de \(-1/12\) : prolongement compatible, régularisation et mémoire modulaire}
\label{chap:p3-final:33:menos_un_doceavo}
\label{chap:p3-83-27-menos_un_doceavo}

\section{Domaine Hensel--Witt, applications de Teichmüller et cible rationnelle \(-1/12\)}

Le NTHW est le paquet de données et d’applications
\[
\NTHW=
\bigl(
\mathscr R_\pi^{\rm micro},
\Gnine,
\Cdoce,
K,
\mathcal W_{12\to24},
\mathcal I_{6\to8}
\bigr),
\]
avec les compatibilités suivantes :
\begin{enumerate}
  \item les régions de Hensel partagent un mot quotient et conservent
  des mémoires distinctes ;
  \item \(\Gnine\) transporte l’orientation et produit une monodromie ;
  \item \(\Cdoce\) organise l’observable signé et le sceau \(K\) ;
  \item la lecture de Witt transforme les supports du même bord en hexades et en
  étoiles ;
  \item l’extension \(W_{12}\to W_{24}\) produit l’interface
  hexade--octade ;
  \item les lecteurs archimédien, exceptionnel et dimensionnel reçoivent
  des invariants du même état.
\end{enumerate}

\begin{figure}[!htbp]
\centering
\resizebox{\textwidth}{!}{\begin{tikzpicture}[font=\small,node distance=1.2cm and 1cm,
  nodebox/.style={draw,rounded corners=2mm,align=center,minimum width=3cm,minimum height=1.05cm},
  arr/.style={-{Stealth[length=2.4mm]},thick}]
  \node[nodebox,draw=hmtblue,fill=hmtlightblue] (hensel) at (-5,2.3) {cylindres henséliens\\cocycle de retenue et \(\MonNine\)};
  \node[nodebox,draw=hmtblue,fill=hmtlightblue] (regions) at (-5,-.1) {cinq régions de \(\pi\)\\régions de \(e,\varphi\)};
  \node[nodebox,draw=hmtgold,fill=hmtlightgold,minimum width=3.8cm,minimum height=2.05cm,font=\bfseries] (core) at (0,1.1) {\(\mathscr T_{\rm HW}\)\\Correspondances\\Hensel--Witt};
  \node[nodebox,draw=hmtviolet,fill=hmtlightviolet] (c12) at (0,-1.8) {\(\Cdoce\), \(\Usgn\), \(K\)\\différences \(D_3,D_4,Q\)};
  \node[nodebox,draw=hmtgreen,fill=hmtlightgreen] (witt) at (5,2.3) {Witt \(W_{12}\)\\supports et \(M_{12}\)};
  \node[nodebox,draw=hmtgreen,fill=hmtlightgreen] (oct) at (5,-.1) {inclusion d’incidence \(6\to8\)\\canaux \(90/120\), normalisation \(-1/12\)};
  \node[nodebox,draw=hmtred,fill=hmtlightred] (readers) at (5,-2.5) {lecteurs coordonnés\\valeur et symétrie};
  \draw[arr,hmtblue] (hensel)--(core);
  \draw[arr,hmtblue] (regions)--(core);
  \draw[arr,hmtgold] (core)--(c12);
  \draw[arr,hmtgreen] (core)--(witt);
  \draw[arr,hmtgreen] (core)--(oct);
  \draw[arr,hmtred] (c12)--(readers);
  \draw[arr,hmtred] (witt.east)--++(1.9,0)|-(readers.east);
  \draw[arr,hmtred] (oct)--(readers);
  \node[align=center,hmtgray,font=\footnotesize] at (0,-3.45) {Compatibilité des cylindres, de la dodécaphase et des incidences sur un même état enrichi.};
\end{tikzpicture}
}
\caption{Le NTHW comme interface typée. La tâche mathématique centrale consiste à faire
commuter le diagramme, et non à accumuler des coïncidences autour d’un nom.}
\label{p3-83-c27-s1:fig:nthw}
\end{figure}

Le NTHW donne un nom formel au lieu auparavant décrit comme « pierre de
Rosette ». La substitution est conceptuelle : une pierre traduit par ressemblance ; une
transduction précise l’entrée, la sortie, la donnée conservée et la perte
d’information.

\section{Application de Teichmüller du code ternaire à l’anneau de Witt}

La frontière réversible d’APP sélectionne, avec des jauges déclarées, une matrice
\(A_W\) sur \(\FF_3\) qui satisfait
\[
A_WA_W^T=2I\pmod3.
\]
Le code graphe
\[
C_W=\{(q,qA_W):q\in\FF_3^6\}
\]
est le code de Golay ternaire étendu. Son énumérateur de poids est
\[
1+264y^6+440y^9+24y^{12}.
\]
Les \(264\) mots de poids six, identifiés au signe près, donnent \(132\)
supports. Chaque sous-ensemble de cinq points appartient à une unique hexade :
\[
W_{12}=S(5,6,12).
\]

Le champ de \(495\) involutions étoile engendre un groupe d’ordre
\(95\,040\), reconnu comme \(M_{12}\). Une étoile isolée n’engendre pas le
groupe. Le repère HMT ordonné possède un stabilisateur trivial et fixe donc une
jauge dans l’action.

\section{Incidence hexade--octade et canaux modulaires \(90/120\)}

Soit \(H\) une hexade et \(O_H\) une octade qui la contient. En marquant un point
et une paire de points de l’hexade, on obtient
\[
F_6(H)=H\times\binom H2,
\qquad
|F_6|=6\binom62=90.
\]
Si le point marqué peut parcourir l’octade :
\[
F_8(H)=O_H\times\binom H2,
\qquad
|F_8|=8\binom62=120.
\]
L’inclusion \(H\hookrightarrow O_H\) induit \(F_6\hookrightarrow F_8\).
Son défaut normalisé est
\[
\boxed{
\frac{90-120}{24\binom62}
=\frac{6-8}{24}
=-\frac1{12}.}
\]
Ce théorème combinatoire précède la lecture angulaire de Barbero. Les
deux objets partagent l’interface \(6\to8\), mais ne sont pas identiques.

\section{4 routes compatibles vers \texorpdfstring{\(-1/12\)}{-1/12}}

La construction contient quatre réalisations qui doivent être exposées ensemble et non
comptées comme quatre preuves indépendantes d’une même provenance.

\subsection{Réalisation angulaire par l’incrément élémentaire dodécaphasique}

Une dent de \(\Cdoce\) mesure \(30^\circ\). Alors
\[
\frac{30}{120}-\frac{30}{90}
=\frac14-\frac13
=-\frac1{12}.
\]
C’est l’ombre angulaire du motif \(90/120\).

\subsection{Pseudo-inverse orientée de l’opérateur de Gram sur le sous-espace survivant}

Dans le saut certifié \(60\to81\), deux fibres ont pour tailles \(0\) et
\(12\). L’opérateur de Gram est
\[
K_{60,81}=\operatorname{diag}(0,12).
\]
Sur le sous-espace vivant, sa pseudo-inverse vaut \(1/12\) ; avec l’orientation
négative de bord,
\[
E_{\rm surv}=-K_{60,81}^{+}
\]
produit \(-1/12\).

\subsection{Extraction d’échelle par 2e différence triadique}

La réalisation par un extracteur d’échelle triadique est la construction par
seconde différence développée intégralement dans
\cref{sec:c27-segunda-diferencia-triadica}. Cette première apparition fixe sa
fonction dans l’inventaire des réalisations ; la résidence indiquée conserve
les définitions de \(T_L\), \(a(L)\), \(A_1(L)\), \(C(L)\), l’indépendance
du résidu à l’égard de \(L\) et le statut exact de la normalisation.

\subsection{Réalisation modulaire par le quotient \(\eta(\tau)/\eta(3\tau)\) et la fonction \(t_3\)}

La fonction êta satisfait
\[
\frac{\eta(\tau)}{\eta(3\tau)}
=q^{-1/12}\prod_{n\ge1}(1+q^n+q^{2n})^{-1}.
\]
L’exposant provient de \(1/24-3/24=-1/12\). En prenant douze copies :
\[
t_3(q)=\left(\frac{\eta(\tau)}{\eta(3\tau)}\right)^{12}
=q^{-1}-12+54q-76q^2-243q^3+\cdots.
\]
Le coefficient \(54\) coïncide avec le défaut APP. Le corridor
\[
j=\frac{(t_3+27)(t_3+243)^3}{t_3^3}
\]
se relie ensuite à la branche modulaire classique.

\begin{remark}[Ce qui a été démontré dans le corridor]
Le défaut APP \(54\), l’exposant êta \(-1/12\), le défaut
hexade--octade et les identités modulaires une fois \(t_3\) fixé sont exacts.
L’affirmation forte est qu’ils soient tous des images naturelles d’un unique opérateur
TPK. Cette naturalité doit être démontrée au moyen du diagramme du NTHW ; elle ne se
déduit pas de la seule répétition des mêmes entiers.
\end{remark}


\section{5 réalisations algébriques de \texorpdfstring{\(-1/12\)}{-1/12} et séparation de leurs domaines}

Outre le défaut normalisé de l’inclusion hexade--octade, apparaissent
cinq réalisations opératorielles du même rationnel. Leurs domaines sont
précisés séparément ; l’égalité des valeurs n’identifie pas les
opérateurs sans un entrelaceur supplémentaire.

\subsection{Différence de périodicités dodécaphasiques}

Une dent de \(\Cdoce\) mesure \(30^\circ\). Alors
\[
\frac{30}{120}-\frac{30}{90}
=\frac14-\frac13
=-\frac1{12}.
\]
C’est la différence orientée entre les périodes de \(90^\circ\) et
\(120^\circ\) dans cette réalisation.

\subsection{Pseudo-inverse de l’opérateur de Gram sur la composante vivante}

Dans le saut exact \(60\to81\), deux fibres ont pour tailles \(0\) et
\(12\). L’opérateur de Gram est
\[
K_{60,81}=\operatorname{diag}(0,12).
\]
Sur la composante non nulle de son image, la pseudo-inverse vaut \(1/12\) ; avec l’orientation
négative de bord,
\[
E_{\rm surv}=-K_{60,81}^{+}
\]
produit \(-1/12\).

\subsection{2e différence entre échelles triadiques}
\label{sec:c27-segunda-diferencia-triadica}

Définissons
\[
T_L=\sum_{n=1}^{L-1}n(L-n)=\frac{L(L^2-1)}6,
\qquad
a(L)=-\frac{T_L}{L^2}=-\frac L6+\frac1{6L}.
\]
La première soustraction d’échelle est
\[
A_1(L)=a(L)-\frac{a(3L)}3=\frac4{27L},
\]
et la seconde
\[
C(L)=LA_1(L)-\frac{3L}{3}A_1(3L)=\frac8{81}.
\]
Le résidu \(8/81\) est indépendant de \(L\). Pour obtenir
\(-1/12\), on applique la normalisation déclarée
\[
R(L)=-\frac{27}{32}C(L)=-\frac1{12}.
\]
Le calcul dérive \(8/81\) ; il ne dérive pas à lui seul le facteur \(-27/32\).

\subsection{Cadena modular}

La fonction êta satisfait
\[
\frac{\eta(\tau)}{\eta(3\tau)}
=q^{-1/12}\prod_{n\ge1}(1+q^n+q^{2n})^{-1}.
\]
L’exposant provient de \(1/24-3/24=-1/12\). En prenant douze copies :
\[
t_3(q)=\left(\frac{\eta(\tau)}{\eta(3\tau)}\right)^{12}
=q^{-1}-12+54q-76q^2-243q^3+\cdots.
\]
Le coefficient \(54\) coïncide avec le défaut APP\@. L’identité rationnelle
\[
j=\frac{(t_3+27)(t_3+243)^3}{t_3^3}
\]
établit la connexion avec la branche modulaire classique.

\Needspace{32\baselineskip}
\subsection{Défaut des projecteurs de différence cyclique}

Soit \(S\) le décalage cyclique sur \(\QQ^{12}\) et définissons
\[
 D_r=I-S^r.
\]
Le noyau de \(D_r\) est formé des vecteurs constants sur chaque orbite
de \(S^r\) ; par conséquent,
\[
 \dim\ker D_r=\gcd(12,r).
\]
Soient \(\Pi_{\ker D_3}\) et \(\Pi_{\ker D_4}\) les projecteurs rationnels sur
\(\ker D_3\) et \(\ker D_4\), et soit
\(\tau(T)=\operatorname{Tr}(T)/12\) la trace normalisée. Alors
\[
 \boxed{
 \tau(\Pi_{\ker D_3}-\Pi_{\ker D_4})
 =\frac{\operatorname{rank}\Pi_{\ker D_3}
       -\operatorname{rank}\Pi_{\ker D_4}}{12}
 =\frac{3-4}{12}
 =-\frac1{12}.}
\]
Cette réalisation exprime le rationnel comme défaut transversal des pas
trois et quatre dans le cycle à douze phases.

\begin{remark}[Domaines des réalisations]
Le défaut APP \(54\), l’exposant êta \(-1/12\), le défaut
combinatoire du relèvement hexade--octade et les
identités modulaires une fois \(t_3\) fixé sont exacts. La valeur
de la pseudo-inverse, l’extracteur normalisé et le défaut des projecteurs sont également exacts.
Chaque égalité conserve son domaine et son application correspondante.
\end{remark}

%
\endgroup

\HMTFlushCurrentChapter
\chapter{Dérivation déterminantielle de l'échelle absolue d'action : forme normale de Smith, conoyau d'ordre \(16\) et valuation \texorpdfstring{\(10^{-34}\)}{10^-34}}%
\begingroup
\renewcommand{\chapter}[2][]{}%
\chapter{Dérivation déterminantale de l’échelle d’action : forme normale de Smith, conoyau d’ordre \(16\) et valuation \(10^{-34}\)}
\label{chap:p3-final:34:escala_accion}
\label{ch:vacuum}

\section{Dérivation de l'échelle d'action à partir de l'état dodécaphasique}

La décennie d'action est déterminée intrinsèquement, sans l'adopter comme
unité héritée. La fonctionnelle déterminantale locale s'applique à la sortie HMT
préalablement construite \(\alpha_{\HMT}\) et définit
\begin{equation}
\Sigma_H=\varphi_{\HMT}\alpha_{\HMT}^{16}.
\label{eq:action-determinantal-reader}
\end{equation}
L'exposant \(16\) est l'ordre du conoyau local de la clôture signée ; sa valuation décimale fixe
\begin{equation}
\operatorname{ord}_{10}(\Sigma_H)=-34.
\label{eq:action-decade}
\end{equation}

L'origine de l'exposant peut être vérifiée dans ce volume. Sur
l'orbite locale à quatre feuillets agit le bloc de Hadamard
\begin{equation}
H_4=
\begin{pmatrix}
1&1&1&1\\
1&1&-1&-1\\
1&-1&1&-1\\
1&-1&-1&1
\end{pmatrix}.
\label{eq:action-hadamard}
\end{equation}

\begin{lemma}[Indice local de la droite d'action]
La forme normale de Smith de \(H_4\) est
\begin{equation}
\operatorname{SNF}(H_4)=\operatorname{diag}(1,2,2,4).
\label{eq:action-smith}
\end{equation}
En conséquence,
\[
\operatorname{coker}(H_4)\simeq
\Z/2\Z\oplus\Z/2\Z\oplus\Z/4\Z,
\qquad
|\operatorname{coker}(H_4)|=16.
\]
\end{lemma}

\begin{proof}
Le plus grand commun diviseur des entrées est un ; celui des mineurs
\(2\times2\) non nuls est deux ; celui des mineurs \(3\times3\) est quatre ; et
\(|\det H_4|=16\). Les quotients successifs de ces diviseurs déterminantaux
sont \(1,2,2,4\), qui prouvent \cref{eq:action-smith}. Le produit des
invariants donne l'ordre du conoyau.
\end{proof}

La norme de la section \(\alpha_{\HMT}\) sur ces seize feuillets
produit \(\alpha_{\HMT}^{16}\) ; l'auto-échelle \(\varphi_{\HMT}\) agit
une seule fois sur la base, et non une fois par feuillet. D'où
\cref{eq:action-determinantal-reader}. Les comparaisons certifiées
\begin{equation}
\alpha_{\HMT}^{16}<10^{-34}
<\varphi_{\HMT}\alpha_{\HMT}^{16}<10^{-33}
\label{eq:action-decade-bounds}
\end{equation}
démontrent, sans introduire l'unité SI, que l'ordre décimal est
exactement \(-34\). L'évaluation terminale commence par
\begin{equation}
\Sigma_H
=1.046243838104756580184229208303089\ldots\times10^{-34}.
\label{eq:action-sigma-value}
\end{equation}
Le remplacement du conoyau par trois copies avec exposant \(16^3\), ou
l'application réitérée de \(\varphi\) dans chaque feuillet, modifierait la
généalogie. Les deux opérations constituent des critères de réfutation et non
des conventions équivalentes.
La chaîne causale qui est conservée est donc
\[
\APP\to\TRIT\to\TPK\to U_{12}^{\rm sign}
\to\alpha_{\HMT}\to\Sigma_H
\to10^{-34}\U_S\to
\{\hbar_{\rm pre},\hbar_{\rm ret}\}.
\]
Ce volume ne redérive pas \(\alpha_{\HMT}\), mais développe la géométrie
métrologique postérieure à la sélection du type et de la décennie par
\(\Sigma_H\). L'entier \(-34\) n'est ni une mantisse ni un ajustement au Système
international, mais la valuation terminale de la fonctionnelle locale.

\section{Sections orientées de la droite d'action}

Les réalisations antérieure et postérieure au retour constituent deux sections
d'une même droite d'action :

\begin{equation}
\hbar_{\rm pre}=H_5\,10^{-34}\U_S,\qquad
\hbar_{\rm ret}=\eta_{\rm ret}\,10^{-34}\U_S.
\label{eq:action-twoway}
\end{equation}

La décennie commune détermine la droite, tandis que l'information
orientée s'exprime par

\begin{equation}
R_{\rm act}=\frac{\hbar_{\rm pre}}{\hbar_{\rm ret}}
=\frac{H_5}{\eta_{\rm ret}},
\qquad
\xi_{\rm act}=\frac12\log R_{\rm act}.
\label{eq:action-rapidity}
\end{equation}

La première quantité fournit une coordonnée multiplicative et la seconde
une coordonnée additive de type rapidité. Elles ne représentent pas deux constantes de
Planck indépendantes, mais les deux orientations d'une même section
par rapport au retour.

Si \(h_a=2\pi\hbar_a\), \(a\in\{\mathrm{pre},\mathrm{ret}\}\), une grandeur proportionnelle à \(h_a^{1/2}\) varie comme \(R_{\rm act}^{1/2}\), une grandeur proportionnelle à \(h_a^{-1/2}\) varie comme \(R_{\rm act}^{-1/2}\), et un quotient où l'action s'annule est un invariant bidirectionnel.


%
\endgroup
% Distribución causal exacta del delta Ley 9; contenido conservado sin síntesis.
\Needspace{28\baselineskip}
\section{Bifurcation causale postérieure à \texorpdfstring{\(\alpha_{\mathrm{HMT}}\)}{alpha HMT}}
\label{sec:l9s102:bifurcacion-accion-angular}

La généalogie de l'action et du couple angulaire forme un diagramme bifurqué.
La branche régionale produit la mantisse réduite \(\eta_{\mathrm{ret}}\) ; la branche
déterminantale produit l'ordre décimal \(\varepsilon_H=-34\). Les deux
convergent dans la droite d'action. Le couple angulaire est construit dans une branche
sœur à partir de \(\alpha_{\mathrm{HMT}}\) et de
\(\eta_{\mathrm{ret}}\) :
\begin{equation}
\begin{tikzcd}[column sep=small,row sep=normal]
\alpha_{\mathrm{HMT}},\pi,\varphi
  \arrow[r] \arrow[d]
& H_5,\mathscr C_\pi
  \arrow[d]
&& \alpha_{\mathrm{HMT}},\varphi,H_4
  \arrow[d] & {}\\
x_A,D_A,\rho_\pi
  \arrow[r]
& \eta_{\mathrm{ret}}
  \arrow[r] \arrow[d]
& \hbar_{\mathrm{ret}}\in\mathcal L_S
& \Sigma_H,\varepsilon_H=-34
  \arrow[l] & {}\\
& (A_{\deg},C^*_{\deg})
  \arrow[r]
& (A_{\mathrm{rad}},C^*_{\mathrm{rad}})
  \arrow[r]
& (q_+,q_-)
  \arrow[r]
& V_{90,120}.
\end{tikzcd}
\label{diag:l9s102:dag-accion-angular}
\end{equation}
L'ordre d'exposition situe la dérivation de Planck avant la publication
angulaire complète ; le diagramme conserve simultanément la dépendance
mathématique exacte de chaque sortie.

\section{Conoyau local et dérivation de l'ordre absolu \texorpdfstring{\(10^{-34}\)}{10^-34}}
\label{sec:l9s102:orden-accion}

La forme normale de Smith du bloc dodécaphasique local est
\begin{equation}
 \operatorname{SNF}(H_4)=\operatorname{diag}(1,2,2,4),
 \qquad
 \operatorname{coker}(H_4)
 \cong\mathbb Z/2\mathbb Z\oplus\mathbb Z/2\mathbb Z
 \oplus\mathbb Z/4\mathbb Z.
 \label{eq:l9s102:snf-h4}
\end{equation}
Le conoyau local est d'ordre \(16\). Pour la section constante
\(s_\alpha(g)=\alpha_{\mathrm{HMT}}\), la norme locale et une action de
\(\varphi\) produisent
\begin{equation}
 \boxed{
 \Sigma_H=\varphi\mathcal N_{G_H}(s_\alpha)
 =\varphi\alpha_{\mathrm{HMT}}^{16}.}
 \label{eq:l9s102:sigma-h}
\end{equation}
Les inégalités exactes
\begin{equation}
 \alpha_{\mathrm{HMT}}^{16}<10^{-34}
 <\varphi\alpha_{\mathrm{HMT}}^{16}<10^{-33}
 \label{eq:l9s102:orden-34}
\end{equation}
déterminent
\begin{equation}
 \boxed{
 \varepsilon_H=\left\lfloor\log_{10}\Sigma_H\right\rfloor=-34.}
 \label{eq:l9s102:epsilon-h}
\end{equation}


\HMTFlushCurrentChapter
\chapter{Caractère régional de degré \(5\), construction de \texorpdfstring{\(\eta_{\mathrm{ret}}\)}{eta ret} et mesure positive d'action sur le quotient des émissions}%
\begingroup
\renewcommand{\chapter}[2][]{}%
\chapter{Mesure de quotient et caractère positif de l’action : hilbertisation canonique, défaut régional--dodécaphasique et loi fonctorielle sur les chemins}
\label{chap:p3-final:35:medida_accion_positiva}
\label{p3-20260826:chap:medida-accion-positiva}
\index{mesure!catalogue d’émissions}
\index{action!cocycle positif}
\index{trace normalisée}

La section précédente identifie isométriquement un secteur positif de rang trois dans
l’interface orientée \(t=7\), la dodécaphase et le module d’incidence de Witt. La construction
suivante dérive l’autre coefficient de la loi angulaire, \(5/468\), et
démontre pourquoi les deux termes se combinent par une somme positive. La
question exige de séparer
trois opérations :
\begin{enumerate}
  \item compter les présentations microscopiques ;
  \item compter les émissions distinctes après formation du quotient TPK ;
  \item attribuer une action additive à une route.
\end{enumerate}
Les deux premières opérations produisent des mesures différentes. La troisième se
caractérise par une propriété universelle et non par la proximité ultérieure
d’un nombre décimal.

Nous noterons \(\Gtr\) le graphe orienté à \(1944\) états du lecteur de
chemins enrichis. Sa catégorie libre de chemins sera notée
\(\mathsf P(\Gtr)\). La définition locale de ses arêtes et de son action est
développée dans la section correspondante.

\section{Espace de germes, quotient d’émissions et fibres régionales}

Soit
\[
 \Omega_{\rm seed}
 =\bigl((\mathbb Z/9\mathbb Z)^2\times D_4\bigr)^2,
 \qquad
 |\Omega_{\rm seed}|=104\,976,
\]
l’ensemble exécutable des paires de germes de l’émetteur fini, et soit
\[
 F:\Omega_{\rm seed}\twoheadrightarrow\mathcal U_6,
 \qquad
 |\mathcal U_6|=468,
\]
l’application qui conserve l’émission sextuple distincte. Le recensement exact des
fibres est
\[
\begin{array}{c|c|c}
|F^{-1}(U)|&\#\{U\text{ de cette taille}\}&\text{germes apportés}\\
\hline
144&432&62\,208\\
432&18&7\,776\\
1\,944&18&34\,992.
\end{array}
\]
Les trois contributions ont pour somme \(104\,976\). Par conséquent, \(F\) n’est pas un
revêtement régulier : des émissions différentes possèdent des nombres différents de
présentations microscopiques.

La microfibre de \(\pi\) contient cinq émissions distinctes. Quatre ont
pour multiplicité \(144\) et une a pour multiplicité \(432\). Cette donnée suffit
à démontrer que la mesure uniforme sur les germes et la mesure uniforme
sur les émissions ne sont pas identiques.

\section{\(2\) mesures exactes sur le quotient des émissions et des mots}

Si chaque germe reçoit le poids \(1/104\,976\), la mesure image par \(F\) est
\[
 \nu_{\rm seed}(U)
 =\frac{|F^{-1}(U)|}{104\,976}.
\]
Ses atomes possibles sont
\[
 \frac1{729},\qquad\frac1{243},\qquad\frac1{54}.
\]
\Needspace{5\baselineskip}
En particulier,
\[
 \nu_{\rm seed}(\mathscr R_\pi^{\rm micro})
 =\frac{4\cdot144+432}{104\,976}
 =\boxed{\frac7{729}}.
\]

Si l’état élémentaire est, en revanche, une émission distincte
\(U\in\mathcal U_6\), l’algèbre diagonale des observables est
\[
 \mathcal D_{468}=\ell^\infty(\mathcal U_6)
\]
et son état de comptage normalisé est
\[
 \tau_{468}(f)
 =\frac1{468}\sum_{U\in\mathcal U_6}f(U).
\]
Soit \(\Pi_\pi^{(468)}\) la fonction caractéristique des cinq régions de la
microfibre. Alors
\[
 \boxed{\tau_{468}(\Pi_\pi^{(468)})=\frac5{468}.}
\]
La différence exacte entre les deux lectures est
\[
 \frac5{468}-\frac7{729}
 =\frac{41}{37\,908}.
\]
Ce ne sont pas deux approximations. La première mesure compte les présentations ; la
seconde compte les classes d’émission.

\begin{theorem}[Unicité de la mesure uniforme du catalogue]
\label{p3-20260826:thm:unicidad-medida-468}
L’unique probabilité sur \(\mathcal U_6\) invariante sous tout
réétiquetage de ses \(468\) atomes est la mesure uniforme.
De manière équivalente, l’unique trace positive et unitaire sur
\(M_{468}(\mathbb C)\) est
\[
 \tau_{468}(A)=\frac1{468}\operatorname{Tr}A.
\]
\end{theorem}

\begin{proof}
Les transpositions envoient tout atome sur tout autre. L’invariance
impose qu’ils aient tous la même masse, et la normalisation la fixe à
\(1/468\). Dans la version matricielle, l’invariance unitaire égalise la valeur
de tous les projecteurs de rang un. La décomposition spectrale et la
linéarité reconstruisent la trace normalisée.
\end{proof}

L’hypothèse décisive est manifeste : le lecteur travaille avec le
\emph{catalogue quotient}. Le théorème n’affirme pas que la mesure uniforme des
germes devienne uniforme lorsqu’on oublie les fibres. La symétrie utilisée
est la naturalité de l’objet fini sans étiquettes, et non une propriété
ergodique attribuée sans preuve à la dynamique microscopique.

\section{Hilbertisation canonique du quotient}

La relation entre germes et émissions peut être formulée sans ambiguïté.
Soient
\[
 H_{\rm seed}=\ell^2(\Omega_{\rm seed}),
 \qquad
 H_U=\ell^2(\mathcal U_6).
\]
L’émetteur détermine
\[
 \boxed{
 J_F\delta_U
 =|F^{-1}(U)|^{-1/2}
 \sum_{\omega\in F^{-1}(U)}\delta_\omega.
 }
\]
Les fibres distinctes sont disjointes et chaque colonne a pour norme un ; par
conséquent,
\[
 J_F^*J_F=I_{H_U}.
\]
De plus, pour tout observable \(f\) du catalogue,
\[
 M_{f\circ F}J_F=J_FM_f.
\]
L’image \(J_FH_U\) est exactement le sous-espace des fonctions constantes
sur chaque fibre.

Si \(F^\sharp:H_U\to H_{\rm seed}\) désigne l’application induite par
précomposition, sans normalisation,
alors
\[
 J_F
 =F^\sharp\bigl((F^\sharp)^*F^\sharp\bigr)^{-1/2}.
\]
C’est le facteur isométrique de la décomposition polaire de cette application. Cette
expression démontre sa naturalité sous les isomorphismes et les réétiquetages de
l’émetteur. Elle n’implique pas une fonctorialité fausse sous toute composition de
surjections non régulières : la normalisation dépend des fibres de
l’application concrète.

Dans le langage des mesures, le relèvement de la trace du quotient qui est
uniforme à l’intérieur de chaque fibre vaut
\[
 \widetilde\tau(\omega)
 =\frac1{468\,|F^{-1}(F\omega)|}.
\]
Chaque fibre reçoit une masse totale \(1/468\). La compensation inverse par
la multiplicité, et non une invocation informelle des orbites, constitue la preuve exacte.

\begin{corollary}[Canonicité de \(5/468\)]
\label{p3-20260826:cor:canonicidad-cinco-468}
Une fois le lecteur régional défini sur les émissions distinctes, la masse
de la microfibre de \(\pi\) est imposée :
\[
 \tau_{468}(\Pi_\pi^{(468)})
 =\frac{\operatorname{rank}\Pi_\pi^{(468)}}{468}
 =\frac5{468}.
\]
\end{corollary}

\section{Défaut positif entre les mesures régionale et dodécaphasique}

Écrivons
\[
 H_{\rm ang}=H_U\otimes V_{11}.
\]
Le scalaire interne
\[
 \Delta_4
 =\frac{\pi-e\log\pi}{270}
 \left(1+\frac\pi{729}\right)
 =0.0001111964226921402\ldots
\]
est construit avec les lectures HMT déjà fixées de \(\pi\) et \(e\), et avec les
entiers \(270\) et \(729\) du registre. Il constitue ici une entrée mathématique
prédimensionnelle : ni unités ni observables physiques n’interviennent.

La construction interne du scalaire ne détermine pas à elle seule quel opérateur
il doit multiplier. Pour garder cette prémisse visible, nous déclarons :
\begin{description}
  \item[\HDef --- défaut de changement de feuillet.] La coordonnée dodécaphasique du défaut
  de changement de feuillet est \(\Delta_4P_3\).
\end{description}
\HAdd fixera la forme additive de l’action, mais ne sélectionne pas \HDef : pour tout
\(c\ge0\), le même argument admettrait la paire
\((\Pi_\pi^{(468)},cP_3)\). Démontrer \HDef directement à partir d’une transition
totale du noyau topologique de phase enrichi est une condition supplémentaire.

Nous définissons la somme de Kronecker
\[
 \boxed{
 D_{\rm switch}^{\rm dod}
 =\Pi_\pi^{(468)}\otimes I_{11}
 +\Delta_4 I_{468}\otimes P_3.
 }
\]
Les deux termes sont positifs, commutent et agissent sur des facteurs différents.
Par conséquent,
\[
 D_{\rm switch}^{\rm dod}\ge0.
\]
L’entrelaceur \(U_W\) du
\cref{thm:entrelazador-dodecafase-witt} produit la carte de Witt :
\[
\begin{aligned}
 D_{\rm switch}^{\rm Witt}
 &=(I_{468}\otimes U_W)
 D_{\rm switch}^{\rm dod}
 (I_{468}\otimes U_W)^*\\
 &=\Pi_\pi^{(468)}\otimes I_{E_{30}}
 +\Delta_4I_{468}\otimes Q_3.
\end{aligned}
\]
Ainsi, le même opérateur positif se lit dans la dodécaphase et dans le module
d’incidence.

La chaîne isométrique orientée identifie également les secteurs actifs
\[
 I_{L_{\rm or}}\otimes P_G,\qquad P_3,\qquad Q_3.
\]
On n’affirme pas une équivalence des compléments complets des trois
modules : \(W_{GK}\) est une équivalence partielle des rangs
sélectionnés. C’est exactement l’information nécessaire pour identifier
la seconde contribution de la trace.

\begin{theorem}[Portée exacte de l’entrelacement et de l’induction de l’action]
\label{p3-20260826:thm:alcance-induccion-accion}
Soient \(W_{GK}\) et \(U_W\) les isométries partielles du
\cref{thm:identificacion-isometrica-tres-proyectores}. Alors :
\begin{enumerate}
  \item pour tout \(c\geq0\),
  \[
   W_{GK}^{*}(cP_3)W_{GK}
   =c(I_{L_{\rm or}}\otimes P_G),
   \qquad
   U_W(cP_3)U_W^{*}=cQ_3;
  \]
  par conséquent, l’entrelacement transporte le coefficient, mais ne le
  sélectionne pas ;
  \item si \(P_{\rm sgn}\) et \(P_{\rm std}\) sont les projecteurs centraux
  de la décomposition
  \[
   P_3V_{11}\big|_{H_*}
   \simeq \operatorname{sgn}_{H_*}\oplus\operatorname{std},
   \qquad
   \operatorname{rank}P_{\rm sgn}=1,
   \quad
   \operatorname{rank}P_{\rm std}=2,
  \]
  tout opérateur positif autoadjoint supporté sur \(P_3V_{11}\) et
  \(H_*\)-équivariant a la forme
  \[
   D_{a,b}=aP_{\rm sgn}+bP_{\rm std},
   \qquad a,b\geq0.
  \]
  Même l’imposition de la trace de \(\Delta_4P_3\) laisse la famille
  \[
   a+2b=3\Delta_4.
  \]
  L’isotropie \(a=b\) ne se déduit pas de la symétrie résiduelle \(S_3\) ;
  \item la probabilité uniforme du catalogue,
  \(\mu_U(U)=1/468\), est la trace canonique du quotient fini et son
  relèvement compensé est celui que détermine \(J_F\). Ce n’est pas la
  mesure image de la probabilité uniforme sur les germes, dont les atomes sont
  \(1/729,1/243,1/54\). En particulier,
  \[
   \mu_U(\mathscr R_\pi^{\rm micro})=\frac5{468},
   \qquad
   F_\#\mu_{\rm seed}(\mathscr R_\pi^{\rm micro})=\frac7{729};
  \]
  \item la loi complète
  \[
   \delta_*=\frac5{468}+\frac3{11}\Delta_4,
   \qquad
   \lambda_*=1-\delta_*,
   \qquad
   a_*(e)=g(e)+\lambda_*s(e)
  \]
  est un théorème relatif aux cinq clauses \HTr--\HAdd--\HDef--\HRes--\HLoc. Le moteur
  exécutable du noyau topologique de phase réduit n’induit pas, par revêtement orienté, l’opérateur ramifié
  \Gtr : la dynamique associée revient à l’identité après \(54\) itérations,
  tandis que chaque état de
  \Gtr possède quatre successeurs distincts et aucune boucle.
\end{enumerate}
Par conséquent, les résultats précédents démontrent la canonicité des identifications
isométriques, de la trace sur le catalogue quotient et du caractère additif
dans le cadre de leurs prémisses. L’affirmation plus forte selon laquelle une transition totale
du noyau topologique de phase enrichi induit nécessairement \(\Delta_4P_3\), la mesure uniforme et la
loi locale exige encore le théorème dynamique supplémentaire spécifié dans
l’analyse de l’induction.
\end{theorem}

\begin{proof}
Les identités du premier point s’obtiennent en multipliant
\(W_{GK}^{*}W_{GK}=I_{L_{\rm or}}\otimes P_G\),
\(W_{GK}W_{GK}^{*}=P_3\) et \(U_WP_3U_W^{*}=Q_3\) par le scalaire \(c\).
Pour le second point, la restriction à \(H_*\simeq S_3\) est une somme sans
multiplicités de deux représentations irréductibles non isomorphes. Le lemme de
Schur diagonalise tout opérateur autoadjoint équivariant dans ces deux blocs ;
la positivité équivaut à \(a,b\geq0\). Sa trace normalisée est
\((a+2b)/11\), ce qui prouve la liberté résiduelle même après fixation de la
trace. L’action complète de \(A_5\) sur le triplet irréductible réduirait
la famille à \(cP_3\), mais ne fixerait pas non plus la valeur de \(c\).

Le troisième point découle du recensement des fibres : \(F\) a pour multiplicités
\(144,432,1944\), de sorte que la mesure image pondère une émission par son
nombre de présentations. La trace \(\tau_{468}\), en revanche, attribue le
même poids à chaque classe ; \(J_F\) réalise exactement son relèvement
normalisé. Le quatrième point est l’enchaînement logique déjà démontré :
\HTr fixe la trace du quotient, \HAdd la somme des traces, \HDef déclare
\(\Delta_4P_3\), \HRes définit la capacité résiduelle et \HLoc la porte sur les arêtes.
L’énumération finie prouve l’obstruction entre le noyau réduit et
\Gtr : un revêtement orienté préserverait l’étoile sortante et une
semi-conjugaison du retour identité exigerait des boucles dans l’image,
en contradiction avec les deux recensements de \Gtr.
\end{proof}

Le spectre complet de \(D_{\rm switch}^{\rm dod}\) est
\[
\begin{array}{c|c}
\text{valeur propre}&\text{multiplicité}\\ \hline
0&(468-5)(11-3)=3704\\
\Delta_4&(468-5)3=1389\\
1&5(11-3)=40\\
1+\Delta_4&5\cdot3=15.
\end{array}
\]
Par conséquent,
\[
\begin{aligned}
 \tau_{468\cdot11}(D_{\rm switch})
 &=\tau_{468}(\Pi_\pi^{(468)})+\Delta_4\tau_{11}(P_3)\\
 &=\boxed{\frac5{468}+\frac3{11}\Delta_4}.
\end{aligned}
\]
La même identité vaut dans la carte de Witt. Dans la carte du secteur orienté, le terme
\(3/11\) est la trace de \(P_G\) sur l’espace homogène à onze branches.

\section{Propriété universelle de la fonctionnelle positive d’action}

Le mot \emph{action} est employé au sens additif : la concaténation de deux
routes reçoit la somme de leurs incréments. Pour séparer ce contenu d’une
formule choisie, nous définissons d’abord la catégorie dans laquelle la loi sera unique.

\begin{definition}[Caractère positif d’action HMT]
\label{p3-20260826:def:caracter-accion-hmt}
Un caractère local d’action sur le produit
régional--dodécaphasique est une application
\[
 \delta:
 M_{468}(\mathbb C)_+\oplus M_{11}(\mathbb C)_+
 \longrightarrow\mathbb R_{\ge0}
\]
qui satisfait :
\begin{enumerate}
  \item additivité exacte dans le cône produit ;
  \item homogeneidad positiva;
  \item invariance par conjugaison unitaire indépendante dans les deux
  facteurs ;
  \item normalisation bifactorielle,
  \[
  \delta(I_{468},0)=\delta(0,I_{11})=1.
  \]
\end{enumerate}
\end{definition}

L’additivité traduit la concaténation ; l’invariance élimine les
coordonnées ; la normalisation bifactorielle fixe une échelle dans chaque terme.
Il ne s’agit pas d’une fonctionnelle unitaire sur l’algèbre somme directe : son unité
est \((I_{468},I_{11})\) et l’additivité donne
\(\delta(I_{468},I_{11})=2\). Aucune clause ne contient la valeur d’un
angle.

\begin{theorem}[Unicité du caractère positif et capacité résiduelle relative]
\label{p3-20260826:thm:ley-accion-hmt}
Dans la classe de la définition précédente, il existe un unique caractère positif
d’action HMT\@. C’est
\[
 \boxed{\delta(A,B)=\tau_{468}(A)+\tau_{11}(B).}
\]
Appliqué au défaut déclaré par \HDef,
\((\Pi_\pi^{(468)},\Delta_4P_3)\), il produit
\[
 \delta_*=\frac5{468}+\frac3{11}\Delta_4.
\]
Si l’on déclare en outre \HRes, selon laquelle la capacité autoduale de changement de
feuillet est normalisée à un et le défaut la réduit par
\(\lambda=1-\delta\), la capacité résiduelle dans le cadre de \HTr--\HAdd--\HDef--\HRes est
\[
 \boxed{
 \lambda_*
 =1-\delta_*
 =1-\frac5{468}-\frac3{11}\Delta_4.
 }
\]
\end{theorem}

\begin{proof}
L’additivité sépare les facteurs :
\[
 \delta(A,B)=\delta(A,0)+\delta(0,B).
\]
Chaque restriction est positive, homogène et invariante par conjugaison. Pour
un projecteur de rang un, l’invariance unitaire donne une valeur commune.
L’additivité sur une décomposition spectrale impose que chaque restriction
soit un multiple de la trace. La normalisation bifactorielle fixe ces multiples à
\(1/468\) et \(1/11\). Il en résulte la somme des traces. La dernière formule
applique la normalisation déclarée de la capacité résiduelle.
\end{proof}

Numériquement,
\[
 \boxed{\lambda_*=0.9892859130191415\ldots},
 \qquad
 0<\lambda_*<1.
\]
La positivité découle déjà de
\(0<\Delta_4<10^{-3}\) et \(5/468<1\). Elle ne dépend pas d’une comparaison
métrologique.

\section{Foncteur d’action positive sur la catégorie des chemins compatibles}

Rappelons le graphe orienté \(\Gtr\) et sa catégorie libre
\(\mathsf P(\Gtr)\). Pour une arête
\(e:x\to y\), nous définissons
\[
 g(e)=\mathbf1_{\{d_y\ne d_x\}},
 \qquad
 s(e)=\mathbf1_{\{m_y\ne m_x\}},
\]
où \(g\) enregistre une rotation et \(s\) un changement de feuillet. La flèche du
scalaire \(\lambda_*\) vers les arêtes ne se déduit pas du
théorème des traces. Elle est déclarée par la clause \HLoc : une rotation a un coût
unitaire, un changement de feuillet a pour coût \(\lambda_*\), et les deux indicateurs
s’additionnent lorsqu’ils apparaissent sur la même arête. Sous \HLoc, l’action locale est
\[
 \boxed{a_*(e)=g(e)+\lambda_*s(e).}
\]
Pour une route \(\gamma=e_1\cdots e_n\),
\[
 \mathcal A_*(\gamma)
 =\sum_{j=1}^na_*(e_j)
 =\#\operatorname{giros}(\gamma)
 +\lambda_*\#\operatorname{cambios\ de\ hoja}(\gamma).
\]
Pour la route vide \(\mathbf1_x\), on fixe
\(\mathcal A_*(\mathbf1_x)=0\). Alors
\[
 \boxed{
 \mathcal A_*(\gamma_1\gamma_2)
 =\mathcal A_*(\gamma_1)+\mathcal A_*(\gamma_2).
 }
\]
Par conséquent,
\[
 \mathcal A_*:\mathsf P(\Gtr)\longrightarrow(\mathbb R_{\ge0},+)
\]
est un foncteur monoïdal, ou cocycle additif de chemins.

Les indicateurs \(g\) et \(s\) sont insensibles au sens de parcours.
L’orientation n’est pas perdue : elle réside dans la droite \(L_{\rm or}\) et dans
l’involution étoile, tandis que le coût positif réside dans les
projecteurs. Confondre ces deux données introduirait des signes négatifs dans
l’action.

Pour \(\beta>0\), l’opérateur de transfert associé est
\[
 (\mathcal L_{\beta,*}f)(y)
 =\sum_{e:x\to y}e^{-\beta a_*(e)}f(x).
\]
Tous ses poids sont strictement positifs. Comme le potentiel dépend seulement
des deux indicateurs locaux, il commute avec l’action grossière
\((\mathbb Z/3\mathbb Z)^3\) et respecte la décomposition en \(27\)
secteurs de Bloch utilisée dans la section suivante.

Une dynamique déterministe finie non pondérée agit par des matrices de
permutation et possède des valeurs singulières égales à un. Le cocycle
\(\mathcal A_*\) transforme l’inventaire des chemins en un opérateur de
transfert qui préserve le cône positif, c’est-à-dire en une matrice à
poids strictement positifs sur les arêtes autorisées. Cette propriété ne
démontre pas à elle seule la contraction, la simplicité de la valeur propre dominante ni
la séparation spectrale ; ces conclusions exigent une normalisation ou un
argument spectral supplémentaire.

\ifdefined\HMTInternalTraceability
\section{Induction par le quotient et extension dynamique}

L’expression \emph{induite par le noyau topologique de phase enrichi} possède deux sens qui doivent rester
séparés.

\begin{enumerate}
  \item \textbf{Induction par l’objet quotient.} L’émetteur exécutable
  détermine \(F\), son facteur polaire \(J_F\), l’espace \(H_U\) et la
  trace normalisée. L’état enrichi apporte
  \(\Pi_\pi^{(468)}\), \(P_3\),
  \(K\), \(B_K\), le repère de Witt et \(\Delta_4\). En ce sens
  catégoriel, \(D_{\rm switch}\) se construit uniquement avec des données internes
  APP--TRIT--TPK\@.
  \item \textbf{Induction par une transition totale.} Cette lecture exigerait
  une transition finie exécutable sur toutes les coordonnées de
  \(\mathcal X_{\mathrm{TPK}}^{\mathrm{enr}}\) et une preuve que sa
  mesure image dynamique
  préserve \(J_FH_U\) et produit le potentiel \(a_*\). Cette transition totale
  n’appartient pas au domaine exécutable construit dans cette section.
\end{enumerate}

Le \cref{p3-20260826:thm:ley-accion-hmt} et la construction des chemins closent la première
lecture au moyen de cinq clauses explicites du lecteur :
\begin{description}
  \item[\HTr --- trace du quotient.] Lorsque la multiplicité de présentation
  reste dans la fibre oubliée, une émission distincte est un atome et
  l’état du catalogue est la trace normalisée de son algèbre finie.
  \item[\HAdd --- cocycle d’action.] Le défaut d’une route est le caractère
  positif, normalisé bifactoriellement, additif et invariant du produit
  régional--dodécaphasique.
  \item[\HDef --- défaut de changement de feuillet.] Le second argument du caractère est
  \(\Delta_4P_3\).
  \item[\HRes --- capacité résiduelle.] La capacité autoduale est normalisée à
  un et le défaut la réduit selon \(\lambda=1-\delta\).
  \item[\HLoc --- action locale.] Sur les arêtes de \Gtr, on fixe
  \(a(e)=g(e)+\lambda s(e)\), et l’action d’une route est la somme sur ses
  arêtes.
\end{description}
\HTr impose \(5/468\). \HAdd impose la somme des traces et exclut les termes
mixtes. \HDef fixe le scalaire qui occupe la coordonnée de commutation. \HRes transforme
le défaut en capacité résiduelle et \HLoc transporte cette capacité vers les arêtes.
Ainsi, \(\delta_*\) est une conséquence de \HTr--\HAdd--\HDef, \(\lambda_*\) l’est de
\HTr--\HAdd--\HDef--\HRes et \(\mathcal A_*\) l’est de
\HTr--\HAdd--\HDef--\HRes--\HLoc. \HTr et \HAdd possèdent des caractérisations universelles ; \HDef, \HRes et
\HLoc sont des identifications constitutives internes ; elles ne se déduisent pas de la
transition finie précédente.

\section{Familles compatibles avec les symétries partielles}

\subsection{Symétrie du cycle \(1/2\) et famille compatible de mesures}

L’échange des deux moitiés d’une émission décompose les \(468\)
émissions en \(234\) paires. La mesure uniforme du catalogue comme les
mesures images des germes sont invariantes sous cet échange. Il en va de même
pour toute attribution constante sur chaque paire. L’invariance du retour,
à elle seule, ne prouve pas l’uniformité.

\subsection{Symétrie de \texorpdfstring{\Gtr}{Gamma tr} et famille de coûts}

Les \(7776\) arêtes de \Gtr se répartissent en \(288\) orbites de taille
\(27\) sous \((\mathbb Z/3\mathbb Z)^3\). Les \(2160\) arêtes de
changement de feuillet forment \(80\) orbites :
\[
 80=5\cdot16.
\]
La symétrie permet d’attribuer des coûts différents à ces orbites. Le coût
scalaire commun provient de \HAdd, et non de la cardinalité.

\subsection{Liberté scalaire de la coordonnée \(2^{\mathrm a}\)}

Pour tout \(c\ge0\), l’opérateur
\[
D_c=\Pi_\pi^{(468)}\otimes I_{11}
+cI_{468}\otimes P_3
\]
est positif et le caractère unique de \HAdd donne
\[
\tau(D_c)=\frac5{468}+\frac3{11}c.
\]
La formule interne de \(\Delta_4\) construit un candidat distingué, mais
la propriété universelle de la trace ne démontre pas à elle seule
\(c=\Delta_4\). Cette égalité est précisément le contenu de \HDef.

\subsection{Positivité, linéarité et termes d’ordre \(2\)}

Les fonctions
\[
 1-u-v+cuv,\qquad(1-u)(1-v),\qquad e^{-(u+v)}
\]
peuvent partager la normalisation et le premier ordre. L’absence de \(uv\) se
déduit de l’additivité exacte. Sans \HAdd, la loi linéaire ne serait pas unique.

\section{Hypothèses d’existence, de positivité, de naturalité et d’extension dynamique}

Le résultat dépend des conditions vérifiables suivantes :
\begin{enumerate}
  \item le catalogue ne contient pas \(468\) émissions ou la microfibre n’a pas
  rang cinq ;
  \item le recensement des fibres n’est pas
  \(432\times144,\ 18\times432,\ 18\times1944\) ;
  \item \(J_F\) n’est pas isométrique ou n’entrelace pas les observables du quotient ;
  \item l’un des \(P_G,P_3,Q_3\) cesse d’être un projecteur de rang trois ;
  \item il existe un autre caractère positif, normalisé bifactoriellement, additif et
  invariant qui n’est pas la somme des traces ;
  \item une transition totale du noyau topologique de phase enrichi, une fois construite, induit un
  coefficient de commutation différent de \(\Delta_4\) ;
  \item \(D_{\rm switch}\) acquiert une valeur propre négative ;
  \item l’action accumulée cesse d’être additive sous concaténation.
\end{enumerate}

Les certificats reconstruisent les fibres, les deux mesures, le facteur polaire,
le spectre du défaut, les orbites d’arêtes et les ablations. La valeur de
\(\lambda_*\) apparaît à la fin comme conséquence de ces objets.
\fi

%
\endgroup
% Distribución causal exacta del delta Ley 9; contenido conservado sin síntesis.
\section{Caractère régional de degré \texorpdfstring{\(5\)}{5} et mantisse réduite d'action}
\label{sec:l9s102:eta-ret}

Le bloc central
\begin{equation}
 \mathcal B_c=\begin{pmatrix}7&2\\2&7\end{pmatrix}
 \label{eq:l9s102:bc-accion}
\end{equation}
possède les valeurs propres \(9\) et \(5\). La longueur \(4\) des orbites de pas
\(3\), le rang \(12\) du cycle dodécaphasique et
\(104976/1944=54\) fixent le caractère
\begin{equation}
 \boxed{
 H_5(\alpha,\varphi)
 =\frac12\sqrt{\frac{1000\alpha}{\varphi}}-\alpha
 +\frac9{16}\alpha^2-\frac59\alpha^3
 +\frac7{48}\alpha^4-\frac1{54}\alpha^5.}
 \label{eq:l9s102:h5}
\end{equation}
La complétion nonadique construit la coordonnée analytique et le déterminant
régional
\begin{equation}
 x_A=\frac{50\pi}{9}\alpha_{\mathrm{HMT}},
 \qquad
 D_A=e^{-2x_A}
 =\exp\!\left(-\frac{100\pi}{9}\alpha_{\mathrm{HMT}}\right).
 \label{eq:l9s102:xa-da}
\end{equation}
La microfibre apporte \(\rho_\pi=169\), et sa série absolument
convergente est
\begin{equation}
 \mathscr C_\pi(\alpha)
 =169\alpha^6+\frac{D_A\alpha^7}{1-D_A\alpha}.
 \label{eq:l9s102:cpi}
\end{equation}
La mantisse réduite postérieure au retour est alors construite par
\begin{equation}
 \boxed{
 \eta_{\mathrm{ret}}
 =H_5-\frac{90}{\pi}\mathscr C_\pi(\alpha_{\mathrm{HMT}})
 \in\mathbb R_{>0}.}
 \label{eq:l9s102:eta-ret}
\end{equation}
Son membre de droite contient exclusivement des objets engendrés auparavant
par APP--TRIT--TPK et par la clôture de \(\alpha_{\mathrm{HMT}}\).



\HMTFlushCurrentChapter
\chapter{Dérivation autoduale de la constante de Planck et de sa droite d'action : périodicité temporelle de \(108\) phases, fermeture de Compton et covariance bidirectionnelle}%
\begingroup
\renewcommand{\chapter}[2][]{}%
\chapter{Dérivation autoduale de la constante de Planck : périodicité temporelle de \(108\) phases, action et longueur de Compton}
\label{chap:p3-final:36:planck_autodual}
\label{chap:planck-autodual}

\section{Réalisation de Compton associée à la périodicité temporelle de \(108\) phases}

Sur le cycle de \cref{eq:realizacion-circular-ct108}, nous réutilisons la section
positive d’action \(\hbar\) déjà engendrée par la branche déterminantale d’action.
Le quantum angulaire associé est
\begin{equation}
E_{108}=\hbar\omega_{108},
\qquad
m_{108}=\frac{E_{108}}{c^2}
=\frac{\hbar\omega_{108}}{c^2}.
\label{eq:cuanto-ct108}
\end{equation}
\(G\) n’a pas encore été utilisé. La géométrie du cycle et la quantification
angulaire suffisent à fixer ses deux lecteurs de Compton.

\begin{theorem}[Identité entre la périodicité temporelle de \(108\) phases et la longueur de Compton]
\label{thm:ct108-compton}
Pour le quantum \eqref{eq:cuanto-ct108},
\begin{equation}
\lambdabarC(m_{108})=\frac{\hbar}{m_{108}c}=R_{108},
\qquad
\lambdaC(m_{108})=\frac{h}{m_{108}c}=C_{108}.
\label{eq:identidad-ct108-compton}
\end{equation}
Par conséquent, le rayon du cycle est sa longueur de Compton réduite et la
longueur d’un tour est sa longueur de Compton complète.
\end{theorem}

\begin{proof}
En substituant \(m_{108}=\hbar\omega_{108}/c^2\),
\[
\frac{\hbar}{m_{108}c}
=\frac{\hbar c^2}{\hbar\omega_{108}c}
=\frac{c}{\omega_{108}}=R_{108}.
\]
Comme \(h=2\pi\hbar\), la deuxième égalité est \(2\pi R_{108}=C_{108}\).
\end{proof}

La coïncidence n’est pas un ajustement numérique. C’est une identité homogène : elle vaut
pour tout \(\hbar,c,t_0>0\) et dépend seulement du fait que la même période détermine
la fréquence, la longueur du cycle et le quantum d’action.

\section{Transducteur gravitationnel et équation de sélection autoduale}

La droite gravitationnelle HMT--MD contient le transducteur
\begin{equation}
G(\theta)=\theta\frac{c^5t_0^2}{\hbar},
\qquad \theta>0.
\label{eq:transductor-g-planck-anexo}
\end{equation}
L’homogénéité est exacte, mais ne sélectionne pas à elle seule \(\theta\). Le
cycle CT108 apporte une condition interne supplémentaire : le lecteur gravitationnel
doit renvoyer la même échelle radiale que le lecteur géométrique et le lecteur de
Compton.

Nous définissons le rayon gravitationnel réduit
\begin{equation}
\rg(\theta):=\frac{G(\theta)m_{108}}{c^2}.
\label{eq:radio-gravitatorio-reducido}
\end{equation}
En substituant \eqref{eq:cuanto-ct108} et
\eqref{eq:transductor-g-planck-anexo},
\begin{equation}
\rg(\theta)
=\theta\frac{\pi}{54}\ell_0.
\label{eq:rg-theta-lineal}
\end{equation}

\begin{theorem}[Sélecteur autodual de la périodicité temporelle de \(108\) phases]
\label{thm:selector-autodual-ct108}
L’équation
\[
\rg(\theta)=R_{108}
\]
possède une unique solution positive :
\begin{equation}
\boxed{\thetaStar=\left(\frac{54}{\pi}\right)^2.}
\label{eq:theta-star}
\end{equation}
\end{theorem}

\begin{proof}
Par \eqref{eq:realizacion-circular-ct108} et
\eqref{eq:rg-theta-lineal}, l’équation est
\[
\theta\frac{\pi}{54}\ell_0
=\frac{54}{\pi}\ell_0.
\]
Comme \(\ell_0>0\), on simplifie et l’on obtient
\(\theta=(54/\pi)^2\). La fonction
\(\theta\mapsto\rg(\theta)\) est strictement croissante sur
\(\RR_{>0}\), de sorte qu’il n’existe aucune autre solution positive.
\end{proof}

\begin{corollary}[Unité de Planck interne]
\label{cor:unidad-planck-interna}
Avec \(G_*=G(\thetaStar)\), les sections de Planck satisfont
\begin{equation}
\ellP:=\sqrt{\frac{\hbar G_*}{c^3}}=R_{108},
\qquad
\tP:=\sqrt{\frac{\hbar G_*}{c^5}}=\frac1{\omega_{108}},
\label{eq:planck-long-time-ct108}
\end{equation}
\begin{equation}
\mP:=\sqrt{\frac{\hbar c}{G_*}}=m_{108},
\qquad
\EP:=\mP c^2=E_{108}.
\label{eq:planck-mass-energy-ct108}
\end{equation}
En particulier,
\begin{equation}
\boxed{
R_{108}=\lambdabarC=\rg=\ellP,
\qquad
C_{108}=\lambdaC=2\pi\rg=2\pi\ellP.}
\label{eq:tres-lectores-planck}
\end{equation}
\end{corollary}

\begin{proof}
On substitue \(G_*=(54/\pi)^2c^5t_0^2/\hbar\). Alors
\[
\ellP=\frac{54}{\pi}ct_0=R_{108},
\qquad
\tP=\frac{54}{\pi}t_0=\frac1{\omega_{108}},
\]
et
\[
\mP=\frac{\pi}{54}\frac{\hbar}{c^2t_0}
=\frac{\hbar\omega_{108}}{c^2}=m_{108}.
\]
Le reste découle du \cref{thm:ct108-compton}.
\end{proof}

\begin{remark}[Trois lecteurs, une section]
La géométrie circulaire publie \(R_{108}\) et \(C_{108}\) ; la quantification de
Compton publie \(\lambdabarC\) et \(\lambdaC\) ; la réalisation gravitationnelle
publie \(\rg\) et \(2\pi\rg\). En \(\thetaStar\), ces trois paires sont les
coordonnées d’une unique section. C’est la raison précise pour laquelle
l’unité de Planck cesse d’apparaître comme une combinaison dimensionnelle accidentelle
dans cette formalisation.
\end{remark}

\section{Équivalence des normalisations sur l’identité autoduale}

Le transducteur \eqref{eq:transductor-g-planck-anexo} admet deux manières de
fixer l’échelle élémentaire, qui ne doivent pas être confondues.

\begin{enumerate}[label=\textup{(\roman*)}]
\item Si l’on fixe \(\theta=1\), alors
\(\sqrt{\hbar G/c^3}=\ell_0\). La transition élémentaire \(ct_0\) est identifiée
à la longueur de Planck de cette carte.
\item Si l’on exige que l’\emph{orbite complète} CT108 ait le rayon de Planck,
la condition sélectionne \(\thetaStar=(54/\pi)^2\) et
\(R_{108}=\ellP\).
\end{enumerate}

Ce ne sont pas des valeurs rivales de \(G\). Ce sont deux jauges de la même équation
homogène : l’une attribue l’unité de Planck au pas élémentaire ; l’autre, au rayon
du cycle complet. L’annexe utilise la seconde parce que sa thèse concerne
l’identité de la clôture CT108.

\section{Conventions gravitationnelles : rayon réduit et rayon de Schwarzschild}

La notation gravitationnelle est fixée sans équivoque :
\[
\rg=\frac{Gm}{c^2},
\qquad
\rs=\frac{2Gm}{c^2}=2\rg.
\]
Le théorème précédent emploie \(\rg\), non \(\rs\). Dans la convention habituelle de
la masse de Planck, l’égalité \(\hbar/(mc)=Gm/c^2\) se produit en \(m=\mP\),
tandis que
\begin{equation}
\frac{\hbar}{mc}=\frac{2Gm}{c^2}
\quad\Longleftrightarrow\quad
m=\frac{\mP}{\sqrt2}.
\label{eq:cruce-schwarzschild-factor-dos}
\end{equation}
C’est pourquoi l’expression « échelle gravitationnelle » désigne dans la thèse directrice le rayon
réduit. Le rayon d’horizon conventionnel demeure un lecteur distinct.

\section{Dédoublement bidirectionnel de la droite d’action}

La droite interne d’action possède deux sections orientées :
\begin{equation}
\hbarpre=H_5\,10^{-34}\mathcal U_S,
\qquad
\hbarret=\eta_{\mathrm{ret}}\,10^{-34}\mathcal U_S,
\qquad
\Ract=\frac{H_5}{\eta_{\mathrm{ret}}}.
\label{eq:hbar-pre-ret-anexo}
\end{equation}
La décade \(10^{-34}\) fixe l’échelle absolue ; le quotient \(\Ract\)
transporte la différence orientée. Cette séparation permet de demander si
l’autodualité de Planck subsiste dans les deux coordonnées.

\begin{theorem}[Covariance réciproque de l’action et de la gravité]
\label{thm:accion-gravedad-reciproca}
Pour \(a\in\{\mathrm{pre},\mathrm{ret}\}\), soit
\[
G_a(\theta)=\theta\frac{c^5t_0^2}{\hbar_a},
\qquad
m_{108,a}=\frac{\hbar_a\omega_{108}}{c^2}.
\]
Alors
\begin{equation}
\frac{G_a(\theta)m_{108,a}}{c^2}
=\theta\frac{\pi}{54}\ell_0
\label{eq:invariante-pre-ret-planck}
\end{equation}
est indépendant de \(a\). De plus,
\begin{equation}
\frac{m_{108,\mathrm{pre}}}{m_{108,\mathrm{ret}}}=\Ract,
\qquad
\frac{G_{\mathrm{pre}}}{G_{\mathrm{ret}}}=\Ract^{-1},
\label{eq:reciprocidad-pre-ret}
\end{equation}
et la longueur de Planck
\(\sqrt{\hbar_aG_a/c^3}=\sqrt\theta\,\ell_0\) est la même dans les deux
orientations.
\end{theorem}

\begin{proof}
Les facteurs \(\hbar_a\) s’annulent dans le produit \(G_am_{108,a}\).
Les deux rapports de \eqref{eq:reciprocidad-pre-ret} sont obtenus directement
à partir des définitions. L’invariance de la longueur de Planck découle de
\(\hbar_aG_a=\theta c^5t_0^2\).
\end{proof}

Une paramétrisation symétrique rend la géométrie visible :
\begin{equation}
\hbar_{\mathrm{geom}}:=\sqrt{\hbarpre\hbarret},
\qquad
\xi:=\frac12\log\Ract,
\qquad
\hbarpre=\hbar_{\mathrm{geom}}e^\xi,
\quad
\hbarret=\hbar_{\mathrm{geom}}e^{-\xi}.
\label{eq:rapidez-accion-pre-ret}
\end{equation}
La mémoire bidirectionnelle réside dans \(\xi\) ; la géométrie de Planck
réside dans le produit invariant. Cela permet une différence réelle entre les
deux coordonnées d’action sans exiger une différence de vitesse de
propagation.

\begin{remark}[Portée physique du théorème]
La réciprocité de \cref{thm:accion-gravedad-reciproca} est exacte au sein du
transducteur. Pour identifier \(G_{\mathrm{pre}}\) et \(G_{\mathrm{ret}}\) à deux
mesures gravitationnelles orientées, il faut spécifier le protocole,
l’observable et la symétrie qui les relie. Le théorème prouve la
compatibilité structurelle ; il n’anticipe pas à lui seul une violation de Lorentz
ni une variation observable de \(G\).
\end{remark}

\label{mdg:sec:identidad-autodual-planck-ct108-rev3}
\index{Planck!identité autoduale}
\index{calendrier d’ordre 108!réalisation circulaire}
\index{Compton!lectures réduite et complète}

La droite d’action construite dans la section précédente admet une
réalisation circulaire compatible avec la périodicité temporelle d’ordre 108.
Cette réalisation coordonne trois objets qui possèdent déjà leurs types propres : la
section d’action, la géométrie d’un tour et la transduction
gravitationnelle. La coordination qui en résulte fournit une identité de
Planck interne et prépare la droite positive sur laquelle agira ensuite
l’opérateur général de masse.

\subsection{Cartes orientées d’action et réalisation circulaire déclarée}

Soient
\[
 \mathcal L_S^+,\quad \mathcal L_T^+,\quad \mathcal L_C^+,
 \quad \mathcal L_L^+,\quad \mathcal L_M^+,\quad \mathcal L_G^+
\]
les demi-droites positives d’action, de temps, de vitesse, de longueur et de masse.
Les coordonnées antérieure et postérieure au retour s’écrivent
\begin{equation}
 \hbar_{\rm pre}=H_5\,10^{-34}\mathcal U_S,
 \qquad
 \hbar_{\rm ret}=\eta_{\rm ret}\,10^{-34}\mathcal U_S,
 \qquad
 R_{\rm act}=\frac{\hbar_{\rm pre}}{\hbar_{\rm ret}}
 =\frac{H_5}{\eta_{\rm ret}}.
 \label{mdg:eq:planck-ct108-cartas-pre-ret-rev3}
\end{equation}
L’indice \(a\in\{\mathrm{pre},\mathrm{ret}\}\) désignera l’une quelconque
des deux cartes. Dans chacune, la distinction exacte est conservée
\begin{equation}
 h_a=2\pi\hbar_a .
 \label{mdg:eq:planck-ct108-h-hbar-rev3}
\end{equation}

Fixons une durée élémentaire \(t_0\in\mathcal L_T^+\), une vitesse
interne \(c_{\rm int}\in\mathcal L_C^+\) et
\(\ell_0=c_{\rm int}t_0\). La réalisation circulaire déclarée de la
périodicité temporelle d’ordre 108 est déterminée par
\begin{equation}
 T_{108}=108t_0,
 \qquad
 \omega_{108}=\frac{2\pi}{T_{108}},
 \qquad
 R_{108}=\frac{c_{\rm int}}{\omega_{108}}
 =\frac{54}{\pi}\ell_0,
 \label{mdg:eq:planck-ct108-realizacion-circular-rev3}
\end{equation}
et par la longueur du tour complet
\begin{equation}
 C_{108}=2\pi R_{108}=c_{\rm int}T_{108}=108\ell_0.
 \label{mdg:eq:planck-ct108-perimetro-rev3}
\end{equation}
Ainsi, \(R_{108}\) est le rayon de la réalisation et \(C_{108}\) son
périmètre. La fréquence angulaire, la fréquence cyclique et les deux
constantes d’action sont reliées par
\(\omega_{108}=2\pi/T_{108}\) et \(h_a=2\pi\hbar_a\).

\subsection{Correspondance entre la périodicité d’ordre 108 et les lectures de Compton}

Dans la carte \(a\), le quantum circulaire et sa section de masse sont
\begin{equation}
 E_{108,a}=\hbar_a\omega_{108},
 \qquad
 m_{108,a}=\frac{E_{108,a}}{c_{\rm int}^{2}}
 =\frac{\hbar_a\omega_{108}}{c_{\rm int}^{2}}.
 \label{mdg:eq:planck-ct108-cuanto-circular-rev3}
\end{equation}

\begin{theorem}[Identité entre la périodicité circulaire et les lectures de Compton]
\label{mdg:thm:planck-ct108-compton-rev3}
Pour chaque carte \(a\in\{\mathrm{pre},\mathrm{ret}\}\), les lectures de
Compton réduite et complète satisfont
\begin{equation}
 \boxed{
 \bar\lambda_{\rm C}(m_{108,a})
 =\frac{\hbar_a}{m_{108,a}c_{\rm int}}=R_{108},
 \qquad
 \lambda_{\rm C}(m_{108,a})
 =\frac{h_a}{m_{108,a}c_{\rm int}}=C_{108}.}
 \label{mdg:eq:planck-ct108-identidad-compton-rev3}
\end{equation}
\end{theorem}

\begin{proof}
La substitution de
\eqref{mdg:eq:planck-ct108-cuanto-circular-rev3} produit
\[
 \frac{\hbar_a}{m_{108,a}c_{\rm int}}
 =\frac{c_{\rm int}}{\omega_{108}}=R_{108}.
\]
La deuxième égalité résulte de l’application simultanée de
\eqref{mdg:eq:planck-ct108-h-hbar-rev3} et de
\eqref{mdg:eq:planck-ct108-perimetro-rev3}.
\end{proof}

Le théorème est homogène en
\((\hbar_a,c_{\rm int},t_0)\) : la même périodicité détermine le rayon,
le périmètre et les deux cartes de Compton. La réalisation circulaire est la
primitive géométrique déclarée ; les égalités précédentes sont ses
conséquences exactes.

\subsection{Sélecteur circulaire et unité interne de Planck}

La droite gravitationnelle est coordonnée à la carte \(a\) par le transducteur
homogène
\begin{equation}
 G_a(\theta)=\theta\,
 \frac{c_{\rm int}^{5}t_0^2}{\hbar_a}
 \in\mathcal L_G^+,
 \qquad \theta\in\mathbb R_{>0}.
 \label{mdg:eq:planck-ct108-transductor-rev3}
\end{equation}
Son rayon gravitationnel réduit, évalué sur le quantum circulaire, est
\begin{equation}
 r_{{\rm g},a}(\theta)
 :=\frac{G_a(\theta)m_{108,a}}{c_{\rm int}^{2}}
 =\theta\frac{\pi}{54}\ell_0.
 \label{mdg:eq:planck-ct108-radio-gravitatorio-rev3}
\end{equation}

\begin{proposition}[Sélecteur circulaire de Planck]
\label{mdg:prop:selector-circular-planck-ct108-rev3}
Au sein de la réalisation circulaire
\eqref{mdg:eq:planck-ct108-realizacion-circular-rev3}, la condition
\(r_{{\rm g},a}(\theta)=R_{108}\) sélectionne une unique coordonnée
positive,
\begin{equation}
 \boxed{\theta_{108}^{\circ}=\left(\frac{54}{\pi}\right)^2.}
 \label{mdg:eq:theta-circular-planck-ct108-rev3}
\end{equation}
La coordonnée est indépendante de la carte orientée \(a\).
\end{proposition}

\begin{proof}
Les équations
\eqref{mdg:eq:planck-ct108-realizacion-circular-rev3} et
\eqref{mdg:eq:planck-ct108-radio-gravitatorio-rev3} réduisent la condition à
\[
 \theta\frac{\pi}{54}\ell_0=\frac{54}{\pi}\ell_0.
\]
La positivité de \(\ell_0\) et la stricte monotonie de l’application
\(\theta\mapsto r_{{\rm g},a}(\theta)\) fournissent la solution et son
unicité.
\end{proof}

Soit \(G_{108,a}^{\circ}:=G_a(\theta_{108}^{\circ})\). Les unités internes
associées à cette carte satisfont
\begin{align}
 \ell_{{\rm P},a}^{\circ}
 &:=\sqrt{\frac{\hbar_a G_{108,a}^{\circ}}{c_{\rm int}^{3}}}
 =R_{108},
 &
 t_{{\rm P},a}^{\circ}
 &:=\sqrt{\frac{\hbar_a G_{108,a}^{\circ}}{c_{\rm int}^{5}}}
 =\omega_{108}^{-1},
 \label{mdg:eq:planck-ct108-longitud-tiempo-rev3}\\
 m_{{\rm P},a}^{\circ}
 &:=\sqrt{\frac{\hbar_a c_{\rm int}}{G_{108,a}^{\circ}}}
 =m_{108,a},
 &
 E_{{\rm P},a}^{\circ}
 &:=m_{{\rm P},a}^{\circ}c_{\rm int}^{2}=E_{108,a}.
 \label{mdg:eq:planck-ct108-masa-energia-rev3}
\end{align}
L’égalité des rayons est donc publiée comme
\begin{equation}
 R_{108}=\bar\lambda_{\rm C}=r_{\rm g}=\ell_{\rm P}^{\circ},
 \qquad
 C_{108}=\lambda_{\rm C}=2\pi r_{\rm g}=2\pi\ell_{\rm P}^{\circ}.
 \label{mdg:eq:planck-ct108-tres-lectores-rev3}
\end{equation}

La convention de ce théorème est le rayon gravitationnel réduit
\(r_{\rm g}=Gm/c_{\rm int}^{2}\). Le rayon de Schwarzschild satisfait
\(r_{\rm s}=2r_{\rm g}\) ; si cette seconde convention est adoptée, son croisement avec
\(\bar\lambda_{\rm C}\) se produit en
\(m=m_{{\rm P},a}^{\circ}/\sqrt2\). La déclaration du rayon fait donc
partie de la carte et évite de transférer par inadvertance un facteur deux
entre les deux identités.

La coordonnée \(\theta_{108}^{\circ}\) appartient à la spécialisation
circulaire déclarée. Le sélecteur gravitationnel général, formulé sur
l’holonomie enrichie et sa réponse spectrale, conserve sa résidence dans
la \cref{mdg:sec:selector-gravitatorio-rev11}. Ainsi, l’identité
circulaire fixe une normalisation interne de Planck et le critère général
fixe les conditions que doit satisfaire la réalisation physique de la
section gravitationnelle.

\subsection{Covariance pré-retour/post-retour}

Le produit des sections réciproques d’action et de gravité conserve la
géométrie circulaire :
\begin{equation}
 \frac{m_{108,\rm pre}}{m_{108,\rm ret}}=R_{\rm act},
 \qquad
 \frac{G_{108,\rm pre}^{\circ}}{G_{108,\rm ret}^{\circ}}
 =R_{\rm act}^{-1},
 \qquad
 \hbar_a G_{108,a}^{\circ}
 =\theta_{108}^{\circ}c_{\rm int}^{5}t_0^2.
 \label{mdg:eq:planck-ct108-covariancia-pre-ret-rev3}
\end{equation}
En particulier, \(\ell_{{\rm P},\rm pre}^{\circ}\) et
\(\ell_{{\rm P},\rm ret}^{\circ}\) coïncident avec \(R_{108}\). La mémoire
bidirectionnelle réside dans le quotient des masses et dans son réciproque
gravitationnel ; l’échelle géométrique réside dans le produit invariant.

La réciprocité admet une paramétrisation symétrique qui sépare l’échelle
commune de l’orientation du retour. Définissons
\begin{equation}
 \begin{aligned}
 \hbar_{\rm geom}
 &:=\sqrt{\hbar_{\rm pre}\hbar_{\rm ret}},
 &
 G_{108,\rm geom}^{\circ}
 &:=\sqrt{G_{108,\rm pre}^{\circ}G_{108,\rm ret}^{\circ}},
 &
 \xi_{\rm act}&:=\frac12\log R_{\rm act}.
 \end{aligned}
 \label{mdg:eq:planck-ct108-parametrizacion-simetrica-rev3}
\end{equation}
Alors
\[
 \hbar_{\rm pre}=\hbar_{\rm geom}e^{\xi_{\rm act}},
 \qquad
 \hbar_{\rm ret}=\hbar_{\rm geom}e^{-\xi_{\rm act}},
\]
tandis que les sections gravitationnelles conjuguées satisfont
\[
 G_{108,\rm pre}^{\circ}
 =G_{108,\rm geom}^{\circ}e^{-\xi_{\rm act}},
 \qquad
 G_{108,\rm ret}^{\circ}
 =G_{108,\rm geom}^{\circ}e^{\xi_{\rm act}}.
\]
La coordonnée \(\xi_{\rm act}\) conserve l’orientation
pré-retour/post-retour ; le produit \(\hbar_aG_{108,a}^{\circ}\)
conserve l’échelle géométrique. Cette paramétrisation des cartes d’action
n’introduit pas à elle seule des vitesses de propagation distinctes.

\subsection{Droite autoduale déployée des masses}

Une carte \(a\) étant fixée, soit l’algèbre déployée
\[
 \mathbb D_{\rm sp}:=\mathbb R[j]/(j^2-1),
 \qquad
 P_\pm:=\frac{I\pm j}{2},
\]
et considérons sa décomposition spectrale
\[
 \mathbb D_{\rm sp}
 =\mathbb R P_+\oplus\mathbb R P_-,
 \qquad
 P_\pm^2=P_\pm,\quad
 P_+P_-=0,\quad
 P_++P_-=I.
\]
La conjugaison déployée échange les projecteurs,
\[
 \overline{uP_++vP_-}=vP_++uP_-,
\]
et sa norme est
\[
 N_{\rm sp}(uP_++vP_-):=uv.
\]

Pour \(m\in\mathcal L_M^+\), définissons la section
\begin{equation}
 Z_a(m)
 :=\frac{m_{{\rm P},a}^{\circ}}{m}P_+
   +\frac{m}{m_{{\rm P},a}^{\circ}}P_- .
 \label{mdg:eq:planck-ct108-seccion-partida-masa-rev3}
\end{equation}
L’unité \(m_{{\rm P},a}^{\circ}\) détermine également l’involution
\begin{equation}
 \iota_{{\rm P},a}(m)
 =\frac{(m_{{\rm P},a}^{\circ})^2}{m}.
 \label{mdg:eq:planck-ct108-involucion-masas-rev3}
\end{equation}

\begin{theorem}[Droite autoduale déployée des masses]
\label{mdg:thm:planck-ct108-recta-partida-masa-rev3}
Pour tout \(m\in\mathcal L_M^+\),
\begin{equation}
 \boxed{
 N_{\rm sp}(Z_a(m))=1,
 \qquad
 \overline{Z_a(m)}
 =Z_a\!\left(\iota_{{\rm P},a}(m)\right).}
 \label{mdg:eq:planck-ct108-norma-conjugacion-rev3}
\end{equation}
La conjugaison possède un unique point fixe positif,
\[
 m=m_{{\rm P},a}^{\circ},
 \qquad
 Z_a(m_{{\rm P},a}^{\circ})=I.
\]
\end{theorem}

\begin{proof}
La norme est le produit des deux coordonnées réciproques :
\[
 \frac{m_{{\rm P},a}^{\circ}}m
 \frac m{m_{{\rm P},a}^{\circ}}=1.
\]
La conjugaison échange les deux coordonnées, ce qui équivaut à remplacer
\(m\) par \((m_{{\rm P},a}^{\circ})^2/m\). Le point fixe positif satisfait
\(m^2=(m_{{\rm P},a}^{\circ})^2\).
\end{proof}

Avec la rapidité
\begin{equation}
 x_a(m)=\log\frac{m}{m_{{\rm P},a}^{\circ}},
 \label{mdg:eq:planck-ct108-rapidez-masa-rev3}
\end{equation}
la section et sa conjuguée s’écrivent
\[
 Z_a(m)=e^{-x_a}P_++e^{x_a}P_-,
 \qquad
 \overline{Z_a(m)}
 =e^{x_a}P_++e^{-x_a}P_-.
\]
Leurs lecteurs pair et impair sont
\begin{equation}
 Q_a(m):=\frac{e^{x_a}+e^{-x_a}}2=\cosh x_a,
 \qquad
 A_a(m):=\frac{e^{x_a}-e^{-x_a}}2=\sinh x_a.
 \label{mdg:eq:planck-ct108-lectores-par-impar-rev3}
\end{equation}
Le premier conserve la distance au point autodual et le second conserve son
orientation.

Les lecteurs physiques réciproques associés à la même coordonnée sont
\[
 r_{{\rm g},a}(m)=\frac{G_{108,a}^{\circ}m}{c_{\rm int}^{2}},
 \qquad
 \bar\lambda_{{\rm C},a}(m)=\frac{\hbar_a}{mc_{\rm int}},
\]
et satisfont
\begin{equation}
 \frac{r_{{\rm g},a}(m)}{\bar\lambda_{{\rm C},a}(m)}
 =e^{2x_a(m)},
 \qquad
 \frac{\bar\lambda_{{\rm C},a}(m)}{r_{{\rm g},a}(m)}
 =e^{-2x_a(m)}.
 \label{mdg:eq:planck-ct108-lectores-exponenciales-rev3}
\end{equation}
La conjugaison échange les deux lecteurs ; l’identité de Planck est leur
équilibre unitaire.

\begin{remark}[Identité de Planck et clôture électronique]
L’identité \(m=m_{{\rm P},a}^{\circ}\) fixe l’origine autoduale de la droite
positive des masses. La section électronique fixe une position différente sur
cette même droite par une généalogie propre. La
\cref{mdg:sec:ley-general-masa-accion-autonoma} construit l’opérateur universel
qui transporte l’échelle absolue, le caractère, les tours et la lecture bidirectionnelle ;
la \cref{mdg:sec:cierre-electron-two-way-rev11} réalise ensuite sa réduction
centrale de rang un. Ainsi, l’unité de Planck fournit l’identité de
la coordonnée et l’électron fournit la première transition matérielle
complète sur celle-ci.
\end{remark}

\begin{proposition}[Réalisation généalogique du centre électronique]
\label{mdg:prop:planck-ct108-realizacion-centro-electron-rev3}
Soit
\[
 F_e^+=\mathbb R_{>0}(P_++P_-)
\]
la demi-droite positive de la fibre centrale de rang un, soit
\(\mathcal G_e\) l’espace des généalogies
électroniques enrichies et soit \(\mathcal L_M^+\) la droite positive des
masses. Un morphisme de réalisation
\[
 \mathfrak R_e:
 F_e^+\times\mathcal G_e\longrightarrow\mathcal L_M^+
\]
peut envoyer l’identité \(I=P_++P_-\) sur une masse différente de
\(m_{{\rm P},a}^{\circ}\) sans modifier le
\cref{mdg:thm:planck-ct108-recta-partida-masa-rev3} : l’identité de Planck fixe
l’origine autoduale de la carte, tandis que la généalogie électronique fixe un
déplacement sur cette carte.
\end{proposition}

\begin{proof}
Si \(\chi_e:\mathcal G_e\to\mathbb R_{>0}\) est le caractère positif construit
par le registre électronique, la règle
\[
 \mathfrak R_e(\lambda I,g)
 =\lambda\,m_{{\rm P},a}^{\circ}\chi_e(g)
\]
est positive et multiplicative dans la coordonnée généalogique. Pour la généalogie
neutre, \(\chi_e=1\), on retrouve l’origine de Planck ; pour une généalogie avec
\(\chi_e(g)\ne1\), la même identité interne occupe une position différente sur
la droite. La clôture électronique concrète conserve son développement antécédent dans la
\cref{mdg:sec:cierre-electron-two-way-rev11}.
\end{proof}

\section{Autodualité de l'unité interne de Planck sous inversion d'orientation}

Le sélecteur CT108 construit

\begin{equation}
\theta_{108}=\left(\frac{54}{\pi}\right)^2,
\qquad
\ell_P=\frac{54}{\pi}\ell_0,
\qquad
t_P=\frac{54}{\pi}t_0,
\label{eq:planck-length-time}
\end{equation}

et

\begin{equation}
E_P=\frac{\hbar_{\rm int}}{t_P}
=\frac{\pi\hbar_{\rm int}}{54t_0}.
\label{eq:planck-energy}
\end{equation}

En comparant \cref{eq:boltzmann-clock,eq:planck-energy}, on obtient

\begin{equation}
\boxed{E_P=k_B\Theta_{\rm clk}\log3.}
\label{eq:planck-trit}
\end{equation}

La relation ne dépend pas d'une coïncidence entre unités : les deux membres
proviennent de la même structure chronologique \(108\).

Une fois la constante gravitationnelle dérivée en interne dans le
\cref{ch:gravity}, la famille de Planck satisfait

\begin{align}
\ell_P^2&=\frac{G\hbar}{c^3}=\theta_{108}\ell_0^2,
\label{eq:planck-area}\\
F_P&=\frac{c^4}{G},
\qquad
P_P=\frac{c^5}{G},
\label{eq:planck-force-power}\\
\rho_P&=\frac{\hbar}{\theta_{108}^2c^5t_0^4}.
\label{eq:planck-density}
\end{align}

Bien que \(G\) et \(\hbar\) se transforment réciproquement entre les sections
antérieure et postérieure au retour, leur produit dans \(\ell_P^2\) conserve la
géométrie. Cette invariance détermine l'aire comme interface entre les
réalisations géométrique, informationnelle et d'intrication.

%
\endgroup
% Distribución causal exacta del delta Ley 9; contenido conservado sin síntesis.
\section{Droite d'action et constante de Planck}
\label{sec:l9s102:recta-planck}

Soit \(\mathcal L_S\) la droite réelle d'action de base positive
\(\mathcal U_S\). La transduction
\begin{equation}
 \tau_S:\mathbb R_{>0}\longrightarrow\mathcal L_S,
 \qquad u\longmapsto u\,10^{\varepsilon_H}\mathcal U_S
 \label{eq:l9s102:transductor-accion}
\end{equation}
publie les sections antérieure et postérieure au retour :
\begin{equation}
 \hbar_{\mathrm{pre}}=H_5\,10^{-34}\mathcal U_S,
 \qquad
 \boxed{
 \hbar_{\mathrm{ret}}=\eta_{\mathrm{ret}}\,10^{-34}\mathcal U_S.}
 \label{eq:l9s102:hbar-pre-ret}
\end{equation}
Les actions de cycle complet sont
\begin{equation}
 h_{\mathrm{pre}}=2\pi\hbar_{\mathrm{pre}},
 \qquad
 h_{\mathrm{ret}}=2\pi\hbar_{\mathrm{ret}}.
 \label{eq:l9s102:h-pre-ret}
\end{equation}
Pour une phase \(\theta\), les cartes angulaire et d'action commutent :
\begin{equation}
 S(\theta)
 =\hbar_{\mathrm{ret}}\theta_{\mathrm{rad}}
 =\hbar_{\mathrm{ret}}\frac{\pi}{180}\theta_{\deg}
 =h_{\mathrm{ret}}\frac{\theta_{\deg}}{360}.
 \label{eq:l9s102:accion-fase}
\end{equation}



% Transición visible de precedencia causal. Reúne propietarios ya activos sin
% repetir sus pruebas y mantiene tipadas las ramas dimensional y compleja.
% Reunión expositiva de precedencia causal entre resultados ya establecidos.
% Los propietarios íntegros continúan gobernando cada flecha.
% Propietarios reunidos por el grafo activo:
% - sucesor_102/deltas_ley9/c27_teorema_generacion_coinductiva_profundidad_arbitraria.tex
% - sucesor_102/deltas_ley9/c34_diagrama_bifurcado_y_decada_accion.tex
% - sucesor_102/deltas_ley9/c36_recta_accion_planck.tex
% - colaboracion/parte_iii/tex/capitulos_finales/32_raiz_menos_uno.tex
% - incorporaciones/parte_iv_ampliaciones_20260824/sources/editor_ready/03_pareja_planck_respuesta.tex
% - integracion_83/parte_iv/chapters/48_electron_primera_realizacion.tex

\section{Généalogie opératoire de l'action, de la géométrie angulaire et de la masse}
\label{sec:l9s102:transicion-constantes-accion-geometria-masas}

La dynamique coinductive APP--TRIT--TPK se prolonge à profondeur arbitraire
sur un même état enrichi et sur la structure discrète conjointe du
continu. À chaque prolongation, elle conserve valeur, généalogie, phase, retenue,
orientation, vacance, chemin, frontière, mémoire et survie. Ses
publications distinguées sont d'abord les trois coordonnées internes et,
par le sceau dodécaphasique, la coordonnée \(\alpha_{\mathrm{HMT}}\) :
\[
 \mathrm{APP}\longrightarrow\mathrm{TRIT}\longrightarrow\mathrm{TPK}
 \longrightarrow\mathcal S_{\infty}^{\mathrm{enr}}
 \longrightarrow\mathcal C_{\mathrm{cont}}^{\mathrm{disc}}
 \longrightarrow(\pi_{\mathrm{HMT}},\varphi_{\mathrm{HMT}},e_{\mathrm{HMT}})
 \longrightarrow\alpha_{\mathrm{HMT}}.
\]
La dernière flèche représente la publication signée de
\(\alpha_{\mathrm{HMT}}\) au sein de cette même dynamique ; elle ne reçoit pas de valeurs
conventionnelles et ne transforme pas les \(\pi,\varphi,e\) conventionnels en générateurs.

Ce n'est qu'ensuite que l'on applique le lecteur déterminantal d'action. Celui-ci publie la
valuation \(10^{-34}\mathcal U_S\), les deux sections
\(\hbar_{\mathrm{pre}},\hbar_{\mathrm{ret}}\) et leurs actions de cycle
complet
\[
 h_{\mathrm{pre}}=2\pi_{\mathrm{HMT}}\hbar_{\mathrm{pre}},
 \qquad
 h_{\mathrm{ret}}=2\pi_{\mathrm{HMT}}\hbar_{\mathrm{ret}}.
\]
À partir de ce point causal, deux branches restent séparées et ne se
sélectionnent pas mutuellement. La branche dimensionnelle--gravitationnelle prolonge l'action et
la section causale jusqu'aux formes déjà démontrées du couple de Planck,
\[
 \ell^{2}=G\hbar c^{-3},
 \qquad
 \frac{\ell}{Q}=Gc^{-4};
\]
la branche complexe publie, à partir de l'opérateur interne de rotation,
\[
 u^2+1=0,
 \qquad u\longmapsto J_{+1},
 \qquad J_{+1}^{2}=-I.
\]
Telle est la réalisation opératoire interne de \(i\), non l'introduction d'un
symbole complexe sans antécédent. Les branches dimensionnelle--gravitationnelle et
complexe sont parallèles sur l'état enrichi commun : la première ne
transforme pas \(i\) en entrée de l'échelle d'action et, réciproquement,
\(h\), \(\hbar\), \(c^{-3}\) et \(c^{-4}\) ne génèrent pas \(i\). Les identités
du couple de Planck restent ainsi séparées de la réalisation complexe.

En écrivant \(\prec_{\mathrm{HMT}}\) pour la précédence causale avec conservation
du type ---et non pour affirmer que chaque objet de droite est fonction de
tous ceux de gauche---, on obtient
\[
\begin{aligned}
 (\pi,\varphi,e,\alpha)_{\mathrm{HMT}}
 &\prec_{\mathrm{HMT}}(10^{-34}\mathcal U_S,\hbar,h)\\
 &\prec_{\mathrm{HMT}}
 \bigl\{G\hbar c^{-3},Gc^{-4};\,J_{+1}^{2}=-I, i\bigr\}\\
 &\prec_{\mathrm{HMT}}(A_{\mathrm{deg}},C^{*}_{\mathrm{deg}})\\
 &\prec_{\mathrm{HMT}}(\text{électron},\text{loi des masses}).
\end{aligned}
\]
La construction angulaire conserve ses entrées effectives
\(\alpha_{\mathrm{HMT}}\) et \(\eta_{\mathrm{ret}}\) : les deux branches parallèles
précédentes fixent la disponibilité conjointe et le typage, mais ne sont pas promues
rétrospectivement au rang de sélecteurs de l'angle ou du contre-angle.

Dans le chemin électronique central, les lecteurs bidirectionnels conservent deux
étapes distinctes d'une seule généalogie :
\[
 \begin{aligned}
 E_e^{(0)}
   &=0.5109989224880660725\ldots\ \mathrm{MeV},
   &&\text{étape basale},\\
 E_e^{\beta}
   &=E_e^{(0)}R_{\mathrm{act}}^{\,80-54\alpha_{\mathrm{HMT}}+6\Delta_4}
     =0.510998950690000808608\ldots\ \mathrm{MeV},
   &&\text{sortie bêta directrice},
 \end{aligned}
\]
avec
\[
 K_D(\Gamma_e)=K_{\Omega}(\Gamma_e)=(80,54,6).
\]
La coïncidence du triplet dans la fibre centrale n'identifie ni les domaines ni
les opérations de \(K_D\) et \(K_{\Omega}\). Dans le registre canonique
historique, \texttt{MD-ELECTRON-001} désigne le candidat étalonné de l'étape
basale. Cette identification conserve sa fonction de provenance, mais ne gouverne pas
la sortie bêta et ne sélectionne pas le chemin ; le basal et la clôture bêta restent
des étapes typées des lecteurs bidirectionnels.

\paragraph{Critère de réfutation.}
La généalogie précédente est réfutée si une valeur conventionnelle ou la sortie
bêta est utilisée pour sélectionner un chemin, si l'étape basale est présentée comme
clôture directrice, si l'une des deux branches parallèles est identifiée sans une
application typée, ou si l'angle, le contre-angle, l'électron ou une masse sont placés
comme entrées du générateur de constantes ou d'action.


\HMTFlushCurrentChapter
\chapter{Opérateur de phase antérieur à la transduction dimensionnelle : décomposition en \(27\) secteurs, rayon spectral et factorisation de la double projection}%
\begingroup
\renewcommand{\chapter}[2][]{}%
\chapter{Opérateur angulaire prédimensionnel : décomposition en \(27\) secteurs, rayon spectral et composition effective de la double projection}
\label{chap:p3-final:37:operador_angular}
\label{p3-20260826:chap:operador-angular-dictamen}
\index{opérateur angulaire}
\index{phase angulaire!sortie spectrale}
\index{clôture!pré-dimensionnelle}

Les deux sections précédentes ont réuni trois éléments qui apparaissaient auparavant
séparés : une identification isométrique entre projecteurs, une mesure de
quotient et un cocycle positif d’action. La présente section étudie la conséquence
spectrale de cette loi et situe le résultat dans l’architecture complète
de HMT\@.

La séparation des statuts est précise. L’identification triple est un théorème
exact relatif à la nouvelle règle variationnelle interne et à un ancrage orienté.
La mesure et le caractère positif sont des théorèmes exacts de \HTr--\HAdd ; le défaut
est en outre relatif à \HDef, la capacité à \HRes et le coût local à \HLoc. La valeur
angulaire est une sortie certifiée par intervalles de l’opérateur \Gtr. Sa
proximité avec la coordonnée angulaire conjuguée canonique est extraordinaire, mais l’intervalle de
la différence exclut l’égalité.

\section{Décomposition de l’opérateur angulaire en \(27\) secteurs de \(72\) états}

Le lecteur de tous les chemins \Gtr possède \(1944\) états microscopiques.
L’action de translation grossière de \((\mathbb Z/3\mathbb Z)^3\) permet la
décomposition
\[
 1944=27\cdot72.
\]
Chaque bloc de \(72\) états est paramétré par
\[
 (r,s,m,d)
 \in(\mathbb Z/3\mathbb Z)^2\times\{0,1\}\times D_4,
\]
où \(m\) est le feuillet et \(d\) la direction précédente. De chaque état
partent quatre continuations cardinales. Le TRIT de la cellule d’arrivée
décide du feuillet suivant.

Avec l’action du \cref{p3-20260826:chap:medida-accion-positiva}, le poids d’une arête
est
\[
 \exp[-g(e)-\lambda_*s(e)].
\]
La transformée en caractères réduit l’opérateur à des matrices complexes
\(L_{k_a,k_b}\) de taille \(72\times72\). Le programme énumère les
arêtes de chaque bloc ; il ne reçoit pas une matrice spectrale précalculée.
L’itération des puissances est conservée à titre exploratoire, mais le résultat qui
suit ne repose plus sur elle : les rayons spectraux sont isolés au moyen d’une
certification rationnelle complète.

Soient \(\rho_{01}\) et \(\rho_{12}\) les rayons spectraux des secteurs
orientés \((0,1)\) et \((1,2)\). Nous définissons la \emph{phase de bande}
\[
 \boxed{y_{\rm band}=\log\rho_{12}-\log\rho_{01}.}
\]
Pour
\[
 \lambda_*
 =1-\frac5{468}-\frac3{11}\Delta_4,
\]
la certification donne
\[
\begin{aligned}
1.460655120862011111
&<\rho_{01}<1.460655120862404024,\\
1.515812008270944328
&<\rho_{12}<1.515812008271195416,
\end{aligned}
\]
et, par propagation par intervalles,
\[
 0.037066226434427529
 <y_{\rm band}<
 0.037066226434862172,
\]
\[
 \boxed{
 2.123738337168943^\circ
 <C_{\rm band}:=\frac{180}{\pi}y_{\rm band}
 <2.123738337193847^\circ.}
\]
La quantité primaire est \(y_{\rm band}\), un logarithme adimensionnel de quotient
spectral. Les degrés apparaissent ensuite lorsque l’on applique la carte
\(180/\pi\). L’opérateur produit d’abord une phase interne et seulement ensuite
une représentation angulaire archimédienne.

\subsection{Calcul exact du rayon spectral de l’opérateur angulaire}

Les blocs sont non normaux. Par conséquent, un petit résidu
\(\|Lv-\mu v\|\) ne suffit pas à démontrer que \(\mu\) est la valeur propre de
plus grand module. La preuve utilisée ici comporte quatre étapes.

On observe d’abord douze lignes structurellement nulles dans chaque bloc. En
développant le polynôme caractéristique suivant ces lignes, on obtient un facteur
\(z^{12}\) et un mineur principal d’ordre soixante. Dans ce dernier apparaissent huit
autres lignes structurellement nulles. Par conséquent,
\[
 \det(zI_{72}-L)=z^{20}\det(zI_{52}-A),
\]
et toute valeur propre non nulle est contenue, avec sa multiplicité, dans un bloc
réduit \(A\) d’ordre cinquante-deux.

Deuxièmement, le certificat fournit des matrices gaussiennes rationnelles
\(S,R\), de dénominateur \(10^{16}\). Le vérificateur calcule exactement
\(E=I-RS\) et obtient
\[
 \|E\|_\infty\le8.5664\times10^{-14}
 \quad\hbox{et}\quad
 \|E\|_\infty\le3.7628\times10^{-14}
\]
pour les secteurs \((0,1)\) et \((1,2)\), respectivement. La série de
Neumann démontre alors que \(S\) est inversible.

Troisièmement, si \(A_0\) est le point médian rationnel du bloc d’intervalles, l’identité
\[
 S^{-1}AS=(I-E)^{-1}RAS
\]
borne la distance à \(C_0=RA_0S\) par
\[
 1.25125\times10^{-13}
 \quad\hbox{et}\quad
 5.7037\times10^{-14}.
\]
Les disques de Gershgorin de \(C_0\), agrandis de ces bornes, contiennent
les disques de la similitude exacte. Dans chaque bloc, un disque est séparé
des cinquante et un autres ; le second théorème de Gershgorin lui attribue
exactement une valeur propre. Son module inférieur dépasse le module supérieur de
tous les autres disques. C’est la preuve que la valeur propre isolée est le
rayon spectral.

Quatrièmement, \(\pi,e,\log\pi,e^{-1},e^{-\lambda_*}\) et \(\sqrt3\) ne sont pas
acceptés comme nombres décimaux de bibliothèque. Le vérificateur les encadre au moyen de
la formule de Machin, des séries factorielle et logarithmique avec reste, de la
série alternée de l’exponentielle et d’un encadrement rationnel quadratique pour
\(\sqrt3\). La vérification emploie uniquement des entiers et des fractions ;
l’argument est clos par la réduction rationnelle, la série de Neumann et
la séparation de Gershgorin exposées dans cette partie. Les développements
décimaux reproductibles à une profondeur de dix mille sont publiés, comme témoins
finis indépendants, dans
\hyperref[app:desarrollos-10000-rev9]{l’annexe des développements}.

\section{Opérateur effectif de bande \texorpdfstring{\(T_{\rm band}\)}{Tband} et isolement de la phase angulaire}

Soit
\[
 x=\frac{50\pi}{9}\alpha_{\rm HMT}
 =0.127362829012869966\ldots
\]
et soit
\[
 R=\begin{pmatrix}1&0\\0&-1\end{pmatrix}.
\]
L’opérateur bicouche effectif associé exclusivement à la phase de bande est
\[
 \boxed{T_{\rm band}=\exp(-xI+y_{\rm band}R).}
\]
Dans la base des secteurs orientés,
\[
 T_{\rm band}
 =\begin{pmatrix}
 e^{-x+y_{\rm band}}&0\\
 0&e^{-x-y_{\rm band}}
 \end{pmatrix}.
\]
On vérifie
\[
 \det T_{\rm band}=e^{-2x},
 \qquad
 \frac{(T_{\rm band})_{11}}{(T_{\rm band})_{22}}
 =e^{2y_{\rm band}}.
\]
Par conséquent,
\[
 -\frac9{100\pi}\log\det T_{\rm band}=\alpha_{\rm HMT},
 \qquad
 \frac{90}{\pi}
 \log\frac{(T_{\rm band})_{11}}{(T_{\rm band})_{22}}
 =C_{\rm band}.
\]

Les deux récupérations ont des statuts distincts. La seconde restitue la
sortie produite par l’écart. La première relit le
\(\alpha_{\rm HMT}\) entré dans \(x\) ; elle démontre la réversibilité de la
carte, et non une seconde dérivation indépendante de \(\alpha\).

La forme exponentielle sépare deux fonctions. Le facteur commun \(e^{-x}\)
décrit la contraction partagée par les deux feuillets. Les facteurs
\(e^{\pm y_{\rm band}}\) distinguent les orientations en conservant le déterminant. Comme
\(I\) et \(R\) commutent, la paire
\[
 \left(\det T_{\rm band},
 \frac{(T_{\rm band})_{11}}{(T_{\rm band})_{22}}\right)
\]
reconstruit univoquement \((x,y_{\rm band})\).

\section{Comparaison ultérieure et échelle de précision}

La section suivante construit une autre phase, la \emph{phase régionale
régularisée} \(y_{\rm reg}\) ---également notée \(y_*\) dans cette section---. Pour
formuler ici une comparaison avec ce résultat ultérieur, nous anticipons la
notation
\[
 C_{\rm reg}:=\frac{180}{\pi}y_{\rm reg}
 =2.123738338968646^\circ.
\]
Cette ligne ne constitue ni une seconde définition ni une entrée de l’opérateur
spectral : elle cite la conclusion du
\cref{p3-20260826:thm:contraangulo-regional}, dont la construction est développée dans la
partie suivante.
Les deux paires d’objets ne sont pas identifiées :
\[
 (y_{\rm band},C_{\rm band})
 \ne
 (y_{\rm reg},C_{\rm reg}).
\]
La différence certifiée par intervalles de la loi linéaire est
\[
 -1.799703\times10^{-9}\ {}^\circ
 <C_{\rm band}-C_{\rm reg}
 <-1.774799\times10^{-9}\ {}^\circ,
\]
avec une différence relative encadrée par
\[
 -8.475\times10^{-10}
 <\frac{C_{\rm band}-C_{\rm reg}}{C_{\rm reg}}
 <-8.356\times10^{-10}.
\]
La loi qui produit \(C_{\rm band}\) ne reçoit ni \(C_{\rm reg}\), ni CKM, ni Barbero,
ni masses ni données SI\@. Cette indépendance à l’égard de la cible fait de la proximité un
candidat de reconnaissance fort et falsifiable. La certification spectrale
élimine l’ancienne lacune numérique ; la différence par intervalles, qui exclut
zéro, fixe désormais la limite exacte : \(C_{\rm band}\) est une sortie spectrale
prédimensionnelle distincte de la réalisation régionale
\(C_{\rm reg}\).

Le terme supplémentaire
\[
 -\frac{\Delta_4^2}{27}
\]
en \(\lambda\) produce
\[
 C_{\rm num}\approx2.12373833894104980{}^\circ,
\]
à \(2.76\times10^{-11}\) degrés de \(C_{\rm reg}\). L’entier \(27\) est interne,
car il compte les caractères de \((\mathbb Z/3\mathbb Z)^3\), mais ce
fait n’impose pas le coefficient du second moment. L’additivité exacte de
\HAdd exclut les termes quadratiques. La variante est conservée comme ablation et non
comme loi principale.

\section{Dépendance fonctionnelle du rayon spectral à l’égard des composantes de l’opérateur}

Toutes les lignes suivantes utilisent le même opérateur ; seule une entrée
de la loi change. Ce sont des diagnostics numériques d’ablation ; l’intervalle certifié
précédent correspond exclusivement à la loi complète
\HTr--\HAdd--\HDef--\HRes--\HLoc (et, en son sein, le défaut positif est fixé dans la
strate \HTr--\HAdd--\HDef).
\[
\begin{array}{l|c|c}
\text{règle}&\lambda&C_{\rm num}\text{ en degrés}\\ \hline
\text{autodual}&1&2.082755794784418\\
\text{microfibre seule}&1-5/468&2.123621809570142\\
\text{trace }5/468+(3/11)\Delta_4
&0.989285913019141&2.123738337181414\\
4/468\text{ au lieu de }5/468
&0.991422665155894&2.115535228138553\\
3/12\text{ au lieu de }3/11
&0.989288440210566&2.123728626433667\\
5/243\text{ au lieu de }5/468
&0.979393542015659&2.161907060464779\\
5/729\text{ au lieu de }5/468
&0.993110963140488&2.109064222037112\\
7/729\text{, mesure des germes}
&0.990367478915522&2.119584299852700.
\end{array}
\]
Les ablations démontrent que la sortie distingue la microfibre complète,
la dimension réduite onze et la mesure des émissions. Elles n’attribuent pas à elles
seules une probabilité au hasard : une telle probabilité exigerait de définir au préalable un
espace de modèles alternatifs et une mesure sur celui-ci.

\section{Composition typée des morphismes archimédien, exceptionnel et dimensionnel}

La construction complète peut s’écrire comme une succession d’applications
aux domaines visibles :
\[
\begin{aligned}
\Omega_{\rm seed}
&\xrightarrow{F}\mathcal U_6
\xrightarrow{J_F}H_U,\\
\mathscr R_\pi^{++}
&\xrightarrow{\beta_\pi}P_GH_G,\\
L_{\rm or}\otimes P_GH_G
&\xrightarrow{W_{GK}}P_3V_{11}
\xrightarrow{U_W}Q_3E_{30},\\
(\Pi_\pi^{(468)},\Delta_4P_3)
&\longmapsto D_{\rm switch}
\xrightarrow{\tau}\delta_*
\xrightarrow{\rm \HRes}\lambda_*,\\
\lambda_*
&\xrightarrow{\rm \HLoc}a_*
\xrightarrow{\mathcal L_{1,*}}y_{\rm band}
\xrightarrow{180/\pi}C_{\rm band},\\
(\alpha_{\rm HMT},y_{\rm band})
&\xrightarrow{\exp(-xI+y_{\rm band}R)}T_{\rm band}.
\end{aligned}
\]
Chaque flèche déclare ce qu’elle conserve :
\begin{itemize}
  \item \(F\) conserve la classe d’émission et laisse dans la fibre sa
  multiplicité de présentation ;
  \item \(J_F\) compense cette multiplicité dans la norme ;
  \item \(\beta_\pi\) conserve la position exceptionnelle ;
  \item \(W_{GK}\) conserve la représentation orientée ;
  \item \(U_W\) conserve l’incidence et la métrique ;
  \item la trace oublie les coordonnées et conserve les rangs ;
  \item l’opérateur de transfert somme les chemins avec le poids de l’action ;
  \item la carte \(180/\pi\) représente une phase interne en degrés.
\end{itemize}

\section{Diagramme des dépendances de l’opérateur angulaire et de ses publications}

\begin{longtable}{>{\raggedright\arraybackslash}p{.31\textwidth}>{\raggedright\arraybackslash}p{.22\textwidth}>{\raggedright\arraybackslash}p{.37\textwidth}}
\toprule
Objet&Nature mathématique&Contenido demostrado\\
\midrule
\endfirsthead
\toprule
Objet&Nature mathématique&Contenido demostrado\\
\midrule
\endhead
\(U_W=I/6:V_{11}\to E_{30}\)
&théorème exact
&isométrie, surjectivité, incidence équivariante et unicité positive\\
\(P_3\leftrightarrow Q_3\)
&théorème exact
&projecteurs conjugués, rang trois et trace \(3/11\)\\
étoile \(R_{B_K}\)
&théorème exact typé
&spectre \(1^{[7]},(-1)^{[4]}\), indice \(3/11\) et obstruction de conjugaison\\
\(P_G\) dans \(t=7\)
&théorème fini exact
&filtration \(11\to3\to1\), projecteur positif et action \(S_3\)\\
\(\mathscr R_\pi^{++}\leftrightarrow P_G\)
&théorème exact
&bijection \(S_3\)-équivariante par position exceptionnelle\\
identification isométrique orientée
&nouveau théorème relatif
&sélecteur unique par \(B_K,u_K\), Reynolds, polaire et réduction \(240\to1\)\\
mesure \(1/468\)
&théorème exact de \HTr
&trace normalisée unique et application polaire induite par les fibres\\
\(D_{\rm switch}\) y \(\delta_*\)
&théorème exact relatif à \HTr--\HAdd--\HDef
&positivité, spectre et trace \(5/468+(3/11)\Delta_4\)\\
\(\lambda_*\)
&théorème interne relatif à \HTr--\HAdd--\HDef--\HRes
&capacité résiduelle obtenue à partir du défaut par la règle \HRes\\
\(\mathcal A_*\)
&théorème interne relatif à \HTr--\HAdd--\HDef--\HRes--\HLoc
&cocycle additif de chemins de coût local \(g+\lambda_*s\)\\
\(T_{\rm band}\) y \(C_{\rm band}\)
&théorème computationnel par intervalles
&réduction \(72\to60\to52\), similitude rationnelle, isolement de
Gershgorin et sortie dont la distance certifiée est comprise entre
\(1.774799\times10^{-9}\) et \(1.799703\times10^{-9}\) degrés de
\(C_{\rm reg}\)\\
\(C_{\rm band}=C_{\rm reg}\)
&affirmation réfutée pour la loi linéaire
&l’intervalle certifié de la différence exclut zéro\\
transition totale
\(\mathcal X_{\mathrm{TPK}}^{\mathrm{enr}}\to \Gtr\)
&réalisateur global requis
&\HTr--\HAdd--\HDef--\HRes--\HLoc définit le lecteur relatif ; la transition
totale exige en outre une mesure image dynamique sur \(\Gtr\)\\
\bottomrule
\end{longtable}

\section{Action prédimensionnelle au centre de la double projection de l’état généalogique}

La double projection part du même état enrichi et produit deux
lectures. Dans la branche archimédienne, les cylindres compatibles se raffinent et
leurs diamètres décroissent ; l’évaluation fixe une valeur limite et oublie une partie
de la généalogie. Dans la branche exceptionnelle, les incidences se prolongent vers
Witt, Golay, Niemeier--Leech, \(V^\natural\), le Monstre et le moonshine.

Le nouveau résultat ajoute un objet commun plus fort qu’une coïncidence de
cardinaux. Le module à onze modes, son secteur positif de rang trois et la
trace de l’action passent de la dodécaphase à Witt par \(U_W\). Le secteur orienté
\(t=7\) s’incorpore par \(W_{GK}\). La microfibre de
\(\pi\) apporte le projecteur régional \(\Pi_\pi^{(468)}\), et la mesure uniforme du
catalogue fixe son poids.

L’action positive se situe ainsi au centre des deux lectures :
\[
 \text{région}
 \quad\oplus\quad
 \text{polarisation orientée}
 \quad\longmapsto\quad
 \text{défaut positif}.
\]
L’évaluation archimédienne et l’évaluation exceptionnelle restent des
applications différentes ; l’entrelaceur identifie un secteur commun, il ne
confond pas leurs codomaines.

\section{Compatibilité de l’opérateur angulaire avec les cylindres du continu discret}

Une histoire HMT détermine des cylindres emboîtés
\[
 C_1\supset C_2\supset\cdots
\]
avec des préfixes compatibles. Lorsque le lecteur archimédien satisfait
\[
 \operatorname{diam}(C_n)\longrightarrow0,
\]
la complétude de \(\mathbb R\) fixe au plus une valeur dans l’intersection.
La nouvelle action ajoute une structure avant cette limite : elle distingue le coût
de la route, la classe d’émission et la multiplicité des présentations que la
valeur finale ne montre plus.

C’est le sens mathématique précis de l’affirmation selon laquelle le continu se lit comme
projection de généalogies discrètes. Le système des cylindres et ses
évaluations se prolongent à toute échelle ; le théorème global de couverture
démontre
\(\operatorname{ev}_{\mathbb R}(\Omega_{\infty,K})=K\) pour chaque compact
\(K\subset\mathbb R\) et, en dirigeant les intervalles, que l’image est tout
\(\mathbb R\). Le théorème fini de cette partie contrôle l’action angulaire
à chaque niveau et s’intègre dans cette couverture globale ; il ne la rouvre ni ne la
conditionne. La descente algébrique de chaque opération conserve sa vérification
propre.

\section{Compatibilité avec le corridor exceptionnel, le Moonshine et les écrans de la théorie M}

L’entrelaceur \(U_W\) aboutit dans le module d’incidence de \(W_{12}\).
À partir du code ternaire, les deux continuations exceptionnelles restent
séparées :
\[
 C_W\longrightarrow W_{12}
 \xrightarrow[\text{relatif à }(D,\ell)]{}W_{24},
\]
et
\[
 C_W\xrightarrow{\ \phi\ } N(A_2^{12})
\xrightarrow{\ v_\alpha\ }\Lambda_{24}
\longrightarrow V_\Lambda
\longrightarrow V^\natural
\longrightarrow\mathbb M
\longrightarrow\text{moonshine}.
\]
L’action positive apporte à cette branche un secteur transporté et une droite
d’orientation. Les flèches de Leech au moonshine sont les constructions
classiques appliquées à l’objet fini obtenu. L’action du Monstre se
réalise sur \(V^\natural\) ; les mots HMT ne sont pas identifiés ici à un
module direct pour cette action.

Les écrans \(12\to11\to10\), \(26\to24\) et l’involution T-duale forment
une interface mathématique exacte avec les dimensions et les dualités utilisées en
théorie M. L’identification physique complète relève de la mécanique
dimensionnelle et est typée dans la seconde partie par l’application supplémentaire vers
\(g_{MN}\), \(C_{MNP}\), \(\psi_M\), l’action et la supersymétrie. Elle n’est pas réservée
à un autre ouvrage et ne se déduit pas de l’action du Monstre.

\section{Factorisation spectrale et isolement des secteurs dominants}

La ligne prédimensionnelle est organisée en quatre strates.
\begin{enumerate}
  \item \textbf{Noyau fini.} APP, TRIT, l’émetteur TPK, le catalogue
  \(104\,976\to468\to243\), les cinq régions de \(\pi\), la
  dodécaphase, \(K\), \(\alpha\), Witt et la branche exceptionnelle possèdent
  des constructions explicites et des certificats.
  \item \textbf{Pont opératoriel.} L’incidence construit \(U_W\) ;
  le secteur orienté \(t=7\) construit \(P_G\) ; \(B_K\) et \(u_K\) sélectionnent
  l’identification isométrique orientée ; les trois projecteurs de rang trois
  sont identifiés relativement à cette règle.
  \item \textbf{Extension d’action.} \HTr fixe la mesure \(1/468\), \HAdd
  fixe le caractère additif et \HDef déclare le couplage
  \(\Delta_4P_3\). \HRes transforme le défaut en capacité et \HLoc porte la
  capacité sur les arêtes. Les cinq clauses produisent la loi locale.
  \item \textbf{Conséquence spectrale.} \Gtr produit
  \(T_{\rm band}\) et une phase angulaire prédimensionnelle certifiée, avec
  une différence relative comprise entre \(8.356\times10^{-10}\) et
  \(8.475\times10^{-10}\) par rapport à \(C_{\rm reg}\).
\end{enumerate}

Cela permet une clôture forte et typée de cette route : l’émetteur fini de la
section coinductive infinie, l’incidence, l’identification des
projecteurs, la mesure et l’action sont clos
relativement aux primitives déclarées, à la règle variationnelle
\(\mathcal E_B\) et aux clauses d’action \HTr--\HAdd--\HDef--\HRes--\HLoc. \HAdd
sélectionne la forme linéaire ; \HDef sélectionne sa seconde coordonnée ; \HRes et \HLoc
précisent comment elle se lit comme coût. Le théorème global de
couverture de tout \(\mathbb R\), l’induction à partir d’une transition totale de
\(\mathcal X_{\mathrm{TPK}}^{\mathrm{enr}}\) et l’équivalence physique avec
la théorie M sont formulés comme des extensions des résultats précédents.

\section{Compatibilité entre les réalisations matricielle, angulaire et réticulaire}

Les représentations construites satisfont simultanément les conditions
de compatibilité suivantes :
\begin{enumerate}
  \item cinq régions de \(\pi\), dont trois positives, une négative et une
  transversale ;
  \item \(P_G,P_3,Q_3\) comme projecteurs positifs et \(R_{B_K}\) comme
  involution d’orientation ;
  \item le secteur homogène orienté \(t=7\), et non le quotient hétérogène de \(t=8\) ;
  \item la droite \(L_{\rm or}\) dans l’identification isométrique triple ;
  \item le caractère relatif de la règle variationnelle qui sélectionne
  \(H_*,c_*,r_*\) ;
  \item \(5/468\) sur les émissions et \(7/729\) sur les germes, sans
  les intervertir ;
  \item la linéarité du caractère comme conséquence de \HAdd et le coefficient
  \(\Delta_4\) comme contenu séparé de \HDef ;
  \item \((y_{\rm band},C_{\rm band})\) comme sortie spectrale certifiée,
  séparée de \((y_{\rm reg},C_{\rm reg})\), la réalisation régionale
  canonique de la section suivante ;
  \item \(\alpha_{\rm HMT}\) comme entrée de \(x\), de sorte que la lecture
  du déterminant est réversible et non indépendante ;
  \item les branches \(W_{12}\to W_{24}\) et
  \(A_2^{12}\to\Lambda_{24}\) comme constructions séparées ;
  \item le lecteur relatif \HTr--\HAdd--\HDef--\HRes--\HLoc distingué d’une
  transition totale, dont la mesure image dynamique constitue une
  construction supplémentaire.
\end{enumerate}

Ces conditions ne sont pas des notes éditoriales. Elles font partie de la structure
logique du résultat.

%
\endgroup

% La introducción heredada del capítulo 38 despejaba primero C*/2-alpha.
% Se conserva como antecedente en la cantera, pero la residencia pública
% presenta la construcción directa eta_ret -> hbar_ret -> (A,C*).
\HMTFlushCurrentChapter
\chapter{Équivalence exacte des cartes globale et analytique de la paire angulaire : degrés, radians, phases \texorpdfstring{\(q_{\pm}\)}{q+/-} et orientation complémentaire}
\label{chap:p3-final:38:angulo_contraangulo}
La composante sans dimension postérieure au retour est construite avant le contre-angle par \(\eta_{\mathrm{ret}}=H_5-(90/\pi)C_\pi(\alpha_{\mathrm{HMT}})\).
Dans la branche déterminantielle parallèle, on obtient la valuation \(10^{-34}\), et les deux branches confluent vers la grandeur \(\hbar_{\mathrm{ret}}=\eta_{\mathrm{ret}}10^{-34}U_S\).
Ce n'est qu'ensuite que sont publiés \(A_{\mathrm{deg}}=10^3\alpha_{\mathrm{HMT}}\) et \(C^*_{\mathrm{deg}}=2(\eta_{\mathrm{ret}}+\alpha_{\mathrm{HMT}})\) ; l'isolement de \(C^*/2-\alpha_{\mathrm{HMT}}\) demeure une vérification inverse.
% Adaptador único del capítulo 38 del sucesor 102.
% Sustituye exclusivamente la jerarquía de encabezados. Los cuerpos de los
% dos propietarios científicos y del delta Ley 9 se conservan línea a línea,
% reordenados conforme al DAG causal aprobado. Véase MAPA_SEGMENTOS_CAPITULO_38.tsv.

% HMT_SOURCE_SEGMENT_BEGIN id=B00 source=colaboracion/parte_iii/source/public_final/base_83/c30_body.tex body_lines=1-2
\label{chap:p3-83-30-angulo_contraangulo}

% HMT_SOURCE_SEGMENT_END id=B00
% HMT_SOURCE_SEGMENT_BEGIN id=A00 source=colaboracion/parte_iii/source/public_final/ascensos_20260826/16f_realizacion_analitica_contraangulo.tex body_lines=1-4
\label{p3-20260826:chap:realizacion-analitica-contraangulo}
\index{contre-angle}
\index{microfibre de pi@microfibre de \(\pi\)!caractère analytique}
\index{caractère déterminantal}
% HMT_SOURCE_SEGMENT_END id=A00
% HMT_SOURCE_SEGMENT_BEGIN id=D00 source=manuscrito/sucesor_102/deltas_ley9/c38_par_angular_cartas_y_canales.tex body_lines=1
% Distribución causal exacta del delta Ley 9; contenido conservado sin síntesis.
% HMT_SOURCE_SEGMENT_END id=D00

\section{Phase angulaire intrinsèque \texorpdfstring{\(\alpha_{\mathrm{HMT}}\)}{alpha HMT}, relèvement nonadique \texorpdfstring{\(10^3\)}{10^3} et précurseur \texorpdfstring{\(x_A\)}{x_A}}
\label{sec:s102:c38:1}

\subsection{Construction du précurseur analytique \texorpdfstring{\(x_A\)}{x_A}, de la contraction \texorpdfstring{\(D_A\)}{D_A} et de la correction régionale \texorpdfstring{\(\mathscr C_\pi(\alpha)\)}{C_pi(alpha)}}
\label{sec:s102:c38:1:1}

% HMT_SOURCE_SEGMENT_BEGIN id=B01 source=colaboracion/parte_iii/source/public_final/base_83/c30_body.tex body_lines=4-32
\label{p3-83-c30-s1:chap:realizacion-analitica-contraangulo}
\index{coordonnée angulaire conjuguée \(C^*\)}
\index{microfibre de pi@microfibre de \(\pi\)!caractère analytique}
\index{caractère déterminantal}

Le catalogue fini et une évaluation analytique répondent à des questions de types
distincts. Le catalogue détermine quelles régions existent, lesquelles partagent un
retour et quelle information chacune conserve. Une évaluation analytique doit
en outre décider comment la longueur d'un mot se transforme en degré, comment
une branche se prolonge après son retour et avec quel caractère scalaire cette
prolongation se lit. Ce chapitre construit cette seconde opération de manière
explicite.

La construction part de quatre données déjà obtenues dans les chapitres
précédents : la microfibre à cinq régions de \(\pi\), la longueur six de
chaque mot régional, la clôture dodécaphasique de \(\alpha\) et le transport
bicouche associé à la coordonnée commune
\(A=10^3\alpha\). Aucune valeur de la coordonnée angulaire conjuguée \(C_*^{\rm HMT}\), de la composante réduite
d'action, de la fonctionnelle de Barbero--Immirzi ou d'une constante métrologique n'est
employée dans l'évaluation. Le résultat est un caractère régional convergent,
une récurrence exacte et une phase orientée dont la carte sexagésimale reproduit
la coordonnée angulaire conjuguée publiée.

La séparation entre la partie finie et la partie analytique est essentielle. L'entier
\(169\) sera un théorème du catalogue. La série qui le prolonge sera le
théorème d'un lecteur analytique dont les règles sont déclarées avant de l'évaluer.
Ainsi, les prémisses supplémentaires ne restent pas dissimulées dans une coïncidence
décimale.

% HMT_SOURCE_SEGMENT_END id=B01
\section{Construction préangulaire de \texorpdfstring{\(\eta_{\mathrm{ret}}\)}{eta_ret} au moyen de \texorpdfstring{\(H_5\)}{H5}, \texorpdfstring{\(D_A\)}{D_A} et \texorpdfstring{\(\mathscr C_\pi(\alpha)\)}{C_pi(alpha)}}
\label{sec:s102:c38:2}

\subsection{Dérivation indépendante de \texorpdfstring{\(\eta_{\mathrm{ret}}=H_5-(90/\pi)\mathscr C_\pi(\alpha)\)}{eta_ret=H5-(90/pi)C_pi(alpha)} au moyen du lecteur analytique régional}
\label{sec:s102:c38:2:1}

% HMT_SOURCE_SEGMENT_BEGIN id=A01 source=colaboracion/parte_iii/source/public_final/ascensos_20260826/16f_realizacion_analitica_contraangulo.tex body_lines=5-38,40-47

Le catalogue fini et une évaluation analytique répondent à des questions de types
distincts. Le catalogue détermine quelles régions existent, lesquelles partagent un
retour et quelle information chacune conserve. Une évaluation analytique doit
décider, en outre, comment la longueur d'un mot se transforme en degré, comment
une branche se prolonge après son retour et avec quel caractère scalaire on lit
cette prolongation. Cette section construit cette seconde opération de manière
explicite.

La construction part de la microfibre de cinq régions de \(\pi\), de la
longueur six de chaque mot régional, de la clôture dodécaphasique de \(\alpha\)
et de l'insertion angulaire déclarée
\[
 \iota_\angle(x):=10^3x\,u_\circ,\qquad
 \Adeg:=10^3\alpha .
\]
Ici, \(u_\circ\) est la base sexagésimale dont le tour complet a pour
coordonnée \(360\).
Le facteur \(10^3\) appartient à cette normalisation de carte et ne se déduit pas de la
clôture arithmétique de \(\alpha\). Le transport bicouche est construit ensuite
à partir de \(\Adeg\). La fonction qui évalue le lecteur ne reçoit comme
argument ni le contre-angle, ni la coordonnée postérieure au retour, ni la fonctionnelle
hexade--octade, ni une constante métrologique. Le résultat est un caractère
régional convergent, une récurrence exacte et une phase orientée dans la carte
sexagésimale déclarée.

La séparation entre la partie finie et la partie analytique est essentielle. L'entier
\(169\) est un théorème du catalogue. La série qui le prolonge appartient
à un lecteur analytique dont les règles sont déclarées avant son évaluation.

La phase est notée \(y_*\) et, lorsque l'on souligne son origine régionale,
\(y_{\rm reg}\) ; les deux symboles désignent le même lecteur construit
ci-après.


L'évaluation reçoit le catalogue de \(468\) émissions, la valeur exacte de
\(\alpha_{\rm HMT}\), \(\varphi=(1+\sqrt5)/2\), la formule fixée de \(H_5\),
la carte angulaire et les cinq règles du lecteur régional. Avec ces données,
elle calcule \(169\), \(D_A\), la récurrence et \(C^*\). Le théorème est exact dans
cette classe déclarée de lecteurs ; il n'utilise pas \(C^*\) comme entrée et n'affirme pas que
les règles de \(H_5\) constituent l'unique classe naturelle possible.

% HMT_SOURCE_SEGMENT_END id=A01
\subsection{Persistance du retour dans les \texorpdfstring{\(5\)}{5} régions de la microfibre}
\label{sec:s102:c38:2:2}

% HMT_SOURCE_SEGMENT_BEGIN id=B02 source=colaboracion/parte_iii/source/public_final/base_83/c30_body.tex body_lines=34-155

Soit
\[
\mathcal F_\pi=\Psi^{-1}(010211)\subset\mathcal U_6
\]
la microfibre relative au catalogue. Rappelons ses cinq éléments, maintenant séparés en bloc initial et bloc terminal :
\[
\begin{aligned}
 &(501,614,498\mid169,272,272),\\
 &(810,923,870\mid169,272,272),\\
 &(810,983,810\mid169,272,272),\\
 &(870,923,810\mid169,272,272),\\
 &(870,923,810\mid418,674,521).
\end{aligned}
\]
Nous définissons la projection de retour
\[
p_{\rm ret}:\mathbb Z^6\longrightarrow\mathbb Z^3,
\qquad
p_{\rm ret}(u_1,\ldots,u_6)=(u_4,u_5,u_6).
\]

\begin{proposition}[Bloc terminal persistant]
\label{p3-83-c30-s1:prop:bloque-terminal-persistente}
Le multiensemble \(p_{\rm ret}(\mathcal F_\pi)\) possède un unique élément de
multiplicité maximale :
\[
b_\pi=(169,272,272),
\qquad
\operatorname{mult}_{\mathcal F_\pi}(b_\pi)=4.
\]
L'unique bloc terminal restant est
\[
b_\pi^-=(418,674,521)
\]
et a une multiplicité de un.
\end{proposition}

\begin{proof}
L'affirmation s'obtient par énumération exacte des cinq lignes du catalogue qui satisfont \(\Psi(U)=010211\). Quatre lignes possèdent le bloc terminal \((169,272,272)\) ; la cinquième, qui porte le retour muté, possède \((418,674,521)\). La multiplicité maximale est quatre et n'est atteinte qu'une seule fois.
\end{proof}

La différence orientée entre le retour muté et le retour persistant est
\[
 \omega_\pi=b_\pi^- - b_\pi=(249,402,249),
 \qquad
 \langle(1,1,1),\omega_\pi\rangle=900.
\]
Cette identité sépare deux objets qu'il ne faut pas confondre. L'entier
\(169\) est une coordonnée d'une donnée terminale, c'est-à-dire une donnée de degré
zéro. Le vecteur \(\omega_\pi\) peut être lu comme une cochaîne de degré un
après orientation de l'arête qui relie les deux blocs. L'appeler cocycle exigerait,
en outre, de fixer un complexe de cochaînes et de vérifier que son cobord s'annule.
La prolongation analytique construite plus loin ne présuppose pas cette
identification.

La sélection de \(169\) n'exige pas de privilégier la première position du bloc.
Le groupe symétrique \(\mathfrak S_3\) agit en permutant ses trois coordonnées.
Le stabilisateur de \(b_\pi\) échange les deux occurrences de \(272\) et
fixe l'unique coordonnée de multiplicité un. Notons
\(\operatorname{sing}(b)\) cette coordonnée lorsque le bloc est de type
\((r,s,s)\), avec \(r\ne s\). De manière équivalente,
\[
\operatorname{sing}(b)
=\sum_{j=1}^3b_j-2\operatorname{moda}(b).
\]
En particulier,
\[
\boxed{\operatorname{sing}(b_\pi)=169.}
\]

\begin{definition}[Sélecteur résiduel de la microfibre]
\label{p3-83-c30-s1:def:selector-residual-pi}
Le sélecteur \(\rho_\pi\) applique successivement deux opérations :
\begin{enumerate}
  \item sélectionne le bloc terminal de multiplicité maximale au sein de \(p_{\rm ret}(\mathcal F_\pi)\) ;
  \item retient la coordonnée fixe sous l'action du stabilisateur de ce bloc.
\end{enumerate}
Par la \cref{p3-83-c30-s1:prop:bloque-terminal-persistente},
\[
\boxed{\rho_\pi=169.}
\]
\end{definition}

\begin{theorem}[Unicité équivariante du retour]
\label{p3-83-c30-s1:thm:unicidad-retorno-169}
Parmi les sélecteurs qui
\begin{enumerate}
  \item sont invariants sous les permutations des cinq régions ;
  \item sont équivariants sous les permutations des trois coordonnées
  terminales ;
  \item sélectionnent le bloc de multiplicité maximale ;
  \item retiennent une coordonnée fixe sous l'action de son stabilisateur,
\end{enumerate}
la valeur de retour est déterminée de manière unique et vaut \(169\).
\end{theorem}

\begin{proof}
L'invariance sous \(\mathfrak S_5\) empêche d'utiliser l'ordre des lignes. La
troisième condition sélectionne \(b_\pi\), qui est l'unique mode d'après la
\cref{p3-83-c30-s1:prop:bloque-terminal-persistente}. Son stabilisateur dans
\(\mathfrak S_3\) possède deux orbites sur les coordonnées : la paire de
valeurs \(272\) et la valeur singleton \(169\). La quatrième condition sélectionne la
seconde orbite. L'équivariance conserve sa valeur lorsque sa position est déplacée.
\end{proof}

Cet argument améliore la lecture positionnelle du retour. Le \(169\) n'est pas
« la première coordonnée » d'une ligne choisie : c'est un invariant du
multiensemble régional et de l'action de ses symétries de réordonnancement.
Il n'est pas non plus sélectionné pour sa rareté globale. Dans les \(468\) émissions du
catalogue, \(169\) apparaît \(84\) fois, uniformément réparti avec
\(14\) occurrences dans chacune des six positions. L'unicité du
\cref{p3-83-c30-s1:thm:unicidad-retorno-169} est donc relationnelle et régionale : elle réside
dans la combinaison de la microfibre, de la persistance modale et du stabilisateur, non dans la
fréquence isolée de l'entier.

% HMT_SOURCE_SEGMENT_END id=B02
\subsection{Construction du lecteur analytique régional et de son domaine de convergence}
\label{sec:s102:c38:2:3}

% HMT_SOURCE_SEGMENT_BEGIN id=B05 source=colaboracion/parte_iii/source/public_final/base_83/c30_body.tex body_lines=261-301

L'extension analytique est fixée par les règles suivantes.

\begin{definition}[Lecteur régional déterminantal]
\label{p3-83-c30-s1:def:lector-regional-determinantal}
Le lecteur régional déterminantal satisfait :
\begin{enumerate}
  \item \emph{compatibilité de degré} : la longueur cylindrique est le degré
  analytique;
  \item \emph{décomposition de renouvellement} : l'impulsion finie de retour et
  la continuation stricte appartiennent à des termes additifs distincts ;
  \item \emph{normalisation de continuation} : après l'enregistrement, une seule
  fois, du résidu \(\rho_\pi\), la branche survivante commence à l'unité de
  la droite déterminant ;
  \item \emph{covariance par concaténation} : chaque prolongation élémentaire
  agit au moyen de
  \(\bigwedge^2\mathcal T=D_A\) ;
  \item \emph{orientation} : dans la convention cardinale fixée pour APP, le retour régional appartient au secteur compensateur et est soustrait de la phase d'action de cinquième degré.
\end{enumerate}
\end{definition}

La troisième règle a un contenu mathématique. La valeur \(169\) est enregistrée
une fois comme amplitude du retour ; elle n'est pas multipliée à nouveau à chaque
prolongation. La continuation stricte compte l'unique droite qui se poursuit et
la normalise par son unité. C'est la différence entre une récurrence de
renouvellement et l'itération homogène de l'impulsion initiale.

La nécessité de déclarer cette normalisation peut être prouvée. Pour tout
\(\lambda\in\mathbb R\), la famille
\[
F_\lambda(z)
=169z^6+\lambda\frac{D_Az^7}{1-D_Az}
\]
conserve le même retour fini et le même rapport déterminantal entre
coefficients consécutifs. Le catalogue et le rapport seuls ne distinguent pas
\(\lambda\). L'unité de la droite déterminant fixe
\(\lambda=1\) avant d'évaluer \(z=\alpha\). De cette manière, la normalisation
ne s'obtient pas en ajustant la valeur finale ; elle fait partie de la définition du
caractère.

% HMT_SOURCE_SEGMENT_END id=B05
\subsection{Résolution analytique de la récurrence de retour}
\label{sec:s102:c38:2:4}

% HMT_SOURCE_SEGMENT_BEGIN id=B06 source=colaboracion/parte_iii/source/public_final/base_83/c30_body.tex body_lines=303-364

Nous définissons les coefficients \(\kappa_n\) par
\[
\kappa_6=169,
\qquad
\kappa_7=D_A,
\qquad
\kappa_{n+1}=D_A\kappa_n\quad(n\ge7).
\]
La première égalité est l'impulsion régionale ; les deux suivantes décrivent la
continuation normalisée.

\begin{theorem}[Réalisation analytique graduée]
\label{p3-83-c30-s1:thm:realizacion-analitica-graduada}
Le germe analytique compatible avec les règles de la
\cref{p3-83-c30-s1:def:lector-regional-determinantal} est unique et vaut
\[
\boxed{
\mathscr C_\pi(z)
=169z^6+\frac{D_Az^7}{1-D_Az}.}
\]
Ses coefficients satisfont
\[
\kappa_n=D_A^{n-6}\qquad(n\ge7).
\]
\end{theorem}

\begin{proof}
Écrivons
\[
F(z)=169z^6+\sum_{k\ge1}c_kz^{6+k}.
\]
La normalisation de continuation et la covariance déterminantale donnent
\[
c_1=D_A,
\qquad
c_{k+1}=D_Ac_k.
\]
Par récurrence, \(c_k=D_A^k\). La somme géométrique produit
\[
\sum_{k\ge1}D_A^kz^{6+k}
=\frac{D_Az^7}{1-D_Az}.
\]
Aucun coefficient libre ne subsiste.
\end{proof}

Comme \(0<\alpha<1\) et
\(0<D_A<1\),
\[
0<D_A\alpha<1.
\]
En conséquence, \(\mathscr C_\pi(\alpha)\) converge absolument et son
rayon de convergence est \(D_A^{-1}>1\). Numériquement,
\[
D_A=0.7751291192471564301888840964802\ldots,
\]
\[
D_A\alpha=0.00565639046986492673885079095469\ldots.
\]
La rapidité de cette contraction explique pourquoi la somme complète et sa
troncature de degré sept possèdent les mêmes quinze premiers chiffres dans la
carte finale.

% HMT_SOURCE_SEGMENT_END id=B06
\subsection{Persistance du retour dans les \texorpdfstring{\(5\)}{5} régions et compatibilité de la composante \texorpdfstring{\(\eta_{\mathrm{ret}}\)}{eta_ret}}
\label{sec:s102:c38:2:5}

% HMT_SOURCE_SEGMENT_BEGIN id=A02 source=colaboracion/parte_iii/source/public_final/ascensos_20260826/16f_realizacion_analitica_contraangulo.tex body_lines=49-170

Soit
\[
\mathcal F_\pi=\Psi^{-1}(010211)\subset\mathcal U_6
\]
la microfibre relative au catalogue. Rappelons ses cinq éléments, maintenant séparés en bloc initial et bloc terminal :
\[
\begin{aligned}
 &(501,614,498\mid169,272,272),\\
 &(810,923,870\mid169,272,272),\\
 &(810,983,810\mid169,272,272),\\
 &(870,923,810\mid169,272,272),\\
 &(870,923,810\mid418,674,521).
\end{aligned}
\]
Nous définissons la projection de retour
\[
p_{\rm ret}:\mathbb Z^6\longrightarrow\mathbb Z^3,
\qquad
p_{\rm ret}(u_1,\ldots,u_6)=(u_4,u_5,u_6).
\]

\begin{proposition}[Bloc terminal persistant]
\label{p3-20260826:prop:bloque-terminal-persistente}
Le multiensemble \(p_{\rm ret}(\mathcal F_\pi)\) possède un unique élément de
multiplicité maximale :
\[
b_\pi=(169,272,272),
\qquad
\operatorname{mult}_{\mathcal F_\pi}(b_\pi)=4.
\]
L'unique bloc terminal restant est
\[
b_\pi^-=(418,674,521)
\]
et a une multiplicité de un.
\end{proposition}

\begin{proof}
L'affirmation s'obtient par énumération exacte des cinq lignes du catalogue qui satisfont \(\Psi(U)=010211\). Quatre lignes possèdent le bloc terminal \((169,272,272)\) ; la cinquième, qui porte le retour muté, possède \((418,674,521)\). La multiplicité maximale est quatre et n'est atteinte qu'une seule fois.
\end{proof}

La différence orientée entre le retour muté et le retour persistant est
\[
 \omega_\pi=b_\pi^- - b_\pi=(249,402,249),
 \qquad
 \langle(1,1,1),\omega_\pi\rangle=900.
\]
Cette identité sépare deux objets qu'il ne faut pas confondre. L'entier
\(169\) est une coordonnée d'une donnée terminale, c'est-à-dire une donnée de degré
zéro. Le vecteur \(\omega_\pi\) peut être lu comme une cochaîne de degré un
après orientation de l'arête qui relie les deux blocs. L'appeler cocycle exigerait,
en outre, de fixer un complexe de cochaînes et de vérifier que son cobord s'annule.
La prolongation analytique construite plus loin ne présuppose pas cette
identification.

La sélection de \(169\) n'exige pas de privilégier la première position du bloc.
Le groupe symétrique \(\mathfrak S_3\) agit en permutant ses trois coordonnées.
Le stabilisateur de \(b_\pi\) échange les deux occurrences de \(272\) et
fixe l'unique coordonnée de multiplicité un. Notons
\(\operatorname{sing}(b)\) cette coordonnée lorsque le bloc est de type
\((r,s,s)\), avec \(r\ne s\). De manière équivalente,
\[
\operatorname{sing}(b)
=\sum_{j=1}^3b_j-2\operatorname{moda}(b).
\]
En particulier,
\[
\boxed{\operatorname{sing}(b_\pi)=169.}
\]

\begin{definition}[Sélecteur résiduel de la microfibre]
\label{p3-20260826:def:selector-residual-pi}
Le sélecteur \(\rho_\pi\) applique successivement deux opérations :
\begin{enumerate}
  \item sélectionne le bloc terminal de multiplicité maximale au sein de \(p_{\rm ret}(\mathcal F_\pi)\) ;
  \item retient la coordonnée fixe sous l'action du stabilisateur de ce bloc.
\end{enumerate}
Par la \cref{p3-20260826:prop:bloque-terminal-persistente},
\[
\boxed{\rho_\pi=169.}
\]
\end{definition}

\begin{theorem}[Unicité équivariante du retour]
\label{p3-20260826:thm:unicidad-retorno-169}
Parmi les sélecteurs qui
\begin{enumerate}
  \item sont invariants sous les permutations des cinq régions ;
  \item sont équivariants sous les permutations des trois coordonnées
  terminales ;
  \item sélectionnent le bloc de multiplicité maximale ;
  \item retiennent une coordonnée fixe sous l'action de son stabilisateur,
\end{enumerate}
la valeur de retour est déterminée de manière unique et vaut \(169\).
\end{theorem}

\begin{proof}
L'invariance sous \(\mathfrak S_5\) empêche d'utiliser l'ordre des lignes. La
troisième condition sélectionne \(b_\pi\), qui est l'unique mode d'après la
\cref{p3-20260826:prop:bloque-terminal-persistente}. Son stabilisateur dans
\(\mathfrak S_3\) possède deux orbites sur les coordonnées : la paire de
valeurs \(272\) et la valeur singleton \(169\). La quatrième condition sélectionne la
seconde orbite. L'équivariance conserve sa valeur lorsque sa position est déplacée.
\end{proof}

Cet argument améliore la lecture positionnelle du retour. Le \(169\) n'est pas
« la première coordonnée » d'une ligne choisie : c'est un invariant du
multiensemble régional et de l'action de ses symétries de réordonnancement.
Il n'est pas non plus sélectionné pour sa rareté globale. Dans les \(468\) émissions du
catalogue, \(169\) apparaît \(84\) fois, uniformément réparti avec
\(14\) occurrences dans chacune des six positions. L'unicité du
\cref{p3-20260826:thm:unicidad-retorno-169} est donc relationnelle et régionale : elle réside
dans la combinaison de la microfibre, de la persistance modale et du stabilisateur, non dans la
fréquence isolée de l'entier.

% HMT_SOURCE_SEGMENT_END id=A02
\subsection{Application du lecteur régional au couple angulaire conjugué}
\label{sec:s102:c38:2:6}

% HMT_SOURCE_SEGMENT_BEGIN id=A05 source=colaboracion/parte_iii/source/public_final/ascensos_20260826/16f_realizacion_analitica_contraangulo.tex body_lines=277-317

L'extension analytique est fixée par les règles suivantes.

\begin{definition}[Lecteur régional déterminantal]
\label{p3-20260826:def:lector-regional-determinantal}
Le lecteur régional déterminantal satisfait :
\begin{enumerate}
  \item \emph{compatibilité de degré} : la longueur cylindrique est le degré
  analytique;
  \item \emph{décomposition de renouvellement} : l'impulsion finie de retour et
  la continuation stricte appartiennent à des termes additifs distincts ;
  \item \emph{normalisation de continuation} : après l'enregistrement, une seule
  fois, du résidu \(\rho_\pi\), la branche survivante commence à l'unité de
  la droite déterminant ;
  \item \emph{covariance par concaténation} : chaque prolongation élémentaire
  agit au moyen de
  \(\bigwedge^2\mathcal T=D_A\) ;
  \item \emph{orientation} : dans la convention cardinale fixée pour APP, le retour régional appartient au secteur compensateur et est soustrait de la phase d'action de cinquième degré.
\end{enumerate}
\end{definition}

La troisième règle a un contenu mathématique. La valeur \(169\) est enregistrée
une fois comme amplitude du retour ; elle n'est pas multipliée à nouveau à chaque
prolongation. La continuation stricte compte l'unique droite qui se poursuit et
la normalise par son unité. C'est la différence entre une récurrence de
renouvellement et l'itération homogène de l'impulsion initiale.

La nécessité de déclarer cette normalisation peut être prouvée. Pour tout
\(\lambda\in\mathbb R\), la famille
\[
F_\lambda(z)
=169z^6+\lambda\frac{D_Az^7}{1-D_Az}
\]
conserve le même retour fini et le même rapport déterminantal entre
coefficients consécutifs. Le catalogue et le rapport seuls ne distinguent pas
\(\lambda\). L'unité de la droite déterminant fixe
\(\lambda=1\) avant d'évaluer \(z=\alpha\). De cette manière, la normalisation
ne s'obtient pas en ajustant la valeur finale ; elle fait partie de la définition du
caractère.

% HMT_SOURCE_SEGMENT_END id=A05
\subsection{Somme fermée et publication de la composante postérieure au retour}
\label{sec:s102:c38:2:7}

% HMT_SOURCE_SEGMENT_BEGIN id=A06 source=colaboracion/parte_iii/source/public_final/ascensos_20260826/16f_realizacion_analitica_contraangulo.tex body_lines=319-380

Nous définissons les coefficients \(\kappa_n\) par
\[
\kappa_6=169,
\qquad
\kappa_7=D_A,
\qquad
\kappa_{n+1}=D_A\kappa_n\quad(n\ge7).
\]
La première égalité est l'impulsion régionale ; les deux suivantes décrivent la
continuation normalisée.

\begin{theorem}[Réalisation analytique graduée]
\label{p3-20260826:thm:realizacion-analitica-graduada}
Le germe analytique compatible avec les règles de la
\cref{p3-20260826:def:lector-regional-determinantal} est unique et vaut
\[
\boxed{
\mathscr C_\pi(z)
=169z^6+\frac{D_Az^7}{1-D_Az}.}
\]
Ses coefficients satisfont
\[
\kappa_n=D_A^{n-6}\qquad(n\ge7).
\]
\end{theorem}

\begin{proof}
Écrivons
\[
F(z)=169z^6+\sum_{k\ge1}c_kz^{6+k}.
\]
La normalisation de continuation et la covariance déterminantale donnent
\[
c_1=D_A,
\qquad
c_{k+1}=D_Ac_k.
\]
Par récurrence, \(c_k=D_A^k\). La somme géométrique produit
\[
\sum_{k\ge1}D_A^kz^{6+k}
=\frac{D_Az^7}{1-D_Az}.
\]
Aucun coefficient libre ne subsiste.
\end{proof}

Comme \(0<\alpha<1\) et
\(0<D_A<1\),
\[
0<D_A\alpha<1.
\]
En conséquence, \(\mathscr C_\pi(\alpha)\) converge absolument et son
rayon de convergence est \(D_A^{-1}>1\). Numériquement,
\[
D_A=0.7751291192471564301888840964802\ldots,
\]
\[
D_A\alpha=0.00565639046986492673885079095469\ldots.
\]
La rapidité de cette contraction explique pourquoi la somme complète et sa
troncature de degré sept possèdent les mêmes quinze premiers chiffres dans la
carte finale.

% HMT_SOURCE_SEGMENT_END id=A06
\section{Confluence de \texorpdfstring{\(\eta_{\mathrm{ret}}\)}{eta_ret} avec la valuation \texorpdfstring{\(\varepsilon_H=-34\)}{epsilon_H=-34} dans la grandeur \texorpdfstring{\(\hbar_{\mathrm{ret}}\)}{hbar_ret}}
\label{sec:s102:c38:3}

\subsection{Dérivation du caractère déterminantal du transport parangulaire}
\label{sec:s102:c38:3:1}

% HMT_SOURCE_SEGMENT_BEGIN id=B04 source=colaboracion/parte_iii/source/public_final/base_83/c30_body.tex body_lines=196-259

La fermeture dodécaphasique fournit \(\alpha\). La complétion des trois
couples nonadiques fournit l'échelle \(10^3\), de sorte que
\[
A=10^3\alpha.
\]
La carte angulaire archimédienne étant choisie, la coordonnée commune en radians est
\[
x=\frac{\pi A}{180}=\frac{50\pi}{9}\alpha.
\]
Soit \(R\) une involution réelle de trace nulle en dimension deux,
\[
R^2=I_2,
\qquad
\operatorname{tr}R=0,
\]
et considérons le transport formel bicouche
\[
\mathcal T(x,y)=\exp(-xI_2-yR).
\]
La variable orientée \(y\) n'a pas encore été évaluée. Cependant, l'action
de \(\mathcal T\) sur la droite déterminante en est indépendante :
\[
\begin{aligned}
\det\mathcal T(x,y)
&=\exp\!\left(\operatorname{tr}(-xI_2-yR)\right)\\
&=e^{-2x}.
\end{aligned}
\]
On définit
\[
\boxed{
D_A=e^{-2x}=e^{-100\pi\alpha/9}.}
\]
L'indice rappelle qu'il s'agit du déterminant de la contraction commune associée à \(A\), et non d'une fonction de \(C_*^{\rm HMT}\).

\begin{proposition}[Caractère de la droite déterminante]
\label{p3-83-c30-s1:prop:caracter-determinante}
L'application
\[
\delta_A:(\mathbb N_0,+)\longrightarrow(\mathbb R_{>0},\cdot),
\qquad
\delta_A(k)=D_A^k,
\]
est l'unique caractère multiplicatif dont la valeur sur une prolongation
élémentaire est \(D_A\).
\end{proposition}

\begin{proof}
La multiplicativité donne
\[
\delta_A(k)
=\delta_A(\underbrace{1+\cdots+1}_{k\text{ fois}})
=\delta_A(1)^k=D_A^k.
\]
La même identité prouve l'existence et l'unicité.
\end{proof}

Le déterminant reste invariant par conjugaison de \(\mathcal T\) et par échange des feuillets \(R\mapsto-R\). La queue construite à partir de \(D_A\) appartient donc au mode commun. L'orientation apparaîtra dans le signe avec lequel le caractère régional corrige la phase.

% HMT_SOURCE_SEGMENT_END id=B04
\subsection{Caractère d'action de degré \texorpdfstring{\(5\)}{5} dans la construction régionale}
\label{sec:s102:c38:3:2}

% HMT_SOURCE_SEGMENT_BEGIN id=B07 source=colaboracion/parte_iii/source/public_final/base_83/c30_body.tex body_lines=366-405

Le bloc central
\[
\mathcal B_c=\begin{pmatrix}7&2\\2&7\end{pmatrix}
\]
possède les valeurs propres \(9\) et \(5\). Le cycle dodécaphasique apporte le rang
\(12\) ; ses orbites de pas trois possèdent une longueur de quatre ; et le recensement complet
apporte le quotient exact
\[
\frac{104976}{1944}=54.
\]
À partir de ces invariants, on fixe le caractère d'action de cinquième degré
\[
\boxed{
H_5(\alpha,\varphi)
=\frac12\sqrt{\frac{1000\alpha}{\varphi}}-\alpha
+\frac9{16}\alpha^2-\frac59\alpha^3
+\frac7{48}\alpha^4-\frac1{54}\alpha^5.}
\]
Les quatre coefficients rationnels peuvent s'écrire
\[
\frac9{16}=\frac{\lambda_+}{4^2},
\qquad
\frac59=\frac{\lambda_-}{\lambda_+},
\qquad
\frac7{48}=\frac{\tfrac12\operatorname{tr}\mathcal B_c}{4\cdot12},
\qquad
\frac1{54}=\frac{1944}{104976},
\]
avec \((\lambda_+,\lambda_-)=(9,5)\). Ces identités certifient la
provenance des nombres utilisés. L'affectation successive de ces
invariants aux degrés deux, trois, quatre et cinq fait partie du caractère
analytique ; elle ne se déduit pas de la seule existence du catalogue fini.

Cette précision élimine une ambiguïté habituelle. Une combinaison
d'invariants de la construction est une définition interne et reproductible ; sa
naturalité exige les règles du lecteur qui spécifient l'ordre, les signes et la
normalisation. Ici, ces règles sont visibles et sont soumises aux mêmes
contrôles négatifs que le reste de la construction.

% HMT_SOURCE_SEGMENT_END id=B07
\subsection{Décomposition exacte de la droite d'action en sections pré-retour et post-retour}
\label{sec:s102:c38:3:3}

% HMT_SOURCE_SEGMENT_BEGIN id=B12 source=colaboracion/parte_iii/source/public_final/base_83/c30_body.tex body_lines=571-617

Le caractère de cinquième degré et la composante d'action postérieure au retour
sont deux valeurs distinctes :
\[
H_5
=1.05457181764615446962582182943306\ldots,
\]
\[
\bar h_{\rm HMT}^{\rm ret}
=1.05457181691503906113369058472430\ldots.
\]
La première est l'action locale avant l'incorporation du caractère régional ; la
seconde est l'action orientée qui reste après le retour. La correction
qui engendre la coordonnée angulaire conjuguée empêche de les identifier.

La comparaison métrologique doit respecter cette distinction. La première valeur donne
la mantisse
\[
2\pi H_5
=6.62607014999998792600111034989947\ldots,
\]
tandis que l'action postérieure donne
\[
2\pi\bar h_{\rm HMT}^{\rm ret}
=6.62607014540625433351074984759288\ldots.
\]
Dans le Système international, la valeur
\[
h=6.62607015\times10^{-34}\;\mathrm{J\,s}
\]
est exacte par définition \cite{BIPMSI}. C'est pourquoi,
\[
2\pi H_5-6.62607015
=-1.2073998889650101\ldots\times10^{-14},
\]
et ce n'est pas une égalité exacte. La proximité de la valeur antérieure au retour est
un contrôle numérique remarquable ; elle n'autorise pas à remplacer la constante définitoire
du SI par la composante postérieure ni à confondre les deux actions HMT\@.

Ce résultat détermine avec précision la situation mathématique. Le lecteur
régional engendre \(C_*^{\rm HMT}\) et, à partir de celui-ci, la fonctionnelle
\(\Gamma_{6\to8}\). Le même calcul sépare l'approximation de Planck
produite par \(H_5\) de l'action réduite qui intervient dans
la coordonnée angulaire conjuguée. La séparation n'affaiblit pas la construction : elle élimine une
identification incompatible avec les chiffres exacts et transforme chaque
comparaison en une affirmation réfutable de type bien défini.

% HMT_SOURCE_SEGMENT_END id=B12
\subsection{Factorisation du caractère déterminantal dans les \texorpdfstring{\(2\)}{2} cartes de phase}
\label{sec:s102:c38:3:4}

% HMT_SOURCE_SEGMENT_BEGIN id=A04 source=colaboracion/parte_iii/source/public_final/ascensos_20260826/16f_realizacion_analitica_contraangulo.tex body_lines=211-275

La clôture dodécaphasique fournit \(\alpha\). L'insertion angulaire déclarée
emploie l'échelle \(10^3\), compatible avec le regroupement de trois paires
nonadiques, de sorte que
\[
\Adeg=10^3\alpha\,{}^\circ.
\]
La carte angulaire archimédienne étant choisie, la coordonnée commune en radians est
\[
x=: \Arad=\frac{\pi\Adeg}{180}=\frac{50\pi}{9}\alpha.
\]
Soit \(R\) une involution réelle de trace nulle en dimension deux,
\[
R^2=I_2,
\qquad
\operatorname{tr}R=0,
\]
et considérons le transport formel bicouche
\[
\mathcal T(x,y)=\exp(-xI_2-yR).
\]
La variable orientée \(y\) n'a pas encore été évaluée. Cependant, l'action
de \(\mathcal T\) sur la droite déterminante en est indépendante :
\[
\begin{aligned}
\det\mathcal T(x,y)
&=\exp\!\left(\operatorname{tr}(-xI_2-yR)\right)\\
&=e^{-2x}.
\end{aligned}
\]
On définit
\[
\boxed{
D_A=e^{-2x}=e^{-100\pi\alpha/9}.}
\]
L'indice rappelle qu'il s'agit du déterminant de la contraction commune
associée à \(A\), non d'une fonction du contre-angle.

\begin{proposition}[Caractère de la droite déterminante]
\label{p3-20260826:prop:caracter-determinante}
L'application
\[
\delta_A:(\mathbb N_0,+)\longrightarrow(\mathbb R_{>0},\cdot),
\qquad
\delta_A(k)=D_A^k,
\]
est l'unique caractère multiplicatif dont la valeur sur une prolongation
élémentaire est \(D_A\).
\end{proposition}

\begin{proof}
La multiplicativité donne
\[
\delta_A(k)
=\delta_A(\underbrace{1+\cdots+1}_{k\text{ fois}})
=\delta_A(1)^k=D_A^k.
\]
La même identité prouve l'existence et l'unicité.
\end{proof}

Le déterminant reste invariant par conjugaison de \(\mathcal T\) et par échange des feuillets \(R\mapsto-R\). La queue construite à partir de \(D_A\) appartient donc au mode commun. L'orientation apparaîtra dans le signe avec lequel le caractère régional corrige la phase.

% HMT_SOURCE_SEGMENT_END id=A04
\subsection{Publication du caractère d'action de degré \texorpdfstring{\(5\)}{5} dans la carte parangulaire}
\label{sec:s102:c38:3:5}

% HMT_SOURCE_SEGMENT_BEGIN id=A07 source=colaboracion/parte_iii/source/public_final/ascensos_20260826/16f_realizacion_analitica_contraangulo.tex body_lines=382-418

Le bloc central
\[
\mathcal B_c=\begin{pmatrix}7&2\\2&7\end{pmatrix}
\]
possède les valeurs propres \(9\) et \(5\). Le cycle dodécaphasique apporte le rang
\(12\) ; ses orbites de pas trois possèdent une longueur de quatre ; et le recensement complet
apporte le quotient exact
\[
\frac{104976}{1944}=54.
\]
Le lecteur en vigueur déclare le caractère d'action de cinquième degré
\[
\boxed{
H_5(\alpha,\varphi)
=\frac12\sqrt{\frac{1000\alpha}{\varphi}}-\alpha
+\frac9{16}\alpha^2-\frac59\alpha^3
+\frac7{48}\alpha^4-\frac1{54}\alpha^5.}
\]
Les quatre coefficients rationnels peuvent s'écrire
\[
\frac9{16}=\frac{\lambda_+}{4^2},
\qquad
\frac59=\frac{\lambda_-}{\lambda_+},
\qquad
\frac7{48}=\frac{\tfrac12\operatorname{tr}\mathcal B_c}{4\cdot12},
\qquad
\frac1{54}=\frac{1944}{104976},
\]
avec \((\lambda_+,\lambda_-)=(9,5)\). Ces identités fixent la
représentabilité arithmétique interne des coefficients du caractère. Leur
ordre, leurs signes et leur normalisation appartiennent à la règle graduée du lecteur
régional et sont appliqués avant toute carte d'unités. Les ablations
certifient que chaque terme intervient effectivement dans la phase résultante.
En particulier, ni la mantisse SI de l'action ni la valeur sexagésimale du
contre-angle ne font partie du domaine de la fonction génératrice.

% HMT_SOURCE_SEGMENT_END id=A07
\section{Application parangulaire \texorpdfstring{\(P(\alpha,\eta_{\mathrm{ret}})\)}{P(alpha,eta_ret)} et clôture bicouche \texorpdfstring{\(C^*=2(\eta_{\mathrm{ret}}+\alpha)\)}{C*=2(eta_ret+alpha)}}
\label{sec:s102:c38:4}

\subsection{Phase orientée et coordonnée angulaire conjuguée}
\label{sec:s102:c38:4:1}

% HMT_SOURCE_SEGMENT_BEGIN id=B08 source=colaboracion/parte_iii/source/public_final/base_83/c30_body.tex body_lines=407-472

La phase de cinquième degré antérieure au retour est
\[
y_5=\frac\pi{90}\bigl(H_5+\alpha\bigr).
\]
Le caractère régional complet produit la correction
\(\mathscr C_\pi(\alpha)\). Avec l'orientation fixée dans la
\cref{p3-83-c30-s1:def:lector-regional-determinantal}, nous définissons
\[
\boxed{
y_*=\frac\pi{90}\bigl(H_5+\alpha\bigr)
-169\alpha^6
-\frac{D_A\alpha^7}{1-D_A\alpha}.}
\]
C'est la formule primaire. L'unité sexagésimale n'intervient qu'à l'application de la carte
\[
C_*^{\rm HMT}=\frac{180}{\pi}y_*.
\]

\begin{theorem}[Coordonnée angulaire conjuguée du lecteur régional déterminantal]
\label{p3-83-c30-s1:thm:contraangulo-regional}
Une fois fixés la clôture dodécaphasique de \(\alpha\),
\(\varphi=(1+\sqrt5)/2\) et les règles du lecteur régional déterminantal, la
phase \(y_*\) et sa coordonnée angulaire conjuguée sont déterminées sans introduire la valeur
recherchée. Leur évaluation est
\[
\boxed{
y_*=0.03706622646583826398493779409230638\ldots,}
\]
\[
\boxed{
C_*^{\rm HMT}
=2.12373833896864572426195138039336\ldots{}^\circ.}
\]
Donc, à quinze décimales,
\[
\boxed{C_*^{\rm HMT}=2.123738338968646^\circ.}
\]
\end{theorem}

\begin{proof}
Le \cref{p3-83-c30-s1:thm:unicidad-retorno-169} détermine \(169\). La clôture de
\(\alpha\) détermine \(A\), \(x\) et
\(D_A=e^{-100\pi\alpha/9}\). Le \cref{p3-83-c30-s1:thm:realizacion-analitica-graduada}
détermine la série complète. La formule de \(H_5\) contient uniquement
\(\alpha\), \(\varphi\) et des invariants structurels déclarés. Substituer
ces données dans l'expression de \(y_*\) donne la valeur imprimée. La valeur décimale
de \(C_*^{\rm HMT}\) n'apparaît dans aucun des membres de droite.
\end{proof}

La composante réduite d'action postérieure au retour est désormais une
sortie :
\[
\boxed{
\bar h_{\rm HMT}^{\rm ret}
=\frac{C_*^{\rm HMT}}2-\alpha
=H_5-\frac{90}{\pi}\mathscr C_\pi(\alpha).}
\]
Numériquement,
\[
\bar h_{\rm HMT}^{\rm ret}
=1.05457181691503906113369058472430\ldots.
\]
L'identité inverse n'a pas été utilisée pour introduire \(C_*^{\rm HMT}\) ; elle
s'obtient après la génération de \(y_*\).

% HMT_SOURCE_SEGMENT_END id=B08
\subsection{Coordonnée conjuguée \texorpdfstring{\(C^*\)}{C*} et réversion de son orientation sous l'involution bicouche}
\label{sec:s102:c38:4:2}

% HMT_SOURCE_SEGMENT_BEGIN id=A08 source=colaboracion/parte_iii/source/public_final/ascensos_20260826/16f_realizacion_analitica_contraangulo.tex body_lines=420-503

La phase de cinquième degré antérieure au retour est
\[
y_5=\frac\pi{90}\bigl(H_5+\alpha\bigr).
\]
Le caractère régional complet produit la correction
\(\mathscr C_\pi(\alpha)\). Avec l'orientation fixée dans la
\cref{p3-20260826:def:lector-regional-determinantal}, nous définissons
\[
\boxed{
\Crad:=y_*=\frac\pi{90}\bigl(H_5+\alpha\bigr)
-169\alpha^6
-\frac{D_A\alpha^7}{1-D_A\alpha}.}
\]
Désormais, nous écrivons aussi
\[
 \boxed{y_{\rm reg}:=y_*}
\]
lorsqu'il est nécessaire de distinguer cette réalisation de la phase spectrale de
\Gtr. Les deux notations désignent ici le même objet régional canonique.
Telle est la formule primaire. L'unité sexagésimale n'intervient qu'à l'application de la carte
\[
\boxed{\Cdeg:=\frac{180}{\pi}\Crad.}
\]
Les symboles \(\Crad\) et \(\Cdeg\) ne désignent pas deux contre-angles : ce sont les
coordonnées radiale et sexagésimale du même élément angulaire.

\begin{theorem}[Contre-angle du lecteur régional déterminantal]
\label{p3-20260826:thm:contraangulo-regional}
La clôture dodécaphasique de \(\alpha\),
\(\varphi=(1+\sqrt5)/2\), le germe \(H_5\) et les règles du lecteur régional
déterminantal étant fixés, la phase \(\Crad\) et son contre-angle sont déterminés sans
introduire la valeur recherchée dans la fonction d'évaluation. Son évaluation est
\[
\boxed{
\Crad=0.03706622646583826398493779409230638\ldots,}
\]
\[
\boxed{
\Cdeg
=2.12373833896864572426195138039336\ldots{}^\circ.}
\]
Donc, à quinze décimales,
\[
\boxed{\Cdeg=2.123738338968646^\circ.}
\]
\end{theorem}

\begin{proof}
Le \cref{p3-20260826:thm:unicidad-retorno-169} détermine \(169\). La clôture de
\(\alpha\) détermine \(\Adeg\), \(\Arad\) et
\(D_A=e^{-100\pi\alpha/9}\). Le \cref{p3-20260826:thm:realizacion-analitica-graduada}
détermine la série complète. La formule fixée de \(H_5\) contient uniquement
\(\alpha\), \(\varphi\) et des invariants structurels déclarés. Substituer
ces données dans l'expression de \(\Crad\) donne la valeur imprimée. La valeur décimale
du contre-angle n'apparaît dans aucun des membres de droite de cette
évaluation. La comparaison avec une mantisse métrologique d'action est effectuée
uniquement après l'obtention de la phase et ne modifie pas la fonction génératrice.
\end{proof}

Soit \(u_\circ\) la base sexagésimale dont le tour complet a pour coordonnée
\(360\), et introduisons la coordonnée angulaire
\(\alpha_{\angle,{\rm deg}}:=\alpha_{\rm HMT}\). La sortie postérieure au
retour est alors typée par
\[
\boxed{
\etaRet
=\frac{\Cdeg}2-\alpha_{\angle,{\rm deg}}
=H_5-\frac{90}{\pi}\mathscr C_\pi(\alpha).}
\]
Sa coordonnée radiale est
\[
 \etaRetrad=\frac{\pi}{180}\etaRet
 =\frac{\Crad}{2}-\frac{\pi}{180}\alpha_{\rm HMT}.
\]
Numériquement,
\[
\etaRet
=1.05457181691503906113369058472430\ldots.
\]
L'identité inverse n'a pas été utilisée pour introduire \(\Cdeg\) ; elle s'obtient
après la génération de \(\Crad\). Le symbole \(\hbar\) reste réservé à
l'élément de la droite d'action construit dans la section suivante.

% HMT_SOURCE_SEGMENT_END id=A08
\subsection{Théorème direct \texorpdfstring{\(C^*=2(\eta_{\mathrm{ret}}+\alpha)\)}{C*=2(eta_ret+alpha)} et caractère inverse de l'identité \texorpdfstring{\(\eta_{\mathrm{ret}}=C^*/2-\alpha\)}{eta_ret=C*/2-alpha}}
\label{sec:s102:c38:4:3}

% HMT_SOURCE_SEGMENT_BEGIN id=A13 source=colaboracion/parte_iii/source/public_final/ascensos_20260826/16f_realizacion_analitica_contraangulo.tex body_lines=750-797

Le caractère de cinquième degré et la coordonnée postérieure au retour
sont deux valeurs distinctes :
\[
H_5
=1.05457181764615446962582182943306\ldots,
\]
\[
\etaRet
=1.05457181691503906113369058472430\ldots.
\]
Le premier est le caractère local avant l'intégration du caractère régional ; le
second est la coordonnée orientée qui subsiste après le retour. La correction
qui génère le contre-angle empêche de les identifier.

La comparaison métrologique doit respecter cette distinction. La première valeur donne
la mantisse
\[
2\pi H_5
=6.62607014999998792600111034989947\ldots,
\]
tandis que la coordonnée ultérieure donne
\[
2\pi\etaRet
=6.62607014540625433351074984759288\ldots.
\]
À titre de contrôle externe uniquement, nous retenons la mantisse \(6.62607015\) de la
valeur définitoire du Système international \cite{BIPMSI}. On n'intègre
pas encore son exposant ni son unité à la droite HMT : la décennie est dérivée dans le
\cref{chap:orden-determinantal-accion} et l'unité appartient à une carte
métrologique ultérieure. La comparaison des mantisses donne
\[
2\pi H_5-6.62607015
=-1.2073998889650101\ldots\times10^{-14},
\]
et ce n'est pas une égalité exacte. La proximité de la valeur antérieure au retour est
un contrôle numérique remarquable ; elle n'autorise pas à remplacer la constante définitoire
du SI par la coordonnée ultérieure ni à confondre une coordonnée angulaire avec
l'élément d'action qui sera construit lors de l'intégration de la décennie.

Ce résultat détermine avec précision la situation mathématique. Le lecteur
régional génère \((\Crad,\Cdeg)\) et, à partir de cet unique élément angulaire, la fonctionnelle
\(\Gamma_{6\to8}\). Le même calcul sépare la proximité externe de la
mantisse produite par \(H_5\) de la coordonnée de retour qui intervient dans le
contre-angle. La séparation n'affaiblit pas la construction : elle élimine une
identification incompatible avec les chiffres exacts et transforme chaque
comparaison en une affirmation falsifiable de type bien défini.

% HMT_SOURCE_SEGMENT_END id=A13
% HMT_SOURCE_SEGMENT_BEGIN id=D01 source=manuscrito/sucesor_102/deltas_ley9/c38_par_angular_cartas_y_canales.tex body_lines=3-20
\label{sec:l9s102:par-angular}

La complétion cubique
\begin{equation}
 10^3=(9+1)^3=9^3+3\cdot9^2+3\cdot9+1
 =729+243+27+1
 \label{eq:l9s102:base-1000}
\end{equation}
détermine la carte globale de la coordonnée directe. La clôture bicouche apporte
le facteur \(2\) de la coordonnée conjuguée :
\begin{equation}
 \boxed{
 A_{\deg}=10^3\alpha_{\mathrm{HMT}},
 \qquad
 C^*_{\deg}=2(\eta_{\mathrm{ret}}+\alpha_{\mathrm{HMT}})
 \quad\text{dans }\mathbb R/(360\mathbb Z).}
 \label{eq:l9s102:a-cstar-deg}
\end{equation}
% HMT_SOURCE_SEGMENT_END id=D01
% HMT_SOURCE_SEGMENT_BEGIN id=D04 source=manuscrito/sucesor_102/deltas_ley9/c38_par_angular_cartas_y_canales.tex body_lines=60-66
L'identité inverse
\begin{equation}
 \eta_{\mathrm{ret}}=\frac12C^*_{\deg}-\alpha_{\mathrm{HMT}}
 \label{eq:l9s102:eta-corolario-inverso}
\end{equation}
est établie comme corollaire algébrique de la construction directe.

% HMT_SOURCE_SEGMENT_END id=D04
\section{Commutativité exacte des cartes en tours, degrés et radians}
\label{sec:s102:c38:5}

\subsection{Dérivation de la longueur cylindrique et du degré analytique du précurseur parangulaire}
\label{sec:s102:c38:5:1}

% HMT_SOURCE_SEGMENT_BEGIN id=B03 source=colaboracion/parte_iii/source/public_final/base_83/c30_body.tex body_lines=157-194

Soit \(\mathcal P\) le module libre engendré par les mots régionaux et
soit \(F^n\mathcal P\) sa filtration par longueur. La graduation associée
enregistre le premier niveau auquel apparaît chaque mot. Pour un paramètre
\(z\), le caractère ordinaire de longueur est
\[
\chi_z([w])=z^{|w|}.
\]
La multiplicativité
\[
\chi_z([uv])=\chi_z([u])\chi_z([v])
\]
exprime que la concaténation additionne les longueurs. Lors de l'évaluation en
\(z=\alpha\), un mot de longueur \(n\) reçoit le degré \(\alpha^n\).

\begin{definition}[Compatibilité de degré]
\label{p3-83-c30-s1:def:compatibilidad-grado}
Le lecteur analytique régional est compatible avec la filtration cylindrique
lorsqu'il envoie la composante homogène de longueur \(n\) sur le degré analytique
\(n\) :
\[
\operatorname{gr}_n\mathcal P\longrightarrow
\alpha^n\mathbb R.
\]
\end{definition}

Les cinq régions appartiennent à \(\mathcal U_6\) ; c'est pourquoi la donnée de
retour apparaît en degré six :
\[
\rho_\pi[U_6]\longmapsto169\alpha^6.
\]
Ce six est la longueur du mot régional. Ce n'est pas la longueur des
marches fermées du \cref{chap:cinco-ciclos-pi}. Ces marches parcourent à
la fois les quatre ancrages
\(U_\pi,U_e,U_\varphi,U_{\rm APP}\) et possèdent une longueur minimale de huit. Un
mot, une arête du graphe de blocs et une marche fermée sont des objets de
types distincts ; la graduation utilisée ici appartient au premier.

% HMT_SOURCE_SEGMENT_END id=B03
\subsection{Publication du degré analytique dans la carte globale de phase}
\label{sec:s102:c38:5:2}

% HMT_SOURCE_SEGMENT_BEGIN id=A03 source=colaboracion/parte_iii/source/public_final/ascensos_20260826/16f_realizacion_analitica_contraangulo.tex body_lines=172-209

Soit \(\mathcal P\) le module libre engendré par les mots régionaux et
soit \(F^n\mathcal P\) sa filtration par longueur. La graduation associée
enregistre le premier niveau auquel apparaît chaque mot. Pour un paramètre
\(z\), le caractère ordinaire de longueur est
\[
\chi_z([w])=z^{|w|}.
\]
La multiplicativité
\[
\chi_z([uv])=\chi_z([u])\chi_z([v])
\]
exprime que la concaténation additionne les longueurs. Lors de l'évaluation en
\(z=\alpha\), un mot de longueur \(n\) reçoit le degré \(\alpha^n\).

\begin{definition}[Compatibilité de degré]
\label{p3-20260826:def:compatibilidad-grado}
Le lecteur analytique régional est compatible avec la filtration cylindrique
lorsqu'il envoie la composante homogène de longueur \(n\) sur le degré analytique
\(n\) :
\[
\operatorname{gr}_n\mathcal P\longrightarrow
\alpha^n\mathbb R.
\]
\end{definition}

Les cinq régions appartiennent à \(\mathcal U_6\) ; c'est pourquoi la donnée de
retour apparaît en degré six :
\[
\rho_\pi[U_6]\longmapsto169\alpha^6.
\]
Ce six est la longueur du mot régional. Ce n'est pas la longueur des
marches fermées du \cref{chap:cinco-ciclos-pi}. Ces marches parcourent à
la fois les quatre ancrages
\(U_\pi,U_e,U_\varphi,U_{\rm APP}\) et possèdent une longueur minimale de huit. Un
mot, une arête du graphe de blocs et une marche fermée sont des objets de
types distincts ; la graduation utilisée ici appartient au premier.

% HMT_SOURCE_SEGMENT_END id=A03
% HMT_SOURCE_SEGMENT_BEGIN id=D02 source=manuscrito/sucesor_102/deltas_ley9/c38_par_angular_cartas_y_canales.tex body_lines=21-34
Le calendrier global présente les trois factorisations
\begin{equation}
 360=12\cdot30=4\cdot90=3\cdot120.
 \label{eq:l9s102:360-factorizaciones}
\end{equation}
La carte analytique s'obtient au moyen de
\begin{equation}
 A_{\mathrm{rad}}=\frac\pi{180}A_{\deg}
 =\frac{50\pi}{9}\alpha_{\mathrm{HMT}},
 \qquad
 C^*_{\mathrm{rad}}=\frac\pi{180}C^*_{\deg}
 =\frac\pi{90}(\eta_{\mathrm{ret}}+\alpha_{\mathrm{HMT}}).
 \label{eq:l9s102:a-cstar-rad}
\end{equation}
% HMT_SOURCE_SEGMENT_END id=D02
% HMT_SOURCE_SEGMENT_BEGIN id=D03 source=manuscrito/sucesor_102/deltas_ley9/c38_par_angular_cartas_y_canales.tex body_lines=35-59

\begin{theorem}[Commutativité des factorisations régionale et d'action]
\label{thm:l9s102:conmutatividad-cstar}
Les deux chemins vers la phase conjuguée produisent la même sortie :
\begin{equation}
 \boxed{
 \frac\pi{180}C^*_{\deg}
 =\frac\pi{90}(\eta_{\mathrm{ret}}+\alpha_{\mathrm{HMT}})
 =\frac\pi{90}(H_5+\alpha_{\mathrm{HMT}})
 -\mathscr C_\pi(\alpha_{\mathrm{HMT}})=y_*.}
 \label{eq:l9s102:cuadrado-cstar}
\end{equation}
\end{theorem}

\begin{proof}
En substituant \cref{eq:l9s102:eta-ret} dans le second membre, on obtient
\[
 \frac\pi{90}
 \left(H_5-\frac{90}{\pi}\mathscr C_\pi+\alpha_{\mathrm{HMT}}\right)
 =\frac\pi{90}(H_5+\alpha_{\mathrm{HMT}})-\mathscr C_\pi.
\]
La dernière expression est la définition régionale de \(y_*\), tandis que la
première égalité procède de \cref{eq:l9s102:a-cstar-deg}.
\end{proof}

% HMT_SOURCE_SEGMENT_END id=D03
\section{Décomposition spectrale \texorpdfstring{\(q_\pm\)}{q+/-} et calcul fonctionnel sur les canaux \texorpdfstring{\(90/120\)}{90/120} préalablement générés}
\label{sec:s102:c38:6}

\subsection{Reconstruction spectrale du couple angulaire et transition hexade--octade}
\label{sec:s102:c38:6:1}

% HMT_SOURCE_SEGMENT_BEGIN id=B11 source=colaboracion/parte_iii/source/public_final/base_83/c30_body.tex body_lines=523-569

La phase commune et la phase orientée déterminent les deux canaux
\[
q_-=e^{-x+y_*},
\qquad
q_+=e^{-x-y_*}.
\]
On retrouve exactement les invariants
\[
q_-q_+=D_A,
\qquad
\alpha=-\frac9{100\pi}\log D_A,
\qquad
C_*^{\rm HMT}=\frac{90}{\pi}\log\frac{q_-}{q_+}.
\]
La première égalité lit la contraction commune ; la troisième, la séparation
orientée des feuillets. L'échange \(R\mapsto-R\) permute
\(q_+\) et \(q_-\), conserve \(D_A\) et inverse le signe de \(y_*\).

Pour
\[
S_{90}(q)=\frac{q^{90}}{1-q^{270}},
\qquad
S_{120}(q)=\frac{q^{120}}{1-q^{360}},
\]
la fonctionnelle de transition hexade--octade s'écrit intégralement dans la
carte des phases :
\[
\Gamma_{6\to8}
=\frac{180}{\pi}xy_*
+12\bigl(S_{90}(q_-)-S_{90}(q_+)\bigr)
-\bigl(S_{120}(q_-)-S_{120}(q_+)\bigr).
\]
L'évaluation produite par la coordonnée \(C_*^{\rm HMT}\) du
\cref{p3-83-c30-s1:thm:contraangulo-regional} est
\[
\boxed{
\Gamma_{6\to8}
=0.274007672694459274747255260774\ldots.}
\]
Ainsi, la valeur interne reconnue en mécanique dimensionnelle comme fonctionnelle de
Barbero--Immirzi cesse de recevoir \(C_*^{\rm HMT}\) comme littéral indépendant :
toutes deux descendent du même couple de phases engendré dans ce chapitre.
L'interprétation géométrique et physique de \(\Gamma_{6\to8}\) appartient à l'œuvre
de mécanique dimensionnelle ; son évaluation mathématique appartient déjà à l'interface
précédemment construite.

% HMT_SOURCE_SEGMENT_END id=B11
\subsection{Transporteur relatif et disjonction spectrale des blocs}
\label{sec:s102:c38:6:2}

% HMT_SOURCE_SEGMENT_BEGIN id=B13 source=colaboracion/parte_iii/source/public_final/base_83/c30_body.tex body_lines=619-675
\label{p3-83-c30-s1:sec:transportador-relativo-split}

Dans la base ordonnée des canaux, soit
\[
 R=\begin{pmatrix}0&1\\1&0\end{pmatrix},
 \qquad P_\pm=\frac{I\pm R}{2},
 \qquad
 B_0=7I+2R,
 \qquad B_1=\Adeg I+\Cdeg R.
\]
La décomposition scindée
\[
 xI+yR=\rho(x,y)e^{\eta(x,y)R},
 \quad \rho(x,y)=\sqrt{x^2-y^2},
 \quad \eta(x,y)=\operatorname{artanh}(y/x)
\]
construit l'unique transporteur
\[
 T_{\rm rel}
 =\frac{\sqrt{(\Adeg)^2-(\Cdeg)^2}}{\sqrt{45}}
   \exp\!\left[
   \left(\operatorname{artanh}\frac{\Cdeg}{\Adeg}
   -\operatorname{artanh}\frac27\right)R\right],
 \qquad T_{\rm rel}B_0=B_1.
\]
Dans la base spectrale, il agit par
\(9\mapsto\Adeg+\Cdeg\) et
\(5\mapsto\Adeg-\Cdeg\).

\begin{theorem}[Disjonction spectrale des réalisations locale et angulaire]
\label{p3-83-c30-s1:thm:no-go-entrelazador-bloques-split}
Les spectres
\[
 \operatorname{Spec}(B_0)=\{9,5\},\qquad
 \operatorname{Spec}(B_1)=\{\Adeg+\Cdeg,\Adeg-\Cdeg\}
\]
sont disjoints. En conséquence,
\[
 \operatorname{Hom}_{B_0,B_1}
 :=\{\Phi:\Phi B_0=B_1\Phi\}=\{0\}.
\]
\end{theorem}

\begin{proof}
Les bornes rationnelles construites pour les coordonnées angulaires donnent
\(7.29<\Adeg<7.30\) et \(2.12<\Cdeg<2.13\) ; donc,
\(9.41<\Adeg+\Cdeg<9.43\) et
\(5.16<\Adeg-\Cdeg<5.18\). Dans la base qui diagonalise \(R\), chaque entrée
d'un entrelaceur est multipliée par la différence entre une valeur propre
de \(B_0\) et une de \(B_1\) ; la disjonction impose qu'elles soient toutes nulles.
\end{proof}

Le transporteur relatif compare deux blocs au sein de l'algèbre scindée ;
les entrelaceurs dodécaphasique--Witt et d'incidence agissent entre des
représentations équivalentes dans leurs espaces de résidence propres. Cette séparation
type positivement les deux opérations et conserve leurs domaines.

% HMT_SOURCE_SEGMENT_END id=B13
\subsection{Évaluation spectrale \texorpdfstring{\(V_{90,120}\)}{V_90,120} des canaux d'incidence préalablement générés par le TPK}
\label{sec:s102:c38:6:3}

% HMT_SOURCE_SEGMENT_BEGIN id=A11 source=colaboracion/parte_iii/source/public_final/ascensos_20260826/16f_realizacion_analitica_contraangulo.tex body_lines=554-572

La phase commune et la phase orientée déterminent les deux canaux
\[
q_-=e^{-\Arad+\Crad},
\qquad
q_+=e^{-\Arad-\Crad}.
\]
On retrouve exactement les invariants
\[
q_-q_+=D_A,
\qquad
\alpha=-\frac9{100\pi}\log D_A,
\qquad
\Cdeg=\frac{90}{\pi}\log\frac{q_-}{q_+}.
\]
La première égalité lit la contraction commune ; la troisième, la séparation
orientée des feuillets. L'échange \(R\mapsto-R\) permute
\(q_+\) et \(q_-\), conserve \(D_A\) et inverse le signe de \(\Crad\).

% HMT_SOURCE_SEGMENT_END id=A11
% HMT_SOURCE_SEGMENT_BEGIN id=D05 source=manuscrito/sucesor_102/deltas_ley9/c38_par_angular_cartas_y_canales.tex body_lines=68-97
\label{sec:l9s102:qpm-posteriores}

Soient \(H^2=I\) et \(P_\pm=(I\pm H)/2\). Le bloc angulaire
\begin{equation}
 B=A_{\mathrm{rad}}I+C^*_{\mathrm{rad}}H
 =(A_{\mathrm{rad}}+C^*_{\mathrm{rad}})P_+
 +(A_{\mathrm{rad}}-C^*_{\mathrm{rad}})P_-
 \label{eq:l9s102:b-angular}
\end{equation}
se publie multiplicativement par
\begin{equation}
 e^{-B}=q_+P_++q_-P_-,
 \qquad
 q_\pm=\exp[-(A_{\mathrm{rad}}\pm C^*_{\mathrm{rad}})].
 \label{eq:l9s102:qpm-angular}
\end{equation}
L'involution \(C^*\mapsto-C^*\) permute \(q_+\) et \(q_-\), et conserve
\begin{equation}
 q_+q_-=e^{-2A_{\mathrm{rad}}}=D_A.
 \label{eq:l9s102:det-angular}
\end{equation}
Les cardinaux \(90/120\), construits préalablement par l'incidence du
TPK, reçoivent maintenant leur évaluation spectrale :
\begin{equation}
 V_{90,120}
 =(I-T^{90})(I-T^{120})^{-1}
 =\sum_{\sigma\in\{+,-\}}
 \frac{1-q_\sigma^{90}}{1-q_\sigma^{120}}P_\sigma.
 \label{eq:l9s102:v90120-posterior}
\end{equation}
% HMT_SOURCE_SEGMENT_END id=D05
\subsubsection{Transporteur relatif scindé}
\label{sec:s102:c38:6:3:1}

% HMT_SOURCE_SEGMENT_BEGIN id=A12 source=colaboracion/parte_iii/source/public_final/ascensos_20260826/16f_realizacion_analitica_contraangulo.tex body_lines=574-748
\label{p3-20260826:sec:transportador-relativo-split}

La comparaison avec le bloc central d'APP peut s'effectuer sans identifier
des représentations distinctes. Fixons la base ordonnée des canaux et
l'involution
\[
 R=\begin{pmatrix}0&1\\1&0\end{pmatrix},
 \qquad R^2=I,
 \qquad
 P_\pm=\frac{I\pm R}{2}.
\]
Dans l'algèbre commutative
\[
 \mathscr B_R:=\{xI+yR:x,y\in\mathbb R\}
\]
considérons, dans la carte sexagésimale déjà déclarée,
\[
 B_0:=7I+2R,
 \qquad
 B_1:=\Adeg I+\Cdeg R.
\]
Le changement orthogonal de base
\[
 S=\frac1{\sqrt2}
 \begin{pmatrix}1&1\\1&-1\end{pmatrix}
\]
diagonalise simultanément l'involution et les deux blocs. Pour \(x>|y|\),
définissons
\[
 \rho(x,y):=\sqrt{x^2-y^2},
 \qquad
 \eta(x,y):=\operatorname{artanh}\!\left(\frac yx\right).
\]
Alors
\[
 xI+yR=\rho(x,y)e^{\eta(x,y)R}.
\]
Si
\[
 \rho_0=\sqrt{45},\quad
 \eta_0=\operatorname{artanh}(2/7),\quad
 \rho_1=\rho(\Adeg,\Cdeg),\quad
 \eta_1=\eta(\Adeg,\Cdeg),
\]
l'unique multiplication à gauche qui envoie \(B_0\) sur \(B_1\) est
\[
 \boxed{
 T_{\rm rel}
 :=\frac{\rho_1}{\sqrt{45}}
 e^{(\eta_1-\eta_0)R},
 \qquad
 T_{\rm rel}B_0=B_1.}
\]

\begin{proposition}[Transporteur relatif des deux blocs]
\label{p3-20260826:prop:transportador-relativo-split}
Une fois fixés \(R\), l'ordre des canaux, la branche positive et la carte
sexagésimale, \(T_{\rm rel}\) est l'unique élément de
\(\mathscr B_R\) qui satisfait \(TB_0=B_1\). Dans la base spectrale, il agit par
\[
 9\longmapsto\Adeg+\Cdeg,
 \qquad
 5\longmapsto\Adeg-\Cdeg.
\]
\end{proposition}

\begin{proof}
La décomposition polaire scindée donne la formule imprimée et prouve l'existence.
Comme \(\det B_0=45\ne0\), toute solution matricielle de \(TB_0=B_1\) satisfait
\(T=B_1B_0^{-1}\) ; elle est donc unique. La diagonalisation au moyen de \(S\)
produit les deux quotients spectraux
\((\Adeg+\Cdeg)/9\) et \((\Adeg-\Cdeg)/5\), équivalents à l'expression
exponentielle de \(T_{\rm rel}\).
\end{proof}

Ce transporteur n'est pas un entrelaceur représentationnel. En effet, un
entrelaceur \(\Phi\) entre les endomorphismes \(B_0\) et \(B_1\) devrait
satisfaire
\[
 \Phi B_0=B_1\Phi.
\]

\begin{theorem}[Obstruction spectrale à l'entrelacement]
\label{p3-20260826:thm:no-go-entrelazador-bloques-split}
Pour les valeurs de \(\Adeg\) et \(\Cdeg\) construites dans cette section,
\[
 \boxed{\Phi B_0=B_1\Phi\quad\Longrightarrow\quad\Phi=0.}
\]
\end{theorem}

\begin{proof}
Les spectres sont
\[
 \operatorname{Spec}(B_0)=\{9,5\},
 \qquad
 \operatorname{Spec}(B_1)
 =\{\Adeg+\Cdeg,\Adeg-\Cdeg\}.
\]
Les bornes rationnelles préalablement démontrées fournissent les intervalles
disjoints
\[
 7.29<\Adeg<7.30,
 \qquad
 2.12<\Cdeg<2.13,
\]
et, par conséquent,
\[
 9.41<\Adeg+\Cdeg<9.43,
 \qquad
 5.16<\Adeg-\Cdeg<5.18.
\]
Dans la base qui diagonalise \(R\), chaque entrée \(\phi_{ij}\) de \(\Phi\)
est multipliée par la différence entre une valeur propre de \(B_0\) et une de
\(B_1\). Toutes ces différences sont non nulles ; donc toutes les entrées de
\(\Phi\) s'annulent.
\end{proof}

Numériquement,
\[
 \rho_1=6.9814819335172378048\ldots,
 \qquad
 \frac{\rho_1}{\sqrt{45}}=1.0407378791354140779\ldots,
\]
\[
 \eta_1-\eta_0
 =0.005796351470571295590\ldots.
\]
La formule conserve l'algèbre scindée, les projecteurs \(P_\pm\) et la forme
lorentzienne à un facteur conforme \(\rho_1^2/45\) près. Elle dépend cependant
de \(B_1\) déjà construit : elle ne constitue pas une dérivation antérieure de
\(\Adeg\) ou \(\Cdeg\) et ne doit pas être confondue avec les entrelaceurs
équivariants entre réalisations d'un même module.

Pour
\[
S_{90}(q)=\frac{q^{90}}{1-q^{270}},
\qquad
S_{120}(q)=\frac{q^{120}}{1-q^{360}},
\]
la fonctionnelle de transition hexade--octade s'écrit intégralement dans la
carte des phases :
\[
\Gamma_{6\to8}
=\frac{180}{\pi}\Arad\Crad
+12\bigl(S_{90}(q_-)-S_{90}(q_+)\bigr)
-\bigl(S_{120}(q_-)-S_{120}(q_+)\bigr).
\]
L'évaluation produite par le contre-angle du
\cref{p3-20260826:thm:contraangulo-regional} est
\[
\boxed{
\Gamma_{6\to8}
=0.274007672694459274747255260774\ldots.}
\]
Les formes équivalentes de la fonctionnelle mathématique interne sont
\[
\boxed{
\Gamma_{6\to8}
=\Arad\Cdeg+K_{6\to8}
=\frac{180}{\pi}\Arad\Crad+K_{6\to8}
=\frac{\pi\Adeg\Cdeg}{180}+K_{6\to8}.}
\]
Ces identités évaluent \(\Gamma_{6\to8}\) après la génération de
\((\Crad,\Cdeg)\) ; elles ne sont pas utilisées en sens inverse pour sélectionner le
contre-angle. La lecture physique dans la connexion d'Ashtekar--Barbero est
représentée ensuite au moyen de
\[
 \Gamma_{6\to8}\dashrightarrow\gamma_{\rm BI}.
\]
La flèche n'introduit pas de troisième fonctionnelle et ne rétroagit pas sur \(C^*\). Les
séries \(S_{90}\), \(S_{120}\), les coefficients \(12\) et \(-1\) et leur
affectation au pas \(6\to8\) appartiennent à la construction mathématique
déclarée. L'identification physique exige en outre la normalisation propre
de la connexion.

% HMT_SOURCE_SEGMENT_END id=A12
\section{Naturalité sous raffinement et séparation d'avec la reconnaissance métrologique}
\label{sec:s102:c38:7}

\subsection{Estimation exacte du reste de la prolongation analytique}
\label{sec:s102:c38:7:1}

% HMT_SOURCE_SEGMENT_BEGIN id=B09 source=colaboracion/parte_iii/source/public_final/base_83/c30_body.tex body_lines=474-499

L'expression initialement trouvée,
\[
y_7=\frac\pi{90}(H_5+\alpha)-169\alpha^6-D_A\alpha^7,
\]
est la troncature de degré sept du caractère complet. Elle produit
\[
C_7=2.12373833896864600265403827103919\ldots{}^\circ.
\]
Le reste omis n'est pas un terme indéterminé ; c'est la queue géométrique
exacte
\[
\mathscr C_\pi(\alpha)
-169\alpha^6-D_A\alpha^7
=\frac{D_A^2\alpha^8}{1-D_A\alpha}.
\]
Par conséquent,
\[
\left|C_7-C_*^{\rm HMT}\right|
=\frac{180}{\pi}\frac{D_A^2\alpha^8}{1-D_A\alpha}
=2.7839208689064583\ldots\times10^{-16}\;{}^\circ.
\]
La récurrence transforme ainsi une approximation de septième degré en une
réalisation analytique complète et fournit une borne close pour chaque
troncature ultérieure.

% HMT_SOURCE_SEGMENT_END id=B09
\subsection{Dépendance fonctionnelle du sélecteur régional et du prolongement analytique}
\label{sec:s102:c38:7:2}

% HMT_SOURCE_SEGMENT_BEGIN id=B10 source=colaboracion/parte_iii/source/public_final/base_83/c30_body.tex body_lines=501-521

La construction change lorsque l'une quelconque de ses composantes actives est retirée.
Les valeurs suivantes s'obtiennent sans réétalonner les autres termes :
\[
\begin{array}{@{}lc@{}}
\toprule
\text{variante}&\text{coordonnée angulaire conjuguée en degrés}\\
\midrule
\rho_\pi=168&2.1237383389772976863904487\ldots\\
\rho_\pi=169&2.1237383389686457242619514\ldots\\
\rho_\pi=170&2.1237383389599937621334541\ldots\\
\text{sans queue déterminantale}&2.1237383389686949415301675\ldots\\
D_A\text{ remplacé par }1&2.1237383389686313409965826\ldots\\
\bottomrule
\end{array}
\]
Ces ablations ne constituent pas à elles seules une estimation probabiliste.
Elles démontrent que le sélecteur équivariant et le caractère déterminantal sont
des composantes effectives de la sortie et qu'ils ne peuvent être supprimés tout en maintenant la
même réalisation.

% HMT_SOURCE_SEGMENT_END id=B10
\subsection{Interprétation opératoire de la coordonnée angulaire conjuguée}
\label{sec:s102:c38:7:3}

% HMT_SOURCE_SEGMENT_BEGIN id=B14 source=colaboracion/parte_iii/source/public_final/base_83/c30_body.tex body_lines=677-715

La procédure complète peut être résumée par une chaîne d'applications :
\[
\begin{aligned}
\mathcal F_\pi
&\longmapsto b_\pi
\longmapsto\rho_\pi=169,\\
\alpha
&\longmapsto A
\longmapsto x
\longmapsto D_A,\\
(\rho_\pi,D_A)
&\longmapsto\mathscr C_\pi(\alpha),\\
(H_5,\mathscr C_\pi(\alpha))
&\longmapsto y_*
\longmapsto C_*^{\rm HMT},\\
(x,y_*)
&\longmapsto(q_+,q_-)
\longmapsto\Gamma_{6\to8}.
\end{aligned}
\]
La première ligne appartient au catalogue fini ; la deuxième, à la clôture commune de
\(\alpha\) ; la troisième, à la prolongation déterminantale ; la quatrième produit la
coordonnée orientée ; la cinquième la transporte vers l'interface dimensionnelle.

Dans cette organisation, \(\alpha\) mesure la contraction commune du transport et
\(C_*^{\rm HMT}\) mesure la séparation orientée de ses deux feuillets. Le retour
régional n'ajoute pas de constante externe : il sélectionne une amplitude finie et la
prolonge au moyen du déterminant déjà fixé par \(\alpha\). La carte en
degrés est ultérieure. L'objet primaire est le couple de phases \((x,y_*)\).

Une vérification exhaustive reconstruit les cinq régions directement
à partir du catalogue, parcourt les actions de
\(\mathfrak S_5\) et \(\mathfrak S_3\), vérifie la récurrence et somme la
série. L'entrée de \(\alpha\) est la sortie exacte de sa clôture
dodécaphasique, non un chiffre métrologique transcrit. La vérification maintient
séparés les théorèmes du catalogue, les règles du lecteur et la comparaison
externe ; c'est pourquoi aucune comparaison expérimentale n'intervient dans la fonction
génératrice.
% HMT_SOURCE_SEGMENT_END id=B14
\subsection{Certification de convergence de la troncature parangulaire}
\label{sec:s102:c38:7:4}

% HMT_SOURCE_SEGMENT_BEGIN id=A09 source=colaboracion/parte_iii/source/public_final/ascensos_20260826/16f_realizacion_analitica_contraangulo.tex body_lines=505-530

L'expression initialement trouvée,
\[
y_7=\frac\pi{90}(H_5+\alpha)-169\alpha^6-D_A\alpha^7,
\]
est la troncature de degré sept du caractère complet. Elle produit
\[
C_7=2.12373833896864600265403827103919\ldots{}^\circ.
\]
Le reste omis n'est pas un terme indéterminé ; c'est la queue géométrique
exacte
\[
\mathscr C_\pi(\alpha)
-169\alpha^6-D_A\alpha^7
=\frac{D_A^2\alpha^8}{1-D_A\alpha}.
\]
Par conséquent,
\[
\left|C_7-\Cdeg\right|
=\frac{180}{\pi}\frac{D_A^2\alpha^8}{1-D_A\alpha}
=2.7839208689064583\ldots\times10^{-16}\;{}^\circ.
\]
La récurrence transforme ainsi une approximation de septième degré en une
réalisation analytique complète et fournit une borne close pour chaque
troncature ultérieure.

% HMT_SOURCE_SEGMENT_END id=A09
\subsection{Naturalité du sélecteur régional sous raffinement et changement de carte}
\label{sec:s102:c38:7:5}

% HMT_SOURCE_SEGMENT_BEGIN id=A10 source=colaboracion/parte_iii/source/public_final/ascensos_20260826/16f_realizacion_analitica_contraangulo.tex body_lines=532-552

La construction change lorsque l'une quelconque de ses composantes actives est retirée.
Les valeurs suivantes s'obtiennent sans réétalonner les autres termes :
\[
\begin{array}{@{}lc@{}}
\toprule
\text{variante}&\text{contre-angle en degrés}\\
\midrule
\rho_\pi=168&2.1237383389772976863904487\ldots\\
\rho_\pi=169&2.1237383389686457242619514\ldots\\
\rho_\pi=170&2.1237383389599937621334541\ldots\\
\text{sans queue déterminantale}&2.1237383389686949415301675\ldots\\
D_A\text{ remplacé par }1&2.1237383389686313409965826\ldots\\
\bottomrule
\end{array}
\]
Ces ablations ne constituent pas à elles seules une estimation probabiliste.
Elles démontrent que le sélecteur équivariant et le caractère déterminantal sont
des composantes effectives de la sortie et qu'ils ne peuvent être supprimés tout en maintenant la
même réalisation.

% HMT_SOURCE_SEGMENT_END id=A10
\subsection{Indépendance métrologique de la construction parangulaire}
\label{sec:s102:c38:7:6}

% HMT_SOURCE_SEGMENT_BEGIN id=A14 source=colaboracion/parte_iii/source/public_final/ascensos_20260826/16f_realizacion_analitica_contraangulo.tex body_lines=799-831

La procédure complète peut être résumée par une chaîne d'applications :
\[
\begin{aligned}
\mathcal F_\pi
&\longmapsto b_\pi
\longmapsto\rho_\pi=169,\\
\alpha
&\longmapsto \Adeg
\longmapsto \Arad
\longmapsto D_A,\\
(\rho_\pi,D_A)
&\longmapsto\mathscr C_\pi(\alpha),\\
(H_5,\mathscr C_\pi(\alpha))
&\longmapsto \Crad
\longmapsto \Cdeg,\\
(\Arad,\Crad)
&\longmapsto(q_+,q_-)
\longmapsto\Gamma_{6\to8}.
\end{aligned}
\]
La première ligne appartient au catalogue fini ; la deuxième, à la clôture commune de
\(\alpha\) ; la troisième, à la prolongation déterminantale ; la quatrième produit la
coordonnée orientée ; la cinquième la transporte vers l'interface dimensionnelle.

Dans cette organisation, \(\alpha\) mesure la contraction commune du transport et
le contre-angle mesure la séparation orientée de ses deux feuillets. Dans
l'évaluation en vigueur, le retour régional sélectionne une amplitude finie et la
prolonge au moyen du déterminant déjà fixé par \(\alpha\), sans insérer de
constante externe supplémentaire. La carte en degrés est ultérieure. L'objet
premier est le couple de phases \((\Arad,\Crad)\). La récurrence, la somme
régionale et les ablations dépendent du catalogue et de
\(\alpha_{\rm HMT}\), mais ne reçoivent pas de constante métrologique supplémentaire.
% HMT_SOURCE_SEGMENT_END id=A14



\HMTFlushCurrentChapter
\chapter{Opérateur constitutif positif du vide : complétion \(3\to4\), reconstruction spectrale et dérivation conjointe de \texorpdfstring{\(\varepsilon,\mu,Z\)}{epsilon, mu, Z} et \(c\)}%
\begingroup
\renewcommand{\chapter}[2][]{}%
\chapter{Opérateur constitutif positif du vide : complétion \(3\to4\), reconstruction spectrale et dérivation conjointe de \(\varepsilon\), de \(\mu\), de \(Z\) et de \(c\)}
\label{chap:p3-final:39:vacio_constitutivo}
\section{Décomposition spectrale bicouche de l'opérateur angulaire}

Soit \(R=R^*=R^{-1}\) l'involution active et

\[
P_\pm=\frac{I\pm R}{2},
\qquad
T=q_+P_+ +q_-P_-,
\qquad 0<q_+<q_-<1.
\]

Avec
\[
x=A_{\rm rad}=\frac{\pi A}{180},\qquad
y=C^*_{\rm rad}=\frac{\pi C^*}{180},
\]
les deux canaux sont construits avant d'évaluer le vide :
\begin{equation}
q_+=e^{-x-y},\qquad q_-=e^{-x+y}.
\label{eq:vacuum-angular-channels}
\end{equation}
L'échange de feuillets inverse \(C^*\), c'est-à-dire \(y\mapsto-y\), et permute \(q_+\leftrightarrow q_-\). Le couple est ainsi la lecture spectrale d'un même antécédent angulaire, et non deux paramètres ajustés séparément.

\begin{definition}[Réponse constitutive]
\begin{equation}
\mathcal V_{90,120}=(I-T^{90})(I-T^{120})^{-1}.
\label{eq:vacuum-operator}
\end{equation}
\end{definition}

Comme \(\|T\|<1\), le dénominateur est positif et inversible. Le calcul fonctionnel donne

\begin{equation}
\mathcal V_{90,120}=r_+P_+ +r_-P_-,
\qquad
r_\pm=\frac{1-q_\pm^{90}}{1-q_\pm^{120}}.
\label{eq:vacuum-eigenvalues}
\end{equation}

L'affirmation d'inversibilité mérite d'être explicitée. Pour chaque feuillet,
\(0<q_\pm^{120}<1\) ; par conséquent
\[
I-T^{120}=(1-q_+^{120})P_++(1-q_-^{120})P_-
\]
est strictement positif et
\[
(I-T^{120})^{-1}
=\frac{P_+}{1-q_+^{120}}+
 \frac{P_-}{1-q_-^{120}}.
\]
La réponse \(\mathcal V_{90,120}\) ne s'obtient pas en résolvant quatre
équations scalaires séparées : elle est le résultat du calcul fonctionnel d'une
seule contraction positive. Cette observation empêche que perméabilité,
permittivité, impédance et propagation perdent l'antécédent qu'elles partagent.

\begin{theorem}[Détermination spectrale conjointe des quatre relations constitutives]
Les quatre caractères constitutifs
\begin{equation}
\widehat\varepsilon=r_+^2,\qquad
\widehat\mu=r_-^2,\qquad
\widehat Z=\frac{r_-}{r_+},\qquad
\widehat c=\frac1{r_+r_-}
\label{eq:vacuum-four}
\end{equation}
sont des fonctions spectrales de \(\mathcal V_{90,120}\). Ils satisfont
\begin{equation}
\widehat\mu\widehat\varepsilon=\widehat c^{-2},
\qquad
\frac{\widehat\mu}{\widehat\varepsilon}=\widehat Z^2.
\label{eq:vacuum-relations}
\end{equation}
\end{theorem}

\begin{proof}
Les identités s'obtiennent en substituant \cref{eq:vacuum-four}. Aucune carte dimensionnelle ni aucune valeur extérieure n'intervient.
\end{proof}

\begin{corollary}[Reconstruction des deux feuillets]
Les quatre caractères ne contiennent pas quatre degrés de liberté. N'importe lequel
des couples \((\widehat\mu,\widehat\varepsilon)\),
\((\widehat Z,\widehat c)\) ou \((\mathcal E,\mathcal B)\) reconstruit les
deux valeurs propres positives :
\begin{align}
r_-&=\sqrt{\widehat\mu}
=\sqrt{\widehat Z/\widehat c},
&r_+&=\sqrt{\widehat\varepsilon}
=\frac1{\sqrt{\widehat Z\widehat c}},
\label{eq:vacuum-reconstruct-r}\\
\log r_-&=\frac{\mathcal E+\mathcal B}{2},
&\log r_+&=\frac{\mathcal E-\mathcal B}{2}.
\label{eq:vacuum-reconstruct-log}
\end{align}
En particulier, l'orientation se perd si l'on conserve uniquement le produit,
mais se retrouve en ajoutant le quotient.
\end{corollary}

\begin{proof}
La première ligne est la racine positive des identités
\cref{eq:vacuum-four,eq:vacuum-relations}. La seconde inverse le système
linéaire
\(\mathcal E=\log r_-+\log r_+\),
\(\mathcal B=\log r_--\log r_+\).
\end{proof}

Ce corollaire exprime une propriété centrale de HMT : produit et quotient sont
des fonctionnelles complémentaires d'un même état. Le produit détermine la
propagation commune, tandis que le quotient conserve l'orientation
relative des feuillets intérieur et extérieur.

Dans le certificat V2,

\begin{align*}
r_+&=0.9999996285482493631786952281\ldots,\\
r_-&=0.9997241371895183641315165853\ldots,\\
\widehat\varepsilon&=0.9999992570966367027604416155\ldots,\\
\widehat\mu&=0.9994483504793269350899815559\ldots,\\
\widehat Z&=0.9997245085389372154555543466\ldots,\\
\widehat c&=1.0002763104861575261063862689\ldots.
\end{align*}

Ces numéraux sont la reconnaissance terminale de \cref{eq:vacuum-operator} ; ils ne définissent pas ses exposants \(90,120\).

\section{Factorisation harmonique des secteurs \(90\) et \(120\)}

Le semi-groupe harmonique est

\[
\langle90,120\rangle=30\langle3,4\rangle.
\]

Si \(s=q^{30}\), alors

\begin{equation}
\frac{1-q^{90}}{1-q^{120}}
=\frac{1-s^3}{1-s^4}
=\frac{1+s+s^2}{1+s+s^2+s^3}.
\label{eq:three-four-completion}
\end{equation}

L'identité sépare également la réponse de son défaut. Si
\[
F_{3\to4}(s)=\frac{1+s+s^2}{1+s+s^2+s^3},
\qquad 0<s<1,
\]
alors
\begin{equation}
d_{3\to4}(s):=1-F_{3\to4}(s)
=\frac{s^3}{1+s+s^2+s^3}>0.
\label{eq:vacuum-defect34}
\end{equation}
Le terme supplémentaire est conservé exactement comme défaut
positif. En outre,
\begin{equation}
F_{3\to4}'(s)
=-\frac{s^2(3+2s+s^2)}{(1+s+s^2+s^3)^2}<0.
\label{eq:vacuum-monotonic34}
\end{equation}
Par conséquent, l'ordre des canaux angulaires détermine univoquement
l'ordre des réponses constitutives. L'attribution de
\(\widehat\mu\) et \(\widehat\varepsilon\) aux feuillets respectifs s'obtient
analytiquement, sans recourir à leurs développements décimaux.

Chaque feuillet du vide constitue la complétion exacte de trois à quatre
termes d'un même mode élémentaire \(30\). En conséquence, perméabilité
et permittivité s'interprètent comme deux fonctionnelles quadratiques d'une même
complétion évaluée sur des feuillets conjugués, et non comme des paramètres
constitutifs indépendants.

Comme la fonction \((1-q^{90})/(1-q^{120})\) décroît dans \(0<q<1\) et \(q_+<q_-\), on obtient

\[
0<r_-<r_+<1,\qquad
0<\widehat\mu<\widehat\varepsilon<1,\qquad
0<\widehat Z<1<\widehat c.
\]

Les limites de la fonctionnelle déterminent sa géométrie. Lorsque \(q\to0^+\),
\(F(q)\to1\) : les quatre lectures s'approchent de l'unité et le défaut de
mémoire disparaît. Lorsque \(q\to1^-\),
\[
F(q)\longrightarrow\frac{90}{120}=\frac34.
\]
La réponse prend donc ses valeurs dans l'intervalle positif \((3/4,1)\). La
borne \(3/4\) n'est pas une constante ajustée : elle est le quotient des deux
longueurs de canal et la réalisation continue de la complétion
trois--à--quatre.

\section{Involution des feuillets et conservation de l'orientation}

L'échange de feuillets \(C^*\mapsto-C^*\) agit par

\begin{equation}
\widehat\mu\longleftrightarrow\widehat\varepsilon,\qquad
\widehat Z\longmapsto\widehat Z^{-1},
\qquad
\widehat c\longmapsto\widehat c.
\label{eq:vacuum-sheet}
\end{equation}

La propagation commune est paire ; l'impédance conserve l'orientation.
Cette distinction réapparaîtra dans Barbero et CKM : toutes les fonctionnelles
associées au changement de feuillet n'ont pas la même parité.

\section{Caractères logarithmiques et compatibilité traciale}

Définissons

\[
\mathcal E=\log(r_+r_-),\qquad
\mathcal B=\log(r_-/r_+).
\]

Avec la trace normalisée \(\tau=\Tr/2\),

\begin{equation}
2\tau(\log\mathcal V)=\mathcal E,
\qquad
-2\tau(R\log\mathcal V)=\mathcal B.
\label{eq:vacuum-traces}
\end{equation}

Le couple \((\mathcal E,\mathcal B)\) constitue un système de coordonnées linéaire du logarithme
opératoriel :
\begin{equation}
\log\mathcal V
=\frac{\mathcal E}{2}I-\frac{\mathcal B}{2}R.
\label{eq:vacuum-log-decomposition}
\end{equation}
La composante scalaire est invariante sous l'échange de feuillets ; la
composante orientée change de signe. Cette décomposition démontre
directement la parité de la propagation et le caractère orienté de
l'impédance.

L'amplification nonadique

\[
\mathcal V_m=\mathcal V\otimes I_{9^m}
\]

est compatible avec les espérances \(E_m=\mathrm{id}\otimes\tau_9\) :

\[
E_m(\mathcal V_{m+1})=\mathcal V_m.
\]

Par conséquent, \(\mathcal E\) et \(\mathcal B\) ne sont pas des particularités d'une matrice \(2\times2\) ; ils subsistent dans la tour UHF et dans sa clôture traciale \(II_1\).

Plus précisément, pour tout polynôme \(p\) puis, par continuité,
pour toute fonction continue sur le spectre,
\begin{equation}
E_m\bigl(p(\mathcal V_{m+1})\bigr)=p(\mathcal V_m).
\label{eq:vacuum-functional-refinement}
\end{equation}
La compatibilité ne préserve pas seulement deux nombres : elle préserve toute
l'algèbre fonctionnelle engendrée par la réponse. En particulier, elle conserve ses
projecteurs spectraux, ses puissances, son logarithme et les formes quadratiques
qui mesurent les défauts \cref{eq:vacuum-defect34}.

\section{Réalisation dimensionnelle des relations constitutives}

La mécanique dimensionnelle détermine

\begin{align}
Z_{\rm vac}&=\widehat Z\,u_Su_Q^{-2},\\
c_{\rm vac}&=\widehat c\,u_C,\\
\mu_{\rm vac}&=\widehat\mu\,u_Su_C^{-1}u_Q^{-2},\\
\varepsilon_{\rm vac}&=\widehat\varepsilon\,u_S^{-1}u_C^{-1}u_Q^2.
\label{eq:vacuum-sections}
\end{align}

Les identités constitutives se maintiennent dans les droites :

\begin{equation}
\mu_{\rm vac}\varepsilon_{\rm vac}=c_{\rm vac}^{-2},
\qquad
\mu_{\rm vac}/\varepsilon_{\rm vac}=Z_{\rm vac}^2.
\end{equation}

La carte SI est postérieure. Il n'est pas nécessaire de déclarer \(c\) primitif et de résoudre ensuite \(\mu\varepsilon=c^{-2}\) : les trois objets descendent conjointement de \(\mathcal V\) et des droites déclarées.

La séparation des niveaux s'exprime par le diagramme
\[
\begin{array}{c}
(A,C^*)\longrightarrow T\longrightarrow\mathcal V
\longrightarrow
(\widehat\mu,\widehat\varepsilon,\widehat Z,\widehat c)
\\[2mm]
\Big\downarrow\text{ cartes de droite}\hspace{36mm}
\Big\downarrow\text{ évaluation finale}
\\[2mm]
(u_S,u_Q,u_C)\longrightarrow
(\mu_{\rm vac},\varepsilon_{\rm vac},Z_{\rm vac},c_{\rm vac})
\longrightarrow\text{ coordonnées SI}.
\end{array}
\]
La première ligne est adimensionnelle et opératorielle ; la deuxième détermine les types de
grandeur ; la troisième sélectionne une carte métrologique. Modifier la carte finale
ne modifie ni les valeurs propres ni les identités constitutives. En ce
sens précis, l'unification mathématique--physique consiste à distinguer
l'invariant généalogique de son système de coordonnées, et non à supprimer la
structure dimensionnelle.

\begin{proposition}[Minimalité spectrale]
Soit \(W\) un opérateur positif à deux feuillets de spectre simple
\(\{r_+,r_-\}\). Si quatre caractères positifs satisfont
\(\mu\varepsilon=c^{-2}\) et \(\mu/\varepsilon=Z^2\), si la propagation
commune est déterminée par \(c^{-1}=\det W=r_-r_+\), et si son orientation
fixe \(Z=r_-/r_+\), alors nécessairement
\[
(\mu,\varepsilon,Z,c)
=\left(r_-^2,r_+^2,\frac{r_-}{r_+},\frac1{r_-r_+}\right).
\]
Par conséquent, \cref{eq:vacuum-four} n'est pas une paramétrisation parmi d'autres :
c'est la réalisation positive minimale compatible avec le produit, le quotient et
l'orientation.
\end{proposition}

\begin{proof}
De \(Z=r_-/r_+\) et \(c^{-1}=\det W=r_-r_+\), on obtient de manière unique les
racines positives. Les deux identités constitutives imposent ensuite
\(\mu=r_-^2\) et \(\varepsilon=r_+^2\).
\end{proof}


La première inversion conceptuelle consiste à concevoir le vide comme une réponse relationnelle positive.\par

Dans HMT, le vide est la réponse dynamique minimale du système lorsque ses deux feuillets doivent coexister, se propager et clore leur différence en conservant l’orientation qui les distingue.\par

Les deux feuillets constituent deux lectures conjuguées du même antécédent angulaire. L’un conserve une orientation et l’autre conserve l’orientation opposée. C’est pourquoi apparaissent deux canaux, \(q_+\) et \(q_-\), comme les deux faces spectrales corrélées d’un même état.\par

L'opérateur\par

\[
\mathcal V_{90,120}
=
(I-T^{90})(I-T^{120})^{-1}
\]

compare deux profondeurs de propagation : \(90\) et \(120\). Ces profondeurs sont déterminées par le même mode élémentaire :\par

\[
90=3\cdot30,
\qquad
120=4\cdot30.
\]

C’est pourquoi la réponse de chaque feuillet peut se réécrire comme\par

\[
\frac{1-q^{90}}{1-q^{120}}
=
\frac{1+s+s^2}{1+s+s^2+s^3},
\qquad
s=q^{30}.
\]

Et c’est ici que l’interprétation devient réellement puissante. Le numérateur contient trois états du mode \(30\) ; le dénominateur contient ces mêmes trois états et un de plus. Le vide est donc une complétion exacte \(3\to4\).\par

Ce quatrième terme intègre et conserve les trois précédents. Il les complète. Ce qui est ajouté est enregistré comme un résidu positif de complétion :\par

\[
1-\frac{1+s+s^2}{1+s+s^2+s^3}
=
\frac{s^3}{1+s+s^2+s^3}>0.
\]

La mémoire du vide est précisément l’accroissement exact qui permet à trois termes de devenir quatre. Elle constitue l’empreinte positive et nécessaire de la complétion.\par

Des deux valeurs propres de cet unique opérateur découlent les quatre relations constitutives :\par

\[
\widehat\varepsilon=r_+^2,
\qquad
\widehat\mu=r_-^2,
\qquad
\widehat Z=\frac{r_-}{r_+},
\qquad
\widehat c=\frac1{r_+r_-}.
\]

Cela change complètement l’interprétation de la perméabilité et de la permittivité.\par

La permittivité est la réponse quadratique d’un feuillet. La perméabilité est la réponse quadratique du feuillet conjugué. L’impédance conserve la relation orientée entre les deux. La vitesse de propagation dépend de leur produit commun.\par

Ainsi, produit et quotient accomplissent deux tâches différentes :\par

\begin{itemize}[leftmargin=*,itemsep=0.35em]
\item le produit oublie quel était chaque feuillet et conserve la propagation commune ;
\item le quotient conserve quelle était leur orientation relative.
\end{itemize}

Le produit répond à la question : « que partagent les deux feuillets ? ». Le quotient répond à la question : « comment se distinguent-ils ? ».\par

Por eso\par

\[
\widehat\mu\,\widehat\varepsilon=\widehat c^{-2}
\]

y\par

\[
\frac{\widehat\mu}{\widehat\varepsilon}
=
\widehat Z^2
\]

sont les deux manières complémentaires de lire la même paire spectrale.\par

L’aspect holographique apparaît ici de façon très nette : quatre grandeurs visibles ne contiennent en réalité que deux degrés de liberté internes. Chacune de certaines paires adéquates permet de reconstruire les deux feuillets. La publication constitutive conserve la généalogie lorsque produit et quotient restent conjointement disponibles.\par

Telle est la réparation profonde de l’ablation de \(\mu\) et de \(\varepsilon\). La réparation réunit leurs valeurs et leur origine opératorielle commune.\par

La carte du Système international peut ensuite transformer ces sections intrinsèques en \(\mu_0\), \(\varepsilon_0\), \(Z_0\) et \(c\). Sa fonction consiste à exprimer, au sein d’une convention métrologique, une structure dont la provenance a déjà été sélectionnée par l’opérateur.\par

Le vide cesse ainsi d’être un réservoir de quatre constantes. Il devient une unique opération spectrale avec deux mémoires conjuguées.\par

Et cette vision a une conséquence philosophique importante : le vide est une relation de compatibilité entre deux manières orientées de réaliser un même état.\par


%
\endgroup

\HMTFlushCurrentChapter
\chapter{Dérivation de la charge et des quanta électriques : résistance de von Klitzing, conductance, constante de Josephson et flux magnétique}%
\begingroup
\renewcommand{\chapter}[2][]{}%
\chapter{Dérivation de la charge et des quanta électriques : résistance de von Klitzing, conductance, constante de Josephson et flux magnétique}
\label{chap:p3-final:40:cuantos_electricos}
\label{ch:metrology}

\section{Dérivation de la section de charge à partir du vide}

\(Z_{\rm vac}\), \(h_a=2\pi\hbar_a\) et \(\alpha_{\HMT}\) étant fixés, la charge du feuillet \(a\in\{\mathrm{pre},\mathrm{ret}\}\) est

\begin{equation}
e_a=\sqrt{\frac{2\alpha_{\HMT}h_a}{Z_{\rm vac}}}.
\label{eq:charge}
\end{equation}

La section \(e_a\) n'est pas introduite au moyen d'une coordonnée du Système
international destinée à reconstruire \(\alpha\). L'ordre causal procède de
\(\alpha_{\HMT}\), de l'action et du vide vers la section de charge.

\section{Constantes de von Klitzing et de Josephson, conductance et flux magnétique}

De \cref{eq:charge} se déduisent

\begin{align}
R_K&=\frac{h_a}{e_a^2}=\frac{Z_{\rm vac}}{2\alpha_{\HMT}},
\label{eq:rk}\\
G_0&=\frac{2e_a^2}{h_a}=\frac{4\alpha_{\HMT}}{Z_{\rm vac}},
\label{eq:g0}\\
\Phi_{0,a}&=\frac{h_a}{2e_a}
=\sqrt{\frac{h_aZ_{\rm vac}}{8\alpha_{\HMT}}},
\label{eq:phi0}\\
K_{J,a}&=\frac{2e_a}{h_a}=\Phi_{0,a}^{-1}.
\label{eq:kj}
\end{align}

En conséquence,

\begin{equation}
R_KG_0=2,\qquad G_0Z_{\rm vac}=4\alpha_{\HMT},
\qquad K_{J,a}\Phi_{0,a}=1.
\label{eq:quantum-electrical-identities}
\end{equation}

La covariance bidirectionnelle prend la forme

\[
R_K^{\rm pre}=R_K^{\rm ret},\qquad
G_0^{\rm pre}=G_0^{\rm ret},
\]

\[
\frac{\Phi_{0,\rm pre}}{\Phi_{0,\rm ret}}=R_{\rm act}^{1/2},
\qquad
\frac{K_{J,\rm pre}}{K_{J,\rm ret}}=R_{\rm act}^{-1/2}.
\]

Les relations de résistance et de conductance sont indépendantes de
l'orientation de l'action. Le couple flux--fréquence conserve cette
orientation et transforme réciproquement ses deux composantes.


%
\endgroup

\HMTFlushCurrentChapter
\chapter{Caractères quadratiques et constantes spectrales : Catalan, Apéry, Euler--Mascheroni et hiérarchie de Stieltjes}%
\begingroup
\renewcommand{\chapter}[2][]{}%
\chapter{Caractères quadratiques et constantes spectrales : Catalan, Apéry et tour d’Euler--Stieltjes}
\label{chap:p3-final:41:catalan_apery_euler}
\label{ch:spectral}

\section{Représentation opératorielle des caractères de Dirichlet}

Les constantes de ce chapitre sont regroupées parce qu'elles partagent une construction
spectrale indépendante de leurs développements décimaux et exempte de valeurs
cibles. Soit

\[
Ne_n=ne_n,\qquad n\in\N,
\]

dans \(\ell^2(\N)\). Un caractère résiduel \(\chi\) définit l'opérateur diagonal \(Q_\chi e_n=\chi(n)e_n\) ; pourvu que la puissance soit de classe trace,

\begin{equation}
\Tr(N^{-s}Q_\chi)=L(s,\chi).
\label{eq:L-trace}
\end{equation}

L'équation conserve le module, l'orientation et la provenance de chaque
caractère ; par conséquent, elle n'identifie pas des objets résiduels distincts.

\section{Construction HMT des caractères résiduels}

L'opérateur diagonal précédent ne constitue pas la donnée algébrique initiale. La
chaîne de sélection est

\[
\begin{aligned}
\APP&\longrightarrow\TRIT\longrightarrow\TPK
\longrightarrow\{\text{résidu modulo }3,\ C_P\}\\
&\longrightarrow\{\chi_{-3},\ J=C_P\bmod3\}
\longrightarrow\{\chi_{-3},\ J^2=-I,\chi_{-4}\}\\
&\longrightarrow Q_3,Q_4.
\end{aligned}
\]

La fonctionnelle résiduelle triadique détermine la différence entre les deux classes
orientées non nulles. Dans le secteur d'incidence, on emploie la matrice de conférence
symétrique de Paley d'ordre six,
\[
C_P^{\mathsf T}=C_P,\qquad C_PC_P^{\mathsf T}=5I_6.
\]
Sa réduction \(J=C_P\bmod3\) satisfait
\[
J^2=2I_6=-I_6
\quad\text{dans }\operatorname{End}_{\mathbb F_3}(\mathbb F_3^6).
\]
Le quart de torsion détermine alors la phase cyclique
\(1,\ii,-1,-\ii\). L'application du théorème chinois des restes conserve
les deux données résiduelles et produit \(Q_{12}=Q_3Q_4\). La trace
\(\Tr(N^{-s}Q)\) est évaluée à l'étape terminale ; ni la série de Dirichlet
ni son expression décimale n'interviennent dans la sélection du caractère.

La construction est explicitée par les tableaux
\begin{equation}
\chi_{-3}(n)=
\begin{cases}
0,&3\mid n,\\
1,&n\equiv1\pmod3,\\
-1,&n\equiv2\pmod3,
\end{cases}
\qquad
\chi_{-4}(n)=
\begin{cases}
0,&2\mid n,\\
1,&n\equiv1\pmod4,\\
-1,&n\equiv3\pmod4.
\end{cases}
\label{eq:characters34-tables}
\end{equation}
En \(\ell^2(\N)\),
\[
Q_3e_n=\chi_{-3}(n)e_n,
\qquad Q_4e_n=\chi_{-4}(n)e_n.
\]
Comme ils sont diagonaux, ils commutent ; le théorème chinois des restes donne
\begin{equation}
(Q_3Q_4)e_n=\chi_{-3}(n)\chi_{-4}(n)e_n
=\chi_{12}(n)e_n.
\label{eq:q12-product-proof}
\end{equation}
Le produit conserve simultanément le résidu triadique et le quart de
torsion. Par conséquent, le caractère modulo douze est déterminé
avant d'effectuer toute évaluation de ses valeurs \(L\).

\begin{remark}[Contrôle négatif]
La chaîne est invalidée si \(Q_3\) ne distingue pas les deux classes orientées,
si \(J^2\ne-I\), si \(Q_3Q_4\) n'a pas le tableau CRT déclaré ou si
l'opérateur est sélectionné à partir de la valeur de la constante.
\end{remark}

\section{Structure complexe interne et constante de Catalan}

La structure complexe interne satisfait

\begin{equation}
J^2=-I.
\label{eq:J-square}
\end{equation}

L'égalité \(J^2=-I\) définit une représentation
\(\rho:C_4\to\operatorname{Aut}_{\mathbb F_3}(\mathbb F_3^6)\),
\(\rho(j)=J\). Pour la réaliser spectralement, on prend le caractère unitaire
\(\upsilon:C_4\to U(1)\), \(\upsilon(j)=\ii\), et l'application de phase
\(n\mapsto j^n\). Sur \(\ell^2(\N)\), la représentation diagonale induite
est
\[
U_4e_n=\upsilon(j^n)e_n=\ii^ne_n.
\]
Sa partie orientée est
\[
Q_4=\frac{U_4-U_4^*}{2\ii},\qquad
Q_4e_n=\sin\!\left(\frac{\pi n}{2}\right)e_n
=\chi_{-4}(n)e_n.
\]
C'est le relèvement explicite entre le générateur fini d'ordre quatre et l'opérateur
résiduel ; on n'identifie pas \(\mathbb F_9\) à \(\mathbb C\).

Sur les entiers impairs, le même caractère peut s'écrire

\[
\chi_J(n)=\ii^{n-1}=\chi_{-4}(n).
\]

De manière équivalente,

\begin{equation}
Q_4=\frac{U_4-U_4^*}{2\ii},
\qquad Q_4e_n=\chi_{-4}(n)e_n.
\label{eq:Q4}
\end{equation}

La constante de Catalan apparaît comme

\begin{equation}
G=\Tr(N^{-2}Q_4)=L(2,\chi_{-4})
=\sum_{n\ge0}\frac{(-1)^n}{(2n+1)^2}.
\label{eq:catalan}
\end{equation}

Le relèvement unitaire du quart de torsion engendre exactement le
caractère dont le moment quadratique est la constante de Catalan.

\begin{proof}[Preuve opératorielle]
L'opérateur \(N^{-2}Q_4\) est de classe trace parce que
\(\sum_{n\ge1}|\chi_{-4}(n)|n^{-2}<\infty\). Sa trace est donc la
somme des éléments diagonaux. Le tableau \cref{eq:characters34-tables}
annule les pairs et alterne les impairs :
\[
\Tr(N^{-2}Q_4)
=1-\frac1{3^2}+\frac1{5^2}-\frac1{7^2}+\cdots.
\]
L'identité avec \(L(2,\chi_{-4})\) est la définition de la série de
Dirichlet ; l'égalité avec \(\operatorname{Im}\operatorname{Li}_2(\ii)\)
s'obtient en prenant la partie imaginaire de
\(\sum_{n\ge1}\ii^n/n^2\).
\end{proof}

La valeur terminale est
\[
G=0.9159655941772190150546035149\ldots.
\]
La structure déterminante est l'identité \(J^2=-I\), qui induit
exactement le signe alterné des classes impaires ; le développement décimal
est son évaluation terminale.

\section{Composition résiduelle triadique--quaternaire}

La fonctionnelle résiduelle modulo trois fournit \(\chi_{-3}\). Son produit
CRT avec \(\chi_{-4}\) définit

\begin{equation}
\chi_{12}=\chi_{-3}\chi_{-4}.
\label{eq:chi12}
\end{equation}

Sur les unités modulo douze,

\[
\chi_{12}(1)=1,\quad
\chi_{12}(5)=-1,\quad
\chi_{12}(7)=-1,\quad
\chi_{12}(11)=1,
\]

et s'annule en dehors de celles-ci. Deux évaluations exactes sont

\begin{align}
L(1,\chi_{12})&=\frac{\log(2+\sqrt3)}{\sqrt3},
\label{eq:L1chi12}\\
L(2,\chi_{12})&=\frac{\pi^2}{6\sqrt3}.
\label{eq:L2chi12}
\end{align}

La première contient l'unité hyperbolique

\begin{equation}
\sqrt3L(1,\chi_{12})=\arcosh2=\log(2+\sqrt3).
\label{eq:pell2}
\end{equation}

Avec \(\log(3+2\sqrt2)=2\log(1+\sqrt2)\) et
\(\log(30+\sqrt{899})\), cette identité détermine une famille d'échelles
de Pell dont les membres conservent leurs provenances algébriques respectives.

Les évaluations terminales certifiées sont

\[
L(1,\chi_{12})=0.7603459963009463475310942549\ldots,
\]

\[
L(2,\chi_{12})=0.9497031262940093952634984917\ldots.
\]

\begin{proof}[Évaluations dodécaphasiques]
En regroupant par classes résiduelles,
\begin{align*}
L(1,\chi_{12})
&=\sum_{n\ge0}\left(
\frac1{12n+1}-\frac1{12n+5}
-\frac1{12n+7}+\frac1{12n+11}\right)\\
&=\int_0^1
\frac{1-x^4-x^6+x^{10}}{1-x^{12}}\,dx
=\int_0^1\frac{1-x^4}{1+x^6}\,dx.
\end{align*}
Une décomposition réelle dans les facteurs quadratiques de \(1+x^6\)
produit
\(\log(2+\sqrt3)/\sqrt3\). Pour le deuxième moment, si
\(\psi_1\) est la fonction trigamma,
\begin{align*}
L(2,\chi_{12})
&=\frac1{144}\bigl[
\psi_1(1/12)-\psi_1(5/12)
-\psi_1(7/12)+\psi_1(11/12)\bigr]\\
&=\frac{\pi^2}{144}\left(
\csc^2\frac\pi{12}-\csc^2\frac{5\pi}{12}\right)
=\frac{\pi^2}{6\sqrt3},
\end{align*}
où l'on a utilisé
\(\psi_1(z)+\psi_1(1-z)=\pi^2\csc^2(\pi z)\). Les deux preuves emploient
le tableau CRT, et non une valeur décimale cible.
\end{proof}

\section{Constantes d'Euler et de Stieltjes et transformée thêta--Mellin}

Soit
\[
\vartheta(t)=\sum_{n\in\Z}e^{-\pi n^2t},\qquad t>0.
\]
Pour \(\Re s>1\), la fonctionnelle thêta--Mellin retrouve l'identité convergente

\[
\pi^{-s/2}\Gamma\!\left(\frac{s}{2}\right)\zeta(s)
=\frac12\int_0^\infty
\bigl(\vartheta(t)-1\bigr)t^{s/2-1}\,dt.
\]

Le prolongement méromorphe s'obtient en scindant l'intégrale en \(t=1\), en utilisant
\(\vartheta(t)=t^{-1/2}\vartheta(1/t)\) et en séparant les termes polaires. Son développement en \(s=1\),

\begin{equation}
\zeta(s)=\frac1{s-1}+\sum_{j\ge0}
\frac{(-1)^j\gamma_j}{j!}(s-1)^j,
\label{eq:stieltjes}
\end{equation}

situe \(\gamma=\gamma_0\) et les constantes de Stieltjes dans une même
famille de moments renormalisés. Cette organisation exprime une
provenance analytique commune, sans attribuer une interprétation physique
identique à tous ses niveaux.


La constante de Catalan apparaît habituellement comme une série alternée :\par

\[
G
=
1-\frac1{3^2}+\frac1{5^2}-\frac1{7^2}+\cdots.
\]

Cette formule donne sa valeur et la généalogie HMT explique pourquoi cette alternance particulière existe.\par

Dans HMT, l’origine se trouve dans le quart de torsion.\par

Il existe un opérateur interne \(J\) tel que\par

\[
J^2=-I.
\]

La racine de \(-1\) émerge de deux applications successives d’un quart de tour. La structure complexe est la mémoire algébrique d’une rotation interne.\par

Cette rotation engendre un caractère d’ordre quatre. Lorsqu’il est publié sur les entiers positifs, il produit exactement le motif résiduel qui alterne entre \(+1\), \(0\), \(-1\), \(0\).\par

La constante de Catalan apparaît alors comme une trace spectrale :\par

\[
G
=
L(2,\chi_{-4})
=
\operatorname{Tr}(N^{-2}Q_4).
\]

Cela change sa signification. Catalan est simultanément la somme d’une série particulière et la réponse spectrale des entiers à un quart de torsion.\par

Chaque terme de la série enregistre la manière dont la rotation interne agit sur une classe résiduelle. L’alternance est l’ombre arithmétique de la rotation.\par

La structure triadique possède son propre caractère modulo \(3\). En le combinant au caractère d’ordre quatre, le théorème chinois des restes produit un caractère modulo \(12\) :\par

\[
\chi_{12}
=
\chi_{-3}\chi_{-4}.
\]

Le module \(12\) apparaît comme l’espace commun minimal dans lequel l’orientation triadique et le quart de torsion peuvent coexister.\par

La dodécaphase est, en ce sens, une réconciliation de \(3\) et de \(4\).\par

Du caractère combiné émergent\par

\[
L(1,\chi_{12})
=
\frac{\log(2+\sqrt3)}{\sqrt3},
\]

y\par

\[
L(2,\chi_{12})
=
\frac{\pi^2}{6\sqrt3}.
\]

La cantidad\par

\[
2+\sqrt3
\]

est l’unité fondamentale du corps réel \(\mathbb Q(\sqrt3)\). Son logarithme représente un déplacement hyperbolique. De nouveau, une structure finie de résidus ouvre une dynamique multiplicative infinie.\par

Ici se rencontrent rotation, torsion, caractère arithmétique, unité de Pell et déplacement hyperbolique.\par

La constante d’Euler–Kronecker du corps \(\mathbb Q(\sqrt3)\) prolonge cette structure. Elle mesure le terme régulier qui demeure après la séparation du pôle principal de la fonction zêta du corps. Dans la généalogie HMT, ce terme constitue l’invariant régulier de mémoire arithmétique fine qui subsiste après le retrait de la divergence universelle.\par

La famille Bendersky–Glaisher prolonge le mécanisme en une tour. Les dérivées spectrales de zêta engendrent cette tour. Apéry et d’autres constantes associées à des valeurs impaires s’organisent comme membres d’une même famille fonctionnelle.\par

L’apport nouveau réside dans la généalogie qui réunit les constantes classiques :\par

\begin{itemize}[leftmargin=*,itemsep=0.35em]
\item le quart de torsion produit la structure complexe ;
\item le résidu triadique produit la structure modulo \(3\) ;
\item leur composition produit le caractère modulo \(12\) ;
\item le caractère engendre des valeurs \(L\) ;
\item les valeurs \(L\) ouvrent le corps quadratique et ses constantes spectrales.
\end{itemize}

La torsion géométrique est devenue arithmétique en conservant intégralement son orientation.\par

%
\endgroup

\HMTFlushCurrentChapter
\chapter{Constantes d'Euler--Kronecker et distributions premières : caractères dodécaphasiques, Meissel--Mertens et séries de nombres premiers jumeaux}%
\begingroup
\renewcommand{\chapter}[2][]{}%
\chapter{Constantes d’Euler--Kronecker et distributions de nombres premiers : caractères dodécaphasiques, Meissel--Mertens et séries de nombres premiers jumeaux}
\label{chap:p3-final:42:euler_kronecker_primos}
\section{Constante d'Euler--Kronecker du caractère dodécaphasique}

La factorisation de Dedekind

\begin{equation}
\zeta_{\mathbb Q(\sqrt3)}(s)=\zeta(s)L(s,\chi_{12})
\label{eq:dedekind-sqrt3}
\end{equation}

donne, en comparant les termes constants en \(s=1\),

\begin{equation}
\gamma_{\mathbb Q(\sqrt3)}
=\gamma+\frac{L'(1,\chi_{12})}{L(1,\chi_{12})}
=1.0539656082484825800\ldots.
\label{eq:euler-kronecker}
\end{equation}

Cette constante représente le terme fini normalisé au pôle de la
fonction zêta du corps réel dodécaphasique. Elle ne constitue pas une variante de
\(\gamma\) d'Euler, mais la correction arithmétique induite par
\(\chi_{12}\).

En effet, en écrivant \(u=s-1\),
\[
\zeta(s)=u^{-1}+\gamma+O(u),
\qquad
L(s,\chi_{12})=L_1+L_1'u+O(u^2),
\]
on obtient
\begin{equation}
\zeta_{\mathbb Q(\sqrt3)}(s)
=\frac{L_1}{u}+\bigl(\gamma L_1+L_1'\bigr)+O(u).
\label{eq:ek-laurent-derivation}
\end{equation}
La constante d'Euler--Kronecker est le quotient du terme fini par
le résidu ; diviser \cref{eq:ek-laurent-derivation} par \(L_1\) prouve
\cref{eq:euler-kronecker}. Cette dérivation explique simultanément le
résidu
\[
\operatorname*{Res}_{s=1}\zeta_{\mathbb Q(\sqrt3)}(s)
=\frac{\log(2+\sqrt3)}{\sqrt3}
\]
et la correction finie. La valeur numérique est certifiée en évaluant
\(L(1,\chi_{12})\) et \(L'(1,\chi_{12})\) avec le même tableau résiduel et une
borne de queue ; elle n'est pas importée d'un tableau de corps quadratiques.

\section{Constantes arithmétiques associées aux nombres premiers}

Une fois les cycles arithmétiques irréductibles sélectionnés, on obtient deux
familles distinctes :

\begin{align}
B_1&=\lim_{x\to\infty}\left(
\sum_{p\le x}\frac1p-\log\log x
\right),
\label{eq:meissel-mertens}\\
C_2&=\prod_{p>2}\frac{p(p-2)}{(p-1)^2}.
\label{eq:twin-prime-constant}
\end{align}

\(B_1\) régularise une somme locale ; \(C_2\) est une densité multiplicative de paires jumelles. Les deux exigent un contrôle de queue. Ils ne se déduisent pas l'un de l'autre parce qu'ils partagent les nombres premiers.


%
\endgroup

\HMTFlushCurrentChapter
\chapter{Famille Bendersky--Glaisher--Kinkelin : produits réguliers, dérivées de la fonction zêta et hiérarchies de normalisation}%
\begingroup
\renewcommand{\chapter}[2][]{}%
\chapter{Famille Bendersky--Glaisher--Kinkelin : produits réguliers, dérivées de la fonction zêta et hiérarchies de normalisation}
\label{chap:p3-final:43:bendersky_glaisher}
\section{Génération de la famille Bendersky--Glaisher--Kinkelin par régularisation zêta}

Soit \(\mathcal Z_3(s):=\zeta(s)\) la réalisation scalaire de la fonctionnelle
zêta triadique dans la normalisation de ce volume,
\(H_k=\sum_{j=1}^k j^{-1}\) (avec \(H_0=0\)) et \(B_{k+1}\) le nombre de
Bernoulli de convention \(B_1=-1/2\). Nous définissons

\begin{equation}
A_k=\exp\!\left(
\frac{H_kB_{k+1}}{k+1}-\mathcal Z_3'(-k)
\right),\qquad k\ge0.
\label{eq:bendersky}
\end{equation}

Les premiers niveaux incluent

\[
A_0=\sqrt{2\pi},
\]

et la constante de Glaisher--Kinkelin en \(A_1\). Pour \(n\ge1\),

\begin{equation}
\zeta(2n+1)=(-1)^{n+1}
\frac{2(2\pi)^{2n}}{(2n)!}\log A_{2n}.
\label{eq:odd-zeta-bendersky}
\end{equation}

La constante d'Apéry, \(\zeta(3)\), est le premier niveau pair non trivial.
Sa présence dans le gaz photonique du \cref{ch:metrology} constitue une
réalisation physique ultérieure de cette valeur spectrale.

La première partie de la tour peut se voir de manière synoptique :
\begin{center}
\small
\begin{tabular}{@{}cll@{}}
\toprule
\(k\) & identidad estructural & valeur terminale\\
\midrule
0 & \(A_0=\exp[-\zeta'(0)]=\sqrt{2\pi}\)
  & \(2.5066282746310005024\ldots\)\\
1 & \(\log A_1=1/12-\zeta'(-1)\)
  & \(1.2824271291006226369\ldots\)\\
2 & \(\log A_2=-\zeta'(-2)=\zeta(3)/(4\pi^2)\)
  & \(1.0309167521973921\ldots\)\\
3 & \(\log A_3=-11/720-\zeta'(-3)\)
  & \(0.9795555269428446\ldots\)\\
4 & \(\log A_4=-\zeta'(-4)\)
  & \(0.9920479745250403\ldots\)\\
\bottomrule
\end{tabular}
\end{center}

\begin{proof}[Relation avec les valeurs impaires]
L'équation fonctionnelle de \(\zeta\), développée en \(s=-2n\), donne
\begin{equation}
\zeta'(-2n)=(-1)^n
\frac{(2n)!}{2(2\pi)^{2n}}\zeta(2n+1).
\label{eq:zeta-derivative-even-negative}
\end{equation}
Comme \(B_{2n+1}=0\) pour \(n\ge1\), la définition
\cref{eq:bendersky} réduit \(\log A_{2n}=-\zeta'(-2n)\). Isoler
\(\zeta(2n+1)\) prouve \cref{eq:odd-zeta-bendersky}. L'indice \(2n\)
conserve le niveau de dérivation : ce n'est pas une étiquette éditoriale.
\end{proof}


%
\endgroup

\HMTFlushCurrentChapter
\chapter{Dynamique de Gauss et constantes de Khinchin, Lévy et Lochs : cofinalité des cylindres, opérateur de transfert et spectre de Gauss--Kuzmin--Wirsing}%
\begingroup
\renewcommand{\chapter}[2][]{}%
\chapter{Dynamique de Gauss et constantes de Khinchin, Lévy et Lochs : cofinalité des cylindres, opérateur de transfert et spectre de Gauss--Kuzmin--Wirsing}
\label{chap:p3-final:44:gauss_khinchin_levy_lochs}
\label{ch:gauss}

\section{Cofinalité entre cylindres HMT et cylindres de fractions continues}

Une section HMT irrationnelle détermine une famille emboîtée de cylindres compacts. Son évaluation archimédienne est un point \(x\in(0,1)\setminus\mathbb Q\). La fraction continue de \(x\) détermine des cylindres classiques

\[
I(a_1,\ldots,a_n).
\]

\begin{proposition}[Cofinalité]
Pour tout cylindre fini \(I(a_1,\ldots,a_n)\) qui contient l'évaluation \(x\), il existe une profondeur HMT \(m(n)\) telle que le cylindre HMT de niveau \(m\ge m(n)\) soit contenu dans \(I(a_1,\ldots,a_n)\).
\end{proposition}

\begin{proof}
\(n\) étant fixé, le cylindre de fraction continue \(I(a_1,\ldots,a_n)\) est un voisinage ouvert de l'irrationnel \(x\). Il existe alors \(\varepsilon_n>0\) avec \((x-\varepsilon_n,x+\varepsilon_n)\subset I(a_1,\ldots,a_n)\). Les cylindres HMT \(C_m\) contiennent \(x\), sont emboîtés et satisfont \(\operatorname{diam}C_m\to0\). Choisissons \(m(n)\) avec \(\operatorname{diam}C_{m(n)}<\varepsilon_n\). Comme \(x\in C_{m(n)}\), tout point de ce cylindre est à une distance de \(x\) inférieure à \(\varepsilon_n\), donc \(C_{m(n)}\subset I(a_1,\ldots,a_n)\). L'emboîtement donne l'inclusion pour tout niveau ultérieur.
\end{proof}

La démonstration utilise la compacité, l'emboîtement et l'unicité de
l'évaluation. Le résultat établit le transport des approximations, mais
n'implique pas la conjugaison du décalage de mémoire TPK avec l'application de
Gauss sur toutes les fibres profondes.

La généalogie complète de cette réalisation est
\[
\APP\to\TRIT\to\TPK
\to (C_m,\rho_{m+1,m})_{m\ge0}
\to x=\bigcap_m C_m
\to \bigl(I(a_1,\ldots,a_n)\bigr)_{n\ge1}.
\]
Les cylindres HMT conservent la route, la phase et la retenue ; l'application terminale projette
ces coordonnées sur les préfixes de fraction continue. La cofinalité
démontre que cette projection converge vers le même point. Un entrelaceur
dynamique exigerait en outre un relèvement \(\widetilde T_G\) avec des
carrés commutatifs sur les fibres de valuation, propriété qui ne se
déduit pas de la cofinalité.

\section{Mesure de Gauss et observables de Khinchin et de Lévy}

L'application de Gauss \cite{Khinchin,IosifescuKraaikamp}

\[
T_G(x)=\frac1x-\left\lfloor\frac1x\right\rfloor
\]

préserve

\begin{equation}
d\mu_G(x)=\frac{dx}{\log2(1+x)}.
\label{eq:gauss-measure}
\end{equation}

La constante de Khinchin est définie par

\begin{equation}
\log K=\int_0^1\log a_1(x)\,d\mu_G(x).
\label{eq:khinchin}
\end{equation}

La distribution du premier chiffre se calcule directement :
\begin{equation}
\mu_G(a_1=k)
=\frac1{\log2}\int_{1/(k+1)}^{1/k}\frac{dx}{1+x}
=\log_2\!\left(\frac{(k+1)^2}{k(k+2)}\right).
\label{eq:gauss-digit-law}
\end{equation}
Par le théorème ergodique,
\begin{equation}
K=\prod_{k\ge1}
k^{\,\log_2((k+1)^2/[k(k+2)])}
=2.685452001065306445309714835\ldots
\label{eq:khinchin-product}
\end{equation}
pour \(\mu_G\)-presque toute section. Le produit ne s'applique pas à chaque point :
la section \(\varphi^{-1}\), par exemple, a tous ses chiffres égaux
à un et constitue un critère de réfutation immédiat d'une lecture point par point. Une
certification de \cref{eq:khinchin-product} tronque la somme des logarithmes et
borne la queue au moyen de l'estimation
\[
\mu_G(a_1=k)=O(k^{-2}),
\qquad
\sum_{k>N}\mu_G(a_1=k)\log k
=O\!\left(\frac{\log N+1}{N}\right).
\]

La constante de Lévy est

\begin{equation}
L=-\int_0^1\log x\,d\mu_G(x)
=\frac{\pi^2}{12\log2}.
\label{eq:levy}
\end{equation}

En effet,

\[
\frac1{1+x}=\sum_{n\ge0}(-1)^nx^n,\qquad
\int_0^1(-\log x)x^n\,dx=\frac1{(n+1)^2},
\]

de sorte que l'intégrale est \(\eta(2)/\log2=\pi^2/(12\log2)\).

La constante de Khinchin détermine la moyenne logarithmique des quotients
partiels, tandis que la constante de Lévy détermine le taux exponentiel
de croissance des dénominateurs. Ce sont deux observables distinctes d'un
même système mesuré.

\section{Entropie métrique, exposant de Lyapunov et constante de Lochs}

Comme

\[
\log|T_G'(x)|=-2\log x,
\]

l'exposant de Lyapunov et l'entropie métrique sont

\begin{equation}
\chi_G=h_G=2L=\frac{\pi^2}{6\log2}.
\label{eq:gauss-entropy}
\end{equation}

Le rapport de Lochs entre les fractions continues et un développement en base \(b\) est

\begin{equation}
\Lambda_b=\frac{\log b}{h_G}
=\frac{6\log2\log b}{\pi^2}.
\label{eq:lochs}
\end{equation}

En particulier,

\begin{align*}
\Lambda_{10}&=0.9702701143920339257402560192\ldots,\\
\Lambda_{729}&=\frac{36\log2\log3}{\pi^2}
=2.777618966374305777705782976\ldots.
\end{align*}

La base \(729=3^6\) est pertinente pour l'espace ambiant \(\mathbb F_3^6\), mais \cref{eq:lochs} reste un théorème presque sûr de \(\mu_G\) ; et non une égalité point par point pour toute route HMT.

La valeur \(\Lambda_{729}\) représente le taux asymptotique typique de
quotients partiels obtenus par bloc complet de six trits. Elle n'identifie pas
l'espace ambiant de \(729\) mots aux autres objets de cardinal \(729\) du
corpus. L'égalité est un taux d'information entre partitions, et non une
bijection entre états.

\section{Opérateur de transfert de Gauss--Kuzmin--Wirsing}

L'opérateur de Gauss--Kuzmin--Wirsing agit par

\begin{equation}
(\mathcal L f)(x)=\sum_{n\ge1}\frac1{(x+n)^2}
f\!\left(\frac1{x+n}\right).
\label{eq:gkw-operator}
\end{equation}

Sa valeur propre dominante correspond à la densité invariante. La constante GKW est le module de la valeur propre sous-dominante sur le sous-espace de moyenne nulle. Un schéma HMT compatible utilise des partitions de \(9^m\) cylindres, des matrices de transfert avec intervalles et une borne de queue.

Pour promouvoir le calcul, il faut certifier : l'isolement de la valeur propre, l'invariance du sous-espace, l'erreur de troncature, la convergence des projecteurs spectraux et l'absence d'une autre valeur propre de même module. La valeur décimale connue ne peut pas être introduite comme centre de l'intervalle initial.

Une déformation conjointe contient Khinchin et Lévy dans un seul objet :
\begin{equation}
(\mathcal L_{s,t}f)(x)
=\sum_{n\ge1}n^t(n+x)^{-2s}
 f\!\left(\frac1{n+x}\right),
\qquad
P(s,t)=\log r(\mathcal L_{s,t}).
\label{eq:gauss-pressure-operator}
\end{equation}
Au point \((s,t)=(1,0)\), la perturbation spectrale donne
\begin{equation}
\partial_tP(1,0)=\log K,
\qquad
-\partial_sP(1,0)=\chi_G=2L.
\label{eq:pressure-derivatives}
\end{equation}
La hessienne de \(P\) détermine les variances et covariances asymptotiques de
\(\log a_1\) et \(-2\log x\). Par conséquent, les constantes de Khinchin et de
Lévy sont deux dérivées transversales de la même fonction de pression.

La reconnaissance numérique de la valeur propre sous-dominante est
\[
\lambda_2\simeq-0.3036630028987326586,\qquad
\kappa_{\rm GKW}=|\lambda_2|,\qquad
-\log\kappa_{\rm GKW}\simeq1.19183673556055.
\]
Dans ce volume, ces décimales constituent des valeurs de référence pour la
comparaison, et non un certificat spectral.
La preuve HMT proposée utilise Galerkin ou Ulam sur les cylindres cofinaux de
profondeur \(9^m\), une inégalité de Lasota--Yorke, l'arithmétique
d'intervalles pour isoler \(\lambda_2\) et une borne rigoureuse de la queue. Un
intervalle contenant deux valeurs propres de même module réfute la
promotion, sans pour autant affecter Khinchin, Lévy ou Lochs.


\section{Certificat cofinal du spectre de Gauss--Kuzmin--Wirsing sur la hiérarchie nonadique}
\label{sec:p3-final:cierre-gkw}

Pour Gauss--Kuzmin--Wirsing, soit \(P_N\) la famille de troncatures
certifiée sur l’espace fonctionnel déclaré, et définissons
\[
 \Pi_m=P_{9^m}.
\]
La hiérarchie \((\Pi_m)\) est une sous-suite cofinale ; les bornes uniformes, la
convergence des projecteurs et l’absence de pollution spectrale sont
héritées de la famille \((P_N)\). Une discrétisation cylindrique distincte devra
fournir son entrelaceur avec \((P_{9^m})\).

%
\endgroup

% L9-059. Propietario analítico íntegro de los funcionales espectrales
% triádicos (SHA-256 c53c5efa08ecacfed1a5cd01f20eaa372a29a56adab4873bcb442709963496f4).
% Se consume una sola vez y se subordina únicamente su jerarquía visible.
\begin{HMTScientificOwnerSubordinated}
\chapter{Construction analytique des fonctionnelles de la tour triadique}
\label{chap:funcionales-espectrales-triadicos}

La synthèse précédente a situé les constantes spectrales au sein de la
tour triadique. Ce bloc contient leur construction analytique primaire :
suite d'opérateurs autoadjoints sur des cycles d'ordre \(3^m\), limite de
traces de chaleur et transformée de Mellin. Les constantes classiques apparaissent
comme évaluations, parties finies ou dérivées de la fonctionnelle méromorphe
résultante, non comme des préfixes reconnus dans un catalogue.

La normalisation de l'opérateur $D_m$ contient explicitement $\pi$ dans l'échelle
qui relie le laplacien discret au noyau de chaleur gaussien. Cette
sous-section constitue donc une extension analytique de la branche
archimédienne déjà construite. Les constantes spectrales ultérieures s'obtiennent
simultanément à partir de la fonctionnelle résultante, avec cette normalisation pour hypothèse
de départ.

\section{Limites inductives et échelle de raffinement}

Le prolongement spectral exige de distinguer deux niveaux analytiques. Le
premier correspond à des familles bornées et compatibles ; le second, à des
opérateurs non bornés, pour lesquels l'entrelacement algébrique ne
résout pas à lui seul les problèmes de domaine.

\begin{theorem}[Limite inductive bornée]
\label{pf84:E03-06}
Soient \(J_m:\mathcal H_m\to\mathcal H_{m+1}\) des isométries et
\(A_m\in\mathcal B(\mathcal H_m)\) des opérateurs autoadjoints tels que
\[
 A_{m+1}J_m=J_mA_m,
 \qquad
 \sup_m\lVert A_m\rVert<\infty.
\]
Il existe un unique opérateur borné autoadjoint \(A\) sur
\(\mathcal H_\infty=\varinjlim(\mathcal H_m,J_m)\) dont la restriction à chaque
niveau est \(A_m\).
\end{theorem}

\begin{proof}
Sur l'union algébrique dense, on définit
\[
 A[J_{m,\infty}\xi]:=J_{m,\infty}A_m\xi.
\]
L'entrelacement rend la définition indépendante du représentant. La
borne uniforme permet de prolonger \(A\) par continuité à
\(\mathcal H_\infty\) ; la symétrie sur le sous-espace dense et le caractère borné
impliquent son autoadjonction.
\end{proof}

\begin{remark}[Contrôle pour les opérateurs non bornés]
\label{pf84:E03-07}
Si les \(A_m\) ne sont pas bornés, l'identité
\(A_{m+1}J_m=J_mA_m\) ne définit qu'une action sur une union de domaines.
Pour obtenir un opérateur fermé ou autoadjoint, il faut déclarer les domaines
transportés, prouver la fermabilité et contrôler les indices de
déficience, ou bien démontrer l'autoadjonction essentielle sur un cœur commun
compatible. Omettre ces obligations constitue une erreur de type, non une
abréviation de la preuve.
\end{remark}

Soit maintenant \(\mathcal S_\infty\) un ensemble non vide de sections
survivantes compatibles et soit \(\mathcal P_m\) l'ensemble fini de leurs
préfixes de longueur \(m\). Nous écrivons
\[
 \mathcal H_m:=\ell^2(\mathcal P_m),
 \qquad
 E_m(\gamma,\gamma')=
 \begin{cases}
 1,&\gamma'\text{ prolonge }\gamma,\\
 0,&\text{sinon},
 \end{cases}
\]
et
\[
 d_m(\gamma):=
 \#\{\gamma'\in\mathcal P_{m+1}:\gamma'\mapsto\gamma\},
 \qquad
 D_m^{\deg}:=E_mE_m^*=\operatorname{diag}(d_m(\gamma)).
\]
Nous nous restreignons au sous-espace vivant \(d_m(\gamma)>0\), et chaque préfixe de
niveau \(m+1\) possède un parent unique.

\begin{theorem}[Isométrie normalisée de prolongement]
\label{pf84:C01}
L'opérateur
\[
 J_m:=E_m^*(D_m^{\deg})^{-1/2}:
 \mathcal H_m\longrightarrow\mathcal H_{m+1}
\]
est une isométrie, et
\[
 J_me_\gamma=
 \frac1{\sqrt{d_m(\gamma)}}
 \sum_{\gamma'\mapsto\gamma}e_{\gamma'}.
\]
\end{theorem}

\begin{proof}
L'unicité du parent donne
\[
 J_m^*J_m
 =(D_m^{\deg})^{-1/2}E_mE_m^*(D_m^{\deg})^{-1/2}=I.
\]
La formule sur \(e_\gamma\) s'obtient en développant la colonne
correspondante de \(E_m^*\).
\end{proof}

En posant
\[
 \mathcal K_m:=\mathcal H_{m+1}\ominus J_m\mathcal H_m,
 \qquad
 \mathcal H_{\mathrm{surv}}
 \cong\mathcal H_0\oplus\bigoplus_{m\ge0}\mathcal K_m,
\]
on définit l'opérateur de raffinement
\[
 \mathcal R|_{\mathcal H_0}=0,
 \qquad
 \mathcal R|_{\mathcal K_m}=3^{m+1}I_{\mathcal K_m},
\]
de domaine maximal
\[
 \Dom(\mathcal R)=
 \mathcal H_0\oplus
 \left\{(\xi_m)_{m\ge0}:
 \sum_{m\ge0}3^{2m+2}\lVert\xi_m\rVert^2<\infty\right\}.
\]

\begin{theorem}[Opérateur de raffinement]
\label{pf84:C02}
Si chaque \(\mathcal K_m\) est de dimension finie, \(\mathcal R\) est positif,
autoadjoint et à résolvante compacte. Son spectre mesure la profondeur de
raffinement ; il ne constitue pas un spectre de nombres premiers ni de zéros.
\end{theorem}

\begin{proof}
La décomposition orthogonale réduit \(\mathcal R\) à des blocs scalaires réels
sur son domaine maximal, ce qui prouve la positivité et l'autoadjonction. Les
valeurs propres \(3^{m+1}\) s'échappent à l'infini et leurs multiplicités sont finies ;
la résolvante est donc compacte. Un bloc de dimension infinie à
valeur propre bornée réfuterait cette dernière conclusion.
\end{proof}

\section{Opérateurs entier, premier et orbital}

Soit \(\mathscr W_{\mathrm{bal}}\) l'ensemble des mots ternaires
équilibrés réduit au moyen du théorème
\cref{pf84:C03}. Sur \(\ell^2(\mathscr W_{\mathrm{bal}})\), l'opérateur
\[
 Be_w=b(w)e_w
\]
est pris sur son domaine maximal. Sa restriction au sous-espace des entiers
positifs est notée \(Q\).

\begin{theorem}[Opérateurs entier et positif]
\label{pf84:C04}
L'opérateur \(B\) est autoadjoint, possède un spectre simple \(\ZZ\) et une
résolvante compacte. L'opérateur \(Q\) possède un spectre simple
\(\NN_{\ge1}\) et une résolvante compacte.
\end{theorem}

\begin{proof}
La bijection de \cref{pf84:C03} énumère chaque entier exactement une fois. Un
opérateur diagonal réel sur son domaine maximal est autoadjoint ; comme ses entrées
s'échappent à l'infini hors des ensembles finis, sa résolvante est compacte. La
restriction positive conserve ces propriétés.
\end{proof}

La preuve classique omise de l'inversion thêta, du prolongement et de l'équation
fonctionnelle est fusionnée une seule fois avec la réalisation triadique déjà construite
dans \cref{pf84:C05}. On y fixe littéralement
\(\Theta_B=1+2\Theta_+\) et \(Z=\mathcal Z_3\) ; ce bloc ne
définit donc pas une seconde origine de \(\zeta\).

Soit \(\mathbb P\) l'ensemble des nombres premiers ordinaires,
\(\mathfrak h_{\mathbb P}:=\ell^2(\mathbb P)\) et
\[
 H_{\mathbb P}e_p=(\log p)e_p.
\]
Dans l'espace de Fock bosonique \(\Gamma_s(\mathfrak h_{\mathbb P})\), les occupations
\((k_p)_p\) sont à support fini.

\begin{theorem}[Représentation de Fock de la factorisation]
\label{pf84:C08}
L'application sur les bases
\[
 U_F\lvert(k_p)_p\rangle=e_{\prod_pp^{k_p}}
\]
se prolonge en un opérateur unitaire
\[
 U_F:\Gamma_s(\mathfrak h_{\mathbb P})
 \longrightarrow\ell^2(\NN_{\ge1})
\]
et satisfait
\[
 U_F\,d\Gamma(H_{\mathbb P})\,U_F^*=\log Q.
\]
\end{theorem}

\begin{proof}
La factorisation unique donne une bijection entre les occupations finies et les entiers
positifs, donc une identification unitaire des bases orthonormales.
De plus,
\[
 \sum_pk_p\log p=\log\!\left(\prod_pp^{k_p}\right),
\]
ce qui prouve l'entrelacement. La structure de départ explicite est le
choix des modes premiers et de l'espace de Fock symétrique.
\end{proof}

Sur
\[
 \mathcal H_{\mathrm{orb}}
 :=\ell^2\{(p,k,\varepsilon):
 p\in\mathbb P,\ k\ge1,\ \varepsilon=\pm1\},
\]
nous définissons
\[
 D_{\mathrm{orb}}e_{p,k,\varepsilon}
 =\varepsilon k\log p\,e_{p,k,\varepsilon},
 \qquad
 W_{\log}e_{p,k,\varepsilon}
 =(\log p)e_{p,k,\varepsilon}.
\]

\begin{proposition}[Trace orbitale]
\label{pf84:C09}
L'opérateur \(D_{\mathrm{orb}}\) est autoadjoint et à résolvante compacte. Si
\(f\in C_c(\RR)\) est bornée, alors
\[
 \Tr\!\left(
 W_{\log}e^{-|D_{\mathrm{orb}}|/2}f(D_{\mathrm{orb}})
 \right)
 =
 \sum_{p,k}(\log p)p^{-k/2}
 \bigl[f(k\log p)+f(-k\log p)\bigr].
\]
\end{proposition}

\begin{proof}
La diagonale est réelle et ses longueurs s'échappent de tout compact avec
multiplicité finie. Le support compact de \(f\) fait que le produit
régulé est de rang fini ; sa trace est la somme des entrées diagonales.
L'indice \(p\) parcourt ici les nombres premiers ordinaires. Le résultat réalise le
côté premier classique, mais n'identifie pas \(D_{\mathrm{orb}}\) à un opérateur
dual des zéros.
\end{proof}

\section{Forme complétée de Guinand--Weil}

Nous fixons
\[
 \mathcal D_{\mathrm W}:=C_c^\infty(\RR),
 \qquad
 \widehat\psi(t)=\int_{\RR}\psi(x)e^{-itx}\,\dd x,
 \qquad
 \psi(x)=\frac1{2\pi}\int_{\RR}\widehat\psi(t)e^{itx}\,\dd t,
\]
et, pour \(\psi,\phi\in\mathcal D_{\mathrm W}\),
\[
 h_{\psi,\phi}(z)
 :=\widehat\psi(z)\,
 \overline{\widehat\phi(\overline z)}.
\]
Sur la droite réelle \(h_{\psi,\psi}=|\widehat\psi|^2\).

\begin{definition}[Classe de test]
\label{pf84:E01}
La classe de test est \(\mathcal D_{\mathrm W}\), stable par dérivation,
réflexion et translation. On utilise l'inverse
\[
 \check h(x)=\frac1{2\pi}\int_{\RR}h(t)e^{itx}\,\dd t.
\]
La régularité et le support compact fournissent une décroissance rapide dans les bandes
et un support compact pour \(\check h\), conditions qui garantissent la
convergence des blocs suivants.
\end{definition}

Sea \(\psi_\Gamma=\Gamma'/\Gamma\). Definimos
\[
 \Irr_\#:=
 \left\{\operatorname{ev}_9(w):
 w\in\mathscr W_9^+,\ 
 P_{\Irr,\#}e_w=e_w\right\}.
\]
Par \cref{pf84:C07}, cette image coïncide avec les nombres premiers ordinaires ; elle n'est
utilisée ici que pour typer l'indice du bloc orbital, sans répéter cette
preuve.
\begin{align}
 \mathcal P(h)&:=h(i/2)+h(-i/2),\\
 \mathcal G(h)&:=-\log\pi\,\check h(0)
 +\frac1{2\pi}\int_{\RR}h(t)
 \Re\psi_\Gamma\!\left(\frac14+\frac{it}{2}\right)\dd t,\\
 \mathcal O_\#(\check h)&:=
 \sum_{p\in\Irr_\#}\sum_{k\ge1}
 (\log p)p^{-k/2}
 \bigl[\check h(k\log p)+\check h(-k\log p)\bigr].
\end{align}

\begin{definition}[Fonctionnelle complétée]
\label{pf84:E02}
Pour \(h=h_{\psi,\phi}\), on fixe
\[
 \mathcal W(h):=\mathcal P(h)+\mathcal G(h)
 -\mathcal O_\#(\check h).
\]
Le bloc orbital provient du sélecteur irréductible du
\cref{chap:producto-nativo} ; les blocs polaire et gamma appartiennent à la
complétion classique de zêta. Cette différence de provenance est conservée dans
la notation et dans la preuve.
\end{definition}

\begin{theorem}[Formule explicite de Guinand--Weil]
\label{pf84:E03} Pour \(\psi,\phi\in C_c^\infty(\RR)\) et \(h=h_{\psi,\phi}\),
\[
 \lim_{T\to\infty}
 \sum_{\substack{\rho:\,\zeta(\rho)=0\\|\Im\rho|\le T}}
 h\!\left(\frac{\rho-1/2}{i}\right)
 =\mathcal W(h),
\]
où les zéros non triviaux sont comptés avec multiplicité et la limite est
prise avec la symétrisation classique.
\end{theorem}

\begin{proof}
Paley--Wiener place \(h_{\psi,\phi}\) dans la classe admissible de la formule
explicite. Les pôles de
\(\Lambda_\zeta(s)=\pi^{-s/2}\Gamma(s/2)\zeta(s)\) produisent
\(\mathcal P\), le facteur archimédien produit \(\mathcal G\) et le produit
d'Euler produit la somme sur les puissances premières. L'identification de cette
dernière somme à \(\mathcal O_\#\) s'effectue uniquement au moyen de
\cref{pf84:E04} ; ainsi, la preuve de primalité native n'est pas dupliquée.
\end{proof}

\begin{proposition}[Réindexation du bloc orbital]
\label{pf84:E04}
Si les irréductibles du produit natif coïncident exactement avec les nombres premiers
ordinaires et que les poids
\((\log p)p^{-k/2}\) et les deux orientations \(\pm k\log p\) sont conservés, alors
\[
 \mathcal O_\#(\check h)=\mathcal O(\check h).
\]
\end{proposition}

\begin{proof}
La bijection entre irréductibles natifs et nombres premiers réindexe la somme sans
modifier les indices, les poids ni les arguments. Un irréductible non premier, un nombre premier
absent ou un changement de poids réfuterait l'identité.
\end{proof}

\begin{remark}[Séparation méromorphe--entière]
\label{pf84:E05}
La fonction \(\Lambda_\zeta\) possède des pôles simples en \(0\) et \(1\), qui
engendrent le bloc polaire. En revanche,
\[
 \xi(s)=\frac12s(s-1)\Lambda_\zeta(s)
\]
est entière : les facteurs \(s(s-1)\) annulent précisément ces pôles. On
n'attribue donc pas de pôles en \(0\) ou \(1\) à \(\xi\).
\end{remark}

\begin{definition}[Forme hermitienne et préopérateur]
\label{pf84:E06}
On définit
\[
 \mathcal B_{\mathrm W}(\psi,\phi)
 :=\mathcal W(h_{\psi,\phi}),
 \qquad
 \mathcal V_{\mathrm W}:=
 \mathcal D_{\mathrm W}/\operatorname{Rad}\mathcal B_{\mathrm W},
\]
et, sur le quotient algébrique, la règle candidate
\[
 D_{\mathrm W,0}[\psi]:=[-i\psi'].
\]
On ne présuppose pas que \(\mathcal B_{\mathrm W}\) soit positive ni que
\(\mathcal V_{\mathrm W}\) soit déjà un espace de Hilbert. La règle n'est
bien définie qu'après avoir prouvé l'invariance du radical.
\end{definition}

\begin{theorem}[Symétrie algébrique et invariances]
\label{pf84:E07}
Le radical est invariant sous \(-i\,\dd/\dd x\) ;
\(D_{\mathrm W,0}\) est bien défini et
\(\mathcal B_{\mathrm W}\)-symétrique. Les translations
\(T_a\psi(x)=\psi(x-a)\) préservent la forme et la réflexion
\(J\psi(x)=\psi(-x)\) satisfait
\[
 JD_{\mathrm W,0}J=-D_{\mathrm W,0}.
\]
\end{theorem}

\begin{proof}
La convention de Fourier donne
\[
 h_{-i\psi',\phi}(z)=h_{\psi,-i\phi'}(z).
\]
L'identité vaut avant l'application de l'un quelconque des trois blocs de
\(\mathcal W\), de sorte que la dérivée préserve le radical et que le
préopérateur est symétrique. Les phases opposées des translations s'annulent
dans \(h\) ; la réflexion change le signe de la dérivée et laisse invariants les
blocs polaire, gamma et orbital. Cette symétrie algébrique n'équivaut pas encore
à l'autoadjonction dans un Hilbert positif.
\end{proof}

\section{Géométrie locale de la droite critique}

\begin{lemma}[Coordonnée de Möbius de la droite critique]
\label{pf84:RH-01}
Pour \(s\in\CC\setminus\{0,1\}\), soit
\[
 \tau(s):=i\frac{s-1}{s}.
\]
Avec \(F(s)=1-s\) et
\(C(s)=\overline s\),
\[
 \tau(Fs)=-\frac1{\tau(s)},
 \qquad
 \tau(Cs)=-\overline{\tau(s)},
\]
et
\[
 \tau(Fs)=\tau(Cs)
 \iff |\tau(s)|=1
 \iff \Re s=\frac12.
\]
\end{lemma}

\begin{proof}
Les deux identités s'obtiennent par substitution directe. Leur égalité
équivaut à \(\tau(s)^{-1}=\overline{\tau(s)}\), c'est-à-dire
\(|\tau(s)|=1\) ; développer cette dernière condition donne
\(\Re s=1/2\). Pour les zéros de zêta, les symétries
fonctionnelle et conjuguée sont en outre requises. L'équivalence paramètre la droite ; elle ne place pas à elle
seule les zéros sur celle-ci.
\end{proof}

\begin{lemma}[Signe local de Hermite--Biehler]
\label{pf84:RH-04}
Soit \(A\) une fonction entière réelle ayant un zéro réel simple \(a\). Pour
\(z=a+iy\), \(y>0\) suffisamment petit,
\[
|A(z)-iA'(z)|^2-|A(z)+iA'(z)|^2
 =-4A'(a)^2y+O(y^2)<0.
\]
La convention locale associée à \(A\) est
\[
 E_A:=A+iA'.
\]
Ce n'est qu'en spécialisant \(A=\Xi\) et en supposant en outre que \(a\) est un zéro réel
simple de \(\Xi\) que l'on obtient
\[
 E_\Xi=\Xi+i\Xi'.
\]
\end{lemma}

\begin{proof}
On développe \(A\) et \(A'\) en séries de Taylor autour de \(a\), en utilisant
\(A(a)=0\) et \(A'(a)\ne0\), et l'on conserve les termes du premier ordre en
\(y\). Le résultat est strictement local : il ne démontre pas que \(E_A\) possède la
propriété globale de Hermite--Biehler.
\end{proof}

\begin{proposition}[Obstruction par les lois de comptage]
\label{pf84:E14}
Soit
\[
 N_A(T):=\Tr\mathbf1_{[-T,T]}(A).
\]
Les opérateurs \(Q\), \(B\), \(D_{\mathrm{orb}}\) et \(\mathcal R\) ne
possèdent pas à eux seuls la loi de comptage exigée d'un éventuel opérateur dual
des zéros.
\end{proposition}

\begin{proof}
Para \(T\ge1\),
\[
 N_Q(T)=\lfloor T\rfloor,
 \qquad
 N_B(T)=2\lfloor T\rfloor+1.
\]
Le théorème des nombres premiers donne
\[
 N_{D_{\mathrm{orb}}}(T)
 =2\sum_{p\le e^T}\left\lfloor\frac{T}{\log p}\right\rfloor
 \ge2\pi(e^T)\sim\frac{2e^T}{T}.
\]
Enfin, \(N_{\mathcal R}(T)\) est constante sur chaque intervalle
multiplicatif \([3^m,3^{m+1})\), car \(\mathcal R\) n'a que des niveaux
\(3^{m+1}\) de multiplicités finies.

Pour
\(\Xi(z)=\xi(1/2+iz)\), l'équation fonctionnelle implique
\(\Xi(-z)=\Xi(z)\) ; les zéros \(\pm\gamma\) ont la même multiplicité et
\[
 N_{\mathrm{sym}}(T)=2N_+(T)+O(1).
\]
Par conséquent, la convention bilatérale exige
\[
 N_{\mathrm{sym}}(T)\sim\frac{T}{\pi}\log T,
\]
tandis que la convention unilatérale équivalente, sous cette symétrie et cet appariement des
multiplicités, exige
\[
 N_+(T)\sim\frac{T}{2\pi}\log T.
\]
Aucune des quatre fonctions précédentes n'a cette asymptotique.
\end{proof}

\section{Laplaciens sur les cycles d'ordre \texorpdfstring{$3^m$}{puissances de trois}}

Pour $m\geq 1$, soient

\[
 N_m=3^m,
 \qquad
 C_m=\ZZ/N_m\ZZ,
\]

et considérons le laplacien combinatoire

\[
 (\Delta_m f)(j)=2f(j)-f(j+1)-f(j-1),
 \qquad f\in\ell^2(C_m).
\]

Les caractères

\[
 e_{m,n}(j)=N_m^{-1/2}
 \exp\!\left(\frac{2\pi i n j}{N_m}\right),
 \qquad n\in C_m,
\]

forment une base orthonormale et satisfont

\[
 \Delta_m e_{m,n}
 =4\sin^2\!\left(\frac{\pi n}{N_m}\right)e_{m,n}.
\]

Comme $N_m$ est impair, chaque classe de $C_m$ possède un unique représentant
$\widetilde n$ avec

\[
 -\frac{N_m-1}{2}\leq\widetilde n\leq\frac{N_m-1}{2}.
\]

Nous définissons l'opérateur autoadjoint $D_m$ par

\begin{equation}
 D_m e_{m,n}
 =\frac{N_m}{\sqrt\pi}
 \sin\!\left(\frac{\pi\widetilde n}{N_m}\right)e_{m,n}.
 \label{eq:operador-espectral-Dm}
\end{equation}

Alors

\begin{equation}
 D_m^2=\frac{N_m^2}{4\pi}\,\Delta_m.
 \label{eq:Dm-laplaciano}
\end{equation}

Soit $P_0$ la projection orthogonale sur le caractère constant. Pour toute
fonction paire bornée $F$ sur le spectre de $D_m$, nous introduisons la trace
positive

\begin{equation}
 \operatorname{Tr}^{+}_m\!\bigl(F(D_m)\bigr)
 :=\frac12\operatorname{Tr}
 \!\left(F(D_m)P_0^{\perp}\right).
 \label{eq:traza-positiva}
\end{equation}

Le facteur $1/2$ sélectionne une contribution par paire spectrale
$\{n,-n\}$ ; l'élimination de $P_0$ exclut la valeur propre nulle.

\begin{theorem}[Limite thêta des traces de chaleur]
\label{thm:limite-theta-triadico}
Pour tout $x>0$,
\[
 \lim_{m\to\infty}
 \operatorname{Tr}^{+}_m\!\left(e^{-xD_m^2}\right)
 =\Theta_+(x)
 :=\sum_{n\geq1}e^{-\pi n^2x}.
\]
\end{theorem}

\begin{proof}
La trace bilatérale, une fois le mode constant retiré, est
\[
 \sum_{\substack{n\in C_m\\n\neq0}}
 \exp\!\left[
 -\frac{N_m^2x}{\pi}
 \sin^2\!\left(\frac{\pi\widetilde n}{N_m}\right)
 \right].
\]
Pour chaque entier fixé $n$, le terme converge vers $e^{-\pi n^2x}$. De plus,
pour $|\widetilde n|\leq N_m/2$, l'inégalité
$\sin y\geq 2y/\pi$, valable sur $0\leq y\leq\pi/2$, fournit une
majorante gaussienne sommable et indépendante de $m$. La convergence
dominée appliquée à la somme discrète implique
\[
 \lim_{m\to\infty}
 \operatorname{Tr}\!\left(e^{-xD_m^2}P_0^\perp\right)
 =\sum_{n\in\ZZ\setminus\{0\}}e^{-\pi n^2x}.
\]
La division par deux établie dans \eqref{eq:traza-positiva} conserve une
contribution de chaque paire $\{n,-n\}$ et donne $\Theta_+(x)$.
\end{proof}

\section{Transformée thêta--Mellin}

Pour $\Re s>1$, nous définissons

\begin{equation}
 \mathcal Z_{3}(s)
 :=
 \frac{\displaystyle
 \int_0^\infty\Theta_+(x)x^{s/2-1}\,\dd x}
 {\pi^{-s/2}\Gamma(s/2)}.
 \label{eq:funcional-theta-mellin}
\end{equation}

La convergence absolue permet d'intégrer terme à terme. Comme

\[
 \int_0^\infty e^{-\pi n^2x}x^{s/2-1}\,\dd x
 =\Gamma(s/2)(\pi n^2)^{-s/2},
\]

on obtient

\begin{equation}
 \mathcal Z_3(s)=\sum_{n\geq1}n^{-s}=\zeta(s),
 \qquad \Re s>1.
 \label{eq:Z3-zeta}
\end{equation}

L'égalité se prolonge méromorphiquement par le prolongement de la fonction
zêta de Riemann \cite{Apostol1976}. Dans la suite, toutes les évaluations
proviennent de \eqref{eq:funcional-theta-mellin}, de ses dérivées ou de la
décomposition résiduelle de sa série de Dirichlet.

\begin{theorem}[Inversion thêta et équation fonctionnelle]
\label{pf84:C05}
Avec la notation opératorielle précédente,
\[
 \Theta_B(x):=\Tr(e^{-\pi xB^2})
 =1+2\Theta_+(x),
 \qquad
 Z(s):=\Tr(Q^{-s})=\mathcal Z_3(s)
 \quad(\Re s>1).
\]
La sommation de Poisson et la transformée de Mellin produisent
\[
 \Theta_B(x)=x^{-1/2}\Theta_B(1/x)
\]
et le prolongement de
\[
 \xi(s)=\frac12s(s-1)\pi^{-s/2}\Gamma(s/2)Z(s),
 \qquad
 \xi(s)=\xi(1-s).
\]
\end{theorem}

\begin{proof}
L'égalité \(\Theta_B=1+2\Theta_+\) sépare le mode zéro et apparie
\(\pm n\) ; l'identité \(Z=\mathcal Z_3\) est
\eqref{eq:Z3-zeta}. La sommation de Poisson appliquée à la gaussienne donne
l'inversion de \(\Theta_B\). Dans le demi-plan de convergence, l'intégration
terme à terme fournit
\[
 \pi^{-s/2}\Gamma(s/2)Z(s)
 =\int_0^\infty x^{s/2-1}\Theta_+(x)\,\dd x.
\]
En divisant l'intégrale en \((0,1)\) et \((1,\infty)\), en appliquant l'inversion
thêta à la première partie et en séparant les deux pôles classiques, on obtient le
prolongement et l'équation fonctionnelle. Cette unique preuve classique incorporée
ne localise pas les zéros de \(\zeta\).
\end{proof}

\section{Décomposition résiduelle modulo trois}

Pour $r\in\{0,1,2\}$, soit

\begin{equation}
 Z_r(s)=
 \sum_{\substack{n\geq1\\n\equiv r\pmod3}}n^{-s}.
 \label{eq:canales-residuales-Zr}
\end{equation}

En termes de la zêta de Hurwitz,

\[
 Z_0(s)=3^{-s}\zeta(s),
 \qquad
 Z_1(s)=3^{-s}\zeta\!\left(s,\frac13\right),
 \qquad
 Z_2(s)=3^{-s}\zeta\!\left(s,\frac23\right).
\]

Les trois termes ont pour résidu $1/3$ en $s=1$. Nous écrivons leurs
développements de Laurent sous la forme

\begin{equation}
 Z_r(1+z)-\frac1{3z}=\sum_{n\geq0}c_{r,n}\,z^n,
 \qquad
 \zeta(1+z)-\frac1z=\sum_{n\geq0}a_n\,z^n.
 \label{eq:coeficientes-Laurent-residuales}
\end{equation}

\begin{theorem}[Décomposition des coefficients de Laurent]
\label{thm:descomposicion-laurent-triadica}
Pour tout $n\geq0$,
\[
 a_n=c_{0,n}+c_{1,n}+c_{2,n},
 \qquad
 \gamma_n=(-1)^n n!\,a_n.
\]
Si $\chi_{-3}$ désigne le caractère réel non principal modulo trois, alors
\[
 \frac{L^{(n)}(1,\chi_{-3})}{n!}=c_{1,n}-c_{2,n}.
\]
De plus, avec \(\lambda=\log3\),
\[
 c_{0,n}=\frac13\left[
 \frac{(-\lambda)^{n+1}}{(n+1)!}
 +\sum_{k=0}^{n}a_k\frac{(-\lambda)^{n-k}}{(n-k)!}
 \right].
\]
Ces trois identités déterminent de manière unique
\((c_{0,n},c_{1,n},c_{2,n})\) à partir de
\((a_n,L^{(n)}(1,\chi_{-3}))\) et des coefficients agrégés d'ordre
inférieur.
\end{theorem}

\begin{proof}
La partition des entiers positifs en classes résiduelles fournit
$\zeta=Z_0+Z_1+Z_2$. En soustrayant les pôles en $s=1$ et en comparant les
coefficients de \eqref{eq:coeficientes-Laurent-residuales}, on obtient la
première identité. La deuxième coïncide avec la convention usuelle pour les
constantes de Stieltjes. De plus,
$Z_1-Z_2=L(s,\chi_{-3})$ ; les pôles de $Z_1$ et $Z_2$ s'annulent et la
comparaison des coefficients donne la troisième relation. Enfin,
\[
 Z_0(1+z)=\frac13e^{-\lambda z}\zeta(1+z).
\]
Le produit des deux séries de Laurent fournit la formule explicite
pour \(c_{0,n}\). La somme et la différence retrouvent alors
\(c_{1,n}\) et \(c_{2,n}\), ce qui prouve l'unicité affirmée.
\end{proof}

La somme des trois termes retrouve le secteur agrégé ; la différence des
classes $1$ et $2$ conserve l'orientation du caractère ; la comparaison
avec la classe divisible par trois enregistre le changement d'échelle. Ces trois
combinaisons sont linéairement indépendantes et déterminent les trois
coefficients $c_{r,n}$.

\section{Euler--Mascheroni et constantes de Stieltjes}

Soit

\[
 C_r:=\operatorname*{FP}_{s=1}Z_r(s).
\]

Les identités classiques pour la fonction digamma donnent
\cite{Apostol1976,Lagarias2013}

\begin{equation}
 C_0=\frac{\gamma-\log3}{3},
 \label{eq:C0-parte-finita}
\end{equation}

\begin{equation}
 C_1=\frac\gamma3+\frac{\log3}{6}+\frac{\pi}{6\sqrt3},
 \qquad
 C_2=\frac\gamma3+\frac{\log3}{6}-\frac{\pi}{6\sqrt3}.
 \label{eq:C1-C2-partes-finitas}
\end{equation}

\begin{corollary}
\label{cor:gamma-log3-L1}
Les parties finies satisfont
\[
 \boxed{\gamma=C_0+C_1+C_2},
\]
\[
 \boxed{L(1,\chi_{-3})=C_1-C_2=\frac{\pi}{3\sqrt3}},
\]
\[
 \boxed{\log3=C_1+C_2-2C_0}.
\]
\end{corollary}

\begin{proof}
À l'ordre zéro, le \cref{thm:descomposicion-laurent-triadica} donne
$a_0=c_{0,0}+c_{1,0}+c_{2,0}$. Par définition,
$a_0=\gamma$ et $c_{r,0}=C_r$, ce qui prouve la première égalité. La
différence des deux classes orientées est
$C_1-C_2=L(1,\chi_{-3})$. En substituant
\eqref{eq:C0-parte-finita}--\eqref{eq:C1-C2-partes-finitas}, on obtient les
deux autres identités.
\end{proof}

Une unique extraction des coefficients au moyen de la formule intégrale de
Cauchy appliquée à \eqref{eq:coeficientes-Laurent-residuales} détermine
simultanément $\gamma_0,\ldots,\gamma_6$. Les discrétisations à 96 et
192 nœuds coïncident jusqu'à l'ordre $10^{-87}$ ; l'identité des
coefficients provient toutefois du
\cref{thm:descomposicion-laurent-triadica} et non de cette comparaison numérique.

\section{Évaluation de la constante d'Apéry}

\begin{proposition}
\label{prop:evaluacion-apery}
La fonctionnelle \eqref{eq:funcional-theta-mellin} satisfait
\[
 \boxed{
 \mathcal Z_3(3)=\zeta(3)
 =1.202056903159594285399738161511449\ldots}.
\]
\end{proposition}

\begin{proof}
L'égalité exacte est la spécialisation $s=3$ de \eqref{eq:Z3-zeta}.
Pour certifier tout préfixe décimal sans modifier la fonctionnelle,
considérons $S_N=\sum_{n=1}^{N}n^{-3}$. Le critère intégral donne
\[
 S_N+\frac1{2(N+1)^2}
 <\zeta(3)<
 S_N+\frac1{2N^2}.
\]
Les intervalles rationnels emboîtés résultants déterminent les chiffres
imprimés.
\end{proof}

L'irrationalité de $\zeta(3)$ est le théorème classique d'Apéry
\cite{Apery1979} ; elle ne se déduit pas de la seule évaluation de la fonctionnelle.

\section{Constante de Glaisher--Kinkelin}

La dérivée spectrale en $s=-1$ définit

\begin{equation}
 \boxed{
 A_{\mathrm{GK}}
 =\exp\!\left(\frac1{12}-\mathcal Z_3'(-1)\right)}.
 \label{eq:glaisher-kinkelin-espectral}
\end{equation}

C'est la formule du déterminant régularisé associé au produit
hyperfactoriel \cite{Vardi1988}. Une dérivation centrale à cinq points,
évaluée à deux échelles, fournit

\[
 A_{\mathrm{GK}}
 =1.2824271291006226368753425688697917\ldots.
\]

Le terme $1/12$ appartient à la régularisation de l'hyperfactorielle et n'est pas
introduit comme paramètre d'ajustement.

\section{Constante d'Euler--Kronecker du corps d'Eisenstein}

Le caractère $\chi_{-3}$ détermine le corps quadratique imaginaire

\[
 F_{-3}=\QQ(\sqrt{-3}),
 \qquad
 \zeta_{F_{-3}}(s)=\zeta(s)L(s,\chi_{-3}).
\]

Sa constante d'Euler--Kronecker est

\begin{equation}
 \boxed{
 \gamma_{F_{-3}}
 =\gamma+\frac{L'(1,\chi_{-3})}{L(1,\chi_{-3})}}.
 \label{eq:euler-kronecker-eisenstein}
\end{equation}

Les coefficients résiduels de
\eqref{eq:coeficientes-Laurent-residuales} donnent

\[
 L(1,\chi_{-3})
 =0.6045997880780726168646927525473852\ldots,
\]
\[
 L'(1,\chi_{-3})
 =0.2226629869686015094866602627647443\ldots,
\]

et, par \eqref{eq:euler-kronecker-eisenstein},

\[
 \gamma_{F_{-3}}
 =0.9454972808716807032397499941581890\ldots.
\]

Le choix de $F_{-3}$ provient du caractère résiduel modulo trois, et non
d'une recherche sur les valeurs décimales de constantes de corps
quadratiques.

\section{Produit local associé aux nombres premiers jumeaux}

Pour un motif fini $H=\{h_1,\ldots,h_k\}\subset\ZZ$, nous définissons

\[
 \nu_p(H):=\left|\{-h_1,\ldots,-h_k\}\bmod p\right|,
 \qquad
 \sigma_p(H)=
 \frac{1-\nu_p(H)/p}{(1-1/p)^k}.
\]

Pour $H=\{0,2\}$, les facteurs locaux conduisent à

\begin{equation}
 C_2^{\mathrm{twin}}
 =\prod_{p>2}\frac{p(p-2)}{(p-1)^2}.
 \label{eq:constante-primos-gemelos}
\end{equation}

Soit \(P_X\) le produit restreint aux nombres premiers \(2<p\leq X\). Pour
\(p>X\), tous les facteurs appartiennent à \((0,1)\) ; en étendant le produit
de la queue des nombres premiers à tous les entiers, on obtient la borne
télescopique
\[
 \prod_{n>X}\left(1-\frac1{(n-1)^2}\right)
 =\frac{X-1}{X}
 \leq
 \prod_{p>X}\left(1-\frac1{(p-1)^2}\right)<1.
\]
Le produit \(P_{2\times10^6}\), évalué avec arrondi décimal dirigé, et
cette inégalité rationnelle pour la queue donnent

\begin{equation}
 0.6601615071224244826403933551
 <C_2^{\mathrm{twin}}<
 0.6601618372035081248537566591.
 \label{eq:intervalo-primos-gemelos}
\end{equation}

Le décompte direct jusqu'à $10^6$ contient 8169 paires de nombres premiers jumeaux. La
convergence du produit local et ce recensement fini n'impliquent pas
l'existence d'une infinité de paires ; cette dernière est une proposition arithmétique
distincte.

\section{Constante de Meissel--Mertens et hamiltonien premier}

La constante de Meissel--Mertens est définie par

\begin{equation}
 B_1=\lim_{x\to\infty}
 \left(\sum_{p\leq x}\frac1p-\log\log x\right)
 \label{eq:meissel-mertens-def}
\end{equation}

et satisfait

\begin{equation}
 B_1=\gamma+
 \sum_p\left(\log\!\left(1-\frac1p\right)+\frac1p\right).
 \label{eq:meissel-mertens-euler}
\end{equation}

Avec la notation de \cref{pf84:C08}, nous posons
\[
 \mathcal F_P:=\Gamma_s(\mathfrak h_{\mathbb P}),
 \qquad
 H_P:=d\Gamma(H_{\mathbb P}).
\]
Ni l'espace de Fock bosonique des modes premiers ni son hamiltonien ne sont redéfinis ici.
L'opérateur déjà construit a pour fonction de partition

\begin{equation}
 \operatorname{Tr}_{\mathcal F_P}(e^{-sH_P})
 =\prod_p\frac1{1-p^{-s}}
 =\zeta(s),
 \qquad \Re s>1.
 \label{eq:hamiltoniano-primo-zeta}
\end{equation}

Soit \(\mathcal H_P^{(1)}\) le sous-espace à une particule, engendré par
\(\{|p\rangle:p\text{ premier}\}\). Alors
\[
 \operatorname{Tr}_{\mathcal H_P^{(1)}}\!\left(
 e^{-H_P}\mathbf 1_{(-\infty,\log x]}(H_P)\right)
 =\sum_{p\leq x}\frac1p.
\]
Ainsi, \(B_1\) est la partie finie, lorsque \(x\to\infty\), de cette trace
spectrale tronquée après soustraction de \(\log\log x\). L'égalité
\eqref{eq:hamiltoniano-primo-zeta} présuppose la décomposition de Fock sur
les nombres premiers ; elle constitue donc une représentation opératorielle du produit
d'Euler, non une déduction de la primalité à partir du laplacien
triadique.

\section{Dépendance à l'égard de l'opérateur}

Les évaluations précédentes partagent l'opérateur
\eqref{eq:operador-espectral-Dm}, la trace positive
\eqref{eq:traza-positiva} et la transformée
\eqref{eq:funcional-theta-mellin}. D'autres constantes exigent des dynamiques
distinctes. La constante de Catalan est associée au caractère modulo quatre
$\chi_{-4}$, tandis que la constante de Khinchin appartient à l'opérateur de
Gauss et à sa mesure invariante. Un préfixe de l'une quelconque d'entre elles peut
apparaître dans un catalogue fini sans que la fonctionnelle triadique la sélectionne.
Cette dépendance opératorielle distingue l'évaluation spectrale de la
capacité purement combinatoire de codage.

\end{HMTScientificOwnerSubordinated}

\HMTFlushCurrentChapter
\chapter{Criticité unimodale de la famille dodécaphasique : profil analytique exact, approximants de Feigenbaum et certificat hyperbolique de renormalisation}%
\begingroup
\renewcommand{\chapter}[2][]{}%
\chapter{Criticité unimodale de la famille dodécaphasique : profil analytique exact, approximants de Feigenbaum et réduction de l’universalité au certificat hyperbolique}
\label{chap:p3-final:45:feigenbaum}
\label{mab:rm:ch:chaos}

\section{Construction intrinsèque de la famille dodécaphasique}

Les constantes de Feigenbaum caractérisent l’universalité du
doublement de période. Leur dérivation dans HMT exige de construire au préalable
une famille analytique unimodale, de démontrer sa criticité quadratique et
d’étudier ensuite l’action de l’opérateur de renormalisation.

L’application initiale ne procède pas d’un choix harmonique externe. Le retour
de \(\TPK\) détermine la roue dodécaphasique orientée
\(U_{12}^{\rm sign}\). Dans le secteur uniforme \(\eta=0\), la trace
normalisée de ses douze secteurs fixe \(w_m=1/12\), et la fonctionnelle de phase
\begin{equation}
\mathcal P_{\rm crit}:U_{12}^{\rm sign}\longrightarrow\R/2\pi\Z,
\qquad
m\longmapsto \frac{(30m-10)\pi}{180}
=\pi\frac{3m-1}{18}
\label{mab:rm:eq:critical-phase-reader}
\end{equation}
conserve l’ordre et l’orientation. La composition causale est
\[
\APP\to\TRIT\to\TPK\to U_{12}^{\rm sign}
\xrightarrow{\mathcal P_{\rm crit}}
\{(\phi_m,w_m)\}_{m=1}^{12}
\to u_0\to f_\lambda.
\]
Toute modification de \(\mathcal P_{\rm crit}\), des poids ou de l’ordre des
phases altère la famille et constitue un critère de réfutation. Ces
choix doivent être fixés avant la confrontation avec \(\delta_F\).

La famille dodécaphasique résultante part de

\[
a_m=\frac{3m-1}{18},\qquad
\phi_m=\pi a_m,\qquad m=1,\ldots,12,
\]

et du deuxième moment

\[
D_0=\frac1{12}\sum_{m=1}^{12}a_m^2.
\]

L’identité

\begin{equation}
\sum_{m=1}^{12}(3m-1)^2=5394=6\cdot899
\label{mab:rm:eq:chaos-sum2}
\end{equation}

produit

\[
D_0=\frac{899}{648},\qquad
s_0=\frac{36}{\pi\sqrt{899}}.
\]

Par conséquent,

\begin{equation}
s_0\phi_m=\frac{2(3m-1)}{\sqrt{899}}.
\label{mab:rm:eq:pi-cancellation}
\end{equation}

Le facteur \(\pi\) s’annule avant d’effectuer les itérations. La famille ne
dépend ni de \(\delta_F\) ni de \(\alpha_F\) comme données d’entrée.

L’annulation est une propriété structurelle de la normalisation. Avant de
normaliser,
le deuxième moment angulaire est
\[
\frac1{12}\sum_{m=1}^{12}\phi_m^2
=\pi^2D_0.
\]
La condition de criticité quadratique fixe
\(s_0^2\pi^2D_0=2\), et chaque phase normalisée dépend donc seulement de
l’arithmétique entière de \(3m-1\). La fonctionnelle conserve la provenance angulaire de
la roue, mais l’itération reçoit une famille réelle intrinsèque. C’est
précisément ce qui permet de distinguer une dérivation d’un étalonnage
par \(\pi\).

\section{Représentation analytique exacte et expression trigonométrique fermée}

On définit

\begin{equation}
u_0(x)=1-\frac1{12}\sum_{m=1}^{12}
\cos\!\left(\frac{2(3m-1)}{\sqrt{899}}x\right).
\label{mab:rm:eq:chaos-profile}
\end{equation}

Le développement à l’origine est

\begin{equation}
u_0(x)=x^2-\frac{715769}{2424603}x^4+O(x^6).
\label{mab:rm:eq:chaos-taylor}
\end{equation}

La somme finie admet l’expression fermée suivante. Si \(y=x/\sqrt{899}\),

\begin{equation}
u_0(x)=1-\frac{\cos(37y)\sin(36y)}{12\sin(3y)}
=1-\frac{\sin(73y)-\sin y}{24\sin(3y)}.
\label{mab:rm:eq:chaos-closed}
\end{equation}

La singularité apparente en \(y=0\) est éliminable et reproduit \cref{mab:rm:eq:chaos-taylor}.

\begin{proof}[Dérivation de la forme fermée]
Les douze arguments forment une progression arithmétique :
\[
2(3m-1)y=(6m-2)y,\qquad m=1,\ldots,12.
\]
La somme des cosinus d’une progression satisfait
\[
\sum_{m=1}^{N}\cos(a+md)
=\frac{\sin(Nd/2)}{\sin(d/2)}
 \cos\!\left(a+\frac{N+1}{2}d\right).
\]
En prenant \(N=12\), \(d=6y\) et \(a=-2y\), on obtient
\[
\sum_{m=1}^{12}\cos((6m-2)y)
=\frac{\sin36y}{\sin3y}\cos37y.
\]
L’identité de transformation produit--en--somme
\(2\cos37y\sin36y=\sin73y-\sin y\) donne la deuxième expression de
\cref{mab:rm:eq:chaos-closed}.
\end{proof}

Le coefficient quartique n’est pas non plus ajusté. Les sommes entières
\begin{equation}
\sum_{m=1}^{12}(3m-1)^2=5394,\qquad
\sum_{m=1}^{12}(3m-1)^4=4294614
\label{mab:rm:eq:chaos-moments24}
\end{equation}
insérées dans le développement de \(\cos z\) produisent exactement
\[
\frac1{24\cdot12}\sum_m
\left(\frac{2(3m-1)}{\sqrt{899}}\right)^4
=\frac{715769}{2424603}.
\]
Ainsi, tant la normalisation quadratique que la première correction non
linéaire sont des moments finis de la même roue.

\section{Appartenance à la classe analytique unimodale quadratique}

Considérons

\[
f_\lambda(x)=1-\lambda u_0(x),\qquad \lambda>0.
\]

Pour \(0<x\le1\), tous les arguments

\[
\frac{2(3m-1)}{\sqrt{899}}x
\]

appartiennent à \((0,\pi)\). Par conséquent,

\[
u_0'(x)>0,\qquad u_0'''(x)<0.
\]

Comme \(u_0\) est paire, \(x=0\) est son unique point critique et \(u_0''(0)=2\). La dérivée schwarzienne de \(f_\lambda\) est

\begin{equation}
S(f_\lambda)
=\frac{u_0'''}{u_0'}
-\frac32\left(\frac{u_0''}{u_0'}\right)^2<0
\qquad(0<|x|\le1).
\label{mab:rm:eq:negative-schwarzian}
\end{equation}

\begin{theorem}[Appartenance exacte à la classe unimodale]
Pour
\[
0<\lambda\le \frac{2}{u_0(1)},
\]
la famille \(f_\lambda=1-\lambda u_0\) est une application réelle--analytique de \([-1,1]\) dans lui-même, paire, avec un unique point critique quadratique à l’origine et une dérivée schwarzienne négative en dehors de ce point. Elle appartient donc, sans utiliser les valeurs de Feigenbaum comme valeur cible, à la classe \(S\)-unimodale quadratique sur laquelle agit le renormalisateur de doublement de période \cite{rm:Feigenbaum1978,rm:Lanford1982}. Pour \(\lambda\) en dehors de cet intervalle, le germe analytique et sa forme locale sont conservés, mais l’application de l’intervalle dans lui-même n’est pas automatiquement affirmée.
\end{theorem}

Le théorème établit la contribution intrinsèque de HMT préalable au théorème
d’universalité : le germe analytique et son type de criticité.

\begin{proof}[Détails du signe et de l’application dans soi-même]
En dérivant \cref{mab:rm:eq:chaos-profile},
\begin{align*}
u_0'(x)&=\frac1{12}\sum_m k_m\sin(k_mx),\\
u_0'''(x)&=-\frac1{12}\sum_m k_m^3\sin(k_mx),
\qquad k_m=\frac{2(3m-1)}{\sqrt{899}}.
\end{align*}
Pour \(0<x\le1\), \(0<k_mx<\pi\), donc tous les termes de la première
ligne sont positifs et ceux de la seconde négatifs. La parité étend la
monotonie aux deux côtés de l’origine. Comme \(u_0(0)=0\) et
\(0\le u_0(x)\le u_0(1)\), la borne
\(0<\lambda\le2/u_0(1)\) implique
\(-1\le1-\lambda u_0(x)\le1\). Enfin,
\(u_0''(0)=1/12\sum k_m^2=2\), de sorte que le point critique est
quadratique.
\end{proof}

\section{Suites finies de paramètres superstables}

La suite superstable certifiée jusqu’au niveau huit produit

\[
\delta_8=4.66921218915\ldots,
\qquad
\alpha_8=2.50296847988\ldots.
\]

Son statut est celui de \emph{Certificat fini exempt de valeurs cibles} : chaque
approximant est obtenu à partir de la famille préalablement définie, sans introduire
comme donnée la limite universelle. La coïncidence décimale fournit un
indicateur de convergence, mais ne remplace pas le théorème d’hyperbolicité.

Une extension numérique indépendante de la même suite, avec un pas
adapté à la séparation entre les racines et un contrôle de la période exacte,
atteint provisoirement :

\begin{center}
\small
\begin{tabular}{@{}ccll@{}}
\toprule
Niveau & \(\lambda_n\) & \(\delta_n\) & \(\alpha_n\)\\
\midrule
11 & 1.87403828195439908045 & 4.66920170694424982745 & 2.50290847517165181764\\
12 & 1.87403836411023460451 & 4.66920163036308158848 & 2.50290800379001546915\\
13 & 1.87403838170549870372 & 4.66920161361839219913 & 2.50290790268748193108\\
14 & 1.87403838547386545431 & 4.66920161007472591023 & 2.50290788100969754274\\
\bottomrule
\end{tabular}
\end{center}

Ces niveaux ont le statut de \emph{Nouveau certificat numérique
en attente d’un enregistrement reproductible} ; ils ne remplacent pas le certificat persistant
jusqu’au niveau huit. Une extrapolation d’Aitken fournit
\(\lambda_\infty^{\HMT}\simeq1.8740383865008918080\) et
\(\alpha_F\simeq2.50290787509\) comme candidats à la limite, sous réserve de
la démonstration de l’hyperbolicité et de la transversalité.

Pour éviter une procédure numérique non explicitée, la méthode est formulée comme
suit. Définissez
\begin{equation}
F_n(\lambda)=f_\lambda^{\,2^n}(0).
\label{mab:rm:eq:superstable-equation}
\end{equation}
La racine superstable \(\lambda_n\) est la première racine de \(F_n\) à
droite de \(\lambda_{n-1}\) qui satisfait simultanément
\[
f_{\lambda_n}^{\,2^j}(0)\ne0
\quad(0\le j<n).
\]
Cette inégalité exclut les périodes diviseurs. Avec des intervalles disjoints
\(I_n\ni\lambda_n\), on calcule
\begin{equation}
\delta_n=
\frac{\lambda_{n-1}-\lambda_{n-2}}
     {\lambda_n-\lambda_{n-1}},
\qquad
\alpha_n=-\frac{d_{n-1}}{d_n},
\label{mab:rm:eq:feigenbaum-approximants}
\end{equation}
où \(d_n=f_{\lambda_n}^{\,2^{n-1}}(0)\) est la distance orientée du
dernier point nouveau au point critique. Un certificat persistant doit conserver
les intervalles de racine, l’exclusion des périodes inférieures et la propagation
de l’erreur de \cref{mab:rm:eq:feigenbaum-approximants} ; un simple nombre décimal imprimé ne
remplit pas ces trois conditions.

La règle de discrétisation constitue également un contrôle négatif. Une
borne inférieure fixe supérieure à la séparation
\(|\lambda_n-\lambda_{n-1}|\) peut manquer la première racine et produire
néanmoins une suite numériquement plausible. C’est pourquoi l’incrément
doit être proportionné à la séparation précédente et chaque candidat doit être vérifié
par sa période exacte.

\section{Conditions de certification de l’opérateur de renormalisation}

Sur un espace de fonctions paires analytiques, nous définissons

\begin{equation}
\mathcal R(f)(x)=-a(f)^{-1}f\bigl(f(-a(f)x)\bigr),
\qquad a(f)=-f(1).
\label{mab:rm:eq:renormalization}
\end{equation}

La clôture complète de \(\delta_F\) et de \(\alpha_F\) exige de certifier conjointement :

\begin{enumerate}
\item l’existence d’un point fixe \(g=\mathcal R(g)\) dans une boule analytique déclarée ;
\item l’unicité locale de ce point fixe ;
\item une unique valeur propre réelle instable de \(D\mathcal R_g\), égale à \(\delta_F\) ;
\item un rayon spectral strictement inférieur à un sur le complément stable ;
\item la transversalité de la courbe \(\lambda\mapsto f_\lambda\) par rapport à la variété stable de \(g\).
\end{enumerate}

Une implémentation HMT naturelle emploie des troncatures de Chebyshev d’ordre
\(9^m\), de l’arithmétique d’intervalles et des polynômes de rayons. La preuve doit
fournir une boule, un inverse approché, une borne du défaut et une borne
de Lipschitz satisfaisant \(Y+Z(r)<r\). Ainsi, la profondeur
nonadique paramètre un certificat de point fixe et non une simple suite
d’approximations décimales.

Dans une base paire de Chebyshev,
\[
g(x)=\sum_{j=0}^{N}g_{2j}T_{2j}(x),\qquad N=9^m,
\]
le résidu fini est \(F_N(g)=\Pi_N(\mathcal Rg-g)\). Si \(A_N\) approche
\((DF_N(\bar g))^{-1}\), les quantités à certifier sont
\begin{align*}
Y&=\|A_NF_N(\bar g)\|,\\
Z_0&=\|I-A_NDF_N(\bar g)\|,\\
Z_1(r)&=\sup_{\|h\|\le r}
\|A_N(DF_N(\bar g+h)-DF_N(\bar g))\|.
\end{align*}
Une inégalité
\(Y+(Z_0+Z_1(r))r<r\) enferme un unique point fixe. Le bloc spectral
ultérieur doit isoler une direction instable et contracter le complément. L’usage
de \(9^m\) intervient ainsi comme loi de raffinement du certificat, non comme
décoration numérologique.

\begin{corollary}[Confluence exacte du mode trente]
Le mode \(30\) admet trois réalisations simultanées et typologiquement
distinctes :
\begin{align*}
30&\longmapsto(3\cdot30,4\cdot30)=(90,120)
&&\text{dans la complétion constitutive},\\
30&\longmapsto30^2-1=899=29\cdot31
&&\text{dans la normalisation critique},\\
30&\longmapsto
M_{30}=\begin{pmatrix}30&899\\1&30\end{pmatrix}
&&\text{dans la transformation hyperbolique de Pell}.
\end{align*}
Les trois identités sont exactes et sont obtenues sans prescrire de valeurs cibles. L’invariant conservé est
l’entier central \(30\) avec son orientation : extension \(3\to4\),
défaut quadratique \(-1\) et unité de norme un. Ce qui n’est pas affirmé est
l’égalité entre l’opérateur de vide et le renormalisateur de doublement.
\end{corollary}

\begin{remark}[Contrôle négatif]
Une deuxième valeur propre en dehors du disque unité, une boule de rayons sans solution, une perte d’unicité ou la tangence de la famille HMT à la variété stable réfutent la promotion universelle. Aucune ne réfute le germe exact ni le théorème de la dérivée schwarzienne.
\end{remark}

\begin{proposition}[Résultat principal]
La signification structurelle de la constante de doublement réside dans
l’annulation dodécaphasique de \(\pi\), dans l’identité
\(899=30^2-1\), dans la forme trigonométrique fermée et dans l’appartenance
démontrée à la classe analytique sur laquelle agit le renormalisateur. Son
développement décimal constitue exclusivement l’évaluation terminale.
\end{proposition}

\subsection*{Portée logique et domaine de validité}
Profil, Taylor, normalisateur, unité de Pell, unimodalité et dérivée schwarzienne : \emph{Théorème interne exact}. Approximants jusqu’au niveau huit : \emph{Certificat fini persistant}. Niveaux 11--14 : \emph{Nouveau certificat numérique en attente de pérennisation}. Hyperbolicité et transversalité du renormalisateur : \emph{Programme ouvert avec critère de réfutation exécutable}.


\subsection{Déformation harmonique pondérée et sélection de la branche superstable}
\label{mab:sec:extension-dodecafasica-ponderada-armonica}

Le profil uniforme possède une extension paramétrique qui conserve les douze
phases et sépare la déformation harmonique du paramètre dynamique. Pour
\(1\le m\le12\), soient

\begin{equation}
 \phi_m=\frac{(30m-10)\pi}{180},
 \qquad
 p(\phi)=12\cos(3\phi)-\cos(4\phi).
 \label{mab:eq:fases-y-deformacion-armonica-armonizada}
\end{equation}

Pour tout \(\eta\) réel tel que les poids restent positifs, nous définissons

\begin{equation}
 w_m(\eta)=\frac{1+\eta p(\phi_m)}{12}.
 \label{mab:eq:pesos-dodecafasicos-armonizados}
\end{equation}

Les sommes des harmoniques de rang trois et quatre sur les douze phases s’annulent ;
par conséquent,

\[
 \sum_{m=1}^{12}w_m(\eta)=1.
\]

Soit

\begin{equation}
 s_\eta^2=\frac{2}{\sum_{m=1}^{12}w_m(\eta)\phi_m^2},
 \qquad
 u_\eta(x)=1-\sum_{m=1}^{12}w_m(\eta)
 \cos(s_\eta\phi_m x).
 \label{mab:eq:familia-ponderada-armonizada}
\end{equation}

Le développement du cosinus et la définition de \(s_\eta\) donnent

\[
 u_\eta(x)=x^2+O(x^4)\qquad(x\to0).
\]

Par conséquent, \(u_\eta\) est réelle--analytique, paire et conserve un point
critique quadratique à l’origine. Nous introduisons le paramètre dynamique par

\begin{equation}
 f_{\lambda,\eta}(x)=1-\lambda u_\eta(x).
 \label{mab:eq:familia-feigenbaum-ponderada-armonizada}
\end{equation}

Un paramètre superstable de période exacte \(2^n\) satisfait

\[
 f_{\lambda,\eta}^{\,2^n}(0)=0,
 \qquad
 f_{\lambda,\eta}^{\,2^k}(0)\ne0\quad(0\le k<n).
\]

L’équation peut posséder plusieurs branches. On fixe une branche
\(\mathscr B\) en parcourant, immédiatement après le paramètre du
niveau précédent, un maillage adaptatif ; on prend le premier intervalle avec changement
de signe et l’on applique une dichotomie. La notation
\(\lambda_n^{\mathscr B}(\eta)\) incorpore ce sélecteur et n’attribue pas
d’unicité globale à l’équation superstable. Le premier niveau est déterminé
exactement :

\[
 \lambda_1^{\mathscr B}(\eta)=\frac1{u_\eta(1)},
\]

si \(u_\eta(1)>0\), parce que
\(f_{\lambda,\eta}^{\,2}(0)=f_{\lambda,\eta}(1)=0\).
Étant donnés trois niveaux consécutifs et

\[
 x_n^{\mathscr B}(\eta)
 :=f_{\lambda_n^{\mathscr B}(\eta),\eta}^{\,2^{n-1}}(0),
\]

on forme

\begin{equation}
 \delta_n^{\mathscr B}(\eta)
 =\frac{\lambda_{n-1}^{\mathscr B}(\eta)-
         \lambda_{n-2}^{\mathscr B}(\eta)}
        {\lambda_n^{\mathscr B}(\eta)-
         \lambda_{n-1}^{\mathscr B}(\eta)},
 \qquad
 \alpha_n^{\mathscr B}(\eta)
 =\frac{|x_{n-1}^{\mathscr B}(\eta)|}
        {|x_n^{\mathscr B}(\eta)|}.
 \label{mab:eq:aproximantes-ponderados-armonizados}
\end{equation}

Pour \(\eta=0\), on retrouve exactement le profil uniforme. L’extension
pondérée introduit un paramètre transversal : elle permet d’étudier
la transversalité de la famille par rapport à la variété stable du
renormalisateur sans transformer une coïncidence finie en une preuve
d’universalité.

\begin{remark}[Contrôle négatif]
L’extension est réfutée si les poids perdent leur positivité dans le domaine
déclaré, si la normalisation cesse de produire un coefficient quadratique
unitaire, si le sélecteur n’exclut pas les périodes diviseurs ou si une branche est choisie
en utilisant comme entrée les valeurs universelles terminales.
\end{remark}

\section{Réduction de l’universalité de Feigenbaum au certificat hyperbolique de bassin et de transversalité}
\label{sec:p3-final:cierre-feigenbaum}

La famille de Feigenbaum HMT possède un profil analytique, un maximum quadratique,
une dérivée schwarzienne négative et des approximants superstables certifiés jusqu’au niveau
\(8\). Sa promotion à l’universalité exige le certificat
\[
 \ell_u(\partial_t f_t)\ne0
\]
ainsi que l’appartenance au bassin hyperbolique du point fixe. La tangence,
la sortie du bassin ou la présence d’une deuxième valeur propre instable sont
des réfutateurs complets.

%
\endgroup

\HMTFlushCurrentChapter
\chapter{Arithmétique de la famille modulaire dodécaphasique de période \(30\): identités de Rogers--Ramanujan, normalisateur \(899\) et relèvement nonadique de Pell}%
\begingroup
\renewcommand{\chapter}[2][]{}%
\chapter{Arithmétique du mode \(30\) : identités de Rogers--Ramanujan, normalisateur \(899\) et relèvement nonadique de Pell}
\label{chap:p3-final:46:modo_30_pell}
\section{Normalisateur \(899\), unité de Pell et structure arithmétique du mode \(30\)}

Le normalisateur possède la factorisation arithmétique

\begin{equation}
899=30^2-1=29\cdot31.
\label{mab:rm:eq:899}
\end{equation}

Le même mode \(30\) qui engendre les secteurs \(90\) et \(120\) détermine
simultanément la paire de nombres premiers jumeaux \(29,31\). De plus,

\begin{equation}
\varepsilon_{30}=30+\sqrt{899},\qquad
\varepsilon_{30}^{-1}=30-\sqrt{899},\qquad
\log\varepsilon_{30}=\arcosh30.
\label{mab:rm:eq:pell30}
\end{equation}

Cette unité de Pell détermine l’échelle hyperbolique associée au
normalisateur et relie la phase dodécaphasique, le mode \(30\), les nombres premiers
jumeaux et la géométrie hyperbolique. Cette relation n’identifie pas
l’opérateur constitutif du vide au renormalisateur de doublement ; une telle
identification exigerait un entrelaceur explicite.

L’unité agit de manière entière par
\begin{equation}
M_{30}=\begin{pmatrix}30&899\\1&30\end{pmatrix},
\qquad \det M_{30}=1.
\label{mab:rm:eq:pell-matrix30}
\end{equation}
Sur \(\Z[\sqrt{899}]\), cette matrice représente la multiplication par
\(30+\sqrt{899}\). Ses valeurs propres sont \(\varepsilon_{30}^{\pm1}\) ; par
conséquent, \(\log\varepsilon_{30}\) est la longueur de translation de la transformation
hyperbolique associée. Le même entier central \(30\) réapparaît dans le vide
parce que \(90=3\cdot30\) et \(120=4\cdot30\), mais les deux opérateurs continuent
d’avoir des domaines distincts. La confluence démontrée est celle du mode et de ses
complétions ; l’identité des dynamiques n’est pas postulée.


\section{Relèvement nonadique de l'unité de Pell}

Pour \(k\ge1\), soit l'anneau quadratique non ramifié
\begin{equation}
\mathcal R_k=(\Z/3^k\Z)[r]/(r^2-2),
\label{eq:pell-ring}
\end{equation}
et soit

\[
u=3+2r.
\]

La conjugaison \(r\mapsto-r\) donne
\[
N_{\mathcal R_k/(\Z/3^k\Z)}(u)
=(3+2r)(3-2r)=9-8=1,
\]
de sorte que \(u\) est une unité de norme un. En outre,
\[
u^2=17+12r,\qquad
u^4-1=576+408r=3(192+136r).
\]
Le second facteur est une unité modulo trois parce qu'il se réduit à \(r\ne0\). Par
conséquent, la valuation \(3\)-adique satisfait
\(v_3(u^4-1)=1\). Le lemme de relèvement pour une unité congrue à un
modulo trois produit, pour tout \(j\ge0\),
\[
v_3\!\left(u^{4\cdot3^j}-1\right)=j+1.
\]
Comme \(u\bmod3=2r\) et \((2r)^2=-1\), son ordre modulo trois est quatre.
En conséquence,

\begin{equation}
\operatorname{ord}(u)=4\cdot3^{k-1},
\label{eq:pell-order}
\end{equation}

et la clôture profinie des groupes cycliques compatibles est
\begin{equation}
\varprojlim_k\langle u\bmod3^k\rangle
\simeq C_4\times\mathbb Z_3.
\label{eq:pell-profinite}
\end{equation}
Sa réalisation archimédienne est

\[
3+2\sqrt2=(1+\sqrt2)^2.
\]

Cette construction démontre la conservation de la profondeur nonadique
par une unité hyperbolique. Les unités \(3+2\sqrt2\),
\(2+\sqrt3\) et \(30+\sqrt{899}\) forment une famille de Pell de
provenances distinctes ; toute identification entre elles exige des
applications explicites entre les extensions correspondantes.


Le mode \(30\) est l’un des points où l’architecture commence à montrer une unité déterminée par la généalogie opératorielle et arithmétique.\par

Il apparaît d’abord dans le vide parce que\par

\[
90=3\cdot30,
\qquad
120=4\cdot30.
\]

C’est la cadence élémentaire qui permet de comparer les secteurs \(90\) et \(120\). Le vide lit cette cadence comme une complétion \(3\to4\).\par

Mais le même \(30\) réapparaît dans la construction du germe de criticité. Là, l’identité centrale est\par

\[
899=30^2-1=29\cdot31.
\]

La normalisation des douze phases produit \(\sqrt{899}\), et le profil analytique peut s’écrire sous forme fermée. Ce qui est réellement frappant, c’est que \(\pi\) disparaît avant le début de l’itération.\par

Cela doit être lu attentivement. La géométrie angulaire de \(\pi\) est absorbée dans l’organisation des phases. Une fois la configuration dodécaphasique normalisée, le germe dynamique est gouverné par l’arithmétique interne de \(899\).\par

C’est pourquoi la construction commence par engendrer intérieurement un objet dynamique propre :\par

\begin{itemize}[leftmargin=*,itemsep=0.35em]
\item es par;
\item il est analytique ;
\item il a un point critique quadratique ;
\item il appartient à la classe unimodale correspondante ;
\item il possède une dérivée schwarzienne négative ;
\item il peut être soumis à l’opérateur de renormalisation.
\end{itemize}

Les rapports superstables commencent alors à s’approcher des constantes universelles. Le germe a été produit en premier et l’universalité apparaît comme comportement terminal de sa dynamique.\par

C’est une différence épistémologique profonde. Une coïncidence numérique dirait : « notre calcul ressemble à Feigenbaum ». La construction HMT dit : « voici l’objet exact qui entre dans la classe d’universalité, et voici les quotients qui émergent lorsqu’on l’itère ».\par

Le \(30\) réapparaît encore une troisième fois dans la mémoire modulaire. Si\par

\[
N=N_{90}+N_{120}
\]

mesure l’occupation totale et\par

\[
D=90N_{90}+120N_{120}
\]

mesure l’occupation pondérée par la provenance, la matrice qui transforme \((N_{90},N_{120})\) en \((N,D)\) a pour déterminant \(30\).\par

Cela permet d’inverser le processus :\par

\[
N_{90}=4N-\frac D{30},
\qquad
N_{120}=\frac D{30}-3N.
\]

De sorte que le même \(30\) remplit trois fonctions :\par

\begin{itemize}[leftmargin=*,itemsep=0.35em]
\item dans le vide, il rend commensurables les propagations \(90\) et \(120\) ;
\item dans la criticité, il détermine la normalisation \(899=30^2-1\) ;
\item dans la mémoire modulaire, il est le déterminant qui permet de retrouver l’origine de chaque occupation.
\end{itemize}

L’opérateur du vide, le germe de chaos et le hamiltonien modulaire restent typologiquement distincts. La confluence est plus subtile : tous trois conservent une partition commune de la généalogie.\par

Nous pourrions dire que le mode \(30\) est un engrenage. Dans une machine, il mesure la propagation ; dans une autre, il normalise la criticité ; dans une autre, il rend réversible la perte apparente de provenance. L’engrenage est le même, mais chaque machine remplit une fonction différente.\par

L’apparition de l’unité de Pell\par

\[
30+\sqrt{899}
\]

renforce cette interprétation. Le \(30\) organise simultanément une somme finie et une croissance hyperbolique. La même arithmétique qui clôt le germe ouvre une dynamique d’échelle.\par

C’est peut-être la raison profonde pour laquelle le résultat est si beau : le mode qui rend possible la complétion finie du vide est aussi le mode qui prépare une universalité infinie de la criticité.\par

Fini et infini s’articulent par la répétition cohérente d’une structure finie qui conserve la mémoire ; de celle-ci naît l’infini dynamique.\par


%
\endgroup
\begin{HMTSpiralChapterBody}
% Corte mecánico íntegro: parte_ii.tex:323--419.
% Residencia material del ensamblaje sucesor de 102 capítulos.
% Residencia causal: capítulo 46.

\chapter{Vis spectrale de Pell et suspension logarithmique}

\section{Loi de vis}

Le corpus contient l'identité exacte
\begin{equation}
 V^{-1}\mathscr S V=(2+\sqrt3)\ii\,\mathscr S.
 \label{espiral:eq:tornillo-pell}
\end{equation}
Soit \(\varepsilon=2+\sqrt3\). L'orbite transverse
\[
 z_n=z_0(\varepsilon\ii)^n
\]
satisfait
\[
 r_n=r_0\varepsilon^n,
 \qquad
 \theta_n=\theta_0+\frac{\pi n}{2}.
\]
Eliminando \(n\),
\begin{equation}
 r(\theta)=r_0
 \exp\!\left[
 \frac{2\log(2+\sqrt3)}{\pi}
 (\theta-\theta_0)
 \right].
 \label{espiral:eq:espiral-log-pell}
\end{equation}
La spirale logarithmique est une sortie de l'opérateur : son équation n'a pas été choisie pour ajuster ensuite \(\varepsilon\).

\begin{proof}[Dérivation de \cref{espiral:eq:espiral-log-pell}]
De \(\theta_n-\theta_0=n\pi/2\) découle \(n=2(\theta_n-\theta_0)/\pi\). En substituant dans \(r_n=r_0\varepsilon^n\),
\[
 r_n=r_0\exp\left(
 \frac{2\log\varepsilon}{\pi}(\theta_n-\theta_0)
 \right).
\]
L'interpolation continue est l'équation indiquée.
\end{proof}

\section{Involution conjuguée de contraction}

L'extension bilatérale \(n\in\Z\) produit
\[
 n\to+\infty:\ r_n\to\infty,
 \qquad
 n\to-\infty:\ r_n\to0.
\]
L'inversion \(n\mapsto-n\) échange expansion et contraction. Comme \(\varepsilon^{-1}=2-\sqrt3\), les deux branches appartiennent à la même unité de Pell. Ce ne sont pas deux spirales indépendantes, mais deux orientations de la même vis spectrale.

\begin{figure}[htbp]
\centering
\begin{tikzpicture}[scale=.82]
 \draw[->,HMTGray] (-5.1,0)--(5.1,0) node[right] {$x$};
 \draw[->,HMTGray] (0,-3.1)--(0,3.1) node[above] {$y$};
 \draw[HMTBlue,very thick,domain=-7:5.2,samples=240,smooth,variable=\t]
   plot ({.27*exp(.22*\t)*cos(deg(\t))},{.27*exp(.22*\t)*sin(deg(\t))});
 \draw[HMTRed,thick,domain=-7:5.2,samples=240,smooth,variable=\t]
   plot ({.27*exp(.22*\t)*cos(deg(-\t))},{.27*exp(.22*\t)*sin(deg(-\t))});
 \foreach \k in {-4,-3,...,4}{
   \pgfmathsetmacro{\ang}{1.57079632679*\k}
   \fill[HMTNavy]
     ({.27*exp(.22*\ang)*cos(deg(\ang))},
      {.27*exp(.22*\ang)*sin(deg(\ang))}) circle (.8pt);
 }
 \node[HMTBlue,font=\small\sffamily] at (3.5,2.4) {expansion};
 \node[HMTRed,font=\small\sffamily] at (3.5,-2.4) {orientation conjuguée};
\end{tikzpicture}
\caption{Projection transverse de la vis de Pell. Les marques représentent des quarts de tour.}
\label{espiral:fig:tornillo-pell}
\end{figure}

\section{Covariance d'échelle et conservation énergétique}

La dilatation \(\varepsilon\) ne doit pas être interprétée comme une croissance gratuite de l'amplitude énergétique d'une onde dans le vide. Sa réalisation cohérente agit sur l'échelle spatiale ou spectrale :
\[
 (r,\varphi,\lambda)
 \longmapsto
 (\varepsilon r,\varphi+\pi/2,\varepsilon\lambda).
\]
Cette distinction sera essentielle dans la réalisation électromagnétique de la \cref{espiral:chap:realizacion-em}.

\begin{resultbox}[title={Compatibilité exacte entre la vis spectrale de Pell et la suspension phase--mémoire}]
Le relèvement phase--mémoire minimal est une hélice circulaire. La vis spectrale combine un quart de tour et une dilatation de Pell ; sa projection est une spirale logarithmique. Tous deux proviennent du registre HMT, mais ne sont pas la même périodicité et ne s'obtiennent pas l'un à partir de l'autre par oubli des types.
\end{resultbox}

\end{HMTSpiralChapterBody}

\HMTFlushCurrentChapter
\chapter{Théorie arithmétique des vacances nonadiques : défaut \texorpdfstring{\(10^3-3^6\)}{10^3-3^6}, réseau radical \(3\to6\to9\), loi puissance--vacance et torsion de bord}%
\begingroup
\renewcommand{\chapter}[2][]{}%
\chapter{Théorie arithmétique des lacunes nonadiques : défaut \(10^3-3^6\), réseau radical \(3\to6\to9\), loi puissance--lacune et torsion de bord}
\label{chap:p3-final:47:vacancias_nonadicas}
\section{Classification des cycles arithmétiques irréductibles du réseau radical nonadique}

Dans la réalisation arithmétique, un nombre premier correspond à un cycle
irréductible qui ne se factorise pas en cycles positifs plus petits. La périodicité
nonadique peut sélectionner des candidats et organiser des classes résiduelles ; la
primalité reste définie par la propriété multiplicative exacte

\[
p=ab\Longrightarrow a=1\ \text{ou}\ b=1.
\]

Une fois le coproduit des cycles arithmétiques premiers sélectionné, on
obtient \cref{eq:meissel-mertens,eq:twin-prime-constant}. La structure
hélicoïdale et la racine numérique paramètrent la phase et la retenue, mais ne remplacent pas
le critère de factorisation.

\section{Densité, terme régularisé et holonomie des lacunes}

La densité positive de base est

\begin{equation}
\varepsilon_{\rm vac}=\log_{10}\frac{10}{9}.
\label{eq:vacancy-density}
\end{equation}

La séquence se déduit de la complétion
\(1000=729+271\) détermine
\begin{equation}
\varepsilon_{\rm vac}
=1-\log_{1000}(729)
=\log_{10}\frac{10}{9},
\label{eq:vacancy-density-two-forms}
\end{equation}
et le mot mécanique exact
\begin{equation}
z_n=\lfloor(n+1)\varepsilon_{\rm vac}\rfloor
-\lfloor n\varepsilon_{\rm vac}\rfloor,
\qquad n\ge1.
\label{eq:vacancy-mechanical-word}
\end{equation}
Comme \(\varepsilon_{\rm vac}\) est irrationnel, le mot est sturmien :
il a une complexité \(p(N)=N+1\), un équilibre un et une entropie topologique nulle.
Cette dynamique diffère du décalage complet sur l'alphabet
triadique : la lacune détermine une frontière ordonnée d'entropie nulle.

La somme télescopique produit
\begin{equation}
\sum_{n=1}^{N}z_n
=\lfloor(N+1)\varepsilon_{\rm vac}\rfloor,
\qquad
\sum_{n=1}^{N}(z_n-\varepsilon_{\rm vac})
=\varepsilon_{\rm vac}
-\{(N+1)\varepsilon_{\rm vac}\}.
\label{eq:vacancy-discrepancy}
\end{equation}
Les sommes partielles de la discrépance ont un module inférieur à un. C'est
l'estimation qui soutient les constantes harmoniques ultérieures.

La fraction continue commence
\[
\varepsilon_{\rm vac}
=[0;21,1,5,1,6,2,3,\ldots].
\]
C'est pourquoi les uns de \(z_n\) sont séparés par \(21\) ou \(22\) pas.
Si
\[
\sigma=22-\varepsilon_{\rm vac}^{-1}
=1-T_G(\varepsilon_{\rm vac}),
\]
alors \(a_1(\sigma)=6\) et les retours de la séparation courte sont de
longueur \(6\) ou \(7\). Les couples \(21/22\) et \(6/7\) sont deux échelles d'une
même rotation, et non des coïncidences d'entiers.

Pour le mot exact \cref{eq:vacancy-mechanical-word}, la constante régularisée est

\begin{equation}
\gamma_{\rm vac}^{+}
=\lim_{N\to\infty}\left(
\sum_{n\le N}\frac{z_n}{n}
-\varepsilon_{\rm vac}\log N
\right).
\label{eq:vacancy-gamma}
\end{equation}

L'holonomie associée est

\begin{equation}
\rho_\varepsilon(k)=\exp(2\pi\ii k\varepsilon_{\rm vac}).
\label{eq:vacancy-holonomy}
\end{equation}

Les trois quantités ont des types distincts : densité, terme fini et phase. La tour de Stieltjes des lacunes prolonge \cref{eq:vacancy-gamma} :

\begin{equation}
\gamma_{\rm vac,j}^{+}
=\varepsilon_{\rm vac}\gamma_j
+\sum_{n\ge1}\frac{(z_n-\varepsilon_{\rm vac})(\log n)^j}{n}.
\label{eq:vacancy-stieltjes}
\end{equation}

Pour \(j=0\), on retrouve \(\gamma_{\rm vac}^{+}\). La convergence exige la borne de discrépance du développement intégral des lacunes.

La convergence est prouvée ici. Pour \(j\ge0\), la suite
\(b_n=(\log n)^j/n\) tend vers zéro et est décroissante à partir de
\(n>e^j\). La borne uniforme \cref{eq:vacancy-discrepancy} et le critère de
Dirichlet prouvent la convergence de
\[
\sum_{n\ge1}(z_n-\varepsilon_{\rm vac})b_n.
\]
Pour \(j=0\), séparer
\(\varepsilon_{\rm vac}\sum_{n\le N}n^{-1}\) donne exactement
\cref{eq:vacancy-gamma}. Une sommation d'Abel borne de manière conservative la
queue après \(N\) par \(2/(N+1)\) : un terme provient du bord et
un autre de la variation totale de \(1/n\). L'exécution exhaustive déclarée utilise
\(N=10^8\), énumère les positions de Beatty des uns, vérifie les
ensembles d'écarts \(\{21,22\}\) et de retours \(\{6,7\}\), et obtient comme
partie finie
\[
-0.11207107505876366224207105242526397929\ldots.
\]
En combinant la borne analytique avec les gardes d'arrondi et de frontière
enregistrées, le certificat fini encadre
\begin{equation}
-0.112071094143613634
<\gamma_{\rm vac}^{+}
<-0.112071055516338732.
\label{eq:vacancy-gamma-interval}
\end{equation}
Cet intervalle est une valeur terminale avec erreur contrôlée. Il ne définit pas le
mot ; le mot et sa discrépance le précèdent.

Une seconde évaluation par arithmétique multiprécision et contrôle dirigé
de l'arrondi fournit, à dix-huit décimales, l'intervalle certifié
\[
-0.112071086058763562
<\gamma_{\rm vac}^{+}
<-0.112071064058763762,
\]
strictement contenu dans \cref{eq:vacancy-gamma-interval}. L'emboîtement des
deux intervalles vérifie la stabilité de la partie finie par rapport à la
troncature, aux gardes de frontière et au régime de précision. Le
résultat dépend uniquement des récurrences, des bornes de queue et de
l'arithmétique d'intervalles exposées dans ce chapitre.

\section{Contenu opératoriel de la représentation hélicoïdale}

La représentation hélicoïdale code une phase périodique accompagnée d'un
incrément de mémoire. Elle possède un contenu opératoriel lorsqu'il existe un cocycle

\[
c(\gamma\eta)=c(\gamma)+c(\eta)
\]

qui transporte la retenue et dont la monodromie s'entrelace avec l'opérateur
arithmétique. En l'absence de cette application, l'hélice ne constitue qu'une
paramétrisation de phase ; avec l'entrelaceur, elle distingue le retour de phase
du retour d'état.

\begin{remark}[Contrôle négatif]
La cofinalité des cylindres n'implique pas la conjugaison dynamique. Un
contre-exemple à la commutation entre le décalage TPK et \(T_G\)
empêche l'entrelaceur, mais n'affecte pas
\cref{eq:levy,eq:gauss-entropy}. Dans le secteur des lacunes, la divergence
de \cref{eq:vacancy-stieltjes} sous la borne déclarée réfute la famille de
constantes régularisées.
\end{remark}

\subsection*{Portée logique et domaine de validité}
Cofinalité, identités de Gauss, Lévy, entropie et Lochs : \emph{Théorème interne exact dans le système transporté}. GKW et entrelaceur dynamique : \emph{Programme ouvert avec critère de réfutation}. Densité, constante et holonomie de lacune : \emph{Théorème interne exact avec contrôle de discrépance}.

\section{Théorème directeur du réseau radical et de la loi puissance--lacune}

L’apparition de \(369\) et de \(693\) ne se réduit pas à une ressemblance graphique. Deux mécanismes exacts convergent :

\begin{enumerate}
\item \textbf{mécanisme algébrique nonadique :} \(3,6,9\) sont les représentants visuels du radical nilpotent \(\mathfrak m=(3)\subset\mathbb Z/9\mathbb Z\) ; les mots \(369\), \(693\) et \(936\) sont les trois phases de son orbite multiplicative ;
\item \textbf{mécanisme archimédien des lacunes :} la différence entre la longueur décimale de \(9^n\) et la frontière \(10^n\) est gouvernée exactement par la même rotation irrationnelle \(\delta_{\rm vac}=\log_{10}(10/9)\) qui organise déjà le calendrier des lacunes.
\end{enumerate}

Pour \(n=9^9\), les deux mécanismes se rencontrent dans l’identité

\[
9^9-\left\lceil 9^9\log_{10}\frac{10}{9}\right\rceil
=369\,693\,099,
\]

de sorte que

\[
\#\operatorname{cifras}\bigl(9^{9^9}\bigr)=369\,693\,100.
\]

Le bloc \(369\mid693\) est donc une sortie exacte du lecteur logarithmique de la puissance nonadique ; ce n’est pas le préfixe décimal de \(9^{9^9}\). L’unité finale provient de l’opérateur universel qui transforme l’ordre décimal en longueur décimale.

Cette identité ouvre une lecture structurelle de plus grande portée : somme, produit, puissance, soustraction, division, racine et logarithme peuvent être organisés comme une hiérarchie de transports et de quotients. Chaque montée d’opération compactifie des histoires ; le TPK conserve les coordonnées qui permettent de les reconstruire.


\section{Structure radicale de l’APP sur \(\mathbb Z/9\mathbb Z\)}

\subsection{Anneau local nonadique \(\mathbb Z/9\mathbb Z\) et radical nilpotent}

Soit

\[
R=\mathbb Z/9\mathbb Z.
\]

Son groupe des unités et son idéal maximal sont

\[
R^\times=\{1,2,4,5,7,8\},
\qquad
\mathfrak m=(3)=\{0,3,6\}=\{9,3,6\}_{\rm vis}.
\]

De plus,

\[
\mathfrak m^2=0.
\]

L’égalité précédente est une égalité d’idéaux : tout produit de deux éléments de \(\mathfrak m\) s’annule modulo neuf. Elle ne doit pas être confondue avec le produit cartésien \(\mathfrak m\times\mathfrak m\), qui contient neuf positions de l’APP.

\subsection{Suite exacte du multiplicateur \(3\)}

La multiplication par trois,

\[
M_3:R\longrightarrow R,
\qquad x\longmapsto3x,
\]

a

\[
\ker M_3=\mathfrak m,
\qquad
\operatorname{im}M_3=\mathfrak m.
\]

Elle induit donc un isomorphisme d’espaces sur \(\mathbb F_3\),

\[
\overline M_3:R/\mathfrak m\xrightarrow{\sim}\mathfrak m,
\qquad
\overline x\longmapsto3x.
\]

En représentants visuels,

\[
\overline0\longmapsto9,
\qquad
\overline1\longmapsto3,
\qquad
\overline2\longmapsto6.
\]

Ainsi, la bande \(3,6,9\) n’est pas un ornement décimal. C’est une copie interne du quotient ternaire logée dans le radical nonadique.

\subsection{Orbite cyclique \(369\to693\to936\to369\)}

La ligne de la table multiplicative correspondant à \(3\) est

\[
3,6,9,3,6,9,3,6,9.
\]

Le changement de phase \(j\mapsto j+1\) produit

\[
369\longmapsto693\longmapsto936\longmapsto369.
\]

Ces trois mots sont les trois lectures cycliques de l’application

\[
\mathbb F_3\longrightarrow\mathfrak m,
\qquad
\bar j\longmapsto3j.
\]

La ligne \(6\), en satisfaisant \(6\equiv-3\pmod9\), inverse l’orientation du même cycle. La marque \(9\) remplit deux fonctions typées :

\[
9\bmod9=0,
\qquad
\operatorname{dr}_9(9)=9.
\]

Le module publie la classe nulle ; la racine numérique positive conserve la marque de clôture. C’est précisément cette différence de lecteurs qui permet à \(9\) d’être à la fois zéro résiduel et clôture visible.

\subsection{Croix radicale de l’APP et réduction coordonnée à \(\mathbb F_3^2\)}

Sur

\[
X_9=R^2,
\]

l’union

\[
\mathscr C_{369}
=(\mathfrak m\times R)\cup(R\times\mathfrak m)
\]

est la croix formée par les lignes et les colonnes \(3,6,9\). Elle contient

\[
3\cdot9+3\cdot9-3\cdot3=45
\]

positions. Son intersection centrale \(\mathfrak m\times\mathfrak m\) contient neuf positions. Dans l’observable multiplicative, la somme de la croix est

\[
297=216+81=3(72+27).
\]

La réduction coordonnée produit

\[
X_9/(\mathfrak m\times\mathfrak m)\cong\mathbb F_3^2
\]

comme ensemble quotient typé par les classes modulo trois. Le relèvement inverse exige de conserver la coordonnée à l’intérieur de chaque fibre de trois éléments.

\subsection{Noyau \(7227\), décomposition spectrale et lien avec la croix radicale}

Le bloc central de l’APP est

\[
B_c=
\begin{pmatrix}
7&2\\
2&7
\end{pmatrix}
=7I+2R
=9P_++5P_-.
\]

Sa lecture par lignes produit \(72\mid27\) ; sa lecture par diagonales produit \(77\mid22\). On vérifie

\[
72+27=77+22=99,
\]

\[
72-27=45=\det B_c,
\]

\[
77-22=55=54+1.
\]

La relation avec la croix radicale n’est pas nominale :

\[
3\cdot72=216,
\qquad
3\cdot27=81,
\qquad
3\cdot99=297.
\]

Le bloc unitaire central et le réseau non unitaire \(3,6,9\) sont ainsi reliés par une identité locale–globale exacte. La hiérarchie de la différence somme–produit conserve ses types :

\[
4\quad\text{(saut spectral local)},\qquad
2\quad\text{(couplage hors diagonale)},
\]

\[
6\quad\text{(différence comprimée modulo trois)},\qquad
54\quad\text{(déploiement bidimensionnel)}.
\]


\section{Loi puissance--lacune pour les échelles \texorpdfstring{\(9^n\)}{9 à la puissance n} et \texorpdfstring{\(10^n\)}{10 à la puissance n}}

\subsection{Fonctions de longueur décimale et de lacune cumulée}

Pour \(n\ge1\), soit

\[
D_9(n)=\#\operatorname{cifras}(9^n)
=\left\lfloor n\log_{10}9\right\rfloor+1.
\]

La frontière décimale \(10^n\) possède \(n+1\) chiffres. Nous définissons la lacune cumulée de la puissance nonadique par

\[
V_9(n)=(n+1)-D_9(n).
\]

Nous introduisons les deux nombres complémentaires

\[
\rho_{\rm vac}=\log_{10}9,
\qquad
\delta_{\rm vac}=1-\rho_{\rm vac}
=\log_{10}\frac{10}{9}.
\]

\subsection{Identité exacte entre longueur décimale et rotation des lacunes}

\textbf{Théorème.} Pour tout entier \(n\ge1\),

\[
\boxed{V_9(n)=\left\lceil n\delta_{\rm vac}\right\rceil},
\]

et, par conséquent,

\[
\boxed{D_9(n)=n-\left\lceil n\delta_{\rm vac}\right\rceil+1}.
\]

\textbf{Démonstration.} Comme \(\delta_{\rm vac}\) est irrationnel, \(n\delta_{\rm vac}\notin\mathbb Z\). Alors

\[
\begin{aligned}
D_9(n)
&=\lfloor n(1-\delta_{\rm vac})\rfloor+1\\
&=\lfloor n-n\delta_{\rm vac}\rfloor+1\\
&=n-\lceil n\delta_{\rm vac}\rceil+1.
\end{aligned}
\]

En soustrayant de \(n+1\), on obtient la première identité. \(\square\)

\subsection{Encodage mécanique des lacunes par les puissances nonadiques}

Les accroissements de la lacune sont

\[
Z_n=V_9(n+1)-V_9(n)
=\lceil(n+1)\delta_{\rm vac}\rceil
-\lceil n\delta_{\rm vac}\rceil
\in\{0,1\}.
\]

C’est exactement le mot mécanique de pente \(\delta_{\rm vac}\) utilisé par le calendrier HMT des lacunes. Ce calendrier admet donc une lecture supplémentaire entièrement exacte :

\begin{quote}
\(Z_n=1\) si et seulement si, en passant de \(9^n\) à \(9^{n+1}\), on perd un nouveau rang décimal par rapport à la frontière \(10^n\to10^{n+1}\).
\end{quote}

La puissance n’ajoute pas une coïncidence à la théorie des lacunes : elle fournit sa réalisation naturelle comme compteur de perte de rang décimal.

\subsection{Loi \(21/22\) comme cas du théorème puissance--lacune}

Puisque

\[
\delta_{\rm vac}^{-1}
=21.85434532678283256\ldots,
\]

les indices pour lesquels \(Z_n=1\) sont séparés par \(21\) ou \(22\). Les vingt et un premiers exposants satisfont

\[
D_9(n)=n,\qquad1\le n\le21,
\]

tandis que

\[
D_9(22)=21.
\]

En particulier,

\[
9^9=387\,420\,489
\]

possède exactement neuf chiffres. Cette conservation initiale du rang n’est pas une règle indéfinie : le calendrier des lacunes détermine exactement quand elle cesse d’être vérifiée.

\subsection{Évaluation exacte en \(n=9^9\) et longueur de \(9^{9^9}\)}

Soit

\[
n_*=9^9=387\,420\,489.
\]

Alors

\[
n_*\delta_{\rm vac}
=17\,727\,389.3684296412564569\ldots,
\]

et, par conséquent,

\[
\boxed{V_9(n_*)=17\,727\,390}.
\]

La coordonnée conservée après déduction des lacunes est

\[
n_*-V_9(n_*)
=387\,420\,489-17\,727\,390
=369\,693\,099.
\]

Par conséquent,

\[
\boxed{\operatorname{ord}_{10}\bigl(9^{9^9}\bigr)=369\,693\,099},
\]

\[
\boxed{D_9(9^9)=369\,693\,100}.
\]

On obtient en outre les congruences

\[
9^9\equiv0\pmod9,\qquad
V_9(9^9)\equiv0\pmod9,
\]

\[
369\,693\,099\equiv0\pmod9,\qquad
369\,693\,100\equiv1\pmod9.
\]

Les sommes de chiffres correspondantes sont

\[
s_{10}(9^9)=45,\qquad
s_{10}(V_9(9^9))=36,
\]

\[
s_{10}(369\,693\,099)=54,\qquad
s_{10}(369\,693\,100)=37.
\]

Le retour \(9\to1\) apparaît ici avec des types précis :

\begin{enumerate}
\item \(9^9\) et la lacune cumulée appartiennent à la classe nulle modulo neuf ;
\item leur différence, l’ordre décimal, reste dans la classe nulle et possède une somme de chiffres \(54\) ;
\item l’opérateur universel \(D=\operatorname{ord}+1\) ajoute l’unité de clôture et publie la classe positive \(1\).
\end{enumerate}

\subsection{Domaine de validité de la loi puissance--lacune}

Il est démontré :

\begin{itemize}
\item le théorème puissance–lacune pour tout \(n\) ;
\item l’identité avec le mot mécanique des lacunes ;
\item l’évaluation exacte en \(n=9^9\) ;
\item l’ordre et la longueur décimale de \(9^{9^9}\) ;
\item les congruences et sommes de chiffres précédentes.
\end{itemize}

On n’affirme pas que \(9\) soit l’unique base ayant une propriété analogue dans tout domaine, que chaque apparition externe de \(369\) provienne de cette tour, ni que la tour remplace la construction des neuf phases ou le certificat de génération de \(\pi,e,\varphi\).


\section{Hiérarchie opératorielle, réversibilité et mémoire généalogique}

\subsection{Translation, dilatation et exponentiation comme niveaux opératoriels}

Les opérations élémentaires peuvent être vues comme des actions :

\[
T_a(x)=x+a,\qquad
M_a(x)=ax,\qquad
P_k(x)=x^k.
\]

Leurs inverses partielles sont, respectivement, la soustraction, la division et l’extraction de racines. Le passage d’une opération à la suivante modifie la structure des fibres.

\subsection{Réversibilité stratifiée dans \(\mathbb Z/9\mathbb Z\)}

\textbf{Proposition.} Dans l’anneau nonadique :

\begin{enumerate}
\item toute translation \(T_a\) est bijective ;
\item \(M_a\) est bijective si et seulement si \(a\in R^\times\) ;
\item si \(a\in\mathfrak m\), la multiplication possède des fibres non triviales et son image est contenue dans \(\mathfrak m\) ;
\item pour \(x\in\mathfrak m\), \(x^k=0\) pour tout \(k\ge2\) ;
\item sur \(R^\times\cong C_6\), l’application \(P_k\) est un automorphisme si et seulement si \(\gcd(k,6)=1\).
\end{enumerate}

\textbf{Démonstration.} La première affirmation utilise la structure de groupe additif. La deuxième est la caractérisation des unités. La troisième découle de \(aR\subseteq\mathfrak m\). La quatrième procède de \(\mathfrak m^2=0\). La cinquième est la condition pour que la multiplication par \(k\) soit un automorphisme du groupe cyclique d’ordre six. \(\square\)

Le résultat produit une hiérarchie exacte :

\[
\text{somme : réversibilité globale},
\]

\[
\text{produit : réversibilité restreinte aux unités},
\]

\[
\text{puissance : dynamique cyclique sur les unités et effondrement nilpotent dans }\mathfrak m.
\]

Les racines ne sont des fonctions que lorsque la puissance correspondante est injective sur le domaine choisi ; sur le domaine total, ce sont des correspondances à plusieurs branches ou des relations vides.

\subsection{Lecture du Topological Phase Kernel sur la hiérarchie opératorielle}

Cette stratification suggère une définition fonctionnelle :

\begin{quote}
Le TPK n’ajoute pas une opération extérieure à la somme, au produit et à la puissance. Il conserve le feuillet, le quotient, la retenue, l’orientation, la route et la frontière que chaque évaluation ordinaire élimine lorsqu’elle identifie plusieurs histoires à une même sortie.
\end{quote}

La somme et le produit sont les deux premières évaluations sur l’APP. La puissance est une composition répétée du produit ; la racine sélectionne ou reconstruit des branches de cette composition. Le registre distingue la valeur publiée de l’histoire opératorielle qui la produit.

\subsection{Descente logarithmique entre niveaux consécutifs de la hiérarchie}

Dans un domaine positif,

\[
\log(xy)=\log x+\log y,\qquad
\log(x^k)=k\log x,\qquad
\log(a^x)=x\log a.
\]

Le logarithme transforme le produit en somme, la puissance en produit et l’exponentielle en dilatation. Le défaut apparaît lorsque le lecteur ultérieur discrétise :

\[
\mathcal L_{10}(a^n)
=\lfloor n\log_{10}a\rfloor+1.
\]

La partie entière introduit le calendrier des franchissements de frontière. Pour \(a=9\), ce calendrier est précisément le mot des lacunes.

\subsection{Conjugaison affine des translations et des dilatations}

Les translations et les dilatations satisfont

\[
M_aT_bM_a^{-1}=T_{ab}
\]

lorsque \(a\) est inversible. Cette identité exprime que le produit réorganise les unités de somme. Lorsque \(a\) cesse d’être inversible, la conjugaison se rompt parce que des branches disparaissent. Le radical \(3,6,9\) localise exactement cet échec.


\section{Algèbre des chemins : composition multiplicative, action additive et superposition}

Dans une algèbre des chemins, la somme agrège les alternatives et le produit concatène les trajectoires. Pour des poids locaux \(w(e)\),

\[
K(x,y)
=\sum_{\gamma:x\to y}
\prod_{e\in\gamma}w(e).
\]

Si l’on fixe une branche logarithmique cohérente,

\[
\prod_{e\in\gamma}w(e)
=\exp\!\left(\sum_{e\in\gamma}\log w(e)\right),
\]

et, avec l’action discrète appropriée,

\[
K(x,y)=\sum_{\gamma:x\to y}e^{-S(\gamma)/\hbar}.
\]

L’équation réunit les opérations :

\begin{itemize}
\item le produit compose les pas successifs ;
\item le logarithme accumule l’action locale ;
\item l’exponentielle reconstruit le poids de la route ;
\item la somme agrège toutes les alternatives.
\end{itemize}

L’APP est le support minimal sur lequel somme et produit agissent sur les mêmes histoires. Le TPK ajoute les données qui empêchent la somme sur les chemins de devenir une somme sur des valeurs sans généalogie.

Dans ce cadre, un chiffre est la sortie d’un lecteur ; une valeur est la projection archimédienne d’une section ; un nombre HMT conserve la famille compatible de chemins qui soutient cette valeur ; une particule conserve cette généalogie sous une réalisation spectrale persistante.


\section{Racine interne de \(-1\), structure exponentielle et géométries elliptique, parabolique et hyperbolique}

\subsection{Racine interne de \(-1\) dans la branche finie}

La matrice d’incidence de Paley fournit un opérateur \(A_W\) avec

\[
A_W^2=-I.
\]

Cela fait de \(\mathbb F_3^6\) un espace tridimensionnel sur

\[
\mathbb F_9=\mathbb F_3[\iota]/(\iota^2+1).
\]

La racine de \(-1\) n’est pas adjointe de l’extérieur comme une notation complexe. Elle émerge comme endomorphisme d’un support ternaire déjà construit.

\subsection{Réalisation continue du type quadratique induit}

Dans la branche réelle elliptique, le générateur \(J_{\rm e}\) satisfait

\[
J_{\rm e}^2=-I,
\]

et

\[
e^{tJ_{\rm e}}=\cos t\,I+\sin t\,J_{\rm e}.
\]

En particulier,

\[
e^{\frac\pi2J_{\rm e}}=J_{\rm e},\qquad
e^{\pi J_{\rm e}}=-I,\qquad
e^{2\pi J_{\rm e}}=I.
\]

Ainsi, \(J_{\rm e}\) est le quart de tour interne et l’équation d’Euler est sa clôture au demi-tour.

\subsection{Support opératoriel commun de \(3\) lectures géométriques}

Soit

\[
\mathbb D=\{z\in\mathbb C:|z|<1\}.
\]

Ce même ensemble contient le diamètre réel \((-1,1)\), l’intérieur complexe, la frontière trigonométrique \(S^1\) et la métrique hyperbolique de Poincaré

\[
ds_{\rm P}^2=\frac{4\ell^2|dz|^2}{(1-|z|^2)^2}.
\]

L'application

\[
R_{\pi/2}(z)=iz
\]

préserve \(|z|\) et \(|dz|\). Par conséquent, c’est simultanément :

\begin{itemize}
\item une multiplication complexe par \(\sqrt{-1}\) ;
\item une rotation trigonométrique de \(90^\circ\) ;
\item une isométrie euclidienne du disque ;
\item une isométrie hyperbolique de Poincaré ;
\item un automorphisme de sa frontière circulaire.
\end{itemize}

La coïncidence est littérale : la même application agit sur le même porteur ; c’est le lecteur métrique qui change.

\subsection{Géométries elliptique, parabolique et hyperbolique de l’équation d’Euler étendue}

Les générateurs

\[
J_{+1}^2=-I,\qquad
J_0^2=0,\qquad
J_{-1}^2=I
\]

produisent

\[
e^{tJ_{+1}}=\cos t\,I+\sin t\,J_{+1},
\]

\[
e^{tJ_0}=I+tJ_0,
\]

\[
e^{tJ_{-1}}=\cosh t\,I+\sinh t\,J_{-1}.
\]

La même série exponentielle se sépare en trois régimes selon le carré du générateur : elliptique, parabolique et hyperbolique. La spécialisation

\[
e^{2\log\varphi\,J_{-1}}=F_{\rm av}^{\,2}
\]

intègre l’auto-échelle de Fibonacci. Euler et Fibonacci apparaissent comme des spécialisations d’une famille exponentielle quadratique commune.

\subsection{Porteur commun, opérateur commun et famille exponentielle quadratique}

La fusion géométrique se formule à trois niveaux :

\begin{enumerate}
\item \textbf{porteur commun :} le disque unité contient le diamètre réel, l’intérieur complexe et le bord circulaire ;
\item \textbf{opérateur commun :} \(z\mapsto iz\) est un quart de tour, une multiplication complexe et une isométrie hyperbolique ;
\item \textbf{famille commune :} l’équation exponentielle tritique réunit les branches elliptique, parabolique et hyperbolique.
\end{enumerate}

Le résultat n’identifie pas les métriques euclidienne et hyperbolique. Il affirme que toutes deux sont des réalisations distinctes sur un même support et partagent l’automorphisme issu de la racine interne de \(-I\).


\section{Restriction de la rotation intérieure, holonomie de bord et torsion minimale}

La construction contient deux objets positifs et différenciés :

\begin{enumerate}
\item le générateur interne \(J_{\rm e}\), avec \(J_{\rm e}^2=-I\) ;
\item l’invariant global
\end{enumerate}
   \[
   \Delta_4=
   \frac{\pi-e\log\pi}{270}
   \left(1+\frac{\pi}{729}\right),
   \]
   appelé torsion minimale et utilisé avant la sélection d’espèce.

La thèse de l’auteur ajoute que la dynamique rotationnelle interne, lorsqu’elle est publiée au bord, est absorbée comme torsion minimale. La relation ne doit pas s’écrire comme une simple égalité \(J_{\rm e}=\Delta_4\) : les objets ont des types différents. Sa forme naturelle exige

\[
\text{rotation intérieure}
\xrightarrow{\operatorname{Res}_{\partial}}
\text{holonomie de bord}
\xrightarrow{\operatorname{Def}}
\text{défaut scalaire}.
\]

Soit \(q_\partial\) le lecteur de frontière, \(\nabla_\partial\) la connexion induite et \(e_\partial\) le repère discret. La torsion de bord a pour type

\[
T_\partial=d_{\nabla_\partial}e_\partial.
\]

L’affirmation de l’auteur acquiert une forme probatoire au moyen d’une fonctionnelle

\[
\Lambda_\partial:
\operatorname{Tors}(\partial X)
\longrightarrow\mathbb R
\]

tel que

\[
\Lambda_\partial(T_\partial)=\Delta_4
\]

et d’un théorème qui fasse procéder \(T_\partial\) de la restriction de \(J_{\rm e}\) par les opérateurs de bord du TPK.

La racine interne de \(-I\), le quart de tour, l’équation d’Euler étendue et la formule de \(\Delta_4\) sont démontrés dans leurs domaines. Le passage de la rotation intérieure à la torsion minimale de bord exige le morphisme \(\operatorname{Res}_{\partial}\) et la fonctionnelle \(\Lambda_{\partial}\) déclarés ; il n’est pas identifié par une égalité scalaire.


\section{Triplicité générationnelle, conjugaison et mélange sur l’état enrichi}

La phrase d’Isidor Isaac Rabi peut servir de référence historique, mais elle ne doit pas régir le titre ni constituer une section autonome. La question scientifique est :

\begin{quote}
Quel mécanisme interne distingue trois ordres stables de réalisation fermionique et conserve cette distinction sous couleur, conjugaison et mélange ?
\end{quote}

La réponse HMT n’est pas « parce que TRIT a trois valeurs ». La structure sépare

\[
\operatorname{flav}(X)
=\bigl(\sigma_{ud}(X),g(X),\chi_{\rm fina}(X)\bigr),
\]

où \(\sigma_{ud}\) distingue les branches, \(g\) distingue la génération physique et \(\chi_{\rm fina}\) conserve l’orientation et la mémoire au sein d’une branche. Dans les quarks, l’orbite de couleur est ajoutée après la saveur et avant le caractère.

La triplicité est obtenue comme conséquence de la structure enrichie ; ses effets sur la couleur, la conjugaison et le mélange sont développés ci-dessous :

\begin{enumerate}
\item couleur, saveur, génération et masse sont des lecteurs distincts du même état enrichi ;
\item la conjugaison particule–antiparticule agit sur l’orientation sans effacer la classe générationnelle ;
\item CKM et PMNS transportent des bases d’états déjà typés ; ils ne créent pas la triplicité ;
\item la classification neutrinique doit découler de la représentation, de l’occupation et du mélange, non d’un changement de nom des profondeurs \(G0,\ldots,G4\).
\end{enumerate}

Le fait que les neutrinos soient des leptons de spin demi-entier et d’obéissance fermionique est une classification physique connue. L’apport structurel HMT consiste à obtenir intérieurement pourquoi le secteur neutre occupe cette représentation, pourquoi il admet trois réalisations physiques et comment il se couple à l’opérateur de mélange sans lire ses angles externes.


\section{Mémoire opératorielle du TPK : conservation du résidu, du quotient, de la retenue, du feuillet, de l’orientation, de la route et de la frontière sous projection}

\subsection{Perte de réversibilité et relèvement généalogique}

Dans l’APP, différentes histoires peuvent produire le même résidu. Le produit replie davantage d’histoires que la somme ; la puissance en replie davantage que le produit ; la racine tente de reconstruire des branches déjà identifiées. Le TPK conserve la donnée qui permet d’inverser cette perte :

\[
\text{résidu},\quad
\text{quotient},\quad
\text{retenue},\quad
\text{feuillet},\quad
\text{orientation},\quad
\text{route},\quad
\text{frontière},\quad
\text{mémoire}.
\]

L’opération algébrique peut être caractérisée non seulement par sa valeur, mais par le type de généalogie qu’elle conserve ou détruit.

\subsection{Projection des histoires compatibles sur le continu archimédien}

La droite réelle apparaît lorsqu’une famille de cylindres compatibles se contracte sous l’évaluation archimédienne. La valeur réelle est l’ombre de la section ; non sa généalogie complète. La projection exceptionnelle conserve une autre partie comme incidence et symétrie.

La somme sur les chemins exprime le même principe :

\[
\text{alternatives}\xrightarrow{\sum}\text{amplitude globale},
\qquad
\text{étapes successives}\xrightarrow{\prod}\text{poids de route}.
\]

Le logarithme transforme la composition multiplicative en action additive ; l’exponentielle la ramène à une amplitude. La structure discrète du continu peut être formulée comme conservation généalogique de toutes les routes avant leur projection sur \(\mathbb R\).

\subsection{Distinction entre section numérique et réalisation persistante de particule}

Un nombre et une particule partagent des invariants ontologiques minimaux : mémoire, orientation, persistance, clôture, signature et position dans une famille de prolongements.

Le nombre publie une coordonnée archimédienne. La particule ajoute une réalisation spectrale, une occupation, une propagation et une stabilité. Le passage HMT \(\to\) MD change le lecteur appliqué à une même généalogie enrichie ; il ne change pas arbitrairement d’objet.


\section{Clôture conjointe du réseau radical, de la loi puissance--lacune, de la géométrie d’Euler et de la hiérarchie opératorielle}

\subsection{Théorème du réseau radical nonadique}

Le réseau \(3,6,9\) de l’APP est la représentation visuelle du radical nilpotent \((3)\subset\mathbb Z/9\mathbb Z\). Ses phases \(369\), \(693\) et \(936\) réalisent l’isomorphisme \(R/(3)\cong(3)\) induit par la multiplication par trois. Les bandes localisent exactement le défaut d’inversibilité du produit.

\subsection{Théorème puissance--lacune}

Pour tout \(n\ge1\), le nombre de rangs décimaux perdus par \(9^n\) par rapport à \(10^n\) est \(\lceil n\log_{10}(10/9)\rceil\). Ses accroissements forment le mot mécanique des lacunes. En \(n=9^9\), la coordonnée retenue est \(369\,693\,099\) et la longueur résultante est \(369\,693\,100\).

\subsection{Théorème géométrique de la famille exponentielle d’Euler}

La racine interne de \(-I\) fournit un quart de tour. Dans le disque unité, la même application \(z\mapsto iz\) est une multiplication complexe, une rotation trigonométrique et une isométrie de Poincaré. L’équation exponentielle tritique prolonge ce régime aux régimes parabolique et hyperbolique, en incorporant dans le dernier l’auto-échelle de Fibonacci.

\subsection{Principe de conservation opératorielle entre niveaux}

La hiérarchie somme–produit–puissance est une hiérarchie de fibres et de perte de réversibilité. Le TPK est le relèvement qui conserve les coordonnées nécessaires pour distinguer les histoires identifiées par les lecteurs ordinaires.


\section{Hypothèses de typage et contrôles négatifs}

\begin{enumerate}
\item Ne pas identifier \(369\,693\,100\) à un mot TPK sans déclarer le lecteur entre les domaines.
\item Ne pas confondre l’idéal \(\mathfrak m^2=0\) avec \(\mathfrak m\times\mathfrak m\).
\item Ne pas confondre modulo neuf et racine numérique : \(9\mapsto0\) dans le premier et \(9\mapsto9\) dans la seconde.
\item Ne pas utiliser la triplicité de TRIT comme démonstration suffisante de trois générations physiques.
\item Ne pas transformer \(J^2=-I\) en \(\Delta_4\) sans l’application de bord et la fonctionnelle de torsion.
\item Ne pas présenter \(369\) comme un préfixe de \(9^{9^9}\) : il appartient à son ordre et à sa longueur décimale.
\item Ne pas réduire les dix mille chiffres à une preuve de grandeur : ils certifient des troncatures étendues d’un processus paramétrable dont la thèse est le prolongement indéfini.
\end{enumerate}

\section{Conséquences conjointes pour le réseau radical, la mémoire opératorielle et la torsion de bord}

La découverte principale n’est pas qu’une tour contienne visuellement les chiffres \(3,6,9\). La puissance nonadique et le calendrier des lacunes sont deux présentations du même défaut logarithmique, tandis que l’APP localise \(3,6,9\) comme le radical où la multiplication perd son inversibilité. La tour relie les deux extrémités :

\[
\text{radical nonadique}
\longrightarrow
\text{puissance}
\longrightarrow
\text{logarithme}
\longrightarrow
\text{vacance}
\longrightarrow
\text{ordre décimal}
\longrightarrow
\text{clôture }9\to1.
\]

L’équation d’Euler étendue ajoute la deuxième direction : la racine interne de \(-I\) transforme la structure discrète en dynamique de rotation et réunit, sur un support commun, la droite réelle, le cercle trigonométrique et le disque hyperbolique. Le passage ultérieur à la torsion minimale doit conserver l’application de bord déjà développée dans le corpus, non la remplacer par une identification nominale.

%
\endgroup

\HMTFlushCurrentChapter
\chapter{Constantes thermiques et rayonnement : Boltzmann, Wien, Stefan--Boltzmann et thermodynamique spectrale du gaz photonique}%
\begingroup
\renewcommand{\chapter}[2][]{}%
\chapter{Constantes thermiques et rayonnement : Boltzmann, Wien, Stefan--Boltzmann et thermodynamique spectrale du gaz photonique}
\label{chap:p3-final:48:constantes_termicas}
\section{Relation thermique intrinsèque de l'échelle \(108\)}

La structure chronologique \(108\) fixe

\begin{equation}
k_B\Theta_{\rm clk}\log3
=\frac{h_{\rm int}}{108t_0}
=\frac{2\pi\hbar_{\rm int}}{108t_0}.
\label{eq:boltzmann-clock}
\end{equation}

Cette égalité ne détermine pas séparément une coordonnée numérique de \(k_B\) et une autre de \(\Theta_{\rm clk}\) sans normalisation supplémentaire. Elle détermine bien leur produit comme section d'énergie. Le facteur \(\log3\) a trois antécédents compatibles :

\begin{equation}
\log3=C_1+C_2-2C_0
=\frac{S_{\TRIT}}{k_B}
=\frac{E_P}{k_B\Theta_{\rm clk}}.
\label{eq:log3-triple}
\end{equation}

Les trois égalités correspondent aux réalisations résiduelle, informationnelle et
thermodynamique d'un même scalaire. La coïncidence de valeur n'identifie pas
leurs domaines respectifs.

\section{Formulations spectrales de la loi de déplacement de Wien}

Pour la densité spectrale par longueur d'onde, le maximum satisfait

\[
5(1-\ee^{-x_5})=x_5,
\qquad
x_5=5+W_0(-5\ee^{-5}).
\]

Avec \cref{eq:boltzmann-clock},

\begin{equation}
\lambda_{\max}(\Theta_{\rm clk})
=\frac{108\log3}{x_5}\,\ell_0.
\label{eq:wien-lambda}
\end{equation}

Pour la densité par fréquence,

\[
3(1-\ee^{-x_3})=x_3,
\qquad
x_3=3+W_0(-3\ee^{-3}),
\]

et

\begin{equation}
\nu_{\max}(\Theta_{\rm clk})
=\frac{x_3}{108t_0\log3}.
\label{eq:wien-frequency}
\end{equation}

Les maxima ne sont pas réciproques, puisque les densités spectrales se
transforment avec le jacobien correspondant. Le quotient \(x_5/x_3\)
appartient à la structure de la paramétrisation et ne représente pas une
incohérence.

\section{Loi de Stefan--Boltzmann et thermodynamique du gaz de photons}

Évaluer la loi de rayonnement en \(\Theta_{\rm clk}\) produit le flux intrinsèque

\begin{equation}
j_\star(\Theta_{\rm clk})
=\frac{2\pi^5}{15\,108^4(\log3)^4}\,
\frac{h_{\rm int}}{\ell_0^2t_0^2},
\label{eq:stefan-hmt}
\end{equation}

dans la normalisation déclarée. La constante de Stefan--Boltzmann n'est pas insérée comme primitive : elle résulte de l'intégration du spectre bosonique une fois l'action, la propagation et la température fixées.
Les lois spectrales employées sont celles de Planck et de Wien \cite{Planck,Wien} ; HMT fixe ici la section sur laquelle elles sont évaluées, et non leurs constantes par comparaison rétrospective.

En unités de l'échelle de Planck,

\begin{align}
n_\gamma\ell_P^3
&=\frac{2\zeta(3)}{\pi^2\log^3 3},
\label{eq:photon-number}\\
\frac{u_\gamma\ell_P^3}{E_P}
&=\frac{\pi^2}{15\log^4 3},
\label{eq:photon-energy}\\
\frac{s_\gamma\ell_P^3}{k_B}
&=\frac{4\pi^2}{45\log^3 3}.
\label{eq:photon-entropy}
\end{align}

\(\zeta(3)\) représente le moment bosonique d'ordre trois, tandis que
\(\log3\) détermine l'énergie de décision TRIT. L'équation conserve la
distinction entre les deux invariants.

La conductance thermique quantique s'exprime comme

\begin{equation}
\frac{\kappa_Qt_0}{k_B}=\frac{\pi^2}{324\log3}.
\label{eq:thermal-conductance}
\end{equation}

\section{Confluence métrologique des invariants thermiques et radiatifs}

\begin{longtable}{@{}p{0.18\textwidth}p{0.27\textwidth}p{0.45\textwidth}@{}}
\toprule
Objet & Équation interne & Significado estructural\\
\midrule
\endhead
\(R_K\) & \(Z_{\rm vac}/(2\alpha)\) & résistance quantique invariante bidirectionnelle\\
\(G_0\) & \(4\alpha/Z_{\rm vac}\) & conductancia complementaria\\
\(\Phi_0\) & \(\sqrt{hZ/(8\alpha)}\) & flux qui conserve l'orientation de l'action\\
\(K_J\) & \(\Phi_0^{-1}\) & relation réciproque entre fréquence et flux\\
\(E_P\) & \(k_B\Theta\log3\) & énergie du cycle et contenu informationnel TRIT\\
Wien & \(108\log3/x_5\) & extremum spectral à l'échelle chronologique\\
Apéry photonique & \(2\zeta(3)/(\pi^2\log^3 3)\) & densité de nombre bosonique\\
\bottomrule
\end{longtable}

\begin{remark}[Contrôle négatif]
Le bloc est réfuté si un changement de variable spectrale conserve
indûment le même extremum de Wien, si \(\zeta(3)\) remplace
\(\log3\), ou si des valeurs séparées sont attribuées à \(k_B\) et \(\Theta\) sans
spécifier de carte. Les identités
\cref{eq:quantum-electrical-identities,eq:planck-trit} fournissent des contrôles
algébriques directs.
\end{remark}

\subsection*{Portée logique et domaine de validité}
Quanta électriques et confluence Planck--TRIT : \emph{Théorème interne exact}. Wien, Stefan--Boltzmann et gaz photonique : \emph{Composition mathématique exacte} avec le continu bosonique déclaré. La carte SI reste une \emph{Comparaison métrologique externe}.


%
\endgroup

\HMTFlushCurrentChapter
\chapter{Constantes électrofaibles et mélange : déterminant angulaire, relation \(W/Z\), angle de Weinberg et réflexion rationnelle CKM}%
\begingroup
\renewcommand{\chapter}[2][]{}%
\chapter{Constantes électrofaibles et mélange : déterminant angulaire, relation W/Z, angle de Weinberg et réflexion rationnelle CKM}
\label{chap:p3-final:49:sector_electrodebil}
\section{Déterminant angulaire et relation électrofaible de Weinberg}

Le déterminant de \(T\) est

\begin{equation}
\det T=q_+q_-=D_A.
\label{eq:detT}
\end{equation}

Le secteur constitutif du vide évalue \(F(T)^2\), où
\(F(z)=(1-z^{90})/(1-z^{120})\), tandis que le secteur angulaire
électrofaible évalue \(\det T\). Il s'agit de deux fonctions d'un
antécédent commun, et non d'une factorisation horizontale entre constantes.
Cette distinction est développée dans le \cref{ch:ew}.

\begin{remark}[Contrôle négatif]
Le bloc est réfuté si l'une quelconque de ces identités exactes est violée :
positivité de \(I-T^{120}\), égalité
\cref{eq:three-four-completion}, relations \cref{eq:vacuum-relations} ou
compatibilité \(E_m(\mathcal V_{m+1})=\mathcal V_m\). Une coïncidence SI
ne peut pas compenser une défaillance opératorielle.
\end{remark}

\subsection*{Portée logique et domaine de validité}
Opérateur, spectre, complétion \(3\to4\), parité, amplification et sections constitutives : \emph{Théorème interne exact}. La formulation qui les réunit en un seul opérateur est une \emph{Formalisation nouvelle} ; les données angulaires et les droites sont un \emph{Résultat antécédent démontré} du développement HMT précédent et sont déclarées ici comme structure de départ explicite.


\label{ch:ew}

\section{Fonctions constitutive et déterminantale de l'opérateur angulaire}

Rappelons l'opérateur à deux feuillets du \cref{ch:vacuum},
\begin{equation}
\label{eq:ew-angular-operator}
T=q_+P_+ +q_-P_-,
\qquad
q_-=\ee^{-(A_{\rm rad}-C^*_{\rm rad})},
\qquad
q_+=\ee^{-(A_{\rm rad}+C^*_{\rm rad})}.
\end{equation}
Le vide électromagnétique applique la fonction
\begin{equation}
\label{eq:ew-vacuum-reader}
F(T)^2,
\qquad
F(z)=\frac{1-z^{90}}{1-z^{120}},
\end{equation}
tandis que le secteur angulaire électrofaible applique le déterminant
\begin{equation}
\label{eq:ew-determinant-reader}
\boxed{
\det T=q_+q_-=\ee^{-2A_{\rm rad}}=:D_A.}
\end{equation}
Ce sont deux réalisations fonctionnelles de l'antécédent commun \((A,C^*,T)\). Il n'existe pas de
flèche causale \(F(T)^2\to\det T\) ni son inverse. L'unité structurelle
réside dans \(T\), et non dans l'identification ultérieure de la perméabilité, de la permittivité et de
\(\theta_W\).

Dans l'involution des feuillets \(C^*\mapsto-C^*\), les valeurs propres
\(q_+,q_-\) s'échangent. C'est pourquoi \(D_A\) et la vitesse constitutive
sont paires, tandis que \((\widehat\mu,\widehat\varepsilon)\) s'échange et
\(\widehat Z\) s'inverse. La relation électrofaible qui suit utilise
l'invariant pair.

\section{Relation exacte entre les secteurs \(W\) et \(Z\)}

La coordonnée angulaire de durée--perduration est
\begin{equation}
\label{eq:ew-angular-character}
\chi_{\rm dur}(n_A,s_C)
=\exp\!\left(n_AA_{\rm rad}
+\frac{s_C}{6}C^*_{\rm rad}\right).
\end{equation}
Cette coordonnée n'est pas introduite comme une étiquette isolée. Sa provenance est
\begin{equation}
\label{eq:ew-signature-provenance}
\begin{aligned}
\APP&\longrightarrow\TRIT\longrightarrow\TPK
\longrightarrow\text{extension survivante}
\longrightarrow\text{route}\\
&\longrightarrow\text{registre}
\longrightarrow(n_A,s_C,\ldots)
\longrightarrow\chi_{\rm dur}.
\end{aligned}
\end{equation}
Les signatures de \(W\), \(Z\) et \(H\) sont des grandeurs typées obtenues au moyen de
cette chaîne ; le présent chapitre les adopte sans les sélectionner à partir de
leurs masses terminales. Dans la réalisation angulaire établie, les secteurs
chargé et neutre satisfont
\begin{equation}
\label{eq:ew-WZ-signatures}
(n_A,s_C)_W=(94,-1),
\qquad
(n_A,s_C)_Z=(95,-1).
\end{equation}
Cette affirmation concerne la composante angulaire commune. Les
coordonnées supplémentaires de raffinement de la loi générale des masses ne sont pas
éliminées à moins que le transducteur physique ne démontre leur annulation.

\begin{theorem}[Défaut de Weinberg dans le schéma des masses physiques de la carte angulaire]
\label{thm:ew-weinberg}
Si \(W\) et \(Z\) sont évalués dans une même section d'échelle et partagent
le raffinement qui accompagne \cref{eq:ew-WZ-signatures}, alors
\begin{equation}
\label{eq:ew-WZ-ratio}
\boxed{
\frac{m_W^{\angle}}{m_Z^{\angle}}
=\ee^{-A_{\rm rad}}=\sqrt{D_A},
\qquad
\left(\frac{m_W^{\angle}}{m_Z^{\angle}}\right)^2=D_A.}
\end{equation}
Dans le schéma des masses physiques,
\begin{equation}
\label{eq:ew-sin2thetaW}
\boxed{
\sin^2\theta_W^{\rm OS,HMT}=1-D_A
=0.2248708807528435698111159035\ldots .}
\end{equation}
Adem\'as,
\begin{equation}
\label{eq:ew-WZ-terminal-ratio}
\frac{m_W^{\angle}}{m_Z^{\angle}}
=0.8804141748331613670128885459\ldots .
\end{equation}
\end{theorem}

\begin{proof}
Par \cref{eq:ew-angular-character,eq:ew-WZ-signatures},
\[
\frac{\chi_{\rm dur}(94,-1)}{\chi_{\rm dur}(95,-1)}
=\exp(-A_{\rm rad});
\]
la coordonnée \(C^*\) s'annule parce que les deux secteurs partagent
\(s_C=-1\). En élevant au carré et en utilisant
\cref{eq:ew-determinant-reader}, on obtient \(D_A\). La définition
des masses physiques \(s_W^2=1-m_W^2/m_Z^2\) produit
\cref{eq:ew-sin2thetaW}.
\end{proof}

Le contenu principal est l'identité entre deux constructions
internes : le déterminant de l'opérateur angulaire et la différence d'un pas
dans la réalisation \(W/Z\). La promotion en sélecteur physique de jauge exige
un entrelaceur qui transporte les feuillets de \(T\) vers la base neutre
\((W^3,B)\) et qui démontre quels raffinements sont conservés ou annulés.

\section{Couplages de jauge en normalisation rationalisée}

Nous définissons le couplage adimensionnel de jauge
\begin{equation}
\label{eq:ew-egauge}
e_g:=\sqrt{4\pi\alpha_{\HMT}}
=0.3028221208717526427662442689\ldots .
\end{equation}
Il ne faut pas confondre \(e_g\) avec une section dimensionnelle de charge
\(e_{\rm int}\) : la première est un couplage en unités naturelles ; la
seconde vit dans la droite de charge de la mécanique dimensionnelle.

Soit
\begin{equation}
\label{eq:ew-sw-cw}
s_W^2=1-D_A,
\qquad
c_W^2=D_A.
\end{equation}
La convention électrofaible rationalisée au niveau de l'arbre donne
\begin{align}
g&=\frac{e_g}{s_W}
=2\sqrt{\frac{\pi\alpha_{\HMT}}{1-D_A}}
=0.6385883427334574283647003809\ldots,
\label{eq:ew-g}\\
g'&=\frac{e_g}{c_W}
=2\sqrt{\frac{\pi\alpha_{\HMT}}{D_A}}
=0.3439541633108502429362319257\ldots,
\label{eq:ew-gprime}\\
\frac{g'}{g}
&=\sqrt{D_A^{-1}-1}
=0.5386164141966093594996610933\ldots .
\label{eq:ew-gprime-over-g}
\end{align}

\begin{proposition}[Relation custodiale à l'arbre]
\label{prop:ew-rho}
Dans la même interface,
\begin{equation}
\label{eq:ew-rho}
\rho_{\rm EW}^{(0)}
:=\frac{(m_W^{\angle})^2}
{(m_Z^{\angle})^2c_W^2}=1.
\end{equation}
\end{proposition}

\begin{proof}
Par \cref{eq:ew-WZ-ratio,eq:ew-sw-cw}, le numérateur relatif comme
\(c_W^2\) sont \(D_A\).
\end{proof}

Les équations \cref{eq:ew-g,eq:ew-gprime} sont des compositions exactes une
fois déclarés le schéma des masses physiques, la rationalisation et l'échelle. Elles
n'affirment pas qu'un changement de schéma de renormalisation laisse leurs
valeurs numériques inchangées.

\section{Constante de Fermi dans l'approximation à l'arbre}

En unités naturelles, l'intégration du propagateur chargé à basse
énergie produit
\begin{equation}
\label{eq:ew-GF-def}
\frac{G_F^{(0)}}{\sqrt2}=\frac{g^2}{8m_W^2}.
\end{equation}
En substituant \cref{eq:ew-g}, on obtient la composition
\begin{equation}
\label{eq:ew-GF}
\boxed{
G_F^{(0)}
=\frac{\pi\alpha_{\HMT}}
{\sqrt2(1-D_A)m_W^2}.}
\end{equation}
En adoptant la masse obtenue précédemment
\begin{equation}
\label{eq:ew-W-output}
m_W^{\HMT}=80.381592550221\ldots\ \mathrm{GeV},
\end{equation}
l'évaluation terminale est
\begin{equation}
\label{eq:ew-GF-value}
\boxed{
G_F^{(0)}
=1.1157162817659203186621784\times10^{-5}\ \mathrm{GeV}^{-2}.}
\end{equation}

La différence par rapport à une constante de Fermi renormalisée n'est pas utilisée
pour choisir un coefficient. Dans la convention
\begin{equation}
\label{eq:ew-Delta-r}
G_F
=\frac{\pi\alpha_{\HMT}}
{\sqrt2(1-D_A)m_W^2(1-\Delta r)},
\end{equation}
l'obligation restante est de dériver \(\Delta r\) à partir du secteur interne des
boucles, des états et de l'échelle. Consulter la valeur observée de \(G_F\) pour définir
\(\Delta r\) serait une reconstruction calibrée, et non une prédiction.

\section{Dépendance orientée du couplage de Higgs}

Dans la carte angulaire, les routes pertinentes sont
\begin{equation}
\label{eq:ew-HW-signatures}
(n_A,s_C)_H=(97,9),
\qquad
(n_A,s_C)_W=(94,-1).
\end{equation}
Par conséquent,
\begin{equation}
\label{eq:ew-HW-ratio}
\boxed{
\frac{m_H^{\angle}}{m_W^{\angle}}
=\frac{\chi_{\rm dur}(97,9)}{\chi_{\rm dur}(94,-1)}
=\exp\!\left(3A_{\rm rad}+\frac53C^*_{\rm rad}\right)
=1.55872087212325548564\ldots .}
\end{equation}
Ici, \(C^*\) ne s'annule pas : le secteur scalaire conserve l'orientation
que le déterminant \(D_A=\ee^{-2A_{\rm rad}}\) ne distingue pas.

Prenons la convention du potentiel
\begin{equation}
\label{eq:ew-Higgs-potential}
V(H)=-\mu_H^2H^\dagger H
+\lambda_H(H^\dagger H)^2,
\end{equation}
pour laquelle
\begin{equation}
\label{eq:ew-Higgs-tree-relations}
m_H^2=2\lambda_Hv^2,
\qquad
m_W^2=\frac{g^2v^2}{4}.
\end{equation}

\begin{corollary}[Autocouplage orienté de Higgs]
\label{cor:ew-Higgs-lambda}
Dans l'interface à l'arbre déclarée,
\begin{equation}
\label{eq:ew-Higgs-lambda}
\boxed{
\lambda_H
=\frac{\pi\alpha_{\HMT}}{2(1-D_A)}
\exp\!\left(6A_{\rm rad}+\frac{10}{3}C^*_{\rm rad}\right)
=0.1238479115482466804662\ldots .}
\end{equation}
Adem\'as,
\begin{equation}
\label{eq:ew-GF-Higgs}
G_F^{(0)}m_H^2=\sqrt2\,\lambda_H.
\end{equation}
\end{corollary}

\begin{proof}
Diviser les deux équations de
\cref{eq:ew-Higgs-tree-relations} donne
\(m_H^2/m_W^2=8\lambda_H/g^2\). Avec
\(g^2=4\pi\alpha_{\HMT}/(1-D_A)\) et le carré de
\cref{eq:ew-HW-ratio}, on obtient \cref{eq:ew-Higgs-lambda}. D'autre part,
\(G_F^{(0)}=1/(\sqrt2v^2)\) et
\(m_H^2=2\lambda_Hv^2\) impliquent
\cref{eq:ew-GF-Higgs}.
\end{proof}

La valeur de \(\lambda_H\) est antérieure à l'incorporation des corrections
radiatives et dépend de
la convention et de l'échelle déclarées. Son intérêt structurel est qu'elle
sépare deux mémoires : le secteur commun de jauge emploie l'invariant pair \(D_A\),
tandis que la route de Higgs retient explicitement \(C^*\).

\section{Clôture déterminantale du secteur électrofaible : mélange, couplages et spectre constitutif}

La chaîne complète de ce chapitre est
\begin{equation}
\label{eq:ew-causal-diagram}
\boxed{
\begin{aligned}
(A,C^*)&\longrightarrow T
\longrightarrow
\begin{cases}
F(T)^2&\longrightarrow
(\widehat\mu,\widehat\varepsilon,\widehat Z,\widehat c),\\
\det T=D_A&\longrightarrow
(s_W^2,c_W^2),
\end{cases}\\
(D_A,\alpha_{\HMT})&\longrightarrow(e_g,g,g'),\\
(D_A,\alpha_{\HMT},m_W)&\longrightarrow G_F^{(0)},\\
(A,C^*,D_A,\alpha_{\HMT})&\longrightarrow
\left(\frac{m_H}{m_W},\lambda_H\right).
\end{aligned}}
\end{equation}
Le chiffre apparaît toujours à la fin. Ni \(\sin^2\theta_W\), ni
\(G_F\), ni \(\lambda_H\) ne sont introduits comme valeurs cibles pour sélectionner
\(A,C^*\) ou les signatures de route.

\begin{proposition}[Résultat principal]
Le même opérateur angulaire détermine le vide par une fonction rationnelle
de ses feuillets et le défaut électrofaible par son déterminant. La
différence d'un pas entre \(W\) et \(Z\) transforme ce déterminant en
\(\sin^2\theta_W^{\rm OS}=1-D_A\) ; la route de Higgs, en revanche, conserve
le contre-angle et conserve l'orientation éliminée par le déterminant.
\end{proposition}

\begin{remark}[Contrôle négatif]
La relation de Weinberg est réfutée dans son domaine si les signatures
angulaires \(W/Z\) ne diffèrent pas uniquement par \(n_A:94\to95\), si une coordonnée
de raffinement ne s'annule pas dans l'interface déclarée ou si
\((m_W^{\angle}/m_Z^{\angle})^2\neq D_A\). La promotion de jauge est
bloquée s'il n'existe pas d'entrelaceur avec la base \((W^3,B)\). Les chiffres
de \(g,g',G_F,\lambda_H\) changent si l'on modifie le schéma ou l'échelle sans
l'enregistrer. Utiliser les valeurs observées de \(G_F\) ou \(m_H\) pour sélectionner
\(\Delta r\), les signatures ou \(C^*\) introduit une rétroaction par la
valeur cible et invalide la dérivation. Confondre \(e_g\) avec la charge
dimensionnelle constitue une erreur de type.
\end{remark}

\subsection*{Portée logique et domaine de validité}
Déterminant angulaire, quotient \(W/Z\) et défaut complémentaire de Weinberg :
\emph{Théorème interne exact dans la carte angulaire conditionnelle}. Couplages \(g,g'\),
\(\rho_{\rm EW}^{(0)}\), Fermi et Higgs : \emph{Composition exacte dans
l'interface électrofaible à l'arbre}. L'entrelaceur de base de jauge et la
correction interne \(\Delta r\) : \emph{Frontière ouverte avec critère de réfutation}.
Fermi et Higgs ne sont pas des constantes indépendantes sélectionnées sans valeurs
cibles : ce sont des compositions exactes qui emploient les signatures, la grandeur
\(m_W\) et la
convention à l'arbre déclarée. Les valeurs conventionnelles ultérieures
sont une \emph{Comparaison métrologique externe}.


\section{Involution des feuillets dans \(4\) représentations}

L'involution \(C^*\mapsto-C^*\) organise quatre réponses distinctes :
\begin{equation}
\label{eq:modular-sheet-parities}
\begin{aligned}
S_{\rm bi}(A,-C^*)&=S_{\rm bi}(A,C^*),\\
\Gamma_{6\to8}(A,-C^*)&=-\Gamma_{6\to8}(A,C^*),\\
(\widehat\mu,\widehat\varepsilon)&\longleftrightarrow
(\widehat\varepsilon,\widehat\mu),\\
\widehat c&\longmapsto\widehat c,
\qquad D_A\longmapsto D_A.
\end{aligned}
\end{equation}
L'information bicouche est paire ; l'orientation aréale est impaire ; les deux
constitutives forment un doublet, et les déterminants restent fixes. Ce tableau
des parités précède toute interprétation comme symétrie physique.

\section{Transformation rationnelle involutive dans l'espace des paramètres CKM}

Soit
\begin{equation}
\label{eq:modular-Mckm}
M_{\rm CKM}=
\begin{pmatrix}
2&-5&28\\
0&7&-25/2\\
1&-20&-2/3
\end{pmatrix},
\qquad
\det M_{\rm CKM}=-\frac{3857}{6}.
\end{equation}
L'inverse existe sur \(\mathbb Q\) et est
\begin{equation}
\label{eq:modular-Mckm-inverse}
M_{\rm CKM}^{-1}=
\begin{pmatrix}
1528/3857&3380/3857&801/3857\\
75/3857&176/3857&-150/3857\\
6/551&-30/551&-12/551
\end{pmatrix}.
\end{equation}
Avec
\begin{equation}
\label{eq:modular-ckm-coordinates}
h=\frac{C^*}{6},
\qquad
\Delta=\frac{180}{\pi}\Delta_4,
\qquad
\boldsymbol\theta=M_{\rm CKM}(A,h,\Delta)^T,
\end{equation}
la coordonnée impaire entre dans la direction
\begin{equation}
\label{eq:modular-ch}
c_h=\begin{pmatrix}-5\\7\\-20\end{pmatrix}.
\end{equation}

Si \(c_A=(2,0,1)^T\) et
\(c_\Delta=(28,-25/2,-2/3)^T\), les trois colonnes
\((c_A,c_h,c_\Delta)\) sont précisément la base transportée par
\(M_{\rm CKM}\). La deuxième ligne de
\cref{eq:modular-Mckm-inverse} est le covecteur \(\ell_h\) qui extrait la
coordonnée impaire :
\begin{equation}
\label{eq:modular-dual-coordinate-check}
\ell_hc_A=0,
\qquad
\ell_hc_h=1,
\qquad
\ell_hc_\Delta=0.
\end{equation}

\begin{theorem}[Réflexion CKM induite par le feuillet]
\label{thm:modular-ckm-reflection}
Soit
\begin{equation}
\label{eq:modular-Dh-lh}
D_h=\operatorname{diag}(1,-1,1),
\qquad
\ell_h=\frac1{3857}\begin{pmatrix}75&176&-150\end{pmatrix}.
\end{equation}
L'involution induite dans l'espace de \(\boldsymbol\theta\) est
\begin{equation}
\label{eq:modular-Rckm}
\boxed{
\mathcal R_{\rm CKM}
=M_{\rm CKM}D_hM_{\rm CKM}^{-1}
=I-2c_h\ell_h.}
\end{equation}
Satisfait
\begin{equation}
\label{eq:modular-Rckm-properties}
\mathcal R_{\rm CKM}^2=I,
\qquad
\det\mathcal R_{\rm CKM}=-1,
\qquad
\mathcal R_{\rm CKM}M_{\rm CKM}=M_{\rm CKM}D_h.
\end{equation}
\end{theorem}

\begin{proof}
La multiplication de
\cref{eq:modular-Mckm,eq:modular-Mckm-inverse} donne l'identité ; en
particulier, \cref{eq:modular-dual-coordinate-check} prouve que
\(c_h\ell_h\) est le projecteur de rang un sur \(\mathbb Qc_h\) le
long de \(\operatorname{span}\{c_A,c_\Delta\}\). Le produit exact
\(\ell_hc_h=1\) implique
\[
(I-2c_h\ell_h)^2
=I-4c_h\ell_h+4c_h(\ell_hc_h)\ell_h=I.
\]
L'opérateur fixe point par point l'hyperplan
\(\ker\ell_h=\operatorname{span}\{c_A,c_\Delta\}\) et change le signe de
\(c_h\). Ses valeurs propres sont donc \(1,1,-1\) ; d'où
\(\det\mathcal R_{\rm CKM}=-1\) et
\(\Tr\mathcal R_{\rm CKM}=1\). Enfin,
\[
\mathcal R_{\rm CKM}M_{\rm CKM}
=M_{\rm CKM}D_hM_{\rm CKM}^{-1}M_{\rm CKM}
=M_{\rm CKM}D_h,
\]
ce qui prouve l'identité d'entrelacement.
\end{proof}

La matrice explicite est
\begin{equation}
\label{eq:modular-Rckm-matrix}
\mathcal R_{\rm CKM}=
\begin{pmatrix}
4607/3857&1760/3857&-1500/3857\\
-150/551&199/551&300/551\\
3000/3857&7040/3857&-2143/3857
\end{pmatrix}.
\end{equation}
Bien que \(\mathcal R_{\rm CKM}\) ne soit pas une réflexion orthogonale pour le
produit euclidien des trois coordonnées angulaires, elle l'est pour la
métrique rationnelle transportée
\begin{equation}
\label{eq:modular-ckm-metric}
G_{\rm CKM}=(M_{\rm CKM}^{-1})^TM_{\rm CKM}^{-1},
\qquad
\mathcal R_{\rm CKM}^TG_{\rm CKM}\mathcal R_{\rm CKM}=G_{\rm CKM}.
\end{equation}
En effet, dans la base \((c_A,c_h,c_\Delta)\), la métrique est l'identité
et l'involution est \(D_h\). C'est pourquoi le terme ``réflexion'' a un
contenu métrique précis et ne décrit pas seulement le fait que
\(\mathcal R_{\rm CKM}^2=I\).
La phase \(\delta_{\rm CP}=9A\) reste fixe sous cette réflexion. C'est
pourquoi \(\mathcal R_{\rm CKM}\) n'est pas la conjugaison CP physique, mais le
transport de l'inversion de feuillet HMT à l'espace de mélange.

L'inverse rationnel de \(M_{\rm CKM}\) retrouve
\begin{align}
A&=\frac{1528\theta_{12}+3380\theta_{23}+801\theta_{13}}{3857},
\label{eq:modular-A-from-ckm}\\
C^*&=\frac{6(75\theta_{12}+176\theta_{23}-150\theta_{13})}{3857}.
\label{eq:modular-C-from-ckm}
\end{align}
La troisième coordonnée se retrouve également sans perte :
\begin{equation}
\label{eq:modular-Delta-from-ckm}
\Delta
=\frac{6\theta_{12}-30\theta_{23}-12\theta_{13}}{551}.
\end{equation}
Ces équations sont des changements de coordonnées. Elles ne font pas de CKM la cause
rétrospective de \(A,C^*\) ou de \(\gammaH\).

Comme reconnaissance terminale, le constructeur préalable détermine
\begin{equation}
\label{eq:modular-ckm-values}
(\theta_{12},\theta_{23},\theta_{13},\delta_{\rm CP})
=(13.003313589509^\circ,
2.398056157332^\circ,
0.213977382244^\circ,
65.676173123554^\circ).
\end{equation}
La réflexion rationnelle est démontrée avant l'utilisation de ces décimales.

\begin{proposition}[Résultat principal]
La transition \(6\to8\) est d'abord formalisée par l'inclusion
\(\iota_{6,8}:\mathcal F_6\hookrightarrow\mathcal F_8\), dont on
obtient
\(90,120\) et le défaut \(-1/12\). Le mot de six trits fournit
ensuite la bijection \(729\leftrightarrow729\), et
\(\mathcal J_{6\to8}\) transporte les deux modes collectifs au bord. Sur
cet antécédent, l'aire et le hamiltonien modulaire reconstruisent exactement
les deux familles ; la purification thermochamp double détermine l'entropie
d'intrication, et la première loi
finie conserve le défaut d'entropie relative. Dans CKM, l'inversion de
feuillet devient une réflexion rationnelle de rang impair un.
\end{proposition}

\begin{remark}[Contrôle négatif]
La généalogie initiale échoue si \(\iota_{6,8}\) n'est pas injective, si les
décomptes ne sont pas \(90/120\), si \(\delta_{6,8}^{\rm inc}\neq-1/12\), si
\(\beta_{729}\) n'est pas bijective ou si elle est présentée à tort comme un
isomorphisme additif. Le transport collectif échoue si
\(\mathcal J_{6\to8}D_{\rm inc}\neq
D_{\rm area}\mathcal J_{6\to8}\). La reconstruction échoue si le
déterminant des observables de bord n'est pas trente ou si
\cref{eq:modular-inverse-channels} ne retrouve pas les occupations.
La TFD échoue si sa trace partielle n'est pas \(\rho_A\). La première loi ne
s'étend pas aux changements finis sans l'entropie relative de
\cref{eq:modular-finite-first-law}. Le type
\(III_{\lambda_\partial}\) est invalidé si un canal ne persiste pas, s'il n'y a pas de
produit compatible ou si le groupe des rapports change. La réflexion CKM
est réfutée par n'importe laquelle des conditions
\(\mathcal R^2\neq I\), \(\det\mathcal R\neq-1\) ou
\(\mathcal RM\neq MD_h\). Une formule
\(\boldsymbol\theta=f(\gammaH)\) sans entrelaceur viole l'indépendance
typée des deux réalisations fonctionnelles de l'antécédent commun.
\end{remark}

\subsection*{Portée logique et domaine de validité}
Inclusion \(H\subset O\), décomptes \(90/120\), deux normalisations du
défaut \(-1/12\), bijection de six trits \(729\leftrightarrow729\),
entrelaceur collectif, échelle \(a_0\), opérateurs finis, spectre aréal,
TFD, défaut orienté, première loi locale et finie et réflexion CKM :
\emph{Théorème interne exact avec ses structures de départ explicites}.
Classification
\(III_{\lambda_\partial}\) : \emph{Théorème avec hypothèses explicites de
persistance}. Interprétation de la copie thermochamp double comme extérieur physique :
\emph{Réalisation physique conditionnelle}. La grandeur \(\gammaH\) est adoptée
comme \emph{Résultat antécédent démontré} ; elle n'est pas redérivée dans ce volume.


Le lien électrofaible part du même opérateur angulaire qui a produit les constitutives du vide, mais utilise une autre fonctionnelle.\par

Pour obtenir \(\mu\), \(\varepsilon\), \(Z\) et \(c\), il fallait conserver le spectre complet. Il fallait distinguer les deux feuillets.\par

Le déterminant fait l’inverse :\par

\[
D_A=\det T=q_+q_-.
\]

En multipliant les valeurs propres, le déterminant intègre les deux orientations dans une grandeur angulaire commune et invariante. Il publie la partie paire et réserve l’identité du feuillet dans la mémoire orientée.\par

C’est pourquoi \(D_A\) est pair sous l’inversion\par

\[
C^*\longmapsto-C^*.
\]

Le secteur \(W/Z\) possède précisément cette structure. Les routes de \(W\) et de \(Z\) sont adjacentes dans la coordonnée angulaire principale et partagent la même coordonnée orientée. Lorsque l’on forme le quotient, la contribution de \(C^*\) s’annule.\par

Il reste\par

\[
\left(\frac{m_W}{m_Z}\right)^2
=
D_A,
\]

et donc\par

\[
\sin^2\theta_W
=
1-D_A.
\]

L’identité révèle que l’angle de Weinberg est le défaut complémentaire du déterminant angulaire.\par

Le déterminant mesure la partie commune qui subsiste lorsque les deux feuillets sont rabattus sur leur produit. \(1-D_A\) mesure la partie complémentaire nécessaire au mélange.\par

Nous pourrions le dire ainsi : le vide conserve toute la réponse ; le secteur électrofaible utilise la compression paire de cette réponse.\par

Le même antécédent angulaire produit deux fonctions différentes :\par

\begin{itemize}[leftmargin=*,itemsep=0.35em]
\item le calcul fonctionnel complet engendre les relations constitutives ;
\item le déterminant engendre l’invariant de mélange.
\end{itemize}

Cela explique pourquoi le vide et l’angle de Weinberg sont deux réalisations fonctionnelles distinctes d’un même antécédent angulaire.\par

À partir de là, les couplages \(g\) et \(g'\) sont, au niveau déclaré, les deux projections du même couplage électromagnétique rationalisé sur les directions sinus et cosinus du mélange :\par

\[
g^2
=
\frac{4\pi\alpha_{\rm HMT}}{1-D_A},
\qquad
g'^2
=
\frac{4\pi\alpha_{\rm HMT}}{D_A}.
\]

L’identité custodiale à l’arbre\par

\[
\rho_{\rm EW}^{(0)}=1
\]

exprime que la même partition angulaire organise de manière cohérente masse et couplage.\par

La constante de Fermi apparaît ensuite comme la réponse effective à basse énergie du canal chargé :\par

\[
G_F^{(0)}
=
\frac{\pi\alpha_{\rm HMT}}
{\sqrt2(1-D_A)m_W^2}.
\]

La constante de Fermi apparaît comme une composition de la structure fine, du défaut angulaire et de l’échelle du porteur \(W\).\par

Le secteur de Higgs introduit une différence décisive. Sa route possède une coordonnée orientée différente de celle de la route de \(W\), de sorte que \(C^*\) persiste. L’autocouplage scalaire publie l’orientation que le déterminant réserve.\par

Le secteur de jauge est pair : il utilise \(D_A\).\par
Le secteur scalaire est orienté : il retient \(C^*\).\par

Cette division est physiquement suggestive. Les champs de jauge organisent la partie commune de la transformation ; le secteur scalaire conserve la mémoire de l’orientation brisée.\par

Les corrections radiatives appartiennent à un niveau ultérieur et doivent être construites prospectivement à partir des boucles, états et échelles internes, avec \(G_F\) comme sortie et évaluation terminale. Le squelette causal est déjà présent : vide, mélange, couplages et réponse de Fermi procèdent de différentes lectures du même état angulaire enrichi.\par


\section{Clôture angulaire électrofaible et conditions de compatibilité du transducteur de jauge}
\label{sec:p3-final:cierre-electrodebil}

Soient
\[
 c_W=\sqrt{D_A},
 \qquad
 s_W=\sqrt{1-D_A}.
\]
La rotation
\[
 \begin{pmatrix}Z\\A\end{pmatrix}
 =
 \begin{pmatrix}c_W&-s_W\\s_W&c_W\end{pmatrix}
 \begin{pmatrix}W^3\\B\end{pmatrix}
\]
est orthogonale et préserve l’orientation. Les identités
\[
 \left(\frac{m_W}{m_Z}\right)^2=D_A,
 \qquad
 \sin^2\theta_W=1-D_A
\]
sont closes dans la carte angulaire. La promotion de jauge exige que le
transducteur préserve la forme cinétique, les charges et les identités de Ward ;
la correction \(\Delta r\) appartient à la dynamique radiative et ne s’obtient pas
en l’ajustant à \(G_F\).


%
\endgroup

\HMTFlushCurrentChapter
\chapter{Dérivation intrinsèque de \(G\) et réciprocité contragrédiente \texorpdfstring{\(G\leftrightarrow\hbar\)}{G<->hbar}: fermeture de phase, échelles de Compton et de Planck et invariance de \(G\hbar\)}%
\begingroup
\renewcommand{\chapter}[2][]{}%
\chapter{Dérivation intrinsèque de \(G\) et réciprocité \(G\leftrightarrow\hbar\) : clôture de phase, longueurs de Compton et de Planck et invariance aréale}
\label{chap:p3-final:50:gravitacion_accion}
\label{ch:gravity}

\section{Chaîne constitutive du secteur gravitationnel : échelle chronologique \(108\), réduction tensorielle, holonomie et évaluation métrologique de \(G\)}

La gravité s'organise ici en quatre strates qu'il ne faut pas fusionner :

\begin{equation}
\label{eq:gravity-causal-chain}
\begin{aligned}
\text{structure chronologique CT108 et droites d'action}
&\longrightarrow \text{sélecteur de périodicité }\theta_{108}^{\rm circ}
\longrightarrow G_{108},\\
\text{incidence icosaédrique}
&\longrightarrow V_{12}\xrightarrow{P_5}V_5
\xrightarrow{\Gamma_5}\operatorname{STF}_3
\xrightarrow{\Lambda(n)}\mathcal H_{\rm TT}(n),\\
\text{route holonomique enrichie}
&\dashrightarrow G_{\rm hol},\\
\text{carte métrologique}
&\longrightarrow 6.67430\times10^{-11}\,
\mathrm{m^3\,kg^{-1}\,s^{-2}}.
\end{aligned}
\end{equation}

La première ligne exprime une condition intrinsèque et exacte de
périodicité. La deuxième construit
le support et sa réduction transverse sans trace. La troisième conserve une
frontière de construction explicite. La quatrième est une évaluation dans une
carte décimale déclarée : elle identifie la grandeur, mais ne sélectionne pas le caractère
gravitationnel.

\section{Droite de grandeur gravitationnelle et échelle chronologique \(108\)}

Soient \(c_{\rm int}>0\), \(t_0>0\),
\(\ell_0=c_{\rm int}t_0\) et une section positive
\(\hbar_a\) de la droite d'action, avec
\(a\in\{\mathrm{pre},\mathrm{ret}\}\). Le transducteur dimensionnel est

\begin{equation}
\label{eq:gravity-transducer}
G_a(\theta)
=\theta\frac{c_{\rm int}^{5}t_0^2}{\hbar_a}
=\theta\frac{c_{\rm int}^{3}\ell_0^2}{\hbar_a},
\qquad \theta>0.
\end{equation}

L'égalité des deux expressions utilise seulement
\(\ell_0=c_{\rm int}t_0\), et
\([G]=L^3M^{-1}T^{-2}\). La dimension est fixée par la droite ; le
problème mathématique est de sélectionner le caractère adimensionnel \(\theta\).

Le calendrier CT108 définit

\begin{equation}
\label{eq:gravity-clock108}
T_{108}=108t_0,
\qquad
\omega_{108}=\frac{2\pi}{108t_0},
\qquad
R_{108}=\frac{c_{\rm int}}{\omega_{108}}
=\frac{54}{\pi}\ell_0.
\end{equation}

Pour chaque feuillet d'action, soient

\begin{equation}
\label{eq:gravity-clock-energy-mass}
E_{108,a}=\hbar_a\omega_{108},
\qquad
m_{108,a}=\frac{E_{108,a}}{c_{\rm int}^2}.
\end{equation}

Alors, la longueur de Compton réduite coïncide avec le rayon du cycle :

\begin{equation}
\label{eq:gravity-compton-clock}
\bar\lambda_{C,a}
=\frac{\hbar_a}{m_{108,a}c_{\rm int}}
=R_{108}.
\end{equation}

\begin{theorem}[Unicité du sélecteur sous la condition de périodicité HMT]
\label{thm:gravity-circular-selector}
Avec la convention réduite
\(r_g=Gm/c_{\rm int}^2\) et la prémisse de périodicité HMT qui identifie le rayon
gravitationnel du quantum CT108 à son rayon de phase, la condition de clôture
\begin{equation}
\label{eq:gravity-closing-condition}
r_{g,a}(\theta)=R_{108}
\end{equation}
sélectionne une unique solution positive,
\begin{equation}
\label{eq:gravity-theta108}
\boxed{
\theta_{108}^{\rm circ}=\left(\frac{54}{\pi}\right)^2
=295.4525715010569415303525\ldots .}
\end{equation}
Dans cette réalisation,
\begin{equation}
\label{eq:gravity-four-lengths}
R_{108}=\bar\lambda_{C,a}=r_{g,a}=\ell_P^{\rm circ}.
\end{equation}
\end{theorem}

\begin{proof}
Par \cref{eq:gravity-transducer,eq:gravity-clock-energy-mass},
\[
r_{g,a}(\theta)
=\frac{G_a(\theta)m_{108,a}}{c_{\rm int}^2}
=\theta c_{\rm int}t_0^2\omega_{108}
=\theta\frac{\pi}{54}\ell_0.
\]
En le comparant à
\(R_{108}=(54/\pi)\ell_0\), on obtient nécessairement
\(\theta=(54/\pi)^2\). L'équation est linéaire en \(\theta\) et tous les
facteurs sont positifs, donc la solution est unique.
\end{proof}

L'égalité \(r_g=R_{108}\) ne se déduit pas de la dimensionnalité de
\cref{eq:gravity-transducer}. C'est la condition géométrique explicite de cette
réalisation HMT : action, Compton et retour périodique doivent déterminer une même
longueur. Le théorème démontre que, cette condition structurelle étant acceptée, le
caractère \(\theta\) est fixé sans utiliser la valeur SI de \(G\).

\begin{remark}[Convention de rayon]
Le théorème utilise le rayon gravitationnel réduit \(Gm/c^2\). Si on le remplace
par le rayon de Schwarzschild \(2Gm/c^2\), le sélecteur change d'un facteur
deux. Les deux expressions correspondent à des conditions géométriques distinctes,
c'est pourquoi la convention doit être déclarée avant le calcul.
\end{remark}

\section{Réciprocité \(G\leftrightarrow\hbar\) et conservation de l'aire de Planck}

\begin{corollary}[Famille de Planck associée à la périodicité temporelle de \(108\) phases]
\label{cor:gravity-planck-family}
Pour \(G_a=G_a(\theta_{108}^{\rm circ})\),
\begin{align}
\ell_P&=\frac{54}{\pi}\ell_0,
&t_P&=\frac{54}{\pi}t_0,
\label{eq:gravity-planck-length-time}\\
E_{P,a}&=\frac{\hbar_a}{t_P}
=\frac{\pi\hbar_a}{54t_0}
=\frac{h_a}{108t_0},
&\ell_P^2&=\frac{G_a\hbar_a}{c_{\rm int}^3}
=\theta_{108}^{\rm circ}\ell_0^2.
\label{eq:gravity-planck-energy-area}
\end{align}
Adem\'as,
\begin{equation}
\label{eq:gravity-planck-force-power-density}
F_{P,a}=\frac{c_{\rm int}^4}{G_a},
\qquad
P_{P,a}=\frac{c_{\rm int}^5}{G_a},
\qquad
\rho_{P,a}=\frac{\hbar_a}
{(\theta_{108}^{\rm circ})^2c_{\rm int}^5t_0^4}.
\end{equation}
\end{corollary}

\begin{proof}
Les trois premières identités découlent de
\(\ell_P=R_{108}\), \(t_P=\ell_P/c_{\rm int}\) et
\(h_a=2\pi\hbar_a\). Pour l'aire,
\[
\frac{G_a\hbar_a}{c_{\rm int}^3}
=\theta_{108}^{\rm circ}c_{\rm int}^2t_0^2
=\theta_{108}^{\rm circ}\ell_0^2.
\]
Les expressions restantes sont les compositions dimensionnelles de force,
de puissance et de masse par volume avec la même section.
\end{proof}

Soit
\begin{equation}
\label{eq:gravity-Ract}
R_{\rm act}:=\frac{\hbar_{\rm pre}}{\hbar_{\rm ret}}.
\end{equation}
Comme l'action apparaît au numérateur de la masse chronologique et au
dénominateur de \(G\),
\begin{equation}
\label{eq:gravity-twoway-covariance}
\frac{m_{108,\rm pre}}{m_{108,\rm ret}}=R_{\rm act},
\qquad
\frac{G_{\rm pre}}{G_{\rm ret}}=R_{\rm act}^{-1},
\qquad
G_a\hbar_a
=\theta_{108}^{\rm circ}c_{\rm int}^5t_0^2.
\end{equation}
Par conséquent, \(\ell_P^2=G_a\hbar_a/c^3\) ne change pas de feuillet. La gravité et
l'action s'orientent réciproquement, tandis que l'aire conserve la
géométrie commune. Cette annulation est l'interface de confluence vers l'opérateur
d'aire du \cref{ch:modular}.

La deuxième relation fondamentale apparaît lorsque l’on compare l’action et la gravité.\par

Dans HMT, l’action mesure combien il se produit et conserve simultanément l’orientation de la réalisation. Il existe une section antérieure au retour et une autre postérieure. La phase peut revenir accompagnée d’une histoire de mémoire différente. Cette différence est stockée dans le rapport bidirectionnel des deux sections d’action.\par

La gravité répond de manière réciproque.\par

Le transducteur gravitationnel a la forme\par

\[
G_a(\theta)
=
\theta\frac{c_{\rm int}^{5}t_0^2}{\hbar_a}.
\]

L’action apparaît au dénominateur. Par conséquent, si, lors du changement de feuillet, \(\hbar\) est multiplié par un rapport \(R_{\rm act}\), la constante gravitationnelle est multipliée par le rapport inverse :\par

\[
\frac{\hbar_{\rm pre}}{\hbar_{\rm ret}}
=
R_{\rm act},
\qquad
\frac{G_{\rm pre}}{G_{\rm ret}}
=
R_{\rm act}^{-1}.
\]

L’action et la gravité forment ici une paire contragrédiente : lorsque l’une avance dans une direction, l’autre compense exactement ce déplacement.\par

Ce qui reste invariant est\par

\[
G_a\hbar_a.
\]

Et c’est précisément la combinaison qui détermine l’aire de Planck :\par

\[
\ell_P^2=\frac{G_a\hbar_a}{c_{\rm int}^3}.
\]

L’action mémorise le feuillet. La gravité répond à cette mémoire. L’aire reste invariante.\par

C’est une idée extraordinairement féconde. L’aire de Planck émerge comme l’objet qui subsiste lorsque l’action et la gravité se transforment réciproquement.\par

On pourrait le dire ainsi : \(\hbar\) conserve l’orientation dynamique ; \(G\) conserve la compensation géométrique ; \(\ell_P^2\) conserve ce sur quoi les deux réalisations s’accordent.\par

L’aire est donc une monnaie d’échange entre physique quantique et géométrie, parce que la transformation bidirectionnelle de \(\hbar\) et de \(G\) sélectionne exactement cette combinaison comme invariant.\par

Cette réciprocité permet de comprendre autrement la relation entre information et gravité. L’information exige une distinction entre les états ; l’action mesure le coût de la transition entre eux. Mais si cette transition doit acquérir une géométrie, la réponse gravitationnelle change de manière inverse pour maintenir stable le quantum d’aire.\par

La géométrie incorpore l’information de phase par une compensation exacte.\par

Selon une lecture machienne, c’est déjà un résultat très suggestif. Une grandeur locale de géométrie est déterminée relationnellement par l’orientation et le retour de l’état complet. L’aire invariante restreint cette dépendance relationnelle et en fixe la cohérence.\par

La version HMT du principe relationnel consiste en quelque chose de plus précis que la maxime « la matière détermine la géométrie » : l’orientation de l’action et la réponse gravitationnelle se codéterminent de telle manière que la géométrie élémentaire reste invariante.\par

L’action contient l’histoire. La gravité ajuste la réponse. L’aire conserve l’accord.\par


La dérivation de \(G\) porte cette réciprocité à une condition géométrique beaucoup plus concrète.\par

La mécanique dimensionnelle détermine d’abord quel type d’objet \(G\) peut être. Elle fixe sa droite de grandeur et son homogénéité dimensionnelle. La clôture chronologique sélectionne ensuite sa valeur : la dimension détermine la classe de grandeur et la clôture détermine son coefficient adimensionnel.\par

La sélection procède de la périodicité temporelle de (108) phases.\par

Le cycle \(108\) détermine une fréquence, un rayon de phase et une énergie. À partir de cette énergie apparaît une masse chronologique. Et en construisant la longueur de Compton réduite de cette masse, on obtient exactement le même rayon du cycle :\par

\[
\bar\lambda_C=R_{108}.
\]

La phase et la matière expriment déjà la même longueur avant l’introduction de \(G\).\par

La gravité intervient alors. On construit le rayon gravitationnel réduit\par

\[
r_g=\frac{Gm}{c^2},
\]

et l’on exige que la clôture gravitationnelle identifie le rayon de phase, la longueur de Compton et le rayon gravitationnel comme une seule échelle structurelle :\par

\[
r_g=R_{108}=\bar\lambda_C.
\]

Cette condition sélectionne de manière univoque\par

\[
\theta_{108}
=
\left(\frac{54}{\pi}\right)^2.
\]

Ainsi,\par

\[
R_{108}
=
\bar\lambda_C
=
r_g
=
\ell_P.
\]

La fonction structurelle de \(\theta_{108}\) détermine la signification de son nombre décimal terminal : elle montre ce que fait \(G\).\par

\(G\) apparaît comme le coefficient nécessaire pour faire coïncider les quatre longueurs publiées par ces langages :\par

\begin{itemize}[leftmargin=*,itemsep=0.35em]
\item le langage de la phase produit \(R_{108}\) ;
\item le langage quantique produit \(\bar\lambda_C\) ;
\item le langage gravitationnel produit \(r_g\) ;
\item le langage de Planck produit \(\ell_P\).
\end{itemize}

La constante gravitationnelle est le médiateur qui empêche ces quatre lecteurs de se séparer.\par

Dit de manière plus physique : \(G\) exprime combien la géométrie doit répondre pour qu’un quantum dont la masse provient de l’horloge \(108\) se referme sur sa propre phase.\par

C’est la relation qui rend possible que matière, phase et géométrie soient trois publications compatibles du même événement.\par

Ici, la lecture machienne devient plus forte. L’échelle gravitationnelle locale est sélectionnée par une condition de clôture globale du cycle. Le rayon local d’un quantum acquiert un sens en s’intégrant dans une périodicité complète. L’inertie, la phase et la géométrie se déterminent relationnellement.\par

La portée démontrée fournit un mécanisme mathématique concret pour une lecture machienne : une propriété locale — le couplage gravitationnel — est sélectionnée par la cohérence de l’état global. Les autres formulations historiques du principe de Mach conservent leur domaine propre.\par

Ensuite apparaît la dimension cinq.\par

Le scalaire \(G\) fixe l’intensité gravitationnelle et le porteur spécifie sa structure représentationnelle. HMT construit un espace tensoriel symétrique de trace nulle et de dimension cinq. Ce porteur admet cinq poids hélicoïdaux :\par

\[
+2,\quad +1,\quad 0,\quad -1,\quad -2.
\]

Ce sont les cinq degrés de liberté naturels d’un tenseur symétrique sans trace en trois dimensions avant l’imposition des contraintes dynamiques d’une théorie particulière.\par

Cela clarifie un point souvent confondu. « Cinq composantes » désigne le porteur complet, tandis que les deux hélicités désignent la réduction transverse de masse nulle en relativité générale.\par

Cela signifie que le porteur complet possède cinq coordonnées irréductibles. Ensuite, la dynamique décide quelle partie de ce porteur est physique :\par

\begin{itemize}[leftmargin=*,itemsep=0.35em]
\item dans une réalisation massive de spin deux, les cinq poids peuvent subsister ;
\item dans une réalisation sans masse avec invariance de jauge, la projection transverse sans trace élimine les secteurs \(\pm1\) et \(0\) ;
\item les hélicités \(\pm2\) subsistent.
\end{itemize}

Ainsi, HMT fournit un espace antérieur à la décision dynamique. Elle construit le porteur complet à cinq composantes et contient, par la projection correspondante, sa réduction physique à deux hélicités.\par

La beauté réside dans la coexistence ordonnée de deux affirmations :\par

\begin{itemize}[leftmargin=*,itemsep=0.35em]
\item la géométrie de spin deux possède un porteur à cinq composantes ;
\item le rayonnement gravitationnel sans masse occupe un sous-espace bidimensionnel.
\end{itemize}

Les deux affirmations forment une réduction cohérente et précise.\par

Et cette réduction conserve un autre enseignement central : le scalaire \(G\) fixe l’intensité commune et l’information d’orientation réside dans le porteur tensoriel. L’intensité réside dans la constante, les hélicités résident dans la représentation et la généalogie complète réunit grandeur et porteur.\par


%
\endgroup
% Secciones 1--3 del delta gravitatorio íntegro.
\section{Domaine généalogique de la réalisation gravitationnelle}
\label{sec:l9s102:gravedad-dominio-genealogico}

La réalisation gravitationnelle reçoit comme objet de départ une section du
continu discret déjà construite par la composition
\begin{equation}
 \mathrm{APP}\longrightarrow\mathrm{TRIT}\longrightarrow\mathrm{TPK}
 \longrightarrow X_{\mathrm{TPK}}^{\mathrm{enr}}
 \longrightarrow\mathcal C_{\mathrm{cont}}^{\mathrm{disc}}.
 \label{eq:l9s102:gravedad-cadena-genealogica}
\end{equation}
APP conserve les évaluations additive et multiplicative au moyen du résidu,
du quotient et de la retenue ; TRIT détermine le régime et l'orientation ; le TPK
incorpore le sélecteur tritique de phase sans éliminer aucun des deux feuillets
APP, effectue le transport cardinal, conserve la mémoire et réalise le retour
de phase sans réinitialiser l'état. Le transducteur
dimensionnel agit ensuite sur la section survivante du continu :
\begin{equation}
 \mathfrak G_{108}:\mathcal C_{\mathrm{cont}}^{\mathrm{disc}}
 \times\mathcal L_S^+
 \longrightarrow\mathcal G_{\mathrm{HMT\text{-}MD}},
 \qquad
 (x,\hbar_a)\longmapsto
 \bigl(G_a,R_{108},m_{108,a},\ell_{P,a}^{,2}\bigr).
 \label{eq:l9s102:gravedad-mapa-realizacion}
\end{equation}
Les cardinaux \(90/120\), la périodicité \(108\), l'action et le couple
angulaire entrent avec leur généalogie conservée. Aucune identification
gravitationnelle ultérieure ne sélectionne rétroactivement ces antécédents.

\section{Transducteur gravitationnel et sélecteur chronologique de périodicité}
\label{sec:l9s102:gravedad-transductor}

Soient \(c_{\mathrm{int}}>0\), \(t_0>0\),
\(\ell_0=c_{\mathrm{int}}t_0\) et
\(\hbar_a\in\mathcal L_S^+\), avec
\(a\in\{\mathrm{pre},\mathrm{ret}\}\). Le transducteur dimensionnel est
\begin{equation}
 G_a(\theta)
 =\theta\frac{c_{\mathrm{int}}^5t_0^2}{\hbar_a}
 =\theta\frac{c_{\mathrm{int}}^3\ell_0^2}{\hbar_a},
 \qquad \theta>0.
 \label{eq:l9s102:g-transductor}
\end{equation}
La chronologie cyclique de \(108\) états détermine
\begin{equation}
 T_{108}=108t_0,
 \qquad
 \omega_{108}=\frac{2\pi_{\rm HMT}}{108t_0},
 \qquad
 R_{108}=\frac{c_{\mathrm{int}}}{\omega_{108}}
 =\frac{54}{\pi_{\rm HMT}}\ell_0.
 \label{eq:l9s102:g-reloj108}
\end{equation}
Le quantum chronologique et sa masse associée sont
\begin{equation}
 E_{108,a}=\hbar_a\omega_{108},
 \qquad
 m_{108,a}=\frac{E_{108,a}}{c_{\mathrm{int}}^2}.
 \label{eq:l9s102:g-cuanto-masa}
\end{equation}
La longueur de Compton réduite satisfait
\begin{equation}
 \bar\lambda_{C,a}
 =\frac{\hbar_a}{m_{108,a}c_{\mathrm{int}}}=R_{108}.
 \label{eq:l9s102:g-compton}
\end{equation}

\begin{theorem}[Unicité du sélecteur gravitationnel de périodicité]
\label{thm:l9s102:g-selector}
Avec la convention réduite \(r_g=Gm/c_{\mathrm{int}}^2\), la condition
\(r_{g,a}(\theta)=R_{108}\) possède l'unique solution positive
\begin{equation}
 \boxed{
 \theta_{108}^{\mathrm{circ}}=\left(\frac{54}{\pi_{\rm HMT}}\right)^2.}
 \label{eq:l9s102:g-theta}
\end{equation}
La réalisation correspondante identifie structurellement
\begin{equation}
 R_{108}=\bar\lambda_{C,a}=r_{g,a}=\ell_P^{\mathrm{circ}}.
 \label{eq:l9s102:g-cuatro-longitudes}
\end{equation}
\end{theorem}

\begin{proof}
Par \cref{eq:l9s102:g-transductor,eq:l9s102:g-cuanto-masa},
\[
 r_{g,a}(\theta)
 =\frac{G_a(\theta)m_{108,a}}{c_{\mathrm{int}}^2}
 =\theta\frac{\pi_{\rm HMT}}{54}\ell_0.
\]
L'égalité avec \(R_{108}=(54/\pi_{\rm HMT})\ell_0\) impose
\(\theta=(54/\pi_{\rm HMT})^2\). La linéarité en \(\theta\) et la positivité des
facteurs déterminent l'unicité.
\end{proof}

Dans la carte où \(c_{\mathrm{int}}=(\mu\varepsilon)^{-1/2}\), le
transducteur adopte la forme
\begin{equation}
 \boxed{
 G_a=\left(\frac{54}{\pi_{\rm HMT}}\right)^2
 \frac{t_0^2}{\hbar_a(\mu\varepsilon)^{5/2}}.}
 \label{eq:l9s102:g-mu-epsilon}
\end{equation}
Le produit constitutif \(\mu\varepsilon\) gouverne le cône causal et
l'échelle gravitationnelle ; le quotient \(Z^2=\mu/\varepsilon\) conserve
l'orientation d'impédance.

\section{Covariance bidirectionnelle de la gravité et de l'action}
\label{sec:l9s102:g-hbar}

On définit
\begin{equation}
 R_{\mathrm{act}}=\frac{\hbar_{\mathrm{pre}}}
 {\hbar_{\mathrm{ret}}}.
 \label{eq:l9s102:r-act}
\end{equation}
La masse chronologique et le couplage gravitationnel se transforment
de manière contragrédiente :
\begin{equation}
 \frac{m_{108,\mathrm{pre}}}{m_{108,\mathrm{ret}}}=R_{\mathrm{act}},
 \qquad
 \frac{G_{\mathrm{pre}}}{G_{\mathrm{ret}}}=R_{\mathrm{act}}^{-1}.
 \label{eq:l9s102:g-hbar-reciprocidad}
\end{equation}
Leur produit demeure fixe :
\begin{equation}
 \boxed{
 G_a\hbar_a
 =\theta_{108}^{\mathrm{circ}}c_{\mathrm{int}}^5t_0^2,
 \qquad
 \ell_P^2=\frac{G_a\hbar_a}{c_{\mathrm{int}}^3}
 =\theta_{108}^{\mathrm{circ}}\ell_0^2.}
 \label{eq:l9s102:g-hbar-invariante}
\end{equation}
L'action conserve l'orientation du feuillet ; l'aire de Planck conserve la
géométrie commune aux deux orientations.

\begin{theorem}[Indépendance de carte des sections \(h\) et \(G\)]
\label{thm:l9s102:covariancia-unidades-g}
Soient \(u_L,u_T,u_S\) des bases positives des droites dimensionnelles de
longueur, de temps et d'action, et remplaçons-les par
\[
 u_L'=\lambda_Lu_L,\qquad
 u_T'=\lambda_Tu_T,\qquad
 u_S'=\lambda_Su_S,
 \qquad \lambda_L,\lambda_T,\lambda_S>0.
\]
Les coordonnées des mêmes sections vérifient
\begin{equation}
 c_{\rm int}'=\frac{\lambda_T}{\lambda_L}c_{\rm int},\qquad
 t_0'=\frac{t_0}{\lambda_T},\qquad
 \hbar_a'=\frac{\hbar_a}{\lambda_S},\qquad
 G_a'=\frac{\lambda_S\lambda_T^3}{\lambda_L^5}G_a.
 \label{eq:l9s102:cambio-carta-g}
\end{equation}
Pour la base induite
\(u_G=u_L^5u_T^{-3}u_S^{-1}\), on obtient
\begin{equation}
 \hbar_a'u_S'=\hbar_au_S,
 \qquad G_a'u_G'=G_au_G.
 \label{eq:l9s102:secciones-h-g-invariantes}
\end{equation}
Par conséquent,
\(h_a=2\pi_{\rm HMT}\hbar_a\),
\(G_a=\theta_{108}^{\rm circ}c_{\rm int}^5t_0^2/\hbar_a\) et
\(G_a\hbar_a=\theta_{108}^{\rm circ}c_{\rm int}^5t_0^2\)
sont des identités de sections, non des égalités dépendant d'un choix
d'unités. La carte métrologique publie leurs coordonnées numériques après
la construction HMT--MD.
\end{theorem}

\begin{proof}
Les trois premières transformations de
\eqref{eq:l9s102:cambio-carta-g} sont la loi contragrédiente des coordonnées
dans \(L\otimes T^{-1}\), \(T\) et \(S\). Leur substitution dans
\eqref{eq:l9s102:g-transductor} produit la quatrième. La base de la droite
gravitationnelle se transforme par le facteur inverse, ce qui prouve
\eqref{eq:l9s102:secciones-h-g-invariantes}. Les scalaires
\(\pi_{\rm HMT}\) et \(\theta_{108}^{\rm circ}\) sont adimensionnels et
demeurent invariants.
\end{proof}


\HMTFlushCurrentChapter
\chapter{Dérivation HMT du paramètre de Barbero--Immirzi et de l'aire minimale : incidence hexade--octade, canaux \(90/120\) et spectre de l'opérateur d'aire}%
\begingroup
\renewcommand{\chapter}[2][]{}%
\chapter{Dérivation HMT du paramètre de Barbero--Immirzi et de l’aire minimale : incidence hexade--octade, canaux \(90/120\) et spectre de l’opérateur d’aire}
\label{chap:p3-final:51:barbero_area}
\section{Morphisme d’incidence hexade--octade et conservation des canaux orientés \(90/120\)}

La dodécade et l’alignement du
\cref{sec:w12-w24-elevacion} étant fixés, soit \(H\) une hexade dans les coordonnées HMT,
soit \(\ell(H)\subset D\) son image et soit
\(O_H=\iota_D(\ell(H))\) son unique relèvement. En marquant un point
et une paire de points de l’hexade résiduelle, on obtient
\[
F_6(H)=\ell(H)\times\binom{\ell(H)}2,
\qquad
|F_6|=6\binom62=90.
\]
Si le point marqué peut parcourir l’octade :
\[
F_8(H)=O_H\times\binom{\ell(H)}2,
\qquad
|F_8|=8\binom62=120.
\]
L’inclusion \(\ell(H)\hookrightarrow O_H\) induit le morphisme d’incidence
\[
 \mathcal I_{6\to8}:F_6(H)\hookrightarrow F_8(H).
\]
En normalisant la différence par le facteur
choisi \(24\binom62\), on obtient
\[
\boxed{
\frac{90-120}{24\binom62}
=\frac{6-8}{24}
=-\frac1{12}.}
\]
L’égalité est un défaut normalisé de l’interface combinatoire
\(6\to8\). Le facteur \(24\binom62\) fait partie de la définition de cette
normalisation ; l’inclusion à elle seule détermine la différence \(-30\).
L’opérateur \(\mathcal I_{6\to8}\) a pour domaine les configurations
marquées d’incidence ; ce n’est ni l’extracteur transversal de l’état
dodécaphasique ni la famille des moments angulaires.

% Construcción dodecafásica del funcional de retorno.
\section{Chaîne fondamentale dodécaphasique et fonctionnelle hexadique--octadique}

L’incidence active se construit avant toute réalisation gravitationnelle.
Soit
\[
 H_{\mathrm{act}}=\{1,2,4,5,7,8\},
 \qquad
 \mathcal P_2(H_{\mathrm{act}})
 =\{A\subset H_{\mathrm{act}}:|A|=2\}.
\]
Alors
\[
 |H_{\mathrm{act}}|=6,
 \qquad
 |\mathcal P_2(H_{\mathrm{act}})|=15,
\]
et les incidences antérieure et postérieure à l’extension ont pour cardinaux
\[
 \mathcal F_{90}=H_{\mathrm{act}}\times
 \mathcal P_2(H_{\mathrm{act}}),
 \qquad |\mathcal F_{90}|=90,
\]
\[
 \mathcal F_{120}=\{1,\ldots,8\}\times
 \mathcal P_2(H_{\mathrm{act}}),
 \qquad |\mathcal F_{120}|=120.
\]
Ces deux cardinaux sont des sorties de la règle de phase ; ce ne sont ni des angles ni
des paramètres reçus de la réalisation physique.

Notons \(C_{12}=\mathbb Z/12\mathbb Z\) le registre temporel du TPK et
\(\partial_{\rm ret}\) son unique couture de retour par tour. Dans le
module libre des événements orientés, on considère la chaîne fondamentale
\begin{equation}
 c_{12}^{+}
 =\sum_{j\in C_{12}}[j,\mathcal F_{90}]
  -[\partial_{\rm ret},\mathcal F_{120}].
 \label{intsuccpass:eq:c12-fundamental-chain}
\end{equation}
La somme contient exactement les douze secteurs de la roue ; le second
terme est unique parce qu’un tour possède une seule couture de clôture. Son signe
est le signe de frontière. L’involution bidirectionnelle envoie
\(c_{12}^{+}\) sur \(c_{12}^{-}=-c_{12}^{+}\) : elle change l’orientation de la
chaîne, mais n’altère ni ses multiplicités absolues ni ses types
d’incidence.

Pour \(m\in\{90,120\}\), soit
\[
 S_m(q)=\frac{q^m}{1-q^{3m}}.
\]
Le lecteur de retour \(\mathscr R\) est défini sur les générateurs de
\eqref{intsuccpass:eq:c12-fundamental-chain} par
\[
 \mathscr R([j,\mathcal F_{90}])=S_{90},
 \qquad
 \mathscr R([\partial_{\rm ret},\mathcal F_{120}])=S_{120}.
\]
La première égalité est constante sur \(C_{12}\) par covariance sous la
rotation de phase ; la seconde a un domaine distinct et n’agit que sur la
couture étendue. Par linéarité,
\begin{equation}
 \boxed{\mathscr R(c_{12}^{+})=12S_{90}-S_{120}.}
 \label{intsuccpass:eq:dodecaphase-forces-12minus1}
\end{equation}
Ainsi, les coefficients \(12:-1\) sont l’image de la chaîne fondamentale : douze
cellules de phase et une frontière orientée. Ils ne sont sélectionnés ni au moyen de la valeur
conventionnelle du paramètre de Barbero--Immirzi, ni au moyen d’une comparaison
décimale, ni en imposant au préalable le coefficient de pôle qui sera calculé
ci-après.

\begin{theorem}[Coefficient de pôle de la chaîne dodécaphasique]
La fonctionnelle \(F(q)=\mathscr R(c_{12}^{+})(q)
=12S_{90}(q)-S_{120}(q)\), déterminée par les douze secteurs
et leur couture orientée, a pour coefficient \(1/24\) dans
\((1-q)^{-1}\). Son résidu complexe en \(q=1\) est \(-1/24\).
Le coefficient de pôle est une sortie de la fonctionnelle déjà construite.
\end{theorem}

\begin{proof}
Pour une fonctionnelle ayant un pôle simple en \(q=1\), nous écrivons
\[
 p_1(F):=\lim_{q\uparrow1}(1-q)F(q).
\]
L’identité
\[
 1-q^{3m}=(1-q)\sum_{j=0}^{3m-1}q^j
\]
implique
\[
 p_1(S_m)
 =\lim_{q\uparrow1}\frac{q^m}{\sum_{j=0}^{3m-1}q^j}
 =\frac1{3m}.
\]
Par linéarité,
\[
 p_1\bigl(\mathscr R(c_{12}^{+})\bigr)
 =\frac{12}{270}-\frac1{360}
 =\frac{48-3}{1080}=\frac1{24}.
\]
Ces égalités s’effectuent dans \(\mathbb Q\). Le résidu complexe
est le coefficient de \((q-1)^{-1}\), de sorte que
\[
 \operatorname*{Res}_{q=1}\mathscr R(c_{12}^{+})(q)
 =-p_1\bigl(\mathscr R(c_{12}^{+})\bigr)=-\frac1{24}.
\]
\end{proof}

L’équation diophantienne
\[
 \frac a{270}-\frac b{360}=\frac1{24}
 \quad\Longleftrightarrow\quad
 4a-3b=45
\]
reste un certificat arithmétique ultérieur. À elle seule, elle admet la
famille \((a,b)=(12+3t,1+4t)\) ; la chaîne
\eqref{intsuccpass:eq:c12-fundamental-chain} élimine cette ambiguïté avant le calcul,
car modifier \((12,1)\) altère la multiplicité de l’orbite ou de sa couture
et ne représente donc plus un tour fondamental du TPK.

Avec les canaux conjugués \(q_\pm\) déjà engendrés par la branche angulaire, on
obtient enfin
\[
 \mathscr R(c_{12}^{+})(q_-)-
 \mathscr R(c_{12}^{+})(q_+)
 =12\Delta S_{90}-\Delta S_{120}.
\]
C’est l’entrée interne de la fonctionnelle hexadique--octadique. Son évaluation
comme paramètre de Barbero--Immirzi et son insertion dans l’opérateur d’aire sont
des publications ultérieures : elles ne participent pas à la sélection de
\(c_{12}^{+}\), de \(90/120\), de \(12:-1\) ni de \(1/24\).

\paragraph{Falsificateurs.}
La construction est réfutée si le recensement de phase ne donne pas douze secteurs, si
un tour contient plus d’une couture, si les incidences n’ont pas pour
cardinaux \(90\) et \(120\), si l’involution n’inverse pas la chaîne complète,
ou si le coefficient exact de \((1-q)^{-1}\) de son image diffère de \(1/24\).

\label{mdg:chap:barbero}

La transition d’une hexade à une octade combine trois structures de types
différents : une incidence finie, deux séries de Lambert associées aux
canaux \(q_\pm\) et une identification au paramètre réel de l’action de
Holst. Les trois s’articulent par l’incidence, la fonctionnelle angulaire et
le dictionnaire aréal vers la géométrie de la connexion
d’Ashtekar--Barbero développé dans cette même résidence.

\section{Cardinalités d’incidence 90 et 120}

Soit \(H\) un ensemble de six points et soit \(O\) un ensemble de huit points
contenant une copie distinguée de \(H\). Considérons
\[
 \mathcal F_6=H\times\binom H2,
 \qquad
 \mathcal F_8=O\times\binom H2.
\]
L’inclusion distinguée \(H\hookrightarrow O\) induit
\[
 \iota_{6,8}:\mathcal F_6\hookrightarrow\mathcal F_8,
 \qquad
 (h,\{h_1,h_2\})\longmapsto(h,\{h_1,h_2\}).
\]
Ce morphisme d’incidences est l’antécédent combinatoire des
cardinaux \(90/120\) ; ce n’est ni l’extracteur dodécaphasique ni une série angulaire.

\begin{theorem}[Cardinalités de l’inclusion]
\[
 |\mathcal F_6|=6\binom62=90,
 \qquad
 |\mathcal F_8|=8\binom62=120.
\]
La différence normalisée par vingt-quatre coordonnées est
\[
 \frac{90-120}{24\binom62}=-\frac1{12}.
\]
\end{theorem}
\begin{proof}
Chaque élément de la première coordonnée peut être combiné avec n’importe laquelle des
quinze paires de \(H\). Le principe multiplicatif donne les deux cardinalités ;
la dernière identité résulte de la substitution de \(\binom62=15\).
\end{proof}

\begin{figure}[htbp]
\centering
\begin{tikzpicture}[
  font=\small,
  setnode/.style={circle,draw,minimum size=6mm,inner sep=0pt,
    font=\scriptsize\bfseries},
  box/.style={draw,rounded corners=1.5mm,align=center,
    text width=5.0cm,inner xsep=6pt,inner ysep=5pt},
  arr/.style={-{Stealth[length=2.2mm]},line width=.7pt}]

  % Inclusión geométrica. Las flechas parten del contorno y no atraviesan nodos.
  \node[setnode,draw=hmtblue,fill=hmtlightblue] (h1) at (-2.3,1.25) {1};
  \node[setnode,draw=hmtblue,fill=hmtlightblue] (h2) at (-1.1,1.9) {2};
  \node[setnode,draw=hmtblue,fill=hmtlightblue] (h3) at (.1,1.25) {3};
  \node[setnode,draw=hmtblue,fill=hmtlightblue] (h4) at (.1,-.15) {4};
  \node[setnode,draw=hmtblue,fill=hmtlightblue] (h5) at (-1.1,-.8) {5};
  \node[setnode,draw=hmtblue,fill=hmtlightblue] (h6) at (-2.3,-.15) {6};
  \node[setnode,draw=hmtgreen,fill=hmtlightgreen] (o7) at (-3.2,.55) {7};
  \node[setnode,draw=hmtgreen,fill=hmtlightgreen] (o8) at (1.0,.55) {8};
  \node[draw=hmtblue,dashed,rounded corners=4mm,
    fit=(h1)(h2)(h3)(h4)(h5)(h6),inner sep=4mm,
    label={[hmtblue,font=\bfseries]above:{hexada \(H\), \(|H|=6\)}}] (H) {};
  \node[draw=hmtgreen,line width=.8pt,rounded corners=6mm,
    fit=(H)(o7)(o8),inner sep=5mm,
    label={[hmtgreen,font=\bfseries]below:{octada \(O\), \(|O|=8\), \(H\subset O\)}}] (O) {};

  \node[box,draw=hmtblue,fill=hmtlightblue] (f6) at (5.3,1.35)
    {Incidences définies sur \(H\)\\[1mm]
     \(\mathcal F_6=H\times\binom H2\),\quad
     \(|\mathcal F_6|=6\binom62=90\)};
  \node[box,draw=hmtgreen,fill=hmtlightgreen] (f8) at (5.3,-1.0)
    {Extension des incidences à \(O\)\\[1mm]
     \(\mathcal F_8=O\times\binom H2\),\quad
     \(|\mathcal F_8|=8\binom62=120\)};
  \draw[arr,hmtblue] (O.east) -- ++(.45,0) |- (f6.west);
  \draw[arr,hmtgreen] (O.east) -- ++(.45,0) |- (f8.west);
  \draw[arr,hmtgold] (f6.south) -- node[right,align=left,font=\footnotesize]
    {inclusion \(\iota_{6,8}\)} (f8.north);

  \node[box,draw=hmtgold,fill=hmtlightgold,text width=10.9cm]
    (def) at (1.45,-3.80)
    {Défaut combinatoire orienté\\[-.2mm]
     avec la normalisation déclarée :\\[1mm]
     \(\displaystyle
       \frac{|\mathcal F_6|-|\mathcal F_8|}{24\binom62}
       =\frac{90-120}{24\cdot15}=-\frac1{12}.\)};
  \draw[arr,hmtgold] (f8.south) -- (f8.south |- def.north);

  \node[align=center,text=hmtgray,font=\footnotesize,text width=11.2cm]
    at (1.45,-5.50)
    {Cette identité appartient à l’application d’incidences. Elle n’identifie pas
     l’extracteur dodécaphasique, le semi-groupe de tours ni la fonctionnelle angulaire.};
\end{tikzpicture}

\caption{Inclusion d’un sous-ensemble distingué de six éléments dans un ensemble de
huit éléments et dénombrements d’incidence. Les cardinaux \(90\) et \(120\) ne fusionnent ni l’extracteur
dodécaphasique, ni la fonctionnelle angulaire, ni le semi-groupe des tours.}
\label{mdg:fig:incidencia-seis-ocho-rev9}
\end{figure}

Les indices initiaux de la hiérarchie angulaire sont ainsi associés à deux
espaces d’incidence : 90 configurations sur l’hexade et 120 sur
l’octade ambiante. Le quotient \(-1/12\) est une identité combinatoire de cette
normalisation.

\section{Séries de Lambert associées aux 2 secteurs d’incidence}

Pour \(0<q<1\), nous définissons simultanément
\[
 S_{90}(q)=\frac{q^{90}}{1-q^{270}},
 \qquad
 S_{120}(q)=\frac{q^{120}}{1-q^{360}},
 \qquad
 \Delta S_m=S_m(q_-)-S_m(q_+).
\]
Les développements géométriques
\[
 \Delta S_{90}=\sum_{j\geq0}s_{90+270j},
 \qquad
 \Delta S_{120}=\sum_{j\geq0}s_{120+360j}
\]
situent les deux termes dans le semi-groupe global du
\cref{mdg:chap:torres}.

La fonctionnelle angulaire associée à la transition \(6\to8\) est
\[
 K_{6\to8}=12\Delta S_{90}-\Delta S_{120},
\]
et sa précédence a été fixée par la chaîne fondamentale dodécaphasique :
les douze secteurs et l’unique couture orientée produisent d’abord le poids
\(12:-1\). Le coefficient de \((1-q)^{-1}\), égal à \(1/24\), est ensuite calculé comme sortie exacte ;
l’équation diophantienne qui suit certifie la paire déjà engendrée et ne la
sélectionne pas.
La coordonnée angulaire résultante est
\[
\Gamma_{6\to8}
=\frac{\pi A\Cstar}{180}+K_{6\to8}.
\]
Au point HMT,
\[
\begin{aligned}
 \Delta S_{90}&=2.95169439297457621139847987861\times10^{-4},\\
 \Delta S_{120}&=1.96835112502951720093738706217\times10^{-5},\\
 \frac{\pi A\Cstar}{180}&=0.270485322934140078465586458790\ldots,
\end{aligned}
\]
et, par conséquent,
\[
 \boxed{\Gamma_{6\to8}
 =0.27400767269445927474725526077398041940\ldots.}
\]
Les tableaux calculés avec la coordonnée angulaire conjuguée \(C^*\) arrondie à quinze décimales donnent
\(0.2740076726944593\) en arithmétique à double précision ; la valeur précédente
est l’évaluation multiprécision de la phase régionale complète.

\section{Certificat diophantien ultérieur des poids}

L’identité
\[
 \lim_{q\to1^-}(1-q)\frac{q^m}{1-q^n}=\frac1n
\]
associe à la combinaison \(aS_{90}-bS_{120}\), avant son évaluation
dans les canaux conjugués, le coefficient de \((1-q)^{-1}\)
\[
 \frac a{270}-\frac b{360}.
\]

\begin{theorem}[Poids entiers de coefficient de pôle \texorpdfstring{\(1/24\)}{1/24}]
Les solutions entières positives de
\[
 \frac a{270}-\frac b{360}=\frac1{24}
\]
sont
\[
 (a,b)=(12+3t,1+4t),\qquad t\in\ZZ_{\geq0}.
\]
La paire \((12,1)\) est l’unique solution minimale coordonnée par coordonnée et
l’unique de norme \(\ell^1\) minimale.
\end{theorem}
\begin{proof}
La multiplication par \(1080\) produit \(4a-3b=45\). Modulo trois, on obtient
\(a\equiv0\pmod3\) ; en écrivant \(a=3u\), il en résulte \(4u-b=15\). La
positivité implique \(u=4+t\), \(b=1+4t\), avec \(t\geq0\). Par conséquent,
\(a+b=13+7t\), ce qui démontre la minimalité.
\end{proof}

La primitivité \(\gcd(a,b)=1\) ne sélectionne pas à elle seule la paire minimale, car
il existe d’autres solutions primitives. Le tour fondamental du TPK a déjà
produit le poids \(12:-1\) ; le coefficient de pôle \(1/24\), la positivité et la
minimalité forment son certificat arithmétique ultérieur. Aucun de ces
objets ne se déduit de l’évaluation décimale de \(\Gamma_{6\to8}\).

\newpage
\section{Monotonie de la fonctionnelle par rapport à la coordonnée angulaire conjuguée}

Para \(m\in\{90,120\}\),
\[
 S_m'(q)=\frac{m q^{m-1}(1+2q^{3m})}{(1-q^{3m})^2}>0.
\]
De plus,
\[
 \frac{\dd q_-}{\dd C}=\frac\pi{180}q_-,
 \qquad
 \frac{\dd q_+}{\dd C}=-\frac\pi{180}q_+.
\]

\begin{proposition}[Monotonie stricte]
Pour \(1^\circ\leq C\leq3^\circ\), la fonction
\(C\mapsto\Gamma_{6\to8}(A,C)\) est strictement croissante. Par conséquent,
toute valeur de son image possède au plus un antécédent dans cet intervalle.
\end{proposition}
\begin{proof}
La dérivée de chaque différence \(\Delta S_m\) est positive. Dans l’intervalle
indiqué, on obtient
\(\Delta S'_{120}<5.35\times10^{-4}\), tandis que
\(\pi A/180>0.1273\). Par conséquent,
\(\Gamma'(C)>0.1267\). En \(C=\Cstar\),
\[
 \Gamma'(\Cstar)=0.1328995105618907\ldots.
\]
\end{proof}

L’injectivité locale exclut une seconde coordonnée angulaire proche donnant la
même évaluation de la fonctionnelle. Ce contrôle ne construit ni ne sélectionne
\(C^*\) : la coordonnée angulaire conjuguée a déjà été obtenue dans la branche régionale précédente.
La proposition quantifie uniquement la sensibilité de l’évaluation descendante
une fois \(A\) et \(C^*\) fixés ; elle n’autorise pas à inverser
\(\Gamma_{6\to8}\) ni une valeur physique de \(\gamma_{\mathrm{BI}}\) pour
les redéfinir.

\label{ch:modular}

\section{Inclusion hexade--octade et réalisations fonctionnelles associées}

Fixons un ensemble \(H\) de six points et un ensemble \(O\) de huit
points avec une inclusion distinguée \(j:H\hookrightarrow O\). Dans la
résidence de Witt, \(H\) est une hexade d'une dodécade et \(O\) son octade
relevée ; dans ce chapitre, on utilise seulement l'inclusion déjà fixée, et non la
valeur terminale de Barbero--Immirzi. Nous écrivons
\begin{equation}
\label{eq:modular-pairs-H}
\binom H2:=\{P\subset H:|P|=2\},
\qquad \left|\binom H2\right|=\binom62=15,
\end{equation}
et définissons deux espaces d'incidences de types distincts :
\begin{equation}
\label{eq:modular-F6-F8}
\mathcal F_6:=H\times\binom H2,
\qquad
\mathcal F_8:=O\times\binom H2.
\end{equation}
L'inclusion d'incidences est
\begin{equation}
\label{eq:modular-iota68}
\iota_{6,8}:\mathcal F_6\hookrightarrow\mathcal F_8,
\qquad
(h,P)\longmapsto(j(h),P).
\end{equation}
Elle est injective parce que \(j\) l'est et conserve la paire marquée \(P\).

\begin{theorem}[Décomptes d'incidence et défaut orienté]
\label{thm:modular-incidence-90-120}
Les deux espaces de \cref{eq:modular-F6-F8} satisfont
\begin{equation}
\label{eq:modular-counts-90-120}
|\mathcal F_6|=6\binom62=90,
\qquad
|\mathcal F_8|=8\binom62=120,
\qquad
|\mathcal F_8\setminus\iota_{6,8}(\mathcal F_6)|=30.
\end{equation}
La normalisation par les vingt-quatre coordonnées du couple
dodécade--vingt-quatre et par les quinze paires de \(H\) produit
\begin{equation}
\label{eq:modular-incidence-defect}
\boxed{
\delta_{6,8}^{\rm inc}
=\frac{|\mathcal F_6|-|\mathcal F_8|}
{24\binom62}=-\frac1{12}.}
\end{equation}
La même fraction s'obtient par la fonctionnelle réciproque associée
à la mémoire incorporée,
\begin{equation}
\label{eq:modular-reciprocal-defect}
\boxed{
\delta_{6,8}^{\rm rec}
=\frac{30}{120}-\frac{30}{90}=-\frac1{12}.}
\end{equation}
Ce sont deux réalisations exactes du même défaut orienté, mais elles ont
des domaines distincts et ne sont pas identifiées à \(\zeta(-1)\) sans un
entrelaceur supplémentaire.
\end{theorem}

\begin{proof}
Le principe multiplicatif appliqué à
\cref{eq:modular-F6-F8} donne \(6\cdot15=90\) et \(8\cdot15=120\).
L'image de \(\iota_{6,8}\) a quatre-vingt-dix éléments, de sorte que le
complémentaire en a trente. Substituer dans
\cref{eq:modular-incidence-defect} donne
\(-30/(24\cdot15)=-1/12\), tandis que
\cref{eq:modular-reciprocal-defect} est
\(1/4-1/3=-1/12\).
\end{proof}

La transition hexade--octade admet deux réalisations fonctionnelles qui ne
se déduisent pas l'une de l'autre. Le diagramme algébrique, formulé sur les espaces
qui justifient les décomptes, est
\begin{equation}
\label{eq:modular-68-two-readers}
\begin{aligned}
(H\xhookrightarrow{\,j\,}O)
&\longmapsto
\left(|\mathcal F_6|,|\mathcal F_8|\right)=(90,120)
\longmapsto\Gamma_{6\to8}=\gammaH,\\
(h_6,o_8,t)=(6,8,3)
&\longmapsto-\frac{o_8-h_6}{t}=-\frac23,
\end{aligned}
\end{equation}
où \(-2/3\) est le coefficient de la troisième coordonnée de la ligne
\(13\) du constructeur CKM. La première réalisation transforme
l'incidence en aire et en mémoire modulaire ; la seconde transforme la même
transition en paramètres de mélange angulaire.
On ne postule pas \(\boldsymbol\theta_{\rm CKM}=f(\gammaH)\) : les deux sorties
conservent un antécédent commun.

Le morphisme qui détermine les secteurs \(90/120\) est
\(\iota_{6,8}\). La grandeur \(\gammaH\) a déjà été déterminée en interne
par la construction hexade--octade et est réutilisée ici sans être reconstruite à
partir de sa représentation décimale. Les
conséquences opératorielles ne commencent qu'après fixation de l'échelle
élémentaire \(a_0\), qui est définie ci-dessous sans utiliser l'entropie ni le
produit terminal \(\gammaH a_0\).

\section{Correspondance entre l'intérieur nonadique et les degrés de liberté de bord}

Soit \(\mathbb T_3=\{0,1,2\}\). Tout \(x\in\mathbb Z/9\mathbb Z\) possède
une écriture en chiffres unique
\(x=x_0+3x_1\), avec \(x_i\in\mathbb T_3\), et tout
\(y\in\mathbb Z/27\mathbb Z\) une écriture unique
\(y=y_0+3y_1+9y_2\). C'est pourquoi
\begin{equation}
\label{eq:modular-bulk-boundary-729}
\mathcal B=(\mathbb Z/9\mathbb Z)^3,
\qquad
\mathcal B_\partial=(\mathbb Z/27\mathbb Z)^2,
\qquad
|\mathcal B|=|\mathcal B_\partial|=3^6=729.
\end{equation}
L'égalité des cardinaux s'accompagne de l'application explicite. Si
\(x_j=r_{2j-1}+3r_{2j}\) pour \(j=1,2,3\), nous définissons
\begin{equation}
\label{eq:modular-beta729}
\beta_{729}(x_1,x_2,x_3)
=\bigl(r_1+3r_2+9r_3,\;r_4+3r_5+9r_6\bigr).
\end{equation}

\begin{proposition}[Bijection généalogique \(729\leftrightarrow729\)]
\label{prop:modular-beta729}
L'application \(\beta_{729}:\mathcal B\to\mathcal B_\partial\) est bijective et
préserve le mot ordonné de six trits. Ce n'est pas, en général, un
isomorphisme de groupes additifs : regrouper les chiffres change l'endroit où
la retenue est enregistrée.
\end{proposition}

\begin{proof}
Extraire les six chiffres de trois coordonnées nonadiques et les regrouper en
deux triplets est une permutation des coordonnées de \(\mathbb T_3^6\). La
décomposition positionnelle est unique des deux côtés, de sorte que
l'opération possède un inverse. La mise en garde additive se vérifie, par exemple,
en additionnant des mots qui engendrent une retenue entre \(r_3\) et \(r_4\) : cette
frontière sépare les deux coordonnées d'ordre vingt-sept, mais traverse
deux coordonnées distinctes dans le regroupement nonadique.
\end{proof}

Sur le tore de bord \(\mathcal B_\partial\), nous prenons le bloc de sommets
\begin{equation}
\label{eq:modular-block-A}
A=\{0,1,\ldots,8\}\times\{0,1,\ldots,8\}.
\end{equation}
Ses liens orientés qui traversent le bord se séparent par direction :
\begin{align}
\partial A_{90}
&=\{((-1,j),(0,j)),((8,j),(9,j)):0\leq j\leq8\},
\label{eq:modular-boundary90}\\
\partial A_{120}
&=\{((i,-1),(i,0)),((i,8),(i,9)):0\leq i\leq8\},
\label{eq:modular-boundary120}
\end{align}
où toutes les coordonnées se lisent modulo vingt-sept. Ainsi
\begin{equation}
\label{eq:modular-boundary-counts}
\partial A=\partial A_{90}\sqcup\partial A_{120},
\qquad
|\partial A_{90}|=|\partial A_{120}|=18.
\end{equation}
Les indices \(90,120\) étiquettent la provenance du canal ; ils ne disent pas que
le bord possède quatre-vingt-dix ou cent vingt liens.

\begin{theorem}[Loi aréale de la coupe triadique]
\label{thm:modular-triadic-cut}
Dans le graphe cubique \(N\times N\times N\) à capacité unitaire, toute
coupe entre deux faces opposées possède une capacité d'au moins \(N^2\), et une
couche transversale atteint cette borne. Pour \(N=9\),
\begin{equation}
\label{eq:modular-mincut-entropy}
\mathcal A_{\min}=81,
\qquad
S_{\rm tri}=81\log3.
\end{equation}
\end{theorem}

\begin{proof}
Il existe \(N^2\) droites d'arêtes parallèles à la direction normale, deux à deux
disjointes. Toute coupe doit intercepter chacune d'elles, de sorte que sa capacité
est d'au moins \(N^2\). Une seule couche transversale contient exactement
\(N^2\) arêtes et réalise le minimum. Si chaque lien coupé porte un trit
maximalement mélangé, son entropie est \(\log3\) ; l'additivité sur les
quatre-vingt-un liens donne \cref{eq:modular-mincut-entropy}.
\end{proof}

L'entropie \(S_{\rm tri}\) mesure la capacité de la coupe cubique.
L'entropie modulaire des trente-six liens de
\cref{eq:modular-boundary-counts} mesure un autre état et ne doit pas lui être
égalisée.

\section{Entrelaceur collectif d'incidence entre intérieur et bord}

Dans la somme
\(\ell^2(\mathcal F_6)\oplus\ell^2(\mathcal F_8)\), nous définissons
\begin{equation}
\label{eq:modular-u90-u120}
u_{90}=\frac1{\sqrt{90}}\sum_{f\in\mathcal F_6}(\delta_f,0),
\qquad
u_{120}=\frac1{\sqrt{120}}\sum_{f\in\mathcal F_8}(0,\delta_f).
\end{equation}
Dans
\(\ell^2(\partial A_{90})\oplus\ell^2(\partial A_{120})\), soient
\begin{equation}
\label{eq:modular-v90-v120}
v_{90}=\frac1{\sqrt{18}}\sum_{e\in\partial A_{90}}(\delta_e,0),
\qquad
v_{120}=\frac1{\sqrt{18}}\sum_{e\in\partial A_{120}}(0,\delta_e).
\end{equation}
Les couples \((u_{90},u_{120})\) et \((v_{90},v_{120})\) sont orthonormaux.
Il existe donc une unique isométrie entre les sous-espaces collectifs
qui satisfait
\begin{equation}
\label{eq:modular-J68}
\mathcal J_{6\to8}u_{90}=v_{90},
\qquad
\mathcal J_{6\to8}u_{120}=v_{120}.
\end{equation}
Définissons
\begin{align}
D_{\rm inc}
&=90|u_{90}\rangle\langle u_{90}|
+120|u_{120}\rangle\langle u_{120}|,
\label{eq:modular-Dinc}\\
D_{\rm area}
&=90|v_{90}\rangle\langle v_{90}|
+120|v_{120}\rangle\langle v_{120}|.
\label{eq:modular-Darea}
\end{align}

\begin{theorem}[Transport exact des deux étiquettes]
\label{thm:modular-J68}
L'application \(\mathcal J_{6\to8}\) est unitaire entre les deux plans
collectifs et entrelace les opérateurs d'incidence et d'aire :
\begin{equation}
\label{eq:modular-J-intertwining}
\boxed{
\mathcal J_{6\to8}D_{\rm inc}
=D_{\rm area}\mathcal J_{6\to8}.}
\end{equation}
Par conséquent, toute combinaison polynomiale des étiquettes \(90,120\), en
particulier la combinaison primitive \(12:-1\), est transportée avec les
mêmes coefficients.
\end{theorem}

\begin{proof}
L'orthonormalité montre que \(\mathcal J_{6\to8}\) envoie une base
orthonormale sur une autre. En \(u_{90}\), les deux membres de
\cref{eq:modular-J-intertwining} valent \(90v_{90}\) ; en \(u_{120}\),
ils valent tous deux \(120v_{120}\). L'égalité sur une base prouve l'identité.
Le calcul fonctionnel polynomial conserve l'entrelacement.
\end{proof}

Cette application ne prétend pas injecter quatre-vingt-dix incidences individuelles dans
dix-huit liens : elle transporte exactement les deux modes uniformes et leurs
étiquettes spectrales. L'étendre à des fluctuations non uniformes exigerait
une nouvelle application d'incidence ; l'opérateur modulaire défini ici n'utilise que
le sous-espace collectif qui vient d'être construit.

\section{Détermination de l'échelle élémentaire du spectre de l'opérateur d'aire}

Soit \(C^*\) la coordonnée angulaire conjuguée, exprimée en degrés, et fixons
\begin{equation}
\label{eq:modular-theta-clock-a0}
\theta_{\rm clk}=\frac{C^*}{6}\frac{\pi}{180},
\qquad
\boxed{
a_0=\frac{1+2\cos(10^\circ)}{3}\,\theta_{\rm clk}.}
\end{equation}
Le facteur qui précède \(\theta_{\rm clk}\) n'est pas une correction
décimale. C'est le caractère réel normalisé du triplet de phases
\begin{equation}
\label{eq:modular-c10-character}
\frac13\Tr\operatorname{diag}
\left(1,\ee^{\ii\pi/18},\ee^{-\ii\pi/18}\right)
=\frac{1+2\cos(10^\circ)}3.
\end{equation}
Ainsi, \(a_0\) est la moyenne spectrale de la structure angulaire sur le trit
orienté. Elle est positive, est fixée avant l'opérateur d'aire et n'est pas
sélectionnée à partir de l'entropie. L'identification ultérieure
\(\gammaH=\Gamma_{6\to8}\) réalise la sortie HMT dans la connexion
d'Ashtekar--Barbero \cite{Barbero,Immirzi} ; elle ne rouvre pas sa dérivation.

\begin{remark}[Portée du symbole \(a_0\)]
Dans ce chapitre, \(a_0\equiv a_0^{\rm area}\) désigne exclusivement la
normalisation de \cref{eq:modular-theta-clock-a0}. Ce n'est ni le rayon de Bohr
ni le facteur d'échelle cosmologique initial, bien que ces objets s'écrivent
traditionnellement avec le même symbole. Aucune identité de ce chapitre
ne permet de substituer l'un à l'autre.
\end{remark}

\section{Loi aréale de la coupe triadique : capacité \(81\log 3\), opérateur d’aire, spectre et entrelaceur collectif des canaux \(90/120\)}
\label{hap:sec:area-entrelazamiento-barbero-rev11}
\index{Barbero--Immirzi!opérateur d’aire}
\index{entropie modulaire!Hamiltonien modulaire}

L’identification HMT de la fonctionnelle \(\Gamma_{6\to8}\) à la coordonnée
de Barbero--Immirzi s’effectue sur le même bord triadique qui porte
l’incidence (90/120). L’intérieur et son feuillet de bord ont la même cardinalité,
\[
 \mathcal B=(\mathbb Z/9\mathbb Z)^3,
 \quad |\mathcal B|=729,
 \qquad
 \partial\mathcal B=(\mathbb Z/27\mathbb Z)^2,
 \quad |\partial\mathcal B|=729,
\]
par le regroupement de six trits en trois coordonnées nonadiques ou deux
coordonnées d’ordre vingt-sept.

\begin{theorem}[Loi aréale de la coupe triadique]
\label{hap:thm:ley-areal-corte-triadico-rev11}
Dans le graphe cubique \(N\times N\times N\) de capacité unitaire, la coupe
minimale entre deux faces opposées a pour capacité \(N^2\). Pour \(N=9\),
\[
 \mathcal A_{\min}=81,
 \qquad S_{\rm tri}=\mathcal A_{\min}\log3=81\log3.
\]
\end{theorem}
\begin{proof}
Les \(N^2\) droites parallèles à la direction normale sont des chemins disjoints en
arêtes, de sorte que toute coupe les intercepte. Une couche transversale contient
exactement \(N^2\) arêtes et atteint la borne. Chaque liaison triadique apporte
\(\log3\) à l’entropie d’un état maximalement mélangé.
\end{proof}

Dans le tore de bord \(27\times27\), soit \(A\) le bloc canonique
\(9\times9\). Sa frontière se décompose exactement en
\begin{equation}
 \partial A=\partial A_{90}\sqcup\partial A_{120},
 \qquad |\partial A_{90}|=|\partial A_{120}|=18.
 \label{hap:eq:corte-90-120-rev11}
\end{equation}
Les liaisons horizontales représentent le canal (90), et les verticales le
canal (120). Cette décomposition transforme la paire d’incidences en une
décomposition opératorielle du bord.

L’entropie aréale \(S_{\rm tri}=81\log3\) appartient à l’état maximalement
mélangé des liaisons de la coupe. L’entropie bicouche
\(S_{\rm bi}(y)\) de la \cref{hap:def:entropia-bicapa-rev6} appartient à la
distribution normalisée entre les deux feuillets ; les deux objets conservent
des domaines distincts.

Soit \(N=\operatorname{diag}(0,1,2)\) l’opérateur nombre d’une liaison tritique.
La réalisation aréale déclare l’échelle positive
\[
 \theta_{\rm clk}=\frac{C^*}{6}\frac{\pi}{180},
 \qquad
 a_0=\frac{1+2\cos(10^\circ)}{3}\,\theta_{\rm clk},
 \qquad \gamma_{\rm HMT}=\Gamma_{6\to8}.
\]
La formule de \(a_0\) constitue la structure de départ de cette
réalisation ; une fois cette échelle fixée, l’opérateur et son spectre se déduisent
exactement.

\begin{definition}[Opérateur d’aire HMT]
Sur le feuillet de bord
\(\mathcal H_{\partial}=\bigotimes_{e\in\partial A}\mathbb C^3\), on définit
\begin{equation}
 \widehat{\mathcal A}_{\rm HMT}
 =\gamma_{\rm HMT}a_0\sum_{e\in\partial A}N_e.
 \label{hap:eq:operador-area-hmt-rev11}
\end{equation}
Son spectre dans la coupe canonique est
\[
 \operatorname{Spec}\widehat{\mathcal A}_{\rm HMT}
 =\{n\gamma_{\rm HMT}a_0:0\leq n\leq72\},
\]
avec des multiplicités données par les coefficients de
\((1+z+z^2)^{36}\). Le prolongement par coupes compatibles produit la
famille protégée \(\mathcal A_n^{\rm prot}=n\gamma_{\rm HMT}a_0\).
\end{definition}

L’entrelaceur qui fixe la normalisation se construit dans le sous-espace des
canaux collectifs. Soient \(u_{90},u_{120}\) les vecteurs unitaires uniformes
des incidences hexadique et octadique, et soient \(v_{90},v_{120}\) les
vecteurs unitaires uniformes des deux ensembles de liaisons de
\eqref{hap:eq:corte-90-120-rev11}. L’application
\begin{equation}
 \begin{aligned}
 \mathcal J_{6\to8}:{}
 &\operatorname{span}\{u_{90},u_{120}\}
 \longrightarrow
 \operatorname{span}\{v_{90},v_{120}\},\\
 &\mathcal J_{6\to8}u_{90}=v_{90},
 \qquad
 \mathcal J_{6\to8}u_{120}=v_{120}
 \end{aligned}
 \label{hap:eq:entrelazador-area-barbero-rev11}
\end{equation}
est une isométrie entre ces sous-espaces collectifs. Si
\[
 D_{\rm inc}=90|u_{90}\rangle\langle u_{90}|
 +120|u_{120}\rangle\langle u_{120}|,
 \qquad
 D_{\rm area}=90|v_{90}\rangle\langle v_{90}|
 +120|v_{120}\rangle\langle v_{120}|,
\]
alors
\[
 \boxed{\mathcal J_{6\to8}D_{\rm inc}
 =D_{\rm area}\mathcal J_{6\to8}.}
\]
En particulier, la combinaison primitive \(12:-1\) agit
avec les mêmes coefficients avant et après \(\mathcal J_{6\to8}\). Cette
égalité fixe le dictionnaire de normalisation qui insère
\(\gamma_{\rm HMT}=\Gamma_{6\to8}\) dans la connexion
d’Ashtekar--Barbero.


Le passage de l’hexade à l’octade est une transformation de la structure d’incidence.\par

Lors du passage de \(6\) à \(8\), les paires disponibles, les relations internes et la manière dont l’orientation peut se conserver changent. Cette transition engendre un défaut orienté. La variation cardinale réalise une rotation du mode d’organisation.\par

La même inversion de feuillet s’exprime par\par

\[
C^*\longmapsto -C^*.
\]

Mais des lecteurs différents répondent de manières différentes à cette inversion.\par

La fonctionnelle hexade–octade qui conduit au paramètre de Barbero–Immirzi est impaire : lorsque le feuillet est inversé, elle change de signe. L’orientation géométrique s’inverse.\par

L’entropie bicouche, en revanche, est paire : elle attribue la même entropie à \(C^*\) et à \(-C^*\) lorsque l’on observe la distribution de probabilités. L’information totale peut rester identique même si l’orientation a été inversée.\par

L’espace CKM réalise une troisième réponse. L’inversion induit sur les trois angles une réflexion rationnelle exacte :\par

\[
R_{\rm CKM}
=
M_{\rm CKM}
\operatorname{diag}(1,-1,1)
M_{\rm CKM}^{-1},
\]

con\par

\[
R_{\rm CKM}^2=I,
\qquad
\det R_{\rm CKM}=-1.
\]

Cela signifie que le changement de feuillet est transporté dans l’espace de mélange comme une véritable réflexion. Il existe une direction impaire de rang un et un plan qui reste fixe.\par

Le contre-angle \(C^*\) est précisément la coordonnée qui change. La phase construite à partir de \(A\) reste invariante. La transformation sépare ainsi deux mémoires :\par

\begin{itemize}[leftmargin=*,itemsep=0.35em]
\item une mémoire d’orientation ;
\item une mémoire de phase.
\end{itemize}

Cela offre une lecture très profonde de CKM. Les angles de mélange sont les coordonnées d’un espace sur lequel agit la même involution qui organisait déjà la transition hexade–octade.\par

La rotation \(6\to8\) fournit l’opération d’orientation que CKM réalise dans son propre codomaine ; la causalité physique appartient au transducteur correspondant.\par

Nous avons alors trois publications sœurs du même passage :\par

\begin{enumerate}[leftmargin=*,itemsep=0.65em]
\item \textbf{Publication géométrique.}
Le paramètre de Barbero–Immirzi oriente le quantum d’aire.\par

\item \textbf{Publication informationnelle.}
L’état réduit et son hamiltonien modulaire mesurent quelle information de cette orientation demeure au bord.\par

\item \textbf{Publication de mélange.}
La réflexion CKM transforme les coordonnées de mélange et isole leur composante impaire.\par
\end{enumerate}

\(\gamma_{\rm BI}\), l’entropie et les angles CKM demeurent des réalisations distinctes d’une charnière généalogique commune. L’unité réside dans l’antécédent antérieur à la ramification. Ils partagent la charnière, mais chacun en conserve une partie différente.\par

C’est l’une des plus belles expressions de la méthodologie HMT : l’unité fondamentale est retrouvée en remontant les généalogies jusqu’à trouver l’opération commune qui engendre des résultats terminaux différents.\par

La gravité conserve la rotation comme orientation de l’aire.\par
L’information conserve la rotation comme mémoire de l’état.\par
Le mélange conserve la rotation comme réflexion des coordonnées.\par

La même transformation « respire » de trois manières.\par


\section{Séparation causale entre le spectre d’aire, la mémoire modulaire et la réalisation Cartan--Holst}
La dérivation de $\gamma_{\mathrm{BI}}$, de l’aire minimale et du spectre de l’opérateur d’aire s’achève dans ce chapitre. La reconstruction conjointe de l’occupation et de la provenance par le hamiltonien modulaire de bord, et la réalisation Palatini--Cartan--Holst, conservent leurs domaines, opérateurs et preuves propres dans la Partie IV.
%
\endgroup
% Secciones 4--7 del delta gravitatorio íntegro.
\section{Incidence hexade--octade, paramètre de Barbero--Immirzi et spectre d'aire}
\label{sec:l9s102:barbero-area}

La dérivation surfacique reçoit comme antécédents l'inclusion d'incidence
\(6\to8\), les cardinaux TPK \(90/120\), le défaut normalisé
\(-1/12\) et les canaux spectraux \(q_\pm\). La fonctionnelle
hexade--octade détermine le paramètre de Barbero--Immirzi et la
normalisation de l'opérateur d'aire. Cet emplacement contient la dérivation
HMT et le spectre ; la réalisation Palatini--Cartan--Holst occupe un emplacement
ultérieur dans la mécanique dimensionnelle.

Soient \(H\subset O\), avec \(|H|=6\) et \(|O|=8\), et conservons le couple
marqué de l'hexade dans l'inclusion
\begin{equation}
 \iota_{6,8}:\mathcal F_6=H\times\binom H2
 \lhook\joinrel\longrightarrow
 \mathcal F_8=O\times\binom H2.
 \label{eq:l9s102:barbero-inclusion-incidencia}
\end{equation}
Le comptage interne donne
\begin{equation}
 |\mathcal F_6|=6\binom62=90,
 \qquad |\mathcal F_8|=8\binom62=120,
 \qquad
 \frac{90-120}{24\binom62}=-\frac1{12}.
 \label{eq:l9s102:barbero-conteos}
\end{equation}
Une fois \(\pi_{\rm HMT}\), l'angle \(A\) et le
contre-angle unique \(C^*\) publiés, les deux canaux sont
\begin{equation}
 q_\pm=\exp\!\left[-\frac{\pi_{\rm HMT}}{180}(A\pm C^*)\right],
 \qquad 0<q_+<q_-<1,
 \label{eq:l9s102:barbero-qpm}
\end{equation}
et leurs séries de retour orienté sont
\begin{equation}
 S_{90}(q)=\frac{q^{90}}{1-q^{270}},\qquad
 S_{120}(q)=\frac{q^{120}}{1-q^{360}},\qquad
 \Delta S_m=S_m(q_-)-S_m(q_+).
 \label{eq:l9s102:barbero-series}
\end{equation}

L'ordre causal des coefficients est fixé dans le registre temporel du TPK,
avant d'évaluer le coefficient du pôle. Soit \(C_{12}=\mathbb Z/12\mathbb Z\) la roue
dodécaphasique et soit \(\partial_{\rm ret}\) son unique couture orientée de retour.
Dans le module libre des événements d'incidence, le tour fondamental est
\begin{equation}
 c_{12}^{+}
 =\sum_{j\in C_{12}}[j,\mathcal F_6]
  -[\partial_{\rm ret},\mathcal F_8].
 \label{eq:l9s102:barbero-cadena-dodecafasica}
\end{equation}
Le signe du second terme est le signe de bord. L'involution
bidirectionnelle inverse l'orientation complète,
\(c_{12}^{-}=-c_{12}^{+}\), sans modifier les multiplicités absolues des
deux types d'incidence. Définissons le lecteur de retour sur les
générateurs par
\begin{equation}
 \mathscr R([j,\mathcal F_6])=S_{90},
 \qquad
 \mathscr R([\partial_{\rm ret},\mathcal F_8])=S_{120}.
 \label{eq:l9s102:barbero-lector-retorno}
\end{equation}
La covariance de phase rend la première image constante sur les douze
secteurs ; la seconde agit uniquement sur la couture élargie. Par linéarité,
\begin{equation}
 \boxed{\mathscr R(c_{12}^{+})=12S_{90}-S_{120}.}
 \label{eq:l9s102:barbero-peso-generado}
\end{equation}
Ainsi, le poids \(12:-1\) est l'image du tour fondamental : douze secteurs et
un bord orienté. Il ne provient ni d'une imposition préalable de \(1/24\), ni d'une
comparaison décimale, ni de la valeur conventionnelle de Barbero--Immirzi.

\Needspace{12\baselineskip}
\begin{theorem}[Détermination dodécaphasique de la fonctionnelle de Barbero--Immirzi]
\label{thm:l9s102:barbero-funcional}
La chaîne fondamentale \eqref{eq:l9s102:barbero-cadena-dodecafasica} détermine
le couple orienté \((a,b)=(12,1)\). Le coefficient de \((1-q)^{-1}\) de son image est
alors \(1/24\), comme sortie exacte de la fonctionnelle déjà construite. Par conséquent,
après évaluation des canaux conjugués \(q_\pm\), la sortie adimensionnelle est
\begin{equation}
 \boxed{
 \gamma_{\rm HMT}=\Gamma_{6\to8}
 =\frac{\pi_{\rm HMT}AC^*}{180}
   +12\Delta S_{90}-\Delta S_{120}.}
 \label{eq:l9s102:barbero-gamma}
\end{equation}
Cette formule reçoit le couple angulaire déjà généré et n'utilise pas en entrée
la valeur conventionnelle du paramètre de Barbero--Immirzi.
\end{theorem}

\begin{proof}
La somme de \eqref{eq:l9s102:barbero-cadena-dodecafasica} contient exactement
les douze secteurs de \(C_{12}\), tandis que la couture de retour apparaît une
seule fois et avec une orientation négative. Le lecteur
\eqref{eq:l9s102:barbero-lector-retorno} produit donc
\eqref{eq:l9s102:barbero-peso-generado}, avant tout calcul de pôle. Pour
chaque \(m>0\),
\[
 p_1(S_m)
 =\lim_{q\to1}(1-q)\frac{q^m}{1-q^{3m}}
 =\frac1{3m}.
\]
En conséquence,
\[
 p_1\bigl(\mathscr R(c_{12}^{+})\bigr)
 =\frac{12}{270}-\frac1{360}
 =\frac{48-3}{1080}=\frac1{24}.
\]
Le coefficient \(p_1\) est rapporté à \((1-q)^{-1}\).
Dans la coordonnée complexe \(q-1\),
\[
 \operatorname*{Res}_{q=1}S_m(q)=-\frac1{3m},
 \qquad
 \operatorname*{Res}_{q=1}\mathscr R(c_{12}^{+})(q)=-\frac1{24}.
\]
L'équation diophantienne ultérieure
\[
 \frac a{270}-\frac b{360}=\frac1{24}
 \quad\Longleftrightarrow\quad 4a-3b=45
\]
possède les solutions entières positives
\((a,b)=(12+3t,1+4t)\), \(t\ge0\) ; l'unique solution minimale est
\((12,1)\). Ce calcul certifie arithmétiquement le couple que l'orbite
fondamentale avait déjà engendré : tout \(t>0\) altère la multiplicité des
secteurs ou de la couture et cesse de représenter un tour du TPK.
Enfin, substituer \(q_\pm\) dans l'image orientée et ajouter la composante
angulaire préalablement générée prouve \eqref{eq:l9s102:barbero-gamma}.
\end{proof}

La séquence causale est fixée sans échange des rôles :
\begin{equation}
\begin{aligned}
 (\mathcal F_6\hookrightarrow\mathcal F_8,C_{12},\partial_{\rm ret})
 &\longrightarrow (90,120,-1/12,c_{12}^{+})\\
 &\longrightarrow \mathscr R(c_{12}^{+})=12S_{90}-S_{120},\\
 \mathscr R(c_{12}^{+})
 &\longrightarrow p_1\bigl(\mathscr R(c_{12}^{+})\bigr)=\frac1{24},\\
 \bigl(\mathscr R(c_{12}^{+}),A,C^*,q_\pm\bigr)
 &\longrightarrow \Gamma_{6\to8}=\gamma_{\rm HMT}
 \longrightarrow \widehat{\mathcal A}_{\rm phys}.
\end{aligned}
 \label{eq:l9s102:barbero-cadena-completa}
\end{equation}
Le couple angulaire et ses canaux proviennent de la construction précédente.
L'évaluation du coefficient du pôle constitue une vérification
de la fonctionnelle de retour ; l'évaluation conjuguée emploie les
coordonnées angulaires préalablement générées.
La réalisation ultérieure dans l'action de Holst reconnaît la même section
adimensionnelle ; elle n'intervient pas dans sa sélection.

Pour les deux secteurs de bord, soient
\begin{equation}
 N_m=\sum_{e\in\partial A_m}N_e,
 \qquad m\in\{90,120\},
 \qquad \operatorname{Spec}(N_e)=\{0,1,2\}.
 \label{eq:l9s102:n90-n120}
\end{equation}
L'opérateur physique d'aire s'écrit
\begin{equation}
 \boxed{
 \frac{\widehat{\mathcal A}_{\mathrm{phys}}}{\ell_P^2}
 =\gamma_{\mathrm{HMT}}a_0(N_{90}+N_{120}).}
 \label{eq:l9s102:area-operador}
\end{equation}
Sur \(36\) liens qutrit,
\begin{equation}
 \operatorname{Spec}(N_{90}+N_{120})=\{0,1,\ldots,72\},
 \label{eq:l9s102:area-numero-spectrum}
\end{equation}
et la multiplicité de la valeur propre \(n\) est le coefficient de \(z^n\) dans
\((1+z+z^2)^{36}\). Il en résulte
\begin{equation}
 \operatorname{Spec}\!\left(
 \frac{\widehat{\mathcal A}_{\mathrm{phys}}}{\ell_P^2}\right)
 =\{n\gamma_{\mathrm{HMT}}a_0:0\le n\le72\}.
 \label{eq:l9s102:area-spectrum}
\end{equation}

\section{Géométrie totale et mémoire de provenance}
\label{sec:l9s102:area-hamiltoniano}

Le hamiltonien modulaire du bord est
\begin{equation}
 \boxed{
 K_A=-\log\rho_A
 =cI+\Delta_{\mathrm{mod}}(90N_{90}+120N_{120}).}
 \label{eq:l9s102:hamiltoniano-modular}
\end{equation}
Soient
\begin{equation}
 N=N_{90}+N_{120},
 \qquad
 D=90N_{90}+120N_{120}.
 \label{eq:l9s102:n-d}
\end{equation}
La matrice de passage
\(\left(\begin{smallmatrix}1&1\\90&120\end{smallmatrix}\right)\)
a pour déterminant \(30\), et son inverse reconstruit les secteurs :
\begin{equation}
 \boxed{
 N_{90}=4N-\frac{D}{30},
 \qquad
 N_{120}=\frac{D}{30}-3N.}
 \label{eq:l9s102:reconstruccion-sectores}
\end{equation}
L'aire publie le nombre total d'excitations ; le hamiltonien modulaire
publie leur secteur de provenance. Le couple conjoint reconstruit l'occupation
complète du bord.

\section{État bicouche et information orientée}
\label{sec:l9s102:estado-bicapa}

Soit \(R=R^*=R^{-1}\) l'involution des feuillets et
\(y=C^*_{\mathrm{rad}}\). L'état normalisé est
\begin{equation}
 \boxed{
 \rho_y=\frac{e^{-yR}}{2\cosh y}.}
 \label{eq:l9s102:rho-y}
\end{equation}
Son hamiltonien modulaire et son entropie sont
\begin{equation}
 K_y=-\log\rho_y=\log(2\cosh y)I+yR,
 \label{eq:l9s102:ky}
\end{equation}
\begin{equation}
 S(\rho_y)=\log(2\cosh y)-y\tanh y.
 \label{eq:l9s102:entropia-y}
\end{equation}
L'inversion \(y\mapsto-y\) laisse l'entropie fixe et produit la distance
informationnelle orientée
\begin{equation}
 \boxed{
 D(\rho_y\Vert\rho_{-y})=2y\tanh y.}
 \label{eq:l9s102:entropia-relativa-hojas}
\end{equation}
La composante paire mesure la distribution spectrale ; la composante impaire mesure
le coût informationnel de l'inversion de l'orientation.

\section{Quantum informationnel d'énergie et cellule gravitationnelle de demi-action}
\label{sec:l9s102:horizonte-media-accion}

La périodicité de \(108\) états détermine
\begin{equation}
 t_P=\frac{54}{\pi}t_0,
 \qquad
 E_P=\frac{\hbar}{t_P}=\frac{h}{108t_0}.
 \label{eq:l9s102:ep-tp}
\end{equation}
La réalisation thermique de l'horloge tritique satisfait
\begin{equation}
 \boxed{
 E_P=k_B\Theta_{\mathrm{clk}}\log3.}
 \label{eq:l9s102:planck-trit}
\end{equation}
Pour l'horizon de rayon \(R\),
\begin{equation}
 E_\Lambda(R)=\frac{c^4R}{2G},
 \qquad
 E_\Lambda(R)=T_H(R)S_H(R).
 \label{eq:l9s102:energia-horizonte}
\end{equation}
Dans la cellule de Planck,
\begin{equation}
 \boxed{
 E_\Lambda
 =T_HS_H
 =\frac{E_P}{2}
 =\frac{k_B\Theta_{\mathrm{clk}}\log3}{2}.}
 \label{eq:l9s102:confluencia-horizonte}
\end{equation}
Par conséquent,
\begin{equation}
 \boxed{
 E_\Lambda t_P=\frac{\hbar}{2},
 \qquad
 E_\Lambda T_{108}=\frac h2.}
 \label{eq:l9s102:media-accion}
\end{equation}
L'énergie du vide, l'entropie surfacique, la température chronologique et
l'action partagent une même cellule de confluence avec des domaines typés.

Lorsque \(S_H/k_B=\pi\), la multiplicité effective satisfait
\begin{equation}
 \Omega_{\mathrm{eff}}=e^{S_H/k_B}=e^\pi.
 \label{eq:l9s102:omega-e-pi}
\end{equation}
L'opérateur hyperbolique d'Euler possède
\begin{equation}
 \operatorname{Spec}(e^{\pi H})=\{e^\pi,e^{-\pi}\}.
 \label{eq:l9s102:e-pi-spectrum}
\end{equation}
L'égalité des scalaires établit la confluence spectrale ;
l'opérateur d'entrelacement entre multiplicité microscopique d'horizon et dynamique de
Perron relève de son emplacement physique spécifique.


\HMTFlushCurrentChapter
% Capítulo 52 del sucesor: integración pública del propietario íntegro del radión.
% No carga portada, main, style ni C_estatutos_procedencia del módulo autónomo.

% Adaptador de compatibilidad del propietario íntegro del radión.
%
% El tratado rector ya carga tipografía, geometría, hipervínculos, teoremas,
% tablas, TikZ y tcolorbox.  Este archivo aporta únicamente los nombres y
% entornos exclusivos del propietario, sin importar su preámbulo autónomo ni
% alterar la jerarquía tipográfica general.

\makeatletter
\@ifundefined{color@RadNavy}{\definecolor{RadNavy}{HTML}{102A43}}{}
\@ifundefined{color@RadBlue}{\definecolor{RadBlue}{HTML}{1F5E7A}}{}
\@ifundefined{color@RadTeal}{\definecolor{RadTeal}{HTML}{168C8C}}{}
\@ifundefined{color@RadCopper}{\definecolor{RadCopper}{HTML}{B66A3C}}{}
\@ifundefined{color@RadGold}{\definecolor{RadGold}{HTML}{D4A73A}}{}
\@ifundefined{color@RadInk}{\definecolor{RadInk}{HTML}{1A202C}}{}
\@ifundefined{color@RadMist}{\definecolor{RadMist}{HTML}{EDF4F7}}{}
\@ifundefined{color@RadCream}{\definecolor{RadCream}{HTML}{FAF7F0}}{}
\@ifundefined{color@RadRose}{\definecolor{RadRose}{HTML}{8D3F55}}{}
\@ifundefined{color@RadGreen}{\definecolor{RadGreen}{HTML}{2F7D5D}}{}
\@ifundefined{color@RadGray}{\definecolor{RadGray}{HTML}{66788A}}{}

\providecommand{\TRIT}{\ensuremath{\mathrm{TRIT}}}
\providecommand{\HMT}{\ensuremath{\mathrm{HMT}}}
\providecommand{\Rstate}{\mathcal R_{12}}
\providecommand{\epsR}{\varepsilon_R}
\providecommand{\epsone}{\varepsilon_R^{(1)}}
\providecommand{\pih}{\pi_{\!\HMT}}
\providecommand{\alphah}{\alpha_{\!\HMT}}
\providecommand{\hret}{\hbar_{\mathrm{ret}}}
\providecommand{\hfull}{h_{\mathrm{ret}}}
\providecommand{\diff}{\mathop{}\!\mathrm d}
\providecommand{\e}{\mathrm e}
\providecommand{\R}{\mathbb R}
\providecommand{\C}{\mathbb C}
\providecommand{\F}{\mathbb F}
\providecommand{\Z}{\mathbb Z}
\providecommand{\dd}{\,\mathrm d}
\providecommand{\status}[1]{\textsf{\bfseries #1}}
\providecommand{\sourcepath}[1]{\nolinkurl{#1}}

\@ifundefined{typebox}{%
  \newtcolorbox{typebox}[1][]{enhanced,breakable,colback=white,
    colframe=RadBlue,boxrule=.7pt,arc=1.5mm,left=4mm,right=4mm,
    top=3mm,bottom=3mm,title={Reçu de types},fonttitle=\bfseries,
    coltitle=RadNavy,#1}%
}{}
\@ifundefined{narrativebox}{%
  \newtcolorbox{narrativebox}[1][]{enhanced,breakable,colback=white,
    colframe=RadGold,borderline west={2.4pt}{0pt}{RadGold},
    boxrule=.25pt,arc=0mm,left=5mm,right=4mm,top=3mm,bottom=3mm,#1}%
}{}

% Las once portadillas internas del propietario autónomo se convierten en
% divisiones científicas ordinarias dentro del capítulo 52.  Se evita así
% introducir partes autónomas, gigantismo tipográfico o páginas vacías.
\providecommand{\partopener}[3]{\section{#2}\noindent\emph{#3}\par\medskip}

\providecommand{\BigRadionEquation}{%
  \begin{equation*}
  \boxed{\displaystyle
  \frac{180^{\circ}}{\pih}
  =50^{\circ}+A^{\circ}-\frac{20}{81}\,\Delta_4^{\circ}+\epsR^{\circ}}
  \end{equation*}}
\makeatother


\chapter{Le radion angulaire HMT--MD : raccord avec retenue des sections directe et torsionnelle et clôture exacte de l'unité de phase, d'action et de mémoire}
\label{chap:sucesor-102-radion}

\section{Composition APP--TRIT--TPK de l'état de phase : génération de l'unité orientée avec mémoire}
% Fragmento público generado desde propietarios íntegros.
% Residencia: 52.1.
% Fuente: ../incoming/radion_monografia_20260827/manuscrito/sections/01_resumen_y_guia.tex:1-128
\paragraph{Synthèse mathématique de l'objet et de sa chaîne causale}

Cette monographie reconstruit le radian à partir de la chaîne générative
\[
\APP\longrightarrow\TRIT\longrightarrow\TPK
\longrightarrow X_{\mathrm{enr}}
\longrightarrow\text{structure discrète du continu}
\longrightarrow\MD,
\]
sans prendre comme entrée ni la valeur métrologique de \(180/\pi\), ni le petit
résidu angulaire que l'on souhaite expliquer. Le résultat est une unité enrichie,
le \emph{radion}, défini comme l'unité orientée qui parcourt un radian de
phase et balaie exactement une action \(\hret\) dans l'ellipse positive associée
au même état.

La clôture directrice est
\BigRadionEquation
avec
\[
S_E=2\pih\left(\frac5{36}+\frac{25}{9}\alphah\right),
\qquad
\Delta_4=
\frac{\pih-e_{\HMT}\log\pih}{270}
\left(1+\frac{\pih}{729}\right),
\]
\[
\rho_{\mathrm{tw}}=\frac{2\cdot90}{729}=\frac{20}{81},
\qquad
\epsone=\frac{20}{81}\Delta_4-(S_E-1).
\]
La division APP prouve \(1<S_E<2\), de sorte que le quotient est exactement
un et le reste est \(D_E=S_E-1\). L'égalité
\[
1=S_E-\frac{20}{81}\Delta_4+\epsone
\]
est alors une identité interne. La carte angulaire ultérieure produit
\[
\epsR^\circ
=5.13830167585752820716851646226459474899085744\ldots
\times10^{-8}\,{}^\circ.
\]

Le texte démontre en outre que l'ellipse d'action
\[
Q=a_S\cos\theta,\qquad P=b_S\sin\theta,\qquad
a_Sb_S=2\hret
\]
satisfait
\[
\operatorname{Area}(0,\theta)=\hret\theta.
\]
Par conséquent, un radion balaie \(\hret\) et le tour \(2\pi\) enferme
\(h_{\mathrm{ret}}=2\pi\hret\). La même ellipse possède un bord symplectique
fini et un intérieur d'aire infinie lorsqu'elle est munie de la métrique
de Hilbert--Beltrami--Klein. Telle est la forme exacte de la thèse physique qui guide
l'ouvrage : la surface visible clôt l'action ; la profondeur intérieure
conserve un développement non épuisable par la projection plane.

L'exposition maintient typées quatre réalisations du même état :

\begin{enumerate}[label=\textbf{\arabic*.}]
\item le cercle \(S^1\), qui publie la phase ;
\item l'ellipse positive, qui publie l'action et l'anisotropie ;
\item l'intérieur de Klein, qui mesure la profondeur au moyen du birapport ;
\item l'hélice nonadique, qui conserve retour, mémoire et contraction.
\end{enumerate}

Ce ne sont pas quatre géométries superposées sans critère. Ce sont quatre lecteurs ayant
des domaines, des codomaines et des données conservées différents, appliqués à un unique état
enrichi. Le secteur \(90/120\), la dodécaphase, la retenue du radion, la queue
de Lambert, le noyau de Barbero et les tours globales sont également présentés
comme des objets apparentés mais non interchangeables.

\begin{thesisbox}
Le radion est l'unité orientée qui balaie \(\hret\) dans l'ellipse d'action.
Le cercle publie sa phase, Klein mesure sa profondeur, l'hélice conserve sa
mémoire et la retenue coinductive raccorde exactement ses sections directe et
torsionnelle.
\end{thesisbox}

\paragraph{Ordre d'exposition et dépendances de lecture}

L'ordre n'est pas ornemental. Les Parties I et II produisent l'objet et sa clôture
à partir d'APP--TRIT--TPK. Ce n'est qu'ensuite que sont introduites les réalisations dans le langage
de Hilbert, LC, Maxwell, Minkowski, Hodge, Lorentz, de la modularité et des orbites. Chaque
chapitre distingue trois niveaux :

\begin{description}[style=nextline,leftmargin=4.2cm,labelwidth=3.9cm]
\item[Résultat interne.] Égalité ou construction démontrée dans le domaine
HMT avec ses primitives déclarées.
\item[Formalisation réunie.] Nouvelle application articulant des résultats déjà existants
sans utiliser la cible comme sélecteur.
\item[Interface physique.] Traduction dimensionnelle ultérieure reconnaissant une
structure conventionnelle sans en faire l'origine de la sélection.
\end{description}

Les encadrés de \emph{statut et falsificateur} indiquent quelle observation concrète
invaliderait chaque promotion forte. Cette discipline n'affaiblit pas la thèse ; elle la
rend auditable. L'objectif de la narration étendue est qu'aucun mot
comme \emph{cercle}, \emph{ellipse}, \emph{torsion}, \emph{charge} ou
\emph{mémoire} ne masque un changement de type.

\paragraph{Résultats directeurs et portée probatoire}

\begin{enumerate}[label=\textbf{C\arabic*.}]
\item Clôture indépendante de la valeur finale du radion par deux publications du même état et
soustraction avec retenue APP.
\item Dérivation du coefficient \(20/81\) comme deux feuillets du secteur actif
\(90\), normalisés par la fibre \(3^6=729\).
\item Égalisateur typé entre la couture \(\Re s=1/2\), la face APP de cent
marques et la phase dodécaphasique \(\theta_2=50^\circ\).
\item Théorème d'aire du radion : \(\operatorname{Area}(\theta)=\hret\theta\).
\item Étude focale exacte de l'ellipse et de ses invariants euclidiens,
hyperboliques, symplectiques et modulaires.
\item Théorème bord--intérieur : action totale finie \(h\) et intérieur de Klein
d'aire hyperbolique infinie.
\item Réalisation exacte par covariance minimale et oscillateur LC, suivie de
la publication électromagnétique \(90/120\).
\item Atlas de non-identifications pour tous les objets \(90/120\), y compris
la séparation entre \(\epsR\), queue spectrale et noyau de Barbero \(12:-1\).
\item Caractérisation structurelle de \(\pi\) dans la carte du radion, distincte
du générateur coinductif primaire.
\item Synthèse ontologique : dimension comme rang minimal de mémoires qu'une
projection de phase ne peut plus conserver.
\end{enumerate}
% Fuente: ../incoming/radion_monografia_20260827/manuscrito/sections/01b_mapa_resultados.tex:1-133
\paragraph{Correspondance entre résultats, types et résidences causales}

La correspondance suivante ordonne les affirmations directrices et leurs démonstrations,
en maintenant visible la hiérarchie qui empêche que
la richesse physique efface la clôture mathématique qui la soutient.

% HMT_RESULT_BEGIN:RDN-01-GENEALOGIA:theorem
\begin{exactbox}[title={Généalogie APP--TRIT--TPK du radion}]
Le radion naît de la composition effective
\(\APP\to\TRIT\to\TPK\to X_{\mathrm{enr}}\to\MD\) ; il ne s'obtient pas
en redéfinissant rétrospectivement le radian conventionnel.
\end{exactbox}
% HMT_RESULT_END:RDN-01-GENEALOGIA:theorem

% HMT_RESULT_BEGIN:RDN-02-CIERRE:theorem
\begin{exactbox}[title={Clôture exacte et indépendance prospective}]
La retenue APP produit intérieurement
\[
1=S_E-\frac{20}{81}\Delta_4+\varepsilon_R^{(1)},
\qquad
\varepsilon_R^{(1)}=\frac{20}{81}\Delta_4-(S_E-1),
\]
et la carte angulaire transporte cette même identité vers
\(180^\circ/\pi_{\!\HMT}\) sans introduire le résidu cible.
\end{exactbox}
% HMT_RESULT_END:RDN-02-CIERRE:theorem

% HMT_RESULT_BEGIN:RDN-03-BISAGRA-50:theorem
\begin{exactbox}[title={Coordonnée critique de \(50^\circ\)}]
La coordonnée critique \(\Re s=1/2\), publiée sur la face APP de cent
marques et sur la roue dodécaphasique, sélectionne la phase
\(\theta_2=50^\circ\). C'est une égalité de lecteurs typés du même état.
\end{exactbox}
% HMT_RESULT_END:RDN-03-BISAGRA-50:theorem

% HMT_RESULT_BEGIN:RDN-04-COEFICIENTE-2081:theorem
\begin{exactbox}[title={Densité bidirectionnelle nonadique \(20/81\)}]
Le coefficient torsionnel n'est pas ajusté sur le réseau \(1/81\) : il est construit comme
\(\rho_{\mathrm{tw}}=2N_6/3^6=2\cdot90/729=20/81\), en conservant feuillet,
orientation et normalisation de la fibre nonadique.
\end{exactbox}
% HMT_RESULT_END:RDN-04-COEFICIENTE-2081:theorem

% HMT_RESULT_BEGIN:RDN-05-AREA-ACCION:theorem
\begin{exactbox}[title={Loi surfacique du radion}]
Pour l'ellipse positive
\(Q=a_S\cos\theta\), \(P=b_S\sin\theta\),
\(a_Sb_S=2\hbar_{\mathrm{ret}}\), l'aire orientée balayée est
\(\mathcal A(\theta)=\hbar_{\mathrm{ret}}\theta\). Un radion balaie exactement
\(\hbar_{\mathrm{ret}}\), et le tour \(2\pi\) enferme
\(h_{\mathrm{ret}}=2\pi\hbar_{\mathrm{ret}}\).
\end{exactbox}
% HMT_RESULT_END:RDN-05-AREA-ACCION:theorem

% HMT_RESULT_BEGIN:RDN-06-INVARIANTES-FOCALES:theorem
\begin{exactbox}[title={Invariants focaux de l'opérateur positif}]
Le couple publié \(A\pm C^\ast\) fixe les demi-axes, l'excentricité, la
séparation focale et le paramètre de compression symplectique. En particulier,
\((d_1/d_0)^2=C^\ast/2\), identité adimensionnelle qui sépare la forme de l'échelle.
\end{exactbox}
% HMT_RESULT_END:RDN-06-INVARIANTES-FOCALES:theorem

% HMT_RESULT_BEGIN:RDN-07-BORDE-INTERIOR:theorem
\begin{exactbox}[title={Finitude symplectique du bord et divergence de l'aire hyperbolique intérieure}]
Le bord symplectique de l'ellipse enferme l'action finie
\(h_{\mathrm{ret}}\), tandis que l'aire de Hilbert--Beltrami--Klein de l'intérieur
diverge lorsqu'elle atteint la frontière. Phase visible et profondeur ne sont pas la même
mesure.
\end{exactbox}
% HMT_RESULT_END:RDN-07-BORDE-INTERIOR:theorem

% HMT_RESULT_BEGIN:RDN-08-HELICE-MEMORIA:theorem
\begin{exactbox}[title={Relèvement hélicoïdal nonadique avec retour de phase et avance de mémoire}]
La monodromie nonadique ramène la phase sans réinitialiser l'état. L'hélice
relève le cercle en conservant retenue et mémoire ; la clôture \(2\pi/4\pi\)
distingue retour projectif, inversion d'orientation et retour complet.
\end{exactbox}
% HMT_RESULT_END:RDN-08-HELICE-MEMORIA:theorem

% HMT_RESULT_BEGIN:RDN-09-OSCILADOR-EM:development
\begin{narrativebox}[title={Réalisation oscillatoire et constitutive électromagnétique}]
La même anisotropie se réalise par une covariance minimale et un oscillateur
LC : une coordonnée stocke l'aspect électrique-élastique et sa conjuguée
l'aspect magnétique-rotationnel. La fréquence et l'impédance sont typées
séparément.
\end{narrativebox}
% HMT_RESULT_END:RDN-09-OSCILADOR-EM:development

% HMT_RESULT_BEGIN:RDN-10-TIPOS-90120:theorem
\begin{exactbox}[title={Séparation typée des opérateurs \(90/120\)}]
Extracteur dodécaphasique, retenue du radion, queue de Lambert, fonctionnelle
hexade--octade, noyau de Barbero et semi-groupe global de tours sont des objets
distincts. Le poids \(12:-1\) procède des douze secteurs de la chaîne
fondamentale dodécaphasique et de son unique couture orientée. Le coefficient de
\((1-q)^{-1}\), égal à \(1/24\), et le contre-terme entier minimal vérifient
ultérieurement le couple produit.
\end{exactbox}
% HMT_RESULT_END:RDN-10-TIPOS-90120:theorem

% HMT_RESULT_BEGIN:RDN-11-GEOMETRIAS-CAMPO:development
\begin{narrativebox}[title={Réalisations complexe, lorentzienne et duale de Hodge}]
Les régimes elliptique, parabolique et hyperbolique du TRIT sont publiés comme
rotation, seuil et transformation hyperbolique. L'étoile de Hodge, les lois constitutives et la force
de Lorentz agissent ensuite sur l'état déjà orienté ; elles ne sélectionnent pas ses
constantes.
\end{narrativebox}
% HMT_RESULT_END:RDN-11-GEOMETRIAS-CAMPO:development

% HMT_RESULT_BEGIN:RDN-12-MODULAR-ORBITAL:development
\begin{narrativebox}[title={Caractères modulaire, orbital et arithmétique de frontière}]
L'involution modulaire, les orbites elliptiques, l'avance-retard et l'arithmétique
3--6--9 reconnaissent différentes mémoires du même état. La fonction modulaire et
les vacances sont utilisées comme lecteurs ultérieurs, jamais comme substituts du
générateur APP--TRIT--TPK.
\end{narrativebox}
% HMT_RESULT_END:RDN-12-MODULAR-ORBITAL:development

% HMT_RESULT_BEGIN:RDN-13-SINTESIS-ONTOLOGICA:conclusion
\begin{thesisbox}[title={Synthèse opératoire de phase, d'action, de profondeur et de mémoire}]
Le cercle est la projection de phase ; l'ellipse est la section d'action ; Klein
mesure la profondeur ; l'hélice conserve l'histoire. Le radion est l'unité
orientée qui maintient ces publications couplées sans effacer le résidu qui
les distingue.
\end{thesisbox}
% HMT_RESULT_END:RDN-13-SINTESIS-ONTOLOGICA:conclusion

% HMT_RESULT_BEGIN:RDN-14-CONTROL-ALCANCE:control
\begin{statusbox}[title={Délimitation des domaines et continuité causale}]
Les lois complètes des masses, le raffinement électronique, CKM/CP et le
couloir exceptionnel conservent leurs propres démonstrations directrices. Cette monographie les
cite seulement là où ils partagent un antécédent ou une carte ; elle ne transfère pas au radion leurs
sélecteurs spécifiques et ne confond pas leurs résidences probatoires.
\end{statusbox}
% HMT_RESULT_END:RDN-14-CONTROL-ALCANCE:control


\subsection{Support bivarié d'APP, orientation ternaire du TRIT et actualisation de l'état enrichi par le TPK}
% Fragmento público generado desde propietarios íntegros.
% Residencia: 52.1.1.
% Fuente: ../incoming/radion_monografia_20260827/manuscrito/sections/02_genealogia.tex:1-180
\noindent La généalogie opératoire précède la mesure angulaire et détermine les variables conservées par chaque lecteur.\par\medskip

\subsubsection{Support APP--TRIT--TPK et construction de l'état enrichi}

\paragraph{Antériorité de l'état enrichi par rapport à la publication circulaire}

La définition scolaire du radian commence par un cercle déjà donné, un centre déjà
donné et une longueur d'arc déjà mesurable. Si le rayon et l'arc ont la
même longueur, l'angle sous-tendu mesure un radian. La définition est correcte
au sein de la géométrie euclidienne, mais ne répond pas à une question antérieure :
quelle structure permet à un tour de posséder phase, orientation, action et
mémoire, et pourquoi l'unité naturelle est-elle liée à \(2\pi\) ?

HMT inverse l'ordre causal. Elle construit d'abord l'état fini, puis sa
prolongation coinductive, ensuite le tour et seulement à la fin sa publication
angulaire. Le radion ne corrige pas le radian conventionnel ; il le relève vers son domaine
généalogique.

\begin{typebox}
\[
\APP\to\TRIT\to\TPK\to X_{\mathrm{enr}}
\to\mathcal C^{\mathrm{disc}}_{\mathrm{cont}}
\to\bigl(\pih,e_{\HMT},\varphi_{\HMT},\alphah,A,C^*,\Delta_4,\hret\bigr)
\to\mathfrak r_\sigma.
\]
Les mathématiques conventionnelles n'apparaissent qu'ensuite comme langage de réalisation :
\(S^1\), ellipse symplectique, métrique de Klein, espace de Hilbert, circuit
LC, écran de Minkowski et équations de champ.
\end{typebox}

\paragraph{Support bivarié d'APP : feuillets, résidus, quotients et retenue}

La plateforme d'arithmétique positionnelle (APP) conserve simultanément deux
réalisations des symboles \(1,\dots,9\) : un feuillet additif et un feuillet
multiplicatif. Un nombre n'est pas réduit à sa valeur publiée. Avec la valeur
subsistent feuillet, orientation, quotient, reste, retenue, frontière, chemin et
survie.

Le recensement de base comporte trois strates qui ne doivent jamais être identifiées par simple
coïncidence de cardinalité :
\[
104,976\longrightarrow468\;U_6\longrightarrow243\;w_6
\subset\F_3^6,\qquad |\F_3^6|=729.
\]
Les \(468\) émissions décimales, les \(243\) mots visibles et l'espace ambiant
de \(729\) mots sont des objets distincts. L'espace ambiant sera décisif pour la
normalisation \(20/81\).

En base cubique \(B=1000\), la soustraction orientée de deux mots de triades
s'effectue de droite à gauche :
\begin{equation}
u_i=t_i-v_i+c_{i+1},\qquad
r_i=u_i\bmod1000,\qquad
c_i=\frac{u_i-r_i}{1000}.
\label{eq:subacarreo}
\end{equation}
Les opérandes et la retenue distante étant fixés, quotient et reste sont uniques. Tel est
l'opérateur qui, plus loin, publiera les chiffres de \(\epsone\). Il n'y a
aucune décimale choisie après avoir vu le radian.

La complétion
\[
10^3=729+243+27+1
\]
n'est pas une curiosité décimale. Elle exprime la conservation de l'unité lors du passage
de l'espace ambiant ternaire à une chambre de mille positions. La différence
\(1000-729=271\) est la couronne complétante ; l'identité correcte est
\(999=729+270\), non \(729+271\).

\paragraph{Orientation ternaire et classification des 3 régimes opératoires du TRIT}

Le TRIT attribue à chaque état local l'un de trois régimes orientés. Dans la
convention temporelle active,
\[
\begin{aligned}
+1&\longleftrightarrow\text{retour, mémoire, passé},\\
0&\longleftrightarrow\text{seuil, présent},\\
-1&\longleftrightarrow\text{ouverture bidirectionnelle, futur}.
\end{aligned}
\]
La réalisation quadratique distingue :
\begin{equation}
J^2=-I,\qquad N^2=0,\qquad R^2=I.
\label{eq:tres-regimenes}
\end{equation}
Ainsi apparaissent, respectivement, la rotation elliptique, le seuil parabolique et
la séparation hyperbolique. On ne mélange pas trois métriques. On conserve trois
actions typées sur un antécédent commun.

La racine interne de \(-1\) n'est pas importée en supposant d'avance le plan
complexe. Dans le bloc APP, deux projections \(P,Q\) de rapport
\(r=3/4=90/120\) produisent
\[
J_{\mathrm{comm}}=\frac{[P,Q]}{\sqrt{r(1-r)}},\qquad J_{\mathrm{comm}}^2=-I,
\]
tandis que
\[
R_{\mathrm{comm}}=2P-I,\qquad R_{\mathrm{comm}}^2=I,\qquad
R_{\mathrm{comm}}J_{\mathrm{comm}}=-J_{\mathrm{comm}}R_{\mathrm{comm}}.
\]
Le couple génère la forme minimale de \(\mathrm{Cl}_{1,1}\simeq M_2(\R)\) :
rotation et ouverture émergent ensemble, mais restent des opérations différentes.

\paragraph{Composition dynamique du TPK : sélection, transport, mémoire et retour}

Le Topological Phase Kernel n'est pas le nom d'une boîte d'objets. C'est la
composition qui sélectionne la phase, transporte la cardinalité, relève les feuillets,
enregistre la mémoire et exécute le retour. La dépendance fondamentale est
\[
U_t\longrightarrow\Gamma_9\longrightarrow K_{\mathrm{ph}}.
\]
L'holonomie nonadique \(\Gamma_9\) agit sur l'état enrichi ; la
publication \(K_{\mathrm{ph}}\) est ultérieure et conserve une mémoire réduite. L'identité
\(K_{\mathrm{ph}}^9|_r=E_+E_\times\) ne transforme pas \(K_{\mathrm{ph}}\) en
générateur.

Le certificat nonadique conserve la chaîne complète
\[
w_6\to w_{12}\to w_{18}\to w_{24}\to w_{30}
\to R_{36}\to G_9,
\]
avec \(396\) niveaux, \(2376\) trits, \(43\) tours, \(387\) paires
\((K,K+9)\) par canal et cinq régions de \(\pi\). La phase satisfait
\[
g(k+9)=g(k),
\]
mais l'état complet ne revient pas :
\[
B_{k+9}\ne B_k.
\]
Tel est le germe mathématique de l'hélice : la projection angulaire se ferme ; la
mémoire avance.

\paragraph{Extension temporelle non scindée \(C_{12}\to C_{108}\to C_9\)}

La mémoire de douze pas s'organise au moyen de
\begin{equation}
0\longrightarrow C_{12}\longrightarrow C_{108}
\longrightarrow C_9\longrightarrow0.
\label{eq:extension-108}
\end{equation}
Le cocycle
\[
c(r,r')=\left[\left\lfloor\frac{r+r'}9\right\rfloor\right]_{12}
\]
mesure la retenue entre retour de phase et avance de mémoire. Neuf pas ferment
la phase observable \(C_9\) ; ils ne réinitialisent pas le registre dodécaphasique.

Il convient de séparer quatre objets :
\begin{center}
\begin{tabularx}{.95\textwidth}{>{\bfseries}l X}
\toprule
\(C_{12}\) & groupe cyclique de douze positions ;\\
\(\Rstate\) & registre enrichi \((E_m,\tau_m,\sigma_m,\kappa_m,\Lambda_m)_{m=1}^{12}\) ;\\
\(\Theta_{108}\) & publication brute dans le calendrier de 108 ;\\
\(\Gamma_{108}\) & histoire ou holonomie relevée.\\
\bottomrule
\end{tabularx}
\end{center}
C'est pourquoi l'expression vague \enquote{post-\(C_{12}\)} sera remplacée par
\enquote{postérieur au registre dodécaphasique \(\Rstate\), dans sa lecture
sexadique} lorsque telle sera l'application effective.

\paragraph{Émergence corrélative des 5 constructions du continu à partir de l'état unique}

La structure discrète du continu ne s'obtient pas en assemblant cinq théories
indépendantes. Ses cinq aspects internes ---spectral, solénoïdal, de jauge,
cohomologique et déterminantal--- émergent corrélativement du même état
APP--TRIT--TPK. Cet ouvrage ne démontre pas de nouveau le traité intégral, mais
préserve le fait causal dont le radion a besoin : la section coinductive qui
génère \(\pi,e,\varphi\) et \(\alpha\) ne vit pas sur une droite réelle nue ; elle vit
dans un objet doté de fibres, de relations, d'orientation, de retenue, de mémoire, de terminal et de
lecteur.

\begin{statusbox}
\textbf{Falsificateur de généalogie.} Si une formule ultérieure reçoit
\(180/\pi\), \(\epsR\), une masse ou une table métrologique pour choisir une phase,
un coefficient ou une branche, elle cesse d'être une génération HMT-first. La confrontation
externe peut reconnaître une sortie ; elle ne peut pas la sélectionner.
\end{statusbox}


\subsection{Indépendance prospective de la clôture par rapport à la valeur finale}
% Fragmento público generado desde propietarios íntegros.
% Residencia: 52.1.2.
% Fuente: ../incoming/radion_monografia_20260827/manuscrito/sections/02_genealogia.tex:181-272

\subsubsection{Généalogie des constantes et construction du couple angulaire}

\paragraph{Section coinductive et unicité de la valeur générée de \(\pi\)}

Soit \(x_\infty\) la section globale du système projectif. Les lecteurs
coinductifs produisent
\[
\operatorname{ev}_\pi(x_\infty)=\pih,\qquad
\operatorname{ev}_e(x_\infty)=e_{\HMT},\qquad
\operatorname{ev}_\varphi(x_\infty)=\varphi_{\HMT}.
\]
L'indice HMT ne désigne pas un autre nombre. Il conserve la généalogie de la même
constante. En particulier,
\[
\pih=3.14159265358979323846264338327950288419716939937510\ldots.
\]
Le tour positif est
\begin{equation}
T_+(x_\infty)=2\pih.
\label{eq:vuelta-plus}
\end{equation}
Sa publication angulaire d'une unité est
\begin{equation}
R_{\HMT}^{\circ}=\frac{360^\circ}{T_+}=\frac{180^\circ}{\pih}.
\label{eq:radian-publicado}
\end{equation}
L'équation ne définit pas encore le radion enrichi ; elle définit la carte visible
du radian une fois le tour généré.

\paragraph{Bidirectionnalité des chemins et prolongations conjuguées}

La bidirectionnalité ne s'énonce pas comme un slogan. L'involution concrète
\begin{equation}
J_\alpha(T)=\frac{\alphah^2}{T},\qquad
T_-=J_\alpha(T_+),\qquad T_+T_-=\alphah^2
\label{eq:bidireccional-turn}
\end{equation}
conserve un produit et inverse le feuillet. Expansion et contraction sont deux
prolongations conjuguées du même état. La racine finie d'un jet
quadratique peut approcher \(T_+\), mais ne remplace pas la section
coinductive : la différence certifiée est
\[
T_{P_9}-T_+=-5.1110566290335\ldots\times10^{-25}.
\]
Le contrôle fini est un témoin ; le fondement est la limite.

\paragraph{Dépendance causale entre \(\alpha\), action et publication angulaire}

La clôture dodécaphasique génère \(\alphah\) à partir de l'état signé et de ses
cartes réversibles. Ensuite, le lecteur déterminantal local fixe la décennie
d'action au moyen de
\[
\Sigma_H=\varphi_{\HMT}\alphah^{16},
\qquad \operatorname{ord}_{10}\Sigma_H=-34.
\]
La section de retour produit \(\hret\). Ni \(\alpha\) ni \(\hbar\) ne se
\enquote{mesurent en degrés}. Ils sont publiés ultérieurement au moyen du couple angulaire
\begin{equation}
P(\alpha,\eta_{\mathrm{ret}})=
\left([10^3\alpha]_{360},[2(\eta_{\mathrm{ret}}+\alpha)]_{360}\right)
=(A,C^*).
\label{eq:parangular}
\end{equation}
Dans la branche canonique,
\[
A=7.297352569283800997285105472380663\ldots{}^\circ,
\]
\[
C^*=2.1237383389686457242619513803933607937\ldots{}^\circ.
\]
Alors seulement, la carte analytique
\[
x=\frac{\pi A}{180},\qquad y=\frac{\pi C^*}{180}
\]
ouvre les canaux \(q_\pm=e^{-x\mp y}\).

L’ordre causal complet est
\[
\alpha\to(\eta\parallel-34)\to\hbar
\to P\to\kappa_{360}\to q_\pm
\to\text{réalisations MD}.
\]
La relation \(C^*/2-\alpha\) peut servir de vérification inverse. Elle n'est pas utilisée
pour générer \(\hret\) rétroactivement.

\begin{narrativebox}
La carte angulaire ne transforme pas une constante adimensionnelle ou une action en un
angle par décret. Elle publie deux coordonnées de l'état sur une roue finie.
Le degré est le langage de cette publication ; il n'est pas la substance ontologique de
\(\alpha\) ni de \(\hbar\).
\end{narrativebox}
% Fuente: ../incoming/radion_monografia_20260827/manuscrito/sections/03_cierre_exacto.tex:254-285
\paragraph{Indépendance prospective par rapport à la valeur finale}

Le vérificateur de l'opérateur de profondeur interdit expressément trois entrées :
\[
\frac{180}{\pi},\qquad \epsR,\qquad
\text{le décimal cible}.
\]
Il construit d'abord \(S_E\) et \(S_T\), puis exécute la soustraction APP, et ne
compare la clôture qu'à la fin. Les contrôles donnent
\[
\texttt{target\_epsilon\_used=False},\qquad
\texttt{target\_180\_over\_pi\_used=False}.
\]
De même :
\begin{itemize}
\item un seul feuillet donne un défaut d'ordre \(10^{-5}\) ;
\item remplacer \(\theta_2\) par \(\theta_1\) ou \(\theta_3\) déplace le
résultat d'environ \(\pm\pi/6\) ;
\item éliminer le canal \(e\) laisse un défaut d'ordre \(10^{-3}\) ;
\item la variante historique \(\Delta_4(A,C^*)\) produit un autre nombre ;
\item la queue spectrale postérieure à \(\Rstate\) et l'opérateur harmonique à
quatre termes ne coïncident pas non plus avec \(\epsone\).
\end{itemize}

\begin{statusbox}
\textbf{Statut.} La clôture est une \status{formalisation nouvelle} démontrée
avec une structure de départ explicite et un certificat informatique indépendant de la valeur finale.
L'égalité est exacte à la limite coinductive. Le profil court de 36 chiffres
de \(\alpha\) génère une variante qui ne diverge qu'après de nombreux chiffres ;
il est conservé comme témoin historique, non comme valeur canonique.
\end{statusbox}



\section{Sections directe et torsionnelle : retenue APP et clôture exacte du radion}
\subsection{Construction indépendante des sections \(S_E\) et \(S_T\) avant la définition du cocycle \(\varepsilon_R=S_T-S_E\)}
% Fragmento público generado desde propietarios íntegros.
% Residencia: 52.2.1.
% Fuente: ../incoming/radion_monografia_20260827/manuscrito/sections/03_cierre_exacto.tex:1-50
\noindent La clôture exacte s'obtient au moyen de deux sections construites avant le cocycle qui les relie.\par\medskip

\subsubsection{Construction indépendante des sections directe et torsionnelle}

\paragraph{Section directe du système temporel dodécaphasique}

La roue dodécaphasique a pour centres
\begin{equation}
\theta_m=30m-10^\circ,\qquad m=1,\ldots,12.
\label{eq:rueda-dodecafasica}
\end{equation}
La deuxième phase est \(\theta_2=50^\circ\). Avec la publication
\(A=10^3\alphah\), elle forme la section directe
\begin{align}
S_E
&=\frac{T_+}{360}\,(50+10^3\alphah)\\
&=2\pih\left(\frac5{36}+\frac{25}{9}\alphah\right).
\label{eq:seccion-electrica}
\end{align}
L'indice \(E\) signifie ici \emph{directe ou commune} ; sa réalisation
électrique viendra après le lecteur constitutif. On ne postule pas
littéralement \(A=E\).

Les cylindres coinductifs prouvent
\begin{equation}
1<S_E<2.
\label{eq:SE-intervalo}
\end{equation}
La division APP impose alors
\begin{equation}
q_E=\lfloor S_E\rfloor=1,
\qquad
D_E=\operatorname{rem}_1(S_E)=S_E-1.
\label{eq:DE}
\end{equation}
L'entier un n'est pas le nombre de tours de l'hélice. C'est le quotient unique
de la section directe au sein de la cellule d'unité.

\paragraph{Section torsionnelle déterminée par l'opérateur \(\Delta_4\)}

La seconde section utilise une sortie coinductive antérieure :
\begin{equation}
\Delta_4=
\frac{\pih-e_{\HMT}\log\pih}{270}
\left(1+\frac{\pih}{729}\right).
\label{eq:delta4}
\end{equation}
Le mot \emph{torsion} désigne ici le quantum global de la carte HMT. Il ne
l'identifie pas automatiquement à la torsion de Cartan
\(T^I=D_\omega e^I\). Le pont vers Cartan exige une autre flèche typée.

\subsection{Densité orientée du secteur actif et détermination unique de \(20/81\)}
% Fragmento público generado desde propietarios íntegros.
% Residencia: 52.2.2.
% Fuente: ../incoming/radion_monografia_20260827/manuscrito/sections/03_cierre_exacto.tex:51-103

\paragraph{Densité orientée \(20/81\) du secteur actif \(90/120\)}

Le sélecteur tritique de phase sépare
\[
H_+=\{1,4,7\},\qquad H_-=\{2,5,8\},\qquad H_0=\{3,6,9\}.
\]
Il y a six phases actives et quinze paires non ordonnées :
\[
H_{\mathrm{act}}=H_+\cup H_-,\qquad
\left|\binom{H_{\mathrm{act}}}2\right|=\binom62=15.
\]
Le mode actif et son extension à huit positions sont
\begin{equation}
N_6=6\binom62=90,\qquad
N_8=8\binom62=120.
\label{eq:90-120-incidencia}
\end{equation}
Deux feuillets conjugués du secteur actif, normalisés par la fibre ambiante,
donnent
\begin{equation}
\rho_{\mathrm{tw}}
=\frac{2N_6}{|\F_3^6|}
=\frac{2\cdot90}{729}
=\frac{180}{729}
=\boxed{\frac{20}{81}}.
\label{eq:rho-twoway}
\end{equation}

\begin{proposition}[Unicité relative de la densité]
Le secteur actif \(N_6=90\), les deux feuillets conjugués et la
normalisation par \(|\F_3^6|=729\) étant fixés, la densité orientée est nécessairement
\(20/81\).
\end{proposition}
\begin{proof}
La cardinalité transportée est \(2N_6=180\). La normalisation par l'espace ambiant
donne \(180/729\) ; diviser le numérateur et le dénominateur par \(9\) produit \(20/81\).
Aucune comparaison avec \(180/\pi\) ni avec \(\epsR\) n'intervient.
\end{proof}

La lecture historique comme projection sur un réseau \(1/81\) est un contrôle de
rigidité utile, mais n'est plus le fondement : le coefficient s'obtient avant
la clôture. L'ablation d'un feuillet produit \(10/81\), et la sortie change à
l'ordre \(10^{-5}\) ; le facteur deux n'est donc pas décoratif.

La section torsionnelle est
\begin{equation}
\Phi_T=\rho_{\mathrm{tw}}\Delta_4,
\qquad
S_T=1+\Phi_T.
\label{eq:seccion-torsional}
\end{equation}


\subsection{La retenue \(\varepsilon_R\) comme cocycle de transition entre sections}
% Fragmento público generado desde propietarios íntegros.
% Residencia: 52.2.3.
% Fuente: ../incoming/radion_monografia_20260827/manuscrito/sections/03_cierre_exacto.tex:104-253
\subsubsection{Cocycle de transition entre les sections directe et torsionnelle}

\paragraph{Définition opératoire de la retenue \(\varepsilon_R\)}

\begin{definition}[Retenue réelle du radion]
La sortie unitaire du radion est le cocycle de transition
\begin{equation}
\epsone:=S_T-S_E
=\Phi_T-D_E
=\frac{20}{81}\Delta_4-operatorname{rem}_1(S_E).
\label{eq:epsilon-def}
\end{equation}
\end{definition}

L'équation contient une soustraction, mais ce n'est pas un ajustement à la cible. Ses deux
opérandes \(S_T\) et \(S_E\) sont générés séparément. L'indépendance
démontrée est une indépendance par rapport à \(180/\pi\) et à un
\(\varepsilon\) cible ; on n'affirme pas d'indépendance algébrique entre une
différence et ses termes.

\begin{theorem}[Clôture exacte unitaire]
Avec \(S_E\), \(\Phi_T\) et \(\epsone\) définis par
\eqref{eq:seccion-electrica}, \eqref{eq:delta4}, \eqref{eq:rho-twoway} et
\eqref{eq:epsilon-def}, on a exactement
\begin{equation}
\boxed{1=S_E-\Phi_T+\epsone.}
\label{eq:cierre-unitario}
\end{equation}
\end{theorem}
\begin{proof}
Par \eqref{eq:SE-intervalo}, \(D_E=S_E-1\). Par définition,
\(\epsone=\Phi_T-D_E\). Alors
\[
S_E-\Phi_T+\epsone
=S_E-\Phi_T+\Phi_T-(S_E-1)=1.
\]
La force de la preuve ne réside pas dans cette annulation élémentaire, mais dans le fait que la
généalogie préalable fixe les deux opérandes sans utiliser la clôture comme entrée.
\end{proof}

\paragraph{Publication de la clôture dans la carte angulaire}

La carte du tour multiplie une fraction unitaire par
\(360/T_+=180/\pih\). En particulier,
\[
\epsR^\circ=\frac{360^\circ}{T_+}\epsone.
\]
Multiplier \eqref{eq:cierre-unitario} par \(360^\circ/T_+\) et utiliser
\[
\frac{360}{T_+}S_E
=360\left(\frac5{36}+\frac{25}{9}\alphah\right)
=50+1000\alphah=50+A
\]
produit la clôture angulaire.

\begin{theorem}[Équation exacte du radion]
\BigRadionEquation
\end{theorem}

L'équation dit quelque chose de plus précis que \enquote{le radian est
\(57.2957\ldots{}^\circ\)}. Elle décompose cette publication en quatre opérations
avec démonstration de référence :
\begin{center}
\begin{tabularx}{.96\textwidth}{>{\bfseries}l X}
\toprule
\(50^\circ\) & deuxième phase de la roue ; charnière critique dans une carte déclarée ;\\
\(A^\circ\) & coordonnée commune de la publication parangulaire ;\\
\(-\frac{20}{81}\Delta_4^\circ\) & transport orienté de la torsion globale ;\\
\(+\epsR^\circ\) & retenue réelle qui raccorde les deux relèvements.\\
\bottomrule
\end{tabularx}
\end{center}

\paragraph{Convergence selon la profondeur du raffinement}

À la profondeur \(N\), le générateur fournit des cylindres rationnels
\(p_N\to\pih\) et \(e_N\to e_{\HMT}\). On définit
\[
S_{E,N}=2p_N\left(\frac5{36}+\frac{25}{9}\alphah\right),
\]
\[
\Delta_{4,N}=\frac{p_N-e_N\log p_N}{270}
\left(1+\frac{p_N}{729}\right),\qquad
S_{T,N}=1+\frac{20}{81}\Delta_{4,N},
\]
et
\[
\mathcal E_N=S_{T,N}-S_{E,N}.
\]
La continuité de la fonctionnelle en \(3<p<4\), \(2<e<3\), jointe aux
dérivées monotones, transporte chaque couple de cylindres vers un intervalle
certifié. À la profondeur de \(45\) blocs, sa largeur est inférieure à
\(4.81\times10^{-130}\).

Chaque retour de neuf niveaux contient six trits par niveau. C'est pourquoi le taux
de raffinement est
\begin{equation}
729^{-9}=3^{-54}.
\label{eq:contraccion-54}
\end{equation}
Ce chiffre contrôle la largeur des intervalles ; ce n'est pas la valeur de
\(\epsR\).

La soustraction de triades stabilise
\begin{center}
\ttfamily
000|000|000|896|802|822|044|562|981|999|050|695|020|651|286|413
\end{center}
et le préfixe de retenues
\begin{center}
\ttfamily
0,0,0,0,0,-1,0,-1,-1,-1,0,-1,0.
\end{center}
Tel est le mécanisme concret de mémoire décimale : non une queue choisie, mais
la récurrence \eqref{eq:subacarreo} appliquée à des opérandes préalablement
construits.

\paragraph{Évaluation des valeurs canoniques}

La sortie coinductive est
\begin{align}
\epsone
&=8.9680282204456298199905069502065128641372736205009\ldots
\times10^{-10},\\
\epsR^\circ
&=5.1383016758575282071685164622645947489908574400054\ldots
\times10^{-8}\,{}^\circ.
\label{eq:epsilon-valores}
\end{align}
La publication complète donne
\[
\frac{180^\circ}{\pih}
=57.2957795130823208767981548141051703324054724665643\ldots{}^\circ.
\]

\begin{exactbox}
\begin{align*}
50^\circ+A^\circ
&=57.29735256928380099728510547238066299956\ldots{}^\circ,\\
\frac{20}{81}\Delta_4^\circ
&=0.00157310758449687906223272996065728980465\ldots{}^\circ,\\
50^\circ+A^\circ-\frac{20}{81}\Delta_4^\circ
&=57.29577946169930411822287274242000570976\ldots{}^\circ,\\
\epsR^\circ
&=0.00000005138301675857528207168516462264594749\ldots{}^\circ,\\
\text{somme finale}
&=57.29577951308232087679815481410517033241\ldots{}^\circ.
\end{align*}
\end{exactbox}



\section{Coordonnée critique et système temporel dodécaphasique orienté : publication affine de la phase de \(50^\circ\)}
\subsection{Application affine \(j_{100}\) de la carte APP et unicité de la phase \(m=2\) dans \(j_{100}(1/2)=\theta_m=50^\circ\)}
% Fragmento público generado desde propietarios íntegros.
% Residencia: 52.3.1.
% Fuente: ../incoming/radion_monografia_20260827/manuscrito/sections/04_cincuenta_critica.tex:1-105
\noindent La coordonnée critique relie la carte de Cayley, le registre APP et la phase dodécaphasique par des applications typées.\par\medskip

\subsubsection{Deux réalisations exactes de la coordonnée angulaire \(50^\circ\)}

\paragraph{Deuxième phase du système temporel dodécaphasique}

Par \eqref{eq:rueda-dodecafasica}, la suite des centres commence
\[
20^\circ,\ 50^\circ,\ 80^\circ,\ 110^\circ,\ldots.
\]
Résoudre
\[
30m-10=50
\]
donne uniquement \(m=2\). Par conséquent, \(50^\circ\) n'est ni un arrondi de \(A\),
ni un demi-tour ni le centre d'un cercle : c'est une phase typée du calendrier
dodécaphasique. La fraction de tour est
\[
\frac{50}{360}=\frac5{36},
\]
exactement le premier terme de \(S_E/T_+\).

\paragraph{Égalisation de la coordonnée critique de Cayley}

L'involution fonctionnelle
\[
\rho\longmapsto1-\bar\rho
\]
fixe
\[
\operatorname{Fix}(\rho\mapsto1-\bar\rho)
=\{\rho:\Re\rho=1/2\}.
\]
Le nombre \(1/2\) est la coordonnée réelle autoduale de l'intervalle fonctionnel
\(0\leq\Re s\leq1\). Ce n'est pas \enquote{la moitié d'une droite} au sens
longitudinal, et il ne contient pas la hauteur spectrale \(\gamma\). Pour
\[
\rho_\gamma=\frac12+i\gamma,
\]
la transformation de Cayley peut s'écrire
\begin{equation}
\tau_\gamma
=\frac{\rho_\gamma-1}{\rho_\gamma}
=\frac{\gamma+i/2}{\gamma-i/2}
\in S^1,
\qquad
\gamma=\frac12\cot\frac{\theta_\gamma}{2}.
\label{eq:cayley-critica}
\end{equation}
Ainsi, \(1/2\) fixe la couture, \(\gamma\) détermine le caractère angulaire et
\(\Gamma_9\) conserve la profondeur de ses itérations.

\paragraph{Injection affine \(j_{100}\) du registre APP dans la carte angulaire}

La complétion cubique dispose une chambre de dix marques par coordonnée. Une
face contient \(10^2=100\) positions. Nous déclarons la carte affine
\begin{equation}
j_{100}:[0,1]\longrightarrow[0,100^\circ],
\qquad j_{100}(\sigma)=100\sigma^\circ.
\label{eq:j100}
\end{equation}
Alors
\[
j_{100}\left(\frac12\right)=50^\circ.
\]
En exigeant simultanément \(j_{100}(1/2)=\theta_m\), il vient
\[
100\cdot\frac12=30m-10
\quad\Longleftrightarrow\quad m=2.
\]

\begin{theorem}[Égalisateur critique--dodécaphasique]
Dans la carte APP de cent marques, la coordonnée fixe de l'involution critique
se publie exactement dans la deuxième phase de la roue :
\begin{equation}
\boxed{j_{100}(1/2)=\theta_2=50^\circ.}
\label{eq:igualador-50}
\end{equation}
La phase \(m=2\) est unique.
\end{theorem}

Le théorème est exact une fois \eqref{eq:j100} déclarée. Son statut est
\status{formalisation nouvelle}: le demi critique, la complétion \(10^3\), la
face de cent marques et \(\theta_2\) étaient déjà construits ; le choix de cette
face comme carte critique affine réunit pour la première fois leurs domaines.

\paragraph{Compatibilité de la valeur critique avec la carte de Cayley}

La transformation de Cayley n'envoie pas \(1/2\) sur \(50^\circ\). Si
\(\gamma=0\), \eqref{eq:cayley-critica} donne
\[
\tau_0=-1=e^{i\pi},
\]
c'est-à-dire \(180^\circ\). Les \(50^\circ\) appartiennent à la carte APP de cent
marques et au calendrier dodécaphasique ; l'angle de Cayley dépend de \(\gamma\).
Cette distinction empêche de réduire tous les zéros de la fonction zêta à une seule
phase.

\begin{falsifier}
L'identification \(1/2\leftrightarrow50^\circ\) échoue si l'on élimine la
carte \eqref{eq:j100}. L'égalité scalaire \(100(1/2)=50\) ne suffit pas
à identifier à elle seule deux objets. Le théorème dépend de la face APP
distinguée et de la roue dodécaphasique.
\end{falsifier}


\subsection{Compatibilité de la phase critique avec le retour nonadique et l'avance de mémoire}
% Fragmento público generado desde propietarios íntegros.
% Residencia: 52.3.2.
% Fuente: ../incoming/radion_monografia_20260827/manuscrito/sections/04_cincuenta_critica.tex:106-171
\subsubsection{Séparation typée des 2 objets de valeur \(1/2\)}

APP contient en outre un mode singulier avec
\[
s=\frac12,\qquad s^2=\frac14,\qquad
1-s^2=\frac34=\frac{90}{120},\qquad
s\sqrt{1-s^2}=\frac{\sqrt3}{4}.
\]
Ce \(s\) naît des angles principaux des feuillets somme et produit. La
couture critique \(\Re\rho=1/2\) naît d'une involution fonctionnelle. Que les deux
coordonnées valent \(1/2\) est une coïncidence scalaire dotée d'un contenu possible,
mais n'autorise pas à les identifier sans un carré typé.

La manière correcte de conserver la découverte est un diagramme d'égalisateurs :
\[
\begin{tikzpicture}[baseline=(current bounding box.center),node distance=14mm and 23mm,
 every node/.style={font=\small}]
\node[draw=RadBlue,rounded corners,fill=RadMist] (crit) {couture critique \(\Re s=1/2\)};
\node[draw=RadTeal,rounded corners,fill=RadCream,right=of crit] (face) {carte \(j_{100}\)};
\node[draw=RadCopper,rounded corners,fill=RadCream,right=of face] (theta) {phase \(\theta_2=50^\circ\)};
\node[draw=RadRose,rounded corners,fill=white,below=of face] (mode) {mode APP \(s=1/2\), \(1-s^2=90/120\)};
\draw[-{Latex},thick,RadBlue] (crit)--node[above]{\(j_{100}\)}(face);
\draw[-{Latex},thick,RadTeal] (face)--node[above]{égalisateur}(theta);
\draw[-{Latex},dashed,RadRose] (mode)--node[right,align=left]{même chiffre;\\domaine distinct}(face);
\end{tikzpicture}
\]
Le diagramme ne simule pas la flèche absente. Il montre ce qui est démontré, ce qui a
été formalisé et quelle égalité nécessite encore un transport supplémentaire si
l'on vise une identification opératoire globale.

\subsubsection{Compatibilité entre cercle spectral et mémoire hélicoïdale}

Le lecteur de Cayley--Pontryagin transporte une fibre spectrale au moyen de
\[
\widetilde U_{\Gamma_9^n}J\psi_\gamma
=\tau_\gamma^nJ\psi_\gamma.
\]
La publication circulaire est \(\tau_\gamma^n\in S^1\). Le relèvement
\[
(\tau_\gamma^n,n)\in S^1\times\Z
\]
conserve le nombre de retours. Sur le feuillet contractant apparaît
\[
\widetilde M_9^nJ\psi_\gamma
=\left(\frac59\right)^n\tau_\gamma^nJ\psi_\gamma.
\]
Nous avons donc trois images typées :
\[
\text{cercle de phase},\qquad
\text{hélice de mémoire},\qquad
\text{spirale contractante}.
\]

Un zéro n'est pas un tour, et un neuf n'est pas un zéro particulier de la fonction
zêta. Le zéro fixe un caractère \(\tau_\gamma\) qui persiste à travers les
tours ; le retour nonadique incrémente la mémoire. Le neuf fonctionne comme zéro
résiduel dans \(\Z/9\Z\), opération arithmétique distincte de l'énumération
des zéros spectraux.

\begin{narrativebox}
La droite critique devient un cercle lorsque nous oublions la hauteur linéaire et
conservons le caractère unitaire. Elle devient une hélice lorsque nous réinsérons
la mémoire du nombre de fois où le caractère a été transporté. La profondeur
n'est pas une troisième coordonnée décorative : c'est l'information que la projection
circulaire ne peut pas retenir.
\end{narrativebox}


\section{Couple angulaire et action réduite : construction de l'ellipse positive et de l'aire par radion}
\subsection{Loi surfacique du radion : \(\mathcal A(\theta)=\hbar_{\mathrm{ret}}\theta\) et \(\mathcal A(2\pi)=h_{\mathrm{ret}}\)}
% Fragmento público generado desde propietarios íntegros.
% Residencia: 52.4.1.
% Fuente: ../incoming/radion_monografia_20260827/manuscrito/sections/05_geometria_accion.tex:1-217
\noindent La géométrie spectrale détermine la forme et la loi surfacique détermine l'action transportée par la phase.\par\medskip

\subsubsection{Équivalence entre les cartes spectrale et géométrique}

\paragraph{Opérateur spectral positif et valeurs propres conjuguées}

Le couple angulaire produit l'opérateur bicouche
\begin{equation}
B_1=AI+C^*R,\qquad R^2=I,\qquad
P_\pm=\frac{I\pm R}{2}.
\label{eq:B1}
\end{equation}
Ses valeurs propres sont \(A\pm C^*\). Comme \(0<C^*<A\), les deux sont positives.
Dans la carte \(\kappa=\pi/180\), l'ellipse spectrale a pour demi-axes
\begin{equation}
a_1^2=\kappa(A+C^*),\qquad
b_1^2=\kappa(A-C^*).
\label{eq:semiejes-espectrales}
\end{equation}
La décomposition est réversible :
\begin{equation}
A=\frac{a_1^2+b_1^2}{2\kappa},\qquad
C^*=\frac{a_1^2-b_1^2}{2\kappa}.
\label{eq:reconstruccion-AC}
\end{equation}
L'ellipse n'est pas dessinée pour illustrer un nombre. Elle reconstruit exactement le
couple qui l'a générée.

Définissons
\begin{equation}
\rho_1=\sqrt{A^2-C^{*2}},\qquad
\eta=\operatorname{artanh}\frac{C^*}{A}
=\frac12\log\frac{A+C^*}{A-C^*}.
\label{eq:eta}
\end{equation}
Alors
\[
A\pm C^*=\rho_1e^{\pm\eta},
\]
et
\begin{equation}
a_1=\sqrt{\kappa\rho_1}\,e^{\eta/2},
\qquad
b_1=\sqrt{\kappa\rho_1}\,e^{-\eta/2}.
\label{eq:semiejes-eta}
\end{equation}
Le produit conserve l'échelle ; le quotient conserve la forme :
\[
a_1b_1=\kappa\rho_1,
\qquad
\frac{b_1}{a_1}=e^{-\eta}
=\sqrt{\frac{A-C^*}{A+C^*}}.
\]

\paragraph{Invariants focaux de l'opérateur positif}

La demi-distance focale est
\begin{equation}
c_1^2=a_1^2-b_1^2=2\kappa C^*,
\qquad
F_\pm=(\pm c_1,0).
\label{eq:focos}
\end{equation}
La distance totale vaut
\begin{equation}
d_1=2c_1=2\sqrt{2\kappa C^*}
=0.544545509325626623406299779866023\ldots.
\label{eq:d1}
\end{equation}
Le nombre isolé n'est pas la thèse. L'identité \(c_1^2=2\kappa C^*\) dit que
\(C^*\) est littéralement la scission focale quadratique dans la carte
distinguée. Les foyers sont les extrémités conjuguées de l'ouverture spectrale ;
ils ne sont pas, à eux seuls, les pôles électrique et magnétique. Cette lecture nécessite
le transducteur constitutif.

L'excentricité est
\begin{equation}
e_f=\frac{c_1}{a_1},\qquad
e_f^2=1-e^{-2\eta}=\frac{2C^*}{A+C^*}.
\label{eq:excentricidad}
\end{equation}
Numériquement,
\[
\begin{aligned}
\eta&=0.299689683921630799684926562863927\ldots,\\
e_f&=0.671451895550191986916277055864332\ldots.
\end{aligned}
\]

\paragraph{Bloc APP de référence et transport normalisé}

Le bloc central
\[
B_0=7I+2R=9P_++5P_-
\]
possède
\[
\rho_0=\sqrt{45},\qquad
\eta_0=\frac12\log\frac95,\qquad
e_0=\frac23.
\]
Sa distance focale est
\[
d_0=4\sqrt\kappa
=0.528443639680801468446003835349358\ldots.
\]
La comparaison exacte est
\begin{equation}
\boxed{\left(\frac{d_1}{d_0}\right)^2=\frac{C^*}{2}},
\qquad
d_1=d_0\sqrt{\frac{C^*}{2}}.
\label{eq:invariante-focal-relativo}
\end{equation}
Telle est la généalogie défendable de \(0.544545\ldots\) : une échelle focale APP
de référence multipliée par un invariant relatif. Sa proximité visuelle avec
\(0.54\) n'est pas une identité avec le défaut APP \(54\).

Le transport depuis \(B_0\) se sépare en échelle et anisotropie :
\begin{equation}
T_{\mathrm{rel}}=\frac{\rho_1}{\sqrt{45}}
\exp\bigl((\eta-\eta_0)R\bigr).
\label{eq:transporte-rel}
\end{equation}
Sur les feuillets,
\[
t_+=\frac{A+C^*}{9}=1.0467878786947163\ldots,
\qquad
t_-=\frac{A-C^*}{5}=1.0347228460630311\ldots.
\]
La déformation n'est pas une seule décimale. Sa partie isotrope est
\(\sqrt{t_+t_-}\) ; sa partie hyperbolique est
\(\tfrac12\log(t_+/t_-)\).

\subsubsection{Ellipse positive associée au couple angulaire}

\paragraph{Construction au moyen de demi-axes conjugués}

L'échelle spectrale précédente est dépourvue d'unités d'action. La mécanique
dimensionnelle publie la même forme sur la droite interne d'action :
\begin{equation}
a_S=\sqrt{2\hret e^\eta},
\qquad
b_S=\sqrt{2\hret e^{-\eta}}.
\label{eq:semiejes-accion}
\end{equation}
Immédiatement,
\begin{equation}
a_Sb_S=2\hret,\qquad
\frac{a_S}{b_S}=e^\eta.
\label{eq:producto-forma}
\end{equation}
Le produit fixe l'action ; le quotient fixe l'anisotropie. L'ellipse est
paramétrée par
\begin{equation}
Q(\theta)=a_S\cos\theta,\qquad
P(\theta)=b_S\sin\theta.
\label{eq:elipse-param}
\end{equation}

Les ellipses spectrale et d'action sont homothétiques :
\begin{equation}
\lambda_{\mathrm{act}}^2=\frac{2\hret}{\kappa\rho_1},
\qquad
(a_S,b_S,c_S)=\lambda_{\mathrm{act}}(a_1,b_1,c_1).
\label{eq:homotecia}
\end{equation}
Elles partagent la forme \(e^{-\eta}\), non la dimension physique.

\paragraph{Loi surfacique de l'action par radion}

La forme symplectique canonique est \(\omega=\diff Q\wedge\diff P\). L'aire
orientée balayée de \(0\) à \(\theta\) est
\begin{align}
\mathcal A(\theta)
&=\frac12\int_0^\theta
\bigl(Q(t)P'(t)-P(t)Q'(t)\bigr)\,\diff t\\
&=\frac12a_Sb_S\int_0^\theta
(\cos^2t+\sin^2t)\,\diff t.
\end{align}
En utilisant \eqref{eq:producto-forma}, nous obtenons le résultat central.

\begin{theorem}[Action par radion]
Pour l'ellipse \eqref{eq:elipse-param},
\begin{equation}
\boxed{\mathcal A(\theta)=\hret\theta.}
\label{eq:area-theorem}
\end{equation}
En particulier,
\begin{equation}
\boxed{\mathcal A(1)=\hret},
\qquad
\boxed{\mathcal A(2\pi)=2\pi\hret=\hfull}.
\label{eq:area-radion-vuelta}
\end{equation}
\end{theorem}

Cette égalité permet de définir le radion sans faire d'abord appel à la longueur
d'arc.

\begin{definition}[Radion orienté]
\begin{equation}
\mathfrak r_\sigma=(\theta=1,\sigma),
\qquad \sigma\in\{+1,-1\},
\qquad
\mathcal A(\mathfrak r_\sigma)=\sigma\hret.
\label{eq:radion-oriented}
\end{equation}
\end{definition}

Le radian est la projection qui oublie \(\sigma\), le feuillet, la mémoire et
l'aire. Le radion conserve cette information. D'où le nom \emph{rad--ion} :
l'orientation précède la charge physique, et la charge apparaît plus tard comme
\[
Q(x)=n_Q(x)e_{\mathrm{int}},\qquad
e_{\mathrm{int}}=\sqrt{\frac{2\alphah\hfull}{Z_{\mathrm{int}}}}.
\]


\subsection{Demi-axes conjugués, invariants focaux et orientation de l'anisotropie}
% Fragmento público generado desde propietarios íntegros.
% Residencia: 52.4.2.
% Fuente: ../incoming/radion_monografia_20260827/manuscrito/sections/05_geometria_accion.tex:218-290
\paragraph{Correspondance entre phase, action, aire et orientation}

\begin{figure}[H]
\centering
\begin{tikzpicture}[scale=1.12]
  \def\aa{4.0}
  \def\bb{2.95}
  \def\cc{2.70}
  \draw[->,RadGray] (-4.7,0)--(4.7,0) node[right]{\(Q\)};
  \draw[->,RadGray] (0,-3.5)--(0,3.5) node[above]{\(P\)};
  \fill[RadTeal!8] (0,0) ellipse[x radius=\aa,y radius=\bb];
  \draw[RadTeal,line width=1.35pt] (0,0) ellipse[x radius=\aa,y radius=\bb];
  \fill[RadCopper] (-\cc,0) circle (2pt) node[below=2mm]{\(F_-\)};
  \fill[RadCopper] (\cc,0) circle (2pt) node[below=2mm]{\(F_+\)};
  \draw[RadGold,line width=1.1pt,-{Latex}] (\aa,0)
     arc[start angle=0,end angle=57.3,x radius=\aa,y radius=\bb];
  \node[RadGold] at (3.7,2.3) {un radion};
  \draw[dashed,RadBlue] (0,0)--(\aa,0) node[midway,below]{\(a_S\)};
  \draw[dashed,RadBlue] (0,0)--(0,\bb) node[midway,left]{\(b_S\)};
  \node[align=center,RadNavy] at (0,-3.25)
    {aire totale \(\pi a_Sb_S=2\pi\hret=\hfull\)};
\end{tikzpicture}
\caption{L'ellipse d'action. Un arc paramétrique d'un radian balaie
\(\hret\) ; un tour enferme \(\hfull\). Les foyers codent l'ouverture
\(C^*\) dans la carte distinguée.}
\label{fig:elipse-accion}
\end{figure}

\paragraph{Relation entre l'action réduite, l'incertitude et l'aire}

La relation \(\hbar=h/(2\pi)\) est souvent présentée comme une commodité
de notation : lorsqu'on décrit des oscillations, des rotations et des transformées de Fourier,
les facteurs \(2\pi\) disparaissent si l'on utilise \(\hbar\). L'ellipse d'action
révèle une lecture plus forte. La division par \(2\pi\) est exactement le
passage de l'action d'un tour à l'action par unité de phase :
\[
h=\mathcal A(2\pi),\qquad \hbar=\mathcal A(1).
\]
La constante réduite n'est pas seulement \(h\) avec une notation plus commode ; c'est la
densité surfacique d'action par rapport à la phase.

Heisenberg, Dirac et Jordan ont établi le langage des variables conjuguées,
des commutateurs et des transformations canoniques. On ne leur attribue pas ici une
\enquote{ellipse ontologique} littérale sans source. La contribution HMT--MD consiste à
identifier la figure positive, à construire ses demi-axes à partir de \((A,C^*)\) et à
démontrer \eqref{eq:area-theorem}. La généalogie historique ouvre le langage ;
le théorème actuel fixe l'objet.

\begin{narrativebox}
L'action ne s'ajoute pas ensuite au mouvement. Dans l'ellipse du radion,
l'action est l'aire que le mouvement génère. Un tour ne transporte pas une
constante préfabriquée : il l'enferme.
\end{narrativebox}

\Needspace{10\baselineskip}
\paragraph{Contrôles algébriques des invariants focaux}

Les ressemblances décimales sont soumises à un contrôle :
\begin{center}
\begin{tabularx}{.96\textwidth}{l r X}
\toprule
Comparaison & Résidu & Verdict\\
\midrule
\(d_1\) face à \(0.54\) & \(4.5455\times10^{-3}\) & pas d'égalité ; le \(54\) ne se déduit pas ;\\
\(d_1^2\) face à \(8/27\) & \(2.3352\times10^{-4}\) & proximité sans application ;\\
\(\eta\) frente a \(3/10\) & \(-3.1032\times10^{-4}\) & no igualdad;\\
\(e_f\) face à \(2/3\) & \(4.7852\times10^{-3}\) & \(2/3\) appartient exactement à \(B_0\) ;\\
\(e^{-\eta}\) face à \(20/27\) & \(3.0740\times10^{-4}\) & no igualdad.\\
\bottomrule
\end{tabularx}
\end{center}
Les invariants solides sont \eqref{eq:eta}, \eqref{eq:excentricidad} et
\eqref{eq:invariante-focal-relativo}. Le critère HMT ne consiste pas à récompenser
une ressemblance décimale, mais à exiger objet, opération et conservation.


\section{Retour nonadique et intérieur projectif : disque de Klein et profondeur hélicoïdale de mémoire}
\subsection{Normalisation au disque et intérieur hyperbolique de Hilbert--Beltrami--Klein}
% Fragmento público generado desde propietarios íntegros.
% Residencia: 52.5.1.
% Fuente: ../incoming/radion_monografia_20260827/manuscrito/sections/06_interior_helice.tex:1-108
\noindent La géométrie projective distingue l'action finie de la profondeur hyperbolique et de la mémoire du retour.\par\medskip

\subsubsection{Intérieur hyperbolique de Hilbert--Beltrami--Klein}

\paragraph{Normalisation projective dans le disque de Klein}

L'application affine
\[
u=\frac{Q}{a_S},\qquad v=\frac{P}{b_S}
\]
envoie l'ellipse d'action sur le disque unité
\[
D=\{(u,v):u^2+v^2<1\}.
\]
Sur chaque corde, la distance de Hilbert est définie par le birapport
de ses extrémités de frontière. La métrique infinitésimale résultante est la forme
de Beltrami--Klein
\begin{equation}
\diff s_K^2=
\frac{(1-r^2)(\diff u^2+\diff v^2)+(u\diff u+v\diff v)^2}
{(1-r^2)^2},
\qquad r^2=u^2+v^2.
\label{eq:klein-metric}
\end{equation}
Avec cette normalisation, la courbure gaussienne est constante :
\begin{equation}
K=-1.
\label{eq:klein-curvature}
\end{equation}
Les géodésiques sont des cordes euclidiennes, mais leur longueur hyperbolique diverge
à l'approche de la frontière.

\paragraph{Finitude de la frontière et infinitude de la profondeur hyperbolique}

Le déterminant métrique de \eqref{eq:klein-metric} produit
\begin{equation}
\diff A_K=\frac{\diff u\,\diff v}{(1-r^2)^{3/2}}.
\label{eq:klein-area-density}
\end{equation}
L'aire du sous-disque \(r<R<1\) est
\begin{align}
A_K(R)
&=\int_0^{2\pi}\int_0^R
\frac{r\,\diff r\,\diff\theta}{(1-r^2)^{3/2}}\\
&=2\pi\left(\frac1{\sqrt{1-R^2}}-1\right).
\label{eq:klein-area-R}
\end{align}
\Needspace{4\baselineskip}
Par conséquent,
\[
\lim_{R\to1^-}A_K(R)=+\infty.
\]
Simultanément, la même frontière enferme l'aire symplectique
\[
\pi a_Sb_S=2\pi\hret=\hfull<\infty.
\]

\begin{theorem}[Finitude de l'action et infinitude de la profondeur]
L'ellipse du radion possède une action de tour finie \(\hfull\), tandis que
son intérieur, muni de la métrique de Hilbert--Beltrami--Klein, a une aire
hyperbolique infinie :
\begin{equation}
\boxed{
\operatorname{Area}_{\mathrm{simp}}(\partial\mathbb E_S)=\hfull<\infty,
\qquad
\operatorname{Area}_{K}(\mathbb E_S)=+\infty.}
\label{eq:finite-infinite}
\end{equation}
\end{theorem}

Il n'y a pas de contradiction : \(\operatorname{Area}_{\mathrm{simp}}\) et
\(\operatorname{Area}_K\) sont des mesures différentes. La première compte l'action
orientée en phase ; la seconde mesure la profondeur projective intérieure. La thèse
\enquote{bord fini, intérieur infini} est désormais une égalité et une limite,
non une métaphore.

\paragraph{Coordonnées focales dans l'intérieur de Klein}

Dans le disque normalisé, les foyers se situent en \((\pm e_f,0)\). La distance
du centre à un foyer est
\begin{equation}
u_f=\operatorname{artanh}e_f
=\operatorname{arcosh}(e^\eta)
=0.813382331175342333760889731088228\ldots.
\label{eq:klein-focal}
\end{equation}
La distance foyer--foyer est
\[
d_K(F_-,F_+)=2u_f
=1.62676466235068466752177946217646\ldots.
\]
On a
\begin{equation}
e_f=\tanh u_f,\qquad
e^{-\eta}=\operatorname{sech}u_f,\qquad
\eta=\log\cosh u_f.
\label{eq:focal-hyperbolic-chain}
\end{equation}
La région jusqu'au rayon focal a pour aire
\[
A_K(e_f)=2\pi(e^\eta-1)
=2.19559620884061869541206115980360\ldots.
\]

Toutes les ellipses ouvertes sont projectivement isométriques en tant que domaines de
Hilbert. C'est pourquoi la courbure \(-1\) seule ne mémorise pas l'excentricité. La
forme HMT réside dans la structure conjointe : carte affine distinguée, forme
symplectique, étiquetage des feuillets et paramétrisation de frontière.

\subsection{Différence entre retour de phase et retour de l'état enrichi}
% Fragmento público generado desde propietarios íntegros.
% Residencia: 52.5.2.
% Fuente: ../incoming/radion_monografia_20260827/manuscrito/sections/06_interior_helice.tex:109-201
\subsubsection{Relèvement hélicoïdal du retour nonadique}

\paragraph{Décomposition nonadique et quotient de mémoire}

L'opération élémentaire
\begin{equation}
n=9q_9(n)+\rho_9(n),
\qquad
H_9(n)=(\rho_9(n),q_9(n))
\label{eq:H9}
\end{equation}
satisfait
\begin{equation}
H_9(n+9)=H_9(n)+(0,1).
\label{eq:H9-helix}
\end{equation}
La première coordonnée se ferme ; la seconde s'élève. Cette identité est la
version arithmétique minimale d'une vis.

Pour une fibre spectrale,
\[
n\longmapsto(\tau_\gamma^n,n)
\]
est une hélice sur \(S^1\). Ajouter la contraction \((5/9)^n\) produit une
spirale intérieure. Le radion participe aux deux lectures : une unité de phase
à la frontière et une unité orientée d'action dans le relèvement.

\paragraph{Rang de mémoire comme lecteur dimensionnel}

Supposons qu'une projection visible conserve la valeur de phase, mais perde
\(d\) cocycles de mémoire linéairement indépendants. Tout relèvement fidèle
nécessite au moins \(d\) coordonnées supplémentaires pour les reconstruire. Cela
motive une lecture précise de la mécanique dimensionnelle :
\begin{equation}
\begin{aligned}
\text{dimension émergente}
&=\text{rang minimal de mémoire indépendante}\\
&\phantom{=}\text{non publiable sur le feuillet visible}.
\end{aligned}
\label{eq:dimension-memory}
\end{equation}
On n'affirme pas que toute dimension physique soit un compteur. On affirme que, dans le
domaine TPK, chaque mémoire indépendante qui doit survivre impose une
coordonnée du relèvement.

\begin{narrativebox}
La dimension apparaît lorsque la projection ne peut plus conserver toute
l'histoire. Le cercle suffit à dire où se trouve la phase ; il ne
suffit pas à dire combien de fois elle est revenue, quel feuillet elle a transporté, quel intervalle
s'est contracté ou quelle retenue a survécu.
\end{narrativebox}

\paragraph{Assemblage de cylindres, de spirales et de trajectoires hélicoïdales}

Les cylindres coinductifs ne sont pas des tubes physiques ajoutés à une spirale. Ce sont
des intervalles emboîtés qui conservent les préfixes et resserrent la valeur. La spirale
décrit la contraction transversale ; l'hélice décrit la mémoire longitudinale ;
le cylindre décrit l'ensemble des valeurs compatibles à une profondeur
finie. À la limite, l'intersection des cylindres sélectionne une section
unique.

L'image spatiale réunit, sans confondre :
\[
\begin{array}{ccl}
\text{angle} &\leftrightarrow& \text{phase visible},\\
\text{hauteur} &\leftrightarrow& \text{retour/mémoire},\\
\text{rayon de cylindre} &\leftrightarrow& \text{incertitude de raffinement},\\
\text{feuillet} &\leftrightarrow& \text{orientation bidirectionnelle}.
\end{array}
\]
Le taux \(3^{-54}\) contracte le cylindre à chaque retour ; il ne s'ajoute pas à l'angle et
ne remplace pas la retenue réelle \(\epsR\).

\subsubsection{Clôtures de phase d'ordres \(2\pi\) et \(4\pi\)}

Dans la projection circulaire, \(2\pi\) restaure la phase. Dans un relèvement orienté
à deux feuillets, un tour peut se fermer sur le feuillet conjugué et nécessiter un
second tour pour restaurer l'orientation :
\[
U_{2\pi}=-I,\qquad U_{4\pi}=I.
\]
La structure coïncide avec la périodicité spinorielle lorsque l'on applique le
foncteur à une représentation \(SU(2)\) ; on n'identifie pas le relèvement TPK à
\(SU(2)\) sans cette application.

La lecture surfacique est compatible :
\[
\mathcal A(2\pi)=\hfull,\qquad
\mathcal A(4\pi)=2\hfull.
\]
Dans un ruban de Möbius, un circuit visible inverse le feuillet et deux circuits
restaurent l'orientation. Le \(2\pi/4\pi\) n'est pas une duplication accidentelle ;
c'est la différence entre retour de phase et retour de l'état enrichi.


\section{Ellipse orientée et variables conjuguées : incertitude minimale et oscillation électromagnétique}
\subsection{Variables symplectiques conjuguées et détermination de l'incertitude minimale}
% Fragmento público generado desde propietarios íntegros.
% Residencia: 52.6.1.
% Fuente: ../incoming/radion_monografia_20260827/manuscrito/sections/07_incertidumbre_oscilador.tex:1-111
\noindent La forme symplectique coordonne action, incertitude et dynamique oscillatoire sans modifier la généalogie de l'état.\par\medskip

\subsubsection{Réalisation de l'état dans un espace de Hilbert}

\paragraph{Matrice de covariance minimale}

Une fois construite l'ellipse d'action, elle peut se réaliser dans un couple
d'opérateurs conjugués
\[
[\widehat Q,\widehat P]=i\hret.
\]
On définit
\begin{equation}
V_\eta=\frac{\hret}{2}
\begin{pmatrix}
e^\eta&0\\[2pt]0&e^{-\eta}
\end{pmatrix}.
\label{eq:covariance}
\end{equation}
Son déterminant est
\begin{equation}
\det V_\eta=\frac{\hret^2}{4}.
\label{eq:cov-det}
\end{equation}
Par conséquent,
\begin{equation}
(\Delta Q)^2=\frac{\hret}{2}e^\eta,
\qquad
(\Delta P)^2=\frac{\hret}{2}e^{-\eta},
\qquad
\Delta Q\,\Delta P=\frac{\hret}{2}.
\label{eq:uncertainty}
\end{equation}
L'inégalité de Robertson--Schrödinger est saturée.

L'annihilateur comprimé
\begin{equation}
\widehat a_\eta=
\frac{e^{-\eta/2}\widehat Q+i e^{\eta/2}\widehat P}
{\sqrt{2\hret}}
\label{eq:a-eta}
\end{equation}
satisfait \([\widehat a_\eta,\widehat a_\eta^\dagger]=1\). Son vide a pour
fonction d'onde
\begin{equation}
\psi_\eta(Q)=
(\pi\hret e^\eta)^{-1/4}
\exp\left(-\frac{Q^2}{2\hret e^\eta}\right).
\label{eq:gaussian}
\end{equation}

\paragraph{Ellipse comme niveau de la fonctionnelle de Mahalanobis}

L’ensemble
\begin{equation}
\mathbb E_{\mathrm{act}}
=\{x\in\R^2:x^{\mathsf T}V_\eta^{-1}x\le4\}
\label{eq:mahalanobis}
\end{equation}
a pour demi-axes
\begin{equation}
a_S=2\Delta Q,\qquad b_S=2\Delta P.
\label{eq:two-sigma}
\end{equation}
Son aire est
\[
4\pi\sqrt{\det V_\eta}=2\pi\hret=\hfull.
\]
L'ellipse HMT construite précédemment coïncide exactement avec le contour à deux
écarts de l'état gaussien minimal.

\begin{proposition}[Chaîne unique du quotient de forme]
Après avoir déclaré les interfaces spectrale, symplectique et de Hilbert,
\begin{equation}
\boxed{
e^{-\eta}
=\sqrt{\frac{A-C^*}{A+C^*}}
=\frac{b_1}{a_1}
=\frac{b_S}{a_S}
=\frac{\Delta P}{\Delta Q}.}
\label{eq:shape-chain-hilbert}
\end{equation}
\end{proposition}

L'ellipse d'action correspond au niveau classique \(H=\hret\omega\) et au
contour \(2\sigma\) du vide comprimé. Elle ne doit pas être confondue avec l'énergie
moyenne du vide, \(\hret\omega/2\). Cette différence d'un facteur deux est un
contrôle physique obligatoire.

\paragraph{Conditions d'invariance de l'ellipse comme objet opératoire}

Dans la formulation conventionnelle, une ellipse d'incertitude peut être considérée comme
un outil graphique pour représenter une matrice de covariance. Ici, l'ordre
causal est inversé : l'ellipse positive existe déjà comme publication de
\((A,C^*,\hret)\) ; l'espace de Hilbert la reconnaît au moyen de
\eqref{eq:covariance}. La réalité ontologique revendiquée ne signifie pas qu'une
particule classique parcourt physiquement un ovale dans le plan \((Q,P)\). Elle signifie
que l'état conserve une forme positive, une action totale et un quotient
d'orientation avant qu'une théorie probabiliste ne les interprète.

La fonction d'onde ne crée pas la forme ; elle la représente. L'incertitude ne crée pas
\(\hbar\) ; elle publie l'impossibilité de comprimer simultanément les deux
directions sans modifier le produit symplectique.

\begin{narrativebox}
La lecture probabiliste est une réalisation possible de l'ellipse, non sa cause.
L'objet premier est la conservation conjointe du produit et de l'orientation. La
probabilité apparaît lorsque cette géométrie est publiée comme covariance d'un
état sur Hilbert.
\end{narrativebox}


\subsection{Échange \(\mu\leftrightarrow\varepsilon\), inversion \(Z\leftrightarrow Z^{-1}\) et conservation de la vitesse causale}
% Fragmento público generado desde propietarios íntegros.
% Residencia: 52.6.2.
% Fuente: ../incoming/radion_monografia_20260827/manuscrito/sections/07_incertidumbre_oscilador.tex:156-197
\paragraph{Échange énergétique entre les secteurs électrique et magnétique}

Les deux termes de \eqref{eq:HLC} deviennent
\begin{equation}
U_E(\theta)=\hret\omega\cos^2\theta,
\qquad
U_B(\theta)=\hret\omega\sin^2\theta.
\label{eq:energy-exchange}
\end{equation}
Par conséquent,
\[
U_E+U_B=\hret\omega.
\]
La polarité électrique--magnétique ne consiste pas à étiqueter arbitrairement
les foyers. Elle consiste en deux réserves conjuguées qui s'échangent pendant
l'orbite tout en maintenant constantes la fréquence et l'action totale.

Le quotient commun se prolonge :
\begin{equation}
\boxed{
e^{-\eta}
=\frac{b_S}{a_S}
=\frac{\Delta P}{\Delta Q}
=\frac{Z_\eta}{Z_*}
=\sqrt{\frac{A-C^*}{A+C^*}}.}
\label{eq:shape-chain-LC}
\end{equation}
Les produits conservés sont
\[
a_Sb_S=2\hret,\qquad
\det V_\eta=\frac{\hret^2}{4},
\qquad
L_\eta C_\eta=\omega^{-2}.
\]

\begin{exactbox}
L'oscillateur LC est la réalisation physique minimale de la double lecture : le
produit conserve l'action et la fréquence ; le quotient conserve l'orientation et
l'impédance. Un cycle complet échange électricité et magnétisme sans
perdre l'aire \(\oint P\,\diff Q=\hfull\).
\end{exactbox}


\subsection{Couple inductif--capacitif \((L_\eta,C_\eta)\) : invariance de \(L_\eta C_\eta=\omega^{-2}\) et orientation de l'impédance}
% Fragmento público generado desde propietarios íntegros.
% Residencia: 52.6.3.
% Fuente: ../incoming/radion_monografia_20260827/manuscrito/sections/07_incertidumbre_oscilador.tex:112-155
\subsubsection{Oscillateur inductif--capacitif associé au couple angulaire}

\paragraph{Construction indépendante des 2 réservoirs conjugués}

Étant donnés \(\omega>0\) et une impédance de référence \(Z_*>0\), nous définissons
\begin{equation}
L_\eta=\frac{Z_*}{\omega}e^{-\eta},
\qquad
C_\eta=\frac{1}{Z_*\omega}e^\eta.
\label{eq:LC-def}
\end{equation}
Alors
\begin{equation}
L_\eta C_\eta=\omega^{-2},
\qquad
\sqrt{\frac{L_\eta}{C_\eta}}=Z_*e^{-\eta}.
\label{eq:LC-invariants}
\end{equation}
La fréquence reste fixe ; l'impédance conserve l'orientation de la
forme.

L'hamiltonien est
\begin{equation}
H_{LC}=\frac{q^2}{2C_\eta}+\frac{\Phi^2}{2L_\eta}.
\label{eq:HLC}
\end{equation}
Avec les variables normalisées
\begin{equation}
Q=\sqrt{Z_*}\,q,\qquad
P=\frac{\Phi}{\sqrt{Z_*}},
\label{eq:LC-map}
\end{equation}
on obtient
\begin{equation}
H_{LC}=\frac\omega2
\left(e^{-\eta}Q^2+e^\eta P^2\right).
\label{eq:HLC-QP}
\end{equation}
Sur \eqref{eq:elipse-param},
\begin{equation}
H_{LC}=\hret\omega.
\label{eq:LC-level}
\end{equation}

% Fuente: ../incoming/radion_monografia_20260827/manuscrito/sections/07_incertidumbre_oscilador.tex:198-239
\paragraph{Oscillateur LC comme unité dynamique élémentaire}

L'oscillateur harmonique traverse presque toute la physique parce qu'il linéarise le
voisinage d'un équilibre, décompose les champs en modes, organise les quanta et
permet de lire la propagation comme une superposition de phases. Dans le radion, cette
universalité reçoit une généalogie concrète : le mode ne commence pas par une
équation différentielle externe, mais par une forme positive dont l'action de tour est
fixe. La dynamique conventionnelle apparaît lorsqu'on choisit une horloge \(\omega\) et un
transducteur \(Z_*\).

La fréquence ne modifie pas l'ellipse d'action : elle change l'énergie du niveau
au moyen de \(E=\hret\omega\). Cette séparation est cruciale. La géométrie fixe
l'action ; l'horloge convertit l'action en énergie.

\subsubsection{Réversion temporelle et propagations conjuguées}

Soit \(H=H^*\) et soit \(J_E\) une structure complexe avec \(J_E^2=-I\).
Supposons un opérateur involutif \(C\) tel que
\[
CJ_EC^{-1}=-J_E,\qquad CHC^{-1}=H.
\]
L'évolution
\begin{equation}
U(t)=\exp\left(-\frac{J_EHt}{\hret}\right)
\label{eq:evolution-J}
\end{equation}
satisfait exactement
\begin{equation}
CU(t)C^{-1}=U(-t).
\label{eq:time-reversal}
\end{equation}
Telle est la forme opératoire minimale de la bidirectionnalité : la conjugaison
échange les prolongations temporellement opposées.

La théorie de l'absorbeur de Wheeler--Feynman et la lecture des antiparticules
comme trajectoires temporellement conjuguées offrent une réalisation historique
ultérieure. L'identité \eqref{eq:time-reversal} est démontrée ; la transformer
en une théorie complète des propagateurs avancés et retardés exige de déclarer
l'action, les conditions de frontière, la causalité et les interactions de l'absorbeur.
La monographie conserve l'intuition physique sans transformer une correspondance
en une égalité de théories.



\section{3 régimes du TRIT : polarités électrique--magnétique et élastique--rotationnelle}
\subsection{Coordination des réalisations elliptique, parabolique et hyperbolique}
% Fragmento público generado desde propietarios íntegros.
% Residencia: 52.7.1.
% Fuente: ../incoming/radion_monografia_20260827/manuscrito/sections/08_complejo_electromagnetismo.tex:1-58
\noindent Les 3 régimes du TRIT admettent des réalisations complexe, électromagnétique et mécanique avec des types différenciés.\par\medskip

\subsubsection{Opérateur quadratique de racine interne de \(-1\)}

\paragraph{Représentation exponentielle du régime elliptique \(J^2=-I\)}

Avec \(J^2=-I\),
\begin{equation}
T_{\mathrm{ell}}(x,y)=e^{-xI-yJ}
=e^{-x}(\cos y\,I-\sin y\,J).
\label{eq:elliptic-exp}
\end{equation}
Sous l'identification générée par \(J\),
\[
T_{\mathrm{ell}}(x,y)\longleftrightarrow e^{-x-iy}.
\]
Le scalaire \(x\) contrôle la dilatation ; \(y\) contrôle la rotation orientée. La forme
\(a+bi\) se reconnaît comme élasticité plus rotation, non comme deux quantités
collées arbitrairement.

\paragraph{Représentation exponentielle du régime parabolique \(J^2=0\)}

Avec \(N^2=0\), l'exponentielle se tronque :
\begin{equation}
T_{\mathrm{par}}(x,y)=e^{-x}(I-yN).
\label{eq:parabolic-exp}
\end{equation}
Le régime décrit le seuil où il n'y a ni rotation périodique ni ouverture
exponentielle. Dans la lecture temporelle active, c'est le présent euclidien, mais
il conserve la mémoire des branches qui le délimitent.

\paragraph{Représentation exponentielle du régime hyperbolique \(J^2=+I\)}

Avec \(R^2=I\) et
\(P_\pm=(I\pm R)/2\),
\begin{equation}
T_{\mathrm{hyp}}(x,y)=e^{-xI-yR}
=q_+P_++q_-P_-,
\qquad q_\pm=e^{-x\mp y}.
\label{eq:hyperbolic-exp}
\end{equation}
Les invariants élémentaires sont
\begin{equation}
q_+q_-=e^{-2x},
\qquad
\frac{q_-}{q_+}=e^{2y}.
\label{eq:q-product-ratio}
\end{equation}
Le produit conserve le mode commun ; le quotient conserve la séparation
orientée.

Appliqué au couple angulaire,
\[
x=\frac{\pi A}{180},\qquad y=\frac{\pi C^*}{180}.
\]
C'est pourquoi \(A\) est la coordonnée commune de la publication et \(C^*\) son
contraste orienté. Ce n'est qu'après l'application d'un lecteur physique qu'ils sont reconnus
comme pôles électriques et magnétiques.


\subsection{Séparation typée des polarités constitutives et mécaniques}
% Fragmento público generado desde propietarios íntegros.
% Residencia: 52.7.2.
% Fuente: ../incoming/radion_monografia_20260827/manuscrito/sections/08_complejo_electromagnetismo.tex:59-158
\subsubsection{Opérateur constitutif sur les canaux \(90/120\)}

\paragraph{Décomposition en composantes commune et orientée}

On définit
\begin{align}
\mathcal E
&=\log\bigl[(1-q_+^{90})(1-q_-^{90})\bigr]
-\log\bigl[(1-q_+^{120})(1-q_-^{120})\bigr],
\label{eq:Ecommon}\\
\mathcal B
&=\log\frac{1-q_-^{90}}{1-q_+^{90}}
-\log\frac{1-q_-^{120}}{1-q_+^{120}}.
\label{eq:Borient}
\end{align}
Sous \(C^*\mapsto-C^*\), \(q_+\leftrightarrow q_-\) s'échangent. Par
conséquent, \(\mathcal E\) est pair et \(\mathcal B\) est impair.

Les lois constitutives normalisées sont
\begin{equation}
\widehat\mu=e^{\mathcal E+\mathcal B},
\qquad
\widehat\varepsilon=e^{\mathcal E-\mathcal B},
\qquad
\widehat Z=e^{\mathcal B},
\qquad
\widehat c=e^{-\mathcal E}.
\label{eq:constitutives}
\end{equation}
On vérifie exactement
\begin{equation}
\widehat Z=\sqrt{\frac{\widehat\mu}{\widehat\varepsilon}},
\qquad
\widehat c=(\widehat\mu\widehat\varepsilon)^{-1/2}.
\label{eq:constitutive-identities}
\end{equation}
La conjugaison produit
\begin{equation}
C^*\mapsto-C^*:\qquad
\widehat\mu\leftrightarrow\widehat\varepsilon,
\quad
\widehat Z\leftrightarrow\widehat Z^{-1},
\quad
\widehat c\mapsto\widehat c.
\label{eq:EM-conjugation}
\end{equation}

\begin{theorem}[Polarité constitutive]
Le lecteur \(90/120\) transforme le mode commun et le contraste orienté du
couple \((A,C^*)\) en deux lois constitutives conjuguées. L'inversion de feuillet
échange perméabilité et permittivité, inverse l'impédance et conserve la
vitesse normalisée.
\end{theorem}

Telle est la formulation précise de l'intuition \enquote{partie électrique et
partie magnétique}. On n'écrit ni \(A=E\) ni \(C^*=B\). On écrit l'application
\[
(A,C^*)\to(x,y)\to(q_+,q_-)
\to(\mathcal E,\mathcal B)
\to(\widehat\varepsilon,\widehat\mu,\widehat Z,\widehat c).
\]

\paragraph{Réalisation élastique--rotationnelle de la polarité constitutive}

La même information peut être publiée dans le régime complexe au moyen de
\begin{equation}
Z_{\mathrm{em}}=e^{\mathcal E/2}
\left(\cos\frac{\mathcal B}{2}\,I
-\sin\frac{\mathcal B}{2}\,J\right).
\label{eq:Zem}
\end{equation}
Le facteur \(e^{\mathcal E/2}\) est une dilatation élastique ; le facteur généré
par \(J\) est une rotation. L'écriture complexe n'ajoute pas une dimension
imaginaire mystérieuse à un monde réel préalablement donné : elle enregistre deux modes
opératoires que le TRIT a déjà séparés.

\paragraph{Transduction de l'orientation en charge effective}

Le feuillet \(\sigma\) du radion fixe une orientation avant d'introduire une
unité de charge. La mécanique dimensionnelle dispose ensuite le transducteur
\[
e_{\mathrm{int}}=\sqrt{\frac{2\alphah\hfull}{Z_{\mathrm{int}}}},
\qquad
Q=n_Qe_{\mathrm{int}}.
\]
La chaîne causale est
\[
\text{feuillet}\to\text{orientation}\to\text{constitutive}
\to\text{transducteur dimensionnel}\to\text{charge physique}.
\]
Appeler l'unité \emph{radion} n'identifie pas prématurément l'orientation
au coulomb ; cela rend visible le lieu où le signe de charge peut résider.

\begin{narrativebox}
L'électricité et le magnétisme ne sont pas collés au cercle depuis l'extérieur. Ils émergent
lorsque le même état est lu à travers le produit et le quotient : le produit
retient le mode commun, le quotient retient l'orientation. L'oscillateur LC
réalise l'échange ; le lecteur \(90/120\) publie les lois constitutives ; la
carte dimensionnelle attribue des unités.
\end{narrativebox}


\section{Séparation des publications \(90/120\) : incidence, registre, spectre, noyau de Barbero et tours}
\subsection{Séparation typée des publications du secteur \(90/120\)}
% Fragmento público generado desde propietarios íntegros.
% Residencia: 52.8.1.
% Fuente: ../incoming/radion_monografia_20260827/manuscrito/sections/09_sector_90120.tex:176-244
\subsubsection{Hiérarchie harmonique globale et semi-groupe \(\langle90,120\rangle\)}

Pour éviter une collision de signes, nous fixons
\[
R_{\mathrm{tow}}:=-R_{\mathrm{ch}}.
\]
Alors
\begin{equation}
T=e^{-xI+yR_{\mathrm{tow}}}
=q_-P_+^{\mathrm{tow}}+q_+P_-^{\mathrm{tow}}.
\label{eq:tower-operator}
\end{equation}
Les tours paires et impaires sont
\begin{align}
c_n&=\operatorname{Tr}(T^n)
=q_-^n+q_+^n
=2e^{-nx}\cosh(ny),
\label{eq:cn}\\
s_n&=\operatorname{Tr}(R_{\mathrm{tow}}T^n)
=q_-^n-q_+^n
=2e^{-nx}\sinh(ny).
\label{eq:sn}
\end{align}
Sous \(C^*\mapsto-C^*\), \(c_n\) est pair et \(s_n\) est impair. On a
\begin{equation}
c_n^2-s_n^2=4e^{-2nx},
\label{eq:tower-invariant}
\end{equation}
et les deux suites satisfont la récurrence d'ordre deux déterminée par les
valeurs propres \(q_\pm\).

Le semi-groupe
\begin{equation}
\langle90,120\rangle=30\langle3,4\rangle
\label{eq:semigroup}
\end{equation}
est la hiérarchie harmonique d'un unique opérateur bicouche. Ce n'est pas une liste de huit
corrections latérales. Sa complétude signifie que tout rapport positif
admet un développement convergent ; elle ne signifie pas qu'un développement obtenu à
partir de la cible soit un sélecteur physique prospectif.

\subsubsection{Classification typée et séparation des réalisations \(90/120\)}

\begin{longtable}{>{\bfseries}p{.19\textwidth} p{.24\textwidth} p{.45\textwidth}}
\toprule
Objet & Opération & Ce qu'il conserve / pourquoi ce n'est pas un autre objet\\
\midrule
\endfirsthead
\toprule
Objet & Opération & Ce qu'il conserve / pourquoi ce n'est pas un autre objet\\
\midrule
\endhead
Incidence \(N_6,N_8\) & dénombrement fini sur les phases actives & cardinalité et orientation ; produit \(90,120\), non une série ;\\
\(-1/12\) & différence normalisée \eqref{eq:minus-one-twelve} & scalaire d'incidence ; non le poids de Barbero ;\\
\(E_{\mathrm{dod}}\) & extracteur harmonique de la roue & phase dodécaphasique ; non la retenue réelle ;\\
\(h=C^*/6\) & reconstruction de six paires & ouverture par paire ; non une normalisation métrologique ;\\
\(F_{\mathrm{spec}}\) & queue logarithmique \(j\ge2\) & publication spectrale complète ; diffère de \(\epsR^\circ\) ;\\
\(\epsone\) & soustraction APP entre \(S_T\) et \(S_E\) & raccord réel de deux relèvements ; non une série de Lambert ;\\
\(K_{6\to8}\) & \(12\Delta S_{90}-\Delta S_{120}\) & saut marqué et évaluation conjuguée ; coefficient de pôle \(1/24\) vérifié ; sortie de Barbero ;\\
\(12:-1\) & douze secteurs et une couture orientée & poids du noyau ; certificat ultérieur \eqref{eq:pole-condition} ; distinct de \(-1/12\) ;\\
\(\langle90,120\rangle\) & semi-groupe des puissances de \(T\) & hiérarchie harmonique globale ; non un résidu par espèce.\\
\bottomrule
\end{longtable}

\begin{thesisbox}
Tous ces objets descendent du même état APP--TRIT--TPK. Cette généalogie
commune explique pourquoi ils partagent des nombres et des symétries. Leur typage explique pourquoi
ils ne peuvent pas se remplacer mutuellement.
\end{thesisbox}

\subsection{Conservation de l'incidence et de l'orientation entre le noyau tritique et ses réalisations ultérieures}
% Fragmento público generado desde propietarios íntegros.
% Residencia: 52.8.2.
% Fuente: ../incoming/radion_monografia_20260827/manuscrito/sections/09_sector_90120.tex:1-30
\noindent La hiérarchie typée sépare incidence, registre, évaluation spectrale, fonctionnelle de Barbero et tours harmoniques.\par\medskip

\subsubsection{Incidence tritique, registre dodécaphasique et lecture sexadique}

\paragraph{Dérivation combinatoire primaire de \(90\) et \(120\)}

Les dénombrements
\[
90=6\binom62,
\qquad
120=8\binom62
\]
naissent du sélecteur de phase et de ses incidences actives. L'orientation
normalisée satisfait
\begin{equation}
\frac{30}{120}-\frac{30}{90}=-\frac1{12}.
\label{eq:minus-one-twelve}
\end{equation}
Ce scalaire est un résultat d'incidence. Ce n'est pas le poids \(12:-1\) du
noyau de Barbero et ce n'est pas une retenue dodécaphasique.

L'extracteur dodécaphasique peut être publié au moyen d'une combinaison
d'harmoniques telle que
\[
E_{\mathrm{dod}}(\phi)=12\cos(3\phi)-\cos(4\phi).
\]
Son domaine est la phase de la roue ; sa sortie est un sélecteur harmonique. Bien que
les nombres \(12\) et \(-1\) apparaissent, il n'opère pas sur des séries de Lambert et ne
produit pas \(\epsR\).


\subsection{Origine dodécaphasique du poids \(12:-1\) et vérification diophantienne de son coefficient de pôle}
% Fragmento público generado desde propietarios íntegros.
% Residencia: 52.8.3.
% Fuente: ../incoming/radion_monografia_20260827/manuscrito/sections/09_sector_90120.tex:115-175
\subsubsection{Fonctionnelle hexade--octade de Barbero--Immirzi}

\paragraph{Séries orientées et saut d'incidence}

On définit
\begin{equation}
S_{90}(q)=\frac{q^{90}}{1-q^{270}},
\qquad
S_{120}(q)=\frac{q^{120}}{1-q^{360}},
\label{eq:barbero-series}
\end{equation}
et
\[
\Delta S_m=S_m(q_-)-S_m(q_+).
\]
La fonctionnelle d'inclusion marquée est
\begin{equation}
K_{6\to8}=12\Delta S_{90}-\Delta S_{120}.
\label{eq:barbero-núcleo}
\end{equation}

\paragraph{Vérification diophantienne du poids dodécaphasique \(12:-1\)}

La chaîne fondamentale dodécaphasique et son lecteur de retour, construits dans
\eqref{eq:l9s102:barbero-cadena-dodecafasica}--%
\eqref{eq:l9s102:barbero-peso-generado}, produisent douze termes de type
\(S_{90}\) et un terme de frontière de type \(-S_{120}\). Le couple orienté
est déterminé avant le calcul du coefficient de pôle.

Considérons \(aS_{90}-bS_{120}\). La vérification du coefficient de
\((1-q)^{-1}\) prend la forme
\begin{equation}
\frac{a}{270}-\frac{b}{360}=\frac1{24}.
\label{eq:pole-condition}
\end{equation}
Multiplier par \(1080\) donne
\[
4a-3b=45.
\]
L'unique couture du tour fondamental a une multiplicité absolue de un.
Sa vérification diophantienne correspond au contre-terme entier minimal
\(|b|=1\). Pour \(b=1\), \(a=12\) ; pour \(b=-1\), \(a=21/2\) n'est pas entier.
Par conséquent,
\begin{equation}
\boxed{a:b=12:1,\qquad aS_{90}-bS_{120}=12S_{90}-S_{120}.}
\label{eq:12minus1}
\end{equation}

\begin{proposition}[Unicité de la vérification diophantienne]
Le couple généré \((12,1)\) est l'unique solution entière de
\eqref{eq:pole-condition} avec canal principal positif et contre-terme
minimal \(|b|=1\).
\end{proposition}

La réalisation hexade--octade conserve les canaux \(90/120\)
préalablement construits ; la chaîne dodécaphasique apporte les multiplicités
et l'orientation de frontière ; le lecteur de retour attribue les séries aux
générateurs. Ces opérations conjointes produisent \(12S_{90}-S_{120}\).
L'équation \eqref{eq:pole-condition}, la positivité et la minimalité
conservent leur fonction de vérification arithmétique ultérieure.

La chaîne de Barbero est ultérieure :
\[
(A,C^*)\to(x,y)\to q_\pm
\to K_{6\to8}\to\gamma_{\mathrm{BI}},
\]
avec l'évaluation
\[
\gamma_{\mathrm{BI}}=0.274007672694459\ldots.
\]
Barbero--Immirzi ne génère pas rétroactivement \(C^*\), \(\Delta_4\) ni
\(\epsR\).


\subsection{Extraction de la composante finie postérieure au système dodécaphasique orienté \(R_{12}\) : \(C^*\mapsto C^*/6\), reste \(j\ge2\) et reconstruction de \(6\) sections}
% Fragmento público generado desde propietarios íntegros.
% Residencia: 52.8.4.
% Fuente: ../incoming/radion_monografia_20260827/manuscrito/sections/09_sector_90120.tex:31-54
\paragraph{Extraction de la composante \(C^*/6\) du registre dodécaphasique}

Les douze secteurs sont regroupés en six paires opposées :
\begin{equation}
p_i=e_i+e_{i+6},\qquad
a_i=e_i-e_{i+6},\qquad i=1,\ldots,6.
\label{eq:sexadic-pairs}
\end{equation}
La reconstruction a pour profil \(1+5+6=12\). L'ouverture par paire est
\begin{equation}
h=\frac{C^*}{6}.
\label{eq:h-C6}
\end{equation}
Ce n'est ni une normalisation extérieure ni le sixième d'une cible : c'est l'indice
interne de la reconstruction sexadique du registre \(\Rstate\).

On définit
\begin{equation}
q_{h,-}=\exp\left[-\frac\pi{180}(A-h)\right],
\qquad
q_{h,+}=\exp\left[-\frac\pi{180}(A+h)\right].
\label{eq:q-h}
\end{equation}


\subsection{Décomposition quantitative des publications de retenue et spectrale : \(F_{\mathrm{spec}}=\varepsilon_R^\circ+\chi_R\)}
% Fragmento público generado desde propietarios íntegros.
% Residencia: 52.8.5.
% Fuente: ../incoming/radion_monografia_20260827/manuscrito/sections/09_sector_90120.tex:55-114
\subsubsection{Publication spectrale de la queue d'ordres \(j\ge2\)}

\paragraph{Définition de la fonctionnelle spectrale}

Pour \(m\in\{90,120\}\), soit
\begin{equation}
L_m^{(2)}(q)=\sum_{j\ge2}\frac{q^{mj}}{j}
=-\log(1-q^m)-q^m.
\label{eq:L2}
\end{equation}
La lecture d'une paire est
\begin{equation}
T_h^{(2)}=\frac{180}{\pi}
\Bigl[
L_{90}^{(2)}(q_{h,-})-L_{90}^{(2)}(q_{h,+})
-L_{120}^{(2)}(q_{h,-})+L_{120}^{(2)}(q_{h,+})
\Bigr].
\label{eq:T-h}
\end{equation}
La reconstruction des six paires donne
\begin{equation}
F_{\mathrm{spec}}=6T_h^{(2)}
=5.1499235153094574410\ldots\times10^{-8}\,{}^\circ.
\label{eq:Fspec}
\end{equation}

La retenue du radion est
\[
F_{\mathrm{acarreo}}=\epsR^\circ
=5.1383016758575282072\ldots\times10^{-8}\,{}^\circ.
\]
Ils ne coïncident pas :
\begin{equation}
\chi_R:=F_{\mathrm{spec}}-F_{\mathrm{acarreo}}
=1.16218394519292338\ldots\times10^{-10}\,{}^\circ.
\label{eq:chiR}
\end{equation}

L’égalité
\[
F_{\mathrm{acarreo}}=F_{\mathrm{spec}}-\chi_R
\]
est un changement de carte valable une fois les deux sorties construites. Il ne doit pas
être présenté comme génération indépendante de \(\epsR^\circ\) si \(\chi_R\) est
défini par la soustraction elle-même. La queue est clôturée et significative ; ce qui est faux,
c'est de l'identifier directement à la retenue.

\begin{falsifier}
Après conversion de la queue exprimée en degrés dans l'unité angulaire interne,
\[
\frac{\pi}{180^\circ}F_{\mathrm{spec}}-\epsone
=2.0283936357433839\ldots\times10^{-12}\ne0.
\]
Par conséquent, ni le log complet ni la queue \(j\ge2\), même après application de la
carte d'unités, ne sont \(\epsone\). Une
seconde factorisation spectrale autonome nécessiterait de générer \(\chi_R\) sans
utiliser la différence comme définition ; la clôture APP--TPK du radion ne dépend pas
de cette promotion.
\end{falsifier}



\section{Holonomie orientée et géométrie physique : réalisations ultérieures de Minkowski, Hodge, Lorentz et Cartan--Holst}
% Fragmento público generado desde propietarios íntegros.
% Residencia: 52.9.
% Fuente: ../incoming/radion_monografia_20260827/manuscrito/sections/10_minkowski_hodge_lorentz.tex:1-55
\noindent L'holonomie orientée admet des réalisations causales et géométriques ultérieures dans Minkowski, Hodge, Lorentz et Cartan--Holst.\par\medskip

\paragraph{Construction de la métrique lorentzienne à partir du bloc APP}

\paragraph{Matrice APP de paramètres \(7\) et \(2\)}

La matrice centrale
\begin{equation}
B_c=\begin{pmatrix}7&2\\2&7\end{pmatrix}
=9P_++5P_-
\label{eq:Bc}
\end{equation}
est normalisée par son déterminant :
\begin{equation}
M_c=\frac{B_c}{\sqrt{45}}
=\exp\left(\frac12\log\frac95\,R\right)
\in SO^+(1,1).
\label{eq:Mc}
\end{equation}
La rapidité interne est
\[
\eta_c=\frac12\log\frac95.
\]
La publication projective associée satisfait
\begin{equation}
P_B^n=P_++\left(\frac59\right)^nP_-.
\label{eq:PB-contract}
\end{equation}
Le même rapport \(5/9\) qui apparaît comme contraction spectrale est reconnu ici
comme séparation d'un mode stable et d'un autre contractant. L'identité est
opératoire ; elle ne transforme pas toutes les occurrences de \(5/9\) en la même application.

\paragraph{Incidence \(6\to8\) et quotient de signature causale}

Soit \(M\) l'incidence typée entre six faces et huit sommets. La relation
spectrale récupérée est
\begin{equation}
MM^{\mathsf T}=12P_u+4P_-,
\label{eq:MMt}
\end{equation}
où \(P_u\) projette sur le mode uniforme et \(P_-\) sur le sous-espace
orienté. Le quotient survivant de dimension quatre se décompose comme
\begin{equation}
Q_4\simeq\R u\oplus E_-.
\label{eq:radion-Q4}
\end{equation}
En déclarant négatif le mode uniforme et positifs les trois modes orientés,
on obtient
\begin{equation}
B_{\mathrm{caus}}=\operatorname{diag}(-1,1,1,1).
\label{eq:minkowski}
\end{equation}
L'écran de Minkowski ne sélectionne pas l'état HMT. C'est la réalisation
causale ultérieure du quotient que l'incidence a déjà construit.


\subsection{Dualité de Hodge sur les \(2\)-formes lorentziennes}
% Fragmento público generado desde propietarios íntegros.
% Residencia: 52.9.1.
% Fuente: ../incoming/radion_monografia_20260827/manuscrito/sections/10_minkowski_hodge_lorentz.tex:56-100
\subsubsection{Dualité de Hodge et opérateur de quart de tour}

La signature \((-+++ )\) et une orientation étant fixées, l'étoile de Hodge sur
les deux-formes satisfait
\begin{equation}
\star:\Lambda^2Q_4^*\to\Lambda^2Q_4^*,
\qquad \star^2=-I.
\label{eq:hodge-star}
\end{equation}
De cette manière, le quart de tour qui, dans le plan de phase, s'exprimait au moyen de
\(J^2=-I\) réapparaît dans l'espace des champs comme structure complexe sur
\(\Lambda^2\). Les domaines sont différents, mais la généalogie d'orientation
est conservée par le foncteur.

Soit \(A_{\mathrm{em}}\) un potentiel et
\begin{equation}
F=\diff A_{\mathrm{em}}.
\label{eq:F-dA}
\end{equation}
La décomposition de \(F\) par rapport à un observateur temporel produit des champs
électrique et magnétique ; \(\star F\) échange leurs rôles avec le signe fixé
par l'orientation et la signature. La polarité constitutive de
\eqref{eq:EM-conjugation} acquiert ainsi une réalisation différentielle.

\paragraph{Chaîne de transduction électromagnétique et causale}

Le transport qui doit rester visible est
\begin{equation}
\Gamma_9
\longrightarrow\text{hélice de mémoire}
\longrightarrow(Q_4,B_{\mathrm{caus}},\star)
\longrightarrow A_{\mathrm{em}}
\longrightarrow F=\diff A_{\mathrm{em}}
\longrightarrow
m\nabla_u u=q(\iota_uF)^\sharp.
\label{eq:lorentz-chain}
\end{equation}
Chaque flèche ajoute de la structure : l'holonomie conserve le retour ; le quotient fixe
la causalité ; Hodge fixe la dualité ; le potentiel fixe le champ ; le couplage
avec la charge produit la dynamique.

L'équation finale est la force de Lorentz sous forme géométrique. L'hélice TPK
est son antécédent de mémoire, non une orbite physique déjà mue par un champ. La
dynamique n'apparaît qu'à l'introduction de \(F\), \(q\) et de la connexion \(\nabla\).


\subsection{Séparation entre mémoire torsionnelle interne et réalisation géométrique de Cartan--Holst}
% Fragmento público generado desde propietarios íntegros.
% Residencia: 52.9.2.
% Fuente: ../incoming/radion_monografia_20260827/manuscrito/sections/10_minkowski_hodge_lorentz.tex:101-186
\subsubsection{Réalisation lorentzienne des trajectoires hélicoïdales}

Dans un champ magnétique uniforme, une particule de charge \(q\), de masse \(m\) et de
facteur relativiste \(\gamma\) possède
\begin{equation}
\omega_c=\frac{|q|B}{\gamma m},
\qquad
r_L=\frac{p_\perp}{|q|B},
\qquad
\Delta z=\frac{2\pi p_\parallel}{|q|B}.
\label{eq:lorentz-helix}
\end{equation}
La phase cyclotronique se ferme tandis que le mouvement parallèle avance : l'orbite
est hélicoïdale. Le parallélisme avec \eqref{eq:H9-helix} est structurel et peut
être formalisé par un transducteur qui préserve phase, orientation et pas. On
n'identifie pas le quotient \(q_9\) à la coordonnée spatiale \(z\) ; on reconnaît
la même architecture de retour plus avance.

La force de Lorentz ne se dérive pas de l'image d'une spirale par ressemblance.
Elle se dérive en complétant \eqref{eq:lorentz-chain}. L'image aide à comprendre
pourquoi phase circulaire et mémoire longitudinale sont compatibles ; la forme de
champ détermine la dynamique.

\subsubsection{Confluence entre action, aire, information et gravité}

\paragraph{Réalisation Cartan--Holst de l'holonomie}

La promotion gravitationnelle s'exprime au moyen d'un morphisme d'holonomies
\begin{equation}
\operatorname{Hol}_{\mathcal A}
\bigl(\mathfrak R_{\mathrm{grav}}(\gamma)\bigr)
=\rho_C\bigl(\operatorname{Hol}_{\TPK}(\gamma)\bigr).
\label{eq:holonomy-gravity}
\end{equation}
Le membre de droite contient la mémoire de chemin ; \(\rho_C\) la réalise dans une
connexion gravitationnelle. La torsion de Cartan
\[
T^I=D_\omega e^I
\]
apparaît dans le codomaine. La torsion globale \(\Delta_4\) est une coordonnée
en amont qui peut alimenter l'application, mais n'est pas renommée \(T^I\) sans
\eqref{eq:holonomy-gravity}.

\paragraph{Transduction de l'aire d'action en aire informationnelle}

Le radion fixe un quantum d'action par phase. Le lecteur hexade--octade de
Barbero utilise ultérieurement les canaux \(q_\pm\), l'incidence \(90/120\) et
la fonctionnelle \eqref{eq:barbero-núcleo}. La direction causale est
\[
(A,C^*,\hret)\to q_\pm\to K_{6\to8}
\to\gamma_{\mathrm{BI}}\to\text{aire/information}.
\]
L'alphabet tritique apporte l'unité d'entropie \(\log3\) et une coupure de
\(81\) canaux apporte \(81\log3\). Cette entropie ne se confond ni avec
\(\log2\), ni avec \(\log\varphi\), ni avec une entropie microcanonique, ni avec une
entropie de von Neumann sans spécifier l'état.

La synthèse physique est affirmative mais typée : le radion apporte l'unité
d'action orientée qu'une holonomie peut transporter ; Barbero publie la
normalisation d'aire ; l'écran causal et Hodge permettent de réaliser des champs ;
Cartan--Holst reçoit l'holonomie en géométrie gravitationnelle.

\begin{statusbox}
L'équation scalaire du radion n'est pas à elle seule une équation d'Einstein. Sa
contribution à la gravité consiste à être une section exacte de l'architecture
d'action, d'orientation et d'holonomie que la réalisation gravitationnelle consomme.
\end{statusbox}

\subsubsection{Transformation de Wick comme changement de lecteur analytique}

Dans la formulation conventionnelle, la rotation de Wick relie une variable
temporelle lorentzienne à une coordonnée euclidienne par un prolongement
complexe. Dans HMT, la racine \(J^2=-I\), l'opérateur involutif \(R^2=I\) et le seuil
\(N^2=0\) existent déjà auparavant. Le changement de lecteur
\[
J\longleftrightarrow R
\]
transporte un couple commun/orienté entre publication elliptique et
publication hyperbolique. Le prolongement complexe reconnaît ensuite cette
opération en coordonnées analytiques.

Cela explique pourquoi cercle, ellipse d'action et intérieur hyperbolique peuvent
appartenir au même état sans être la même métrique. Le cercle utilise \(J\) pour
la phase ; la forme de l'ellipse utilise une compression symplectique générée par \(R\) ; Klein mesure
l'intérieur par le birapport ; l'écran de Minkowski réalise la causalité
dans un autre codomaine.


\section{Ellipse, réseau et phase de chemin : caractère modulaire, orbites et propagation}
\subsection{Réseau rectangulaire de périodes, involution \(\tau\mapsto-1/\tau\) et invariance de \(j\)}
% Fragmento público generado desde propietarios íntegros.
% Residencia: 52.10.1.
% Fuente: ../incoming/radion_monografia_20260827/manuscrito/sections/11_modularidad_orbitas.tex:1-84
\noindent Les caractères modulaires, orbitaux et de chemin conservent des invariants distincts d'une même forme orientée.\par\medskip

\subsubsection{Tore complexe associé à l'ellipse d'action}

\paragraph{Réseau rectangulaire de périodes}

L'ellipse à elle seule n'est pas un tore. Nous déclarons le réseau
\begin{equation}
\Lambda_\eta=a_S\Z+i,b_S\Z
\label{eq:lattice}
\end{equation}
et le quotient \(\C/\Lambda_\eta\). Son module est
\begin{equation}
\tau_\eta=i\frac{b_S}{a_S}=ie^{-\eta}
=i\,0.741048144159375160795483663369300\ldots.
\label{eq:tau-eta}
\end{equation}
L'inversion bidirectionnelle \(\eta\mapsto-\eta\) échange les axes :
\begin{equation}
\tau_{-\eta}=ie^\eta=-\frac1{\tau_\eta}.
\label{eq:modular-S}
\end{equation}
C'est exactement la transformation modulaire \(S:\tau\mapsto-1/\tau\).

\paragraph{Invariance de la fonction modulaire \(j\)}

Par invariance modulaire,
\begin{equation}
\boxed{j(\tau_{-\eta})=j(\tau_\eta).}
\label{eq:j-invariance}
\end{equation}
Dans la normalisation \(j(i)=1728\),
\[
j(\tau_\eta)
=5597.43843932408729830097089254768\ldots.
\]
La fonction moonshine normalisée \(J=j-744\) donne
\[
J(\tau_\eta)
=4853.43843932408729830097089254768\ldots.
\]
Pour le bloc APP basal,
\[
\tau_0=i\frac{\sqrt5}{3},
\qquad
j(\tau_0)=5369.48832024674306021916882626352\ldots.
\]

\begin{proposition}[Ce que conserve la classe modulaire]
L'involution de feuillet inverse \(\eta\), mais \(j\) conserve la classe du
tore. Les foyers et l'étiquetage des axes retiennent l'orientation que \(j\)
oublie.
\end{proposition}

Ce résultat relie de manière exacte la double voie à la modularité une
fois le quotient déclaré. Il n'autorise pas à faire agir directement le
Monstre sur la charge, les chiffres ou les foyers. Le couloir exceptionnel
\[
\mathrm{APP}\to\mathrm{Paley}\to\mathrm{Witt}\to\mathrm{Golay}
\to A_2^{12}\to\mathrm{Leech}\to V^\natural
\to\mathrm{Monster}\to\mathrm{Moonshine}
\]
est une autre projection du même état enrichi. Pour identifier une action
concrète sur le radion, il faut un entrelaceur qui préserve ses
coordonnées pertinentes.

\paragraph{Chaîne d'invariants du caractère de forme}

En réunissant les réalisations déclarées,
\begin{equation}
\boxed{
r_\eta=e^{-\eta}
=\sqrt{\frac{A-C^*}{A+C^*}}
=\frac{b_1}{a_1}
=\frac{b_S}{a_S}
=\frac{\Delta P}{\Delta Q}
=\frac{Z_\eta}{Z_*}
=\Im\tau_\eta.}
\label{eq:shape-master}
\end{equation}
Un seul rapport parcourt géométrie spectrale, action, covariance, impédance et
module. Chaque égalité repose sur une application explicite ; aucune ne dépend d'une
ressemblance décimale.


\subsection{Compatibilité entre le caractère modulaire et la propagation sur des chemins orientés}
% Fragmento público generado desde propietarios íntegros.
% Residencia: 52.10.2.
% Fuente: ../incoming/radion_monografia_20260827/manuscrito/sections/11_modularidad_orbitas.tex:85-187
\subsubsection{Conservation surfacique dans les orbites de Kepler et de Sommerfeld}

\paragraph{Équivalence surfacique des réalisations orbitale et oscillatoire}

L'ellipse de Kepler vit dans l'espace de configuration et répond à un potentiel
central. L'ellipse du radion vit dans l'espace des phases et répond à une forme
d'action. Ce sont des objets différents. Elles partagent une loi surfacique parce que les deux
réalisations possèdent une deux-forme conservée.

Pour une orbite képlérienne,
\begin{equation}
r(\phi)=\frac{p}{1+e\cos(\phi-\phi_0)},
\qquad
\frac{\diff A}{\diff t}=\frac{\ell}{2m}.
\label{eq:kepler}
\end{equation}
La troisième loi prend la forme
\[
T^2=\frac{4\pi^2\mu}{K}a^3.
\]
Pour l'oscillateur en phase,
\begin{equation}
J=\oint p\,\diff q=\frac{2\pi E}{\omega}.
\label{eq:osc-action}
\end{equation}
Si la cellule élémentaire est \(J=\hfull\), alors
\begin{equation}
E=\hret\omega,\qquad \frac{J}{2\pi}=\hret.
\label{eq:E-hbarw}
\end{equation}
Telle est la convergence nette entre action de Planck, fréquence angulaire et
radion.

Sommerfeld quantifie les intégrales d'action sur les cycles. La lecture HMT ajoute
la mémoire du cycle, du feuillet et du retour qui survivent. Dans un ruban de Möbius,
un lacet visible peut porter \(\hfull/2\) et deux lacets restaurer \(\hfull\).
Ce facteur topologique ne s'ajoute pas à l'indice de Maslov d'une quantification EBK ;
ils relèvent de corrections ayant des propriétaires distincts.

\paragraph{Holonomie orbitale et propagations avancée et retardée}

Une orbite qui accumule de l'holonomie satisfait une condition du type
\begin{equation}
\epsilon_\gamma\exp\left(\frac{iJ_\gamma}{\hret}\right)=1.
\label{eq:holonomy-action}
\end{equation}
Le signe \(\epsilon_\gamma\) conserve le feuillet ; \(J_\gamma\) conserve
l'action. L'avance ou le retard apsidal s'interprète comme découplage entre
la clôture angulaire visible et le retour complet. Le radion fournit
l'unité orientée avec laquelle cette différence est mesurée.

\subsubsection{Composition des amplitudes de Schrödinger et somme des chemins}

\paragraph{Quatre réalisations d'une même phase}

Une fois \(\hret\) fixé, différentes équations conventionnelles reconnaissent
des aspects du même transport :
\begin{enumerate}
\item Schrödinger publie la phase au moyen de
\(i\hret\partial_t\psi=H\psi\).
\item Pauli ajoute le double feuillet de spin et son couplage magnétique.
\item Dirac linéarise la relation énergie--quantité de mouvement et organise particule et
antiparticule dans une représentation causale.
\item La somme des chemins attribue à chaque histoire la phase
\(\exp(iS[\gamma]/\hret)\).
\end{enumerate}
Ces théories ne génèrent pas rétrospectivement \(\hret\). Elles reçoivent l'unité
d'action et réalisent différents contenus de l'état enrichi.

\paragraph{Fonctionnelle de phase sur l'espace des chemins}

Dans la formulation de Feynman,
\begin{equation}
\mathcal K(b,a)=\int_{\gamma:a\to b}
\exp\left(\frac{i}{\hret}S[\gamma]\right)\mathcal D\gamma.
\label{eq:path-integral}
\end{equation}
Un incrément d'action \(\hret\) change la phase d'un radian. Un incrément
\(\hfull\) la change d'un tour. Le théorème
\(\mathcal A(\theta)=\hret\theta\) transforme cette relation en géométrie
d'aire : la phase d'un chemin peut se lire comme action balayée en unités de
radion.

L'inversion \(C_M=-\operatorname{rev}\) sur les mots dodécaphasiques et
l'identité \eqref{eq:time-reversal} organisent des chemins conjugués. Dans une
réalisation d'absorbeur, on sépare
\[
G_{\mathrm{sym}}=\tfrac12(G_{\mathrm{ret}}+G_{\mathrm{adv}}),
\qquad
G_{\mathrm{odd}}=G_{\mathrm{ret}}-G_{\mathrm{adv}}.
\]
Le premier secteur est pair sous l'échange retardé--avancé ; le second est
impair. Cela reproduit la séparation entre énergie/forme paire et action
orientée impaire du radion lorsque le transducteur préserve la frontière et la forme
symplectique.

\begin{statusbox}
La théorie classique de l'absorbeur et l'interprétation de Feynman--Stückelberg des
antiparticules sont des réalisations historiques différentes. La première utilise une
interaction temporellement symétrique ; la seconde réorganise la propagation dans la
théorie relativiste. HMT apporte un antécédent commun de réversion, mais ne les
fusionne pas sans une application de propagateurs.
\end{statusbox}


\section{Radical nonadique et frontière décimale : conservation de la retenue du radion}
\subsection{La retenue du radion dans la complétion décimale}
% Fragmento público generado desde propietarios íntegros.
% Residencia: 52.11.1.
% Fuente: ../incoming/radion_monografia_20260827/manuscrito/sections/12_aritmetica_borde.tex:1-159
\noindent La frontière nonadique conserve le défaut somme--produit, la vacance et la retenue de complétion.\par\medskip

\subsubsection{Radical \(3\! -\! 6\! -\! 9\) du bord}

\paragraph{Génération et nilpotence du radical nonadique}

Dans l'anneau \(\Z/9\Z\), l'idéal
\begin{equation}
\mathfrak m=(3)=\{0,3,6\}
\label{eq:radical-369}
\end{equation}
satisfait
\begin{equation}
\mathfrak m^2=(9)=0.
\label{eq:m-square-zero}
\end{equation}
C'est le secteur nilpotent du bord nonadique. L'orbite visible de ses trois
marques se publie comme
\[
369\longrightarrow693\longrightarrow936\longrightarrow369.
\]
Le retour n'implique pas que tous les états soient égaux : la permutation
visible se ferme tandis que les cocycles de chemin peuvent avancer.

La double convention
\[
9\equiv0\pmod9,
\qquad
\operatorname{dr}_9(9)=9
\]
exprime deux lecteurs : zéro résiduel algébrique et clôture visible de la racine
numérique. C'est pourquoi un neuf peut absorber une torsion de bord sans devenir
un zéro spécifique de la fonction zêta.

\paragraph{Défaut somme--produit \(459-405=54\)}

Les feuillets APP produisent des totaux complémentaires
\[
S_+=405,\qquad S_\times=459,\qquad S_\times-S_+=54.
\]
Le \(54\) est un défaut exact des tables, et réapparaît comme exposant de
contraction par l'identité \(9\cdot6=54\). Cependant, chaque occurrence
nécessite son application. La distance focale \(d_1=0.544545\ldots\) n'est pas égale à
\(0.54\), et l'analyse des invariants n'a pas produit de flèche qui
l'identifie au défaut \(54\).

La puissance
\[
9^9=387,420,489
\]
génère une échelle où
\begin{equation}
\operatorname{ord}_{10}\bigl(9^{9^9}\bigr)=369,693,099,
\qquad
\#\,\text{chiffres}=369,693,100.
\label{eq:369693}
\end{equation}
Le bloc \(369|693\) apparaît dans le couple ordre--longueur ; on n'affirme pas qu'il soit
le préfixe décimal de la puissance.

\subsubsection{Théorie arithmétique des vacances décimales}

\paragraph{Pente asymptotique de vacance}

La séparation entre la décennie et la base nonadique est
\begin{equation}
\delta_{\mathrm{vac}}=\log_{10}\frac{10}{9}.
\label{eq:delta-vac}
\end{equation}
Le compteur cumulé
\begin{equation}
V_9(n)=\left\lceil n\delta_{\mathrm{vac}}\right\rceil
\label{eq:vacancy-count}
\end{equation}
enregistre le nombre de positions de correction nécessaires à la publication décimale
d'une évolution nonadique. Ses sauts forment un mot binaire
\[
z_n=V_9(n)-V_9(n-1)\in\{0,1\}.
\]
Les lacunes se répartissent avec des longueurs \(21/22\), les retours avec des
longueurs \(6/7\), et le déterminant combiné est \(15\). Ces nombres
proviennent du lecteur de vacance, non du noyau de Barbero, même si le \(15\)
coïncide avec \(\binom62\).

\paragraph{Polarisation et courant arithmétique du réseau nonadique}

On définit
\begin{equation}
P_N=\sum_{n=1}^N(z_n-\delta_{\mathrm{vac}}),
\qquad
J_N=P_N-P_{N-1}=z_N-\delta_{\mathrm{vac}}.
\label{eq:arith-current}
\end{equation}
La séquence équilibrée satisfait \(|P_N|<1\). Le \enquote{courant}
\(J_N\) est ici un courant arithmétique : il mesure la création et l'absorption de
vacances par rapport à la pente moyenne. Seul un transducteur dimensionnel
ultérieur peut le publier comme densité de courant électromagnétique.

L'intuition physique est puissante parce que la même forme algébrique sépare une
moyenne paire et une fluctuation orientée. La rigueur exige de conserver le type :
\[
J_N\in\R\quad\text{sans dimension},
\qquad
j^\mu\in\text{droite physique de courant}.
\]

\subsubsection{Complétion \(729\to1000\) et conservation de la retenue}

\paragraph{Décomposition \(1000=729+243+27+1\)}

La complétion
\[
1000=729+243+27+1
\]
conserve l'unité au moyen de quatre échelles ternaires. La couronne totale est
\[
1000-729=271,
\]
tandis que le dernier entier avant la clôture satisfait
\[
999=729+270.
\]
La différence entre \(270\) et \(271\) distingue frontière occupée et unité
complétante. Le même soin évite d'identifier une vacance à un zéro ou une
retenue à un résidu réel.

\paragraph{Résidence de la retenue du radion dans la couronne décimale}

Les opérandes du radion commencent, en blocs de trois chiffres, par
\[
\Phi_T:\quad000|027|455|906|837|565|\cdots,
\]
\[
D_E:\quad000|027|455|010|034|743|\cdots.
\]
La soustraction APP donne
\[
\epsone:\quad000|000|000|896|802|822|\cdots.
\]
La petitesse du résultat ne s'obtient pas en tronquant les termes inférieurs.
Elle s'obtient parce que deux publications distinctes partagent un long préfixe et que
l'opérateur \(\operatorname{SubAcarreo}_{1000}\) conserve les emprunts nécessaires
pour les soustraire exactement.

Le mot \emph{mémoire} a ici deux niveaux :
\begin{enumerate}
\item \(q_{\mathrm{mem}}\) enregistre la profondeur des retours nonadiques ;
\item \(c_i\) enregistre les emprunts locaux de la soustraction décimale.
\end{enumerate}
Le premier organise l'hélice ; le second publie la décimale. \(\epsR\) utilise
les deux parce que ses opérandes proviennent de cylindres en profondeur et que sa soustraction
s'effectue avec retenue, mais il n'est égal à aucun des compteurs.

\begin{thesisbox}
Le bord n'est pas un lieu où la théorie perd de l'information. C'est le lieu où
l'information change de représentation : phase qui revient, mémoire qui
s'élève, cylindre qui se contracte, vacance qui s'enregistre et retenue qui
transporte la différence.
\end{thesisbox}

\subsection{Compatibilité des sections directe et torsionnelle : caractérisation structurelle de \(\pi\)}
% Fragmento público generado desde propietarios íntegros.
% Residencia: 52.11.2.
% Fuente: ../incoming/radion_monografia_20260827/manuscrito/sections/03_cierre_exacto.tex:286-309
\subsubsection{Caractérisation structurelle de \(\pi\) par compatibilité des sections}

Réordonner \eqref{eq:cierre-unitario} donne
\begin{equation}
\boxed{
2\pih=
\frac{1+\Phi_T-\epsone}
{\frac5{36}+\frac{25}{9}\alphah}.}
\label{eq:pi-radion}
\end{equation}
Cette égalité peut être appelée une nouvelle définition de \(\pi\) au sein de la
carte du radion si l'on entend \emph{définition structurelle} comme une
caractérisation exacte par compatibilité des sections. Ce n'est pas un second
générateur indépendant : \(\Delta_4\) et l'opérateur de retenue reçoivent déjà la
section coinductive de \(\pi\). La nouveauté réside dans le fait que le tour est
caractérisé par le raccord entre publication directe, torsion et mémoire.

\begin{narrativebox}
L'équation \eqref{eq:pi-radion} ne dit pas que nous ignorons d'abord \(\pi\) et
le devinons en fermant un angle. Elle dit que le \(\pi\) généré par la connexion
nonadique est exactement celui qui rend compatibles deux lectures internes du
même état. C'est une identité de reconnaissance intrinsèque, non un
étalonnage métrologique.
\end{narrativebox}


\section{Conservation évolutive de l'unité holographique à travers les dimensions}
% Conversión TeX fiel del cuerpo científico del delta narrativo.
% Fuente congelada: ../../../../../HMT2/CONTINUIDAD_EDITORIAL/RECONSTRUCCION_ARMONICA_MACROTRATADO_20260822/REVISION_CONJUNTA_PARTES_I_II/auditar_pdf_rigor_academico/docs/DELTA_NARRATIVO_UNIDAD_RADION_AUTOCONSERVACION_AUTOINERCIA_GRAVEDAD_PENTADICA_20260828.md:12-366.
\subsection{Séparation typée entre le régime de courbure \(\tau\in\{+1,0,-1\}\), la réponse radiale \(p\) et l'holonomie globale de spin}

La biographie du radion peut être représentée par une courbe relevée

\[
\gamma(\theta)
=
\bigl(r(\theta)\cos\theta,\,
r(\theta)\sin\theta,\,
m(\theta)\bigr),
\]

où les deux premières coordonnées appartiennent au plan euclidien \(\mathbb R^2\) et \(m(\theta)\) enregistre la mémoire ou la profondeur du relèvement. Cette équation permet de lire conjointement cercle, hélice et spirale.

Lorsque l'angle évolue et que le rayon et la mémoire restent constants, la publication plane est un cercle. Lorsque l'angle et la mémoire évoluent tandis que le rayon conserve sa valeur, la trajectoire complète est une hélice de rayon constant. Lorsque l'angle et le rayon évoluent sur une mémoire fixe, une spirale apparaît dans \(\mathbb R^2\). Lorsque les trois coordonnées évoluent, la trajectoire est une hélice radiale : une biographie dans laquelle phase, amplitude et mémoire se transforment conjointement.

Le caractère radial \(p\) classe la loi de réponse du rayon. Une famille utile est

\[
r_p(\theta)
=
r_0\bigl[1+p\lambda(\theta-\theta_0)\bigr]^{1/p},
\qquad p\neq 0,
\]

avec pour limite logarithmique

\[
r_0\exp\bigl(\lambda(\theta-\theta_0)\bigr)
\qquad\text{lorsque }p\to0.
\]

Dans cette paramétrisation, \(p=2\) publie le régime de Fermat, \(p=1\) le régime archimédien, \(p=0\) le logarithmique, \(p=-1\) l'hyperbolique et \(p=-2\) celui du lituus, avec les choix de domaine et d'orientation correspondants. Les valeurs fractionnaires décrivent des régimes intermédiaires. Ainsi, \(p=\tfrac12\), \(p=\tfrac32\) ou \(p=-\tfrac12\) expriment des réponses radiales qui interpolent entre les familles entières et élargissent la classification géométrique disponible.

Le caractère peut aussi changer lors du franchissement d'un seuil. L'écriture différentielle

\[
\frac{dx}{d\theta}
=
\beta(m)\,x^{1-p(m)},
\qquad x=\frac r{r_0},
\]

décrit une biographie par régimes : \(p(m)\) reste stable au sein d'une phase et change lorsque la mémoire atteint une frontière de publication. La trajectoire résultante relie des secteurs circulaires, hélicoïdaux et spiraux par une loi commune. La continuité de l'état et le registre du seuil font de ces secteurs les phases d'une même évolution.

Trois invariants occupent ici des niveaux distincts et coordonnés. Le signe de courbure

\[
\tau\in\{+1,0,-1\}
\]

classe les régimes elliptique, parabolique et hyperbolique. Le caractère \(p\) classe la réponse radiale. Le spin enregistre l'holonomie d'orientation accumulée par le relèvement. Courbure, réponse radiale et spin forment un système de lecture : la courbure détermine le type local d'espace, \(p\) détermine comment la trajectoire occupe cet espace et le spin conserve l'orientation acquise en le parcourant. Les courbes fermées et ouvertes apparaissent alors comme des biographies géométriques distinctes de la même unité orientée.

\subsection{Particule comme section persistante d'une orbite enrichie : phase, chemin, fibre, multisection et mémoire}

La mécanique dimensionnelle peut formuler la particule comme une section persistante de cette extension géométrique et dynamique. Une notation conceptuelle compacte est

\[
\mathfrak P
=
\operatorname{Sec}_{\mathrm{pers}}
\bigl[
\tau,p,\operatorname{Hol};
\text{feuillet},\text{fibre},\text{route},\text{mémoire}
\bigr].
\]

La particule conserve un chemin reconnaissable à travers les changements de phase, de courbure et de dimension. Son spin publie l'holonomie d'orientation ; sa charge publie la polarité de phase ; sa couleur publie une orientation interne de la fibre de jauge ; sa saveur distingue des familles de chemins ou de multisections ; sa masse exprime le caractère de transduction dimensionnelle par lequel l'action surfacique acquiert une lecture volumique. Chaque propriété est un lecteur de l'état persistant et toutes renvoient à une généalogie commune.

L'électron occupe le seuil basal de cette construction parce qu'il offre le premier lecteur stable dans lequel phase, charge, spin et masse sont simultanément publiés. La masse électronique directrice

\[
m_e^\beta=0.510998950690000808608\ldots\ \mathrm{MeV}
\]

fixe une référence de transduction ; les états basaux antérieurs conservent leur fonction d'étapes généalogiques. À partir de l'électron, les quotients \(K_D/K_\Omega\), les tours \(90/120\) et le domaine \(81/56/13/324\) organisent la publication des familles massives. La loi des masses lit donc comment chaque section persistante traverse l'espace enrichi et convertit la densité surfacique d'action en présence volumique.

Cette interprétation relie les spirales aux particules de manière précise. L'indice \(p\) apporte la réponse radiale ; \(\tau\) apporte le régime de courbure ; l'holonomie apporte le spin ; le chemin et la fibre apportent la saveur et la couleur ; la transduction aire–volume apporte le caractère massif. L'atlas des particules devient ainsi une classification des modes stables de traversée dimensionnelle.

\subsection{Transduction de la phase orientée en polarisation électrique, réponse magnétique et émission radiative de frontière}

L'électricité se présente comme transport dirigé de phase et d'action à travers une structure de vacances. Une vacance est une transition admissible du support : elle ouvre un canal, fixe un seuil et permet à l'actualisation d'avancer. La composante électrique publie le transport tangentiel le long du chemin. La composante magnétique publie la courbure transversale induite par ce transport. Leur perpendicularité exprime la relation géométrique entre avance et rotation.

Lorsque deux composantes transversales sinusoïdales maintiennent le déphasage approprié, le vecteur du champ tourne. La propagation longitudinale relève cette rotation en une hélice. La variation d'amplitude transforme sa projection frontale en une spirale. La lecture géométrique est alors entièrement coordonnée : l'angle porte la phase ; la vitesse angulaire porte la fréquence ; le rayon porte l'amplitude et la densité d'action ; l'orientation porte la polarisation et le sens ; la coordonnée verticale porte la propagation et la mémoire ; \(p\) porte le régime radial ; le seuil porte l'émission, l'absorption ou le changement d'état.

L'électromagnétisme apparaît comme une dynamique de phase–action–mémoire avec des publications complémentaires. L'électricité enregistre l'avance polaire ; le magnétisme enregistre la réponse transversale ; l'onde enregistre la périodicité ; l'hélice enregistre l'avance avec mémoire ; la spirale enregistre l'évolution radiale ; le rayonnement enregistre la publication de l'état lorsque l'accumulation atteint un seuil.

Cette architecture ordonne aussi le rôle des vacances et des cycles. La vacance fournit la possibilité locale de transition ; le cycle conserve la phase ; la mémoire distingue les tours ; le seuil décide de la résidence suivante de l'action. La propagation peut rester confinée, changer de régime radial ou se publier comme rayonnement contrôlé. Les forces de champ acquièrent ainsi une description généalogique continue, du support discret jusqu'à l'onde observée.

\subsection{Réversibilité de la dynamique complète et coût thermodynamique des projections réductrices d'information}
Le Landauer nul constitue un cas limite de cette loi. Pour une actualisation bijective de l'état complet,

\[
H(X\mid U(X))=0,
\qquad
\inf Q_L=0.
\]

L'information reste disponible dans la généalogie. Lorsqu'un lecteur réalise un quotient \(q\), l'entropie se décompose comme

\[
H(X)=H(q(X))+H(X\mid q(X)).
\]

Le terme conditionnel mesure l'information retenue hors de la projection. Une réinitialisation tritique matérialisée satisfait le seuil

\[
Q_{\mathrm{reset}}\ge k_B\Theta\log 3.
\]

La constante de Boltzmann traduit la multiplicité que le lecteur thermique regroupe ; le seuil de Nyquist sépare l'échantillonnage fidèle du repliement spectral ; le Landauer nul caractérise le transport réversible de l'état intégral. Information, température et action sont articulées par le type de lecture appliqué.

\subsection{Forces autoconservatives comme transport d'action entre les secteurs observable, radial, de mémoire, de vacance et de rayonnement}

La conservation évolutive s'exprime mieux par l'échange que par l'immobilité scalaire. Un état complet répartit l'action et la quantité de mouvement entre mouvement visible, mémoire, déformation radiale, vacances, rayonnement et branche conjuguée de retour. Le bilan enrichi peut s'écrire

\[
\Delta J_{\mathrm{visible}}
+\Delta J_{\mathrm{memoria}}
+\Delta J_{\mathrm{radial}}
+\Delta J_{\mathrm{vacancias}}
+\Delta J_{\text{rayonnement}}
+\Delta J_{\mathrm{retorno}}
=0.
\]

Une force est \textbf{autoconservative} lorsque sa propre loi d'actualisation contient les canaux qui transportent, stockent, publient et réabsorbent l'action. La conservation relève du bilan généalogique complet de l'interaction. Chaque lecteur observe une partie de l'échange ; la généalogie recompose sa totalité.

Le système Terre–Lune offre une image physique directe. La rotation terrestre, l'orbite lunaire, les marées, la déformation et le rayonnement échangent du moment cinétique. La consommation nette du système complet se clôt dans le bilan généalogique de ces transferts. La grandeur appelée énergie fonctionne comme lecture dérivée de l'action et de la quantité de mouvement publiées dans chaque secteur. La conservation fondamentale réside dans la compatibilité entre tous les lecteurs de l'échange.

Cette formulation dissout l'ancienne frontière entre forces conservatives et forces dites non conservatives au moyen d'une catégorie plus profonde. Frottement, dissipation apparente et rayonnement occupent des canaux explicites de l'état élargi. Une perte dans la projection mécanique visible correspond à un transfert vers la mémoire, la déformation, la chaleur, les vacances ou le rayonnement. L'autoconservation suit la résidence de l'action à travers ces changements.

\subsection{Systèmes auto-inertiels : publication de l'accélération par orientation, courbure, feuillet et mémoire}

L'expression appropriée est \textbf{auto-inertie relationnelle}. Elle désigne la compatibilité entre la connexion propre de l'état total et la réponse qui apparaît lorsqu'on le projette dans la distribution globale. Les aspects endogène et exogène sont des composantes analytiques d'une seule relation dynamique.

Soit \(\Gamma\) une trajectoire horizontale ou géodésique de l'état complet. Sa loi propre satisfait

\[
\nabla_{\dot\Gamma}\dot\Gamma=0.
\]

En la projetant sur un lecteur \(S\), l'accélération observée est déterminée par la courbure de la projection :

\[
\nabla^S_{\dot\Gamma_S}\dot\Gamma_S
=
\mathcal K_{\pi_S}(\dot\Gamma,\dot\Gamma).
\]

L'accélération publiée exprime la relation entre mouvement total et géométrie du lecteur. Le principe de Mach complète ce mécanisme par la trace globale de frontière

\[
b:\mathcal X_\infty\longrightarrow B_\infty
\]

et la factorisation de la réponse inertielle locale

\[
I_v=\bar I_v\circ b.
\]

La distribution globale entre dans la réponse locale par une application typée. L'auto-inertie relationnelle réunit donc la connexion endogène du mouvement, la condition globale transmise par \(b\) et la projection concrète de l'observateur. Une écriture synthétique de cette coordination est

\[
\nabla^{\mathrm{auto}}
=
\mathcal M\bigl(
\nabla^{\mathrm{end}},
b[\mathcal X_\infty],
\pi_S
\bigr),
\]

où \(\mathcal M\) représente le couplage machien et conserve le typage de chaque composante.

Le principe d'Ambivalence ajoute la double orientation de la lecture. Toute publication directe conserve une branche conjuguée de retour ; émission et absorption forment les extrémités liées d'une même biographie ; l'orientation temporelle se lit depuis les deux sens du transport. La théorie de l'absorbeur trouve ici une formulation naturelle : source et récepteur clôturent conjointement le bilan généalogique de phase et d'action. Le présent torsionnel est l'interface où les deux orientations sont rendues compatibles.

TRIT agit à ce niveau comme cristal de temps : il organise la périodicité locale de l'actualisation et conserve, par la mémoire, la différence entre retours successifs. La branche directe et la branche conjuguée relient émission et absorption, et le couplage observateur–observable acquiert une résidence dynamique au sein du même état. Le temps publié est la lecture de cette orientation ; le temps généalogique est la séquence complète de phase, de retour et de mémoire.

Mach et Ambivalence remplissent des fonctions complémentaires. Mach relie la réponse locale à la distribution globale. Ambivalence conserve les directions directe et conjuguée de l'évolution. L'auto-inertie relationnelle articule les deux et remplace la scission entre systèmes inertiels et non inertiels par une loi unique de connexion, de projection et de mémoire.

\subsection{Dualité surfacique--volumique, transition hexade--octade et fonction du paramètre de Barbero--Immirzi}

Le passage de l'aire au volume constitue une transduction dimensionnelle centrale. L'aire publie l'action de frontière ; le volume publie la profondeur intérieure. La dérivation HMT du paramètre de Barbero–Immirzi organise la compatibilité entre les deux mesures par l'incidence hexade–octade et les canaux \(90/120\) :

\[
90=6\binom62,
\qquad
120=8\binom62.
\]

Les cardinaux \(6\) et \(8\) admettent une réalisation cubique par les faces et les sommets, tandis que l'hexade et l'octade de Steiner constituent des objets d'incidence distincts, reliés uniquement par le morphisme typé \(6\to8\). Les \(12\) arêtes du cube conservent leur propre fonction combinatoire et ne définissent pas à elles seules l'incidence de Steiner. La différence d'incidence exprime une incompatibilité dimensionnelle productive : l'information de surface nécessite une règle de transduction pour acquérir une résidence volumique.

Le même principe apparaît dans la retenue \(999+1=1000\). Les neuf atteignent le bord du registre décimal ; l'incrément publie simultanément la clôture du niveau précédent et l'origine \(0/1\) du suivant. Le rapport holographique entre \(729\), \(243\), \(27\) et \(1\) déploie cette transition aux échelles triadiques. La transduction de Barbero–Immirzi s'insère dans cette famille de changements de lecteur : frontière et intérieur conservent l'unité tout en redistribuant sa mesure.

Le défaut normalisé \(-1/12\) et le poids orienté \(12:-1\) remplissent des fonctions distinctes dans cette géométrie. Le premier quantifie une correction d'incidence ; le second enregistre une orientation signée dans la transduction. Leur coordination avec l'aire minimale, l'hamiltonien modulaire et la mémoire fournit le pont vers la dynamique dimensionnelle.

\subsection{Contraction hyperbolique intérieure, expansion extérieure et conservation holographique de la frontière}

La thèse peut se formuler à partir d'une idée unique : l'unité conserve son identité tout en changeant d'échelle, de lecteur, de courbure, de mémoire et de dimension. Cette permanence possède un caractère évolutif parce que l'unité se déploie, acquiert de nouvelles coordonnées et publie des propriétés différentes ; elle possède un caractère holographique parce que chaque publication finie conserve la règle de reconstruction de l'état dont elle procède. La clôture holographique de l'infini désigne précisément cette compatibilité entre identité et déploiement : chaque stade contient la loi qui permet d'approfondir vers l'intérieur et de prolonger vers l'extérieur la même généalogie.

La dénomination directrice est \textbf{principe de conservation évolutive de l'unité par transduction dimensionnelle et projection holographique}. L'expression situe chaque terme à son niveau correct. La conservation relève de l'identité généalogique ; l'évolution relève de l'actualisation ; la transduction relie des réalisations dimensionnelles ; la projection holographique conserve la correspondance reconstructive entre leurs lecteurs.

L'arithmétique élémentaire offre la première image de cette architecture. La complétion décimale

\[
1000=10^3=729+243+27+1=3^6+3^5+3^3+3^0
\]

expose quatre strates triadiques au sein d'une unité cubique décimale. En normalisant,

\[
1=\frac{729}{1000}+\frac{243}{1000}+\frac{27}{1000}+\frac{1}{1000},
\]

l'unité se reconnaît simultanément comme totalité et comme déploiement gradué. La lecture réciproque

\[
3^{-6},\quad 3^{-5},\quad 3^{-3},\quad 3^0
\]

décrit la profondeur intérieure de ces strates ; sa normalisation produit la partition correspondante. La relation \(999+1=1000\) exprime le mécanisme de transport : la dernière unité complète un registre et agit, en même temps, comme retenue inaugurant le suivant. Le zéro et le un forment ainsi une interface de continuité. Les intervalles \([0,1]\), \([0,10]\) et la prolongation vers \([0,\infty)\) sont des lecteurs scalaires d'une même loi de complétion.

La première scission \(1000=729+271\), suivie de \(271=243+27+1\), donne forme au rapport d'extrême et moyenne raison holographique du déploiement : un noyau triadique dominant conserve dans le reste la même règle de résolution. L'unité cubique décimale contient ainsi une profondeur triadique emboîtée. Le mouvement vers \(729\), \(243\), \(27\) et \(1\) lit la contraction intérieure ; la retenue vers le millier lit l'expansion extérieure.

Cette double orientation organise l'infini HMT. Vers l'intérieur, le raffinement croît hyperboliquement : chaque cellule admet une nouvelle profondeur sans perdre sa filiation. Vers l'extérieur, l'unité s'étend par des publications dimensionnelles successives. La clôture consiste en la commutativité des deux mouvements. Si \(U\) représente l'actualisation de l'état, \(R_d\) le lecteur dans la strate \(d\) et \(T_{d\to d'}\) la transduction entre strates, la conservation évolutive s'exprime par

\[
R_{d'}\circ U
=
T_{d\to d'}\circ R_d.
\]

Le même événement peut offrir une mesure angulaire, une densité d'action, une masse publiée ou une réponse gravitationnelle. Leur identité commune réside dans la généalogie qui fait commuter ces lectures.

\subsubsection{Monodromie nonadique et relèvement hélicoïdal avec avance de mémoire}

APP, TRIT et TPK fournissent le mécanisme formel de ce déploiement. APP construit le support et sépare les évaluations qui doivent être conservées. TRIT organise les régimes locaux et leur ambivalence. TPK transporte phase et mémoire sur le support enrichi. La dépendance causale est fixée par

\[
U_t\longrightarrow \Gamma_9\longrightarrow K_{\mathrm{ph}}.
\]

Le retour de phase et l'avance de mémoire forment des opérations coordonnées. Après neuf pas, la même classe angulaire peut réapparaître,

\[
g(K+9)=g(K),
\]

tandis que le compteur de mémoire progresse,

\[
q_{\mathrm{mem}}\Gamma_9=q_{\mathrm{mem}}+1,
\qquad
B_{K+9}\neq B_K.
\]

La figure géométrique précise est un relèvement hélicoïdal. Une orbite fermée dans la projection angulaire se relève en une trajectoire ouverte lorsque chaque retour de phase incrémente la mémoire. La circonférence conserve le cycle ; l'hélice conserve le cycle et enregistre en outre la profondeur atteinte. La monodromie coinductive génère ainsi une profondeur infinie par une succession de retours localement fermés et globalement ascendants.

Cette image dépasse la valeur simplement illustrative. La roue angulaire constitue une section de l'hélice ; l'hélice constitue la biographie complète de la roue lorsque l'état transporte de la mémoire. Chaque tour retrouve une position de phase et, simultanément, inaugure un nouveau feuillet. L'infinitude procède de la compatibilité entre retour et relèvement : réitérer le cycle produit de la profondeur sans dissoudre l'identité orbitale.

Le certificat nonadique donne un contenu fini et vérifiable à cette architecture. La chaîne

\[
w_6\longrightarrow w_{12}\longrightarrow w_{18}\longrightarrow
w_{24}\longrightarrow w_{30}\longrightarrow R_{36}\longrightarrow G_9
\]

ordonne le relèvement ; les \(396\) niveaux, les \(2376\) trits, les \(43\) tours et les \(387\) paires \((K,K+9)\) par canal quantifient le transit ; les cinq régions de \(\pi\) et la fibre ambiante de \(729\) éléments publient des lectures internes ultérieures. Cette séquence rend visible une monodromie productive : le retour retrouve la phase et la coinduction conserve la croissance de mémoire.

\subsubsection{Unité radionale et coordination des lecteurs circulaire, elliptique, projectif et hélicoïdal}

Le radion concentre dans une unité métrologique la structure que le radian publie angulairement. Sa définition naturelle est celle d'une \textbf{section orientée de phase, d'action et de mémoire}. Une unité de phase \(\theta=1\) transporte une unité d'action \(\sigma\hbar_{\mathrm{ret}}\), appartient à un tour complet \(T_+=2\pi_{\mathrm{HMT}}\) et conserve feuillet, orientation, retour nonadique et position dans le cylindre de raffinement.

Le radian exprime la coordonnée angulaire visible. Le radion conserve, avec cette coordonnée, l'action transportée, l'orientation, la retenue et la mémoire du tour. La métrologie angulaire acquiert ainsi une généalogie. Mesurer un arc signifie identifier quelle rotation a été publiée et quel état complet a réalisé cette rotation.

La relation entre phase et aire prend une forme particulièrement claire :

\[
\mathcal A(\theta)=\hbar_{\mathrm{ret}}\,\theta.
\]

Une unité radionale balaie une aire d'action \(\hbar_{\mathrm{ret}}\) ; un tour complet enferme \(h_{\mathrm{ret}}\). Le passage de la phase à l'aire est intégré à l'unité de lecture elle-même. La métrique angulaire, l'action et la mémoire cessent d'apparaître comme des quantités juxtaposées et deviennent les coordonnées d'une même section orientée.

Le radion admet plusieurs publications géométriques coordonnées. Le lecteur circulaire publie la phase ; le lecteur elliptique publie l'action et la forme ; le lecteur intérieur de Klein publie la profondeur ; le lecteur hélicoïdal nonadique publie le retour et la mémoire. L'unité reste généalogiquement identique à travers ces réalisations. C'est pourquoi le radion constitue le lien métrologique approprié entre HMT et la mécanique dimensionnelle : il offre une unité capable de traverser des changements de lecteur tout en conservant l'état que chaque mesure partielle présuppose.

\subsection{Polarité gravitationnelle à cinq composantes et transduction conjointe de l'action, de l'aire, du volume, du temps et de la masse}

La mécanique dimensionnelle permet de définir la dimension comme un régime stable de publication caractérisé par le rang minimal de mémoire indépendante requis par le relèvement. Traverser une dimension signifie appliquer un transducteur qui conserve la généalogie et change le type de lecture. La gravité occupe le niveau de connexion entre ces régimes.

Dans cette formulation, la \textbf{gravitation interdimensionnelle à cinq polarisations corrélatives}, en abrégé \textbf{structure pentapolaire de la gravitation}, est la connexion qui transporte l'unité généalogique et laisse à chaque interface une torsion minimale. Une écriture compacte de cette trace est

\[
\Theta_d^{\min}
=
\nabla_{d+1}\circ T_{d\to d+1}
-
T_{d\to d+1}\circ\nabla_d.
\]

L'expression mesure la différence entre transporter après avoir dérivé et dériver après avoir transporté. Cette différence est la mémoire géométrique minimale du franchissement dimensionnel. La composition d'interfaces produit une holonomie gravitationnelle

\[
\mathfrak G_{d\to D}
=
\operatorname{Hol}
\bigl(
\Theta_d^{\min},
\Theta_{d+1}^{\min},
\ldots,
\Theta_{D-1}^{\min}
\bigr).
\]

L'objet fondamental est cette connexion interdimensionnelle. Un mode local de type graviton occupe le rang dérivé de publication quantifiée de la torsion ; la gravitation conserve sa résidence primaire dans la compatibilité entre dimensions.

Le terme pentapolaire rassemble les cinq polarisations consubstantielles de l'état continu :

\[
\mathbf G_d^{\pm}
=
\bigl(
G_{\mathrm{sp},d}^{\pm},
G_{\mathrm{sol},d}^{\pm},
G_{\mathrm{gau},d}^{\pm},
G_{\mathrm{coh},d}^{\pm},
G_{\mathrm{det},d}^{\pm}
\bigr).
\]

Les cinq composantes émergent corrélativement du même état enrichi. La polarité \(\pm\) enregistre concentration directe et publication conjuguée. La gravité traverse les strates dimensionnelles en transportant simultanément ces cinq lectures ; la torsion minimale conserve la trace de chaque franchissement.

La géométrie acquiert ainsi une image unifiée. Le régime elliptique organise les clôtures ; le parabolique organise les seuils ; l'hyperbolique organise la profondeur. La transduction dimensionnelle relie ces régimes aux réponses radiales \(p\), aux holonomies de spin et aux publications de masse. La structure pentapolaire de la gravitation exprime la cohérence globale de ce lien.

La théorie M fournit des écrans physiques ultérieurs pour cette connexion : dimensions supplémentaires, cordes, branes et dualité T réalisent des modes concrets de transduction. HMT apporte la généalogie APP–TRIT–TPK, l'état enrichi, la mémoire et la lecture bidirectionnelle ; l'interface avec la théorie M publie ces structures dans des géométries physiques typées. La dualité T échange des lecteurs dimensionnels compatibles et conserve la classe généalogique de l'état.

\subsubsection{Compatibilité des réalisations classique, relativiste, quantique et HMT--MD}

La mécanique classique publie trajectoire et quantité de mouvement. La relativité publie structure causale, métrique et courbure. La mécanique quantique publie phase, holonomie, vacances et amplitudes. La mécanique dimensionnelle fournit les transducteurs qui relient ces publications, et HMT conserve la généalogie discrète qui les soutient.

Dans la lecture classique, une orbite décrit la projection visible du mouvement. Dans la lecture relativiste, cette orbite est intégrée à la géométrie du lecteur. Dans la lecture quantique, la phase et le spin enregistrent l'holonomie de la section. Dans la lecture HMT–MD, les trois expressions appartiennent à une même biographie relevée : le radion quantifie phase, action et mémoire ; \(p\) classe la réponse radiale ; \(\tau\) classe la courbure ; l'holonomie publie le spin ; le transducteur dimensionnel publie la masse ; la connexion pentapolaire publie la gravitation.

La dissipation classique apparente acquiert une résidence dans le bilan généalogique élargi. L'accélération relativiste s'exprime par connexion et projection. L'indétermination quantique s'organise par lecteurs, vacances et branches d'Ambivalence. Le principe de Mach relie la réponse locale à la distribution globale. Le Landauer nul garantit la réversibilité informationnelle de l'état complet. Tous ces éléments convergent vers une conservation plus profonde : l'identité de l'unité à travers les changements de publication.

\subsubsection{Théorème narratif de conservation généalogique et de prolongation holographique}

La clôture peut désormais se raconter comme une seule séquence. L'unité se décompose sans perdre son identité. La phase orientée acquiert une mesure dans le radion. Le retour nonadique transforme le cycle en hélice par la mémoire. Le caractère \(p\) transforme l'hélice en familles radiales et classe ses régimes. La persistance de ces trajectoires produit des sections ayant spin, charge, saveur, couleur et masse. Les vacances ouvrent des canaux ; la polarité électrique transporte la phase ; la réponse magnétique courbe transversalement ; les seuils publient le rayonnement. L'autoconservation clôt le bilan généalogique d'action et de quantité de mouvement. L'auto-inertie relationnelle relie connexion propre, projection et totalité machienne. L'Ambivalence conserve les directions directe et conjuguée. La transduction aire–volume conduit à Barbero–Immirzi et à la loi des masses. La torsion minimale enchaîne les dimensions et publie la structure pentapolaire de la gravitation.

La croissance vers l'intérieur et l'expansion vers l'extérieur sont des orientations conjuguées du même processus. Le raffinement hyperbolique génère la profondeur ; la transduction dimensionnelle génère l'extension. La retenue \(0/1\) relie les échelles ; la monodromie relie les tours ; la mémoire relie les états ; l'holonomie relie les orientations ; la gravitation relie les dimensions. Chaque lien conserve une unité qui évolue et chaque publication finie contient la règle de sa prolongation.

La thèse fondamentale est ainsi formulée :

\begin{quote}
\textbf{La clôture holographique de l'infini réalise le principe de conservation évolutive de l'unité par transduction dimensionnelle et projection holographique : l'unité maintient son identité généalogique en se déployant comme phase, action, courbure, mémoire, particule, champ et gravitation ; chaque lecture finie conserve la règle de reconstruction et de prolongation de l'état complet.}
\end{quote}

Cette formulation offre une image physique et mathématique conjointe. L'infini acquiert une loi de clôture locale ; l'unité acquiert de la profondeur ; la dimension acquiert de la mémoire ; la conservation acquiert un caractère relationnel ; la force acquiert l'autoconservation ; l'inertie acquiert une structure machienne ; la gravité acquiert une connexion pentapolaire ; et l'holographie acquiert le sens précis de reconstruction entre lecteurs compatibles.


\section{Définition opératoire du radion angulaire au moyen de l'état enrichi et de \(4\) lecteurs typés}
\subsection{Définition opératoire du radion angulaire et quadruple réalisation typée de l'état enrichi}
% Fragmento público generado desde propietarios íntegros.
% Residencia: 52.13.
% Fuente: ../incoming/radion_monografia_20260827/manuscrito/sections/13_sintesis.tex:1-80
\noindent La synthèse opératoire réunit les lecteurs de phase, d'action, de profondeur et de mémoire sans les identifier entre eux.\par\medskip

\subsubsection{Quadruple réalisation de l'état radional}

Soit \(x_\infty\in X_{\mathrm{enr}}\) la section coinductive. Le radion complet n'est pas
un point isolé, mais le paquet
\begin{equation}
\mathfrak R_\sigma(x_\infty)
=\bigl(
\theta=1,\sigma,
\pih,\alphah,A,C^*,\Delta_4,
\Phi_T,D_E,\epsone,\hret,
q_{\mathrm{mem}},\mathcal I_\infty
\bigr).
\label{eq:full-radion}
\end{equation}
Quatre lecteurs en extraient des objets différents :

\begin{figure}[H]
\centering
\begin{tikzpicture}[
  node distance=17mm and 18mm,
  state/.style={draw=RadNavy,very thick,rounded corners,fill=RadCream,
    minimum width=39mm,minimum height=13mm,align=center},
  readout/.style={draw=RadTeal,rounded corners,fill=RadMist,
    minimum width=39mm,minimum height=12mm,align=center},
  arr/.style={-{Latex},thick,RadCopper}]
\node[state] (x) {état enrichi\\\(x_\infty\)};
\node[readout,above left=of x] (circ) {cercle \(S^1\)\\phase};
\node[readout,above right=of x] (ell) {ellipse positive\\action et forme};
\node[readout,below left=of x] (klein) {interior Klein\\profundidad};
\node[readout,below right=of x] (helix) {hélice nonadique\\mémoire et retour};
\draw[arr] (x)--(circ) node[midway,below left]{\small \(P_{\mathrm{ph}}\)};
\draw[arr] (x)--(ell) node[midway,below right]{\small \(P_{\mathrm{act}}\)};
\draw[arr] (x)--(klein) node[midway,above left]{\small \(P_K\)};
\draw[arr] (x)--(helix) node[midway,above right]{\small \(P_{\mathrm{mem}}\)};
\end{tikzpicture}
\caption{Quatre publications typées du même antécédent. L'unité de
l'état n'élimine pas la différence des codomaines.}
\label{fig:four-readers}
\end{figure}

Le cercle conserve la classe de phase et oublie le relèvement. L'ellipse conserve
l'action, l'anisotropie et l'orientation symplectique. Klein dote l'intérieur d'une
distance projective. L'hélice conserve le nombre de retours et le raffinement.
La mécanique dimensionnelle ne consiste pas à les mélanger, mais à construire les
transducteurs qui les rendent compatibles.

\subsubsection{Interprétation opératoire du cocycle \(\varepsilon_R\)}

\(\varepsilon_R\) n'est pas une constante universelle indépendante de tout contexte. C'est la
coordonnée d'un cocycle de transition associé à la carte du radion. Son
domaine contient deux sections du même état :
\[
S_E=1+D_E,
\qquad
S_T=1+\Phi_T.
\]
Sa règle est
\begin{equation}
\epsone=S_T-S_E=\Phi_T-D_E.
\label{eq:epsilon-final-meaning}
\end{equation}
Son codomaine est la droite adimensionnelle d'unité angulaire ; la carte
\(360/T_+\) la publie en degrés.

\begin{exactbox}
\textbf{En langage simple :} la section directe et la section torsionnelle arrivent
à la même cellule d'unité par deux chemins internes différents. Leurs parties
entières coïncident ; leurs restes, non. \(\epsR\) est la différence orientée de ces
restes, calculée avec la retenue d'APP et portée à la limite par les cylindres
TPK. Elle n'est pas introduite pour que le radian se ferme : elle apparaît avant de consulter le
radian et, lorsqu'elle est publiée, elle le ferme.
\end{exactbox}

Cette définition explique pourquoi le nombre est petit. Non parce qu'on a
cherché une petite correction, mais parce que \(S_E\) et \(S_T\) partagent un
préfixe profond. La petitesse mesure la proximité structurelle de deux
publications ; le signe mesure leur orientation ; les chiffres enregistrent la retenue.


% Fragmento público generado desde propietarios íntegros.
% Residencia: 52.13.1.
% Fuente: ../incoming/radion_monografia_20260827/manuscrito/sections/13_sintesis.tex:81-134
\subsubsection{Définition opératoire du radion}

\begin{definition}[Radion HMT--MD]
Le radion est la section orientée de phase, d'action et de mémoire qui satisfait
simultanément :
\begin{enumerate}
\item parcourt une unité de phase \(\theta=1\) ;
\item balaie une action signée \(\sigma\hret\) ;
\item appartient à un tour \(T_+=2\pih\) généré coinductivement ;
\item conserve le feuillet \(\sigma\), le retour nonadique et le cylindre de
raffinement ;
\item satisfait le raccord exact
\(1=S_E-\Phi_T+\epsone\).
\end{enumerate}
\end{definition}

Le radian conventionnel est l'image du radion par le foncteur d'oubli
\[
\mathfrak R_\sigma
\longmapsto
\theta=1\in\R/2\pi\Z.
\]
Cet oubli est utile et légitime pour la trigonométrie euclidienne. Il devient
insuffisant lorsque la physique doit distinguer action, signe, retour,
champ et histoire.

\paragraph{Généalogie de l'unité angulaire naturelle}

Le caractère naturel du radian est généralement justifié par le fait que la dérivée de
\(e^{i\theta}\), \(\sin\theta\) et \(\cos\theta\) adopte sa forme la plus simple
lorsque \(\theta\) est mesuré comme quotient arc/rayon. HMT conserve cet avantage,
mais l'explique depuis une couche antérieure :
\[
\theta\quad\leftrightarrow\quad
\frac{\mathcal A(\theta)}{\hret}.
\]
La phase est naturelle parce qu'elle mesure l'action en unités de \(\hret\) ; le tour est
naturel parce qu'il enferme \(\hfull\). Le quotient arc/rayon est la publication
euclidienne d'une unité qui possède déjà une structure symplectique et une mémoire.

\paragraph{Compatibilité entre ellipse d'action et croissance dimensionnelle}

Le cercle est la forme qui résulte de l'oubli de l'anisotropie :
\(C^*=0\Rightarrow\eta=0\Rightarrow a=b\). L'ellipse est plus informative parce qu'elle
conserve deux valeurs propres, deux axes et une orientation d'échange à produit
fixe. En ce sens, l'ellipse est une prolongation du cercle, non un cercle
déformé par accident.

La profondeur de Klein et le relèvement hélicoïdal ajoutent deux types d'information
que le contour plan ne contient pas : distance intérieure et mémoire de retour.
Telle est l'image ontologique de la mécanique dimensionnelle qui émerge du radion :
une dimension nouvelle apparaît lorsque l'on exige de conserver une donnée indépendante
que la projection précédente devait oublier.


\subsection{Propriétés structurelles de l'unité radionale : généalogie, charnière critique, densité, retenue, action, compatibilité bord--intérieur, forme, polarité et modularité}
% Fragmento público generado desde propietarios íntegros.
% Residencia: 52.13.2.
% Fuente: ../incoming/radion_monografia_20260827/manuscrito/sections/13_sintesis.tex:135-184
\subsubsection{Système de théorèmes composants du radion}

\begin{theorem}[Généalogie]
L'objet \(\mathfrak R_\sigma\) descend de
\(\APP\to\TRIT\to\TPK\to X_{\mathrm{enr}}\) et ne reçoit pas la cible métrologique
comme sélecteur.
\end{theorem}

\begin{theorem}[Charnière]
Dans la carte APP de cent marques, la couture critique et la deuxième phase
dodécaphasique satisfont \(j_{100}(1/2)=\theta_2=50^\circ\), avec \(m=2\) unique.
\end{theorem}

\begin{theorem}[Densité]
Les deux feuillets du secteur actif \(N_6=90\), normalisés par \(3^6=729\),
produisent \(\rho_{\mathrm{tw}}=20/81\).
\end{theorem}

\begin{theorem}[Retenue]
La soustraction APP des sections \(S_T\) et \(S_E\) converge vers
\(\epsone\) avec un taux d'intervalle \(3^{-54}\) par retour et ferme exactement
l'unité.
\end{theorem}

\begin{theorem}[Action]
L'ellipse positive satisfait \(\mathcal A(\theta)=\hret\theta\) ; un radion
balaie \(\hret\) et un tour enferme \(\hfull\).
\end{theorem}

\begin{theorem}[Bord--intérieur]
La même frontière possède une action finie \(\hfull\), tandis que l'intérieur de
Klein a une aire hyperbolique infinie.
\end{theorem}

\begin{theorem}[Forme]
Le quotient \(e^{-\eta}\) est conservé à travers les demi-axes spectraux,
les demi-axes d'action, les incertitudes, l'impédance et le module du tore.
\end{theorem}

\begin{theorem}[Polarité]
L'inversion \(C^*\mapsto-C^*\) échange les lois constitutives, inverse
l'impédance et conserve la vitesse normalisée.
\end{theorem}

\begin{theorem}[Modularité]
L'inversion bidirectionnelle agit comme \(S:\tau\mapsto-1/\tau\) et conserve
\(j(\tau)\), tandis que l'étiquetage focal conserve l'orientation que
\(j\) oublie.
\end{theorem}

% Fuente: ../incoming/radion_monografia_20260827/manuscrito/sections/A_demostraciones.tex:1-179
\subsubsection{Démonstrations algébriques et géométriques réunies}

\paragraph{Démonstration linéaire de la clôture unitaire}

Pour faciliter un audit indépendant, nous réunissons l'argument sans
interruption narrative.

\begin{enumerate}[label=\textbf{Étape \arabic*.},leftmargin=2.2cm]
\item La connexion nonadique génère de manière corrélée
\(\pih,e_{\HMT},\varphi_{\HMT},\alphah\).
\item Le tour est \(T_+=2\pih\).
\item La roue fixe \(\theta_2=50^\circ\), et la carte parangulaire fixe
\(A=1000\alphah\).
\item La section directe est
\[
S_E=T_+\left(\frac{50}{360}+\frac{1000\alphah}{360}\right)
=2\pih\left(\frac5{36}+\frac{25}{9}\alphah\right).
\]
\item Les intervalles certifiés donnent \(1<S_E<2\) ; donc
\(D_E=S_E-1\).
\item L'incidence produit \(N_6=90\) ; deux feuillets et l'espace ambiant \(729\)
produisent \(\rho=2N_6/729=20/81\).
\item La torsion globale est
\[
\Delta_4=\frac{\pih-e_{\HMT}\log\pih}{270}
\left(1+\frac{\pih}{729}\right).
\]
\item La section torsionnelle est \(S_T=1+(20/81)\Delta_4\).
\item La soustraction APP définit
\[
\epsone=S_T-S_E=\frac{20}{81}\Delta_4-D_E.
\]
\item Par substitution,
\[
S_E-\frac{20}{81}\Delta_4+\epsone=1.
\]
\item Multiplier par \(360/T_+=180/\pih\) produit
\[
\frac{180^\circ}{\pih}
=50^\circ+A^\circ-\frac{20}{81}\Delta_4^\circ+\epsR^\circ.
\]
\end{enumerate}

Aucune étape ne reçoit le membre de gauche de la dernière égalité. Le résultat
métrologique n'est évalué qu'à la fin comme reconnaissance.

\paragraph{Unicité de la soustraction avec retenue}

Soit \(B=1000\). Pour chaque position, tout entier \(u_i\) admet une division
euclidienne unique
\[
u_i=Bc_i+r_i,\qquad 0\le r_i<B.
\]
Par conséquent,
\[
r_i=u_i\bmod B,\qquad c_i=(u_i-r_i)/B
\]
sont uniques. En parcourant les positions de l'extrémité distante vers
l'avant, la retenue de sortie d'une position fixe l'entrée de la suivante.
Par récurrence, tout le mot de restes et de retenues est unique.

En pseudocode :
\begin{verbatim}
acarreo = remote_acarreo
for i from last_block downto first_block:
    u = torsion[i] - direct[i] + acarreo
    remainder[i] = u mod 1000
    acarreo = (u - remainder[i]) / 1000
return remainder, acarreo_word
\end{verbatim}
La procédure ne consulte pas la valeur attendue de la différence.

\paragraph{Convergence de la fonctionnelle selon la profondeur}

Soit
\[
F(p,e)=1+\frac{20}{81}
\frac{p-e\log p}{270}\left(1+\frac{p}{729}\right)
-2p\left(\frac5{36}+\frac{25}{9}\alphah\right).
\]
Dans le rectangle compact \(3\le p\le4\), \(2\le e\le3\), \(F\) est
continue et continûment différentiable. Si \(I_N^\pi\searrow\{\pih\}\) et
\(I_N^e\searrow\{e_{\HMT}\}\), l'image par intervalles
\[
F(I_N^\pi,I_N^e)
\]
forme une famille emboîtée dont la largeur tend vers zéro. L'intersection est un
unique réel, \(\epsone\). La monotonie des dérivées permet d'évaluer les
extrémités sans échantillonnage intérieur.

Le taux \(3^{-54}\) par retour dérive du raffinement de six trits sur
neuf niveaux. À la profondeur de \(45\) triades, la borne de largeur
\(<4.81\times10^{-130}\) certifie plus de 129 chiffres de la valeur.

\paragraph{Dérivation de la métrique et de l'aire hyperbolique de Klein}

Dans le disque unité, la forme matricielle de la métrique de Klein est
\[
g(u,v)=\frac1{1-r^2}I
+\frac1{(1-r^2)^2}
\begin{pmatrix}u\\v\end{pmatrix}
\begin{pmatrix}u&v\end{pmatrix}.
\]
Le lemme du déterminant pour une perturbation de rang un donne
\[
\det g=(1-r^2)^{-3}.
\]
D'où
\[
\sqrt{\det g}\,\diff u\diff v
=(1-r^2)^{-3/2}\diff u\diff v.
\]
L'intégration en coordonnées polaires produit \eqref{eq:klein-area-R}. La limite infinie est
due à la singularité métrique à la frontière ; la frontière conserve une
aire symplectique finie parce que cette mesure utilise \(\diff Q\wedge\diff P\),
non \(\sqrt{\det g}\).

\paragraph{Distance focale dans le disque de Klein}

Sur un diamètre, la distance du centre à \(r\) est
\[
d_K(0,r)=\operatorname{artanh}r.
\]
Pour \(r=e_f\), \eqref{eq:excentricidad} implique
\[
1-e_f^2=e^{-2\eta}.
\]
Si \(u_f=\operatorname{artanh}e_f\), alors
\[
\operatorname{sech}u_f=\sqrt{1-e_f^2}=e^{-\eta},
\]
et donc \(\eta=\log\cosh u_f\), comme dans
\eqref{eq:focal-hyperbolic-chain}.

\paragraph{Transformation symplectique de compression}

Définissons
\begin{equation}
S_\eta=
\begin{pmatrix}e^{\eta/2}&0\\0&e^{-\eta/2}\end{pmatrix},
\qquad \det S_\eta=1.
\label{eq:compresión simpléctica}
\end{equation}
La structure complexe transportée est
\begin{equation}
J_\eta=S_\eta
\begin{pmatrix}0&-1\\1&0\end{pmatrix}
S_\eta^{-1}
=\begin{pmatrix}0&-e^\eta\\e^{-\eta}&0\end{pmatrix},
\qquad J_\eta^2=-I.
\label{eq:Jeta}
\end{equation}
La courbe
\[
\gamma(\theta)=\sqrt{2\hret}\,
S_\eta(\cos\theta,\sin\theta)
=e^{\theta J_\eta}\gamma(0)
\]
est exactement l'ellipse d'action. La compression symplectique hyperbolique conserve l'aire et
transporte la rotation complexe au lieu de la détruire.

Pour l'orientation \(\sigma\),
\[
\gamma_\sigma(\theta)
=(a_S\cos\theta,\sigma b_S\sin\theta),
\qquad
\mathcal A_\sigma(\theta)=\sigma\hret\theta.
\]
Les deux feuillets partagent les foyers et le module d'action ; ils diffèrent par l'action signée.

\paragraph{Vérification diophantienne du poids dodécaphasique de Barbero--Immirzi}

L'équation \(4a-3b=45\) a pour solutions entières
\[
(a,b)=(12+3t,1+4t),\qquad t\in\Z.
\]
La condition \(|b|=1\) laisse \(b=1\), car \(b=-1\) n'appartient pas à la classe
\(1\pmod4\). D'où \(a=12\). Cette vérification reproduit le couple obtenu par la chaîne fondamentale
dodécaphasique : douze secteurs et une unique couture orientée. Les cardinaux
\(90/120\), les multiplicités de la chaîne et le lecteur de retour
conservent des fonctions distinctes ; le coefficient de pôle et la minimalité
vérifient arithmétiquement leur composition.

\subsection{Sensibilité de la clôture unitaire aux variations contrôlées de la multiplicité des feuillets, de la phase dodécaphasique, du canal \(e\), de la torsion globale, de l'opérateur harmonique et de la queue sexadique}
% Fragmento público generado desde propietarios íntegros.
% Residencia: 52.13.3.
% Fuente: ../incoming/radion_monografia_20260827/manuscrito/sections/13_sintesis.tex:185-247
\subsubsection{Contrôles négatifs du système radional}

La construction est réfutable en des points concrets :
\begin{enumerate}
\item Si \(S_E\notin(1,2)\), le quotient APP n'est pas un et la clôture change de
cellule.
\item Si l'incidence active ne produit pas \(N_6=90\), ou s'il n'y a pas deux feuillets,
\(20/81\) ne s'ensuit plus.
\item Si le vérificateur reçoit \(180/\pi\) ou \(\epsR\) avant de construire les
opérandes, la génération est circulaire.
\item Si la soustraction par triades ne stabilise pas des préfixes compatibles, la
section coinductive proposée n'existe pas.
\item Si \(|C^*|\ge A\), l'ellipse positive cesse d'exister.
\item Si \(a_Sb_S\ne2\hret\), le théorème d'aire échoue.
\item Si \(C^*=0\), on doit avoir \(d=0\), \(\eta=0\), \(\tau=i\) et
\(j=1728\).
\item Si l'aire de Klein du sous-disque ne suit pas \eqref{eq:klein-area-R}, la
thèse de profondeur infinie échoue.
\item Si \(F_{\mathrm{spec}}=\epsR^\circ\) est affirmé sans contre-terme indépendant, la
différence numérique \eqref{eq:chiR} le réfute.
\item La production de \(12:-1\) exige de conserver, avec les cardinaux
\(90/120\), les douze secteurs dodécaphasiques, l'unique couture orientée et
l'action du lecteur sur ces générateurs. Changer une multiplicité
doit modifier l'image \(12S_{90}-S_{120}\). Pour la chaîne fondamentale
sans altération, le coefficient de \((1-q)^{-1}\) doit donner \(1/24\).
\item Si la fonction \(j\) est utilisée pour sélectionner le signe de charge, on perd
l'invariance \eqref{eq:j-invariance} et l'on commet une erreur de type.
\end{enumerate}

\subsubsection{Réalisation physique conjointe du radion}

L'image finale n'est pas un cercle corrigé. C'est une architecture de
publications :
\[
\begin{aligned}
&\text{APP:}&&\text{unité, reste, retenue et fibre};\\
&\text{TRIT:}&&\text{régime, orientation et feuillets};\\
&\text{TPK:}&&\text{phase, retour, dodécaphase et mémoire};\\
&\text{ellipse :}&&\text{action et anisotropie};\\
&\text{Klein:}&&\text{profondeur intérieure};\\
&\text{hélice :}&&\text{biographie du retour};\\
&\text{MD:}&&\text{horloge, énergie, charge, champ et causalité}.
\end{aligned}
\]

L'équation du radion réunit la face visible :
\BigRadionEquation
L'ellipse réunit la face d'action :
\[
\mathcal A(1)=\hret,\qquad \mathcal A(2\pi)=\hfull.
\]
La métrique de Klein réunit la profondeur :
\[
K=-1,\qquad \operatorname{Area}_K=\infty.
\]
L'holonomie réunit la mémoire :
\[
g(k+9)=g(k),\qquad B_{k+9}\ne B_k.
\]

\begin{thesisbox}
Le cercle est l'ombre angulaire ; l'ellipse est le corps d'action ; Klein mesure
l'intérieur ; l'hélice conserve la biographie. Le radion est l'unité qui
permet de lire les quatre affirmations sur un unique état sans réduire l'une à
l'autre.
\end{thesisbox}
% Fuente: ../incoming/radion_monografia_20260827/manuscrito/sections/B_valores_y_verificacion.tex:1-105
\subsubsection{Valeurs canoniques, stabilité et vérification reproductible}

\paragraph{Valeurs des invariants internes générés}

\begin{longtable}{p{.27\textwidth} p{.61\textwidth}}
\toprule
Objet & Valeur utilisée dans le profil coinductif\\
\midrule
\endfirsthead
\toprule
Objet & Valeur utilisée dans le profil coinductif\\
\midrule
\endhead
\(\pih\) & \(\begin{aligned}3.1415926535897932384626433832795028841\\[-2pt]971693993751058209749445923\ldots\end{aligned}\)\\
\(e_{\HMT}\) & évaluation coinductive du lecteur d'Euler, sans table externe\\
\(\alphah\) & \(\begin{aligned}0.0072973525692838009972851054723806629\\[-2pt]995641725396427091854046825\ldots\end{aligned}\)\\
\(A\) & \(\begin{aligned}7.2973525692838009972851054723806629995\\[-2pt]641725396427\ldots{}^\circ\end{aligned}\)\\
\(C^*\) & \(2.1237383389686457242619513803933607937\ldots{}^\circ\)\\
\(\eta\) & \(0.299689683921630799684926562863927\ldots\)\\
\(\hret\) & \(1.05457181691503906113369058472430\ldots\times10^{-34}U_S\)\\
\bottomrule
\end{longtable}

\paragraph{Convergence quantitative selon la profondeur}

\begin{center}
\small
\begin{tabular}{r l l}
\toprule
Profondeur & \(\epsR^\circ\) & lecture\\
\midrule
12 & \(5.1681282186551375597\times10^{-8}\) & préasymptotique\\
18 & \(5.1383017764316644535\times10^{-8}\) & 7 cifras estables\\
24 & \(5.1383016758575394560\times10^{-8}\) & 14 cifras estables\\
30 & \(5.1383016758575282072\times10^{-8}\) & 20 cifras estables\\
36 & \(5.138301675857528207168517\times10^{-8}\) & perfil largo\\
45 & \(5.13830167585752820716851646226459\times10^{-8}\) & largeur \(<4.81\times10^{-130}\)\\
\bottomrule
\end{tabular}
\end{center}

Les retours de profondeur \(9,18,27,36,45\) conservent la phase \(8\) et
augmentent la mémoire \(0,1,2,3,4\). Le rapport entre largeurs consécutives tend
vers
\[
3^{-54}=1.7196982334549856127610399316004871\ldots\times10^{-26}.
\]

\paragraph{Analyse de sensibilité structurelle par suppression contrôlée}

\begin{longtable}{p{.46\textwidth} p{.34\textwidth}}
\toprule
Opération & Sortie ou différence unitaire\\
\midrule
\endfirsthead
\toprule
Opération & Sortie ou différence unitaire\\
\midrule
\endhead
Opérateur canonique & \(8.96802822044562981999\times10^{-10}\)\\
Un feuillet, \(\rho=90/729\) & \(-1.3727056615960678\times10^{-5}\)\\
Phase \(\theta_1=20^\circ\) & \(+5.2359877649510169\times10^{-1}\)\\
Phase \(\theta_3=80^\circ\) & \(-5.2359877470149605\times10^{-1}\)\\
Sans canal \(e\) & \(1.8065350748904653\times10^{-3}\)\\
Torsion historique \(\Delta_4(A,C^*)\) & \(1.0517084608768816\times10^{-9}\)\\
Opérateur harmonique \(M_4\) moins canonique & \(1.0525500222895070\times10^{-17}\)\\
Queue sexadique moins canonique & \(2.0283936357433839\times10^{-12}\)\\
\bottomrule
\end{longtable}

La séparation des échelles montre une sensibilité structurelle : la clôture dépend
de la phase, des deux feuillets, des canaux \(\pi/e\), de la torsion globale et de la retenue.
Aucune ablation ne reproduit la valeur.

\paragraph{Certificats informatiques reproductibles}

L'ouvrage utilise les propriétaires exécutables suivants :
\begin{enumerate}
\item \sourcepath{output/RADION_TWOWAY_NONADICO_20260827/verificar_emergencia_radion_por_profundidad.py};
\item \sourcepath{output/RADION_TWOWAY_NONADICO_20260827/verificar_radion_helice_critica_tpk.py};
\item \sourcepath{output/RADION_TWOWAY_NONADICO_20260827/verificar_operador_tpk_target_free_profundidad.py};
\item le vérificateur consolidé de cette monographie, situé dans
\sourcepath{certificados/verificar_monografia_radion.py}.
\end{enumerate}
Les sorties attendues comprennent
\[
\texttt{PASS\_EMERGENCIA\_RADION\_POR\_PROFUNDIDAD},
\]
\[
\texttt{PASS\_RADION\_HELICE\_CRITICA\_TPK},
\qquad
\texttt{PASS\_MONOGRAFIA\_RADION\_HMT\_MD}.
\]

\paragraph{Distinction entre précision informatique et exactitude mathématique}

Qu'un programme affiche zéro à 120 ou 200 chiffres ne transforme pas à lui seul une
identité résiduelle en prédiction. L'exactitude mathématique procède de
\eqref{eq:cierre-unitario} ; la précision arbitraire permet d'évaluer ses
opérandes et de vérifier la convergence. La non-circularité procède de l'ordre
de construction et de l'interdiction des cibles, non du nombre de chiffres.

Le calcul Decimal peut laisser des résidus d'ordre \(10^{-218}\) lorsqu'il combine
des branches évaluées avec des précisions différentes. Ce nombre est un arrondi de
l'implémentation ; non une correction physique supplémentaire.
% Fuente: ../incoming/radion_monografia_20260827/manuscrito/sections/D_historia_fuentes.tex:1-169
\subsubsection{Généalogie conceptuelle et sources primaires}

\paragraph{Quantum d'action de Planck}

Planck introduit \(h\) dans la loi de rayonnement comme constante reliant énergie
et fréquence. La donnée historique décisive pour cet ouvrage est dimensionnelle :
\(h\) est une action. Le radion ne modifie pas ce rôle ; il explique pourquoi l'action
d'un tour et l'action par unité de phase sont liées par \(2\pi\).

\paragraph{Quantification angulaire de Bohr et de Sommerfeld}

Bohr utilise \(h/2\pi\) dans la quantification du moment cinétique. Sommerfeld étend
la quantification aux intégrales d'action et aux orbites képlériennes. Il n'est donc pas
historiquement exact de dire que Dirac a inventé la constante réduite. Sa
contribution appartient à une généalogie dans laquelle l'action angulaire était déjà
présente.

La nouveauté HMT--MD est différente : elle ne prend pas \(h/2\pi\) comme un quotient externe,
mais construit une ellipse dont le tour a pour aire \(h\) et dont le secteur d'un
radian a pour aire \(\hbar\).

\paragraph{Formalisme matriciel de Born et de Jordan et oscillateur quantique}

Born et Jordan placent le principe variationnel, les phases et l'oscillateur
harmonique au cœur de la mécanique matricielle ; le champ électromagnétique se
décompose en une collection d'oscillateurs. La réalisation LC du radion
prolonge cette intuition avec une donnée supplémentaire : le quotient de forme
\(e^{-\eta}\) détermine l'impédance relative tandis que produit, fréquence
et action restent fixes.

\paragraph{Fréquence angulaire et action dans la formulation de Dirac}

En 1925, Dirac écrit les fréquences en radians par unité de temps et utilise
des relations de la forme
\[
xy-yx=\frac{ih}{2\pi}[x,y],
\qquad
\Delta E=\frac{h}{2\pi}\omega.
\]
L'équation \(E=\hret\omega\) relie explicitement action réduite et
fréquence angulaire. Dans le radion, cette relation reçoit une géométrie : \(\hret\)
est l'aire balayée par une unité de phase.

\paragraph{Relations d'indétermination de Heisenberg, Robertson et Schrödinger}

Heisenberg formule la limitation conjointe des variables canoniquement
conjuguées. Robertson généralise l'inégalité ; Schrödinger intègre la
covariance. La matrice \eqref{eq:covariance} relève plus directement de cette
prolongation Robertson--Schrödinger que de l'article de Heisenberg pris isolément.

Aucun de ces auteurs ne formule l'ellipse généalogique complète de cet ouvrage.
L'affirmation historiquement défendable est qu'ils fournissent le langage dans
lequel l'ellipse d'action peut être reconnue comme un contour minimal ; la
construction du couple \((A,C^*)\), les foyers, Klein, l'hélice et la retenue constituent la
contribution HMT--MD.

\paragraph{Fonctionnelle de phase de Feynman}

La formulation spatio-temporelle attribue à chaque chemin une phase
\(e^{iS/\hbar}\). L'unité de radion rend la correspondance littérale :
incrémenter l'action de \(\hbar\) incrémente la phase d'un radian. La somme des
chemins reconnaît, dans un langage intégral, le même principe que l'ellipse
exprime comme aire.

\paragraph{Propagations avancée et retardée chez Wheeler--Feynman}

La théorie de l'absorbeur construit une interaction classique temporellement
symétrique au moyen de combinaisons de champs retardés et avancés. L'identité
HMT \(CU(t)C^{-1}=U(-t)\) et l'inversion des chemins fournissent un
antécédent opératoire. La promotion physique exige que l'application entre chemins
et propagateurs préserve action, orientation, frontière et forme symplectique.

La théorie de l'absorbeur n'est pas la même affirmation que l'interprétation de
Feynman et de Stückelberg des antiparticules comme propagation temporellement
inversée. Les deux peuvent reconnaître la bidirectionnalité ; leurs objets physiques et
leurs conditions de frontière sont différents.

\paragraph{Dualité de Hodge et séparation d'avec la cohomologie de Hochschild}

L'opération pertinente dans ce manuscrit est l'étoile de Hodge
\(\star^2=-I\) sur les deux-formes en signature lorentzienne. On n'a pas localisé dans
le corpus actif du radion de propriétaire d'homologie de Hochschild qui
agisse sur l'ellipse, la retenue ou le secteur \(90/120\). Introduire Hochschild
par ressemblance nominale dégraderait le typage. Si une prolongation future
construit un cocycle de déformation avec ses propres domaine et codomaine, il devra
entrer comme résultat nouveau, non comme rétroattribution.

\paragraph{Bibliographie primaire sélectionnée}

\begin{enumerate}
\renewcommand{\labelenumi}{[\arabic{enumi}]}
\item\hypertarget{radionbib:Planck1901}{}
M.~Planck,
\emph{Ueber das Gesetz der Energieverteilung im Normalspectrum},
Annalen der Physik 309 (1901), 553--563.
\url{https://doi.org/10.1002/andp.19013090310}

\item\hypertarget{radionbib:Bohr1913}{}
N.~Bohr,
\emph{On the Constitution of Atoms and Molecules},
Philosophical Magazine 26 (1913), 1--25.
\url{https://doi.org/10.1080/14786441308634955}

\item\hypertarget{radionbib:Sommerfeld1916}{}
A.~Sommerfeld,
\emph{Zur Quantentheorie der Spektrallinien, I},
Annalen der Physik 356 (1916), 1--94.
\url{https://doi.org/10.1002/andp.19163561702}

\item\hypertarget{radionbib:BornJordan1925}{}
M.~Born et P.~Jordan,
\emph{Zur Quantenmechanik},
Zeitschrift für Physik 34 (1925), 858--888.
\url{https://doi.org/10.1007/BF01328531}

\item\hypertarget{radionbib:Dirac1925}{}
P.~A.~M.~Dirac,
\emph{The Fundamental Equations of Quantum Mechanics},
Proceedings of the Royal Society A 109 (1925), 642--653.
\url{https://doi.org/10.1098/rspa.1925.0150}

\item\hypertarget{radionbib:Heisenberg1927}{}
W.~Heisenberg,
\emph{Über den anschaulichen Inhalt der quantentheoretischen Kinematik und Mechanik},
Zeitschrift für Physik 43 (1927), 172--198.
\url{https://doi.org/10.1007/BF01397280}

\item\hypertarget{radionbib:Dirac1928}{}
P.~A.~M.~Dirac,
\emph{The Quantum Theory of the Electron},
Proceedings of the Royal Society A 117 (1928), 610--624.
\url{https://doi.org/10.1098/rspa.1928.0023}

\item\hypertarget{radionbib:Robertson1929}{}
H.~P.~Robertson,
\emph{The Uncertainty Principle},
Physical Review 34 (1929), 163--164.
\url{https://doi.org/10.1103/PhysRev.34.163}

\item\hypertarget{radionbib:Schrodinger1930}{}
E.~Schrödinger,
\emph{Zum Heisenbergschen Unschärfeprinzip},
Sitzungsberichte der Preussischen Akademie der Wissenschaften (1930), 296--303.

\item\hypertarget{radionbib:WheelerFeynman1945}{}
J.~A.~Wheeler et R.~P.~Feynman,
\emph{Interaction with the Absorber as the Mechanism of Radiation},
Reviews of Modern Physics 17 (1945), 157--181.
\url{https://doi.org/10.1103/RevModPhys.17.157}

\item\hypertarget{radionbib:Feynman1948}{}
R.~P.~Feynman,
\emph{Space-Time Approach to Non-Relativistic Quantum Mechanics},
Reviews of Modern Physics 20 (1948), 367--387.
\url{https://doi.org/10.1103/RevModPhys.20.367}

\item\hypertarget{radionbib:Feynman1949}{}
R.~P.~Feynman,
\emph{The Theory of Positrons},
Physical Review 76 (1949), 749--759.
\url{https://doi.org/10.1103/PhysRev.76.749}
\end{enumerate}

\begin{narrativebox}
Les fondateurs contiennent les éléments : action, fréquence angulaire, variables
conjuguées, oscillateur, incertitude, phase des chemins et réversion temporelle.
La contribution du radion consiste à les réunir dans un objet généré : une ellipse
d'aire \(h\), un secteur d'aire \(\hbar\), un intérieur de profondeur infinie,
une orientation de feuillet et une retenue exacte de compatibilité.
\end{narrativebox}
% Fuente: ../incoming/radion_monografia_20260827/manuscrito/sections/E_conclusion.tex:1-107
\subsubsection{Synthèse finale : unité angulaire avec généalogie de phase, d'action et de mémoire}

Le radian ne disparaît pas. Il devient entièrement visible.

Dans son usage ordinaire, un radian est un rapport qui n'a pas besoin de se souvenir de la façon dont il a été
produit. L'arc et le rayon s'annulent dimensionnellement ; l'unité semble
nue. Cette nudité est une vertu pour la trigonométrie, mais c'est un oubli
pour une théorie qui entend expliquer la structure du continu et l'origine
de l'action. Le radion récupère l'oublié sans modifier la valeur publiée.

La généalogie commence dans APP, où une sortie conserve la valeur et aussi
le feuillet, le quotient, le reste, la retenue et la frontière. Le TRIT introduit régime et
orientation. Le TPK fait revenir la phase sans réinitialiser l'état. La connexion
nonadique construit la section coinductive de \(\pi\), et la dodécaphase publie
la charnière de \(50^\circ\). Le couple angulaire apporte un mode commun \(A\) et une
ouverture orientée \(C^*\). L'incidence de six phases actives produit
\(90\), ses deux feuillets sur l'espace ambiant \(729\) produisent \(20/81\), et la
torsion globale apporte la seconde section.

Alors apparaît \(\varepsilon_R\). Non comme une décimale implorée par la
clôture, mais comme la différence que la structure était tenue de
transporter :
\[
\epsone=\frac{20}{81}\Delta_4-(S_E-1).
\]
La retenue publie ses chiffres ; la profondeur certifie sa limite ; la carte
angulaire le reconnaît comme
\[
5.13830167585752820716851646226459474899\ldots
\times10^{-8}\,{}^\circ.
\]
L'équation du radian est clôturée parce que les deux membres sont des publications du
même état :
\BigRadionEquation

L'histoire ne s'arrête pas à l'égalité. Le couple \((A,C^*)\) construit une
ellipse positive. Son produit de demi-axes fixe \(2\hret\) ; son quotient fixe la
forme. Le théorème d'aire démontre qu'une unité de phase balaie \(\hret\) et
que le tour enferme \(\hfull\). C'est ici que la division
par \(2\pi\) acquiert un contenu physique : \(h\) est l'action de tour ; \(\hbar\) est l'action par radion.

Les foyers cessent d'être des ornements d'une figure. Leur séparation reconstruit
\(C^*\) dans la carte spectrale, et leur distance de Klein mesure une profondeur
hyperbolique finie au sein d'un domaine dont l'aire hyperbolique totale est
infinie. La même frontière contient une action finie. Le paradoxe apparent
est la thèse centrale de la géométrie : une mesure clôt l'action ; une autre ouvre la
profondeur.

Le cercle n'est donc pas faux. C'est la projection exacte qui conserve la phase
et oublie l'anisotropie, l'intérieur et la mémoire. L'ellipse conserve l'action et la forme.
Klein conserve la profondeur. L'hélice conserve le retour et la biographie. Aucune
image ne remplace les autres ; toutes naissent du même état.

La réalisation de Hilbert reconnaît l'ellipse comme contour de covariance
minimale. L'oscillateur LC la reconnaît comme échange de réserves électrique et
magnétique. Le lecteur \(90/120\) publie perméabilité, permittivité,
impédance et vitesse. L'incidence six--huit produit un écran causal ;
Hodge convertit l'orientation en dualité des champs ; Lorentz convertit champ et
charge en dynamique ; Cartan--Holst reçoit l'holonomie en géométrie
gravitationnelle. L'ordre importe : l'objet HMT sélectionne ; la physique
conventionnelle réalise.

Le mot \emph{radion} se trouve ainsi justifié. Il ne signifie pas qu'un foyer soit
une charge positive et l'autre une charge négative. Il signifie que l'unité de
phase possède une orientation signée avant la publication électrique :
\[
\mathcal A(\mathfrak r_+)=+\hret,
\qquad
\mathcal A(\mathfrak r_-)=-\hret.
\]
La charge apparaît lorsque la mécanique dimensionnelle applique le transducteur. Le signe
n'est pas ajouté à la fin ; il est conservé depuis le feuillet.

On explique aussi pourquoi le secteur \(90/120\) semblait cacher de nombreuses
réponses. Incidence, extracteur dodécaphasique, queue spectrale, noyau de
Barbero, tours et retenue partagent une origine. C'est précisément parce qu'ils partagent une origine
qu'ils peuvent être comparés. C'est précisément parce qu'ils ont des opérations distinctes qu'ils ne doivent pas
être fusionnés. Le typage ne réduit pas l'unité du système : il la rend intelligible.

% REMATE_LOCAL_RADION_CIERRE
\Needspace{28\baselineskip}
La conclusion la plus compacte de cet ouvrage peut s'écrire en quatre lignes :
\[
\boxed{
\begin{aligned}
\text{phase :}\quad&\theta=1,\\
\text{action :}\quad&\mathcal A=\sigma\hret,\\
\text{mémoire :}\quad&H_9(n+9)=H_9(n)+(0,1),\\
\text{compatibilité :}\quad&1=S_E-\Phi_T+\epsone.
\end{aligned}}
\]

Le radion est cette unité quadruple. Le radian est son ombre exacte dans la carte
angulaire. La différence n'est pas dans la valeur de \(\pi\) ; elle réside dans la quantité de
structure que nous sommes disposés à conserver lorsque nous disons qu'un tour s'est
fermé.

\vspace{2\baselineskip}
\begin{center}
\begin{minipage}{.82\textwidth}
\centering\large\itshape\color{RadNavy}
Le cercle publie la phase.\\
L'ellipse contient l'action.\\
Klein ouvre la profondeur.\\
L'hélice se souvient du chemin.\\[.7em]
Le radion est l'unité qui n'oblige pas à choisir entre elles.
\end{minipage}
\end{center}



\HMTFlushCurrentChapter
\chapter{Classification généalogique des nombres : sections survivantes, invariants et atlas de constantes}%
\begingroup
\renewcommand{\chapter}[2][]{}%
\chapter{Classification généalogique des nombres : sections survivantes, invariants et atlas des constantes}
\label{chap:p3-final:52:atlas_numero}
\label{chap:p3-83-44-clasificacion_numero}


\label{chap:cobertura-taxonomia-numeros}
\index{nombre HMT!taxonomie}
\index{transducteur de Hensel!décalage complet}
\index{nombre rationnel!périodicité}

Le \cref{chap:numero-medicion} a défini un nombre HMT comme une section
compatible du système inverse de survie, et non comme une suite de
chiffres dépourvue de généalogie. Il restait cependant deux questions qu'il
convient de séparer avec précision. La première est une question de
\emph{représentation} : quelles valeurs réelles le transducteur
3-adique peut-il représenter avant d'imposer les filtres particuliers de clôture ? La seconde est
une question de \emph{sélection} : quelles sections de l'espace brut sont
choisies par APP, TRIT et TPK lorsque l'on exige simultanément leurs
compatibilités ? Cette section résout la première, formule la seconde comme
une restriction explicite du domaine et extrait de la carte de base mille une
taxonomie complète des valeurs visibles. Les outils de base sont
l'arithmétique \(p\)-adique, la dynamique symbolique et la caractérisation positionnelle
des rationnels \cite{Koblitz1984,LindMarcus1995,HardyWright2008}.

La distinction est structurelle. Un langage capable de représenter tout ruban ne
sélectionne pas de ce seul fait le ruban de \(\pi\), de \(e\), de \(\varphi\) ou
de \(\alpha\). À l'inverse, une loi de sélection ne peut pas être évaluée si
l'espace dans lequel elle agit n'a pas été construit auparavant. Représentation et sélection
sont deux théorèmes consécutifs, et non deux noms pour une même étape.

\section{Transducteur de Hensel non filtré et décalage complet sur les fibres ternaires}

Soit \(p\) un nombre premier, soit \(d\geq1\) et soit
\(A\in\operatorname{GL}_d(\mathbb Z)\). Pour un vecteur ligne
\(x_t\in\mathbb Z_p^d\), nous définissons
\begin{equation}
 y_t=x_tA,
 \qquad
 u_t=y_t\bmod p\in\mathbb F_p^d,
 \qquad
 x_{t+1}=\frac{y_t-u_t}{p}.
 \label{eq:transductor-hensel-general}
\end{equation}
Le bloc \(u_t\) est l'émission visible et \(x_{t+1}\) est la mémoire qui
subsiste après retrait de ce chiffre.

\begin{theorem}[Conjugaison de Hensel avec le décalage complet]
\label{thm:shift-hensel-completo}
Pour tout \(T\geq1\), l'application
\[
 \Psi_T:(\mathbb Z/p^T\mathbb Z)^d
 \longrightarrow(\mathbb F_p^d)^T,
 \qquad
 x_0\longmapsto(u_0,\ldots,u_{T-1}),
\]
est bijective. Son inverse est
\begin{equation}
 x_0\equiv
 \sum_{t=0}^{T-1}p^t u_tA^{-(t+1)}
 \pmod {p^T}.
 \label{eq:inversa-hensel-finita}
\end{equation}
En passant à la limite inverse, on obtient un homéomorphisme
\[
 \Psi:\mathbb Z_p^d\xrightarrow{\ \sim\ }
 (\mathbb F_p^d)^{\mathbb N}.
\]
Si \(S_A\) désigne la mise à jour \(x_t\mapsto x_{t+1}\) et
\(\sigma\) supprime le premier bloc d'un ruban, alors
\[
 \Psi\circ S_A=\sigma\circ\Psi.
\]
\end{theorem}

\begin{proof}
De \cref{eq:transductor-hensel-general}, on obtient
\[
 x_t=(u_t+px_{t+1})A^{-1}.
\]
L'itération de cette identité donne
\[
 x_0=
 \sum_{t=0}^{T-1}p^t u_tA^{-(t+1)}
 +p^Tx_TA^{-T}.
\]
Le dernier terme s'annule modulo \(p^T\), ce qui prouve qu'un mot de
\(T\) blocs détermine au plus une classe initiale. Réciproquement, étant donné
un mot quelconque, on fixe \(x_T=0\) et l'on reconstruit à rebours au moyen de
\(x_t=(u_t+px_{t+1})A^{-1}\). L'égalité
\(x_tA=u_t+px_{t+1}\) vérifie que la classe obtenue émet le mot
prescrit. Par conséquent, \(\Psi_T\) est bijective.

Les bijections sont compatibles avec la réduction modulo \(p^T\). Leurs limites
inverses produisent l'homéomorphisme \(\Psi\) ; de manière équivalente, la série de
\cref{eq:inversa-hensel-finita} converge dans \(\mathbb Z_p^d\). Appliquer
\(S_A\) retire exactement le bloc \(u_0\), ce qui correspond donc à
appliquer \(\sigma\) au ruban.
\end{proof}

Dans l'instance HMT, \(p=3\), \(d=6\) et la matrice publiée a pour
déterminant un. C'est pourquoi
\begin{equation}
 \mathbb Z_3^6\simeq(\mathbb F_3^6)^{\mathbb N},
 \qquad |\mathbb F_3^6|=3^6=729.
 \label{eq:hmt-shift-729}
\end{equation}
Chaque ruban de blocs de six trits possède une unique mémoire 3-adique initiale.
La matrice détermine la mise en coordonnées et la dynamique de lecture ;
l'universalité de l'alphabet démontre qu'elle n'exclut aucun ruban.

\begin{corollary}[Entropie et mesure de Haar]
La mise à jour de Hensel brute est conjuguée au décalage unilatéral
complet de
\(729\) symboles. Son entropie topologique est \(6\log3\). La mesure de Haar de
\(\mathbb Z_3^6\) est transportée vers la mesure produit uniforme sur les
blocs.
\end{corollary}

\begin{proof}
Il existe \(729^T=3^{6T}\) cylindres de longueur \(T\), et chaque classe modulo
\(3^T\) a une masse de Haar \(3^{-6T}\). La conjugaison du
\cref{thm:shift-hensel-completo} transporte les deux propriétés vers le
décalage.
\end{proof}

La conclusion probabiliste concerne l'espace brut muni de Haar. Elle
n'affirme pas que les sections publiées des constantes ont été échantillonnées
selon cette mesure ni qu'elles sont algorithmiquement aléatoires.

\section{Carte ternaire équilibrée des entiers}

La carte 3-adique brute et la carte archimédienne n'épuisent pas les types numériques
de la classification HMT. Pour l'opérateur entier, on utilise le mot vide ainsi que les
mots réduits
\[
 \mathscr W_{\mathrm{bal}}
 :=\{\varnothing\}\cup
 \{a_1\cdots a_m:
 a_j\in\{-1,0,+1\},\ a_1\ne0\}.
\]
Nous définissons récursivement
\[
 b(\varnothing)=0,
 \qquad
 b(wa)=3b(w)+a.
\]

\begin{theorem}[Représentation ternaire équilibrée unique]
\label{pf84:C03}
L'évaluation
\[
 b:\mathscr W_{\mathrm{bal}}\longrightarrow\ZZ
\]
est bijective. Les mots dont le premier trit est \(+1\) représentent les
entiers positifs et ceux qui commencent par \(-1\), les négatifs.
\end{theorem}

\begin{proof}
Pour \(n\ne0\), le dernier trit \(a\) est le représentant unique de
\(n\bmod3\) dans \(\{-1,0,+1\}\). Le parent
\[
 \frac{n-a}{3}
\]
a une valeur absolue strictement inférieure. L'itération se termine en \(0\), de
sorte que tout entier possède un développement fini. L'unicité du résidu
équilibré à chaque pas prouve l'unicité du mot. Le signe est
fixé par le premier trit non nul.
\end{proof}

\begin{remark}
Cette bijection est une carte arithmétique exacte. Elle n'identifie pas tous les
mots équilibrés à des sections survivantes TPK ni ne remplace la
généalogie, la retenue ou l'orientation d'une route enrichie.
\end{remark}

\section{Surjectivité de la carte archimédienne non filtrée sur \(\mathbb R\)}

Concaténons les coordonnées des blocs
\(u_t\in\mathbb F_3^6\) pour obtenir un ruban
\((a_n)_{n\geq1}\in\{0,1,2\}^{\mathbb N}\). Son évaluation ternaire est
\[
 \operatorname{ev}_3((a_n))=
 \sum_{n\geq1}\frac{a_n}{3^n}\in[0,1].
\]
Nous adoptons la convention canonique qui exclut la représentation avec une queue
ultimement constante de deux seulement lorsque le même point possède un autre développement
dans la carte. L'extrémité \(1\) conserve donc son unique représentation
fractionnaire \(0.222\ldots{}_3\).

\begin{theorem}[Couverture archimédienne brute]
\label{thm:cobertura-arquimediana-bruta}
L'application
\[
 P_{\rm raw}:\mathbb Z_3^6
 \xrightarrow{\Psi}(\mathbb F_3^6)^{\mathbb N}
 \xrightarrow{\operatorname{flat}}\{0,1,2\}^{\mathbb N}
 \xrightarrow{\operatorname{ev}_3}[0,1]
\]
est continue et surjective. Par conséquent,
\[
 \mathbb Z_3^6/\!\sim_{P_{\rm raw}}
 \ \cong\ [0,1]
\]
comme espace de valeurs, où \(x\sim_{P_{\rm raw}}y\) si leurs évaluations
coïncident.
\end{theorem}

\begin{proof}
Tout chiffre ternaire appartient à un unique bloc consécutif de six chiffres,
de sorte que \(\operatorname{flat}\) est une bijection d'espaces de rubans.
Toute \(r\in[0,1]\) possède un développement ternaire ; le
\cref{thm:shift-hensel-completo} fournit l'unique mémoire 3-adique qui
émet ses blocs. La continuité se vérifie par cylindres : fixer \(T\)
blocs fixe \(6T\) chiffres et enferme la valeur dans un intervalle de diamètre
\(3^{-6T}\). L'identification du quotient à l'image est la factorisation
canonique d'une application surjective par ses fibres.
\end{proof}

Une seule carte 3-adique est compacte et ne peut donc pas couvrir de manière
continue une droite non bornée. La globalisation correcte utilise une famille
dénombrable de cartes.

\begin{corollary}[Couverture archimédienne dénombrable de \(\mathbb R\)]
\label{cor:atlas-hmt-reales}
Soit
\[
 \widehat X_{\rm raw}=\mathbb Z\times\mathbb Z_3^6,
 \qquad
 \widehat P_{\rm raw}(m,x)=m+P_{\rm raw}(x).
\]
Alors \(\widehat P_{\rm raw}\) est surjective sur \(\mathbb R\).
En conséquence, tout réel admet au moins une généalogie de Hensel brute et
\[
 \widehat X_{\rm raw}/\!\sim_{\widehat P_{\rm raw}}
 \ \cong\ \mathbb R
\]
comme ensemble de valeurs.
\end{corollary}

\begin{proof}
Pour \(r\in\mathbb R\), prenons \(m=\lfloor r\rfloor\) et appliquons le
\cref{thm:cobertura-arquimediana-bruta} à \(r-m\in[0,1)\).
\end{proof}

Ce résultat précise l'affirmation selon laquelle \(\mathbb R\) est l'image
archimédienne quotient d'une
structure généalogique : la projection oublie la mémoire 3-adique et conserve
le point évalué. L'affirmation démontrée ici concerne la complétion
\emph{brute}. Si \(X_{\rm surv}\subseteq\mathbb Z_3^6\) désigne le sous-espace
qui satisfait en outre une famille concrète de clôtures APP--TRIT--TPK, alors
\[
 P_{\rm raw}(X_{\rm surv})=[0,1]
\]
est un théorème supplémentaire. Le transducteur garantit une capacité complète ; la
survie détermine quelle partie de cette capacité est réalisée par le système
restreint.

On ne déduit pas non plus ici que la somme et le produit natifs de TPK descendent à
tout le quotient global. On peut transporter l'arithmétique usuelle depuis
\(\mathbb R\), mais démontrer qu'elle coïncide avec des opérations construites avant
l'évaluation exige le théorème de descente correspondant.

Dans cette construction, les cylindres ne sont pas des infinitésimaux non nuls de
\(\mathbb R\). Chaque cylindre de profondeur finie a un diamètre positif,
et la suite des diamètres converge vers zéro. Deux sections qui coïncident jusqu'à
la profondeur \(T\) et diffèrent ensuite forment une \emph{perturbation de
queue} ; leur différence archimédienne est bornée par le diamètre du cylindre
commun. C'est la formulation exacte des ``cylindres qui tendent vers zéro''
dans la droite réelle ordinaire.

\section{Carte positionnelle en base \(10^3\) et classification des sections publiées}

La complétion cubique
\(1000=729+243+27+1\) rend naturel le regroupement des chiffres décimaux en blocs
\(b_n\in\{0,\ldots,999\}\). Pour un ruban de blocs, nous définissons
\[
 \operatorname{ev}_{1000}(b_1,b_2,\ldots)
 =\sum_{n\geq1}\frac{b_n}{1000^n}.
\]
Le développement canonique exclut une queue ultimement constante de blocs
\(999\). Cette convention résout la duplicité des extrémités de cylindre.

\begin{theorem}[Taxonomie visible en base mille]
\label{thm:taxonomia-base-mil}
Soit \(x\in[0,1)\) et soit \((b_n)\) son développement canonique en base mille.
Alors :
\begin{enumerate}
\item \(x\) possède un développement fini si et seulement si, écrit comme
fraction irréductible \(a/q\), on a \(q\mid1000^N\) pour un certain \(N\) ;
de manière équivalente, les seuls nombres premiers qui divisent \(q\) sont \(2\) et \(5\).
\item \(x\) est rationnel si et seulement si \((b_n)\) est ultimement périodique.
\item \(x\) est irrationnel si et seulement si \((b_n)\) n'est pas ultimement
périodique.
\end{enumerate}
\end{theorem}

\begin{proof}
Un développement qui se termine à la profondeur \(N\) a la forme
\(D/1000^N\) ; une fois réduit, son dénominateur ne contient que les nombres premiers \(2\) et
\(5\). Réciproquement, si \(q\mid1000^N\), alors \(a/q\) s'écrit avec
\(N\) blocs.

Si, après un préfixe de longueur \(r\), se répète une période de longueur
\(s\) dont l'entier en base mille est \(P\), la queue est une série géométrique :
\[
 \frac{P}{1000^r(1000^s-1)},
\]
et la valeur est rationnelle. À l'inverse, en développant \(a/q\), les restes de la
division successive ne peuvent prendre que \(q\) valeurs. Un reste nul produit un
développement fini ; la répétition d'un reste produit une période. La troisième
affirmation est la négation exacte de la seconde.
\end{proof}

Le mot \emph{fraction} désigne une présentation \(a/q\) d'un rationnel,
et non une espèce numérique supplémentaire. Il ne faut pas non plus confondre périodicité et
algébricité. Un réel est algébrique lorsqu'il annule un polynôme non nul de
\(\mathbb Q[X]\) ; il est transcendant lorsqu'il ne le fait pas. Cette seconde
classification est indépendante de la base utilisée pour écrire les chiffres.
Par exemple,
\[
 \varphi^2-\varphi-1=0
\]
fait de \(\varphi\) un irrationnel algébrique, tandis que la transcendance
de \(e\) et \(\pi\) découle de théorèmes classiques externes à HMT\@.

\section{Concaténation, réduction nonadique et retenue}

La compatibilité entre la carte cubique et la racine numérique est une identité
arithmétique, et non une ressemblance visuelle. Si \(u_j\in\{0,\ldots,999\}\) et
\[
 D_n=1000D_{n-1}+u_n,
 \qquad D_0=0,
\]
alors \(D_n\) est l'entier obtenu en concaténant les \(n\) premiers
blocs.

\begin{theorem}[Congruence nonadique de contrôle de la carte cubique]
\label{thm:checksum-nonadico-base-mil}
Pour tout \(n\geq1\),
\[
 D_n\equiv\sum_{j=1}^n u_j\pmod9.
\]
Par conséquent, concaténation en base mille, réduction modulo neuf et racine numérique
sont compatibles.
\end{theorem}

\begin{proof}
Comme \(1000\equiv1\pmod9\), la récurrence donne
\(D_n\equiv D_{n-1}+u_n\pmod9\). L'itération à partir de \(D_0=0\) produit la
formule.
\end{proof}

Dans la clôture dodécaphasique de \(\alpha\), l'équation de retenue par bloc
a la forme
\[
 P_m+E_m-\Phi_m-K_m+c_{m+1}=A_m+1000c_m.
\]
Réduite modulo neuf,
\[
 A_m\equiv
 P_m+E_m-\Phi_m-K_m+c_{m+1}-c_m
 \pmod9.
\]
Si les retenues extrêmes s'annulent, la somme télescopique élimine la mémoire
intermédiaire. Le résidu global de la sortie coïncide alors avec le résidu
de la combinaison des entrées. Cette loi explique pourquoi le registre nonadique
contrôle la concaténation décimale ; elle ne sélectionne pas à elle seule les blocs qui
doivent être concaténés.

\section{Coordonnée nonadique relevée et retour de phase}

La réduction modulo neuf ne doit pas être confondue avec l'identification des
entiers qui possèdent le même résidu. Soit
\[
 I_9=\{1,2,\ldots,9\}.
\]
Pour chaque entier positif \(n\), nous définissons sa \emph{phase nonadique} et son
\emph{indice de tour} par
\begin{equation}
 \rho_9(n)=1+((n-1)\bmod9),
 \qquad
 \kappa_9(n)=\frac{n-\rho_9(n)}9.
 \label{eq:coordenada-nonadica-levantada}
\end{equation}
La première coordonnée est la racine numérique écrite dans l'alphabet
\(I_9\), de sorte que la classe résiduelle zéro est représentée par la phase \(9\).
La seconde enregistre combien de tours complets ont précédé la phase
visible.

\begin{theorem}[Décomposition nonadique avec mémoire]
\label{thm:descomposicion-nonadica-memoria}
L'application
\[
 \Theta_9:\mathbb N_{>0}\longrightarrow\mathbb N_0\times I_9,
 \qquad
 n\longmapsto\bigl(\kappa_9(n),\rho_9(n)\bigr),
\]
est bijective, d’inverse
\[
 \Theta_9^{-1}(q,r)=9q+r.
\]
Dans ces coordonnées, le successeur arithmétique est conjugué à la dynamique
\begin{equation}
 \mathcal S_9(q,r)=
 \begin{cases}
  (q,r+1),&1\leq r<9,\\
  (q+1,1),&r=9.
 \end{cases}
 \label{eq:sucesor-nonadico-levantado}
\end{equation}
En particulier, le pas visible \(9\to1\) est un retour de phase accompagné
par l'incrément irréversible \(q\to q+1\) ; ce n'est pas une égalité entre les
entiers \(9\) et \(10\).
\end{theorem}

\begin{proof}
La division euclidienne de \(n-1\) par \(9\) fournit de manière unique
\(n-1=9q+s\), avec \(q\in\mathbb N_0\) et \(0\leq s<9\). En prenant
\(r=s+1\), on obtient \(n=9q+r\), avec \(r\in I_9\), ce qui démontre à la fois
l'existence, l'unicité et la formule inverse. Si \(r<9\), ajouter une unité
incrémente seulement \(r\). Si \(r=9\), l'égalité
\(9q+9+1=9(q+1)+1\) donne la seconde branche de
\cref{eq:sucesor-nonadico-levantado}.
\end{proof}

Ce relèvement sépare trois notions que le chiffre isolé mélange :
\begin{enumerate}
\item l'unité d'avancée, réalisée par le successeur ;
\item la phase visible \(r\), qui parcourt le cycle nonadique ;
\item la mémoire accumulée \(q\), qui distingue les retours d'une même phase.
\end{enumerate}
La projection \((q,r)\mapsto r\) oublie le nombre de tours et satisfait
\[
 \rho_9(n+9)=\rho_9(n),
 \qquad
 \kappa_9(n+9)=\kappa_9(n)+1.
\]
C'est le modèle arithmétique élémentaire du mécanisme qui, dans l'état TPK
enrichi, conserve en outre les quotients de Hensel, l'orientation, la signature et la donnée
de frontière. Un retour de mot visible peut ainsi coexister avec une
généalogie distincte.

Zéro n'appartient pas à \(I_9\) : on l'ajoute comme point de base séparé avant
d'appliquer \(\Theta_9\). Pour les entiers négatifs, on conserve le module
positif de la décomposition et l'on joint une orientation
\(\varepsilon\in\{-1,+1\}\). En conséquence, ni le signe ni zéro ne se
déduisent de la phase résiduelle. Ce typage évite de transformer une carte
cyclique en une identité arithmétique fausse.

\begin{theorem}[Entier orienté et signe généalogique]
\label{thm:entero-orientado-signo-genealogico}
L'application
\[
\Theta_9^\pm:
\mathbb Z\setminus\{0\}
\longrightarrow
\{-1,+1\}\times\mathbb N_0\times I_9,
\]
\[
n\longmapsto
\bigl(\operatorname{sgn}(n),
      \kappa_9(|n|),
      \rho_9(|n|)\bigr)
\]
est bijective, d’inverse
\[
(\varepsilon,q,r)\longmapsto\varepsilon(9q+r).
\]
La négation arithmétique agit par
\[
(\varepsilon,q,r)\longmapsto(-\varepsilon,q,r):
\]
inverse l'orientation et conserve la grandeur, la phase et la mémoire des tours.
\end{theorem}

\begin{proof}
Le \cref{thm:descomposicion-nonadica-memoria} fournit de manière unique
\(|n|=9q+r\), avec \(q\in\mathbb N_0\) et \(r\in I_9\). Le signe d'un entier
non nul est également unique. La formule inverse recompose \(n\), et remplacer
\(n\) par \(-n\) ne change que \(\varepsilon\).
\end{proof}

Ainsi, un négatif n'est ni une phase résiduelle spéciale ni une quantité
``impossible''. C'est la même grandeur nonadique transportée avec une orientation
généalogique opposée. Cette structure sera prolongée en mécanique dimensionnelle
par l'antisection particule--antiparticule.

Itérer \(\kappa_9\) produit une tour finie pour tout entier ordinaire ; dans
la complétion inverse, une suite compatible de quotients peut être
infinie. La différence entre les deux cas anticipe la taxonomie HMT : la valeur
visible répond au développement positionnel, tandis que le type généalogique
répond à l'évolution complète de ses mémoires.

\section{2 axes de classification : valeur visible et généalogie de transport}

La taxonomie précédente classe la valeur projetée. HMT ajoute un second
axe que l'écriture décimale ordinaire écarte. Pour une section
\(x\in\Omega_\infty^{\rm Arch}\), nous distinguons :
\[
 \begin{array}{rcl}
 \text{type visible}&=&
 \text{fini, périodique ou non périodique ; algébrique ou transcendant},\\
 \text{type généalogique}&=&
 \text{classe de retenue, orientation, fermeture et monodromie de }x.
 \end{array}
\]
Deux sections peuvent partager le même réel et différer selon le second axe. Un
réel rationnel possède un développement visible ultimement périodique, mais une
présentation enrichie redondante peut conserver une mémoire interne non
périodique que l'évaluation oublie. La périodicité de la \emph{valeur} se
rapporte toujours à son développement canonique ; la périodicité de la
\emph{généalogie} se rapporte à l'état enrichi de TPK.

Cette séparation fournit une définition concise :
\begin{equation}
 \boxed{
 \text{nombre HMT}
 =
 \text{section compatible avec mémoire};
 \qquad
 \text{valeur réelle}
 =
 \text{sa classe d'évaluation}.}
 \label{eq:numero-hmt-y-valor-real}
\end{equation}
Mesurer dans cette carte consiste à coupler la section au lecteur, à fixer un cylindre et à
conditionner les prolongements compatibles. Augmenter la précision signifie
rétrécir le cylindre, et non consommer la généalogie qui reste dans la fibre.

\section{Mémoire de profondeur arbitraire dans les systèmes d'émission finie}

La récurrence visible de neuf phases ne peut coexister avec une sortie
irrationnelle que si l'état complet conserve une information qui ne se réduit pas à
ces neuf étiquettes. Le résultat suivant rend cette nécessité exacte.

\begin{theorem}[Caractérisation par émetteurs finis]
\label{thm:emisor-finito-racional}
Soit \(Q\) un ensemble fini, \(F:Q\to Q\) une transition déterministe,
\(o:Q\to\{0,\ldots,999\}\) une sortie et \(q_0\in Q\). Le ruban
\(b_n=o(F^{n-1}(q_0))\) est ultimement périodique et son évaluation est rationnelle.
Réciproquement, tout ruban ultimement périodique admet un émetteur déterministe
fini.
\end{theorem}

\begin{proof}
Parmi \(|Q|+1\) états consécutifs, deux états égaux apparaissent. Dès la première
répétition, le déterminisme impose la répétition de tout le futur, de sorte que le
ruban est ultimement périodique. Le \cref{thm:taxonomia-base-mil} donne la
rationalité. Pour la réciproque, il suffit de construire une chaîne d'états pour le
préfixe et un cycle pour la période.
\end{proof}

Par conséquent, une loi qui engendre exactement \(e\), \(\pi\) ou tout autre
irrationnel ne peut pas consister uniquement en un automate déterministe à
états finis qui se répète sans mémoire supplémentaire. Elle peut utiliser une limite
inverse, une mémoire 3-adique, une profondeur croissante ou une entrée
apériodique ; tous ces mécanismes fournissent un état effectif non borné.
Dans HMT, le retour visible
\(g_{\rm vis}(n+9)=g_{\rm vis}(n)\) n'impose pas le retour de la retenue, de
l'orientation ni de la frontière. Cette différence est précisément l'endroit où
peut résider un développement non périodique.

\section{Le carré local de \texorpdfstring{\(e\)}{e} et la rupture de mémoire}

Pour un mot décimal \(a_1\cdots a_n\), nous appellerons \emph{carré local}
de longueur \(\ell\) à la position \(i\) une factorisation
\[
 a_i\cdots a_{i+2\ell-1}=vv,
 \qquad |v|=\ell.
\]
Cette notion relève de la combinatoire des mots. Elle n'affirme pas que la queue
infinie possède une période.

\begin{theorem}[Caractérisation finie du carré \(1828^2\)]
\label{thm:cuadrado-local-e}
Considérons les \(243\) mots du catalogue et le lecteur fixé
\(w_6\mapsto w_{30}\mapsto d_1d_2d_3d_4\), avec quatre triades en base mille.
Il existe un unique mot dont la lecture contient un carré de longueur quatre
commençant au deuxième chiffre décimal. C'est
\[
 w_e=201101,
 \qquad
 718281828459=7\,\underbrace{1828\,1828}_{(1828)^2}\,459.
\]
La répétition ne se prolonge pas en une troisième copie : le chiffre qu'exigerait cette
continuation serait \(1\), tandis que le chiffre effectif est \(4\). Dans tout le
catalogue, il n'existe qu'un autre carré de longueur quatre,
\[
 575157518472=\underbrace{5751\,5751}_{\text{carré local}}\,8472,
\]
qui commence au premier chiffre.
\end{theorem}

\begin{proof}
L'énumération exhaustive évalue les \(243\) mots par arithmétique
entière. Le profil des carrés pour les longueurs \(1,\ldots,6\) contient
\[
 240,\ 21,\ 3,\ 2,\ 0,\ 0
\]
occurrences, réparties entre
\[
 153,\ 19,\ 3,\ 2,\ 0,\ 0
\]
mots. Les deux occurrences de longueur quatre sont exactement celles de
l'énoncé. Deux énumérations entières indépendantes reproduisent le même recensement
et leurs sorties complètes coïncident terme à terme.
\end{proof}

Le mot \(w_e\) possède en outre la chaîne de relèvement
\[
201101\mid121221\mid102011\mid012222\mid102011.
\]
Les troisième et cinquième blocs visibles coïncident. En revanche,
les états qui les transportent ne coïncident pas. Avant les deux émissions, leurs réductions
modulo \(27\) sont
\[
(25,11,7,17,8,16),
\qquad
(7,8,4,23,26,7),
\]
et les quotients de Hensel ultérieurs sont
\[
(6,13,9,2,13,1),
\qquad
(15,0,3,19,23,25).
\]
Par conséquent, le retour \(102011\) est un retour de la projection visible, et non
de l'état enrichi. C'est la formulation exacte de la répétition locale
qui se rompt : le quotient conserve une généalogie que le bloc résiduel
n'enregistre pas.

En conséquence, le début
\[
e=2.718281828459\ldots
\]
ne présente pas de ``semi-périodicité'', mais un carré local de mot avec
retour visible et rupture de mémoire. Le résultat est exhaustif dans le cadre du
lecteur fixé ; son contenu est combinatoire et dynamique, et non une inférence
d'irrationalité à partir de douze chiffres.

\section[2 mesures finies de l'émetteur 468 à 243]{2 mesures finies de l'émetteur \texorpdfstring{\(468\to243\)}{468→243} : fréquence uniforme et priorité structurelle}

Le catalogue public distingue \(468\) émissions et \(104\,976\) germes.
Soit
\[
 q:E_{468}\longrightarrow W_{243}
\]
l'application qui oublie la signature d'émission et conserve le mot de six
trits. Il existe deux mesures naturelles, parce qu'il existe deux univers finis de dénombrement :
\[
 \mu_{468}(w)=\frac{|q^{-1}(w)|}{468},
 \qquad
 \mu_{\rm seed}(w)=
 \frac{\sum_{e\in q^{-1}(w)}m(e)}{104\,976},
\]
où \(m(e)\) est la multiplicité des germes de l'émission \(e\).

\begin{theorem}[Mesures canoniques relatives à l'univers dénombré]
\label{thm:medidas-emisor-468}
La mesure uniforme sur \(E_{468}\) est l'unique probabilité invariante sous
toutes les permutations de ses émissions ; sa mesure image est \(\mu_{468}\).
La mesure uniforme sur les \(104\,976\) germes induit
\(\mu_{\rm seed}\). Leurs histogrammes exacts sont
\[
\begin{array}{c|rrrrrr}
|q^{-1}(w)|&1&2&3&4&5&6\\ \hline
\#\{w\}&132&66&18&3&6&18
\end{array}
\]
et
\[
\begin{array}{c|rrrrrrrrr}
\text{poids des germes}
&144&288&432&576&864&1008&1440&1944&2088\\ \hline
\#\{w\}
&120&54&18&9&15&6&3&12&6.
\end{array}
\]
Pour les trois mots distingués par le lecteur, on obtient
\[
\begin{array}{c|cc}
&\mu_{468}&\mu_{\rm seed}\\ \hline
\pi&5/468&7/729\\
e&1/468&1/729\\
\varphi&1/468&1/729.
\end{array}
\]
\end{theorem}

\begin{proof}
L'invariance sous le groupe symétrique rend égales les masses de tous les
points de \(E_{468}\) ; la normalisation les fixe à \(1/468\). La formule de
la mesure image résulte de la sommation sur chaque fibre. Le même argument sur
l'ensemble des germes produit la seconde mesure. Les histogrammes s'obtiennent
par énumération exhaustive du catalogue fini de \(468\) émissions et de \(243\)
mots ; leurs sommes vérifient
\[
132+66+18+3+6+18=243,
\]
\[
132+2\cdot66+3\cdot18+4\cdot3+5\cdot6+6\cdot18=468
\]
et la somme totale des poids est \(104\,976\).
\end{proof}

La non-uniformité de la mesure image est un théorème fini : des mots distincts
ont des nombres distincts de réalisations. En revanche, il n'existe pas de
probabilité uniforme normalisable sur tout \(\mathbb R\), et ces mesures
n'ordonnent pas intrinsèquement tous les réels. Le mot de \(\pi\) appartient
à l'une des six fibres de taille cinq et de poids \(1008\) ; il n'est le maximum
unique d'aucun des deux histogrammes. Sa singularité démontrée procède
de la structure composée de ses cinq régions, de sa partition orientée
\(3+2\) et de ses compatibilités ultérieures, et non de la déclaration a priori d'une
probabilité absolue sur le continu.

\section{Sélecteur orbital fini antérieur au lecteur archimédien}
\label{sec:selector-orbital-constantes}

Les mesures précédentes révèlent des multiplicités, mais un mot isolé de
six trits dépend encore d'une orientation. Le catalogue possède une symétrie
interne qui permet de sélectionner d'abord des classes non orientées, puis de fixer
un représentant. Écrivons un mot comme \(abc\mid def\) et définissons
\[
 r(abc\mid def)=cab\mid efd,
 \qquad
 s(abc\mid def)=def\mid abc.
\]
Ces permutations agissent également sur les six coordonnées de la donnée
\(U_6\). Dans le catalogue complet, elles satisfont
\[
 r^3=1,
 \qquad s^2=1,
 \qquad srs=r^{-1},
\]
de sorte qu'elles engendrent une action que nous notons \(D_3^{\rm cat}\). L'action utilise
l'étiquetage TRIT fixe
\[
 \Psi(U_6)=(-U_6\bmod3)
\]
et la partition publiée des six coordonnées en deux triades.

Pour un mot \(w\), soient \(f(w)=|q^{-1}(w)|\) le nombre d'émissions,
\(m(w)\) sa masse de germes et
\[
 \nu(w)=(n_0(w),n_1(w),n_2(w))
\]
son recensement ordonné de trits. Soit en outre \(d(e)\) la classe diagonale additive d'une
émission et \(o(e)\in\{\mathrm{ES},\mathrm{NO}\}\) son orientation
publiée.

\begin{theorem}[Sélection orbitale interne des trois mots]
\label{thm:selector-orbital-constantes}
L'action \(D_3^{\rm cat}\) préserve l'ensemble des \(468\) émissions,
la projection \(q:E_{468}\to W_{243}\) et la multiplicité des germes. Son
action sur \(W_{243}\) possède exactement \(43\) orbites : \(38\) de taille
six et cinq de taille trois. Dans ce quotient, les unicités suivantes
sont vérifiées.

\begin{enumerate}
\item Il existe une unique orbite \(\mathcal O_\pi\) de taille six telle que, pour
tout \(w\in\mathcal O_\pi\),
\[
 f(w)=5,
 \qquad m(w)=1008,
 \qquad
 \{m(e):e\in q^{-1}(w)\}=\{432,144,144,144,144\}.
\]
C'est
\[
\mathcal O_\pi=
\{001112,010211,100121,112001,121100,211010\},
\]
et tous ses mots ont le recensement \(\nu=(2,3,1)\).

\item Considérons les orbites de taille six dont les mots ont une seule
émission, de multiplicité \(144\), et dont le support diagonal conjoint est
\(\{0,3,6\}\). Il y en a exactement six. Parmi elles, il existe une unique orbite
avec le même recensement \((2,3,1)\) que \(\mathcal O_\pi\) :
\[
\mathcal O_e=
\{011120,012110,101201,110012,120011,201101\}.
\]

\item Dans le même secteur, il existe une unique orbite équilibrée, c'est-à-dire avec
\(\nu=(2,2,2)\) :
\[
\mathcal O_\varphi=
\{002112,020211,112002,121200,200121,211020\}.
\]
\end{enumerate}

Enfin, la clause d'orientation
\begin{equation}
 \exists e\in q^{-1}(w):
 \qquad d(e)=6,
 \qquad o(e)=\mathrm{ES}
 \label{eq:gauge-orbital-constantes}
\end{equation}
sélectionne un unique représentant de chaque orbite et produit, respectivement,
\[
 \boxed{010211,\qquad201101,\qquad121200.}
\]
\end{theorem}

\begin{proof}
Les relations de groupe se vérifient directement comme des identités de
permutations de six coordonnées. L'énumération des \(468\) lignes
vérifie que \(r\) et \(s\) envoient chaque \(U_6\) sur une autre ligne du catalogue de
même multiplicité, et que la projection ternaire commute avec les deux
actions. Par conséquent, les orbites peuvent se calculer sur les \(243\) mots
sans perdre les fibres.

La partition exhaustive donne l'histogramme des tailles
\[
 38\cdot6+5\cdot3=243.
\]
En filtrant ces \(43\) orbites par le premier prédicat, il en reste exactement une ;
en filtrant par fibre à un point, multiplicité minimale et support diagonal
\(\{0,3,6\}\), il en reste six. Le recensement \((2,3,1)\) ne laisse qu'une orbite dans le
premier cas et le recensement équilibré \((2,2,2)\) n'en laisse qu'une dans le second.
Enfin, la recherche d'un relèvement avec \(d=6\) et orientation
\(\mathrm{ES}\) laisse un mot dans chacune des trois orbites. Toutes les
filtrations sont réalisées sur des descripteurs structurels du catalogue avant de consulter le
lecteur décimal.
\end{proof}

Les prédicats du théorème peuvent être réunis en un invariant de sélection.
Pour une orbite \(\mathcal O\), nous définissons son \emph{profil orbital}
\[
 \Sigma_{\rm cat}(\mathcal O)=
 \bigl(
 |\mathcal O|,\ f,\ m,\ \mathcal M_{\rm fib},\
 \mathcal D,\ \nu
 \bigr),
\]
où \(f\) est la taille de la fibre, \(m\) la masse des germes,
\(\mathcal M_{\rm fib}\) le multiensemble des multiplicités de ses
relèvements, \(\mathcal D\) le support diagonal lorsqu'il est employé et \(\nu\)
le recensement des trits. Les trois classes distinguées occupent des positions différentes
dans cet espace d'invariants :
\[
\begin{array}{c|c|c|c|c|c}
\mathcal O&|\mathcal O|&f&m&\mathcal M_{\rm fib}&(\mathcal D,\nu)\\ \hline
\mathcal O_\pi&6&5&1008&\{432,144,144,144,144\}&(-,(2,3,1))\\
\mathcal O_e&6&1&144&\{144\}&(\{0,3,6\},(2,3,1))\\
\mathcal O_\varphi&6&1&144&\{144\}&(\{0,3,6\},(2,2,2)).
\end{array}
\]
La singularité démontrée ne signifie pas que ces nombres reçoivent une
probabilité absolue sur \(\mathbb R\). Elle signifie qu'au sein de l'univers
fini produit par l'émetteur et avant l'évaluation des chiffres, leurs classes sont
séparables par des invariants internes et que la classe de \(\pi\) possède une fibre
régionale d'un type distinct. C'est la formulation mathématique d'une
priorité structurelle dans HMT\@.

Le contrôle négatif est informatif. Si l'on omet le support diagonal
\(\{0,3,6\}\), il existe quatre orbites à fibre unique équilibrées et
\(\mathcal O_\varphi\) cesse d'être unique. Avant de fixer la jauge, il y a six
orientations possibles dans chaque orbite. Après sa fixation, le mot de
\(\pi\) conserve trois relèvements positifs de \(U_6\), et non une seule région.
Ainsi, le théorème sélectionne le mot orienté mais respecte la multiplicité
géométrique des cinq régions.

La non-uniformité se renforce lors du passage aux orbites :
\[
\begin{array}{c|cc}
&\text{émissions}&\text{germes}\\ \hline
\mathcal O_\pi&5/78&14/243\\
\mathcal O_e&1/78&2/243\\
\mathcal O_\varphi&1/78&2/243.
\end{array}
\]
Ces quotients restent des masses du catalogue fini, et non des probabilités
attribuées aux nombres réels évalués.

L'avancée causale est précise. Les trois mots cessent d'être des entrées
nommées : ils sont retrouvés à partir de la symétrie et des invariants finis
d'APP--TPK, relativement à l'action \(D_3^{\rm cat}\), à l'étiquetage TRIT,
à la répartition \(3+3\) et à la jauge d'orientation de
\cref{eq:gauge-orbital-constantes}. La preuve n'utilise ni \(L_0\), ni \(L_1\),
ni chiffres décimaux, ni \(\alpha\), ni CODATA, ni l'état dodécaphasique. Restent des
affirmations ultérieures l'unicité absolue de cette action parmi tous les
automorphismes possibles, la déduction de la jauge au lieu de sa déclaration, la
canonicité du lecteur \(L_0/L_1\) indépendante de la valeur cible et l'émission coinductive des
développements infinis.

\section{Position de \texorpdfstring{\(\pi,e,\varphi\) et \(\alpha\)}{pi, e, phi et alpha} dans la classification}

Le \cref{thm:shift-hensel-completo} explique pourquoi tout préfixe
décimal peut être réémis à partir d'un état 3-adique : la capacité brute est
universelle. Par conséquent, la simple opération d'aller-retour sur un ruban ne sélectionne pas une
constante. Le contenu distinctif de la construction publiée se trouve
dans les objets antérieurs et postérieurs à cette réémission :
\begin{enumerate}
\item les fibres finies et les cinq régions distinctes de \(\pi\) ;
\item le lecteur archimédien typé qui associe mots, cylindres et blocs ;
\item la dodécaphase signée et la clôture réversible
\[
U_{12}^{\rm sgn}
\longleftrightarrow K
\longleftrightarrow(D_3K,D_4K,Q)
\longleftrightarrow\alpha_{\rm HMT};
\]
\item les critères de comparaison qui séparent une équivalence de cartes d'un
sélecteur antérieur à l'évaluation.
\end{enumerate}

Les douze triplets décimaux publiés de \(\pi,e,\varphi\) sont des coordonnées
finies exactes du lecteur enrichi. La chaîne imprimée de trente-six
chiffres est la projection dodécaphasique rationnelle
\(\alpha_{12}=\pi_{12}(\alpha_{\rm HMT})\) :
\[
\frac{7297352569283800997285105472380663}{10^{36}}.
\]
L'identité exacte dans cette fenêtre est
\[
\alpha_{12}
:=\pi_{12}(\alpha_{\rm HMT})
=\frac{7297352569283800997285105472380663}{10^{36}}.
\]
La récurrence compatible à profondeur arbitraire construit la section
infinie
\[
 \alpha_{\rm HMT}:=\varprojlim_N A^{(N)},
\]
et chaque fenêtre rationnelle est une projection exacte de ce même objet. Le
noyau \(E_{21/22}\) construit, dans une autre carte, la racine positive simple
finie \(\alpha_E\) ; le transport gradué du TPK prolonge ce noyau en
une famille compatible dont la racine limite est \(\alpha_B\). Si
\(\alpha_A:=\varprojlim_N\rho_N(\alpha_{12})\), les carrés de troncature
des deux familles commutent et déterminent à chaque profondeur le même
cylindre survivant. Le théorème exact des deux voies établit donc
\[
 \boxed{\alpha_A=\alpha_B=\alpha_{\rm HMT}}.
\]
L'égalité de \(\operatorname{Tr}_9(\alpha_E^{-1})\) et de
\(\operatorname{Tr}_9(\alpha_{12}^{-1})\) est le premier certificat fini de
cette identité ; elle ne la définit pas, ne limite pas sa portée et ne sélectionne pas les
prolongements. La projection typée conserve la distinction entre la section
coinductive, sa fenêtre rationnelle et la racine analytique sans séparer leur identité
à la limite.

La racine carrée de \(-1\) appartient à une autre carte. Dans le
\cref{chap:algebras-cuadraticas-orden-nueve}, on construit
\[
\mathbb F_9=\mathbb F_3[\iota]/(\iota^2+1),
\qquad \iota^2=-1,
\]
et l'on démontre que la matrice de Witt réalise cette structure quadratique. Il s'agit
d'une extension algébrique finie exacte, et non d'une catégorie visible de
développements réels.

\section{Carte de Hensel non filtrée, survie et frontière globale}

Les contrôles suivants sont des falsificateurs des résultats finis
précédents :
\begin{enumerate}
\item que deux états distincts modulo \(3^T\) émettent le même mot de
\(T\) blocs, ou qu'il existe un mot sans préimage ;
\item que l'inverse explicite de
\cref{eq:inversa-hensel-finita} échoue ;
\item qu'il apparaisse un rationnel dont le développement canonique ne soit pas ultimement
périodique, ou un irrationnel dont le développement le soit ;
\item qu'un émetteur déterministe fini produise un ruban non ultimement
périodique ;
\item que les recensements du catalogue ne totalisent pas \(468\), \(243\) et \(104\,976\),
ou que les mesures des trois mots ne coïncident pas avec l'énumération.
\end{enumerate}

La carte de Hensel brute
\[
\mathbb Z_3^6\longrightarrow[0,1]
\]
est continue et surjective ; en ajoutant la partie entière, son atlas brut couvre
\(\mathbb R\). Ce résultat exact n'implique pas que toute fibre brute contienne
un représentant qui survive aux filtres supplémentaires. Les fenêtres
finies publiées et les lecteurs nommés de \(\pi,e,\varphi\) sont vérifiés
dans HMT, relativement à leurs germes, régions,
orientations et règles enrichies.

\ifdefined\HMTInternalTraceability
Pour un préfixe filtré, soit \(T_u\) l'arbre des prolongements compatibles.
Si \(T_u\) est à ramification finie et possède un nœud à toute profondeur, le
lemme de König produit une branche infinie ; si, en outre, chaque niveau stable est
unitaire, la branche est unique. Cet énoncé est conditionnel : la non-vacuité universelle de \(T_u\)
pour tout préfixe n'a pas été démontrée. Restent donc
comme \textsc{programme ouvert} la couverture de toutes les fibres par des
survivants filtrés et la descente bien définie de la somme et du produit
natifs au quotient archimédien.
\fi




\label{chap:taxonomia-monodromia}

La représentation positionnelle d'un nombre enregistre une suite de symboles,
mais ne conserve pas nécessairement l'état qui détermine sa continuation.
L'holographie modulaire triadique distingue les deux niveaux. La phase nonadique indique
la position d'une lecture dans le cycle de neuf pas ; l'état
enrichi conserve, en outre, l'orientation, le quotient euclidien, le transport
d'entraînement et la condition de frontière. Le retour de la phase n'implique donc pas
le retour de l'état.

Soit \(S\) un espace d'états enrichis, soit
\(T:S\to S\) une transition et soit
\[
g:S\longrightarrow\mathbb Z/9\mathbb Z
\]
la phase, avec \(g(Ts)=g(s)+1\). Pour une base \(b\geq2\), un lecteur
\(d:S\to\{0,\ldots,b-1\}\) produit
\[
x(s)=\sum_{n\geq0}\frac{d(T^ns)}{b^{n+1}}.
\]
Lorsqu'une valeur admet les deux développements positionnels usuels, on adopte la
représentation canonique qui ne se termine pas par une queue constante \(b-1\). Cette
convention sépare une ambiguïté d'écriture de la multiplicité proprement
HMT, qui procède de sections enrichies distinctes ayant une même
évaluation.
L'opérateur de retour nonadique est la restriction
\[
M=T^9:g^{-1}(r)\longrightarrow g^{-1}(r).
\]
La différence entre \(T^9s=s\) et \(g(T^9s)=g(s)\) formalise la différence
entre périodicité et monodromie.

La transition de l'état enrichi de TPK peut être donnée par une relation et non par une
fonction. Dans ce cas, la notation précédente s'applique au décalage sur
l'espace des histoires infinies ou à une réalisation déterministe dont l'état
a été étendu avec l'information nécessaire pour fixer la continuation ; on ne
présuppose pas qu'un bloc visible possède un successeur unique.

\begin{definition}[Nombre HMT et présentation arithmétique]
Un nombre HMT est uniquement une section compatible
\(\boldsymbol x\) du système inverse d'états survivants. Un
lecteur positionnel \(L\) dont le domaine contient \(\boldsymbol x\) étant choisi, le couple
\((\boldsymbol x,L)\) est une présentation ou une mesure du nombre. La valeur
archimédienne publiée par cette présentation est l'intersection ponctuelle des
cylindres associés lorsque leurs diamètres tendent vers zéro. Deux sections ayant la
même valeur ne sont pas identifiées avant l'application du quotient archimédien.
\end{definition}

\begin{definition}[Périodicités visible et enrichie]
La sortie est éventuellement périodique lorsqu'il existe \(N\geq0\) et \(p>0\)
tels que
\[
d(T^{n+p}s)=d(T^ns),\qquad n\geq N.
\]
L'état est éventuellement périodique lorsque
\[
T^{N+p}s=T^Ns.
\]
La seconde condition implique la première. La réciproque exige que le lecteur
sépare les queues de l'orbite considérée, c'est-à-dire que
\[
 \bigl(d(T^{n+k}s)\bigr)_{k\geq0}
 =\bigl(d(T^{m+k}s)\bigr)_{k\geq0}
 \quad\Longrightarrow\quad T^ns=T^ms.
\]
\end{definition}

\begin{theorem}[Classification par retour de phase]
\label{thm:clasificacion-retorno-fase}
Soit \(s\in S\).
\begin{enumerate}
  \item Si l'orbite de \(s\) atteint un état dont la sortie future est
  constamment nulle, alors \(x(s)\) possède un développement fini et est
  rationnel.
  \item Si l'état est éventuellement périodique, alors la sortie est
  éventuellement périodique et \(x(s)\in\mathbb Q\).
  \item Si le lecteur sépare les queues de l'orbite et que l'orbite de \(s\) n'est pas
  éventuellement périodique, alors la sortie n'est pas éventuellement périodique
  et \(x(s)\notin\mathbb Q\).
  \item La non-périodicité ne distingue pas l’irrationalité algébrique de la
transcendance. Cette distinction exige un certificat polynomial ou un
théorème de transcendance supplémentaire.
\end{enumerate}
\end{theorem}

\begin{proof}
Les deux premiers points sont la caractérisation classique des rationnels
par des développements positionnels finis ou ultimement périodiques. Dans
le troisième, une périodicité ultime de la sortie produirait deux queues
égales. L’hypothèse de séparation identifierait alors les états
correspondants, ce qui contredirait la non-périodicité de l’orbite. Le quatrième
point découle de l’existence d’irrationnels algébriques et transcendants
dont les développements ne sont pas ultimement périodiques.
\end{proof}

\begin{corollary}[Monodromie nonadique et apériodicité]
Supposons que la phase revienne tous les neuf pas, que l’orbite de retour
\[
 \{M^ks:k\geq0\},\qquad M=T^9,
\]
ne soit pas ultimement périodique et que le lecteur sépare les queues. Alors la
sortie n’est pas ultimement périodique et le nombre produit est irrationnel. La
périodicité de la phase coexiste dans ce cas avec une évolution apériodique de
l’état enrichi.
\end{corollary}

L’hypothèse ne peut être remplacée par une seule inégalité \(M(s)\ne s\), ni
par un nombre fini de retours distincts. Relativement à l’état initial,
à la transition et au lecteur déclarés, la fenêtre actuelle certifie neuf
retours de phase avec changement d’état ; elle démontre une monodromie finie, mais non
à elle seule l’apériodicité de toute l’orbite.

\ifdefined\HMTRepetirResultadosTaxonomicos
\section{Le carré local de \texorpdfstring{\(e\)}{e} et la rupture de mémoire}

Pour un mot décimal \(a_1\cdots a_n\), nous appellerons \emph{carré local}
de longueur \(\ell\) à la position \(i\) une factorisation
\[
 a_i\cdots a_{i+2\ell-1}=vv,
 \qquad |v|=\ell.
\]
Cette notion appartient à la combinatoire des mots et ne doit pas être confondue avec
une période de la queue infinie.

\begin{theorem}[Caractérisation finie du carré \(1828^2\)]
Considérons les \(243\) mots du catalogue N33 et le lecteur figé
\(w_6\mapsto w_{30}\mapsto d_1d_2d_3d_4\), avec quatre triades en base mille.
Il existe un unique mot dont la lecture contient un carré de longueur quatre
commençant au deuxième chiffre décimal. C’est
\[
 w_e=201101,
 \qquad
 718281828459=7\,\underbrace{1828\,1828}_{(1828)^2}\,459.
\]
La répétition ne se prolonge pas en une troisième copie : le chiffre suivant requis
serait \(1\), tandis que le chiffre effectif est \(4\). Dans tout le catalogue,
il n’existe qu’un autre carré de longueur quatre,
\[
 575157518472=\underbrace{5751\,5751}_{\text{carré de mot}}\,8472,
\]
mais il commence au premier chiffre.
\end{theorem}

\begin{proof}
L’énumération exhaustive évalue les \(243\) mots par arithmétique
entière exacte. Le profil des carrés pour les longueurs \(1,\ldots,6\) contient,
respectivement,
\[
240, 21, 3, 2, 0, 0
\]
occurrences, réparties entre
\[
153, 19, 3, 2, 0, 0
\]
mots. Les deux occurrences de longueur quatre sont exactement celles écrites
dans l’énoncé. Deux énumérations entières indépendantes reproduisent le même
recensement et leurs sorties complètes coïncident terme à terme.
\end{proof}

Le mot \(w_e\) possède en outre la chaîne de relèvement
\[
201101\mid121221\mid102011\mid012222\mid102011.
\]
Le troisième et le cinquième blocs visibles coïncident. En revanche,
les états qui les transportent ne coïncident pas. Avant les deux émissions, leurs réductions
modulo \(27\) sont
\[
(25,11,7,17,8,16),
\qquad
(7,8,4,23,26,7),
\]
et les quotients de Hensel ultérieurs sont
\[
(6,13,9,2,13,1),
\qquad
(15,0,3,19,23,25).
\]
Par conséquent, le retour \(102011\) est un retour de la projection visible, non
de l’état enrichi. C’est la formulation exacte de l’intuition d’une
répétition locale qui se rompt : le quotient mémorise une généalogie que le
bloc résiduel ne conserve pas.

Le résultat appartient au lecteur enrichi déjà fixé par le système
APP--TRIT--TPK\@. Le recensement fini caractérise exhaustivement l’apparition du
carré local ; le prolongement coinductif de la section appartient à la
connexion nonadique construite dans le
\cref{chap:numero-heisenberg-d108}. La propriété combinatoire du préfixe et
la génération du développement complet sont donc deux niveaux
compatibles d’une même généalogie.

Cette formulation fournit le cadre exact pour interpréter le début
décimal de \(e\),
\[
e=2.718281828459\ldots.
\]
L’expression correcte n’est plus « semi-périodicité », mais \emph{carré local
de mot avec retour visible et rupture de mémoire}.

\section{2 mesures finies de l’émetteur \(468\to243\) : fréquence et priorité structurelle}

Il n’existe pas de distribution uniforme sur tous les nombres réels. Il
existe en revanche une mesure canonique sur le catalogue fini lorsque l’on munit d’une loi uniforme
l’ensemble des \(104\,976\) germes TPK\@. Pour un mot visible \(w\), on
définit
\[
\mu_{\mathrm{TPK}}(w)
=\frac{\#\{\text{germes se projetant sur }w\}}{104\,976}.
\]
La mesure n’est pas uniforme. Les \(243\) mots possèdent des poids compris
entre \(144\) et \(2088\), avec neuf valeurs de multiplicité distinctes.
L’histogramme exact, où le poids est écrit comme base et le nombre de mots
comme exposant, est
\[
144^{120},\ 288^{54},\ 432^{18},\ 576^9,\ 864^{15},\
1008^6,\ 1440^3,\ 1944^{12},\ 2088^6.
\]

Pour les trois mots distingués, on obtient
\[
\mu_{\mathrm{TPK}}(010211)=\frac{1008}{104\,976}=\frac7{729},
\]
\[
\mu_{\mathrm{TPK}}(201101)
=\mu_{\mathrm{TPK}}(121200)
=\frac{144}{104\,976}=\frac1{729}.
\]
Le mot associé à \(\pi\) a cinq antécédents \(U_6\) et un poids sept
fois supérieur à celui des mots associés à \(e\) et \(\varphi\). Ce résultat
démontre une hiérarchie combinatoire dans le catalogue ; il n’identifie pas à lui
seul les mots aux constantes classiques et n’établit pas de probabilité
sur \(\mathbb R\).

Il ne définit pas non plus, à lui seul, une priorité totale : six mots ont un poids
\(1008\), tandis que \(e\) et \(\varphi\) appartiennent à la classe minimale de
poids \(144\), partagée par \(118\) autres mots. La distinction HMT doit
s’exprimer par un profil structurel, par exemple
\[
\mathcal P(w)=
\bigl(m_{\rm semilla}(w),m_{U_6}(w),r_{\rm supervivencia}(w),
s_{\rm frontera}(w),s_{\rm excepcional}(w)\bigr),
\]
où chaque coordonnée est définie avant l’application d’étiquettes externes. Un scalaire
de « priorité » exige en outre de fixer, également à l’avance, la manière d’ordonner ou de
pondérer ces coordonnées.
\fi

\section{Classification des constantes selon l’antécédent généalogique, l’opération génératrice et le codomaine}

La classification HMT sépare des propriétés que la représentation ordinaire
réunit sous une seule valeur :
\begin{center}
\begin{tabular}{p{0.25\textwidth}p{0.64\textwidth}}
\toprule
Classe & Condition structurelle \\
\midrule
Fini & La sortie atteint une queue nulle. \\
Rationnel périodique & La sortie est ultimement périodique. \\
Rationnel monodromique & La sortie est périodique, mais l’état enrichi
parcourt une fibre non triviale que le lecteur visible oublie. \\
Irracional coinductivo & La sortie n’est pas ultimement périodique et les
cylindres déterminent une unique valeur. \\
Recurrente local & Un mot fini contient des retours ou des carrés, sans que
cela impose une périodicité à la queue. \\
Algébrique & La valeur satisfait un polynôme non nul à coefficients entiers ;
la route doit être accompagnée de ce certificat. \\
Transcendant & La valeur ne satisfait aucun polynôme entier non nul ; la
classification exige un théorème de transcendance. \\
Distinguido HMT & La section satisfait un prédicat structurel figé de
fibre, de frontière, de survie ou de symétrie ; c’est une classe transversale aux
précédentes. \\
\bottomrule
\end{tabular}
\end{center}

Ainsi, HMT ne modifie pas les définitions arithmétiques correctes de la rationalité,
de l’algébricité et de la transcendance. Elle les affine en ajoutant la généalogie de l’état,
la monodromie de phase et la multiplicité des fibres que l’évaluation
archimédienne avait oubliées.

En particulier, \(\varphi\) est algébrique parce qu’il satisfait
\(X^2-X-1=0\), tandis que la transcendance de \(e\) et de \(\pi\) découle de
théorèmes classiques externes. Ces propriétés sont compatibles avec le fait que les trois
sections soient distinguées par HMT\@. La clôture dodécaphasique construit plus
loin la projection rationnelle exacte
\(\alpha_{12}=\pi_{12}(\alpha_{\mathrm{HMT}})\), et le prolongement compatible
construit la section coinductive \(\alpha_{\mathrm{HMT}}\), dans le
\cref{chap:alpha-rev7} ; aucune d’elles n’est utilisée dans cette taxonomie comme entrée
ni publiée ici comme une seconde résidence. Le noyau \(E_{21/22}\)
produit une racine analytique distincte, dont la concordance avec la voie entière est
démontrée dans le codomaine de \(\operatorname{Tr}_9\). Section, projection
rationnelle, racine analytique, lecteur et valeur conservent ainsi leurs types propres.

Enfin, une « direction infinitésimale » exige un typage propre. Pour un
lecteur fini \(r_m\), deux sections distinctes dans une même fibre de \(r_m\)
définissent une différence invisible à la profondeur \(m\), quoique active dans un
prolongement ultérieur. Elles ne constituent pas pour autant un réel infinitésimal non nul.
Si deux sections coïncident sous toutes les troncatures fidèles du système
inverse, elles sont la même section. Un infinitésimal arithmétique véritable exigerait de
construire une extension ordonnée ou valuée non archimédienne et de démontrer que les
opérations HMT y descendent.



\section[Nombre comme section compatible]{Nombre comme section
compatible : classification arithmétique, évaluation archimédienne et
généalogie}
\label{p3-83-c44-s1:sec:ontologia-numeros-rev11}
\index{nombre!valeur et généalogie}
\index{taxonomie numérique}

La classification numérique de HMT conserve les espèces arithmétiques
ordinaires et ajoute une coordonnée que l’évaluation positionnelle écarte : la
généalogie de la section qui publie la valeur. Le résultat est une
classification conjointe. Le premier axe répond à la question du nombre obtenu ; le
second enregistre par quelle suite compatible d’états, de retours,
de retenues et de clôtures il est obtenu. Cette distinction permet de formuler dans un même
langage les entiers, les développements rationnels, les irrationnels, l’unité
imaginaire, les régularisations et les constantes spectrales sans identifier
des objets qui ne partagent qu’un scalaire.

\subsection{L’évaluation archimédienne sépare la section de sa valeur}

Soit
\[
 X_{\infty}^{\rm surv}
 =\varprojlim_n X_n^{\rm surv}
\]
un système inverse d’états survivants avec ses applications de
restriction. Une section compatible est une famille
\(\boldsymbol x=(x_n)_{n\geq0}\) qui satisfait
\(\rho_{n+1,n}(x_{n+1})=x_n\) à toute profondeur. Un lecteur positionnel
\(L\) associe à chaque préfixe un cylindre \(I_n(\boldsymbol x;L)\). Lorsque les
cylindres sont compatibles, que leurs diamètres convergent vers zéro et que la
convention canonique des extrémités est fixée, on définit
\begin{equation}
 \operatorname{Val}_{L}(\boldsymbol x)
 :=\text{le point de }
 \bigcap_{n\geq0}\overline{I_n(\boldsymbol x;L)}.
 \label{p3-83-c44-s1:eq:valor-lector-seccion-rev11}
\end{equation}

\HMTClaim{HMT-R-NUMBER-SECTION-VALUE}
\begin{definition}[Nombre, présentation et valeur dans HMT]
\label{p3-83-c44-s1:def:numero-presentacion-valor-rev11}
Un \emph{nombre HMT} est une section compatible
\(\boldsymbol x\in X_{\infty}^{\rm surv}\). La paire
\((\boldsymbol x,L)\) est une \emph{présentation} du nombre et
\(\operatorname{Val}_{L}(\boldsymbol x)\) est la valeur publiée par cette
présentation. La fibre
\[
 \operatorname{Val}_{L}^{-1}(r)
\]
conserve les généalogies que le quotient archimédien représente par
la même valeur \(r\).
\end{definition}

La capacité de la carte et la sélection d’une section sont des résultats
typés distincts. Le \cref{thm:cobertura-arquimediana-bruta} construit une
application continue et surjective
\[
 P_{\rm raw}:\mathbb Z_3^6\longrightarrow[0,1],
\]
et le \cref{cor:atlas-hmt-reales} la globalise en
\[
 \widehat P_{\rm raw}:\mathbb Z\times\mathbb Z_3^6
 \longrightarrow\mathbb R.
\]
Ces applications prouvent la capacité représentationnelle complète de la carte brute.
Le prédicat de survie et les lecteurs qui distinguent
\(\pi,e,\varphi\) sélectionnent, au sein de cette capacité, les généalogies
structurelles correspondantes.

\subsection{Terminaison et périodicité déterminent la rationalité}

Pour le développement canonique en base mille
\[
 x=\sum_{n\geq1}\frac{b_n}{1000^n},
 \qquad b_n\in\{0,\ldots,999\},
\]
on adopte la convention qui élimine la duplication produite par une queue
à partir d’un certain rang constante de blocs \(999\). Le
\cref{thm:taxonomia-base-mil} fournit alors la classification exacte
suivante.

\HMTClaim{HMT-R-NUMBER-DECIMAL-CLASSIFICATION}
\begin{theorem}[Classification visible dans la carte décimale cubique]
\label{p3-83-c44-s1:thm:clasificacion-visible-rev11}
Soit \(x=a/q\in[0,1)\) une fraction irréductible lorsque \(x\) est rationnel.
Alors :
\begin{enumerate}
\item le développement se termine si et seulement si \(q\mid1000^N\) pour un certain
      \(N\), équivalence qui restreint les facteurs premiers de \(q\) à
      \(2\) et \(5\) ;
\item le développement est périodique à partir d’un certain rang si et seulement si
      \(x\in\mathbb Q\) ;
\item le développement n’est pas périodique à partir d’un certain rang si et seulement si
      \(x\in\mathbb R\setminus\mathbb Q\).
\end{enumerate}
L’algébricité et la transcendance constituent un second axe : un réel est
algébrique s’il annule un polynôme non nul de \(\mathbb Q[X]\), et il est
transcendant s’il n’en annule aucun.
\end{theorem}

\begin{proof}
Un développement de longueur \(N\) a pour dénominateur un diviseur de \(1000^N\), et
la réciproque s’obtient en exprimant \(a/q\) avec ce dénominateur. Pour un
rationnel, la division successive possède un ensemble fini de restes : un reste
nul produit la terminaison et la répétition d’un reste produit la périodicité.
Une queue périodique est, réciproquement, une série géométrique rationnelle. La
classification polynomiale est indépendante de la base positionnelle.
\end{proof}

La dynamique enrichie distingue en outre la périodicité de sortie et la
périodicité d’état. Si \(T\) transporte l’état et \(d\) publie un
bloc, les conditions
\[
 d(T^{n+p}s)=d(T^n s),
 \qquad
 T^{n+p}s=T^n s
\]
définissent, respectivement, le retour visible et le retour de l'état. Le second
implique le premier. L'implication inverse exige que le lecteur sépare les
queues de l'orbite, conformément au
\cref{thm:clasificacion-retorno-fase}. C'est pourquoi une phase nonadique périodique
peut transporter une mémoire non périodique.

\begin{longtable}{@{}>{\raggedright\arraybackslash}p{.20\textwidth}
                    >{\raggedright\arraybackslash}p{.30\textwidth}
                    >{\raggedright\arraybackslash}p{.42\textwidth}@{}}
\caption{Taxonomie conjointe de la valeur publiée et de l'état enrichi.}
\label{p3-83-c44-s1:tab:taxonomia-conjunta-numeros-rev11}\\
\toprule
\textbf{Classe} & \textbf{Condition de la valeur} &
\textbf{Condition généalogique ou démonstration requise}\\
\midrule
\endfirsthead
\toprule
\textbf{Classe} & \textbf{Condition de la valeur} &
\textbf{Condition généalogique ou démonstration requise}\\
\midrule
\endhead
Entier & Élément de \(\mathbb Z\). & Mot ternaire équilibré fini et
unique ; le signe et la grandeur restent typés.\\
Decimal finito & Rationnel dont le dénominateur réduit divise une puissance
de \(1000\). & La sortie atteint une queue nulle.\\
Rationnel périodique & Développement ultimement périodique. & Un état
éventuellement périodique suffit ; une sortie périodique peut conserver une
fibre interne supplémentaire.\\
Rationnel monodromique & Même valeur rationnelle que la ligne précédente. & La
sortie revient tandis que l'état enrichi parcourt une fibre que le lecteur
visible oublie.\\
Irracional coinductivo & Développement non ultimement périodique. & Une famille
infinie de cylindres compatibles fixe une valeur unique ; l'apériodicité
globale exige une preuve sur toute l'orbite, et non une fenêtre finie.\\
Algébrique & Racine d'un polynôme non nul de \(\mathbb Q[X]\). & On joint
le polynôme et l'on vérifie son annulation.\\
Transcendant & Réel extérieur à la clôture algébrique de \(\mathbb Q\). & On
joint un théorème de transcendance ; la non-périodicité décimale à elle seule
ne démontre que l'irrationalité.\\
Distinguido HMT & L'une quelconque des classes arithmétiques précédentes. & La
section satisfait un prédicat figé de fibre, de frontière, de survie,
de symétrie ou de clôture.\\
\bottomrule
\end{longtable}

Ce tableau place sur des axes différents deux questions auxquelles il faut répondre
séparément. La classification ordinaire détermine les propriétés de la valeur ;
le lecteur HMT détermine la chaîne structurelle qui engendre la section. Par
exemple,
\[
 \varphi^2-\varphi-1=0
\]
classe \(\varphi\) comme irrationnel algébrique, tandis que la
transcendance de \(e\) et de \(\pi\) relève de théorèmes classiques externes.
Les trois peuvent, simultanément, être des sections distinguées par des prédicats
HMT différents. Leurs troncatures imprimées de mille ou dix mille chiffres sont
des rationnels finis qui démontrent les préfixes calculés d'une sortie prolongeable ; elles ne
remplacent pas l'objet infini qu'elles représentent.

La clôture dodécaphasique définit, quant à elle, la projection rationnelle finie
\begin{equation}
 \alpha_{12}=\pi_{12}(\alpha_{\rm HMT})
 =\frac{7297352569283800997285105472380663}{10^{36}}.
 \label{p3-83-c44-s1:eq:alpha-racional-taxonomia-rev11}
\end{equation}
Cette projection est exacte et rationnelle. L'objet
\(\alpha_{\rm HMT}=\varprojlim_NA^{(N)}\) est la section infinie construite
par la récurrence compatible ; ses troncatures rationnelles ne la remplacent pas.
La dérivation structurelle de la limite et la classification de chaque
projection finie sont des affirmations complémentaires.

\subsection{La carte ternaire équilibrée représente les entiers
bijectivement}

Soit
\[
 \mathscr W_{\rm bal}
 =\{\varnothing\}\cup
 \{a_1\cdots a_m:
 a_j\in\{-1,0,+1\},\ a_1\neq0\}
\]
et définissons
\[
 b(\varnothing)=0,
 \qquad
 b(wa)=3b(w)+a.
\]
Le \cref{pf84:C03} démontre que
\begin{equation}
 b:\mathscr W_{\rm bal}\xrightarrow{\ \sim\ }\mathbb Z
 \label{p3-83-c44-s1:eq:carta-entera-balanceada-rev11}
\end{equation}
est une bijection. Le dernier trit est le représentant unique de la
classe modulo trois et le quotient \((n-a)/3\) réduit strictement la valeur
absolue jusqu'à atteindre zéro. Le premier trit non nul fixe le signe. Cette carte
est une représentation arithmétique exacte ; l'appartenance d'un mot à une
section survivante incorpore en outre les conditions de route,
d'orientation et de frontière.

Pour \(n>0\), les coordonnées
\[
 \rho_9(n)=1+((n-1)\bmod9),
 \qquad
 \kappa_9(n)=\frac{n-\rho_9(n)}9
\]
réalisent la décomposition \(n=9\kappa_9(n)+\rho_9(n)\). La phase
\(\rho_9\in\{1,\ldots,9\}\) publie la classe résiduelle zéro comme \(9\), et
\(\kappa_9\) conserve le nombre de tours. Avec le signe
\(\varepsilon\in\{-1,+1\}\),
\[
 n\longleftrightarrow
 \bigl(\varepsilon,\kappa_9(|n|),\rho_9(|n|)\bigr)
\]
est bijectif sur \(\mathbb Z\setminus\{0\}\) ; zéro reste le point
de base. Le retour visible \(9\to1\) est ainsi une transition de phase avec
incrément de mémoire, et non une identification d'entiers consécutifs.

\subsection{Structure complexe finie induite par l'endomorphisme de Witt \texorpdfstring{\(J^2=-I_6\)}{J²=-I₆}}

L'unité imaginaire ordinaire est l'élément
\(i\in\mathbb C\) qui satisfait \(i^2=-1\) ; il est algébrique de degré deux sur
\(\mathbb R\) et sur \(\mathbb Q\). La réalisation HMT appartient à un
domaine fini différent. Si \(C_P\) est la matrice de conférence de Paley et
\(J=A_W\) sa réduction modulo trois dans
\(\operatorname{End}_{\mathbb F_3}(\mathbb F_3^6)\), alors
\begin{equation}
 \boxed{J^2=2I_6=-I_6.}
 \label{p3-83-c44-s1:eq:raiz-interna-menos-uno-ontologia-rev11}
\end{equation}
Comme \(X^2+1\) est irréductible dans \(\mathbb F_3[X]\), l'extension
\[
 \mathbb F_9=\mathbb F_3[\iota]/(\iota^2+1)
\]
sépare les espaces propres de \(J\) associés à \(\pm\iota\). Le résultat,
démontré dans \cref{sec:raiz-interna-menos-uno-rev6}, dote le module ternaire d'une
structure complexe finie et organise deux orientations conjuguées. Sa force
est \emph{exacte interne, démontrée avec la matrice de Paley--Witt déclarée} ; la
comparaison avec \(\mathbb C\) conserve la distinction entre les deux corps.

\subsection{\(6\) opérateurs réalisent le rationnel
\texorpdfstring{\( -1/12 \)}{-1/12} dans des domaines distincts}

La régularisation zêta attribue à la suite des entiers positifs la partie
finie obtenue par prolongement méromorphe :
\begin{equation}
 \sum_{n\geq1}^{[\zeta]}n
 :=\zeta(-1)=-\frac1{12}.
 \label{p3-83-c44-s1:eq:regularizacion-zeta-menos-doce-rev11}
\end{equation}
L'exposant \([\zeta]\) fait partie du type de l'expression. La somme
partielle ordinaire satisfait
\(\sum_{n=1}^{N}n=N(N+1)/2\) et croît sans borne ; c'est pourquoi
\eqref{p3-83-c44-s1:eq:regularizacion-zeta-menos-doce-rev11} est une valeur régularisée, et non
la limite de ces sommes partielles.

HMT contient en outre six réalisations algébriques exactes du même
rationnel. La coïncidence scalaire permet de les comparer ; chaque objet conserve son
domaine, son application et sa normalisation.

\begin{longtable}{@{}>{\raggedright\arraybackslash}p{.23\textwidth}
                    >{\raggedright\arraybackslash}p{.43\textwidth}
                    >{\raggedright\arraybackslash}p{.26\textwidth}@{}}
\caption{Réalisations typées de \(-1/12\).}
\label{p3-83-c44-s1:tab:realizaciones-menos-doce-rev11}\\
\toprule
\textbf{Domaine} & \textbf{Application et évaluation} &
\textbf{Conclusion et normalisation}\\
\midrule
\endfirsthead
\toprule
\textbf{Domaine} & \textbf{Application et évaluation} &
\textbf{Conclusion et normalisation}\\
\midrule
\endhead
Incidence marquée hexade--octade &
\(F_6\hookrightarrow F_8\), avec
\(\lvert F_6\rvert=90\), \(\lvert F_8\rvert=120\), et
\[
 (90-120)/(24\binom62)=-1/12.
\]
& Défaut combinatoire exact pour la normalisation déclarée.\\
Périodicités dodécaphasiques &
Un pas de \(30^\circ\) produit
\[
 \begin{aligned}
 \frac{30}{120}-\frac{30}{90}
 &=\frac14-\frac13,\\
 &=-\frac1{12}.
 \end{aligned}
\]
& Différence orientée exacte de périodes.\\
Opérateur de Gram &
Pour \(K_{60,81}=\operatorname{diag}(0,12)\), la pseudo-inverse sur le
rang non nul vaut \(1/12\) ; l'orientation négative publie \(-1/12\). &
Identité opératorielle exacte sur la composante non nulle.\\
Deuxième différence d'échelle &
Le calcul produit \(C(L)=8/81\) et l'extracteur déclaré
\(R(L)=-(27/32)C(L)=-1/12\). & Exact après fixation du facteur de
normalisation \(-27/32\) ; ce facteur est une primitive explicite.\\
Quotient modulaire êta &
\[
 \begin{aligned}
 \frac{\eta(\tau)}{\eta(3\tau)}
 &=q^{-1/12}\\[-1mm]
 &\quad{}\times
 \prod_{n\geq1}(1+q^n+q^{2n})^{-1}.
 \end{aligned}
\]
& Exposant modulaire exact, provenant de \(1/24-3/24\).\\
Projecteurs du cycle de douze phases &
Pour \(D_r=I-S^r\),
\[
 \begin{aligned}
 \tau(\Pi_{\ker D_3}-\Pi_{\ker D_4})
 &=\frac{3-4}{12},\\
 &=-\frac1{12}.
 \end{aligned}
\]
& Défaut transversal exact entre les pas trois et quatre.\\
\bottomrule
\end{longtable}

Les six réalisations sont démontrées dans
\cref{sec:menos-doce-realizaciones-rev6}. La régularisation zêta, le défaut
d'incidence, la pseudo-inverse, l'extracteur d'échelle, le quotient êta et la
trace des projecteurs sont des applications différentes. Une identification de leurs
domaines nécessiterait un entrelaceur supplémentaire. De même, le terme
\(+1/12\) de la formule de Glaisher--Kinkelin relève de la régularisation
de l'hyperfactorielle et conserve ce type propre.

\subsection{Décomposition multiplicative des nombres premiers comme modes irréductibles du TPK}
\index{nombre premier!irréductible}
\index{horloge première}

Un nombre premier ordinaire est un entier \(p>1\) dont les factorisations
\(p=ab\), avec \(a,b\in\mathbb N\), contiennent nécessairement une unité. Le
coproduit réduit
\[
 \Delta_{\times}^{\rm red}e_n
 =\sum_{\substack{ab=n\\a,b>1}}e_a\otimes e_b
\]
convertit cette définition en un sélecteur spectral. Selon le
\cref{pf84:C07}, l'opérateur positif
\[
 N_{\times}^{\rm red}
 =(\overline{\Delta_{\times}^{\rm red}})^*
   \overline{\Delta_{\times}^{\rm red}}
\]
satisfait
\(N_{\times}^{\rm red}e_n=d_{\rm red}(n)e_n\), et son noyau, après
retrait de \(e_1\), sélectionne exactement les nombres premiers. L'entrelaceur
multiplicatif transporte le même sélecteur au produit natif d'APP\@. Ce
sélecteur est la projection spectrale sur
\(\ker N_{\times}^{\rm red}\ominus\mathbb C e_1\), qui coïncide exactement
avec le sous-espace engendré par les nombres premiers et ne dépend pas de la reconnaissance de listes
décimales.

La factorisation unique fournit en outre la réalisation de Fock
\[
 U_F:\Gamma_s(\ell^2(\mathbb P))
 \xrightarrow{\ \sim\ }\ell^2(\mathbb N_{\geq1}),
 \qquad
 U_F\,d\Gamma(H_{\mathbb P})\,U_F^*=\log Q,
\]
où \(H_{\mathbb P}e_p=(\log p)e_p\). Chaque nombre premier est un mode
multiplicativement irréductible ; un entier composé est une occupation finie
de ces modes.

Pour \(n\) premier avec \(10\), son horloge décimale est l'ordre multiplicatif
\[
 \tau_n=\operatorname{ord}_n(10).
\]
Si \(p\) est premier et \(p\notin\{2,3,5\}\), le
\cref{thm:cierre-nonadico-reloj-primo-rev6} démontre
\begin{equation}
 \tau_p\mid p-1,
 \qquad
 M_p:=\frac{10^{\tau_p}-1}{p}\in\mathbb N,
 \qquad
 M_p\equiv0\pmod9.
 \label{p3-83-c44-s1:eq:reloj-primo-cierre-nonadico-rev11}
\end{equation}
Pour un nombre composé,
\[
 \operatorname{ord}_{n}(10)
 =\operatorname{lcm}_{p^a\parallel n}\operatorname{ord}_{p^a}(10).
\]
Par conséquent, l'horloge composée synchronise des horloges de puissances premières, tandis
que le nombre premier apporte une composante irréductible. La congruence nonadique
\eqref{p3-83-c44-s1:eq:reloj-primo-cierre-nonadico-rev11} décrit la clôture de phase ; le
sélecteur de primalité réside dans le coproduit réduit. La séparation de
ces deux fonctions évite d'attribuer une capacité sélective à la seule période décimale.

\subsection{Chaque constante est associée à son opérateur d'extraction}

Les cycles d'ordre \(3^m\), la limite de leurs traces de chaleur et la
transformée thêta--Mellin construisent la fonctionnelle
\begin{equation}
 \mathcal Z_3(s)
 =\frac{\displaystyle
 \int_0^\infty\Theta_+(x)x^{s/2-1}\,\mathrm dx}
 {\pi^{-s/2}\Gamma(s/2)}
 =\zeta(s),
 \qquad \Re s>1.
 \label{p3-83-c44-s1:eq:funcional-triadico-resumen-rev11}
\end{equation}
L'inversion thêta prolonge méromorphiquement la fonctionnelle. La partition
résiduelle
\[
 \zeta(s)=Z_0(s)+Z_1(s)+Z_2(s)
\]
sépare agrégat, échelle et orientation. Une seule source opératorielle admet
plusieurs applications d'extraction ---partie finie, coefficient de Laurent,
évaluation, dérivée régularisée ou caractère--- ; la coïncidence de source
n'identifie pas les constantes résultantes.

% REMATE_LOCAL_CONSTANTES_TABLA
\par

\begingroup
\small
\fontsize{10pt}{11.2pt}\selectfont
\setlength{\LTpre}{0.5\baselineskip}
\setlength{\LTpost}{0.5\baselineskip}
\begin{longtable}{@{}>{\raggedright\arraybackslash}p{.15\textwidth}
                    >{\raggedright\arraybackslash}p{.245\textwidth}
                    >{\raggedright\arraybackslash}p{.305\textwidth}
                    >{\raggedright\arraybackslash}p{.195\textwidth}@{}}
\caption{Définition, application d'extraction et portée mathématique des
constantes numériques complémentaires.}
\label{p3-83-c44-s1:tab:constantes-operadores-ontologia-rev11}\\
\toprule
\textbf{Objet} & \textbf{Définition standard} &
\textbf{Application et référence mathématique} & \textbf{Conclusion et condition}\\
\midrule
\endfirsthead
\toprule
\textbf{Objet} & \textbf{Définition standard} &
\textbf{Application et référence mathématique} & \textbf{Conclusion et condition}\\
\midrule
\endhead
Euler--Mascheroni \(\gamma\) &
Pour \(H_N=\sum_{n\leq N}n^{-1}\),
\(\gamma=\lim_{N\to\infty}(H_N-\log N)\), équivalente à la partie finie de
\(\zeta\) en \(s=1\). &
Si \(C_r=\operatorname*{FP}_{s=1}Z_r(s)\), alors
\(\gamma=C_0+C_1+C_2\), par
\eqref{eq:residual-three-covectors}. &
Identité exacte de la fonctionnelle résiduelle démontrée avec l'opérateur et la
normalisation déclarés.\\
Constantes de Stieltjes \(\gamma_n\) &
Si \(\zeta(1+z)-z^{-1}=\sum_{n\geq0}a_nz^n\), elles sont définies par
\(\gamma_n=(-1)^n n!a_n\). &
Les coefficients résiduels \(c_{r,n}\) satisfont
\(\gamma_n=(-1)^n n!\sum_r c_{r,n}\) et
\(L^{(n)}(1,\chi_{-3})/n!=c_{1,n}-c_{2,n}\) : c'est la décomposition
coefficient par coefficient de \eqref{eq:stieltjes} selon les trois classes
résiduelles de \eqref{eq:residual-three-covectors}. &
Décomposition exacte ; les extractions numériques vérifient les
coefficients, et non leur définition.\\
Ap\'ery \(\zeta(3)\) &
Valeur de la zêta de Riemann en \(s=3\). &
\(\mathcal Z_3(3)=\zeta(3)\), avec des préfixes bornés par des intervalles rationnels ;
sa résidence spectrale dans la tour graduée est fixée dans
\cref{eq:bendersky,eq:odd-zeta-bendersky}. &
Évaluation exacte de la fonctionnelle ; son irrationalité découle du théorème
classique d'Apéry \cite{Apery1979}.\\
Glaisher--Kinkelin \(A_{\rm GK}\) &
Constante de la croissance régularisée de l'hyperfactorielle. &
\(A_{\rm GK}=\exp(1/12-\mathcal Z_3'(-1))\), par \cref{eq:bendersky,eq:odd-zeta-bendersky}. &
Dérivée régularisée exacte ; \(1/12\) relève de cette normalisation
hyperfactorielle.\\
Euler--Kronecker \(\gamma_{F_{-3}}\) &
Partie finie logarithmique de la zêta de Dedekind du corps
\(F_{-3}=\mathbb Q(\sqrt{-3})\). &
Le caractère résiduel \(\chi_{-3}\) sélectionne le corps et produit
\(\gamma_{F_{-3}}=\gamma+L'(1,\chi_{-3})/L(1,\chi_{-3})\), par
\eqref{eq:euler-kronecker}. &
Identité exacte ; le choix du corps procède du caractère modulo trois.\\
Meissel--Mertens \(B_1\) &
Pour \(S_P(x)=\sum_{p\leq x}p^{-1}\),
\(B_1=\lim_{x\to\infty}(S_P(x)-\log\log x)\). &
C'est la partie finie de la trace première tronquée sur le secteur à une particule
du hamiltonien de Fock, par \eqref{eq:meissel-mertens} ; le sélecteur de
modes premiers est fixé avant d'effectuer la régularisation. &
Représentation opératorielle exacte une fois les modes premiers fixés ; le
sélecteur d'irréductibles est antérieur.\\
Feigenbaum \((\delta,\alpha_{\rm F})\) &
Constantes universelles d'échelle du point fixe de renormalisation pour la
cascade quadratique de doublement de période. &
La famille dodécaphasique déclarée produit au niveau huit
\(\delta^{\mathscr B}_8(0)=4.66921218915\ldots\) et
\(\alpha^{\mathscr B}_8(0)=2.50296847988\ldots\), par
\cref{sec:p3-final:cierre-feigenbaum}. &
Calcul fini exact jusqu'à \(n=8\) ; la convergence vers le point fixe universel
exige de démontrer la convergence de l'opérateur de renormalisation.
\(\alpha_{\rm F}\) et \(\alpha_{\rm HMT}\) sont
des objets distincts.\\
Catalan \(G\) &
\(G=L(2,\chi_{-4})=\sum_{n\geq0}(-1)^n/(2n+1)^2\). &
Le lecteur naturel est le caractère impair modulo quatre évalué en \(2\),
par \cref{hmtctx:catalan:sec:combinaciones-lectores-vecinos-rev6}. &
Identité classique exacte avec opérateur localisé ; ce n'est pas une sortie de la
fonctionnelle résiduelle modulo trois.\\
Khinchin \(K\) &
Moyenne géométrique presque sûre des chiffres de fraction continue pour le
système de Gauss. &
Le lecteur est la moyenne ergodique de \(\log a_1\) pour
\(T(x)=x^{-1}-\lfloor x^{-1}\rfloor\) et sa mesure invariante, par
\cref{hmtctx:khinchin:sec:combinaciones-lectores-vecinos-rev6}. &
Théorème ergodique classique incorporé ; il exige une dynamique différente de
la tour spectrale triadique.\\
\bottomrule
\end{longtable}
\endgroup

Le tableau fixe une règle de lecture : une constante est définie par son
objet, son opérateur et son application d'extraction, et non par l'apparition d'un
préfixe décimal dans un catalogue. Euler--Mascheroni et Stieltjes relèvent des
parties finies et des jets de Laurent ; Apéry, d'une évaluation ;
Glaisher--Kinkelin, d'une dérivée régularisée ; Euler--Kronecker, du caractère
du corps ; Meissel--Mertens, d'une trace première ; Feigenbaum, d'une famille
dynamique finie ; Catalan et Khinchin, d'opérateurs voisins distincts.

\subsection{Domaines, prémisses et conditions nécessaires de la classification généalogique du nombre}

Les définitions arithmétiques d'entier, de rationnel, d'irrationnel, d'algébrique et de
transcendant sont standard. La carte ternaire équilibrée, la séparation entre
section et valeur, le sélecteur natif d'irréductibles, la racine de Paley--Witt de
moins un, les six réalisations de \(-1/12\) et les lecteurs spectraux se
déduisent des applications citées dans chaque partie. Cette section les réunit sans
introduire de constante ni de sélecteur numérique supplémentaire.

La force de chaque affirmation est déterminée par son application :
\begin{enumerate}
\item les bijections, identités matricielles, coproduits, traces et
      évaluations indiqués sont exacts dans leurs domaines déclarés ;
\item la transcendance de \(e\) et \(\pi\), l'irrationalité de
      \(\zeta(3)\) et le théorème ergodique de Khinchin sont des théorèmes classiques
      incorporés comme contrôle ou classification ;
\item les approximants de Feigenbaum démontrent la branche et la profondeur
      calculées ; l'universalité exige l'opérateur de renormalisation et
      son théorème de convergence ;
\item l'égalité de deux scalaires provenant de domaines différents
      autorise une comparaison et exige un entrelaceur pour identifier les
      objets complets ;
\item une troncature décimale démontre le préfixe calculé, tandis que la
      classe arithmétique de la valeur infinie exige, selon le cas, une équation
      polynomiale ou un théorème d'irrationalité ou de transcendance.
\end{enumerate}

Avec ces règles, l'ontologie HMT du nombre acquiert une formulation
précise : la valeur archimédienne est la coordonnée publiée ; la section
compatible conserve la chaîne qui l'engendre ; l'opérateur détermine quelle
propriété est lue ; et ses hypothèses déterminent exactement quelle conclusion s'en
déduit.


\section[Clôture généalogique du concept de nombre]{Clôture généalogique du concept de nombre : compatibilité, mémoire et évaluation archimédienne}
\label{p3-83-c44-s2:chap:cierre-genealogico-numero}
\index{nombre!clôture généalogique}
\index{chiffre!coordonnée visible}
\index{valeur!évaluation archimédienne}
\index{conservation évolutive de l'unité}

La théorie des nombres développée dans HMT culmine dans une distinction de types
qui gouverne la transition ultérieure vers la mécanique dimensionnelle. Un
\emph{chiffre} est une coordonnée visible d'un état fini ; une \emph{valeur}
est l'image archimédienne publiée par un lecteur ; un \emph{nombre HMT} est
la généalogie compatible complète qui conserve, avec ses préfixes,
orientation, quotient, retenue, frontière et mémoire. La fibre d'une valeur
réunit les généalogies distinctes qui publient ce même scalaire.

Cette construction élargit la classification arithmétique ordinaire. Entier,
rationnel, irrationnel, algébrique et transcendant continuent de désigner des
propriétés de la valeur. HMT incorpore un deuxième axe : la manière dont une
section se prolonge, revient, conserve la mémoire et atteint une évaluation
stable. La clôture des Parties~I--III fixe les deux axes et l'application qui les relie.

\subsection{Chiffre, section, valeur et symétrie}

Soit
\[
 \mathcal S_0^{\rm surv}
 \xleftarrow{\rho_{1,0}}
 \mathcal S_1^{\rm surv}
 \xleftarrow{\rho_{2,1}}
 \cdots
\]
le système inverse d'états survivants, où chaque restriction
conserve les données typées du TPK qui appartiennent au niveau précédent, et soit
\[
 \mathcal S_\infty^{\rm surv}
 :=\varprojlim_n\mathcal S_n^{\rm surv}.
\]

\begin{definition}[Chiffre]
\label{p3-83-c44-s2:def:cifra-visible-genealogica}
Un lecteur positionnel \(L\) étant fixé, le chiffre ou bloc visible à la profondeur
\(n\) est
\[
 d_{L,n}(x_n):=\pi^{\rm vis}_{L,n}(x_n).
\]
C'est la coordonnée publiée par la présentation finie \(x_n\) ; la fibre de
\(\pi^{\rm vis}_{L,n}\) conserve les états qui partagent cette coordonnée.
\end{definition}

\begin{definition}[Nombre HMT et présentation]
\label{p3-83-c44-s2:def:numero-hmt-cierre-genealogico}
Un nombre HMT est une section compatible
\[
 \boldsymbol x=(x_n)_{n\geq0}\in\mathcal S_\infty^{\rm surv},
 \qquad
 \rho_{n+1,n}(x_{n+1})=x_n.
\]
Le couple \((\boldsymbol x,L)\) est une présentation du nombre. La section
conserve l'histoire complète ; le lecteur déclare quelles coordonnées de cette
histoire sont publiées.
\end{definition}

\begin{definition}[Valeur archimédienne]
\label{p3-83-c44-s2:def:valor-arquimediano-cierre-genealogico}
Supposons que \(L\) attribue à chaque \(x_n\) un cylindre compact
\(I_{L,n}(\boldsymbol x)\subset\mathbb R\), avec
\[
 I_{L,n+1}(\boldsymbol x)\subseteq I_{L,n}(\boldsymbol x),
 \qquad
 \operatorname{diam}I_{L,n}(\boldsymbol x)\longrightarrow0.
\]
La valeur publiée par la présentation est
\[
 \operatorname{Val}_L(\boldsymbol x)
 :=\text{l'unique élément de }
 \bigcap_{n\geq0}I_{L,n}(\boldsymbol x).
\]
\end{definition}

\begin{proposition}[Existence et unicité de l'évaluation]
\label{p3-83-c44-s2:prop:evaluacion-seccion-genealogica}
Les hypothèses précédentes déterminent une application
\[
 \operatorname{Val}_L:\mathcal S_\infty^{\rm surv}\longrightarrow\mathbb R.
\]
Chaque fibre \(\operatorname{Val}_L^{-1}(r)\) conserve les généalogies que
l'évaluation archimédienne représente par la même valeur \(r\).
\end{proposition}

\begin{proof}
La propriété d'intersection finie de la famille emboîtée de compacts
garantit une intersection non vide. Deux points distincts auraient une
distance positive majorée par tous les diamètres, en contradiction avec la
convergence de ceux-ci vers zéro.
\end{proof}

\begin{definition}[Symétrie généalogique]
\label{p3-83-c44-s2:def:simetria-genealogica-hmt}
Une symétrie du système de nombres HMT est une famille de bijections
\(g=(g_n)_{n\geq0}\) qui préserve la survie et les types, et satisfait
\[
 \rho_{n+1,n}\circ g_{n+1}
 =g_n\circ\rho_{n+1,n}
 \qquad(n\geq0).
\]
L'action
\[
 g_\infty(\boldsymbol x):=(g_n x_n)_{n\geq0}
\]
est un automorphisme de la limite inverse. Un invariant généalogique est une
fonction constante sur ses orbites.
\end{definition}

La symétrie est ainsi caractérisée comme une transformation compatible à toute
profondeur. Cette notion rend comparables la publication décimale,
l'incidence exceptionnelle et la réalisation spectrale ultérieure, en conservant le
type propre de chaque sortie.

\subsection{Conservation évolutive de l'unité}

Pour \(n\geq1\), soient
\begin{equation}
 \rho_9(n)=1+((n-1)\bmod 9),
 \qquad
 \kappa_9(n)=\frac{n-\rho_9(n)}9.
 \label{p3-83-c44-s2:eq:residuo-cociente-conservacion-unidad}
\end{equation}
Alors
\[
 n=9\kappa_9(n)+\rho_9(n),
 \qquad
 \rho_9(n)\in\{1,\ldots,9\}.
\]
La coordonnée \(\rho_9\) publie la phase nonadique et représente la classe nulle
par \(9\) ; \(\kappa_9\) conserve le nombre de tours. La racine numérique
publie la phase. La réduction modulo neuf accompagnée de son quotient publie
phase et mémoire.

La réalisation hélicoïdale associée est
\begin{equation}
 \mathcal H_9(n):=
 \left(
 \cos\frac{2\pi(\rho_9(n)-1)}9,
 \sin\frac{2\pi(\rho_9(n)-1)}9,
 \kappa_9(n)+\frac{\rho_9(n)-1}{9}
 \right).
 \label{p3-83-c44-s2:eq:helice-nonadica-conservacion-unidad}
\end{equation}

\begin{proposition}[Retour de phase et avancée de mémoire]
\label{p3-83-c44-s2:prop:retorno-fase-avance-memoria}
La réalisation \eqref{p3-83-c44-s2:eq:helice-nonadica-conservacion-unidad} satisfait
\[
 \mathcal H_9(n+9)=\mathcal H_9(n)+(0,0,1).
\]
Le retour visible \(9\to1\) clôt la phase et augmente d'une unité la mémoire
de l'état.
\end{proposition}

\begin{proof}
Les identités
\(\rho_9(n+9)=\rho_9(n)\) et
\(\kappa_9(n+9)=\kappa_9(n)+1\) fixent les coordonnées circulaires et
incrémentent exactement la hauteur.
\end{proof}

Cette identité formule la conservation évolutive : l'unité revient comme
phase et le nouvel état transporte la clôture précédente comme mémoire. La
complétion cubique réalise la même loi en dimension trois :
\begin{equation}
 10^3=(9+1)^3
 =9^3+3\cdot9^2+3\cdot9+1
 =729+243+27+1
 =729+271.
 \label{p3-83-c44-s2:eq:emrh-conservacion-evolutiva-unidad}
\end{equation}
L'intérieur possède \(729\) états ; \(243+27\) constitue la couronne
perforée ; le sommet \(+1\) complète l'unité de base mille. L'identité
\(729+270=999\) décrit l'état immédiatement antérieur à la clôture, et
\(729+271=1000\) décrit la complétion. La loi
\[
 10^d=(9+1)^d
 =\sum_{r=0}^{d}\binom dr9^{d-r}
\]
transporte la même unité entre les dimensions et conserve la généalogie de ses
strates.

\subsection{Syntaxe visible et sémantique de prolongement}

Une signature locale décrit la forme actuelle d'une branche. Sa sémantique de
survie est le système de prolongements qu'elle peut soutenir. Pour
\(x_n\in\mathcal S_n^{\rm surv}\), définissons
\[
 \operatorname{Ext}_m(x_n)
 :=\{x_m\in\mathcal S_m^{\rm surv}:\rho_{m,n}(x_m)=x_n\},
 \qquad m\geq n.
\]
Deux états qui satisfont
\[
 d_{L,n}(x_n)=d_{L,n}(y_n)
\]
peuvent posséder des arbres d'extension différents. Dans le témoin fini
\(60\to81\), deux branches ayant la même signature visible possèdent, respectivement,
zéro et douze prolongements. Le témoin établit que la continuité future
ajoute de l'information à la sortie locale. Le nombre HMT incorpore cette continuité
comme partie de l'objet.

La coinduction s'exprime par la compatibilité de tous les préfixes.
Les blocs de mille ou dix mille chiffres sont des publications tronquées d'une règle
uniformément prolongeable. La génération numérique atteint toute
profondeur demandée ; le relèvement à l'état enrichi conserve,
en outre, les hypothèses d'existence, d'unicité et de compatibilité de ses fibres.

\subsection{Classification arithmétique et classification généalogique}

Dans une base \(b\geq2\), la présentation canonique d'un réel sépare les classes
visibles suivantes :
\begin{enumerate}
\item un entier appartient à \(\mathbb Z\) et admet un mot ternaire
      équilibré fini ;
\item un rationnel possède un développement fini ou ultimement périodique ;
\item un irrationnel possède un développement non ultimement périodique ;
\item un algébrique est racine d'un polynôme non nul de \(\mathbb Q[X]\) ;
\item un transcendant appartient au complémentaire de la clôture algébrique de
      \(\mathbb Q\) dans \(\mathbb C\).
\end{enumerate}
La périodicité distingue les rationnels des irrationnels ; l'équation polynomiale
distingue les algébriques des transcendants. Ce sont des axes différents. Par exemple,
\(\varphi\) est irrationnel algébrique parce que
\(\varphi^2-\varphi-1=0\) ; la transcendance de \(e\) et \(\pi\) découle de
théorèmes classiques ; une généalogie HMT distingue, en outre, les
fonctions de propagation, de clôture et d'auto-échelle qui réalisent ces valeurs.

La même séparation régit \(e+\pi\). La somme des deux présentations
définit un objet composé et publie
\[
 \operatorname{Val}(\boldsymbol e)+
 \operatorname{Val}(\boldsymbol\pi)=e+\pi.
\]
Sa classe arithmétique constitue une question propre. La construction HMT
détermine l'application qui forme le composé et conserve les deux généalogies qui
interviennent.

\subsection{La famille opératorielle d'Euler comme grammaire de clôture}

L'équation classique d'Euler appartient à une famille quadratique unique. Soit
\(J_\sigma\) un générateur qui satisfait
\[
 J_\sigma^2=\sigma I,
 \qquad \sigma\in\{-1,0,+1\}.
\]
La séparation des puissances paires et impaires fournit
\begin{equation}
 \exp(tJ_\sigma)=C_\sigma(t)I+S_\sigma(t)J_\sigma,
 \label{p3-83-c44-s2:eq:euler-trigeometrico-cierre-numero}
\end{equation}
avec les réalisations exactes
\[
\begin{array}{ccl}
J_{-1}^2=-I&:&\exp(tJ_{-1})=\cos t\,I+\sin t\,J_{-1},\\[1mm]
J_{0}^2=0&:&\exp(tJ_{0})=I+tJ_{0},\\[1mm]
J_{+1}^2=I&:&\exp(tJ_{+1})=\cosh t\,I+\sinh t\,J_{+1}.
\end{array}
\]
Elliptique, parabolique et hyperbolique sont trois régimes d'une même
exponentielle opératorielle. La classification tritique détermine la géométrie de
l'évolution.

La branche hyperbolique contient l'auto-échelle de Fibonacci. Pour
\[
 F=\begin{pmatrix}0&1\\1&1\end{pmatrix},
 \qquad
 A=F^2=\begin{pmatrix}1&1\\1&2\end{pmatrix},
\]
on a \(\det A=1\), \(\operatorname{tr}A=3\) et
\[
 \eta=\operatorname{arcosh}\frac32=2\log\varphi.
\]
Il existe donc \(J_h^2=I\) tel que
\begin{equation}
 A=\exp(\eta J_h),
 \qquad
 A^n=\cosh(n\eta)I+\sinh(n\eta)J_h.
 \label{p3-83-c44-s2:eq:fibonacci-euler-hiperbolico-cierre-numero}
\end{equation}
La relation \(\varphi^2=\varphi+1\) exprime deux lectures du même processus :
récurrence additive locale et échelle multiplicative spectrale. Cette
compatibilité offre le modèle élémentaire de la double évaluation
somme--produit qu'APP formalise sur son support commun.

\subsection{Unification typée des géométries circulaire, complexe et hyperbolique}
\label{p3-83-c44-s2:sec:unificacion-compleja-circular-hiperbolica-numero}

La motivation géométrique initiale de HMT peut se formuler exactement comme une
architecture commune de support, de frontière et de générateurs. Le plan complexe
\(\mathbb C\) contient le cercle unité
\[
 S^1=\{e^{\mathrm i\theta}:\theta\in\mathbb R\},
\]
qui constitue la réalisation elliptique de
\eqref{p3-83-c44-s2:eq:euler-trigeometrico-cierre-numero}. Le demi-plan supérieur
\(\mathbb H=\{w\in\mathbb C:\operatorname{Im}w>0\}\) est transporté vers le disque
hyperbolique par la transformation de Cayley
\[
 \mathcal C(w)=\frac{w-\mathrm i}{w+\mathrm i},
 \qquad
 \mathcal C:\mathbb H\xrightarrow{\sim}\mathbb D,
 \qquad
 \mathbb D=\{z\in\mathbb C:|z|<1\}.
\]
La frontière idéale de \(\mathbb D\) est le même cercle \(S^1\), tandis que son
intérieur porte la métrique de Poincaré
\[
 \mathrm ds_{\mathbb D}^2
 =\frac{4\,|\mathrm dz|^2}{(1-|z|^2)^2}.
\]
La fusion consiste donc en un diagramme typé : la paramétrisation
trigonométrique publie la frontière circulaire ; la structure complexe
fournit la carte ; la métrique de Poincaré réalise l'intérieur hyperbolique.
Les trois régimes \(J^2=-I,0,+I\) incorporent en outre le seuil parabolique et
situent cette convergence géométrique au sein d'une seule famille opératorielle.

\begin{proposition}[Compatibilité euclidienne et hyperbolique du quart de tour]
\label{p3-83-c44-s2:prop:cuarto-giro-poincare-rev3}
Sur le disque unité
\[
 \mathbb D=\{z\in\mathbb C:|z|<1\},
\]
l'application
\[
 \mathcal R_{\pi/2}:\mathbb D\longrightarrow\mathbb D,
 \qquad
 \mathcal R_{\pi/2}(z)=\mathrm i z,
\]
est simultanément une rotation euclidienne et une isométrie de la métrique de
Poincaré
\[
 \mathrm ds_{\mathbb D}^{2}
 =\frac{4\,|\mathrm dz|^2}{(1-|z|^2)^2}.
\]
Son extension au bord \(S^1=\partial\mathbb D\) est le quart de tour
orienté. Dans la réalisation réelle du
\cref{thm:dos-realizaciones-fuente-cuadratica-rev3}, cette application
correspond à
\(\operatorname{Exp}((\pi/2)J_{\mathrm e})=J_{\mathrm e}\).
\end{proposition}

\begin{proof}
La multiplication par \(\mathrm i\) conserve \(|z|\) et
\(|\mathrm dz|\). Elle conserve donc aussi bien la métrique euclidienne que le
numérateur et le dénominateur de \(\mathrm ds_{\mathbb D}^{2}\). L'identité
exponentielle résulte de \(J_{\mathrm e}^2=-I\).
\end{proof}

Le plan complexe, le cercle trigonométrique et le disque hyperbolique ne sont pas
identifiés comme espaces métriques : ils partagent un support coordonné et
l'automorphisme de quart de tour après déclaration de chaque réalisation.

\subsection{Interface de frontière entre la rotation quadratique et la torsion minimale}
\label{p3-83-c44-s2:sec:interfaz-frontera-rotacion-torsion-minima-rev3}

La fonction de cette interface est de transporter une rotation intérieure déjà
construite vers une donnée de bord sans confondre l'endomorphisme quadratique
avec un défaut scalaire. Fixons un point de base
\(x_0\in\partial X\), un fibré réel
\(E_{\partial}\to\partial X\), une connexion \(\nabla_{\partial}\), un
lecteur de frontière \(q_{\partial}\) et une cotétrade discrète
\(e_{\partial}\in\Omega^1(\partial X;E_{\partial})\). Nous écrivons
\[
 \operatorname{Tors}(\partial X;E_{\partial})
 :=\Omega^2(\partial X;E_{\partial}),
 \qquad
 \operatorname{Hol}_{x_0}(\nabla_{\partial})
 \subseteq\operatorname{Aut}(E_{\partial,x_0}).
\]
Une réalisation de bord de l'interface devra construire les applications
\[
 \operatorname{Res}_{\partial}:
 \langle J_{\mathrm e}\rangle
 \longrightarrow\operatorname{Hol}_{x_0}(\nabla_{\partial}),
 \qquad
 \operatorname{Def}_{\partial}:
 \operatorname{Hol}_{x_0}(\nabla_{\partial})
 \longrightarrow\operatorname{Tors}(\partial X;E_{\partial}).
\]
La torsion déterminée par la connexion et la cotétrade a le type
\begin{equation}
 T_{\partial}
 :=d_{\nabla_{\partial}}e_{\partial}
 \in\Omega^2(\partial X;E_{\partial}).
 \label{p3-83-c44-s2:eq:torsion-borde-interfaz-euler-rev3}
\end{equation}
La réalisation physique de la torsion minimale exige en outre une fonctionnelle
\begin{equation}
 \Lambda_{\partial}:
 \operatorname{Tors}(\partial X;E_{\partial})\longrightarrow\mathbb R
 \label{p3-83-c44-s2:eq:funcional-borde-delta4-interfaz-rev3}
\end{equation}
telle que les trois applications satisfassent les conditions d'entrelacement
\begin{equation}
 \operatorname{Def}_{\partial}
 \bigl(\operatorname{Res}_{\partial}(J_{\mathrm e})\bigr)=T_{\partial},
 \qquad
 \Lambda_{\partial}(T_{\partial})=\Delta_4.
 \label{p3-83-c44-s2:eq:condiciones-borde-delta4-interfaz-rev3}
\end{equation}
Ces équations spécifient le comparateur à construire ; elles n'affirment pas
que les trois applications existent. Le scalaire \(\Delta_4\) est défini par
\cref{eq:delta4-electron-rev6}.

\begin{remark}[Condition exacte du comparateur de frontière]
La racine opératorielle \(J_{\mathrm e}^2=-I\), le quart de tour
\(\operatorname{Exp}((\pi/2)J_{\mathrm e})=J_{\mathrm e}\) et le scalaire
\(\Delta_4\) se déduisent dans leurs domaines respectifs. La chaîne de bord
précédente fixe le type des morphismes qui doivent les relier. Tant que
\(\operatorname{Res}_{\partial}\), \(\operatorname{Def}_{\partial}\) et
\(\Lambda_{\partial}\) ne sont pas construits, les conditions
\eqref{p3-83-c44-s2:eq:condiciones-borde-delta4-interfaz-rev3} définissent le comparateur
requis mais n'impliquent pas une réalisation physique ; la proximité angulaire de
ses extrémités ne remplace pas cet opérateur.
\end{remark}

APP, TRIT et TPK fournissent le support discret de cette architecture. Sur
\(X_9=(\mathbb Z/9\mathbb Z)^2\), les évaluations somme et produit induisent des
projecteurs \(P_{\Sigma,d}\) et \(P_{\Pi,d}\) qui satisfont
\[
 \|[P_{\Sigma,2},P_{\Pi,2}]\|=\frac{\sqrt2}{3},
 \qquad
 \|[P_{\Sigma,3},P_{\Pi,3}]\|=\frac{\sqrt3}{4}.
\]
Cette incompatibilité calculée est le précurseur fini de Heisenberg : deux
familles d'observables partagent un support et possèdent des décompositions
spectrales incompatibles. La forme alternée du tore,
\[
 \sigma((a,b),(c,d))=ad-bc\pmod9,
\]
s'intègre dans la représentation nonadique de Weyl--Heisenberg
\[
 X|n\rangle=|n+1\rangle,\qquad
 Z|n\rangle=\omega^n|n\rangle,\qquad
 ZX=\omega XZ,
 \qquad \omega=e^{2\pi\mathrm i/9}.
\]
Le commutateur publie la phase orientée et le groupe central résultant possède
\(9^3=729\) éléments.

Le bloc central d'APP fournit simultanément la réalisation
lorentzienne locale :
\[
 B_c=
 \begin{pmatrix}7&2\\2&7\end{pmatrix}
 =9P_++5P_-,
 \qquad
 M_c:=\frac{B_c}{\sqrt{45}}
 =\exp\!\left(\frac12\log\frac95\,R\right).
\]
Con \(\eta_{1,1}=\operatorname{diag}(1,-1)\),
\[
 M_c^{\mathsf T}\eta_{1,1}M_c=\eta_{1,1},
 \qquad
 \det M_c=1,
 \qquad
 M_c\in SO^+(1,1).
\]
Ainsi, la géométrie de Lorentz et le plan de Minkowski \(1+1\) émergent comme
réalisation spectrale exacte du bloc central. Toute élévation dimensionnelle
ultérieure doit entrelacer ce bloc local avec son nouveau codomaine. La séquence
\[
 \boxed{
 \text{APP--TRIT--TPK}
 \longrightarrow
 \begin{cases}
 \text{famille elliptique--parabolique--hyperbolique},\\
 \text{bloc de Lorentz--Minkowski }1+1,\\
 \text{représentation finie de Weyl--Heisenberg}
 \end{cases}}
\]
expose les développements géométriques et opératoriels qui précèdent la clôture
généalogique du nombre.

\subsection{Opérateurs de génération de constantes : domaines, germes et évaluations}

La fonction d'une constante est fixée au moyen de quatre données : définition
mathématique, domaine, opérateur ou lecteur et conclusion déduite. La
\cref{p3-83-c44-s2:tab:semillas-cierre-genealogico-numero} réunit les objets qui doivent
rester visibles lors de la clôture de la théorie des nombres.

\begingroup
\small
\setlength{\LTpre}{0.5\baselineskip}
\setlength{\LTpost}{0.5\baselineskip}
\begin{longtable}{@{}>{\raggedright\arraybackslash}p{.145\textwidth}
                    >{\raggedright\arraybackslash}p{.235\textwidth}
                    >{\raggedright\arraybackslash}p{.34\textwidth}
                    >{\raggedright\arraybackslash}p{.20\textwidth}@{}}
\caption{Constantes et germes dans la clôture généalogique du nombre.}
\label{p3-83-c44-s2:tab:semillas-cierre-genealogico-numero}\\
\toprule
\textbf{Objet} & \textbf{Définition ou identité} &
\textbf{Domaine et application HMT} & \textbf{Conclusion et condition}\\
\midrule
\endfirsthead
\toprule
\textbf{Objet} & \textbf{Définition ou identité} &
\textbf{Domaine et application HMT} & \textbf{Conclusion et condition}\\
\midrule
\endhead
\(\pi,e,\varphi\) &
Sections publiées par des cylindres numériques compatibles. &
La filtration
\(w_6\to w_{30}\to R_{36}\to G_9\) produit l'état terminal ; les lecteurs
séparent clôture, propagation et auto-échelle. &
Pour chaque profondeur finie prescrite, l'application produit des cylindres unitaires
compatibles ; la section infinie résulte du théorème coinductif.\\
\(\alpha_{\rm HMT}\) &
Clôture dodécaphasique avec retenue sur l'état terminal. &
La double lecture conduit à
\((P,E,\Phi,K,c_{13}=0)\) et à la sortie rationnelle de \(36\) chiffres ; la voie
\(E_{21/22}\) fournit un calcul indépendant. &
La récurrence est exacte pour l'état terminal déclaré et ne reçoit ni
\(\alpha\) ni CODATA ; la déduction minimale APP\(\to K\) exige de construire son
lecteur bidirectionnel complet.\\
\(i\) interno &
\(J^2=-I\) sur le module de Paley--Witt. &
Réalise une structure complexe finie et sépare des orientations conjuguées. &
Sur le corps de base déclaré, on a exactement \(J^2=-I\).\\
\(-1/12\) &
\(\zeta(-1)=-1/12\) et réalisations opératorielles normalisées. &
Incidence \(90\to120\), défaut dodécaphasique et fonctionnelles spectrales,
chacun avec son domaine et sa normalisation. &
Égalité exacte pour chaque application déclarée.\\
\(e+\pi\) &
Somme des présentations de \(e\) et \(\pi\). &
Objet composé de l'algèbre des présentations ; conserve les deux généalogies. &
Construction exacte ; classe arithmétique séparée.\\
Euler--Mascheroni \(\gamma\) et Stieltjes \(\gamma_n\) &
\(\gamma=\operatorname*{FP}_{s=1}\zeta(s)\) ; coefficients de Laurent. &
Décomposition résiduelle modulo trois et jets de Laurent. &
Identités exactes de la fonctionnelle déclarée.\\
Apéry \(\zeta(3)\) &
Évaluation de la zêta en \(s=3\). &
Lecteur \(\mathcal Z_3(3)\) avec intervalles rationnels de préfixe. &
Évaluation exacte ; irrationalité classique.\\
Glaisher--Kinkelin \(A_{\rm GK}\) &
\(A_{\rm GK}=\exp(1/12-\zeta'(-1))\). &
Dérivée régularisée du lecteur spectral. &
Identité exacte avec normalisation explicite.\\
Euler--Kronecker \(\gamma_{F_{-3}}\) &
Partie finie logarithmique de la zêta de Dedekind de
\(F_{-3}=\mathbb Q(\sqrt{-3})\). &
Le caractère \(\chi_{-3}\) sélectionne le corps et organise
\(L'(1,\chi_{-3})/L(1,\chi_{-3})\). &
Identité exacte du caractère déclaré.\\
Meissel--Mertens \(B_1\) &
\(B_1=\lim_{x\to\infty}(\sum_{p\le x}p^{-1}-\log\log x)\). &
Partie finie de la trace sur les modes premiers. &
Représentation opératorielle exacte avec modes fixés.\\
Catalan \(G\) et Khinchin \(K_0\) &
\(G=L(2,\chi_{-4})\) ; moyenne ergodique du système de Gauss. &
Lecteurs de caractère modulo quatre et de fraction continue. &
Théorèmes classiques incorporés avec leurs dynamiques propres.\\
Feigenbaum \((\delta,\alpha_F)\) &
Échelles de l'opérateur de renormalisation en doublement de période. &
La famille dodécaphasique produit des approximants de niveau huit. &
Le calcul détermine exactement le niveau huit ; l'universalité exige
de démontrer la convergence vers le point fixe de renormalisation.\\
Rogers--Ramanujan \(R(q)\) &
Fraction continue de niveau cinq et identité de Klein--Ramanujan pour
\(j\). &
Organise le niveau \(5\), \(A_5\), la limite
\(R(q)\to\varphi^{-1}\) et la réplicabilité de Hecke--Faber. &
Identités modulaires exactes ; la comparaison de coefficients ne couvre que la
profondeur finie calculée.\\
Horloge première \(p\) &
\(\tau_p=\operatorname{ord}_p(10)\),
\(M_p=(10^{\tau_p}-1)/p\equiv0\pmod9\) pour \(p\nmid30\). &
Les nombres premiers sont des périodes irréductibles ; les composés synchronisent des horloges
au moyen du plus petit commun multiple. &
Théorème arithmétique exact dans le domaine déclaré.\\
\bottomrule
\end{longtable}
\endgroup

Partie finie, valeur spéciale, dérivée régularisée, caractère, moyenne
ergodique, échelle de renormalisation et période première sont des applications distinctes. Leur
coexistence au sein de HMT procède d'une famille typée d'opérateurs et
d'applications d'extraction.

\subsection{Double projection d'une même généalogie}

Soit \(\mathfrak X:=X_\infty^{\rm cal}\) l'espace compact des histoires
enrichies avec jauge initiale transportée. Une jauge déclarée
\(c\in\mathscr C_{\rm exc}\) étant fixée, deux évaluations agissent sur une même section :
\[
 \mathfrak R:\mathfrak X\longrightarrow\mathbb R,
 \qquad
 \mathfrak E_c:\mathfrak X\longrightarrow\mathbf{Exc},
 \qquad
 \mathfrak E_c(x):=\mathfrak E(x,c),
\]
où \(\mathfrak R\) contracte les cylindres et \(\mathfrak E_c\) conserve les supports,
les incidences, la polarité et le transport de frontière.

\begin{theorem}[Unité causale de la double projection]
\label{p3-83-c44-s2:thm:unidad-causal-doble-proyeccion-numero}
Pour toute histoire \(\boldsymbol x\in\mathfrak X\), le couple
\[
 \mathfrak D_c(\boldsymbol x)
 :=(\mathfrak R(\boldsymbol x),\mathfrak E_c(\boldsymbol x))
\]
est une publication conjointe du même antécédent. La première coordonnée
détermine la valeur archimédienne ; la seconde détermine la classe d'incidence
exceptionnelle. Lorsque les lecteurs conjoints séparent les points,
\(\mathfrak D_c\) est injective.
\end{theorem}

\begin{proof}
Les deux composantes sont définies sur le même domaine et conservent les
applications de restriction. Si
\(\mathfrak D_c(\boldsymbol x)=\mathfrak D_c(\boldsymbol y)\), la séparation
conjointe des histoires implique \(\boldsymbol x=\boldsymbol y\).
\end{proof}

La branche archimédienne produit les cylindres et les valeurs de
\(\pi,e,\varphi\). La branche d'incidence prolonge les mêmes mots et le
même sceau vers Paley, Witt, Golay, le réseau de Leech,
\(V^\natural\), le groupe Monstre et Moonshine. La relation causale est la
préimage commune : les deux évaluations conservent des invariants différents de la
même histoire. Le réel est l'ombre contractée ; la symétrie exceptionnelle est la
mémoire d'incidence publiée par le même état.

\subsection{Clôture généalogique du concept de nombre}

Le concept de nombre est fixé par la chaîne
\[
 \boxed{
 \text{état fini}
 \longrightarrow\text{prolongement compatible}
 \longrightarrow\text{section généalogique}
 \longrightarrow
 \begin{cases}
   \text{chiffre visible},\\
   \text{valeur archimédienne},\\
   \text{incidence exceptionnelle}.
 \end{cases}}
\]
Le chiffre est une coordonnée ; la valeur est une projection ; le nombre est la
section complète ; la symétrie est l'automorphisme compatible qui conserve sa
mémoire. La conservation évolutive transporte cette identité à travers le
retour nonadique, la retenue hélicoïdale et la complétion
\(729+271=1000\). Ces distinctions étant closes, la mécanique dimensionnelle
réalise la même généalogie dans un nouveau codomaine.

\paragraph{Dépendances mathématiques.}
La limite inverse fournit la section ; le lecteur archimédien produit la
valeur ; la complétion binomiale et l'extension nonadique transportent unité et
mémoire ; les lecteurs complémentaires classent les sorties ; et la double
projection applique au même antécédent les applications archimédienne et exceptionnelle.
La chaîne Witt--Golay--Leech--Moonshine exige, à chaque étape,
l'entrelaceur d'incidence correspondant.

\section{Définition de l'atlas généalogique et conventions d'évaluation métrologique}

Cet atlas n'est pas un tableau de coïncidences. Son noyau mathématique conserve
cinq champs inséparables ; le statut logique et le critère de réfutation sont énoncés ensuite, sans prendre la place d'aucun des cinq :

\begin{equation}
\boxed{
\begin{gathered}
\text{antécédent}\ \big|\ \text{équation}\ \big|\
\ \text{invariant conservé}\\
\text{signification}\ \big|\ \text{valeur terminale}
\end{gathered}}
\qquad
\begin{gathered}
[\text{statut};\\ \text{critère de réfutation}]
\end{gathered}.
\label{eq:atlas-six-fields}
\end{equation}


\section{Corollaire HMT de confluence métrologique}
\label{sec:metrological-confluence}

\begin{definition}[Quintuplet métrologique]
Soit \(x\) une grandeur obtenue par une réalisation HMT--MD. Son
\emph{quintuplet
métrologique} est
\[
 \mathfrak C(x)=
 \bigl(A_x\mid E_x\mid I_x\mid S_x\mid V_x\bigr),
\]
où \(A_x\) est l'antécédent typé, \(E_x\) l'équation de
détermination, \(I_x\) l'invariant conservé par la réalisation,
\(S_x\) la signification
structurelle et \(V_x\) la valeur terminale. Il n'est pas permis de former \(V_x\) avant
d'avoir fixé les quatre autres composantes. Deux sorties ayant la même valeur décimale
ne sont pas égales si leurs antécédents ou leurs invariants diffèrent.
\end{definition}

\begin{corollary}[Confluence métrologique]
\label{cor:metrological-confluence}
Les sorties de vide, d'action, d'information, de température, de gravité,
de criticité, de spectre, de mélange et de cosmologie construites dans les
\cref{ch:vacuum,ch:metrology,ch:chaos,ch:spectral,ch:gauss,ch:gravity,ch:modular,ch:ew,ch:cosmo}
admettent les quintuplets du \cref{tab:metrological-confluence}. Dans chaque
ligne, la valeur de la cinquième colonne est une conséquence des quatre premières, et
non une donnée qui les sélectionne. Les lignes partagent des nœuds généalogiques
---mode \(30\), \(\log3\), déterminant \(D_A\), canaux \(90/120\) et changement
\(6\to8\)---, mais conservent leurs domaines et leurs opérateurs séparés.
\end{corollary}

\begingroup
\setlength{\tabcolsep}{2pt}
\renewcommand{\arraystretch}{1.22}
\scriptsize
\begin{longtable}{@{}
>{\raggedright\arraybackslash}p{25mm}
>{\raggedright\arraybackslash}p{37mm}
>{\raggedright\arraybackslash}p{24mm}
>{\raggedright\arraybackslash}p{29mm}
>{\raggedright\arraybackslash}p{32mm}@{}}
\caption{Corollaire de confluence métrologique HMT : les cinq colonnes forment
une seule unité probatoire.}\label{tab:metrological-confluence}\\
\toprule
Antécédent & Ecuaci\'on & Invariant conservé & Significado estructural &
Valeur terminale\\
\midrule
\endfirsthead
\multicolumn{5}{l}{\small\itshape Suite du corollaire de
confluence métrologique HMT}\\
\toprule
Antécédent & Ecuaci\'on & Invariant conservé & Significado estructural &
Valeur terminale\\
\midrule
\endhead
\midrule
\multicolumn{5}{r}{\itshape Suite à la page suivante}\\
\endfoot
\bottomrule
\endlastfoot

\multicolumn{5}{@{}l}{\bfseries I. Action}\\*
deux réalisations de la droite d'action
(\cref{eq:action-twoway,eq:action-rapidity}) &
\(R_{\rm act}=H_5/\eta_{\rm ret}\),
\(\xi_{\rm act}=\tfrac12\log R_{\rm act}\) &
Décennie et orientation de feuillet &
L'action absolue et son quotient relatif correspondent à des fonctionnelles
distinctes ; la
décennie n'est pas utilisée comme ajustement &
\(R_{\rm act}\), \(\xi_{\rm act}\)\\

\addlinespace
\multicolumn{5}{@{}l}{\bfseries II. Vide électromagnétique}\\*
\(T=q_+P_++q_-P_-\), \(30\to90/120\)
(\cref{eq:vacuum-operator,eq:three-four-completion}) &
\(V=(I-T^{90})(I-T^{120})^{-1}\) &
Positivité, spectre des feuillets et complétion \(3\to4\) &
Le vide est un opérateur bicouche ; \(90=3\cdot30\) et \(120=4\cdot30\)
déterminent le quatrième secteur sans introduire de constante SI &
\(r_+=0.9999996285\ldots\), \(r_-=0.9997241371\ldots\)\\

réponse spectrale
(\cref{eq:vacuum-four,eq:vacuum-traces}) &
\(\widehat\varepsilon=r_+^2\),
\(\widehat\mu=r_-^2\),
\(\widehat Z=r_-/r_+\),
\(\widehat c=(r_+r_-)^{-1}\) &
Produit et quotient des feuillets, stables sous \(V\mapsto V\otimes I_{9^m}\) &
Permittivité, perméabilité, impédance et propagation sont quatre fonctions
du même opérateur positif &
\(0.9999992571\ldots\), \(0.9994483505\ldots\),
\(0.9997245085\ldots\), \(1.0002763105\ldots\)\\

droites constitutives
(\cref{eq:vacuum-sections}) &
\(\mu\varepsilon=c^{-2}\), \(\mu/\varepsilon=Z^2\) &
Homogénéité dans \(L_S,L_C,L_Q\) &
La carte dimensionnelle est une section ultérieure, et non une unité accolée à un
nombre adimensionnel &
\(Z,\mu,\varepsilon,c\) comme sections exactes\\

\((Z_{\rm vac},h_a,\alpha_{\HMT})\)
(\cref{eq:charge,eq:quantum-electrical-identities}) &
\(e_a^2=2\alpha h_a/Z\), \(R_KG_0=2\), \(K_J\Phi_0=1\) &
Produit bidirectionnel et orientation opposée du flux/de la fréquence &
Les quanta électriques descendent du vide et de l'action ; la charge SI
n'agit pas comme sélecteur &
\(e_a,R_K,G_0,\Phi_{0,a},K_{J,a}\)\\

\addlinespace
\multicolumn{5}{@{}l}{\bfseries III. Information, température et rayonnement}\\*
trit uniforme et structure chronologique \(108\)
(\cref{eq:boltzmann-clock,eq:log3-triple}) &
\(k_B\Theta_{\rm clk}\log3=h_{\rm int}/(108t_0)=E_P\) &
Entropie d'une décision ternaire et période périodicité temporelle de (108) phases &
Un même quotient mesure l'information, la température et l'énergie sans identifier
leurs trois espaces d'états &
\(\log3=1.0986122886681098\ldots\)\\

\(G\hbar\) et périodicité temporelle de (108) phases
(\cref{eq:planck-energy,eq:planck-trit}) &
\(E_P=\pi\hbar/(54t_0)=k_B\Theta_{\rm clk}\log3\) &
Produit bidirectionnel \(G\hbar\) et contenu informationnel TRIT &
Le quantum géométrique de Planck et le quantum informationnel associé à un trit
confluent dans une même section d'énergie &
\(E_P=h/(108t_0)\)\\

spectre de Planck
(\cref{eq:wien-lambda,eq:wien-frequency}) &
\(\lambda_{\max}/\ell_0=108\log3/x_5\),
\(\nu_{\max}t_0=x_3/(108\log3)\) &
Même spectre, jacobiens distincts &
Les maxima par longueur et par fréquence sont des réalisations fonctionnelles d'un antécédent commun, et non deux
approximations de la même coordonnée &
\(23.896756779\ldots\), \(0.0237794888\ldots\)\\

continu bosonique
(\cref{eq:stefan-hmt,eq:photon-number,eq:photon-entropy}) &
\(n_\gamma\ell_P^3=2\zeta(3)/[\pi^2(\log3)^3]\) et lois \(T^4,T^3\) &
Occupation bosonique et contenu informationnel par trit &
Radiance, nombre, énergie et entropie forment une famille thermique ; ce ne
sont pas quatre calibrations indépendantes &
Coefficients exacts dans \(\pi,\zeta(3),\log3\)\\

\addlinespace
\multicolumn{5}{@{}l}{\bfseries IV. Gravité, aire et mémoire modulaire}\\*
clôture circulaire périodicité temporelle de (108) phases
(\cref{thm:gravity-circular-selector,eq:gravity-theta108}) &
\(G=\theta_{108}c^5t_0^2/\hbar\),
\(\theta_{108}=(54/\pi)^2\) &
Égalité du rayon de Compton et du rayon gravitationnel &
La constante gravitationnelle est une section sélectionnée par une clôture de
longueurs, et non un nombre greffé &
\(G_{\rm int}=(54/\pi)^2c^5t_0^2/\hbar\)\\

\(A_5/C_5\to P_5\)
(\cref{thm:gravity-p5-coupling,eq:gravity-G5-spectrum}) &
\(\mathcal G_5=GP_5\), \(P_5^2=P_5\), \(\Tr P_5=5\) &
Rang central et réduction TT &
Cinq composantes du support tensoriel précèdent les deux hélicités
physiques ; elles ne se confondent pas avec elles &
\(\Spec\mathcal G_5=\{G^{(5)},0^{(7)}\}\)\\

changement \(6\to8\), dérivation de Barbero--Immirzi et canaux \(90/120\)
(\cref{eq:modular-area-operator,eq:modular-68-two-readers}) &
\(A_{\rm phys}/\ell_P^2=\gammaH a_0(N_{90}+N_{120})\) &
Quantum d'aire, orientation hexade--octade et \(\ell_P^2\) &
\(\gammaH\) fixe le quantum géométrique ; cette monographie adopte cette
grandeur sans rouvrir sa dérivation &
\(\gammaH a_0=0.0016755940749\ldots\)\\

occupation \(N\) et mémoire \(D\)
(\cref{thm:modular-area-memory,eq:modular-inverse-channels}) &
\(D=90N_{90}+120N_{120}\),
\(N_{90}=4N-D/30\), \(N_{120}=D/30-3N\) &
Déterminant \(30\) et provenance du canal &
L'aire quantifie l'occupation totale ; le hamiltonien holographique
modulaire conserve le secteur \(90/120\) de provenance &
\((N_{90},N_{120})\) reconstruit exactement\\

état de bord
(\cref{thm:modular-tfd,eq:modular-oriented-information,thm:modular-first-law}) &
\(I_{\rm or}=36\log3-S(\rho_A)\),
\(\delta S=\beta_{90}\delta\langle\mathcal A_{90}\rangle+
\beta_{120}\delta\langle\mathcal A_{120}\rangle\) &
Spectre de Schmidt et rapport \(\beta_{120}/\beta_{90}=4/3\) &
Entropie d'intrication, information orientée et aire sont trois
lectures compatibles du même bord &
\(I_{\rm or}=0.0016691832\ldots\),
\((\beta_{90},\beta_{120})=(5.97265\ldots,7.96353\ldots)\)\\

réseau local infini
(\cref{thm:modular-typeIII,eq:modular-typeIII-lambda}) &
\(\lambda_\partial=\exp(-30\Dmod)\) &
Groupe des rapports de poids sous persistance \(90/120\) &
La limite est un générateur modulaire GNS, et non une matrice de densité globale
de classe trace dans l'UHF &
\(\lambda_\partial=0.996669646468957\ldots\)\\

\addlinespace
\multicolumn{5}{@{}l}{\bfseries V. Criticité et mode \(30\)}\\*
phases \(3m-1\)
(\cref{eq:pi-cancellation,eq:chaos-closed,eq:899}) &
\(u_0=1-\cos(37y)\sin(36y)/(12\sin3y)\),
\(899=30^2-1\) &
Normalisation quadratique et annulation de \(\pi\) avant l'itération &
La constante de criticité se dérive d'un profil clos dont le mode \(30\)
sépare également \(90\) de \(120\) &
\(u_0(x)=x^2-(715769/2424603)x^4+O(x^6)\)\\

suite superstable
(\cref{eq:negative-schwarzian,eq:renormalization}) &
\(\mathcal R(f)(x)=-a^{-1}f(f(-ax))\) &
Unimodalité, schwarzienne négative et combinatoire de doublement &
Les quotients finis certifient l'entrée dans la classe universelle ;
l'hyperbolicité du point fixe reste une obligation séparée &
\(\delta_8=4.66921218915\ldots\),
\(\alpha_8=2.50296847988\ldots\)\\

\addlinespace
\multicolumn{5}{@{}l}{\bfseries VI. Spectre, arithmétique et cylindres}\\*
quart de torsion
(\cref{eq:J-square,eq:catalan}) &
\(J^2=-I\), \(G=L(2,\chi_{-4})\) &
Orientation complexe et multiplicativité du caractère &
Catalan est le deuxième moment de l'opérateur de quart de torsion, et non un
nombre décimal associé par ressemblance &
\(G=0.915965594177219\ldots\)\\

produit CRT
(\cref{eq:chi12,eq:L1chi12,eq:euler-kronecker}) &
\(\chi_{12}=\chi_{-3}\chi_{-4}\),
\(\gamma_{\mathbb Q(\sqrt3)}=\gamma+L'/L(1,\chi_{12})\) &
Orientation triadique, quart de torsion et terme fini de Dedekind &
Le réseau \(3/4/12\) détermine un corps réel et sa constante de bord &
\(L(1)=0.7603459963\ldots\),
\(\gamma_{\mathbb Q(\sqrt3)}=1.05396560825\ldots\)\\

zêta triadique renormalisée
(\cref{eq:bendersky,eq:odd-zeta-bendersky}) &
\(A_k=\exp(H_kB_{k+1}/(k+1)-\mathcal Z_3'(-k))\) &
Niveau de dérivée zêta et normalisation de Bernoulli &
Bendersky--Glaisher organise \(\zeta(2n+1)\) comme une famille graduée de
déterminations &
\(A_0=\sqrt{2\pi}\), \(A_1=A_{\rm GK}\), famille \(A_k\)\\

cylindres cofinaux
(\cref{eq:khinchin,eq:levy,eq:gauss-entropy,eq:lochs}) &
\(L=\pi^2/(12\log2)\), \(h_G=2L\),
\(\Lambda_b=\log b/h_G\) &
Mesure de Gauss, croissance des dénominateurs et rythme de codage &
Khinchin, Lévy et Lochs sont des observables distinctes du même système
transporté &
\(K=2.6854520011\ldots\), \(L=1.1865691104\ldots\),
\(\Lambda_{10}=0.9702701144\ldots\)\\

opérateur de transfert de Gauss
(\cref{eq:gkw-operator}) &
\(\mathcal L_Gf(x)=\sum_{n\ge1}(n+x)^{-2}
f((n+x)^{-1})\) &
Masse, mesure invariante et complémentaire de moyenne nulle &
La valeur propre sous-dominante mesure la relaxation des cylindres ; elle n'est pas
sélectionnée à partir d'une valeur décimale tabulée &
\(\lambda_{\rm GKW}\) : intervalle certifié à construire\\

cycles irréductibles et lacunes
(\cref{eq:meissel-mertens,eq:twin-prime-constant,eq:vacancy-gamma,eq:vacancy-stieltjes}) &
\(B_1=\lim(\sum_{p\le x}p^{-1}-\log\log x)\),
\(\gamma_{\rm vac,j}^+=\varepsilon_{\rm vac}\gamma_j+\cdots\) &
Factorisation première, discrépance bornée et holonomie de lacune &
Somme locale, densité multiplicative et termes finis conservent des
constructions antécédentes distinctes &
\(B_1=0.2614972128\ldots\), \(C_2=0.6601618158\ldots\),
\(\varepsilon_{\rm vac}=0.04575749056\ldots\)\\

relèvement \(r^2=2\pmod{3^k}\)
(\cref{eq:pell-order}) &
\(u=3+2r\), \(\operatorname{ord}(u)=4\cdot3^{k-1}\) &
Opérateur de quart de torsion compatible à toute profondeur nonadique &
L'unité de Pell conserve simultanément l'orientation finie et la mémoire
profinie &
\(3+2\sqrt2=5.82842712474619\ldots\)\\

\addlinespace
\multicolumn{5}{@{}l}{\bfseries VII. Changement de feuillet, mélange et électrofaible}\\*
\(C^*\mapsto-C^*\), \(6\to8\)
(\cref{eq:modular-sheet-parities,eq:modular-ckm-coordinates}) &
\((A,C^*)\mapsto(A,-C^*)\),
\(-(8-6)/3=-2/3\) &
Parité de feuillet et défaut orienté hexade--octade &
Barbero et CKM sont des réalisations distinctes du changement d'état, et non
des définitions rétrospectives mutuelles &
\(A,C^*\) retrouvés rationnellement à partir des trois angles\\

matrice de coordonnées \(M\)
(\cref{thm:modular-ckm-reflection,eq:modular-Rckm-matrix}) &
\(R_{\rm CKM}=M\operatorname{diag}(1,-1,1)M^{-1}\) &
Involution rationnelle, \(R^2=I\), \(\det R=-1\) &
L'espace de mélange conserve exactement la coordonnée qui a subi
l'inversion de feuillet ;
la phase CP reste fixe dans cette involution &
Matrice rationnelle exacte, \(\det M=-3857/6\)\\

opérateur angulaire \(T\)
(\cref{thm:ew-weinberg,eq:ew-sin2thetaW}) &
\(D_A=\det T=e^{-2A_{\rm rad}}\),
\(\sin^2\theta_W=1-D_A\) &
Déterminant pair et adjacence des routes \(W/Z\) &
Le vide et Weinberg partagent un opérateur angulaire, mais déterminent des fonctions
spectrales différentes &
\(D_A=0.775129119247\ldots\),
\(\sin^2\theta_W=0.224870880753\ldots\)\\

carte de jauge dans l'approximation à l'arbre
(\cref{eq:ew-g,eq:ew-gprime,eq:ew-GF,cor:ew-Higgs-lambda}) &
\(g=2\sqrt{\pi\alpha/(1-D_A)}\),
\(G_F^{(0)}=\pi\alpha/[\sqrt2(1-D_A)m_W^2]\) &
Le même \(D_A\) conserve la séparation du canal neutre/chargé &
Les couplages, la constante de Fermi et l'autocouplage de Higgs
forment une composition typée ; les corrections radiatives ne sont pas incorporées
à la fonctionnelle interne &
\(G_F^{(0)}=1.1157162817\cdot10^{-5}\,\mathrm{GeV}^{-2}\),
\(\lambda_H=0.123847911548\ldots\)\\

fonctionnelles \((\alpha,\varphi,1000,55)\)
(\cref{eq:atlas-anom-mu}) &
\(a_\mu=\alpha/(2\pi)+\alpha(\varphi-1)/1000+\alpha^2/(55000)\) &
Somme de trois contributions à coefficients figés &
Contrôle de composition leptonique ; ne remplace pas la série radiative complète
de QED &
\(a_\mu=0.001165920713\ldots\)\\

\addlinespace
\multicolumn{5}{@{}l}{\bfseries VIII. Cosmologie et contrôles de portée}\\*
rayon \(R_{108}=\ell_P\)
(\cref{eq:cosmo-ds-H-Lambda,eq:cosmo-ds-temperature,eq:cosmo-ds-entropy}) &
\(H_{108}=\pi/(54t_0)\),
\(\Lambda_{108}=\pi^2/(972\ell_0^2)\) &
Un seul rayon sous la réalisation de Sitter &
Expansion, courbure, température et entropie sont quatre grandeurs d'une
même géométrie conditionnelle &
\(T_{\rm dS}/\Theta_{\rm clk}=\log3/(2\pi)\),
\(S_{\rm dS}/k_B=\pi\)\\

échelle de Perron
(\cref{eq:cosmo-perron-a,eq:cosmo-minus-six,eq:cosmo-EC-scaling}) &
\(a_n/a_0=\varphi^{n/3}\),
\(M_n=(a_n/a_0)^{-6}\) &
Mode stable \(\varphi^{-2n}\) et dilution quadratique du spin &
L'exposant \(-6\) représente la mémoire active dans la loi cosmologique ;
l'amplitude d'Einstein--Cartan exige son transducteur &
\(M_na_n^6=a_0^6\),
\(H_\varphi t_0=0.1604039416865\ldots\)\\

contrôles de composition
(\cref{eq:atlas-bohr,eq:atlas-born}) &
\(a_0=\hbar/(m_ec\alpha)\),
\(E_*=h_{\rm int}c_{\rm int}/(e_{\rm int}\ell_0^2)\) &
Homogénéité dimensionnelle et non-alimentation par le nom &
Bohr se compose à partir de grandeurs définies hors du champ de cette
monographie ; Born--Infeld exige encore
une action déterminantale qui sélectionne son échelle &
Sections exactes \(a_0,E_*\), sans promouvoir de nouveau sélecteur\\

\end{longtable}
\endgroup

\begin{proof}
Chaque ligne s'obtient en substituant exclusivement des objets déjà construits dans
l'équation citée dans sa première ou sa deuxième colonne. Dans le vide, le calcul
est spectral : le théorème fonctionnel pour \(V\) produit \(r_\pm\), et les quatre
constitutives sont des fonctions de ces deux valeurs propres. Dans la réalisation
thermique, la structure chronologique \(108\) et l'entropie uniforme d'un
trit produisent le produit
\(k_B\Theta_{\rm clk}\), après quoi Planck, Wien et le rayonnement sont des
compositions. Dans la gravité, la condition circulaire fixe
\(\theta_{108}\), tandis que \(P_5\) et la réduction TT fixent le type du
support. Dans la réalisation modulaire, le système linéaire de déterminant \(30\)
reconstruit les deux canaux et la décomposition spectrale de l'état réduit
produit l'entropie et la première loi. Les constructions de criticité,
de caractères, de Gauss et de
Pell sont prouvées dans leurs opérateurs respectifs. La réflexion CKM s'obtient
par conjugaison d'une involution diagonale, et le secteur électrofaible par
la fonction déterminant de l'opérateur angulaire. Enfin, les deux
réalisations cosmologiques emploient, respectivement, un rayon périodicité temporelle de (108) phases et la
valeur propre stable
de Perron. À aucune étape, la valeur décimale de la cinquième colonne n'intervient pour
choisir l'antécédent. Cela prouve la terminalité et la confluence déclarées.
\end{proof}

\section{Relations structurelles principales de confluence généalogique}

Le corollaire précédent ne se limite pas à énumérer des lignes : il détermine quels
opérateurs partagent un antécédent et le point exact où leurs domaines se
séparent. Les cinq configurations suivantes contiennent les relations
causales de plus grande densité du volume.

\subsection{Mode \(30\) : complétion harmonique, criticité et mémoire}

Le même entier apparaît dans trois constructions exactes :
\begin{equation}
\begin{aligned}
90&=3\cdot30, &120&=4\cdot30,
&\frac{1-q^{90}}{1-q^{120}}
 &=\frac{1+s+s^2}{1+s+s^2+s^3},\qquad s=q^{30},\\
899&=30^2-1=29\cdot31,
&&&D&=90N_{90}+120N_{120}.
\end{aligned}
\label{eq:atlas-mode30-confluence}
\end{equation}
La première ligne est la complétion \(3\to4\) de l'opérateur du vide ; la
deuxième normalise le germe dodécaphasique sans introduire \(\pi\) dans
l'itération ; la troisième conserve le nombre d'occupations arrivées par chaque canal
modulaire. L'invariant commun est la partition du mode \(30\) ; les
codomaines sont différents. C'est pourquoi le croisement est structurel, et non une
identification de \(V\), \(u_0\) et \(\HHM\).

\paragraph{Identités de complétion du mode \(30\), inversibilité et condition d'incompatibilité.}
Les trois identités de \cref{eq:atlas-mode30-confluence} sont exactes. La
confluence est invalidée si \(90/120\) cessent d'être des multiples consécutifs du
même mode, si \(899\ne30^2-1\), ou si l'inversion
\((N,D)\mapsto(N_{90},N_{120})\) perd son déterminant \(30\). Les différences
de spectre ne la réfutent pas : ces différences sont précisément celles qui
maintiennent les trois réalisations typées.

\subsection{\(\log3\) : invariant résiduel, information et température}

La confluence informationnelle se résume dans la chaîne
\begin{equation}
S_{\TRIT}=k_B\log3,\qquad
E_P=k_B\Theta_{\rm clk}\log3,\qquad
I_{\rm or}=36\log3-S(\rho_A).
\label{eq:atlas-log3-confluence}
\end{equation}
Le premier membre mesure une décision ternaire uniforme ; le deuxième détermine
sa grandeur énergétique dans la structure chronologique \(108\) ; le troisième compare la
capacité de \(36\) trits à l'entropie de l'état réduit de bord.
Le scalaire est conservé,
mais l'unité, la température et l'état réduit restent typologiquement
séparés. Cette
séparation permet précisément de lire l'identité
\(E_P=k_B\Theta_{\rm clk}\log3\) : la réalisation de Planck apporte le quantum
d'énergie, la réalisation TRIT son contenu informationnel et la structure
chronologique le transducteur.

\paragraph{Conservation de \(\log 3\), séparation dimensionnelle et condition d'incompatibilité.}
La triple lecture est une composition exacte de
\cref{eq:log3-triple,eq:planck-trit,eq:modular-oriented-information}. Elle est
invalidée si l'on remplace \(\log3\) par une entropie binaire, si
\(\Theta_{\rm clk}\) est fixé à partir du kelvin avant l'identité, ou si
\(S(\rho_A)\) est calculé avec un état distinct de celui obtenu par trace partielle.

\subsection{\texorpdfstring{\(D_A\)}{Déterminant angulaire} : déterminant des feuillets et relation angulaire de Weinberg}

L'opérateur angulaire \(T=q_+P_++q_-P_-\) possède deux familles fonctionnelles :
\begin{equation}
F(T)^2\longrightarrow
(\widehat\varepsilon,\widehat\mu,\widehat Z,\widehat c),
\qquad
\det T=q_+q_-=D_A\longrightarrow
\sin^2\theta_W=1-D_A.
\label{eq:atlas-DA-confluence}
\end{equation}
La première conserve le spectre complet de la réponse du vide ; la
seconde écarte l'orientation individuelle des feuillets et ne conserve que leur
produit. D'où le fait que \(D_A\) soit pair sous \(C^*\mapsto-C^*\) et puisse
déterminer le quotient \(W/Z\), tandis que l'impédance conserve le quotient
des valeurs propres. La relation contient plus d'information qu'une
coïncidence angulaire : c'est
une bifurcation fonctionnelle exacte du même opérateur antécédent.

\paragraph{Exactitude du déterminant angulaire et obstructions de l'interface électrofaible.}
Le côté spectral et l'identité \(D_A=e^{-2A_{\rm rad}}\) sont exacts dans la
carte déclarée. La réalisation de jauge conserve son statut d'interface jusqu'à
la construction de l'entrelaceur de bases. Un quotient de routes non adjacentes, une
dépendance impaire de \(C^*\) ou \((m_W/m_Z)^2\ne D_A\) réfutent cette interface sans
affecter l'opérateur du vide.

\subsection{Secteurs \(90\) et \(120\) : vide, aire et dynamique modulaire}

Les canaux \(90\) et \(120\) ont trois fonctions non interchangeables :
ce sont des exposants de propagation dans \(V\), des multiplicités de l'incidence
\(6\to8\) et des poids spectraux dans \(\HHM\). Leur rapport \(4/3\) réapparaît dans
les températures modulaires,
\[
\frac{120}{90}=\frac43=\frac{\beta_{120}}{\beta_{90}},
\]
mais l'opérateur du vide n'est pas le hamiltonien modulaire et une température
modulaire n'est pas une perméabilité. La confluence consiste en ce que la même
généalogie conserve trois grandeurs : propagation, occupation aréale et
mémoire de canal.

\paragraph{Reconstruction finie des canaux \(90/120\) et hypothèses de la limite modulaire.}
La reconstruction finie et le rapport \(4/3\) sont exacts. La classification
de type III exige la persistance infinie des deux canaux. La disparition de l'un
d'eux à partir d'un certain niveau, la perte de compatibilité des états ou l'obtention
d'un groupe des rapports distinct de
\(\{\lambda_\partial^n:n\in\Z\}\) réfuteraient cette classification, et non
l'inversion finie.

\subsection{Paramètre de Barbero--Immirzi, entropie et CKM comme réalisations du passage \(6\to8\)}

L'interface de confluence complète s'écrit sans transformer une
réalisation en cause des autres :
\begin{equation}
(6\to8,\,90/120,\,C^*\mapsto-C^*)\longrightarrow
\begin{cases}
\gammaH\ell_P^2\longrightarrow A_{\rm phys},\\
\rho_A\longrightarrow(-\log\rho_A,S,I_{\rm or},\delta S),\\
C^*\mapsto-C^*\longrightarrow
R_{\rm CKM}=M\operatorname{diag}(1,-1,1)M^{-1}.
\end{cases}
\label{eq:atlas-bic-entropy-ckm}
\end{equation}
La première réalisation fixe le quantum géométrique ; la deuxième mesure
l'information de l'état réduit et sa réponse aréale ; la troisième enregistre
l'orientation du feuillet dans
l'espace de mélange. En parallèle,
\(E_P=k_B\Theta_{\rm clk}\log3\) fixe le quantum informationnel d'énergie.
Ainsi, gravité, entropie et information sont mises en équations au moyen de
\((\ell_P^2,E_P)\), sans identifier \(\gammaH\), \(k_B\) ou les angles CKM.

\paragraph{Indépendance causale de Barbero--Immirzi, de l'information modulaire et du mélange CKM.}
La sortie \(\gammaH\) est un résultat antécédent démontré de sa dérivation correspondante ;
l'opérateur d'aire, la purification, la première loi locale et la réflexion CKM
sont exacts dans cette monographie. La classification modulaire infinie et la lecture
physique d'horizon conservent leurs hypothèses. Le diagramme est invalidé si
l'on utilise CKM pour recalibrer \(\gammaH\), si l'on assimile une hexade/octade
d'incidence à des faces/sommets sans entrelaceur, ou si
\(R_{\rm CKM}^2\ne I\) ou \(\det R_{\rm CKM}\ne-1\).

\section{Indépendance théorématique du corollaire de confluence métrologique}

La \cref{tab:metrological-confluence} ne renvoie pas à un tableau extérieur pour
compléter l'une quelconque de ses cinq colonnes. Chaque antécédent est défini dans le
corps ; chaque équation est démontrée ou typée dans le localisateur cité ;
chaque invariant déclare ce qui est conservé et aussi ce qui est écarté ; chaque
signification empêche l'identification par la valeur décimale ; et chaque valeur terminale
est obtenue après cette chaîne. Les statuts et critères de réfutation restent dans les
fiches détaillées qui suivent, où ils sont développés sans la compression inévitable
d'un tableau synoptique.

En conséquence, le tableau ne remplace pas le développement : il le complète. Il peut se lire
de manière synoptique pour reconnaître la confluence, mais sa validité procède des
théorèmes et des preuves du corps, et chaque ligne dispose ensuite d'une
fiche autonome avec antécédent, équation, signification, valeur, statut et
contrôle négatif.

\section{Action, vide et constantes quantiques--électriques}

\subsection*{Sections bidirectionnelles de l'action }
\textbf{Antécédent.} Deux réalisations orientées d'une même droite
d'action. \textbf{Équation.}
\[
R_{\rm act}=\frac{\hbar_{\rm pre}}{\hbar_{\rm ret}}
=\frac{H_5}{\eta_{\rm ret}},
\qquad \xi_{\rm act}=\tfrac12\log R_{\rm act}.
\]
\textbf{Structure et signification.} La décennie commune fixe le type ; le quotient
conserve l'orientation. \textbf{Valeur.} \(R_{\rm act}\) et
\(\xi_{\rm act}\) dans la carte interne, sans comparer des mantisses dimensionnelles à des
résidus. \textbf{Statut.} \emph{Théorème interne exact}.
\textbf{Critère de réfutation.} Utiliser la décennie commune comme correcteur relatif ou mélanger
les deux sections.

\subsection*{Opérateur constitutif du vide : spectre positif et reconstruction conjointe de \(\varepsilon,\mu,Z,c\)}
\textbf{Antécédent.} L'opérateur bicouche \(T=q_+P_++q_-P_-\) et la généalogie
\(30\to90/120\). \textbf{Équation.}
\[
V=(I-T^{90})(I-T^{120})^{-1},
\qquad
\frac{1-q^{90}}{1-q^{120}}
=\frac{1+s+s^2}{1+s+s^2+s^3},\quad s=q^{30}.
\]
Si \(\Spec V=\{r_+,r_-\}\),
\[
\widehat\varepsilon=r_+^2,
\quad\widehat\mu=r_-^2,
\quad\widehat Z=r_-/r_+,
\quad\widehat c=(r_+r_-)^{-1}.
\]
\textbf{Structure et signification.} Les quatre constitutives sont des fonctions
d'un seul opérateur positif ; chaque feuillet complète trois segments par un quatrième segment.
L'amplification \(V_m=V\otimes I_{9^m}\) conserve les fonctionnelles sous
l'espérance traciale.
\textbf{Valeur.}
\begin{align*}
\widehat\varepsilon&=0.9999992570966367027604416155\ldots,\quad
\widehat\mu=0.9994483504793269350899815559\ldots,\\
\widehat Z&=0.9997245085389372154555543466\ldots,\quad
\widehat c=1.0002763104861575261063862689\ldots.
\end{align*}
\textbf{Statut.} \emph{Théorème interne exact} ; amplification nonadique :
\emph{Formalisation nouvelle exacte}. \textbf{Critère de réfutation.} Perte de
positivité, résidu dans l'identité \(3\to4\), ou variation des
invariants sous l'espérance traciale.

\subsection*{Sections constitutives de \(\varepsilon\), \(\mu\), \(Z\) et \(c\) dans les cartes orientées}
\textbf{Antécédent.} Les quatre coefficients précédents et le cadre
\((u_S,u_C,u_T,u_Q)\). \textbf{Équation.}
\[
Z=\widehat Zu_Su_Q^{-2},\quad c=\widehat cu_C,\quad
\mu=\widehat\mu u_Su_C^{-1}u_Q^{-2},\quad
\varepsilon=\widehat\varepsilon u_S^{-1}u_C^{-1}u_Q^2.
\]
\textbf{Structure et signification.} La permittivité et la perméabilité ne
reçoivent pas d'unités après un calcul adimensionnel : ce sont des sections de droites
constitutives. \textbf{Valeur.} Les quatre sections précédentes, avec
\(\mu\varepsilon=c^{-2}\) et \(\mu/\varepsilon=Z^2\).
\textbf{Statut.} \emph{Théorème interne exact}. \textbf{Critère de réfutation.} Introduire
\(c,\mu_0,\varepsilon_0\) SI avant l'opérateur ou violer l'homogénéité.

\subsection*{Dérivation conjointe de la charge, de la résistance de von Klitzing, de la conductance, du flux et de la constante de Josephson}
\textbf{Antécédent.} \((Z_{\rm vac},h_a,\alpha_{\HMT})\).
\textbf{Équation.}
\[
e_a=\sqrt{\frac{2\alpha_{\HMT}h_a}{Z_{\rm vac}}},\quad
R_K=\frac{Z_{\rm vac}}{2\alpha_{\HMT}},\quad
G_0=\frac{4\alpha_{\HMT}}{Z_{\rm vac}},
\]
\[
\Phi_{0,a}=\sqrt{\frac{h_aZ_{\rm vac}}{8\alpha_{\HMT}}},
\qquad K_{J,a}=\Phi_{0,a}^{-1}.
\]
\textbf{Structure et signification.} \(R_K\) et \(G_0\) sont des invariants
bidirectionnels ; \(\Phi_0\) et \(K_J\) conservent l'orientation de l'action avec des
exposants opposés. \textbf{Valeur.} Les expressions exactes dans leurs droites ;
\(R_KG_0=2\) et \(K_J\Phi_0=1\). \textbf{Statut.}
\emph{Théorème interne exact}. \textbf{Critère de réfutation.} Utiliser la charge SI comme sélecteur ou
échouer à satisfaire \(\alpha=Ze^2/(2h)\).

\section{Constantes thermiques et rayonnement}

\subsection*{Relation de Boltzmann et contenu entropique TRIT }
\textbf{Antécédent.} Décision ternaire uniforme, section d'action et
structure chronologique \(108\). \textbf{Équation.}
\[
k_B\Theta_{\rm clk}\log3=\frac{h_{\rm int}}{108t_0}
=\frac{2\pi\hbar_{\rm int}}{108t_0},
\qquad
\log3=\frac{S_{\TRIT}}{k_B}=\frac{E_P}{k_B\Theta_{\rm clk}}.
\]
\textbf{Structure et signification.} \(\log3\) est simultanément un défaut
spectral, le contenu d'une décision et un rapport thermique ; les opérateurs
restent
distincts. \textbf{Valeur.} \(\log3=1.0986122886681098\ldots\) et la section
énergie--température précédente. \textbf{Statut.} \emph{Confluence
exacte}. \textbf{Critère de réfutation.} Séparer arbitrairement \(k_B\) de
\(\Theta_{\rm clk}\) ou inférer l'identité des domaines à partir de l'égalité du scalaire.

\subsection*{Loi de déplacement de Wien }
\textbf{Antécédent.} Le spectre de Planck évalué en
\(\Theta_{\rm clk}\). \textbf{Équation.}
\[
x_5=5+W_0(-5e^{-5}),\quad
\frac{\lambda_{\max}}{\ell_0}=\frac{108\log3}{x_5},
\]
\[
x_3=3+W_0(-3e^{-3}),\quad
\nu_{\max}t_0=\frac{x_3}{108\log3}.
\]
\textbf{Structure et signification.} Longueur d'onde et fréquence sont deux
paramétrisations spectrales avec des jacobiens distincts. \textbf{Valeur.}
\(\lambda_{\max}/\ell_0=23.896756779\ldots\) et
\(\nu_{\max}t_0=0.0237794888\ldots\). \textbf{Statut.}
\emph{Théorème interne exact}. \textbf{Critère de réfutation.} Confondre \(x_5\) avec \(x_3\) ou
introduire le kelvin avant le produit thermique.

\subsection*{Loi de Stefan--Boltzmann et gaz de photons }
\textbf{Antécédent.} La même échelle thermique, le continu bosonique et la
fonctionnelle \(\zeta(3)\). \textbf{Équation.}
\[
\mathcal F(\Theta_{\rm clk})
=\frac{2\pi^5}{15\,108^4(\log3)^4}
\frac{h_{\rm int}}{\ell_0^2t_0^2},
\]
\[
n_\gamma\ell_P^3=\frac{2\zeta(3)}{\pi^2(\log3)^3},
\qquad
\frac{u_\gamma\ell_P^3}{E_P}=\frac{\pi^2}{15(\log3)^4}.
\]
\textbf{Structure et signification.} \(\zeta(3)\) compte l'occupation bosonique ;
\(\log3\) fixe le contenu informationnel TRIT par cellule. \textbf{Valeur.}
Les coefficients
exacts précédents. \textbf{Statut.} \emph{Composition exacte avec
structure de départ explicite}. \textbf{Critère de réfutation.} Perte de facteurs
\(2\pi\), dimensions de flux incorrectes ou fusion des provenances mathématiques de
\(\zeta(3)\) et \(\log3\).

\section{Criticité et constantes dynamiques}

\subsection*{Profil dodécaphasique et normalisation par \(899\)}
\noindent\textit{Résultats directeurs : , et
.}\par
\textbf{Antécédent.} Douze phases \(3m-1\) et mode élémentaire \(30\).
\textbf{Équation.} Avec \(y=x/\sqrt{899}\),
\[
u_0(x)=1-\frac1{12}\sum_{m=1}^{12}
\cos\!\bigl(2(3m-1)y\bigr)
=1-\frac{\cos37y\sin36y}{12\sin3y},
\]
\[
899=30^2-1=29\cdot31,\qquad
u_0(x)=x^2-\frac{715769}{2424603}x^4+O(x^6).
\]
\textbf{Structure et signification.} La phase \(\pi\) s'annule avant
l'itération ; le profil résultant est analytique, unimodal quadratique et de
schwarzienne négative sur \([-1,1]\). \textbf{Valeur.} Profil et coefficients
rationnels exacts. \textbf{Statut.} \emph{Théorème interne exact}.
\textbf{Critère de réfutation.} Un deuxième point critique, une schwarzienne non négative ou un
résidu dans l'identité trigonométrique.

\subsection*{Approximants superstables de Feigenbaum et certificat hyperbolique d'universalité}
\textbf{Antécédent.} La suite superstable du germe précédent.
\textbf{Équation.}
\[
\mathcal R(f)(x)=-a(f)^{-1}f\bigl(f(-a(f)x)\bigr),
\qquad a(f)=-f(1).
\]
\textbf{Structure et signification.} Les quotients finis prouvent que la
famille entre dans la route appropriée ; l'universalité exige le point fixe et
son hyperbolicité. \textbf{Valeur.}
\(\delta_8=4.66921218915\ldots\),
\(\alpha_8=2.50296847988\ldots\). \textbf{Statut.} Approximants :
\emph{Certificat fini} ; limite universelle : \emph{Programme avec
critère de réfutation}. \textbf{Critère de réfutation.} Plus d'une valeur propre instable, queue non bornée
ou transversalité nulle.

\section{Constantes spectrales, arithmétiques et de lacune}

\subsection*{Structure complexe interne et constante de Catalan }
\textbf{Antécédent.} \(J^2=-I\) et le caractère
\(\chi_J(n)=i^{n-1}=\chi_{-4}(n)\) sur les impairs. \textbf{Équation.}
\[
G=L(2,\chi_J)=\sum_{n\ge0}\frac{(-1)^n}{(2n+1)^2}.
\]
\textbf{Structure et signification.} Catalan est le moment quadratique du
quart de torsion. \textbf{Valeur.} \(G=0.915965594177219\ldots\).
\textbf{Statut.} \emph{Théorème interne exact}. \textbf{Critère de réfutation.} Absence de
multiplicativité ou désaccord avec \(\chi_{-4}\).

\subsection*{Caractère dodécaphasique et corps quadratique réel }
\textbf{Antécédent.} Produit CRT
\(\chi_{12}=\chi_{-3}\chi_{-4}\). \textbf{Équation.}
\[
L(1,\chi_{12})=\frac{\log(2+\sqrt3)}{\sqrt3},
\qquad L(2,\chi_{12})=\frac{\pi^2}{6\sqrt3},
\]
\[
\gamma_{\mathbb Q(\sqrt3)}
=\gamma+\frac{L'(1,\chi_{12})}{L(1,\chi_{12})}.
\]
\textbf{Structure et signification.} L'orientation triadique et le quart de
torsion produisent le caractère réel de \(\mathbb Q(\sqrt3)\) ;
Euler--Kronecker est
le terme fini de sa zêta de Dedekind. \textbf{Valeur.}
\(0.7603459963009463\ldots\), \(0.9497031262940094\ldots\) et
\(1.05396560824848258\ldots\). \textbf{Statut.} \emph{Théorème interne exact} ;
l'identité composée conserve les hypothèses et les normalisations déclarées.
\textbf{Critère de réfutation.} Motif de résidus distinct de \((+,-,-,+)\) ou échec de
la factorisation de Dedekind.

\subsection*{Famille Bendersky--Glaisher }
\textbf{Antécédent.} Fonctionnelle zêta triadique renormalisée.
\textbf{Équation.}
\[
A_k=\exp\!\left(\frac{H_kB_{k+1}}{k+1}-\mathcal Z_3'(-k)\right),
\quad
\zeta(2n+1)=(-1)^{n+1}\frac{2(2\pi)^{2n}}{(2n)!}\log A_{2n}.
\]
\textbf{Structure et signification.} Apéry et les valeurs impaires de zêta sont des
niveaux d'une tour, et non des coïncidences isolées. \textbf{Valeur.}
\(A_0=\sqrt{2\pi}\), \(A_1=A_{\rm GK}\) et la famille exacte précédente.
\textbf{Statut.} \emph{Théorème interne exact}. \textbf{Critère de réfutation.} Échec dans
\(A_0,A_1\) ou dans l'identité de \(\zeta(2n+1)\).

\subsection*{Système dynamique de Gauss : \(K\), \(L\), \(\lambda_{\mathrm{GKW}}\) et cylindres transportés}
\textbf{Antécédent.} Évaluation des cylindres HMT et mesure
\(d\mu_G=dx/[\log2(1+x)]\). \textbf{Équation.}
\[
\log K=\int\log a_1\,d\mu_G,\quad
L=-\int\log x\,d\mu_G=\frac{\pi^2}{12\log2},
\]
\[
h_G=2L=\frac{\pi^2}{6\log2},\qquad
\Lambda_b=\frac{6\log2\log b}{\pi^2}.
\]
\textbf{Structure et signification.} Khinchin mesure les chiffres ; Lévy mesure la
croissance des dénominateurs ; Lochs compare les rythmes de codage.
\textbf{Valeur.}
\begin{align*}
K&=2.685452001065306\ldots,
&L&=1.18656911041562545\ldots,\\
h_G&=2.37313822083125091\ldots,
&\Lambda_{10}&=0.97027011439203393\ldots,\\
&&\Lambda_{729}&=2.77761896637430578\ldots.
\end{align*}
\textbf{Statut.} Cofinalité et identités : \emph{Exact dans le système
transporté} ; constante GKW : \emph{Programme ouvert avec critère de réfutation}.
\textbf{Critère de réfutation.} Non-cofinalité des cylindres, mesure non invariante ou absence
de trou spectral dans l'espace déclaré.

\subsection*{Cycles arithmétiques de Meissel--Mertens et séries régularisées de nombres premiers}
\textbf{Antécédent.} Sélecteur de cycles irréductibles et factorisation
exacte. \textbf{Équation.}
\[
B_1=\lim_{x\to\infty}\left(\sum_{p\le x}\frac1p-\log\log x\right),
\qquad
C_2=\prod_{p>2}\frac{p(p-2)}{(p-1)^2}.
\]
\textbf{Structure et signification.} \(B_1\) est une somme locale régularisée ;
\(C_2\) une densité multiplicative de paires. \textbf{Valeur.}
\(B_1=0.2614972128476428\ldots\),
\(C_2=0.6601618158468696\ldots\). \textbf{Statut.}
\emph{Démontré avec structure de départ explicite} : sélecteur premier et
borne de queue déclarés. \textbf{Critère de réfutation.} Utiliser des zéros ou des décimales pour choisir
les nombres premiers, ou affirmer la limite sans appliquer la borne de queue.

\subsection*{Lacunes, holonomie et constantes de Stieltjes }
\textbf{Antécédent.} Mot mécanique de lacunes et discrépance bornée.
\textbf{Équation.}
\[
\varepsilon_{\rm vac}=\log_{10}\frac{10}{9},\quad
\gamma_{\rm vac}^+=\lim_N\left(\sum_{n\le N}\frac{z_n}{n}
-\varepsilon_{\rm vac}\log N\right),
\]
\[
\rho_\varepsilon(k)=e^{2\pi ik\varepsilon_{\rm vac}},\qquad
\gamma_{\rm vac,j}^+=\varepsilon_{\rm vac}\gamma_j+
\sum_{n\ge1}\frac{(z_n-\varepsilon_{\rm vac})(\log n)^j}{n}.
\]
\textbf{Structure et signification.} Densité, terme fini et phase sont trois
types ; la tour prolonge le terme fini. \textbf{Valeur.}
\(\varepsilon_{\rm vac}=0.0457574905606751\ldots\) ; les autres sont déterminés
par leurs limites certifiées. \textbf{Statut.} \emph{Théorème interne exact avec
discrépance}. \textbf{Critère de réfutation.} Perte de la borne de discrépance ou
désaccord du niveau \(j=0\).

\subsection*{Relèvement nonadique de Pell }
\textbf{Antécédent.} \(r^2=2\pmod{3^k}\). \textbf{Équation.}
\[
u=3+2r,\qquad \operatorname{ord}(u)=4\cdot3^{k-1},
\qquad \varprojlim\langle u\rangle=C_4\times\mathbb Z_3.
\]
\textbf{Structure et signification.} Le quart de tour persiste à toute
profondeur ; la réalisation archimédienne est \((1+\sqrt2)^2\).
\textbf{Valeur.} \(3+2\sqrt2=5.82842712474619\ldots\).
\textbf{Statut.} \emph{Théorème interne exact}. \textbf{Critère de réfutation.} Ordre distinct
ou relèvements incompatibles.

\section{Gravité, aire et information modulaire}

\subsection*{Constante gravitationnelle et représentation tensorielle à cinq composantes }
\textbf{Antécédent.} Structure chronologique \(108\), action et projecteur
central de rang
cinq. \textbf{Équation.}
\[
G_a(\theta)=\theta\frac{c_{\rm int}^5t_0^2}{\hbar_a},
\qquad \theta_{108}=\left(\frac{54}{\pi}\right)^2,
\qquad \mathcal G_5=G_a(\theta_{108})P_5.
\]
\textbf{Structure et signification.} \(G\) est une section sélectionnée par la
clôture Compton--gravité ; \(P_5\) détermine un support à cinq composantes,
que la réduction TT ramène à deux hélicités sans masse.
\textbf{Valeur.} \(G_{\rm int}=\theta_{108}c^5t_0^2/\hbar\) et
\(\Spec\mathcal G_5=\{G\text{ avec multiplicité }5,
0\text{ avec multiplicité }7\}\).
\textbf{Statut.} Sélecteur : \emph{Théorème interne exact} ; réalisation du support :
\emph{Réalisation physique typée}. \textbf{Critère de réfutation.} Une autre racine positive,
\(\Tr P_5\ne5\), ou confondre cinq composantes avec cinq polarisations
physiques de GR sans masse.

\subsection*{Système de Planck, réciprocité \(G\leftrightarrow\hbar\) et invariance aréale}
\textbf{Antécédent.} Le produit bidirectionnel \(G_a\hbar_a\) et la périodicité temporelle de (108) phases.
\textbf{Équation.}
\[
\ell_P^2=\frac{G\hbar}{c^3}=\theta_{108}\ell_0^2,
\quad t_P=\frac{54}{\pi}t_0,
\quad E_P=\frac{\pi}{54}\frac{\hbar}{t_0},
\]
\[
F_P=\frac{c^4}{G},\qquad P_P=\frac{c^5}{G},\qquad
\rho_P=\frac{\hbar}{\theta_{108}^2c^5t_0^4}.
\]
\textbf{Structure et signification.} L'aire de Planck est invariante sous le
dédoublement de l'action et relie la gravité à l'échelle thermique.
\textbf{Valeur.} Les sections exactes précédentes.
\textbf{Statut.} \emph{Théorème interne exact}. \textbf{Critère de réfutation.} Combiner
\(G\) et \(\hbar\) de feuillets incompatibles ou échouer à satisfaire l'homogénéité.

\subsection*{Opérateur d'aire et hamiltonien modulaire : \(\widehat{\mathcal A}_{\rm HMT}\), \(H_{\rm mod}\) et canaux \(90/120\)}
\textbf{Antécédent.} Incidence \(6\to8\), canaux \(90/120\), sortie
\(\gammaH\) et \(\ell_P^2\). \textbf{Équation.}
\[
\frac{A_{\rm phys}}{\ell_P^2}=\gammaH a_0(N_{90}+N_{120}),
\qquad
\HHM=\Dmod(90N_{90}+120N_{120})+cI.
\]
Si \(N=N_{90}+N_{120}\) et
\(D=90N_{90}+120N_{120}\),
\[
N_{90}=4N-D/30,\qquad N_{120}=D/30-3N.
\]
\textbf{Structure et signification.} L'aire détermine l'occupation totale ; le
hamiltonien conserve la provenance du canal. Ensemble, ils reconstruisent le bord.
\textbf{Valeur.} \(\gammaH a_0=0.0016755940749224991\ldots\) dans la carte
déclarée. \textbf{Statut.} \emph{Théorème interne exact en coupe finie} ; on ne
redérive pas \(\gammaH\). \textbf{Critère de réfutation.} Commutateur non nul, déterminant
\(30\) perdu ou emploi de l'adjectif holographique sans transport au bord.

\subsection*{Entropie relative, purification thermochamp double et première loi modulaire aréale}
\textbf{Antécédent.} État réduit fidèle du bord et sa purification de
Schmidt. \textbf{Équation.}
\[
I_{\rm or}=36\log3-S(\rho_A),
\quad
|\mathrm{TFD}\rangle=\sum_n\sqrt{p_n}|n\rangle|n\rangle,
\]
\[
\delta S=\beta_{90}\delta\langle \mathcal A_{90}\rangle
+\beta_{120}\delta\langle \mathcal A_{120}\rangle,
\qquad \frac{\beta_{120}}{\beta_{90}}=\frac43.
\]
\textbf{Structure et signification.} La purification transforme l'entropie
réduite en intrication ; la première loi met en équations mémoire modulaire et aire.
\textbf{Valeur.} \(I_{\rm or}=0.0016691832338689407\ldots\),
\(\beta_{90}=5.97264854010709\ldots\),
\(\beta_{120}=7.96353138680946\ldots\).
\textbf{Statut.} \emph{Théorème interne exact local}. \textbf{Critère de réfutation.} Trace
partielle distincte de \(\rho_A\), \(I_{\rm or}<0\), ou utilisation pour des variations
finies sans le terme d'entropie relative.

\subsection*{Limite modulaire et paramètre de type III }
\textbf{Antécédent.} États locaux compatibles avec persistance infinie des
deux canaux. \textbf{Équation.}
\[
\HHM^{(m)}=-\log\rho_m,
\qquad
\lambda_\partial=\exp(-30\Dmod).
\]
\textbf{Structure et signification.} À l'infini vit un réseau local ou le
générateur modulaire GNS, et non une matrice de densité globale dans l'UHF.
\textbf{Valeur.}
\(\lambda_\partial=0.9966696464689573786\ldots\).
\textbf{Statut.} Réseau : \emph{Règle de type opératoriel} ; classification :
\emph{Théorème avec hypothèses de persistance}. \textbf{Critère de réfutation.} Absence
asymptotique d'un canal, état non produit ou rapport asymptotique distinct.

\section{Involution des feuillets, mélange et secteur électrofaible}

\subsection*{Involution des feuillets et transformation rationnelle CKM : unitarité, déterminant et caractère de Jarlskog}
\textbf{Antécédent.} Échange \(C^*\mapsto-C^*\) et donnée
\((h_6,o_8,t)=(6,8,3)\). \textbf{Équation.}
\[
R_{\rm CKM}=M\operatorname{diag}(1,-1,1)M^{-1},
\quad R_{\rm CKM}^2=I,\quad\det R_{\rm CKM}=-1,
\]
\[
-\frac{o_8-h_6}{t}=-\frac23,
\quad
A=\frac{1528\theta_{12}+3380\theta_{23}+801\theta_{13}}{3857},
\]
\[
C^*=6\frac{75\theta_{12}+176\theta_{23}-150\theta_{13}}{3857}.
\]
\textbf{Structure et signification.} Barbero et CKM sont des réalisations
fonctionnelles distinctes du changement \(6\to8\), adopté comme antécédent
commun ; l'involution induit une réflexion rationnelle de rang impair un.
\textbf{Valeur.} Matrice rationnelle exacte, \(\det M=-3857/6\) ; la phase CP
reste fixe sous cette involution, qui n'est pas une conjugaison CP.
\textbf{Statut.} \emph{Exact généalogique/interne}.
\textbf{Critère de réfutation.} \(R^2\ne I\), échec de \(RM=MD\), perte de
\(-2/3\), ou présentation de CKM comme cause rétrospective de \(\gammaH\).

\subsection*{Déterminant angulaire, paramètre de Weinberg et couplages électrofaibles}
\textbf{Antécédent.} Deux fonctions du même opérateur angulaire :
\(F(T)^2\) et \(\det T\) ; routes \(W/Z\) adjacentes.
\textbf{Équation.}
\[
D_A=\det T=e^{-2A_{\rm rad}},
\qquad \sin^2\theta_W^{\rm OS}=1-D_A,
\]
\[
g=2\sqrt{\frac{\pi\alpha}{1-D_A}},
\qquad g'=2\sqrt{\frac{\pi\alpha}{D_A}}.
\]
\textbf{Structure et signification.} Le vide et Weinberg sont des réalisations
spectrales distinctes d'un opérateur commun ; le déterminant est pair sous le
changement de feuillet.

\textbf{Valeur.} \(D_A=0.7751291192471564302\ldots\) et
\(\sin^2\theta_W=0.2248708807528435698\ldots\).

\textbf{Statut.} Relation angulaire : \emph{Théorème interne exact dans la carte} ;
réalisation de jauge : \emph{Réalisation conditionnelle}.
\textbf{Critère de réfutation.} Routes non adjacentes, quotient de masses
distinct de \(D_A\), ou schéma de jauge non
déclaré.

\subsection*{Constante de Fermi et couplage de Higgs }
\textbf{Antécédent.} \(\alpha_{\HMT}\), \(D_A\), routes \(W/H\) et carte
électrofaible à l'arbre. \textbf{Équation.}
\[
G_F^{(0)}=\frac{\pi\alpha}{\sqrt2(1-D_A)m_W^2},
\]
\[
\lambda_H=\frac{\pi\alpha}{2(1-D_A)}
\exp\!\left(6A_{\rm rad}+\frac{10}{3}C^*_{\rm rad}\right),
\qquad G_F^{(0)}m_H^2=\sqrt2\lambda_H.
\]
\textbf{Structure et signification.} Fermi est la limite effective de basse
énergie du canal chargé ; Higgs conserve l'orientation que le
déterminant de jauge élimine.
\textbf{Valeur.} \(G_F^{(0)}=1.1157162817\times10^{-5}\,\mathrm{GeV}^{-2}\)
dans la carte déclarée et \(\lambda_H=0.123847911548\ldots\).
\textbf{Statut.} \emph{Interface électrofaible à l'arbre}.
\textbf{Critère de réfutation.} Utiliser \(G_F\) observé pour choisir \(\Delta r\), omettre
le schéma/l'échelle ou sélectionner \(\lambda_H\) à partir de la masse observée.

\subsection*{Anomalie magnétique leptonique }
\textbf{Antécédent.} \((\alpha,\varphi,1000,54+1)\) avec des coefficients
figés avant la comparaison. \textbf{Équation.}
\begin{equation}
a_\mu=\frac{\alpha}{2\pi}
+\frac{\alpha(\varphi-1)}{1000}
+\frac{\alpha^2}{1000(54+1)}.
\label{eq:atlas-anom-mu}
\end{equation}
\textbf{Structure et signification.} C'est une évaluation de trois fonctionnelles HMT,
et non une substitution à la série radiative complète de QED.
\textbf{Valeur.} \(a_\mu=0.001165920713\ldots\).
\textbf{Statut.} \emph{Évaluation exacte avec provenance locale ouverte}.
\textbf{Critère de réfutation.} Choisir \(1000\) ou \(55\) à partir de la valeur expérimentale.

\section{Détermination des paramètres cosmologiques par mémoire torsionnelle et dynamique de Perron}

\subsection*{Réalisation de de Sitter par la périodicité temporelle de \(108\) phases}
\textbf{Antécédent.} Sélecteur circulaire \(R_{108}=\ell_P\) et réalisation de
Sitter. \textbf{Équation.}
\[
H_{108}=\frac{\pi}{54t_0},\quad
\Lambda_{108}=\frac{\pi^2}{972\ell_0^2},\quad
\frac{T_{\rm dS}}{\Theta_{\rm clk}}=\frac{\log3}{2\pi},\quad
\frac{S_{\rm dS}}{k_B}=\pi.
\]
\textbf{Structure et signification.} Un même rayon détermine l'expansion,
la courbure, la température et l'entropie. \textbf{Valeur.} Le quadruplet exact
précédent. \textbf{Statut.} \emph{Réalisation physique conditionnelle}.
\textbf{Critère de réfutation.} Fond non de Sitter ou réalisateur PCH sélectionnant un autre
rayon.

\subsection*{Réalisation de Perron--Einstein--Cartan }
\textbf{Antécédent.} Échelle \(V_n=\varphi^n\), \(D=3\), mode stable
\(\varphi^{-2n}\). \textbf{Équation.}
\[
\frac{a_n}{a_0}=\varphi^{n/3},\qquad
M_n=\left(\frac{a_n}{a_0}\right)^{-6},\qquad
H_\varphi=\frac{\log\varphi}{3t_0}.
\]
\textbf{Structure et signification.} L'exposant \(-6\) est une mémoire active et
coïncide avec la dilution quadratique du spin d'Einstein--Cartan.
\textbf{Valeur.} \(H_\varphi t_0=0.1604039416865344825\ldots\) et la loi exacte
\(M_na_n^6=a_0^6\). \textbf{Statut.} Échelle :
\emph{Théorème interne exact} ; lecture du spin : \emph{Réalisation conditionnelle}.
\textbf{Critère de réfutation.} Exposant distinct, paramètre temporel non
additif ou amplitude physique
choisie à partir d'un rebond.

\section{Séparation typée entre les rayons de Bohr et de Born--Infeld et les échelles engendrées en interne}

\subsection*{Rayon de Bohr comme échelle électrodynamique dérivée de \(\alpha\)}
\textbf{Antécédent.} Action, vide, \(\alpha\) et masse électronique de son
antécédent. \textbf{Équation.}
\begin{equation}
a_0=\frac{\hbar}{m_ec\alpha}.
\label{eq:atlas-bohr}
\end{equation}
\textbf{Signification et valeur.} C'est une composition
descendante, et non un sélecteur indépendant de cette monographie.
\textbf{Statut.} \emph{Hors du champ}. \textbf{Critère de réfutation.} Recalculer
\(m_e\) ici ou utiliser \(a_0\) pour sélectionner \(\alpha\).

\subsection*{Échelle de Born--Infeld et récupération de la limite de Maxwell}
\textbf{Antécédent.} \((h_{\rm int},c_{\rm int},e_{\rm int},\ell_0)\).
\textbf{Équation.}
\begin{equation}
E_*=\frac{h_{\rm int}c_{\rm int}}{e_{\rm int}\ell_0^2}.
\label{eq:atlas-born}
\end{equation}
\textbf{Structure et signification.} L'expression construit une échelle de
champ ; elle ne construit pas encore une action de Born--Infeld.
\textbf{Valeur.} Section exacte \(E_*\), sans détermination d'un
paramètre physique.
\textbf{Statut.} \emph{Programme ouvert}. \textbf{Critère de réfutation.} Identifier
par le nom ou la valeur décimale sans action déterminantale ni sélecteur.

\begin{proposition}[Résultat principal]
Le résultat métrologique n'est pas la valeur décimale isolée. C'est la conservation
d'une généalogie à travers des opérateurs, des droites de grandeur et des cartes, jusqu'à
ce que la valeur réelle apparaisse à la dernière position de la fiche.
\end{proposition}


\section{Théorèmes d'existence, d'unicité et de compatibilité des constantes généalogiques}

La construction détermine un secteur autonome de constantes intrinsèques.
Ses compositions, théorèmes croisés et certificats forment une chaîne
probatoire autonome. L'ordre directeur est

{\small
\begin{equation}
\APP\longrightarrow\TRIT\longrightarrow\TPK
\longrightarrow\text{réalisation fonctionnelle}
\longrightarrow\text{invariant}
\longrightarrow\text{section}
\longrightarrow\text{valeur terminale}.
\label{eq:final-causal-chain}
\end{equation}
}

Avec cet ordre, on obtient cinq clôtures principales.

\begin{enumerate}
\item Le vide électromagnétique cesse d'être une collection de quatre
formules : c'est un opérateur positif à deux feuillets dont le spectre détermine
\(\varepsilon,\mu,Z,c\), avec une amplification nonadique compatible.

\item La criticité cesse de commencer aux décimales de Feigenbaum : la
dodécaphase annule \(\pi\), le mode \(30\) produit \(899\), le germe se clôt
trigonométriquement et entre de manière démontrée dans la classe analytique
unimodale appropriée. L'hyperbolicité universelle est isolée comme un
certificat distinct.

\item La gravité cesse de n'être qu'une carte décimale : la périodicité temporelle de \(108\) phases sélectionne une
section de \(L_G\), le projecteur de rang cinq fixe le support tensoriel et
la réduction TT distingue ce support des deux hélicités sans masse.

\item Le changement hexade--octade acquiert trois réalisations liées :
quantum
aréal, mémoire modulaire et réflexion de mélange. Le hamiltonien holographique
modulaire est central dans cette interface de confluence, mais n'absorbe pas le vide, la criticité,
la gravité, l'arithmétique ni la cosmologie dans un titre artificiellement
unique.

\item Les constantes moins usuelles incluent Catalan, Lévy, Lochs et
Euler--Kronecker. Bendersky--Glaisher, les lacunes et les réalisations cosmologiques
apparaissent également comme des valeurs d'opérateurs et de dynamiques, et non comme une liste
à laquelle on cherche des coïncidences rétrospectives.
\end{enumerate}

\section{Caractère théorématique du corollaire métrologique}

Le \cref{cor:metrological-confluence} et la
\cref{tab:metrological-confluence} constituent la clôture probatoire du
développement. Ce ne sont pas un résumé abrégé du matériau précédent. Pour qu'une ligne
entre dans ce corollaire, elle doit conserver simultanément
\[
\text{antécédent}\mid\text{équation}\mid\text{invariant}
\mid\text{signification}\mid\text{valeur terminale}.
\]
Si l'on élimine l'une quelconque de ces cinq composantes, la ligne cesse d'être une
sortie HMT et devient une simple équivalence numérique. Le statut et
le critère de réfutation ne sont pas des colonnes de substitution : ils auditent de l'extérieur la force et le
domaine du quintuplet complet.

Cette formulation produit trois conséquences structurelles et mathématiques.
Premièrement, la valeur terminale peut être une expression exacte dans une droite de
grandeur ; il n'est pas exigé de la réduire au SI pour la reconnaître. Deuxièmement, un antécédent
peut admettre plusieurs réalisations fonctionnelles ---par exemple, \(T\)
détermine aussi bien la réponse du vide que \(D_A\)--- sans que ces
réalisations soient identifiées.
Troisièmement, un entier partagé comme \(30\), un rapport partagé comme \(4/3\)
ou un scalaire partagé comme \(\log3\) n'établit une confluence que lorsque l'on
exhibe l'application qui le transporte et l'invariant qu'il conserve.

En particulier, l'atlas permet de lire deux équations de clôture qui
restaient auparavant dispersées :
\begin{equation}
\boxed{E_P=k_B\Theta_{\rm clk}\log3},
\qquad
\boxed{
(6\to8)\longrightarrow
\bigl(\gammaH\ell_P^2,\,-\log\rho_A,\,R_{\rm CKM}\bigr).}
\label{eq:status-two-confluences}
\end{equation}
La première fait confluer quantum d'énergie, température et information. La
seconde détermine, à partir d'une interface de confluence commune, un quantum
géométrique, un générateur
modulaire et une réflexion de mélange. Aucune équation n'affirme que ses
composantes sont le même objet ; elle affirme que leurs généalogies ont été
conservées jusqu'au point exact de ramification.

Le contrôle négatif est également immédiat. Si \(E_P\) est utilisé pour sélectionner
\(k_B\), si CKM est utilisé pour recalibrer \(\gammaH\), si \(\log3\) est confondu
avec une matrice de densité ou si l'égalité de deux valeurs décimales remplace
l'entrelaceur, le corollaire ne s'applique pas. L'autonomie de la dérivation consiste en
ce que ces contrôles, comme les dérivations, sont formulés au sein
du développement lui-même.

\section{Stratification logique des identités internes, des réalisations dimensionnelles et des reconnaissances externes}

La force d'une affirmation ne se déduit ni du caractère extraordinaire de sa
conclusion ni de sa proximité avec une valeur connue. Elle se déduit de la chaîne
qui la soutient.

\begin{longtable}{p{36mm}p{96mm}}
\toprule
Statut & Contenu du développement\\
\midrule
\emph{Théorème interne exact}
& Opérateur du vide et ses constitutives ; mode \(30\) ; charge et quotients
quantiques ; échelle thermique ; profil clos du chaos ; réseau
\(3/4/12\) ; identités de Gauss ; sélecteur \(\theta_{108}\) ; famille de
Planck ; réflexion CKM ; loi
\(M_n=(a_n/a_0)^{-6}\) dans sa carte.\\

\emph{Composition mathématique exacte}
& Wien, Stefan--Boltzmann, gaz photonique, Euler--Kronecker,
Bendersky--Glaisher, relation angulaire de Weinberg, reconstruction conjointe
aire--mémoire et autocouplage de Higgs dans l'interface à l'arbre.\\

\emph{Démontré avec structure de départ explicite}
& Résultats qui dépendent d'une carte déclarée, d'une orientation, d'une
mesure, d'un état de bord ou d'une donnée de frontière, sans introduire la valeur
terminale comme sélecteur.\\

\emph{Théorème avec hypothèse explicite}
& Classification de la clôture GNS comme facteur hyperfini
\(III_{\lambda_\partial}\), sous fidélité, compatibilité et persistance
infinie des deux canaux \(90/120\).\\

\emph{Réalisation physique conditionnelle}
& Identification complète de jauge, \(G_F\) renormalisée, réalisation de
Sitter du
rayon CT108, transduction de la mémoire de Perron en densité de spin et
réalisation dynamique du support tensoriel.\\

\emph{Programme ouvert avec critère de réfutation}
& Hyperbolicité complète du renormalisateur de Feigenbaum, valeur propre
sous-dominante de Gauss--Kuzmin--Wirsing et construction d'une action de
Born--Infeld sélectionnée en interne.\\

\emph{Comparaison métrologique externe}
& Cartes SI, comparaisons CODATA/PDG, valeurs classiques des constantes et
classifications opératorielles appliquées uniquement après la construction des poids
HMT.\\
\bottomrule
\end{longtable}

Une interface ouverte ne diminue pas le théorème interne qui la précède. À
l'inverse, un théorème interne ne promeut pas automatiquement une interprétation
physique dont le transducteur n'a pas encore été construit.

\section{Obstructions transversales de typage, de causalité, d'orientation et de compatibilité nonadique}

La construction est soumise à réfutation. Les contrôles suivants
traversent tous les chapitres.

\begin{enumerate}[label=F\arabic*.]
\item \textbf{Typage dimensionnel.} Une identité qui n'est pas homogène dans les
droites \(L_S,L_C,L_T,L_Q\) est écartée, même si elle reproduit une valeur décimale.

\item \textbf{Non-alimentation par la valeur cible.} Aucune valeur SI, CODATA, PDG,
Feigenbaum, GKW, \(H_0\), \(G_F\) ou anomalie magnétique ne peut sélectionner
un opérateur, un coefficient ou une réalisation interne.

\item \textbf{Convention des feuillets.} Inverser silencieusement \(q_+\) et
\(q_-\), ou confondre l'involution du feuillet avec la conjugaison CP, invalide les
parités et les quotients.

\item \textbf{Compatibilité nonadique.} Une fonctionnelle qui change sous
l'amplification \(9^m\), perd l'unité ou ne conserve pas la mémoire du retour
ne possède pas la limite revendiquée.

\item \textbf{Non-fusion par scalaire.} L'égalité de deux valeurs terminales
n'identifie pas leurs domaines. Il faut un entrelaceur qui commute avec les
opérateurs et préserve l'orientation, la retenue et la frontière.

\item \textbf{Séparation fini--infini.} Un certificat jusqu'au niveau huit
ne prouve pas l'hyperbolicité universelle ; un hamiltonien modulaire fini n'est pas
\(-\log\rho_\infty\) dans l'UHF ; une convergence numérique sans borne de
queue ne certifie pas un produit premier.

\item \textbf{Transducteur physique.} La lecture de Sitter, d'Einstein--Cartan,
de Born--Infeld ou de jauge est bloquée si elle ne détermine pas l'objet complet de la
théorie d'arrivée et ses conditions de bord.
\end{enumerate}

\section{Conditions nécessaires pour les extensions mathématiques et les réalisations physiques non équivalentes}

Il ne reste pas de réserve générique du type ``il manque à relier HMT à la
physique''. Il reste des obligations localisées :

\begin{description}[style=nextline,leftmargin=37mm]
\item[CHAOS-OBL-1]
Certifier dans un espace de fonctions analytiques paires le point fixe de
\(\mathcal R\), une unique valeur propre instable, la borne de queue et la
transversalité du germe HMT.

\item[GAUSS-OBL-1]
Fixer l'espace fonctionnel de l'opérateur de Gauss, isoler la valeur propre
sous-dominante par des intervalles et démontrer la stabilité du projecteur sous les
partitions \(9^m\).

\item[HMOD-OBL-1]
Vérifier la compatibilité des états locaux et la persistance avec
densité positive des deux canaux avant d'appliquer la classification
d'Araki--Woods. Ne pas utiliser une densité globale de classe trace inexistante.

\item[EW-OBL-1]
Construire l'entrelaceur entre les feuillets HMT et la base neutre/chargée électrofaible,
et dériver \(\Delta r\) sans employer la valeur expérimentale de \(G_F\).

\item[COS-OBL-1]
Pour de Sitter, démontrer que le réalisateur PCH sélectionne dynamiquement
\(R_{108}\). Pour Einstein--Cartan, construire l'amplitude, le courant, le signe et la
carte dimensionnelle de l'application \(M_n\mapsto s_n^2\).

\item[BORN-OBL-1]
Construire l'action déterminantale de Born--Infeld et démontrer que l'échelle
\(E_*\) est son paramètre, au lieu de l'identifier par les dimensions.

\item[ANOM-OBL-1]
Clore la provenance locale de \(1000\) et \(54+1\) avant de promouvoir
l'évaluation de \(a_\mu\) en dérivation complète.
\end{description}

Chaque obligation possède un objet, un domaine et un critère de réfutation. Aucune n'autorise à
rouvrir les blocs exacts qui la précèdent.

\section{Délimitation formelle du domaine des dérivations généalogiques}

\begin{itemize}
\item Ne redérive pas \(\pi,e,\varphi,\alpha_{\HMT}\) ; les utilise comme entrées déjà
construites lorsqu'une composition ultérieure l'exige.
\item N'entre pas en concurrence avec la Loi générale des masses et ne recalcule pas la masse électronique.
\item Ne redérive pas \(\gammaH\) ; développe sa signification aréale,
informationnelle et modulaire.
\item N'identifie pas le support à cinq composantes à cinq polarisations
d'un graviton sans masse.
\item Ne confond pas une approximation finie de Feigenbaum avec le théorème
universel ni une échelle dimensionnelle avec une action de Born--Infeld.
\item Ne fait pas du hamiltonien holographique modulaire le nom ou
l'explication exclusive de toutes les constantes. C'est une pièce centrale de
l'interface de confluence aire--information--mélange au sein d'un atlas plus vaste.
\end{itemize}

\section{Capacité reconstructive de la généalogie commune sur les constantes mathématiques et physiques}

La singularité métrologique de HMT--MD ne consiste pas à produire beaucoup de
décimales. Elle consiste en ce que la valeur réelle apparaît à la fin d'une séquence qui
conserve la mémoire. La perméabilité et la permittivité sont les deux valeurs propres
quadratiques d'une même réponse ; \(G\) est la section qui clôt un rayon
de Compton avec un rayon gravitationnel ; \(\log3\) est simultanément contenu de
décision et rapport thermique ; Catalan est un moment du quart de torsion ;
\(\gammaH\) oriente le quantum d'aire ; le hamiltonien modulaire conserve quel
canal a produit l'occupation ; la réflexion CKM conserve la coordonnée
responsable du changement de feuillet ;
l'exposant \(-6\) conserve la mémoire stable en déterminant l'échelle
cosmologique.

\begin{proposition}[Résultat principal]
La signification structurelle démontrée dans ce développement s'exprime
par une règle de réalisation : un même antécédent peut admettre
plusieurs fonctionnelles sans perdre sa généalogie, et chacune détermine
simultanément la signification de la constante et la raison de sa valeur
terminale.
\end{proposition}

L'affirmation finale est donc plus précise que ``HMT reproduit des
constantes''. HMT--MD construit des opérateurs, des caractères, des sections et des
transducteurs dont les constantes sont des évaluations. Là où la chaîne
est close, la valeur et sa signification sont démontrées ensemble. Là où subsiste une
interface, le critère de réfutation permet de continuer sans dégrader ce qui est déjà établi ni
présenter comme théorème ce qui est encore un programme.

%
\endgroup

\HMTFlushCurrentChapter
\chapter{Complétude généalogique des sections HMT et décidabilité de chaque niveau fini}%
\begingroup
\renewcommand{\chapter}[2][]{}%
\chapter{Complétude généalogique des sections HMT et décidabilité de chaque niveau fini}
\label{chap:p3-final:53:completitud_finita}
\label{chap:p3-83-45-completitud_finita}

\section{Extension conservative de la complétude finie par des multiplicités projectives}

Une tour HMT concrète est formée d'ensembles finis, de relations de
transition, d'applications de restriction et de lecteurs. Tout cela peut
être codé dans ZF ou ZFC. La collection de toutes les tours constitue une
classe définissable ; il n'est pas nécessaire de postuler que HMT forme déjà une théorie
autonome des classes comparable à NBG.

\begin{proposition}[Codage conservatif]
Chaque objet généalogique fini ou dénombrable utilisé dans ce cahier peut
être représenté par des ensembles et des relations de ZFC. Par conséquent, les
formalisations construites ici ne démontrent pas une contradiction de ZFC.
\end{proposition}

\begin{proof}
Chaque niveau \(X_m\) est un ensemble, chaque restriction \(p_{n,m}\) est un
sous-ensemble de \(X_n\times X_m\), et une tour dénombrable est une fonction de
domaine \(\mathbb N\) qui code ces données. Les types, lecteurs et mémoires
s'ajoutent comme relations ou fonctions. La limite inverse est un sous-ensemble
du produit \(\prod_mX_m\), que ZFC construit par Remplacement, Union,
Ensemble des parties et Séparation.
\end{proof}

Ce que HMT discute n'est pas la légitimité de ce codage, mais la pratique
qui consiste à oublier les applications et à ne conserver qu'une ombre. Si l'on applique, par exemple,
le lecteur
\[
  L:\OmegaInf\longrightarrow\mathbb R
\]
ou le profil cardinal
\[
  \mathbb X\longmapsto(|X_m|)_{m\ge0},
\]
des objets non isomorphes peuvent être confondus. La projection est valide dans son
codomaine ; ce qui serait invalide serait de lui attribuer une fidélité sans preuve.

\begin{remark}[Conditions de validité]
Il n'est pas affirmé que HMT satisfasse en interne tous les schémas de ZF, IZF ou
NBG. Séparation, Remplacement, Ensemble des parties complet et une hiérarchie transfinie générale
n'ont pas été vérifiés conjointement dans une axiomatique HMT autonome.
\end{remark}

\section{Catégorie d'objets généalogiques}

\begin{definition}
Un objet généalogique est
\[
  \mathbb X=
  \bigl((X_m)_{m\ge0},(p_{n,m})_{n\ge m},
  \mathcal T,\mathcal O,\mathcal M\bigr),
\]
où \(\mathcal T\) enregistre les types, \(\mathcal O\) les observables et
\(\mathcal M\) les mémoires. Un morphisme
\(f:\mathbb X\to\mathbb Y\) est une famille \(f_m:X_m\to Y_m\) qui satisfait
\[
  p^Y_{n,m}f_n=f_mp^X_{n,m}
\]
et préserve les données indiquées.
\end{definition}

La catégorie ainsi obtenue est une présentation concrète d'objets
profinis enrichis. Le foncteur des sections globales
\[
  \Gamma(\mathbb X)=\varprojlim_mX_m
\]
retrouve les histoires complètes. La petite catégorie de chemins de
\(\TPK\) admet en outre l'enveloppe de préfaisceaux
\[
  \widehat{\mathcal P_{\TPK}}
  =
  \mathbf{Set}^{\mathcal P_{\TPK}^{op}}.
\]
Cette construction fournit une logique contextuelle et des faisceaux de données, mais
n'autorise pas à déclarer la booléanité, le choix ou l'existence de suffisamment de points sans étudier
le topos concret.

\section{Fondation des préfixes et coinduction des sections}

Le graphe des préfixes est gradué par la profondeur :
\[
  \operatorname{rk}(x_m)=m.
\]
Appliquer le parent diminue strictement le rang. Il n'existe pas de cycles
parentaux
\[
  x_m\to x_{m-1}\to\cdots\to x_m
\]
compatibles avec cette graduation.

\begin{proposition}
La bonne fondation du graphe des préfixes et l'existence coinductive d'une
section infinie sont compatibles.
\end{proposition}

\begin{proof}
La bonne fondation concerne la direction de descente par les parents. La
coinduction concerne une famille \((x_m)\) compatible avec toutes les
descentes. Aucun \(x_m\) n'appartient à lui-même ni ne réapparaît comme le même
nœud ; la section est une fonction sur \(\mathbb N\), et non un cycle
d'appartenance.
\end{proof}

Cette distinction éclaire le retour \(9\to1\). La phase peut se répéter, mais
la coordonnée de mémoire augmente. Une boucle dans la projection de phase n'est pas un
cycle dans le graphe gradué des états.

\section{Fibration des ordinaux par des histoires généalogiques compatibles}

Von Neumann représente chaque ordinal comme l'ensemble de ses prédécesseurs
\citep{vonneumann1923}. HMT peut relever cette représentation sans la modifier.

\begin{definition}[Relèvement généalogique d'un ordinal fini]
Soit \(\boldsymbol x=(x_m)\) une section. On définit
\[
  \widehat n_{\boldsymbol x}
  =
  \{(k,x_k):k<n\}.
\]
\end{definition}

\begin{proposition}
La projection sur la première coordonnée satisfait
\[
  \pi(\widehat n_{\boldsymbol x})=n,
\]
et la suite obéit à
\[
  \widehat{n+1}_{\boldsymbol x}
  =
  \widehat n_{\boldsymbol x}\cup\{(n,x_n)\}.
\]
\end{proposition}

\begin{proof}
Les deux identités s'obtiennent directement de la définition. L'information
ajoutée ne modifie pas l'ordre sous-jacent ; elle distingue les routes qui réalisent le
même ordinal.
\end{proof}

Pour \(\omega\), on écrit
\[
  \widehat\omega_{\boldsymbol x}
  =
  \{(n,x_n):n\in\mathbb N\}.
\]
Le support de chaque histoire reste dénombrable, bien que l'espace de
toutes les histoires puisse avoir la puissance du continu lorsque la ramification
persiste.

\begin{remark}[Extension et domaine ouvert]
L'extension
\[
S_{\alpha+1}=\operatorname{Ext}_\alpha(S_\alpha),
\qquad
S_\lambda=\varprojlim_{\beta<\lambda}S_\beta
\]
est un programme transfini bien typé, et non un résultat déjà démontré pour
tout ordinal. Il faut vérifier la cohérence des limites, la taille des
niveaux et l'existence d'extensions.
\end{remark}

\section{Foncteur de décatégorification et cardinalité des sections publiées}

Deux tours peuvent satisfaire
\[
  |X_m|=|Y_m|\qquad\text{pour tout }m
\]
et ne pas être isomorphes : il suffit que leurs applications de restriction aient des schémas de
ramification distincts. Le profil cardinal est donc un invariant
utile mais non fidèle.

\begin{remark}[Interprétation structurelle]
Compter combien de cases existent ne dit pas quelle case en prolonge quelle autre, quelle
orientation s'inverse, quelle retenue est transportée ou quel état revient. Le
cardinal mesure l'inventaire ; la généalogie mesure l'organisation.
\end{remark}

\begin{remark}[Séparation des résultats et portée]
Le codage ZFC, la section comme limite et la distinction entre égalité
observationnelle et identité d'état sont des
\emph{résultats établis dans les sections précédentes}. La catégorie unifiée, le topos de
préfaisceaux et l'ordinal fibré sont des
\emph{constructions explicites du présent développement}. Il n'est pas proclamé de théorie autonome complète
des classes ou des ordinaux transfinis HMT.
\end{remark}


\section{Limite profinie du système de troncatures généalogiques}

Soit
\[
  X=\varprojlim_mX_m
\]
une limite inverse d'ensembles finis discrets avec des applications surjectives.
Nous munissons \(X\) de la topologie initiale des projections
\(\pi_m:X\to X_m\). Les cylindres
\[
  [A]_m=\pi_m^{-1}(A),
  \qquad A\subseteq X_m,
\]
sont ouverts-fermés et forment une base.

\begin{theorem}[Ensemble des parties cylindriques]
\label{p3-83-c45-s2:thm:clopen-power}
\[
  \Clop(X)
  \cong
  \varinjlim_m\bigl(\mathcal P(X_m),p_{n,m}^{-1}\bigr).
\]
\end{theorem}

\begin{proof}
Tout cylindre est ouvert-fermé. Réciproquement, si \(C\subset X\) est ouvert-fermé, pour
chaque \(x\in C\) il existe un cylindre contenu dans \(C\). La compacité de \(C\)
donne un sous-recouvrement fini. En relevant ces cylindres à une profondeur
commune \(m\), on obtient \(C=\pi_m^{-1}(A)\) pour un certain
\(A\subseteq X_m\). Deux représentants à des profondeurs distinctes définissent le
même ouvert-fermé exactement lorsqu'ils coïncident après un relèvement commun.
\end{proof}

\begin{theorem}[Puissance sous forme close]
\label{p3-83-c45-s2:thm:closed-power}
\[
  \Closed(X)
  \cong
  \varprojlim_m\bigl(\mathcal P(X_m),p_{n,m*}\bigr),
\]
où \(p_*(A)=p(A)\).
\end{theorem}

\begin{proof}
Si \(F\subset X\) est fermé, il est compact. Ses projections
\(F_m=\pi_m(F)\) satisfont \(p_{n,m}(F_n)=F_m\). Réciproquement, une famille
compatible \((F_m)\) définit la limite inverse
\[
  F=\{x\in X:\pi_m(x)\in F_m\text{ pour tout }m\}.
\]
C'est une intersection de fermés et, par la compatibilité surjective des
restrictions \(F_n\to F_m\), ses projections sont exactement les \(F_m\).
Les deux constructions sont inverses.
\end{proof}

\section{Obstruction à la reconstruction de l'ensemble des parties généalogiques à partir de réductions finies}

Dans un espace de Cantor,
\[
 |\Clop(X)|=\aleph_0,
 \qquad
 |\Closed(X)|=2^{\aleph_0},
 \qquad
 |\mathcal P(X)|=2^{2^{\aleph_0}}.
\]
Les trois objets ne peuvent pas être identifiés.

Cette séparation ne présuppose pas l'hypothèse du continu. En écrivant
\(\mathfrak c=2^{\aleph_0}\),
\[
 |\Clop(X)|=\aleph_0,\qquad
 |\Closed(X)|=\mathfrak c,\qquad
 |\mathcal P(X)|=2^{\mathfrak c}.
\]
\(\mathsf{CH}\) identifie seulement \(\mathfrak c\) à \(\aleph_1\) ; elle ne
confond aucune des trois couches. L'indépendance de Gödel--Cohen change quel
cardinal peut occuper \(\mathfrak c\), et non l'inégalité de Cantor
\(\mathfrak c<2^{\mathfrak c}\).

\begin{proposition}[Falsificateur dense]
Si \(D\subsetneq X\) est dense, alors
\[
  \pi_m(D)=X_m
  \qquad\text{pour tout }m,
\]
de sorte que \(D\) et \(X\) possèdent les mêmes ombres finies.
\end{proposition}

\begin{proof}
Chaque fibre cylindrique non vide \(\pi_m^{-1}(a)\) est ouverte. Par densité,
elle rencontre \(D\). Par conséquent, \(a\in\pi_m(D)\) pour tout \(a\in X_m\).
\end{proof}

Les ombres finies reconstruisent l'adhérence d'un sous-ensemble, et non le
sous-ensemble arbitraire. C'est le point exact où une identification
\[
  \varprojlim_m\mathcal P(X_m)
  =
  \mathcal P(\varprojlim_mX_m)
\]
serait fausse.

\begin{figure}[H]
\centering
\begin{tikzpicture}[>=Latex,node distance=10mm and 17mm,
  box/.style={rounded corners,draw=hmtblue,fill=hmtlightblue,
    minimum width=38mm,minimum height=15mm,align=center},
  full/.style={rounded corners,draw=hmtviolet,fill=white,
    minimum width=40mm,minimum height=15mm,align=center}]
  \node[box] (fin) {\(\mathcal P(X_m)\)\\\scriptsize prédicats finis};
  \node[box,right=of fin] (clop) {\(\Clop(X)\)\\\scriptsize limite directe};
  \node[box,below=of clop] (closed) {\(\Closed(X)\)\\\scriptsize limite inverse};
  \node[full,right=of clop] (power) {\(\mathcal P(X)\)\\\scriptsize ensemble complet des parties};
  \draw[->,thick,hmtblue] (fin)--node[above]{\(p^{-1}\)}(clop);
  \draw[->,thick,hmtblue] (fin) to[bend right=24]
    node[below left]{\(p_*\)} (closed);
  \draw[->,thick,hmtgreen] (clop)--node[above]{\(\subset\)}(power);
  \draw[->,thick,hmtgreen] (closed)--node[below]{\(\subset\)}(power);
  \draw[dashed,very thick,hmtred] (closed.south east)
    .. controls +(8mm,-8mm) and +(-8mm,-8mm) ..
    node[midway,below=2mm]{\(\ne\)} (power.south west);
  \node[below=18mm of closed,text width=.82\textwidth,align=center] {
    \small Un sous-ensemble dense propre et \(X\) ont les mêmes projections finies :
    l'histoire finie retrouve l'adhérence, et non l'ensemble complet des parties.};
\end{tikzpicture}
\caption{Trois notions d'ensemble des parties qu'il ne faut pas aplatir.}
\label{p3-83-c45-s2:fig:three-powers}
\end{figure}

\section{\(4\) exigences de choix dans le système projectif}

Soit \(p:X\twoheadrightarrow Y\). Il convient de séparer :

\begin{enumerate}
\item une \emph{section ensembliste} \(s:Y\to X\) avec \(ps=\id_Y\) ;
\item une \emph{famille naturelle} de sections qui commute avec les
restrictions d'une tour ;
\item une \emph{section algébrique} qui préserve en outre les opérations ;
\item une \emph{section effective} calculée par un algorithme total.
\end{enumerate}

Toute section naturelle, algébrique ou effective est, en oubliant la structure, une
section ensembliste. Les trois propriétés supplémentaires sont, en général,
incomparables : un algorithme peut choisir des représentants sans être un homomorphisme ;
un homomorphisme peut ne pas être calculable dans la présentation choisie ; une
famille naturelle peut ne pas conserver l'opération algébrique.

\begin{remark}[Conditions de validité]
L'axiome du choix garantit une sélection ensembliste dans son
domaine. Il ne garantit ni naturalité, ni algébricité, ni effectivité, ni coût
uniforme. L'indépendance de \(\mathsf{AC}\) par rapport à \(\mathsf{ZF}\)
ne transforme pas non plus ces quatre exigences en une seule : même dans un modèle
avec choix, la section structurelle nécessite sa propre preuve.
\end{remark}

\Needspace{36\baselineskip}
\section{Obstruction de Pell à la périodicité uniforme du système de préfixes}

Le générateur actif construit
\[
  K_\infty\cong C_4\times\mathbb Z_3
\]
et des projections finies \(p_m:K_\infty\to K_m\). Il est possible de choisir un
représentant de chaque classe, mais pas une section homomorphe compatible avec le
groupe.

\begin{proposition}
Il n'existe pas d'homomorphisme \(s_m:K_m\to K_\infty\) avec
\(p_ms_m=\id\) pour le quotient qui introduit une torsion ternaire finie dans
\(K_m\).
\end{proposition}

\begin{proof}
Une section homomorphe enverrait un élément d'ordre \(3^r\) de \(K_m\) sur un
élément du même ordre dans la composante \(\mathbb Z_3\). Mais le groupe
additif \(\mathbb Z_3\) est sans torsion : si \(3^rx=0\), alors \(x=0\).
L'image ne peut pas se projeter de nouveau sur un élément non nul d'ordre
\(3^r\).
\end{proof}

L'obstruction ne dit pas « on ne peut pas choisir ». Elle dit : on ne peut pas choisir
en préservant simultanément la projection et la loi de groupe. L'information de la
retenue qui disparaît dans le quotient réapparaît comme obstruction à la section.

\begin{remark}[Séparation des résultats et portée]
Le continu profini et l'obstruction de Pell sont des
\emph{résultats établis dans les sections précédentes}. Les deux théorèmes d'ensemble des parties et le tableau des
choix sont des \emph{constructions explicites du présent développement}. Le falsificateur dense empêche
de transformer la recherche incomplète d'un sous-ensemble en un ensemble des parties déjà
reconstruit.
\end{remark}

%
\endgroup

\HMTFlushCurrentChapter
\chapter{Obstruction de Gödel--Turing à la décision uniforme des multiplicités projectives}%
\begingroup
\renewcommand{\chapter}[2][]{}%
\chapter{Obstruction de Gödel--Turing à la décision uniforme des multiplicités projectives}
\label{chap:p3-final:54:godel_turing}
\label{chap:p3-83-46-godel_turing}

\section{Décidabilité effective de chaque fenêtre finie de multiplicités projectives}

\begin{proposition}
Une profondeur \(N\) étant fixée, si les domaines sont finis et les observables
atomiques sont calculables, toute formule du langage fini
\(L_N\) est sémantiquement décidable.
\end{proposition}

\begin{proof}
Les formules atomiques se calculent ; les connecteurs se décident par des tables de
vérité ; chaque quantificateur parcourt un domaine fini. On conclut par
induction sur la complexité syntaxique.
\end{proof}

La décidabilité de chaque fenêtre ne fournit pas nécessairement un seul
algorithme qui décide une propriété de la tour complète.

\section{Tour effective associée au calcul de multiplicités projectives}

Soit \(T_e\) la machine de Turing d'indice \(e\). Pour chaque \(n\), nous définissons
\[
 X_n^e
 =
 \{b_n\}
 \cup
 \{c_n:T_e\text{ ne s'est pas arrêté dans ses premiers }n\text{ pas}\}.
\]
Les liens conservent les étiquettes existantes :
\[
  p_{n+1,n}(b_{n+1})=b_n,
  \qquad
  p_{n+1,n}(c_{n+1})=c_n.
\]
La tour visible possède un seul symbole \(v_n\), et
\[
  q_n(b_n)=q_n(c_n)=v_n.
\]

\begin{theorem}[Multiplicité globale et arrêt]
\label{p3-83-c46-s1:thm:halt-fibre}
\[
 \#q_\infty^{-1}(\boldsymbol v)
 =
 \begin{cases}
 1,&T_e\text{ s'arrête},\\
 2,&T_e\text{ ne s'arrête jamais}.
 \end{cases}
\]
Par conséquent, il n'existe pas d'algorithme total qui décide l'unicité de la fibre
globale pour toutes les tours effectives de cette famille.
\end{theorem}

\begin{proof}
La branche \(\boldsymbol b=(b_n)\) existe toujours. Si \(T_e\) s'arrête au
pas \(h\), les états \(c_n\) disparaissent à partir de \(h\) et ne forment pas une
section globale. Si elle ne s'arrête jamais, \(\boldsymbol c=(c_n)\) est une deuxième
section. Un décideur de l'unicité résoudrait \(\HALT\), en contradiction avec
\citep{turing1936}.
\end{proof}

On définit
\[
  U_e:\#q_\infty^{-1}(\boldsymbol v)=1.
\]
Alors \(U_e\iff T_e\downarrow\).

\begin{theorem}[Incomplétude uniforme]
\label{p3-83-c46-s1:thm:uniform-incompleteness}
Soit \(\mathsf T\) une théorie récursivement axiomatisable qui formalise les
tours \(X^e\) et est sémantiquement correcte pour \(U_e\) et
\(\neg U_e\). Alors \(\mathsf T\) ne décide pas tous les \(U_e\).
\end{theorem}

\begin{proof}
Si elle décidait chaque \(U_e\), pour un indice \(e\) on énumérerait en parallèle
les preuves de \(U_e\) et \(\neg U_e\) jusqu'à en trouver une. Par correction,
la première apparaîtrait exactement lorsque la machine s'arrête et la seconde
lorsqu'elle ne s'arrête pas. On obtiendrait un décideur de \(\HALT\).
\end{proof}

\begin{figure}[H]
\centering
\begin{tikzpicture}[>=Latex,node distance=11mm and 15mm,
  box/.style={rounded corners,draw=hmtblue,fill=hmtlightblue,
    minimum width=39mm,minimum height=13mm,align=center},
  warn/.style={rounded corners,draw=hmtred,fill=hmtlightred,
    minimum width=39mm,minimum height=13mm,align=center}]
  \node[box] (sim) {simulation finie de \(T_e\)};
  \node[box,below=of sim] (levels) {\(X_n^e\) calculable\\pour chaque \(n\)};
  \node[box,below=of levels] (fibre)
    {\(\#q_\infty^{-1}(v)=1\iff T_e\downarrow\)};
  \node[warn,below left=of fibre] (undec) {sans décideur uniforme};
  \node[warn,below right=of fibre] (inc)
    {théorie effective correcte\\incomplète pour les \(U_e\)};
  \draw[->,thick,hmtblue] (sim)--(levels);
  \draw[->,thick,hmtblue] (levels)--(fibre);
  \draw[->,thick,hmtred] (fibre)--(undec);
  \draw[->,thick,hmtred] (fibre)--(inc);
\end{tikzpicture}
\caption{De la décision finie à l'incomplétude uniforme.}
\label{p3-83-c46-s1:fig:godel-turing}
\end{figure}

\section{Complétude finie des niveaux HMT et domaine exact de l'obstruction de Gödel--Turing}

Le nom académique utilisé est :
\[
\boxed{\text{Théorème HMT d'incomplétude uniforme pour les multiplicités projectives}.}
\]
On emploie également la dénomination abrégée « troisième résultat de Gödel dans
HMT ». La preuve
est de Gödel--Turing : elle ne reproduit pas la diagonale autoréférente de
\cite{goedel1931}, mais réduit HALT à la multiplicité d'une fibre.

\begin{remark}[Conditions de validité]
Le théorème ne prouve pas l'incohérence de ZFC. Il n'empêche pas non plus de démontrer
directement qu'une tour particulière possède une section unique. Il empêche une
procédure universelle d'unicité pour toutes les tours effectives.
\end{remark}

% REMATE_LOCAL_GODEL_CONCLUSION
\Needspace{25\baselineskip}
\section{Relation exacte avec les \(9\) phases du Topological Phase Kernel}

Les \(9\) phases du TPK apportent un algorithme concret qui, avec ses lecteurs déclarés,
partitionne le bloc de raffinement en \(729\) sous-cylindres et sélectionne l'un d'eux
à chaque étape finie.
Le passage à la section infinie utilise un théorème global de non-vacuité,
de compatibilité et d'unitarité.

Le résultat de Gödel--Turing dit :
\[
\text{il n'existe pas de décideur d'unicité pour toute tour effective}.
\]
Il n'affirme pas :
\[
\text{aucune tour concrète ne peut avoir de preuve uniforme propre}.
\]
Et une preuve particulière ne fournit pas le décideur universel interdit.


\begin{remark}[Conclusion structurelle]
Décision finie, existence de la limite, unicité d'une instance et décision
uniforme de toutes les instances sont quatre niveaux. La clôture holographique
exige de les déclarer au lieu de transférer automatiquement une propriété d'un
niveau à un autre.
\end{remark}

\begin{remark}[Séparation des résultats et portée]
La tour réductrice, l'équivalence \(U_e\iff T_e\downarrow\) et le théorème
métamathématique sont des \emph{résultats établis dans les sections précédentes}. La position centrale au sein
de la clôture de l'infini est une organisation de l'exposé qui ne modifie ni le contenu logique ni le domaine du théorème.
\end{remark}

%
\endgroup

\HMTFlushCurrentChapter
\chapter{Gödel--Cohen : cohérence relative, indépendance et forcing sur les espaces de préfixes}%
\begingroup
\renewcommand{\chapter}[2][]{}%
\chapter{Gödel--Cohen : consistance relative, indépendance et forcing sur les espaces de préfixes}
\label{chap:p3-final:55:godel_cohen}
\label{chap:p3-83-47-godel_cohen}

\section{De la cohérence constructible à l'indépendance par forcing}

\begin{theorem}[Gödel, 1940 : moitié constructible]
Dans la métathéorie usuelle,
\[
  \operatorname{Con}(\mathsf{ZF})
  \Longrightarrow
  \operatorname{Con}
  \bigl(\mathsf{ZF}+\mathsf{AC}+\mathsf{GCH}\bigr).
\]
Plus précisément, étant donné un modèle \(M\models\mathsf{ZF}\), son univers
constructible interne \(L^M\) satisfait, au sens interne correspondant,
\[
  \mathsf{ZF}+V=L+\mathsf{AC}+\mathsf{GCH}.
\]
\end{theorem}

\begin{proof}[Idée de la preuve classique]
La hiérarchie constructible est définie par récurrence sur les ordinaux,
en n'admettant à chaque étape que les sous-ensembles définissables sur les étapes
précédentes. La construction conserve les ordinaux, produit un modèle interne
et fournit un bon ordre définissable global, dont se déduisent le choix et
l'hypothèse généralisée du continu \citep{goedel1940}.
\end{proof}

\begin{theorem}[Cohen, 1963--1964 : moitié par forcing]
On a les implications métamathématiques
\[
  \operatorname{Con}(\mathsf{ZFC})
  \Longrightarrow
  \operatorname{Con}(\mathsf{ZFC}+\neg\mathsf{CH}),
\]
et
\[
  \operatorname{Con}(\mathsf{ZF})
  \Longrightarrow
  \operatorname{Con}(\mathsf{ZF}+\neg\mathsf{AC}).
\]
En combinant les deux moitiés, \(\mathsf{CH}\) est indépendante de
\(\mathsf{ZFC}\) et \(\mathsf{AC}\) est indépendant de \(\mathsf{ZF}\),
toujours sous l'hypothèse de cohérence correspondante
\citep{cohen1963,cohen1964}.
\end{theorem}

Pour la première construction, on part d'un modèle de base \(M\), d'un poset
\(\mathbb P\in M\) et d'un filtre \(M\)-générique \(G\subseteq\mathbb P\) :
\[
  G\cap D\ne\varnothing
  \qquad
  \text{pour tout dense }D\subseteq\mathbb P\text{ appartenant à }M.
\]
Le théorème du forcing relie \(p\Vdash_M\varphi\) à la vérité de
\(\varphi\) dans \(M[G]\). Pour obtenir
\(\mathsf{ZF}+\neg\mathsf{AC}\), on prend en outre un sous-modèle symétrique
approprié d'une extension générique ; une extension ordinaire d'un modèle de
\(\mathsf{ZFC}\) conserve le choix \citep{cohen1966}.

\begin{remark}[Conditions de validité]
Ces théorèmes prouvent la cohérence relative et l'indépendance. Ils ne prouvent pas la
cohérence absolue de \(\mathsf{ZF}\) ou \(\mathsf{ZFC}\), ne décident pas
\(\mathsf{CH}\) dans un univers privilégié et ne transforment pas tout filtre
compatible en filtre générique. Ils ne sélectionnent pas non plus les branches HMT ni
n'identifient une sortie numérique.
\end{remark}

\subsection{\(2\) opérations non équivalentes : extension du modèle et enrichissement de l'objet}

Le forcing et HMT n'effectuent pas la même opération. Le premier compare un
univers \(M\) à une extension \(M[G]\) ; la seconde conserve des fibres qu'un
lecteur oublie. Si \(\mathsf{Hist}_{\rm HMT}\) est la catégorie des histoires
enrichies, le foncteur
\[
  U:\mathsf{Hist}_{\rm HMT}\longrightarrow\mathsf{Set},
  \qquad
  (x_m,\mu_m)_{m\ge0}\longmapsto(x_m)_{m\ge0},
\]
peut identifier des histoires distinctes sans changer l'ensemble visible. La
question de Cohen est de savoir quels sous-ensembles existent dans le modèle ; la question
généalogique est de savoir combien de relèvements avec mémoire et opérateurs possède un objet
visible. Ce sont des axes compatibles, mais aucun ne remplace l'autre.

\begin{center}
\begin{tabularx}{.94\textwidth}{>{\bfseries}p{.22\textwidth}XX}
\toprule
Opération & Ce qui change & Ce qu'elle ne prouve pas à elle seule\\
\midrule
\(M\mapsto M[G]\) & Univers d'ensembles, de noms et de vérité forcée &
Unicité ou généalogie d'une section HMT.\\
\(\xi\mapsto U(\xi)\) & Quantité de mémoire visible dans la projection &
\(\mathsf{CH}\), \(\mathsf{AC}\) ou généricité sur \(M\).\\
\bottomrule
\end{tabularx}
\end{center}

\subsection{L'ensemble des parties HMT ne présuppose pas \texorpdfstring{\(\mathsf{CH}\)}{CH}}

Dans chaque modèle où \(X\) est un espace de Cantor,
\[
  |\Clop(X)|=\aleph_0,\qquad
  |\Closed(X)|=\mathfrak c,\qquad
  |\mathcal P(X)|=2^{\mathfrak c},
  \qquad \mathfrak c=2^{\aleph_0}.
\]
\(\mathsf{CH}\) affirme uniquement \(\mathfrak c=\aleph_1\). Par le théorème
de Cantor, \(\mathfrak c<2^{\mathfrak c}\) dans chacun des modèles
comparés. Lors du passage de \(M\) à \(M[G]\), de nouvelles branches, de nouveaux
codes fermés et de nouveaux sous-ensembles peuvent apparaître, bien que les cylindres finis
conservent leur règle de formation. L'ensemble complet des parties est sensible au modèle ; la
distinction ouvert-fermé--fermé--complet ne dépend pas de la décision de \(\mathsf{CH}\).

\section{Poset de préfixes compatibles et limite inverse des extensions}

On définit
\[
  \mathbb P_{\TPK}=\bigsqcup_{n\ge0}\mathcal S_n
\]
et nous ordonnons par prolongement : \(q\le p\) si \(q\) contient plus d'information
que \(p\). Pour chaque \(N\),
\[
  D_N=\{p\in\mathbb P_{\TPK}:\ell(p)\ge N\}.
\]
Si tout préfixe survivant peut être prolongé, chaque \(D_N\) est dense.

\begin{theorem}[Branches et filtres de profondeur]
Les sections de \(\OmegaInf\) sont en bijection avec les filtres compatibles
de \(\mathbb P_{\TPK}\) qui rencontrent tous les \(D_N\).
\end{theorem}

\begin{proof}
Les préfixes d'une section sont compatibles, dirigés et atteignent toute
profondeur. Réciproquement, un filtre qui rencontre chaque \(D_N\) détermine une
restriction unique à chaque niveau ; la propriété de direction exclut les conditions
incompatibles et produit une section.
\end{proof}

Rencontrer tous les denses \(D_N\) exprime une profondeur illimitée, et non
une généricité sur un modèle. Être \(M\)-générique exigerait de rencontrer tout
dense \(D\in M\).

\begin{proposition}[Cas de Cohen conditionnel]
Si \(\OmegaInf\) est compact, métrisable, de dimension zéro, parfait et sans
points isolés, il est homéomorphe à l'ensemble de Cantor ; le forcing des cylindres non vides
a une partie dense équivalente au forcing de Cohen
\citep{cohen1963,cohen1964}.
\end{proposition}

Une section unique produit, au contraire, un forcing atomique ou trivial.
Le mot « forcing » ne sélectionne pas à lui seul une branche.

Pour une section \(x=(x_m)_{m\ge0}\), ses préfixes forment un filtre \(G_x\)
qui traverse les denses de profondeur. Cependant,
\[
  \bigl(\forall N,\;G_x\cap D_N\ne\varnothing\bigr)
  \centernot\Longrightarrow
  \bigl(\forall D\in M\text{ dense},\;G_x\cap D\ne\varnothing\bigr).
\]
La première propriété exprime un prolongement illimité ; la seconde,
une généricité relative à un modèle. Une section compatible produite par
neuf phases du TPK ne devient pas, par nomenclature, un réel de Cohen.

\begin{remark}[Conditions de validité]
Réfutent une identification automatique avec Cohen : un cylindre terminal ou un
point isolé ; un filtre qui traverse tous les \(D_N\) mais omet un
dense de \(M\) ; deux mémoires compatibles sur la même branche visible ; ou
l'absence d'un plongement dense explicite. Le nom partagé « branche » n'est pas
un morphisme.
\end{remark}

\section{Correspondance entre sections généalogiques et classes de Borel de l'espace limite}

Les cylindres engendrent
\[
 \Gamma_{\rm HMT}^{\rm Bor}
 =
 \sigma\bigl(\{[p]_n:p\in\mathcal S_n\}\bigr).
\]
La classe est stable par complémentation, unions dénombrables et images réciproques
d'applications continues certifiées.

Si un jeu infini sur des histoires possède une condition de victoire borélienne,
le théorème de détermination borélienne de Martin garantit une stratégie
gagnante \citep{martin1975}. C'est une application classique à l'espace HMT ;
elle ne dérive ni AD, ni MM, ni une valeur de \(2^{\aleph_0}\), ni la catégoricité de
\(\mathcal P(\mathbb R)\).

Pour \(A,B\subseteq\OmegaInf\), la réduction de Wadge
\[
 A\le_WB
\]
exige une application continue \(f\) avec \(x\in A\iff f(x)\in B\). Si l'on
restreint les applications à des opérateurs TPK certifiés, il apparaît un préordre plus
étroit dont la bonne fondation ou la semi-linéarité nécessite une preuve propre.

\begin{remark}[Conditions de validité]
Une démonstration qui utilise la bonne fondation de Wadge pour prouver la
détermination qui justifierait ensuite cette bonne fondation est circulaire. Les
branches historiques AD/MM sont conservées comme architecture, et non comme théorèmes HMT.
\end{remark}

\section{Optimisation minimax sur des fenêtres finies de préfixes compatibles}

Soit \(P_N=\mathcal S_N\) l'ensemble fini des candidats, \(T_N\) un ensemble
fini de tests adversariaux et
\[
  A_N:P_N\times T_N\to[-1,1]
\]
un gain fixé avant le résultat.

\begin{theorem}[Minimax fini]
\[
 \max_{\sigma\in\Prob(P_N)}
 \min_{\tau\in\Prob(T_N)}
 \mathbb E_{\sigma,\tau}A_N
 =
 \min_{\tau\in\Prob(T_N)}
 \max_{\sigma\in\Prob(P_N)}
 \mathbb E_{\sigma,\tau}A_N.
\]
\end{theorem}

\begin{proof}
C'est la dualité forte entre le programme linéaire du sélecteur et celui de
l'adversaire, sous la forme du théorème minimax de von Neumann
\citep{vonneumann1928}.
\end{proof}

L'équilibre peut être mixte et ne fournit pas une branche pure unique.

\section{Système projectif de troncatures et construction de la section limite}

Soient
\[
 X=\varprojlim P_N,
 \qquad
 Y=\varprojlim T_N
\]
compacts, et soit \(A:X\times Y\to\mathbb R\) continu. Les espaces de
probabilités sont compacts convexes dans la topologie faible. Le théorème de
Sion \citep{sion1958} donne
\[
 \sup_{\mu\in\Prob(X)}\inf_{\nu\in\Prob(Y)}
 \iint A\,d\mu\,d\nu
 =
 \inf_{\nu\in\Prob(Y)}\sup_{\mu\in\Prob(X)}
 \iint A\,d\mu\,d\nu.
\]

\begin{corollary}
Un équilibre global ne sélectionne une branche pure que si la mesure du sélecteur
est une Dirac \(\delta_x\). L'existence d'un équilibre mixte n'implique pas
la concentration.
\end{corollary}

\begin{remark}[Conditions de validité]
Le passage à la limite échoue sans compacité, continuité ou cohérence
projective des mesures. Si le gain est défini en utilisant la cible qu'il doit
sélectionner, le minimax est alimenté par la cible et n'explique pas la route.
\end{remark}

\section{Extensions projectives et matricielles apportées par les entrelaceurs opératoriels}

Le forcing organise les prolongements possibles ; Borel classe les ensembles
d'histoires ; le minimax construit la robustesse face aux falsificateurs. Aucun ne remplace
la loi structurelle qui produit un survivant unique. Leur fonction est
de clarifier la topologie et la décision de l'espace des routes, et non d'introduire une
sélection rétrospective.

\begin{remark}[Séparation des résultats et portée]
Les théorèmes de Gödel et Cohen sont des \emph{théorèmes classiques appliqués}. Le
poset, la comparaison typée entre filtre de profondeur et filtre générique,
la classe de Borel HMT et le jeu sélecteur--falsificateur sont des
\emph{constructions explicites du présent développement}. Martin, von Neumann et Sion s'appliquent
ensuite, avec leurs hypothèses classiques. Les anciennes affirmations PSC/HTT sur
CH, AD et MM ne sont pas promues.
\end{remark}

%
\endgroup

\HMTFlushCurrentChapter
\chapter{Système opératoire d'Euler en HMT : défaut somme--produit, géométries exponentielles triadiques, Principe d'Ambivalence et double projection}
% Cuerpo editor-ready del nuevo capitulo 57.

\section{Défaut somme--produit d'APP et coefficient linéaire du produit modulaire d'Euler}
\label{sec:l9s102:euler-defecto}

\subsection{Identité exacte \texorpdfstring{\(459-405=54=[q]t_3(q)\)}{459-405=54=[q]t3(q)}}

Sur le support de \(81\) cellules, les deux feuillets d'APP produisent
\begin{equation}
 S_+=405,
 \qquad
 S_\times=459,
 \qquad
 S_\times-S_+=54.
 \label{eq:l9s102:euler-405-459}
\end{equation}
La régularisation triadique fournit l'exposant \(-1/12\). Sa
publication dodécaphasique détermine le pôle \(q^{-1}\), et le quotient de
Dedekind de niveau \(3\) s'écrit
\begin{equation}
 t_3(q)
 =\left(\frac{\eta(\tau)}{\eta(3\tau)}\right)^{12}
 =q^{-1}\prod_{n\ge1}(1+q^n+q^{2n})^{-12},
 \qquad q=e^{2\pi i\tau}.
 \label{eq:l9s102:euler-t3}
\end{equation}

\begin{theorem}[Identité somme--produit d'Euler--HMT]
\label{thm:l9s102:euler-54}
Le défaut entre les deux évaluations d'APP coïncide avec le coefficient
linéaire de \(t_3\) :
\begin{equation}
 \boxed{
 S_\times-S_+=54=[q]\,t_3(q).}
 \label{eq:l9s102:euler-54}
\end{equation}
\end{theorem}

\begin{proof}
Chaque facteur admet
\((1+q^n+q^{2n})^{-12}=1-12q^n+O(q^{2n})\). Pour le coefficient de
\(q\) postérieur au pôle \(q^{-1}\), les partitions de degré
\(2\) du produit de niveau \(3\) contribuent. Le développement fini certifié produit
\([q]t_3=54\). APP obtient le même entier par la soustraction exacte
\(459-405\).
\end{proof}

\subsection{Conservation du défaut au moyen du résidu, du quotient et de la retenue nonadique}

L'identité \cref{eq:l9s102:euler-54} relie la différence additive de deux
feuillets finis à un produit modulaire infini. Le lien conserve
l'origine de \(54\) : somme, produit et régularisation demeurent trois
coordonnées distinctes de l'état.

\subsubsection{Complétion EMRH et conservation de l'unité}
\label{sec:l9s102:euler-emrh}

La fibre ternaire de profondeur \(6\) possède \(3^6=729\) éléments. Sa
complétion décimale distingue
\begin{equation}
 729+270=999,
 \qquad
 729+271=1000,
 \qquad
 271=243+27+1.
 \label{eq:l9s102:euler-729-1000}
\end{equation}
Le terme terminal \(+1\) est l'unité de fermeture transmise par la
retenue. L'égalité binomiale
\begin{equation}
 1000=(9+1)^3=9^3+3\cdot9^2+3\cdot9+1
 \label{eq:l9s102:euler-binomial-1000}
\end{equation}
exprime le transport entre la carte nonadique et la publication décimale sans
perdre la généalogie de la complétion.

L'évaluation archimédienne est réalisée sur des cylindres semi-fermés
compatibles. Les deux développements d'un rationnel terminal représentent la
même valeur et conservent des histoires de bord distinctes. Cette structure
type l'identité \(0.999\ldots=1\) comme égalité d'évaluation accompagnée
d'une mémoire de voie.

\section{Algèbres quadratiques du TRIT et loi exponentielle elliptique, parabolique et hyperbolique}
\label{sec:l9s102:euler-algebras-trit}

\subsection{Famille opératorielle \texorpdfstring{\(J_{+1}^2=-I\), \(J_0^2=0\) et \(J_{-1}^2=I\)}{J+1^2=-I, J0^2=0 et J-1^2=I}}

Para \(\tau\in\{+1,0,-1\}\), definimos
\begin{equation}
 \mathbb A_\tau
 =\mathbb R[J_\tau]/(J_\tau^2+\tau I).
 \label{eq:l9s102:euler-algebras}
\end{equation}
La série exponentielle converge dans chaque algèbre et se réduit à
\begin{align}
 e^{tJ_{+1}}&=\cos t\,I+\sin t\,J_{+1},
 \label{eq:l9s102:euler-eliptica}\\
 e^{tJ_0}&=I+tJ_0,
 \label{eq:l9s102:euler-parabolica}\\
 e^{tJ_{-1}}&=\cosh t\,I+\sinh t\,J_{-1}.
 \label{eq:l9s102:euler-hiperbolica}
\end{align}
TRIT assigne respectivement fermeture elliptique, variation du premier ordre et
ouverture hyperbolique. L'identité classique occupe le feuillet \(\tau=+1\) :
\begin{equation}
 \boxed{\operatorname{Exp}(\pi J_{+1})+I=0.}
 \label{eq:l9s102:euler-clasica}
\end{equation}
Le feuillet parabolique conserve le premier jet, et le feuillet hyperbolique publie
un couple de valeurs propres réciproques.

\subsection{Identités \texorpdfstring{\(\operatorname{Exp}(\pi J_{+1})+I=0\)}{Exp(pi J+1)+I=0} et \texorpdfstring{\(\operatorname{Exp}(2\log\varphi\,H_\varphi)=F_{\mathrm{av}}^2\)}{Exp(2 log phi H_phi)=F_av^2}}

Soit
\begin{equation}
 F_{\mathrm{av}}=\begin{pmatrix}0&1\\1&1\end{pmatrix},
 \qquad
 H_\varphi=\frac{2F_{\mathrm{av}}^2-3I}{2\varphi-1}.
 \label{eq:l9s102:euler-hphi}
\end{equation}
Alors \(H_\varphi^2=I\) et
\begin{equation}
 \boxed{
 \operatorname{Exp}(2\log\varphi\,H_\varphi)=F_{\mathrm{av}}^2.}
 \label{eq:l9s102:euler-fav}
\end{equation}
Dans cette coordination, \(\pi\) mesure la fermeture angulaire, \(\varphi\) la
valeur propre expansive de la mémoire et \(e\) l'opération qui transforme un
générateur additif en transport multiplicatif.

\section{Principe d'Ambivalence : involution entre \texorpdfstring{\(\pi/\varphi^2\)}{pi/phi^2} et \texorpdfstring{\(\varphi/\pi^2\)}{phi/pi^2} et accumulation nonadique d'orientation}
\label{sec:l9s102:euler-ambivalencia}

\subsection{Action \texorpdfstring{\(\mathcal Q(x,y)=(x/y^2,y/x^2)\)}{Q(x,y)=(x/y^2,y/x^2)} sur les lecteurs surfacique et volumétrique}

Nous définissons le lecteur quadratique
\begin{equation}
 \mathcal Q(x,y)
 =\left(\frac{x}{y^2},\frac{y}{x^2}\right).
 \label{eq:l9s102:euler-q}
\end{equation}
Dans les coordonnées
\begin{equation}
 P=xy,
 \qquad O=\frac{x}{y},
 \label{eq:l9s102:euler-po}
\end{equation}
son action est
\begin{equation}
 P\circ\mathcal Q=P^{-1},
 \qquad
 O\circ\mathcal Q=O^3.
 \label{eq:l9s102:euler-q-action}
\end{equation}
La seconde itération satisfait
\begin{equation}
 \boxed{
 P\circ\mathcal Q^2=P,
 \qquad
 O\circ\mathcal Q^2=O^9.}
 \label{eq:l9s102:euler-q2}
\end{equation}
Le produit revient et l'orientation conserve une puissance nonadique de
mémoire.

\subsection{Identité \texorpdfstring{\(\varphi^3(\pi/\varphi^2)^2(\varphi/\pi^2)=1\)}{phi^3(pi/phi^2)^2(phi/pi^2)=1} et bidirectionnalité typée}

Pour \((x,y)=(\pi,\varphi)\), les deux orientations sont
\begin{equation}
 \rho_A=\frac{\pi}{\varphi^2},
 \qquad
 \rho_E=\frac{\varphi}{\pi^2},
 \label{eq:l9s102:euler-rho-a-e}
\end{equation}
et obéissent
\begin{equation}
 \boxed{\varphi^3\rho_A^2\rho_E=1.}
 \label{eq:l9s102:euler-ambivalencia-identidad}
\end{equation}
L'involution \((x,y)\mapsto(y,x)\) échange les orientations surfacique et
électronique du même lecteur \(R(x,y)=x/y^2\).

\section{Transformation exponentielle du générateur parangulaire \texorpdfstring{\(A_{\mathrm{rad}}I+C^*_{\mathrm{rad}}H_\varphi\)}{A_rad I+C*_rad H_phi}}
\label{sec:l9s102:euler-accion-par-angular}

La carte d'action
\begin{equation}
 T_{-34}(u)=u\,10^{-34}\mathcal U_S
 \label{eq:l9s102:euler-t34}
\end{equation}
transporte l'identité angulaire
\(C^*_{\deg}=2(\eta_{\mathrm{ret}}+\alpha_{\mathrm{HMT}})\) vers
\begin{equation}
 \boxed{
 T_{-34}(C^*_{\deg})
 =2\hbar_{\mathrm{ret}}
 +2T_{-34}(\alpha_{\angle,\deg}).}
 \label{eq:l9s102:euler-cstar-hbar}
\end{equation}
Chaque terme appartient ainsi à la même droite dimensionnelle
\(\mathcal L_S\).

\subsection{Projecteurs spectraux \texorpdfstring{\(P_\pm=(I\pm H_\varphi)/2\)}{P+/-=(I+/-H_phi)/2} et valeurs propres conjuguées \texorpdfstring{\(q_\pm\)}{q+/-}}

Soit
\begin{equation}
 B=A_{\mathrm{rad}}I+C^*_{\mathrm{rad}}H_\varphi,
 \qquad
 P_\pm=\frac{I\pm H_\varphi}{2}.
 \label{eq:l9s102:euler-b}
\end{equation}
L'exponentiation transforme le générateur additif en deux feuillets
multiplicatifs conjugués :
\begin{equation}
 \boxed{
 e^{-B}=q_+P_++q_-P_-,
 \qquad
 q_\pm=e^{-A_{\mathrm{rad}}\mp C^*_{\mathrm{rad}}}.}
 \label{eq:l9s102:euler-exp-b}
\end{equation}

\subsection{Compatibilité avec la publication parangulaire et les canaux \texorpdfstring{\(90/120\)}{90/120}}

Les cardinaux \(90\) et \(120\) proviennent de l'incidence tritique du TPK.
Le couple \(q_\pm\) ne les engendre pas : il évalue spectralement l'opérateur
\(V_{90,120}\) sur ces canaux préalablement construits.

\section{Bord elliptique d'action et intérieur hyperbolique de Hilbert--Beltrami--Klein}
\label{sec:l9s102:euler-elipse-klein}

\subsection{Coordonnées conjuguées de l'ellipse d'action et loi \texorpdfstring{\(\operatorname{Area}(\theta)=\hbar_{\mathrm{ret}}\theta\)}{Area(theta)=hbar_ret theta}}

La rapidité angulaire
\begin{equation}
 \eta=\operatorname{artanh}\!\left(\frac{C^*}{A}\right)
 \label{eq:l9s102:euler-rapidez}
\end{equation}
définit les demi-axes symplectiques
\begin{equation}
 a_S=\sqrt{2\hbar_{\mathrm{ret}}e^\eta},
 \qquad
 b_S=\sqrt{2\hbar_{\mathrm{ret}}e^{-\eta}}.
 \label{eq:l9s102:euler-semiejes}
\end{equation}
Avec
\begin{equation}
 Q=a_S\cos\theta,
 \qquad
 P=b_S\sin\theta,
 \label{eq:l9s102:euler-qp}
\end{equation}
l'action contenue dans l'arc satisfait
\begin{equation}
 \boxed{\operatorname{Area}(\theta)=\hbar_{\mathrm{ret}}\theta.}
 \label{eq:l9s102:euler-area-accion}
\end{equation}

\subsection{Correspondance \texorpdfstring{\(\rho=\tan(\theta/2)=\tanh(u/2)\)}{rho=tan(theta/2)=tanh(u/2)} entre bord périodique et intérieur hyperbolique}

La coordonnée de Gudermann
\begin{equation}
 \rho=\tan\frac\theta2=\tanh\frac u2
 \label{eq:l9s102:euler-gudermann}
\end{equation}
coordonne les cartes circulaire et hyperbolique.
Le bord porte l'orbite symplectique ; le domaine convexe intérieur porte
la métrique hyperbolique de Hilbert--Beltrami--Klein. La coordonnée de
Gudermann réalise leur correspondance radiale et conserve la distinction des
métriques.

\section{Compatibilité de la fermeture opératorielle d'Euler avec les publications archimédienne et exceptionnelle}
\label{sec:l9s102:euler-doble-proyeccion}

\subsection{Double projection du même état enrichi et conservation de l'incidence}

Soit \(X^{\mathrm{enr}}\) l'état enrichi construit par
APP--TRIT--TPK. Ses deux publications principales sont
\begin{equation}
 X^{\mathrm{enr}}
 \xrightarrow{\ \mathcal P_{\mathrm{arq}}\ }\mathbb R,
 \qquad
 X^{\mathrm{enr}}
 \xrightarrow{\ \mathcal P_{\mathrm{exc}}\ }
 \mathrm{Witt}\to\mathrm{Golay}\to\mathrm{Leech}
 \to V^\natural\to\mathbb M\to J.
 \label{eq:l9s102:euler-doble-proyeccion}
\end{equation}
La première contracte des cylindres compatibles vers la valeur archimédienne ; la
seconde conserve l'incidence de bord jusqu'au corridor exceptionnel.
L'entier \(54\) apparaît dans APP et dans le coefficient modulaire de
\cref{eq:l9s102:euler-54} ; le corridor \(4\to2\to6\to54\) transporte
ensuite cette incidence vers le réseau de Leech, l'algèbre d'opérateurs
de vertex et Moonshine. L'action du Monstre réside dans
\(V^\natural\) ; la généalogie des chiffres et des voies demeure dans le domaine
du TPK et atteint le corridor au moyen des applications d'incidence déclarées.

\subsection{Système macrovectoriel \texorpdfstring{\(\mathfrak E_{\mathrm{HMT}}=0\)}{E_HMT=0} : confluence des identités eulériennes dans un unique opérateur}
\label{sec:l9s102:euler-macrovector}

Les identités précédentes sont réunies dans l'application typée
\begin{equation}
 \boxed{
 \mathfrak E_{\mathrm{HMT}}=
 \begin{pmatrix}
 (S_\times-S_+)-[q]t_3\\[1mm]
 \operatorname{Exp}(I)-eI\\
 \operatorname{Exp}(\pi J_{+1})+I\\
 \operatorname{Exp}(J_0)-I-J_0\\
 \operatorname{Exp}(2\log\varphi\,H_\varphi)-F_{\mathrm{av}}^2\\
 \varphi^3\rho_A^2\rho_E-1\\
 T_{-34}(C^*)-2\hbar_{\mathrm{ret}}
 -2T_{-34}(\alpha_{\angle,\deg})
 \end{pmatrix}=0.}
 \label{eq:l9s102:euler-macrovector}
\end{equation}
Chaque composante conserve son domaine et sa propre démonstration. L'application
agrégée met en évidence la coordination de la somme, du produit, de la fermeture elliptique, du premier
jet, de l'expansion hyperbolique, de l'ambivalence et de l'action.

La composition interne APP--TRIT--TPK transporte déjà somme, produit,
orientation, résidu, quotient, retenue, voie et mémoire jusqu'aux
publications précédentes. On ne postule donc pas d'opérateur d'entrelacement HMT
absent. Une promotion ultérieure ne peut être formulée qu'au moyen du carré
typé spécifique exigé par une représentation ultérieure ---par exemple,
une équivariance complète de groupe jusqu'à Witt---, sans amoindrir la confluence
opératorielle démontrée ici.













\HMTFlushCurrentChapter
